% Generated mechanically from frozen input p07/ko/etale-cohomology.tex.
% Do not edit this generated file; edit the bound locale source instead.

% OK, start here.
%

\setcounter{chapter}{58}
\chapter{에탈 코호몰로지}



\phantomsection
\label{etale-cohomology-section-phantom}



\section{소개} \label{etale-cohomology-section-introduction}

\noindent 이 장은 스킴의 에탈 코호몰로지에 관한 일련의 장 중 첫 번째 장이다. 이번 장에서는 이 위상에서 아벨 층의 에탈 위상과 코호몰로지의 아주 기본적인 내용을 논의한다. 논의된 많은 주제는 첫 독해에서 안전하게 건너뛸 수 있다. 건너뛸 항목을 결정하는 방법에 대해서는 다음 절의 조언을 참조하라.

\medskip\noindent 이 장의 초기 버전은 2009 가을에 Columbia University에서 Johan de Jong이 가르친 에탈 코호몰로지 과정의 첫 번째 부분의 노트로 구성되었다. 원래 메모 작성자는 Thibaut Pugin, Zachary Maddock 및 Min Lee였다. 과정의 두 번째 부분은 Trace Formula 장에서 찾을 수 있다. Trace Formula, 절 \ref{trace-section-introduction}를 참조하라.




\section{첫 독해에서 어떤 절을 건너뛸까요?} \label{etale-cohomology-section-skip}

\noindent 우리는 대수공간, Deligne-Mumford 스택, 대수 스택 등과 관련된 이론 개발을 위해 이 장의 자료를 사용하고 싶다. 따라서 스킴에 대한 에탈 코호몰로지의 원래 설명에 꽤 기술적인 자료를 추가했다. 독자는 ``토포스''라는 단어의 빈도나 집합 이론과 관련된 논의, 또는 사상의 스킴의 매우 일반적인 성질을 다루는 증명을 통해 이 자료를 인식할 수 있다. 이러한 토론 중 일부는 첫 독해에서 건너뛸 수 있다.

\medskip\noindent 특히, 독자가 다음 절을 건너뛰는 것이 좋다. \begin{enumerate} \item 크고 작은 비교 토포스, 절 \ref{etale-cohomology-section-compare}. \item 사상, 절 \ref{etale-cohomology-section-morphisms}를 복구하는 중이다. \item 밀고 당기세요, 절 \ref{etale-cohomology-section-monomorphisms}. \item 성질 (), 절 \ref{etale-cohomology-section-A}. \item 성질 (B), 절 \ref{etale-cohomology-section-B}. \item 성질 (C), 절 \ref{etale-cohomology-section-C}. \item 작은 에탈 사이트, 절 \ref{etale-cohomology-section-topological-invariance}의 위상학적 불변성. \item 정역 보편적 단사 사상, 절 \ref{etale-cohomology-section-integral-universally-injective}. \item 큰 사이트 및 직상, 절 \ref{etale-cohomology-section-big}. \item 완전성 크고 낮은 비명소리, 절 \ref{etale-cohomology-section-exactness-lower-shriek}. \end{enumerate} 이 절 외에도 독자가 위에 제공된 키워드로 인식할 수 있는 건너뛸 수 있는 산발적인 결과가 있다.



%9.08.09
\section{프롤로그} \label{etale-cohomology-section-prologue}

\noindent 이 강의는 또 다른 코호몰로지 이론에 관한 것이다. 비고의 첫 번째 점은 자리스키 위상이 완전히 만족스럽지 않다는 것이다. 우리가 원하는 결과를 제공하지 못하는 주된 이유 중 하나는 $X$가 복소다양체이고 $\mathcal{F}$가 상수층인 경우 $$
H^i(X, \mathcal{F}) = 0, \quad \text{모두를 위해} i > 0.
$$이기 때문이다. 그 이유는 다음과 같습니다. 기약 스킴(특히 다양체)에서는 비어 있지 않은 두 열린 부분집합이 만나므로 상수층의 제한 매핑은 전사이다. 층은 {\it flasque}라고 한다. 이 경우, 모든 상위 Čech 코호몰로지 군이 사라지고 상위 Zariski 코호몰로지 군도 모두 사라집니다. 즉, Zariski 위상에는 이 더 높은 코호몰로지를 감지하기에는 ``충분하지 않습니다'' 열린집합이 있다.

\medskip\noindent 반면, $X$가 매끄러운 투영 복소다양체인 경우 $$
H_{Betti}^{2 \dim X}(X (\mathbf{C}), \Lambda) = \Lambda \quad \text{~을 위한}
\Lambda = \mathbf{Z}, \ \mathbf{Z}/n\mathbf{Z},
$$ 여기서 $X(\mathbf{C})$는 복합체 점의 집합을 의미한다. $X$. 이것은 대수기하학에서 복제하면 좋을 특징이다. 특히 긍정적인 표수에서요.




\section{에탈 위상} \label{etale-cohomology-section-etale-topology}

\noindent 자리스키 위상을 세분하기 위해 열린집합을 단순히 더 ``추가''하는 일은 매우 어렵다. 위상을 정의하는 효율적인 방법 하나는 열린집합뿐 아니라 그 위에 놓인 몇몇 스킴도 함께 고려하는 것이다. 에탈 위상을 정의할 때에는 에탈 사상 $\varphi : U \to X$를 모두 고려한다. $X$가 $\mathbf{C}$ 위의 매끄러운 사영다양체이면, 이는 다음을 뜻한다. \begin{enumerate} \item $U$는 매끄러운 다양체들의 서로소 합이다. \item $\varphi$는 (해석적으로) 국소 동형사상이다. \end{enumerate} 여기서 ``해석적으로''란 $\mathbf{C}$ 위의 통상적인 (초월적) 위상에 관한 말이다. 따라서 두 번째 조건은 $\varphi$의 미분이 모든 점에서 최대 계수를 갖는다는 뜻이다(특히 $U$의 모든 연결성분은 $X$와 같은 차원을 갖는다).

\medskip\noindent 이중 표지 -- 다양체 사이의 유한 차수 $2$ 맵으로 느슨하게 정의됨 -- 예에 대해 $$
\Spec(\mathbf{C}[t])
\longrightarrow
\Spec(\mathbf{C}[t]),
\quad t \longmapsto t^2
$$는 에탈 사상이 아닌 경우 올 단일 점으로 구성된다. 예에서 $t = 0$일 때 이런 일이 발생한다. 다양체와 $\mathbf{C}$ 사이의 유한 맵이 에탈이 되려면 모든 올의 포인트 수가 동일해야 한다. 예의 지도 소스에서 $t = 0$ 지점을 제거하면 사상 에탈이 된다. 하지만 사상의 소스에서 다른 점도 제거할 수 있으며 사상은 여전히 에탈이다. 에탈 위상을 고려하려면 그러한 사상을 모두 살펴봐야 한다. 자리스키 위상와는 달리, 이것들은 항상 $X$의 열린 부분집합일 필요는 없다.

\begin{definition} \label{etale-cohomology-definition-etale-covering-initial} 사상 $\{ \varphi_i : U_i \to X\}_{i \in I}$의 계열은 {\it 에탈 피복}라고 한다. $\varphi_i$는 에탈 사상이고 그 이미지는 $X$, 즉 $X = \bigcup_{i \in I} \varphi_i(U_i)$를 포함한다. \end{definition}

\noindent 이는 에탈 위상을 ``정의''한다. 즉, 이제 층이 무엇인지 말할 수 있다. {\it 에탈 층} $\mathcal{F}$ 의 집합 (각각.\ 아벨 군, 벡터공간, 등) 위의 $X$는 데이터이다: \begin{enumerate} \item 각 에탈 사상 $\varphi : U \to X$ 집합 (각각.\ 아벨 군, 벡터공간 등) $\mathcal{F}(U)$, \item 에탈의 각 쌍 $U, \ U'$에 대해 스킴 위의 $X$, 사상 $U \to U'$ 위의 $X$ (자동으로 에탈) 제한 사상 $\rho^{U'}_U : \mathcal{F}(U') \to \mathcal{F}(U)$ \end{enumerate} 이러한 데이터는 조건, 즉 $\rho^U_U = \text{id}$를 만족해야 한다. 사상 $U \to U$ 그리고 $\rho^{U'}_U \circ \rho^{U''}_{U'} = \rho^{U''}_U$ 사상 $U \to U' \to U''$ 의 스킴의 에탈이 $X$ 위에 있을 때 다음과 같습니다. {\it 층 공리}: \begin{itemize} \item[(*)] 매 에탈 피복 $\{ \varphi_i : U_i \to U\}_{i \in I}$, 다이어그램 $$
\xymatrix{
\mathcal{F} (U) \ar[r] &
\Pi_{i \in I} \mathcal{F} (U_i) \ar@<1ex>[r] \ar@<-1ex>[r] &
\Pi_{i, j \in I} \mathcal{F} (U_i \times_U U_j)
}
$$은 집합의 범주에 있는 등화자 도식입니다(관련:\ 아벨 군, 벡터공간 등). \end{itemize}

\begin{remark} \label{etale-cohomology-remark-i-is-j} 마지막 진술에서 $i = j$가 일반적으로 매우 사소하지 않은 조건인 경우를 잊지 않는 것이 중요합니다(Zariski 위상과 달리). 실제로 중요한 피복에는 요소가 하나만 있는 경우가 많습니다. \end{remark}

\noindent 항등식은 에탈 사상이므로 에탈 층의 대역 절단을 계산할 수 있으며, 코호몰로지는 단순히 해당 오른쪽-유도 함자가 된다. 즉, 일단 이론이 더 발전되고 진술이 정확해지면 코호몰로지를 정의하는 데 아무런 장애가 없을 것이다.




\section{에탈의 위업 위상} \label{etale-cohomology-section-feats}

\noindent 자연수 $n \in \mathbf{N} = \{1, 2, 3, 4, \ldots\}$의 경우 $$
H_\etale^2 (\mathbf{P}^1_\mathbf{C}, \mathbf{Z}/n\mathbf{Z}) =
\mathbf{Z}/n\mathbf{Z}.
$$가 참이다. 보다 일반적으로, $X$가 복소다양체인 경우 해당 에탈 베티 수는 유한에서 계수이다. 체는 $X(\mathbf{C})$의 일반적인 Betti 수치와 일치한다. 즉, $$
\dim_{\mathbf{F}_q} H_\etale^{2i} (X, \mathbf{F}_q) =
\dim_{\mathbf{F}_q} H_{Betti}^{2i} (X(\mathbf{C}), \mathbf{F}_q).
$$ 이는 매우 만족스럽습니다. 그러나 이러한 등식은 일반적으로 적용되지 않고 비틀림 계수에만 적용된다. 정수 계수의 경우 $$
H_\etale^2 (\mathbf{P}^1_\mathbf{C}, \mathbf{Z}) = 0.
$$가 있다. 대조적으로 $H_{Betti}^2(\mathbf{P}^1(\mathbf{C}), \mathbf{Z}) = \mathbf{Z}$는 위상 공간 $\mathbf{P}^1(\mathbf{C})$가 $2$ 구와 동형이다. 곧 설명할 극한 절차를 통해 비틀림에서 비꼬임 계수로 돌아가는 방법이 있다.




\section{계산} \label{etale-cohomology-section-computation}

\noindent
계수를 사용하여 $\mathbf{P}^1_\mathbf{C}$의 코호몰로지를 어떻게 계산합니까?
$\Lambda = \mathbf{Z}/n\mathbf{Z}$?
우리는 Čech 코호몰로지를 사용한다. $\mathbf{P}^1_\mathbf{C}$의 피복이 주어집니다.
두 가지 표준에 따라 $U_0, U_1$이 열립니다. 둘 다
$\mathbf{A}^1_\mathbf{C}$와 동형이고 그 교차점이 동형이다
$\mathbf{A}^1_\mathbf{C} \setminus \{0\} = \mathbf{G}_{m, \mathbf{C}}$로 변경된다.
Mayer-Vietoris 수열은 에탈 코호몰로지에서 성립하는 것으로 나타났습니다.
이는 완전열을 제공한다.
$$
H_\etale^{i-1}(U_0\cap U_1, \Lambda) \to
H_\etale^i(\mathbf{P}^1_\mathbf{C}, \Lambda) \to
H_\etale^i(U_0, \Lambda) \oplus
H_\etale^i(U_1, \Lambda) \to H_\etale^i(U_0\cap U_1,
\Lambda).
$$
우리가 기대하는 답을 얻으려면, 직합
세 번째 용어가 사라집니다. 사실, 일반적인 위상의 경우,
$$
H_\etale^q (\mathbf{A}^1_\mathbf{C}, \Lambda) = 0
\quad \text{임의의 } q \geq 1 \text{에 대하여},
$$
그리고
$$
H_\etale^q (\mathbf{A}^1_\mathbf{C} \setminus \{0\}, \Lambda) = \left\{
\begin{matrix}
\Lambda & \text{(}q = 1\text{일 때)} \\
0 & \text{(}q \geq 2\text{일 때)}.
\end{matrix}
\right.
$$
이러한 결과는 이미 매우 어렵습니다(초급 증명이란 무엇입니까?). 우리를 보자
에탈 코호몰로지의 기계가 다음과 같을 때 이를 어떻게 계산하는지 설명하라.
우리 마음대로.

\medskip\noindent {\bf 더 높음 코호몰로지.} 이는 다음과 같은 일반적인 사실에 의해 처리된다. $X$가 아핀 곡선인 경우 $\mathbf{C}$, 그러면 $$
H_\etale^q (X, \mathbf{Z}/n\mathbf{Z}) = 0 \quad \text{임의의 } q \geq 2 \text{에 대하여}.
$$ 이는 곡선의 일반적인 점을 고려하고 Galois 코호몰로지를 수행하여 증명된다. 따라서 1 정도의 코호몰로지만 걱정하면 된다.

\medskip\noindent
{\bf 코호몰로지 정도 1.} 다음 식별을 사용한다.
\begin{eqnarray*}
H_\etale^1 (X, \mathbf{Z}/n\mathbf{Z}) = \left\{
\begin{matrix}
\text{집합의 층 }\mathcal{F}\text{이고 에탈 사이트 }X_\etale
\text{ 위에서 다음을 갖는 것} \\
\text{작용 }\mathbf{Z}/n\mathbf{Z} \times \mathcal{F} \to \mathcal{F}
\text{이고 }\mathcal{F}\text{가 }\mathbf{Z}/n\mathbf{Z}\text{-토서인 것}
\end{matrix}
\right\}
\Big/ \cong
\\
 = \left\{
\begin{matrix}
\text{유한 에탈 사상 }Y \to X\text{와 다음을 만족하는 자유로운} \\
\mathbf{Z}/n\mathbf{Z}\text{ 작용을 함께 갖는 것:}
X = Y/(\mathbf{Z}/n\mathbf{Z}).
\end{matrix}
\right\}
\Big/ \cong.
\end{eqnarray*}
첫 번째 식별은 매우 일반적입니다(모든 코호몰로지 이론에 해당됨).
사이트)에 있으며 에탈 위상과는 아무런 관련이 없다. 두 번째
식별은 하강 이론의 결과이다. 마지막 집합은
우리가 손에 넣을 수 있는 기하학적 물체의 모음이다.

\medskip\noindent 곡선 $\mathbf{A}^1_\mathbf{C}$에는 중요하지 않은 유한 에탈 피복이 없으므로 $H_\etale^1 (\mathbf{A}^1_\mathbf{C}, \mathbf{Z}/n\mathbf{Z}) = 0$이다. 이는 위상적으로 또는 다음 단락의 인수를 사용하여 볼 수 있다.

\medskip\noindent
유한 에탈 피복
$\varphi : Y \to \mathbf{A}^1_\mathbf{C} \setminus \{0\}$.
$Y$이 다음과 같은 경우를 고려하면 충분한다.
연결로 가정한다. 우리는 $Y$가 무엇인지 알아볼 것이다.
Riemann-Hurwitz 공식을 적용함으로써(물론 이것은 약간 어리석은 일이다.
원한다면 다음 절을 건너뛸 수 있다.
이 사상은 $n$ ~ 1라고 가정하고
투영적 압축
$$
\xymatrix{
{Y\ } \ar@{^{(}->}[r] \ar[d]^\varphi &
{\bar Y} \ar[d]^{\bar\varphi} \\
{\mathbf{A}^1_\mathbf{C} \setminus \{0\}} \ar@{^{(}->}[r] &
{\mathbf{P}^1_\mathbf{C}}
}
$$
$\varphi$는 에탈이고 파급효과가 없더라도 $\bar{\varphi}$는
0 및 $\infty$에서 분기된다. 0의 사전 이미지가 $y_1,
\ldots, y_r$ 와 함께 indices 의 ramification $e_1, \ldots e_r$, 그리고
$\infty$의 사전 이미지는 인덱스가 있는 $y_1', \ldots, y_s'$ 점이다.
파급 효과 $d_1, \ldots d_s$. 특히, $\sum e_i = n = \sum d_j$.
Riemann-Hurwitz 공식을 적용하면 다음을 얻습니다.
$$
2 g_Y - 2 = -2n + \sum (e_i - 1) + \sum (d_j - 1)
$$
따라서 $g_Y = 0$, $r = s = 1$ 및 $e_1 = d_1 = n$이다.
그러므로 $Y \cong {\mathbf{A}^1_\mathbf{C} \setminus \{0\}}$,
저것 좀 봐$\varphi(z) = \lambda z^n$어떤 사람들에게는$\lambda \in \mathbf{C}^*$.
$Y$를 다시 매개변수화한 후 $\lambda = 1$라고 가정할 수 있다. 따라서 우리의
피복은 좌표의 $n$번째 근을 취하여 주어집니다.
$\mathbf{A}^1_{\mathbf{C}} \setminus \{0\}$.

\medskip\noindent
피복을 분류해야 한다는 점을 기억하세요.
${\mathbf{A}^1_\mathbf{C} \setminus \{0\}}$ 무료와 함께
$\mathbf{Z}/n\mathbf{Z}$-그들에 대한 조치.
우리의 경우 그러한 조치는 해당
$Y$의 자동형성으로 $z$을 $\zeta_n z$로 전송한다. 여기서 $\zeta_n$는
원시적 $n$번째 통일근. $\phi(n)$ 그러한 행위가 있습니다
(여기서 $\phi(n)$는 오일러 함수를 의미한다.) 따라서 정확하게는
$\phi(n)$ 연결 유한 에탈 피복 주어진 자유
$\mathbf{Z}/n\mathbf{Z}$-액션, 각각 기본 요소에 해당
$n$통일의 근원. 우리는 독자들에게 그것이 무엇인지 확인하도록 맡깁니다.
연결이 끊어진 유한 에탈 정도 $n$ 피복 의
$\mathbf{A}^1_{\mathbf{C}} \setminus \{0\}$ 주어진 무료
$\mathbf{Z}/n\mathbf{Z}$-동작은 $n$번째와 일대일로 대응된다.
원시적이지 않은 $1$의 루트이다.
즉, 이 계산은 다음을 보여줍니다.
$$
H_\etale^1 (\mathbf{A}^1_\mathbf{C} \setminus \{0\},
\mathbf{Z}/n\mathbf{Z}) =
\Hom(\mu_n(\mathbf{C}), \mathbf{Z}/n\mathbf{Z}) \cong \mathbf{Z}/n\mathbf{Z}.
$$
첫 번째 식별은 정식이고 두 번째 식별은 그렇지 않습니다.
비고 \ref{etale-cohomology-remark-normalize-H1-Gm}.
Riemann-Hurwitz의 증명은 다음의 계산을 사용하지 않으므로
코호몰로지, 위의 내용은 실제로 증명을 구성합니다(단,
소멸 등의 세부정보를 입력하세요.




\section{비꼬임 계수}
\label{etale-cohomology-section-nontorsion}

\noindent
비꼬임 계수를 연구하려면 다음과 같이 정의를 만듭니다.
$$
H_\etale^i (X, \mathbf{Q}_\ell) :=
\left( \lim_n H_\etale^i(X, \mathbf{Z}/\ell^n\mathbf{Z}) \right)
\otimes_{\mathbf{Z}_\ell} \mathbf{Q}_\ell.
$$
$\lim_n$ 기호는 {\it 극한}를 나타냅니다.
코호몰로지 군 $H_\etale^i(X, \mathbf{Z}/\ell^n\mathbf{Z})$ 색인이 생성됨
$n$로, 참조
범주론, 절 \ref{categories-section-posets-limits}.
따라서 우리는 어느 정도 호환성을 만족하는 층 시스템을 연구해야 할 것이다.
정황.




\section{층 이론} \label{etale-cohomology-section-sheaf-theory} %9.10.09

\noindent 이제부터 사이트와 층에 대해 본격적으로 이야기하기 시작한다. 다음 몇 가지 절에 들어갈 수 있는 유용한 추상 자료가 엄청나게 많습니다. 이 자료 중 일부는 사이트, 가군의 사이트, 코호몰로지의 사이트 장과 같은 이전 장에서 해결되었다. 우리는 여기에 너무 많은 자료를 추가하지 않으려고 노력한다. 이 장의 뒷부분에 나오는 자료가 이해가 될 만큼만 추가한다.




\section{전층}
\label{etale-cohomology-section-presheaves}

\noindent 이 절에 대한 참조는 사이트, 절 \ref{sites-section-presheaves}이다.

\begin{definition} \label{etale-cohomology-definition-presheaf} $\mathcal{C}$를 범주로 둡니다. {\it 전층 의 집합} (각각 {\it 아벨리안 전층}) 위의 $\mathcal{C}$는 함자 $\mathcal{C}^{opp} \to
\textit{Sets}$ (관련:\ $\textit{Ab}$)이다. \end{definition}

\noindent
{\bf 용어.} $U \in \Ob(\mathcal{C})$인 경우 다음 요소는
$\mathcal{F}(U)$는 {\it 절}의 $\mathcal{F}$라고 불립니다.
$U$. 을 위한$\varphi : V \to U$~에$\mathcal{C}$,
지도 $\mathcal{F}(\varphi) : \mathcal{F}(U) \to \mathcal{F}(V)$
{\it 제한 사상}라고 하며 종종 $s \mapsto s|_V$로 표시된다.
또는 때로는 $s \mapsto \varphi^*s$. $s|_V$ 표기가 모호한다.
제한 사상은 $\varphi$에 의존하지만 이는 표준이므로
표기법 남용. 우리는 또한 표기법을 사용합니다
$\Gamma(U, \mathcal{F}) = \mathcal{F}(U)$.

\medskip\noindent $\mathcal{F}$가 함자라는 것은 $W \to V \to U$가 $\mathcal{C}$ 및 $s \in \Gamma(U, \mathcal{F})$에서 사상인 경우를 의미한다. $(s|_V)|_W = s |_W$, 방금 본 표기법의 남용이 있다. 또한, ID 사상 $\text{id}_U : U \to U$에 해당하는 제한 매핑이 ID이다.

\medskip\noindent
범주의 전층의 집합 (각각 아벨 전층)에
$\mathcal{C}$는 $\textit{PSh} (\mathcal{C})$로 표시된다. (관련 $\textit{PAb}
(\mathcal{C})$). It 이다 범주 의 함자 로부터 $\mathcal{C}^{opp}$ ~
$\textit{Sets}$ (관련 $\textit{Ab}$), 즉 사상
전층은 함자의 자연변환이다. 우리는 단지
범주 $\textit{PSh}(\mathcal{C})$ 및 $\textit{PAb}(\mathcal{C})$
범주 $\mathcal{C}$가 작을 때. (우리의 관례는 범주
달리 언급하지 않는 한 작습니다. 작지 않다면 작아야 한다.
범주, 비고 \ref{categories-remark-big-categories}에 나열되어 있다.)

\begin{example}
\label{etale-cohomology-example-representable-presheaf}
객체 $X \in \Ob(\mathcal{C})$가 주어지면 규칙을 고려한다.
$$
\begin{matrix}
U & \longmapsto & h_X(U) = \Mor_\mathcal{C}(U, X) \\
V \xrightarrow{\varphi} U & \longmapsto &
(\psi \mapsto \psi \circ \varphi) : h_X(U) \to h_X(V).
\end{matrix}
$$
이는 함자 $h_X : \mathcal{C}^{opp} \to \textit{Sets}$를 정의한다.
따라서 전층이 된다. 이것을
{\it $X$에 연결된 전층을 표현할 수 있다.}
모든 위상에서 표현 가능한 전층이 층라는 것은 사실이 아닙니다.
사이트마다.
\end{example}

\begin{lemma}[요네다]
\label{etale-cohomology-lemma-yoneda}
\begin{slogan}
사상 객체 사이는 자연변환으로 전단사된다.
함자 사이에 표시된다.
\end{slogan}
$\mathcal{C}$를 범주로 하고 $X, Y \in
\Ob(\mathcal{C})$. 자연스러운 전향이 있다
$$
\begin{matrix}
\Mor_\mathcal{C}(X, Y) &
\longrightarrow &
\Mor_{\textit{PSh}(\mathcal{C})} (h_X, h_Y) \\
\psi &
\longmapsto &
h_\psi = \psi \circ - : h_X \to h_Y.
\end{matrix}
$$
\end{lemma}

\begin{proof} 범주, 보조정리 \ref{categories-lemma-yoneda}를 참조하라. \end{proof}




\section{사이트}
\label{etale-cohomology-section-sites}


\begin{definition} \label{etale-cohomology-definition-family-morphisms-fixed-target} $\mathcal{C}$를 범주로 둡니다. {\it 고정 대상} $\mathcal{U} = \{\varphi_i : U_i \to U\}_{i\in I}$가 있는 사상 계열은 \begin{enumerate} \item 개체 $U \in \mathcal{C}$의 데이터이다. \item 집합 $I$ (비어 있을 수 있음) 및 \item 모든 $i\in I$에 대해, 사상 $\varphi_i : U_i \to U$ 대상이 있는 $\mathcal{C}$ $U$. \end{enumerate} \end{definition}

\noindent
사상 계열의 {\it 사상라는 개념이 있다.
대상}. 그 특별한 경우는 {\it 개선}의 개념이다.
해당 자료에 대한 참고자료는
사이트, 절 \ref{sites-section-refinements}.

\begin{definition}
\label{etale-cohomology-definition-site}
 {\it 사이트}\footnote{우리가 사이트라고 부르는 것은 범주라고 불리는 것이다.
\cite[Expos\'e II, D\'efinition 1.3]{SGA4}의 사전 위상.
\cite{ArtinTopologies}에서는 Grothendieck과 함께 범주라고 한다.
위상.}는 범주 $\mathcal{C}$ 및 집합으로 구성된다.
$\text{Cov}(\mathcal{C})$는 고정 대상이 있는 사상 계열로 구성된다.
{\it 피복}라고 한다.
\begin{enumerate}
\item (동형사상) $\varphi : V \to U$가 $\mathcal{C}$의 동형사상인 경우,
그러면 $\{\varphi : V \to U\}$는 피복이다.
\item (지역) $\{\varphi_i : U_i \to U\}_{i\in I}$가 피복이고
모든 $i \in I$에 대해 피복이 제공된다.
$\{\psi_{ij} : U_{ij} \to U_i \}_{j\in I_i}$, 그런 다음
$$
\{
\varphi_i \circ \psi_{ij} : U_{ij} \to U
\}_{(i, j)\in \prod_{i\in I} \{i\} \times I_i}
$$
또한 피복이며,
\item (기본 변경) 만약 $\{U_i \to U\}_{i\in I}$
는 피복이고 $V \to U$는 $\mathcal{C}$에서 사상이다.
\begin{enumerate}
\item 모든 $i \in I$에 대해 올곱
$U_i \times_U V$이 $\mathcal{C}$에 존재하며,
\item $\{U_i \times_U V \to V\}_{i\in I}$는 피복이다.
\end{enumerate}
\end{enumerate}
\end{definition}

\noindent 우리에게 사이트의 기초가 되는 범주는 항상 ``작은''이다. 즉, 객체의 컬렉션은 집합을 형성하고 피복의 컬렉션은 사이트는 집합이다. (위의 정의에서도 마찬가지) 이 장에서는 객체 컬렉션과 피복을 집합으로 줄이는 인수를 대부분 생략할 것이다. 자세한 내용은 사이트, 비고 \ref{sites-remark-no-big-sites}를 참조하라.

\begin{example} \label{etale-cohomology-example-site-topological-space} $X$가 위상 공간인 경우 연관된 사이트 $X_{Zar}$가 다음과 같이 정의된다. $X_{Zar}$의 객체는 열린 부분집합이다. $X$, 이들 사이의 사상은 포함 매핑이고 피복은 일반적인 위상(사사) 피복이다. $U, V \subset W \subset X$가 열린 부분집합인 경우 $U \times_W V = U \cap V$가 존재하는지 확인하세요. 이 범주에는 올곱이 있다. 모든 확인은 간단하며 모든 것이 예상대로 작동한다. \end{example}




\section{층}
\label{etale-cohomology-section-sheaves}

\begin{definition} \label{etale-cohomology-definition-sheaf} 전층 $\mathcal{F}$ 의 집합 (관련 abelian 전층)는 사이트 $\mathcal{C}$에 있다고 한다. {\it 분리된 전층} 모든 피복 $\{\varphi_i : U_i \to U\}_{i\in I} \in \text{Cov} (\mathcal{C})$에 대해 지도 $$
\mathcal{F}(U) \longrightarrow \prod\nolimits_{i\in I} \mathcal{F}(U_i)
$$는 단사형이다. 여기서 지도는 $s \mapsto (s|_{U_i})_{i\in I}$이다. 전층 $\mathcal{F}$는 {\it 층}이다. 모든 피복 $\{\varphi_i : U_i \to U\}_{i\in I} \in \text{Cov} (\mathcal{C})$에 대해 다이어그램은 \begin{equation}
\label{etale-cohomology-equation-sheaf-axiom}
\xymatrix{
\mathcal{F}(U) \ar[r] &
\prod_{i\in I} \mathcal{F}(U_i) \ar@<1ex>[r] \ar@<-1ex>[r] &
\prod_{i, j \in I} \mathcal{F}(U_i \times_U U_j),
}
\end{equation} 여기서 첫 번째 맵은 $s \mapsto (s|_{U_i})_{i\in I}$이고 두 맵은 오른쪽에는 $(s_i)_{i\in I} \mapsto (s_i |_{U_i \times_U U_j})$ 및 $(s_i)_{i\in I} \mapsto (s_j |_{U_i \times_U U_j})$가 있으며, 집합의 범주에 있는 등화자 도식입니다(관련:\ 아벨 군). \end{definition}

\begin{remark} \label{etale-cohomology-remark-empty-covering} 빈 피복(여기서 $I = \emptyset$)의 경우 이는 $\mathcal{F}(\emptyset)$가 해당 범주($\textit{Sets}$ 및 $\textit{Ab}$). \end{remark}

\begin{example} \label{etale-cohomology-example-sheaf-site-space} 위상 공간과 연관된 사이트 $X_{Zar}$에 대해 이 작업을 수행하려면 예 \ref{etale-cohomology-example-site-topological-space}를 참조하라. 층의 일반적인 개념을 제공한다. \end{example}

\begin{definition}
\label{etale-cohomology-definition-category-sheaves}
$\Sh(\mathcal{C})$ (각각.\ $\textit{Ab}(\mathcal{C})$)를 나타냅니다.
$\textit{PSh}(\mathcal{C})$의 전체 하위 카테고리
(각각.\ $\textit{PAb}(\mathcal{C})$) 객체는 층이다. 이것은
{\it 범주 의 층 의 집합} (관련\ {\it 아벨 층}) 켜짐
$\mathcal{C}$.
\end{definition}




\section{G-집합} \label{etale-cohomology-section-G-sets}의 예

\noindent
$G$을 그룹으로 두고 사이트 $\mathcal{T}_G$를 다음과 같이 정의한다.
범주는 $G$-집합의 범주이다. 즉, 해당 개체에는 집합이 부여된다.
왼쪽 $G$ 동작과 사상은 등변 맵이다. 그리고
$\mathcal{T}_G$ 중 피복은 가족이다.
$\{\varphi_i : U_i \to U\}_{i\in I}$ 만족스럽다
$U = \bigcup_{i\in I} \varphi_i(U_i)$.

\medskip\noindent
사이트 $\mathcal{T}_G$에는 특수 객체가 있다. 즉, $G$-집합 $G$이다.
왼쪽 번역에 따라 자연스러운 동작이 부여된다. ${}_G G$로 표시한다.
자연 그룹 동형사상이 있음을 관찰하라.
$$
\begin{matrix}
\rho : & G^{opp} & \longrightarrow & \text{Aut}_{G\textit{-Sets}}({}_G G) \\
& g & \longmapsto & (h \mapsto hg).
\end{matrix}
$$
특히, 전층 $\mathcal{F}$에 대해 집합 $\mathcal{F}({}_G G)$
상속받는다$G$-액션을 통해$\rho$. (반공변성(반공변성)에 의해
$\mathcal{F}$, 집합 $\mathcal{F}({}_G G)$는 다시 왼쪽 $G$-집합이다.) 실제로,
함자
$$
\begin{matrix}
\Sh(\mathcal{T}_G) & \longrightarrow & G\textit{-Sets} \\
\mathcal{F} & \longmapsto & \mathcal{F}({}_G G)
\end{matrix}
$$
범주의 동치이다. 준역수는 함자 $X \mapsto
h_X$. 완전한 증명을 제공하지 않고(다음에서 찾을 수 있음)
사이트, 절 \ref{sites-section-example-sheaf-G-sets})
이것이 왜 사실인지 설명해 보자.
\begin{enumerate}
\item
$S$가 $G$-집합인 경우 궤도로 분해할 수 있다. $S = \coprod_{i\in I}
O_i$. 층 axiom 에 대하여 피복 $\{O_i \to S\}_{i\in I}$가 다음과 같이 말한다.
$$
\xymatrix{
\mathcal{F}(S) \ar[r] &
\prod_{i\in I} \mathcal{F}(O_i) \ar@<1ex>[r] \ar@<-1ex>[r] &
\prod_{i, j \in I} \mathcal{F}(O_i \times_S O_j)
}
$$
이퀄라이저이다. $G\textit{-Sets}$에 있는 섬유제품이
$\textit{Sets}$의 섬유제품에서 유도되었으며,
$\mathcal{F}(\emptyset)$는 $G$-싱글턴이다.
$$
\prod_{i, j \in I} \mathcal{F}(O_i \times_S O_j) = \prod_{i \in I}
\mathcal{F}(O_i)
$$
위의 두 지도는 실제로 동일한다. 그러므로 층 공리는 단지
$\mathcal{F}(S) = \prod_{i\in I} \mathcal{F}(O_i)$라고 한다.
\item
$S$가 $G$-집합 $S= G/H$이고 $\mathcal{F}$가 $\mathcal{T}_G$에서 층인 경우,
그러면 우리는 주장
$$
\mathcal{F}(G/H) = \mathcal{F}({}_G G)^H
$$
특히 $\mathcal{F}(\{*\}) = \mathcal{F}({}_G G)^G$. 이것을 보려면,
를 사용하자층에 대한 공리피복 $\{ {}_G G \to G/H \}$~의$S$. 우리
가지다
\begin{eqnarray*}
{}_G G \times_{G/H} {}_G G & \cong & G \times H \\
(g_1, g_2) & \longmapsto & (g_1, g_1^{-1} g_2)
\end{eqnarray*}
는 ${}_G G$ 복사본의 분리된 결합입니다($G$-집합). 따라서 층 공리는
읽다
$$
\xymatrix{
\mathcal{F} (G/H) \ar[r] &
\mathcal{F}({}_G G) \ar@<1ex>[r] \ar@<-1ex>[r] &
\prod_{h\in H} \mathcal{F}({}_G G)
}
$$
여기서 오른쪽의 두 지도는 $s \mapsto (s)_{h \in H}$ 및 $s \mapsto
(hs)_{h \in H}$. 따라서 $\mathcal{F}(G/H) = \mathcal{F}({}_G G)^H$
주장했다.
\end{enumerate}
이것은 주장된 범주의 동치를 완전히 증명하지는 않지만 다음과 같이 표시된다.
적어도 층 $\mathcal{F}$는 절에 의해 완전히 결정된다.
${}_G G$. 세부 사항(및 집합 이론적 설명)은 다음에서 찾을 수 있다.
사이트, 절 \ref{sites-section-example-sheaf-G-sets}.




\section{층화}
\label{etale-cohomology-section-sheafification}

\begin{definition}
\label{etale-cohomology-definition-0-cech}
$\mathcal{F}$를 사이트 $\mathcal{C}$에서 전층으로 설정하고
$\mathcal{U} = \{U_i \to U\} \in \text{Cov} (\mathcal{C})$.
{\it 0번째 Čech 코호몰로지 군}을 정의한다.
$\mathcal{F}$ $\mathcal{U}$에 대하여
$$
\check H^0 (\mathcal{U}, \mathcal{F}) =
\left\{
(s_i)_{i\in I} \in \prod\nolimits_{i\in I }\mathcal{F}(U_i)
\text{ }
s_i|_{U_i \times_U U_j} = s_j |_{U_i \times_U U_j}
\right\}.
$$
\end{definition}

\noindent
표준지도가 있습니다
$\mathcal{F}(U) \to \check H^0 (\mathcal{U}, \mathcal{F})$,
$s \mapsto (s |_{U_i})_{i\in I}$.
{\it 사상 의 피복} 로부터 피복
$\mathcal{V} = \{V_j \to V\}_{j \in J}$ ~ $\mathcal{U}$은 트리플이다.
$(\chi, \alpha, \chi_j)$, 여기서
$\chi : V \to U$은 사상이다.
$\alpha : J \to I$는 집합의 맵이며, 모든
$j \in J$ 사상 $\chi_j$는 가환 도식에 맞습니다.
$$
\xymatrix{
V_j \ar[rr]_{\chi_j} \ar[d] & & U_{\alpha(j)} \ar[d] \\
V \ar[rr]^\chi & & U.
}
$$
우리가 정의하는 데이터 $\chi, \alpha, \{\chi_j\}_{j \in J}$가 주어지면
\begin{eqnarray*}
\check H^0(\mathcal{U}, \mathcal{F}) & \longrightarrow &
\check H^0(\mathcal{V}, \mathcal{F}) \\
(s_i)_{i\in I} & \longmapsto &
\left(\chi_j^*\left(s_{\alpha(j)}\right)\right)_{j\in J}.
\end{eqnarray*}
그러면 우리는 주장
\begin{enumerate}
\item 지도가 잘 정의되어 있고
\item 이 사상은 $\chi$에만 의존하며 다음 선택과 무관하다.
$\alpha, \{\chi_j\}_{j \in J}$.
\end{enumerate}
첫 번째 사실의 증명을 생략한다.
(2) 부분을 보려면 다른 트리플 $(\psi, \beta, \psi_j)$을 고려하세요.
$\chi = \psi$. 그러면 가환 도식
$$
\xymatrix{
V_j \ar[rrr]_{(\chi_j, \psi_j)} \ar[dd] & & &
U_{\alpha(j)} \times_U U_{\beta(j)} \ar[dl] \ar[dr] \\
& & U_{\alpha(j)} \ar[dr] & &
U_{\beta(j)} \ar[dl] \\
V \ar[rrr]^{\chi = \psi} & & & U.
}
$$
절 $s \in \mathcal{F}(\mathcal{U})$가 주어지면 해당 이미지는
$\mathcal{F}(V_j)$ 제공한 지도 아래
$(\chi, \alpha, \{\chi_j\}_{j \in J})$
$\chi_j^*s_{\alpha(j)}$이며,
$(\psi, \beta, \{\psi_j\}_{j \in J})$에 의해 제공된 지도 아래의 이미지
$\psi_j^*s_{\beta(j)}$이다. 이것들
가정 $s \in \check H^0(\mathcal{U}, \mathcal{F})$에 의해 두 개는 동일한다.
따라서 둘 다 공통 가치의 하락과 동일한다.
$$
s_{\alpha(j)}|_{U_{\alpha(j)} \times_U U_{\beta(j)}} =
s_{\beta(j)}|_{U_{\alpha(j)} \times_U U_{\beta(j)}}
$$
다이어그램의 지도 $(\chi_j, \psi_j)$에 의해 뒤로 당겨졌습니다.

\begin{theorem}
\label{etale-cohomology-theorem-sheafification}
$\mathcal{C}$를 사이트로 하고 $\mathcal{F}$를 전층으로 $\mathcal{C}$로 설정한다.
\begin{enumerate}
\item 규칙
$$
U \mapsto \mathcal{F}^+(U) :=
\colim_{\mathcal{U} \text{가 }U\text{의 피복}}
\check H^0(\mathcal{U}, \mathcal{F})
$$
전층이다. 그리고 여극한은 방향이 있는 것이다.
\item 전층 $\mathcal{F} \to \mathcal{F}^+$의 표준 맵이 있다.
\item $\mathcal{F}$가 분리된 전층인 경우 $\mathcal{F}^+$는 층이다.
(2)의 맵은 단사적이다.
\item $\mathcal{F}^+$는 분리된 전층이다.
\item $\mathcal{F}^\# = (\mathcal{F}^+)^+$는 층이고 표준
지도는 함수 동형사상을 유도한다.
$$
\Hom_{\textit{PSh}(\mathcal{C})}(\mathcal{F}, \mathcal{G}) =
\Hom_{\Sh(\mathcal{C})}(\mathcal{F}^\#, \mathcal{G})
$$
$\mathcal{G} \in \Sh(\mathcal{C})$에 대해.
\end{enumerate}
\end{theorem}

\begin{proof} 사이트, 정리 \ref{sites-theorem-plus}를 참조하라. \end{proof}

\noindent 즉, 자연지도 $\mathcal{F} \to \mathcal{F}^\#$는 망각하는 함자 $\Sh(\mathcal{C}) \to \textit{PSh}(\mathcal{C})$에 대한 왼쪽 수반라는 뜻이다.




\section{코호몰로지}
\label{etale-cohomology-section-cohomology}

\noindent 사이트에서 아벨 층에 대해 코호몰로지를 정의할 수 있는 기본 결과는 다음과 같습니다.

\begin{theorem} \label{etale-cohomology-theorem-enough-injectives} 사이트에 있는 아벨 층의 범주는 충분한 단사를 갖는 아벨 범주이다. \end{theorem}

\begin{proof} 사이트, 보조정리 \ref{sites-modules-lemma-abelian-abelian} 및 분사, 정리 \ref{injectives-theorem-sheaves-injectives}에 대한 가군을 참조하라. \end{proof}

\noindent
따라서 코호몰로지를 오른쪽 유도 함자로 정의할 수 있다.
절 함자: $U \in \Ob(\mathcal{C})$인 경우
$\mathcal{F} \in \textit{Ab}(\mathcal{C})$,
$$
H^p(U, \mathcal{F}) :=
R^p\Gamma(U, \mathcal{F}) =
H^p(\Gamma(U, \mathcal{I}^\bullet))
$$
여기서 $\mathcal{F} \to \mathcal{I}^\bullet$는 단사 분해이다. 해야 할 일
이를 위해서는 함자 $\Gamma(U, -)$가 왼쪽 완전인지 확인해야 한다. 이것은
참이고 범주 $\textit{Ab}(\mathcal{C})$가 아벨인 이유의 일부이다.
보다
가군 사이트, 보조정리 \ref{sites-modules-lemma-abelian-abelian}.
사이트의 코호몰로지에 대한 보다 일반적인 논의는 다음을 포함한다.
대역 절단 함자 그리고 오른쪽 유도 함자), 참조
코호몰로지 사이트, 절 \ref{sites-cohomology-section-cohomology-sheaves}.



\section{fpqc 위상} \label{etale-cohomology-section-fpqc} %9.15.09

\noindent 에탈 코호몰로지를 수행하기 전에 fpqc 위상에 대해 조금 공부한다. 이는 준연접층에 잘 작동하기 때문이다.

\begin{definition}
\label{etale-cohomology-definition-fpqc-covering}
$T$를 스킴으로 설정한다. {\it fpqc 피복}는 $T$의 제품군이다.
$\{ \varphi_i : T_i \to T\}_{i \in I}$ 이렇게
\begin{enumerate}
\item 각 $\varphi_i$는 평탄 사상이고
$\bigcup_{i\in I} \varphi_i(T_i) = T$ 그리고
\item 각 아핀 개방 $U \subset T$에 대해 유한이 존재한다.
집합 $K$, 지도 $\mathbf{i} : K \to I$ 및 아핀이 열립니다.
$U_{\mathbf{i}(k)} \subset T_{\mathbf{i}(k)}$ 이렇게
$U = \bigcup_{k \in K} \varphi_{\mathbf{i}(k)}(U_{\mathbf{i}(k)})$.
\end{enumerate}
\end{definition}

\begin{remark} \label{etale-cohomology-remark-fpqc} 첫 번째 조건은 fp에 해당하며, 이는 프랑스어로 {\it fid\`element plat}, 충실히 평탄을 의미하며, qc에 이어 두 번째, {\it 준소형}. 첫 번째 조건의 두 번째 부분은 두 번째 조건이 유지되는 경우 불필요한다. \end{remark}

\begin{example} \label{etale-cohomology-example-fpqc-coverings} fpqc 피복의 예이다. \begin{enumerate} \item $T$의 모든 Zariski 오픈 피복은 fpqc 피복이다. \item 패밀리 $\{\Spec(B) \to \Spec(A)\}$는 fpqc 피복 필요충분조건은 $A \to B$는 충실히 평탄 환 맵이다. \item $f: X \to Y$가 평면적이고 전사적이며 준압축이면 $\{ f: X\to
Y\}$는 fpqc 피복이다. \item 사상 $\varphi :
\coprod_{x \in \mathbf{A}^1_k} \Spec(\mathcal{O}_{\mathbf{A}^1_k, x})
\to \mathbf{A}^1_k$는 $k$가 체인 경우 단순하고 전사적이다. 준컴팩트하지 않으며 실제로 $\{\varphi\}$ 계열은 fpqc 피복이 아닙니다. \item $\mathbf{A}^2_k = \Spec(k[x, y])$를 씁니다. $i_x : D(x) \to \mathbf{A}^2_k$ 및 $i_y : D(y) \to \mathbf{A}^2_k$ 표준이 열림을 나타냅니다. 그런 다음 $\{i_x, i_y, \Spec(k[[x, y]]) \to \mathbf{A}^2_k\}$ 및 $\{i_x, i_y, \Spec(\mathcal{O}_{\mathbf{A}^2_k, 0}) \to \mathbf{A}^2_k\}$ 계열은 fpqc 피복이다. \end{enumerate} \end{example}

\begin{lemma} \label{etale-cohomology-lemma-site-fpqc} 스킴의 범주에 있는 fpqc 피복의 컬렉션은 사이트의 공리를 충족한다. \end{lemma}

\begin{proof} 위상, 보조정리 \ref{topologies-lemma-fpqc}를 참조하라. \end{proof}

\noindent 이 보조정리를 사용하면 스킴의 범주의 fpqc 사이트를 정의할 수 있는 것 같습니다. 그러나 fpqc 위상을 고려할 때 집합 이론적 문제가 발생한다. 위상, 절 \ref{topologies-section-fpqc}를 참조하라. 이는 사이트가 ``작은''이라는 요구 사항에서 비롯되지만, 스킴의 작은 범주는 비어 있지 않은 주어진 스킴의 fpqc 피복의 동종 시스템을 포함할 수 없다. 엄밀히 말하면 이것이 ``부분'' fpqc 사이트 정의를 막는 것은 아니지만, 그렇게 하는 것이 현명해 보이지는 않습니다. 해결 방법은 fpqc 위상(아래 참조)에 대해 층 개념을 허용하지만 모든 fpqc 층의 범주 고려를 금지하는 것이다.

\begin{definition} \label{etale-cohomology-definition-sheaf-property-fpqc} $S$를 스킴이라 하자. $S$ 위의 스킴들의 범주를 $\Sch/S$로 나타낸다. 함자 $\mathcal{F} : (\Sch/S)^{opp} \to \textit{Sets}$, 즉 집합의 전층을 생각하자. $\mathcal{F}$가 모든 fpqc 피복 $\{U_i \to U\}_{i \in I}$ 중 $S$ 위의 스킴들로 이루어진 것에 대하여 도식 (\ref{etale-cohomology-equation-sheaf-axiom})이 등화자 도식이 되면, {\it fpqc 위상에 대한 층 조건을 만족한다}고 한다. \end{definition}

\noindent
마찬가지로 $\mathcal{F}$
{\it 자리스키 위상에 대한 층 조건을 만족한다}고 한다.
피복 $U = \bigcup_{i \in I} U_i$ 다이어그램을 열 때마다
(\ref{etale-cohomology-equation-sheaf-axiom})는 등화자 도식이다. 보다
스킴, 정의 \ref{schemes-definition-representable-by-open-immersions}.
분명히 이것은 모든 사람에 대해 다음과 같이 말하는 것과 같습니다.스킴 $T$~ 위에$S$그만큼
$\mathcal{F}$의 $T$ 개방에 대한 제한은 (일반적인) 층이다.

\begin{lemma}
\label{etale-cohomology-lemma-fpqc-sheaves}
$\mathcal{F}$를 $\Sch/S$에서 전층으로 설정한다. 그 다음에
$\mathcal{F}$는 fpqc 위상에 대해 층 성질을 충족한다.
필요충분조건은
\begin{enumerate}
\item $\mathcal{F}$는 층 성질을 만족한다.
자리스키 위상, 그리고
\item 모든 충실히 평탄 사상 $\Spec(B) \to \Spec(A)$
$S$에 대한 아핀 스킴의 층 공리는 피복에 대해 유지된다.
$\{\Spec(B) \to \Spec(A)\}$. 즉, 이는 다음을 의미한다.
$$
\xymatrix{
\mathcal{F}(\Spec(A)) \ar[r] &
\mathcal{F}(\Spec(B)) \ar@<1ex>[r] \ar@<-1ex>[r] &
\mathcal{F}(\Spec(B \otimes_A B))
}
$$
등화자 도식이다.
\end{enumerate}
\end{lemma}

\begin{proof} 위상, 보조정리 \ref{topologies-lemma-sheaf-property-fpqc}를 참조하라. \end{proof}

\noindent
전층 $\mathcal{F}$를 생각하는 또 다른 방법
$\Sch/S$는 층 조건을 만족한다.
fpqc 위상 는 다음 데이터와 같습니다.
\begin{enumerate}
\item 각 $T/S$에 대해 평소(예: Zariski) 층 $\mathcal{F}_T$
$T_{Zar}$,
\item 모든 사상 $f : T' \to T$ 중 $S$ 위의 것에 대하여, 제한 사상
$f^{-1}\mathcal{F}_T \to \mathcal{F}_{T'}$
\end{enumerate}
그렇게
\begin{enumerate}
\item[()] 제한 매핑은 기능적이다.
\item[(b)] $f : T' \to T$가 열린 몰입인 경우 제한 사항은 다음과 같습니다.
매핑 $f^{-1}\mathcal{F}_T \to \mathcal{F}_{T'}$은 동형사상이고
\item[(c)] 모든 충실히 평탄 사상
$\Spec(B) \to \Spec(A)$ 오버 $S$, 다이어그램
$$
\xymatrix{
\mathcal{F}_{\Spec(A)}(\Spec(A)) \ar[r] &
\mathcal{F}_{\Spec(B)}(\Spec(B)) \ar@<1ex>[r] \ar@<-1ex>[r] &
\mathcal{F}_{\Spec(B \otimes_A B)}(\Spec(B \otimes_A B))
}
$$
이퀄라이저이다.
\end{enumerate}
데이터 (1) 및 (2) 및 조건 (), (b)는 전층의 데이터를 제공한다.
자리스키 위상에 대해 층 조건을 만족하는 $\Sch/S$에서.
보조정리 \ref{etale-cohomology-lemma-fpqc-sheaves} 조건 (c)에 의해 다음을 얻는 데 충분한다.
층 조건 fpqc 위상에 대한 것이다.

\begin{example}
\label{etale-cohomology-example-quasi-coherent}
전층을 고려해보세요.
$$
\begin{matrix}
\mathcal{F} : & (\Sch/S)^{opp} & \longrightarrow & \textit{Ab} \\
& T/S & \longmapsto & \Gamma(T, \Omega_{T/S}).
\end{matrix}
$$
국소화와 차동 장치의 호환성은 다음을 의미한다.
$\mathcal{F}$는 Zariski 사이트의 층이다.
그러나 fpqc 위상에 대해서는 층 조건을 만족하지 않습니다.
즉, 경우를 고려
$S = \Spec(\mathbf{F}_p)$ 및 사상
$$
\varphi :
V = \Spec(\mathbf{F}_p[v])
\to
U = \Spec(\mathbf{F}_p[u])
$$
$u$을 $v^p$에 매핑하여 제공된다. $\{\varphi\}$ 계열은 fpqc 피복이다.
아직 제한 매핑
$\mathcal{F}(U) \to \mathcal{F}(V)$
생성기 $\text{d}u$를 $\text{d}(v^p) = 0$로 보냅니다.
그것은 제로 맵이고 다이어그램입니다
$$
\xymatrix{
\mathcal{F}(U) \ar[r]^{0} &
\mathcal{F}(V) \ar@<1ex>[r] \ar@<-1ex>[r] &
\mathcal{F}(V \times_U V)
}
$$
이퀄라이저가 아닙니다. $\mathcal{F}$가 실제로 수행된다는 사실은 나중에 살펴보자.
에탈에 층이 생기고 사이트가 부드러워집니다.
\end{example}

\begin{lemma} \label{etale-cohomology-lemma-representable-sheaf-fpqc} $\Sch/S$에서 표현 가능한 모든 전층은 fpqc 위상에 대한 층 조건을 충족한다. \end{lemma}

\begin{proof} 하강, 보조정리 \ref{descent-lemma-fpqc-universal-effective-epimorphisms} 참조. \end{proof}

\noindent 이 사실의 증명은 준연접층에 대해 하강을 사용하므로 이에 대해 나중에 다시 설명하겠습니다. 이에 대해서는 다음 절에서 논의하겠습니다. 보조정리를 표현하는 멋진 방법은 {\it fpqc 위상이 정식 위상}보다 약하거나 fpqc 위상이 다음과 같다고 말하는 것이다. {\it 하위 표준}. 사이트 설정에서 이는 사이트, 절 \ref{sites-section-representable-sheaves}에서 논의된다.

\begin{remark} \label{etale-cohomology-remark-fpqc-finest} fpqc는 Zariski, 에탈, 매끄러운, 신토믹 및 fppf 위상보다 더 미세한다. 따라서 fpqc 위상에 대해 층 조건을 만족하는 모든 전층은 Zariski, 에탈, 매끄러운, 신토믹 및 fppf 사이트에서 층이 된다. 특히 표현 가능한 전층은 예에 대한 스킴의 에탈 사이트에서 층이 된다. \end{remark}

\begin{example}
\label{etale-cohomology-example-additive-group-sheaf}
$S$를 스킴으로 설정한다.
첨가제 그룹 스킴 $\mathbf{G}_{a, S} = \mathbf{A}^1_S$을 고려하라.
$S$ 이상, 참조
그룹형, 예 \ref{groupoids-example-additive-group}.
연관된 표현 가능 전층은 다음과 같이 제공된다.
$$
h_{\mathbf{G}_{a, S}}(T) =
\Mor_S(T, \mathbf{G}_{a, S}) =
\Gamma(T, \mathcal{O}_T).
$$
위의 내용을 통해 이제 이것이 집합의 전층라는 것을 알 수 있다.
층 조건 fpqc 위상에 대한 것이다. 한편으로는 분명히
환의 전층도 마찬가지이다. 따라서 우리는 이것을 함자로 생각할 수 있다.
$$
\begin{matrix}
\mathcal{O} : &
(\Sch/S)^{opp} &
\longrightarrow &
\textit{Rings} \\
&
T/S &
\longmapsto &
\Gamma(T, \mathcal{O}_T)
\end{matrix}
$$
이는 fpqc 위상에 대해 층 조건을 만족한다.
이에 따라 $\mathcal{O}$-가군 등의 개념이 있다.
등등.
\end{example}




\section{충실히 평탄 강하}
\label{etale-cohomology-section-fpqc-descent}

\noindent 이 절에서는 준일관성 가군에 대한 충실히 평탄 강하에 대해 논의한다. 보다 정확하게는 준연접 가군이 fpqc 피복에 대해 유효 하강을 만족한다는 것을 증명할 것이다.

\begin{definition}
\label{etale-cohomology-definition-descent-datum}
$\mathcal{U} = \{ t_i : T_i \to T\}_{i \in I}$를 다음의 가족으로 설정한다.
고정 대상이 있는 스킴 중 사상. {\it 하강 기준}
$\mathcal{U}$에 대한 준연접층은 컬렉션이다.
$((\mathcal{F}_i)_{i \in I}, (\varphi_{ij})_{i, j \in I})$ 여기서
\begin{enumerate}
\item $\mathcal{F}_i$는 $T_i$의 준연접층이고,
\item $\varphi_{ij} : \text{pr}_0^* \mathcal{F}_i \to
\text{pr}_1^* \mathcal{F}_j$는 다음의 동형사상이다. 가군
$T_i \times_T T_j$에서,
\end{enumerate}
{\it 코사이클 조건}는 다음 다이어그램을 보유한다.
$$
\xymatrix{
\text{pr}_0^*\mathcal{F}_i \ar[dr]_{\text{pr}_{02}^*\varphi_{ik}}
\ar[rr]^{\text{pr}_{01}^*\varphi_{ij}} & &
\text{pr}_1^*\mathcal{F}_j \ar[dl]^{\text{pr}_{12}^*\varphi_{jk}} \\
& \text{pr}_2^*\mathcal{F}_k
}
$$
$T_i \times_T T_j \times_T T_k$로 출퇴근하세요.
준일관성이 있는 경우 이 하강 데이텀을 {\it 유효}라고 한다.
층 $\mathcal{F}$ $T$ 및 $\mathcal{O}_{T_i}$-가군 동형사상
$\varphi_i : t_i^* \mathcal{F} \cong \mathcal{F}_i$ 호환 가능
지도 $\varphi_{ij}$, 즉
$$
\varphi_{ij} = \text{pr}_1^* (\varphi_j) \circ \text{pr}_0^* (\varphi_i)^{-1}.
$$
\end{definition}

\noindent 이 글과 다음 절에서는 다음 정리의 증명에 대한 일부 구성 요소와 관련 자료에 대해 논의한다.

\begin{theorem} \label{etale-cohomology-theorem-descent-quasi-coherent} $\mathcal{V} = \{T_i \to T\}_{i\in I}$가 fpqc 피복인 경우 $\mathcal{V}$에 대한 준연접층에 대한 모든 하강 데이터가 유효한다. \end{theorem}

\begin{proof} 하강, 명제 \ref{descent-proposition-fpqc-descent-quasi-coherent} 참조. \end{proof}

\noindent 즉, 준연접층의 파이버 범주는 fpqc 사이트에 스택이다. 정리의 증명은 2단계로 되어 있다. 첫 번째는 Zariski 피복의 경우 층의 표준 접착(층, 절 \ref{sheaves-section-glueing-sheaves} 참조)과 준 일관성의 국소성을 사용하여 이것이 쉽다는 것(또는 잘 알려져 있음)을 깨닫는 것이다. 두 번째 단계는 $\{\Spec(B) \to \Spec(A)\}$ 형식의 fpqc 피복의 경우이다. 여기서 $A \to B$는 충실히 평탄 환 맵이다. 이것은 우리가 지금 제시하는 대수학의 보조정리이다.

\medskip\noindent {\bf 가군의 하강.} $A \to B$가 환 지도인 경우 다음을 고려한다. 복합체 $$
(B/A)_\bullet : B \to B \otimes_A B \to B \otimes_A B \otimes_A B \to \ldots
$$ 여기서 $B$는 0, $B \otimes_A B$는 1 등이며 지도는 \begin{eqnarray*}
b & \mapsto & 1 \otimes b - b \otimes 1, \\
b_0 \otimes b_1 & \mapsto & 1 \otimes b_0 \otimes b_1 - b_0 \otimes 1 \otimes
b_1 + b_0 \otimes b_1 \otimes 1, \\
& \text{이하도 마찬가지이다.}
\end{eqnarray*}

\begin{lemma} \label{etale-cohomology-lemma-algebra-descent} $A \to B$가 충실히 평탄인 경우 복합체 $(B/A)_\bullet$는 양수로 정확하고 $H^0((B/A)_\bullet) = A$이다. \end{lemma}

\begin{proof} 하강, 보조정리 \ref{descent-lemma-ff-exact} 참조. \end{proof}

\noindent
Grothendieck은 이를 세 단계로 증명했다. 첫째, 그는 지도가 $A
\to B$는 절을 가지며, 복합체에 대한 명시적 호모토피를 구성한다.
$A$는 0 정도에서 0이 아닌 유일한 항이다. 둘째, 그는 증명하기 위해 그것을 관찰합니다
결과적으로, 충실히 평탄 기본 변경 $A \to A'$, replacing $B$ 와 함께 $B' = B \otimes_A A'$. 셋째, 그는 다음을 적용한다.
충실히 평탄 기본 변경 $A \to A' = B$ 및 비고 지도
$A' = B \to B' = B \otimes_A B$에는 자연적인 절이 있다.

\medskip\noindent 동일한 전략으로 다음 보조정리가 증명된다.

\begin{lemma} \label{etale-cohomology-lemma-descent-modules} $A \to B$가 충실히 평탄이고 $M$가 $A$-가군인 경우 복합체 $(B/A)_\bullet \otimes_A M$ 양의 각도로 정확하며 $H^0((B/A)_\bullet \otimes_A M) = M$이다. \end{lemma}

\begin{proof} 하강, 보조정리 \ref{descent-lemma-ff-exact} 참조. \end{proof}

\begin{definition}
\label{etale-cohomology-definition-descent-datum-modules}
$A \to B$는 환 지도이고 $N$는 $B$-가군라고 가정한다. {\it 하강 기준}
$N$에 대하여$A \to B$은동형사상
$\varphi : N \otimes_A B \cong B \otimes_A N$ 중 $B \otimes_A B$-가군 등
$B \otimes_A B \otimes_A B$-가군의 다이어그램은
$$
\xymatrix{
{N \otimes_A B \otimes_A B} \ar[dr]_{\varphi_{02}} \ar[rr]^{\varphi_{01}} & &
{B \otimes_A N \otimes_A B} \ar[dl]^{\varphi_{12}} \\
& {B \otimes_A B \otimes_A N}
}
$$
$\varphi_{01} = \varphi \otimes \text{id}_B$에서 통근하며 이와 유사하게
$\varphi_{12}$ 및 $\varphi_{02}$의 경우.
\end{definition}

\noindent
만약에$N' = B \otimes_A M$어떤 사람들에게는$A$-가군M, 그러면 정식 하강을 가집니다.
지도가 제공하는 데이텀
$$
\begin{matrix}
\varphi_\text{can}: & N' \otimes_A B & \to & B \otimes_A N' \\
& b_0 \otimes m \otimes b_1 & \mapsto & b_0 \otimes b_1 \otimes m.
\end{matrix}
$$

\begin{definition}
\label{etale-cohomology-definition-effective-modules}
하강 기준점 $(N, \varphi)$은 {\it 유효}라고 한다.
$A$-가군 $M$ $(N, \varphi) \cong (B \otimes_A M,
\varphi_\text{can})$, 하강 데이터의 동형사상라는 명백한 개념을 사용한다.
\end{definition}

\noindent 정리 \ref{etale-cohomology-theorem-descent-quasi-coherent}는 다음 결과의 결과이다.

\begin{theorem} \label{etale-cohomology-theorem-descent-modules} $A \to B$가 충실히 평탄인 경우 $A\to B$에 대한 하강 데이터가 유효한다. \end{theorem}

\begin{proof} 하강, 명제 \ref{descent-proposition-descent-module} 참조. 증명의 대체 보기는 하강, 비고 \ref{descent-remark-homotopy-equivalent-cosimplicial-algebras}를 참조하라. \end{proof}

\begin{remarks}
\label{etale-cohomology-remarks-theorem-modules-exactness}
가군의 하강 결과는 여러 가지 용도로 사용된다.
\begin{enumerate}
\item Čech 복합체의 완전성은 양수로 표시된다.
피복 $\{\Spec(B) \to \Spec(A)\}$ 여기서 $A \to B$는
충실히 평탄. 그러면 코호몰로지 중 일부 소멸이 제공된다.
\item $(N, \varphi)$가 충실한 하강 기준점인 경우
평면 지도 $A \to B$이면 해당 $A$-가군은 다음과 같이 제공된다.
$$
M = \Ker \left(
\begin{matrix}
N & \longrightarrow & B \otimes_A N \\
n & \longmapsto & 1 \otimes n - \varphi(n \otimes 1)
\end{matrix}
\right).
$$
보다
하강, 명제 \ref{descent-proposition-descent-module}.
\end{enumerate}
\end{remarks}




%9.17.09
\section{준연접층}
\label{etale-cohomology-section-quasi-coherent}

\noindent 가군의 하강을 준연접층 연구에 적용할 수 있다.

\begin{proposition}
\label{etale-cohomology-proposition-quasi-coherent-sheaf-fpqc}
누구에게나준연접층 $\mathcal{F}$~에$S$그만큼전층
$$
\begin{matrix}
\mathcal{F}^a : & \Sch/S & \to & \textit{Ab}\\
& (f: T \to S) & \mapsto & \Gamma(T, f^*\mathcal{F})
\end{matrix}
$$
은 층 조건을 만족하는 $\mathcal{O}$-가군이다.
fpqc 위상.
\end{proposition}

\begin{proof}
이는 다음에서 입증되었다.
하강, 보조정리 \ref{descent-lemma-sheaf-condition-holds}.
여기서는 증명을 표시한다. 에 설립된 바와 같이
보조정리 \ref{etale-cohomology-lemma-fpqc-sheaves},
층 성질을 확인하는 것으로 충분한다.
자리스키의 피복 및 충실히 평탄 사상 아핀 스킴. 그만큼
층 성질 Zariski에 대한 피복은 표준 스킴 이론이다.
$\Gamma(U, i^\ast \mathcal{F}) = \mathcal{F}(U)$ 때
$i : U \hookrightarrow S$은 열린 몰입이다.

\medskip\noindent
$\left\{\Spec(B)\to \Spec(A)\right\}$에 대해 $A\to B$를 충실히 사용
평평하고
$\mathcal{F}|_{\Spec(A)} = \widetilde{M}$
이는 다음 사실에 해당한다.
$M = H^0\left((B/A)_\bullet \otimes_A M \right)$ 즉,
\begin{align*}
0 \to M \to B \otimes_A M \to B \otimes_A B \otimes_A M
\end{align*}
에 의해 정확하다
보조정리 \ref{etale-cohomology-lemma-descent-modules}.
\end{proof}

\noindent
고리 모양의 사이트에는 준연접층라는 추상적 개념이 있다.
여기에 대해 간략하게 소개한다. 자세한 내용은 상담해 주세요
가군 사이트, 절 \ref{sites-modules-section-local}.
$\mathcal{C}$를 범주로 두고, $U$를 $\mathcal{C}$의 객체로 둡니다.
그런 다음 $\mathcal{C}/U$는 $U$에 대한 개체의 범주를 나타냅니다.
범주론, 예 \ref{categories-example-category-over-X}.
$\mathcal{C}$가 사이트이면 $\mathcal{C}/U$도 사이트이다.
$V/U$의 피복은 사상의 $\{V_i/U \to V/U\}$ 계열이다.
고정된 목표를 가진 $\mathcal{C}/U$의
$\{V_i \to V\}$은 $\mathcal{C}$의 피복이다. 게다가, 주어진 어떤
층 $\mathcal{F}$ $\mathcal{C}$ {\it 제한}
$\mathcal{F}|_{\mathcal{C}/U}$ (명확한 방식으로 정의됨)
층이기도 한다. 보다
사이트, 절 \ref{sites-section-localize}
자세한 내용은.

\begin{definition}
\label{etale-cohomology-definition-ringed-site}
$\mathcal{C}$를 {\it 고리형 사이트}, 즉 사이트로 지정한다.
환 $\mathcal{O}$의 층. $\mathcal{O}$-가군 $\mathcal{F}$의 층
$\mathcal{C}$는 {\it 준일관성}이라고 한다.
$U \in \Ob(\mathcal{C})$ 피복이 존재한다.
$\{U_i \to U\}_{i\in I}$의 $\mathcal{C}$ 제한은 다음과 같습니다.
$\mathcal{F}|_{\mathcal{C}/U_i}$는 다음의 여핵과 동형이다.
$\mathcal{O}$ 자유 $\mathcal{O}$-가군의 선형 맵
$$
\bigoplus\nolimits_{k \in K} \mathcal{O}|_{\mathcal{C}/U_i}
\longrightarrow
\bigoplus\nolimits_{l \in L} \mathcal{O}|_{\mathcal{C}/U_i}.
$$
$K$에 대한 직합은 전층과 연관된 층이다.
$V \mapsto \bigoplus_{k \in K} \mathcal{O}(V)$ 다른 쪽도 마찬가지이다.
\end{definition}

\noindent 위와 같이 일반적인 정의를 제공할 수 있으면 유용하지만 이 개념은 일반적으로 잘 작동하지 않습니다.

\begin{remark} \label{etale-cohomology-remark-final-object} $\mathcal{C}$에 최종 객체가 있는 경우(예: \ $S$)에 대해 정의의 조건을 확인하면 충분한다. $U = S$ 위의 진술에서. 사이트, 보조정리 \ref{sites-modules-lemma-local-final-object}의 가군을 참조하라. \end{remark}

\begin{theorem}[준연접층]의 메타 정리
\label{etale-cohomology-theorem-quasi-coherent}
$S$를 스킴으로 설정한다.
$\mathcal{C}$를 사이트로 설정한다. 가정
\begin{enumerate}
\item 기본 범주 $\mathcal{C}$는
$\Sch/S$의 전체 하위 카테고리,
\item 임의의 자리스키 피복 의 $T \in \Ob(\mathcal{C})$
$\mathcal{C}$의 피복으로 정제할 수 있다.
\item $S/S$는 $\mathcal{C}$의 객체이다.
\item $\mathcal{C}$의 피복마다 스킴의 fpqc 피복이다.
\end{enumerate}
그런 다음 전층 $\mathcal{O}$는 $\mathcal{C}$의 층이고
임의의 준일관성$\mathcal{O}$-가군~에$(\mathcal{C}, \mathcal{O})$
일부 준연접층의 경우 $\mathcal{F}^a$ 형식이다.
$\mathcal{F}$ - $S$.
\end{theorem}

\begin{proof} 몇 가지 공식적인 주장을 따르면 이는 정확히 정리 \ref{etale-cohomology-theorem-descent-quasi-coherent}이다. 자세한 내용은 생략했다. 강하에서는 명제 \ref{descent-proposition-equivalence-quasi-coherent} $S$의 큰 Zariski, fppf, 에탈, 부드럽고 신토믹한 사이트에 대한 정리의 더 정확한 버전을 증명한다. 사이트 중 $S$. \end{proof}

\noindent 즉, 스킴 $S$의 준 일관성 가군과 사이트의 준 일관성 $\mathcal{O}$-가군 사이에는 차이가 없다. $\mathcal{C}$는 정리와 같습니다. 크고 작은 사이트 $(\Sch/S)_{fppf}$, $S_\etale$ 등에 대한 더 정확한 설명은 강하, 절 \ref{descent-section-quasi-coherent-sheaves}, \ref{descent-section-quasi-coherent-cohomology} 및 \ref{descent-section-quasi-coherent-sheaves-bis}에서 찾을 수 있다. 이 장에서는 어떤 상황에서도 유지되는 결과를 편리하게 기술하기 위해 때때로 정리 \ref{etale-cohomology-theorem-quasi-coherent}''와 같이 ``사이트를 참조할 것이다.






\section{Čech 코호몰로지} \label{etale-cohomology-section-cech-cohomology}

\noindent
우리의 다음 목표는 하강 이론을 사용하여 다음을 보여주는 것이다.
$H^i(\mathcal{C}, \mathcal{F}^a) = H_{Zar}^i(S, \mathcal{F})$
모두를 위해준연접층 $\mathcal{F}$~에$S$, 그리고
어느사이트 $\mathcal{C}$에서와 같이정리 \ref{etale-cohomology-theorem-quasi-coherent}.
이를 위해 사이트에 Čech 코호몰로지를 소개한다.
\cite{ArtinTopologies} 및
코호몰로지 사이트, 절 \ref{sites-cohomology-section-cech},
\ref{sites-cohomology-section-cech-functor}
그리고 \ref{sites-cohomology-section-cech-cohomology-cohomology}
자세한 내용은

\begin{definition} \label{etale-cohomology-definition-cech-complex} $\mathcal{C}$를 범주로 둡니다. $\mathcal{U} = \{U_i \to U\}_{i \in I}$는 고정 대상이 있는 $\mathcal{C}$의 사상 계열이라고 가정한다. 모든 올곱 $U_{i_0} \times_U \ldots \times_U U_{i_p}$가 $\mathcal{C}$에 존재한다고 가정한다. $\mathcal{F} \in \textit{PAb}(\mathcal{C})$를 아벨 전층으로 둡니다. {\it čech 복합체} $\check{\mathcal{C}}^\bullet(\mathcal{U}, \mathcal{F})$를 $$
\prod_{i_0 \in I} \mathcal{F}(U_{i_0}) \to
\prod_{i_0, i_1 \in I} \mathcal{F}(U_{i_0} \times_U U_{i_1}) \to
\prod_{i_0, i_1, i_2 \in I}
\mathcal{F}(U_{i_0} \times_U U_{i_1} \times_U U_{i_2}) \to \ldots
$$로 정의한다. 여기서 첫 번째 항의 단위는 0이고 지도는 일반적인 것이다. {\it Čech 코호몰로지 군}은 $$
\check{H}^p(\mathcal{U}, \mathcal{F}) =
H^p(\check{\mathcal{C}}^\bullet(\mathcal{U}, \mathcal{F})).
$$ \end{definition}에 의해 정의된다.

\noindent 위의 정의에서는 인덱스 $(i_0, \ldots, i_p)$가 반복되도록 허용하는 것이 필수적이다.

\begin{lemma} \label{etale-cohomology-lemma-cech-presheaves} 정의 \ref{etale-cohomology-definition-cech-complex}와 같은 표기법 및 가정이다. 함자 $\check{\mathcal{C}}^\bullet(\mathcal{U}, -)$는 범주 $\textit{PAb}(\mathcal{C})$와 동일한다. \end{lemma}

\noindent 즉, $0\to \mathcal{F}_1\to \mathcal{F}_2\to \mathcal{F}_3\to 0$가 아벨 군의 전층의 짧은 완전열라면, $$
0 \to \check{\mathcal{C}}^\bullet\left(\mathcal{U}, \mathcal{F}_1\right)
\to\check{\mathcal{C}}^\bullet(\mathcal{U}, \mathcal{F}_2) \to
\check{\mathcal{C}}^\bullet(\mathcal{U}, \mathcal{F}_3)\to 0
$$는 복합체의 짧은 완전열이다.

\begin{proof}
이는 짧은 완전열의 정의에서 한 번에 이어집니다.
전층. 즉, abelian 전층의 범주는 다음의 범주이다.
함자 $\textit{Ab}$ 값이 있는 일부 범주에서 자동으로
abelian 범주: 시퀀스 $\mathcal{F}_1\to \mathcal{F}_2\to \mathcal{F}_3$
모두에 대해 $\textit{PAb}$ 필요충분조건은에서 정확한다.
$U \in \Ob(\mathcal{C})$, 시퀀스
$\mathcal{F}_1(U) \to \mathcal{F}_2(U) \to \mathcal{F}_3(U)$는 정확한다.
$\textit{Ab}$. 따라서 위의 복합체 은 단지 짧은 정확한 결과일 뿐이다.
각 학위의 순서. 또한보십시오
코호몰로지 사이트, 보조정리 \ref{sites-cohomology-lemma-cech-exact-presheaves}.
\end{proof}

\noindent 이는 $\check{H}^\bullet(\mathcal{U}, -)$이 $\delta$-함자임을 나타냅니다. 이제 우리는 그것이 보편적인 $\delta$-함자임을 보여줍니다. 따라서 우리는 그것이 {\it 제거 가능} 함자임을 보여줄 필요가 있다. 요네다 보조정리를 회상하는 것부터 시작한다.

\begin{lemma}[요네다 보조정리]
\label{etale-cohomology-lemma-yoneda-presheaf}
누구에게나전층 $\mathcal{F}$에범주 $\mathcal{C}$
그리고 $U \in \Ob(\mathcal{C})$ 함수형 동형사상이 있다.
$$
\Hom_{\textit{PSh}(\mathcal{C})}(h_U, \mathcal{F}) =
\mathcal{F}(U).
$$
\end{lemma}

\begin{proof} 범주, 보조정리 \ref{categories-lemma-yoneda}를 참조하라. \end{proof}

\noindent
집합 $E$가 주어지면 (이 절에서)
$\mathbf{Z}[E]$ $E$의 무료 아벨 군이다. 수식에서
$\mathbf{Z}[E] = \bigoplus_{e \in E} \mathbf{Z}$, 즉 $\mathbf{Z}[E]$는
무료$\mathbf{Z}$-가군의 요소로 구성된 기초를 가지고$E$.
이 표기법을 사용하여 우리는 자유 아벨 전층을
집합의 전층.

\begin{definition}
\label{etale-cohomology-definition-free-abelian-presheaf}
$\mathcal{C}$를 범주로 설정한다.
집합 $\mathcal{G}$의 전층이 주어지면 다음을 정의한다.
{\it 무료 아벨리안 전층 위의 $\mathcal{G}$},
$\mathbf{Z}_\mathcal{G}$로 표시되며, 규칙에 따라
$$
\mathbf{Z}_\mathcal{G}(U)
=
\mathbf{Z}[\mathcal{G}(U)]
$$
$U \in \Ob(\mathcal{C})$의 경우
$\mathcal{G}$의 제한 사상에 의해 제한 사상이 유도된다.
특별한 경우에는 $\mathcal{G} = h_U$ 간단히 씁니다.
$\mathbf{Z}_U = \mathbf{Z}_{h_U}$.
\end{definition}

\noindent 함자 $\mathcal{G} \mapsto \mathbf{Z}_\mathcal{G}$는 건망증이 있는 함자 $\textit{PAb}(\mathcal{C}) \to \textit{PSh}(\mathcal{C})$에게 왼쪽 수반이다. 따라서 모든 전층 $\mathcal{F}$에 대해 요네다 보조정리의 마지막 평등인 표준 동형사상 $$
\Hom_{\textit{PAb}(\mathcal{C})}(\mathbf{Z}_U, \mathcal{F})
=
\Hom_{\textit{PSh}(\mathcal{C})}(h_U, \mathcal{F})
=
\mathcal{F}(U)
$$가 있다. 특히 다음과 같은 결과를 얻었습니다.

\begin{lemma}
\label{etale-cohomology-lemma-cech-complex-describe}
정의 \ref{etale-cohomology-definition-cech-complex}와 같은 표기법 및 가정.
Čech 복합체 $\check{\mathcal{C}}^\bullet(\mathcal{U}, \mathcal{F})$
다음과 같이 명시적으로 설명할 수 있다.
\begin{eqnarray*}
\check{\mathcal{C}}^\bullet(\mathcal{U}, \mathcal{F})
& = &
\left(
\prod_{i_0 \in I}
\Hom_{\textit{PAb}(\mathcal{C})}(\mathbf{Z}_{U_{i_0}}, \mathcal{F}) \to
\prod_{i_0, i_1 \in I}
\Hom_{\textit{PAb}(\mathcal{C})}(
\mathbf{Z}_{U_{i_0} \times_U U_{i_1}}, \mathcal{F}) \to \ldots
\right) \\
& = &
\Hom_{\textit{PAb}(\mathcal{C})}\left(
\left(
\bigoplus_{i_0 \in I} \mathbf{Z}_{U_{i_0}} \leftarrow
\bigoplus_{i_0, i_1 \in I} \mathbf{Z}_{U_{i_0} \times_U U_{i_1}} \leftarrow
\ldots
\right), \mathcal{F}\right)
\end{eqnarray*}
\end{lemma}

\begin{proof} 이는 위의 공식에 따른 것이다. 사이트, 보조정리 \ref{sites-cohomology-lemma-cech-map-into}의 코호몰로지를 참조하라. \end{proof}

\noindent 이렇게 하면 마지막 $\Hom$의 첫 번째 인수에서 복합체만 연구하게 된다.

\begin{lemma}
\label{etale-cohomology-lemma-exact}
정의 \ref{etale-cohomology-definition-cech-complex}와 같은 표기법 및 가정.
아벨리안 전층의 복합체
\begin{align*}
\mathbf{Z}_\mathcal{U}^\bullet \quad : \quad
\bigoplus_{i_0 \in I} \mathbf{Z}_{U_{i_0}} \leftarrow
\bigoplus_{i_0, i_1 \in I} \mathbf{Z}_{U_{i_0} \times_U U_{i_1}} \leftarrow
\bigoplus_{i_0, i_1, i_2 \in I}
\mathbf{Z}_{U_{i_0} \times_U U_{i_1} \times_U U_{i_2}} \leftarrow
\ldots
\end{align*}
제외한 모든 각도에서 정확한다.$0$~에$\textit{PAb}(\mathcal{C})$.
\end{lemma}

\begin{proof}
$V\in \Ob(\mathcal{C})$에 대해 아벨 군의 복합체
$\mathbf{Z}_\mathcal{U}^\bullet(V)$는
$$
\begin{matrix}
\mathbf{Z}\left[
\coprod_{i_0\in I} \Mor_\mathcal{C}(V, U_{i_0})\right]
\leftarrow
\mathbf{Z}\left[
\coprod_{i_0, i_1 \in I}
\Mor_\mathcal{C}(V, U_{i_0} \times_U U_{i_1})\right]
\leftarrow \ldots = \\
\bigoplus_{\varphi : V \to U}
\left(
\mathbf{Z}\left[\coprod_{i_0 \in I} \Mor_\varphi(V, U_{i_0})\right]
\leftarrow
\mathbf{Z}\left[\coprod_{i_0, i_1\in I} \Mor_\varphi(V, U_{i_0}) \times
\Mor_\varphi(V, U_{i_1})\right]
\leftarrow
\ldots
\right)
\end{matrix}
$$
이다. 여기서
$$
\Mor_{\varphi}(V, U_i)
=
\{ V \to U_i \text{이고 } V \to U_i \to U \text{가 } \varphi \text{와 같은 것}\}.
$$
$S_\varphi = \coprod_{i\in I} \Mor_\varphi(V, U_i)$라 두자. 그러면
$$
\mathbf{Z}_\mathcal{U}^\bullet(V)
=
\bigoplus_{\varphi : V \to U}
\left(
\mathbf{Z}[S_\varphi] \leftarrow
\mathbf{Z}[S_\varphi \times S_\varphi] \leftarrow
\mathbf{Z}[S_\varphi \times S_\varphi \times S_\varphi] \leftarrow
\ldots \right).
$$
따라서 각 $S = S_\varphi$에 대해 복합체
\begin{align*}
\mathbf{Z}[S] \leftarrow
\mathbf{Z}[S \times S] \leftarrow
\mathbf{Z}[S \times S \times S] \leftarrow \ldots
\end{align*}
가 음의 차수에서 완전함을 보이면 충분하다. 이를 위해 명시적인 호모토피를 줄 수 있다.
$s\in S$를 고정하고 $K: n_{(s_0, \ldots, s_p)} \mapsto n_{(s, s_0,
\ldots, s_p)}.$로 정의하라. $K$가 다음 작용소의 영호모토피임은 쉽게 확인된다.
$$
\delta :
\eta_{(s_0, \ldots, s_p)}
\mapsto
\sum\nolimits_{i = 0}^p (-1)^p \eta_{(s_0, \ldots, \hat s_i, \ldots, s_p)}.
$$
자세한 내용은 사이트 코호몰로지, 보조정리
\ref{sites-cohomology-lemma-homology-complex}를 참조하라.
\end{proof}

\begin{lemma}
\label{etale-cohomology-lemma-hom-injective}
정의 \ref{etale-cohomology-definition-cech-complex}와 같은 표기법 및 가정.
$\mathcal{I}$가 $\textit{PAb}(\mathcal{C})$의 단사 대상인 경우,
그 다음에$\check H^p(\mathcal{U}, \mathcal{I}) = 0$모두를 위해$p > 0$.
\end{lemma}

\begin{proof} Čech 복합체는 함자 $\Hom_{\textit{PAb}(\mathcal{C})}(-, \mathcal{I}) $를 복합체 $
\mathbf{Z}^\bullet_\mathcal{U} $에 적용한 결과이다. 즉, $$
\check H^p(\mathcal{U}, \mathcal{I}) = H^p
(\Hom_{\textit{PAb}(\mathcal{C})} (\mathbf{Z}^\bullet_\mathcal{U},
\mathcal{I})).
$$ 그러나 우리는 방금 $\mathbf{Z}^\bullet_\mathcal{U}$는 음수로 정확하고 함자 $\Hom_{\textit{PAb}(\mathcal{C})}(-,
\mathcal{I})$는 정확하므로 $\Hom_{\textit{PAb}(\mathcal{C})}
(\mathbf{Z}^\bullet_\mathcal{U}, \mathcal{I})$는 양수로 정확한다. \end{proof}

\begin{theorem} \label{etale-cohomology-theorem-cech-derived} 정의 \ref{etale-cohomology-definition-cech-complex}와 같은 표기법 및 가정이다. $\textit{PAb}(\mathcal{C})$에서 함자 $\check{H}^p(\mathcal{U}, -)$는 $\check{H}^0(\mathcal{U}, -)$의 오른쪽 유도 함자이다. \end{theorem}

\begin{proof} 보조정리 \ref{etale-cohomology-lemma-hom-injective}에 의해 함자 $\check H^p(\mathcal{U}, -)$는 지울 수 있으므로 보편적인 $\delta$-함자이다. $\check H^0(\mathcal{U}, -)$의 오른쪽 유도 함자도 마찬가지이다. $0$ 정도에 동의하므로 보편적인 $\delta$-함자의 보편적인 성질에 따라 동의한다. 자세한 내용은 사이트, 보조정리 \ref{sites-cohomology-lemma-cech-cohomology-derived-presheaves}의 코호몰로지를 참조하라. \end{proof}

\begin{remark} \label{etale-cohomology-remark-presheaves-no-topology} 앞의 모든 진술은 전층에 관한 것이므로 아직 위상을 사용하지 않았습니다. \end{remark}




\section{체코-to-코호몰로지 스펙트럼 열} \label{etale-cohomology-section-cech-ss}

\noindent 이 스펙트럼 열은 층의 코호몰로지에 대한 기초 결과를 증명하는 데 필수적이다.

\begin{lemma} \label{etale-cohomology-lemma-forget-injectives} 건망증이 있는 함자 $\textit{Ab}(\mathcal{C})\to \textit{PAb}(\mathcal{C})$는 분사를 분사로 변환한다. \end{lemma}

\begin{proof} 이는 건망증이 있는 함자에 왼쪽 수반, 즉 층화가 있다는 사실을 사용하여 형식화한 것이다. 이는 완전 함자이다. 자세한 내용은 사이트, 보조정리 \ref{sites-cohomology-lemma-injective-abelian-sheaf-injective-presheaf}의 코호몰로지를 참조하라. \end{proof}

\begin{theorem}
\label{etale-cohomology-theorem-cech-ss}
$\mathcal{C}$를 사이트로 설정한다. 피복
$\mathcal{U} = \{U_i \to U\}_{i \in I}$ 중 $U \in \Ob(\mathcal{C})$
그리고 어떤아벨 층 $\mathcal{F}$~에$\mathcal{C}$
스펙트럼 열이 있습니다
$$
E_2^{p, q}
=
\check H^p(\mathcal{U}, \underline{H}^q(\mathcal{F}))
\Rightarrow
H^{p+q}(U, \mathcal{F}),
$$
여기서 $\underline{H}^q(\mathcal{F})$는 아벨 전층이다.
$V \mapsto H^q(V, \mathcal{F})$.
\end{theorem}

\begin{proof}
다음에서 단사 분해 $\mathcal{F}\to \mathcal{I}^\bullet$를 선택하세요.
$\textit{Ab}(\mathcal{C})$, 그리고 이중 복합체를 고려하세요.
$\check{\mathcal{C}}^\bullet(\mathcal{U}, \mathcal{I}^\bullet)$
그리고 지도
$$
\xymatrix{
\Gamma(U, I^\bullet) \ar[r] &
\check{\mathcal{C}}^\bullet(\mathcal{U}, \mathcal{I}^\bullet) \\
& \check{\mathcal{C}}^\bullet(\mathcal{U}, \mathcal{F}) \ar[u]
}
$$
여기서 수평 지도는 자연 지도이다.
$\Gamma(U, I^\bullet) \to
\check{\mathcal{C}}^0(\mathcal{U}, \mathcal{I}^\bullet)$
왼쪽 열에 수직 맵이 유도된다.
$\mathcal{F}\to \mathcal{I}^0$ 그리고 맨 아래 줄에 위치한다.
가정에 의해, $\mathcal{I}^\bullet$는 injectives의 복합체이다.
$\textit{Ab}(\mathcal{C})$, 따라서
보조정리 \ref{etale-cohomology-lemma-forget-injectives}, 그것은 다음의 분사의 복합체이다.
$\textit{PAb}(\mathcal{C})$. 따라서 이중 복합체의 행은 다음과 같습니다.
양의 각도로 정확함(보조정리 \ref{etale-cohomology-lemma-hom-injective})
그만큼핵$\check{\mathcal{C}}^0(\mathcal{U}, \mathcal{I}^\bullet)
\to \check{\mathcal{C}}^1(\mathcal{U}, \mathcal{I}^\bullet)$
는 다음과 같다
$\Gamma(U, \mathcal{I}^\bullet)$, 이후 $\mathcal{I}^\bullet$
층의 복합체이다. 특히, 전체 복합체 중 코호몰로지
표준이다
대역 절단 함자 $H^0(U, \mathcal{F})$의 코호몰로지.

\medskip\noindent
수직 방향의 경우,$q$일코호몰로지그룹$p$번째 열은
$$
\prod_{i_0, \ldots, i_p}
H^q(U_{i_0} \times_U \ldots \times_U U_{i_p}, \mathcal{F})
=
\prod_{i_0, \ldots, i_p}
\underline{H}^q(\mathcal{F})(U_{i_0} \times_U \ldots \times_U U_{i_p})
$$
항목 $E_1^{p, q}$에 있다. 따라서 이것은 표준 이중 복합체 스펙트럼이다.
순서이며 $E_2$ 페이지는 규정에 따릅니다. 자세한 내용은 다음을 참조하라.
코호몰로지 사이트에,
보조정리 \ref{sites-cohomology-lemma-cech-spectral-sequence}.
\end{proof}

\begin{remark} \label{etale-cohomology-remark-grothendieck-ss} 이것은 함자 $$
\textit{Ab}(\mathcal{C}) \longrightarrow
\textit{PAb}(\mathcal{C}) \xrightarrow{\check H^0} \textit{Ab}.
$$ \end{remark}의 구성에 대한 그로텐디크 스펙트럼 열이다.








\section{스킴} \label{etale-cohomology-section-big-small}의 크고 작은 사이트

\noindent $S$를 스킴으로 둡니다. $\tau$를 우리가 논의할 위상 중 하나로 가정하겠습니다. 따라서 $\tau \in \{fppf, syntomic, smooth, \etale, Zariski\}$. 물론 에탈 위상에만 관심이 있다면 전체적으로 $\tau = \etale$를 간단히 가정할 수 있다. 또한, 에탈 사상, 에탈 피복, 에탈 사이트에 대해서는 절 \ref{etale-cohomology-section-etale-site}부터 좀 더 자세히 다루겠습니다. 준연접층의 코호몰로지에 대한 논의를 진행하기 위해서는 큰 $\tau$-사이트를 도입하는 것이 편리하고, $\tau \in \{\etale, Zariski\}$의 경우에는 작은 $\tau$-사이트를 도입하는 것이 편리한다. $S$. 이를 위해 먼저 $\tau$-피복 개념을 도입한다.

\begin{definition} \label{etale-cohomology-definition-tau-covering} (위상, 정의 \ref{topologies-definition-fppf-covering}, \ref{topologies-definition-syntomic-covering}, \ref{topologies-definition-smooth-covering}, \ref{topologies-definition-etale-covering} 및 \ref{topologies-definition-zariski-covering}를 참조하라.) \hfil\break $\tau \in \{fppf, syntomic, smooth, \etale, Zariski\}$로 둡니다. 고정 대상이 있는 스킴 $\{f_i : T_i \to T\}_{i \in I}$의 사상 계열을 {\it $\tau$-피복}라고 한다. 필요충분조건은 각 $f_i$는 유한 표현, 구문론적, 부드러움, 에탈, 각각.\ 및 열린 몰입의 플랫이며 $\bigcup f_i(T_i) = T$가 있다. \end{definition}

\noindent 모든 $\tau$-피복의 클래스는 정의의 공리 (1), (2) 및 (3)를 충족한다. \ref{etale-cohomology-definition-site} (정의의 사이트), 위상, 기본정리 \ref{topologies-lemma-fppf}, \ref{topologies-lemma-syntomic}, \ref{topologies-lemma-smooth}, \ref{topologies-lemma-etale} 및 \ref{topologies-lemma-zariski}를 참조하라.

\medskip\noindent 우리가 작업하게 될 사이트를 소개하겠습니다. \cite{SGA4}에서 일어나는 일과는 반대로, 우리는 우주를 선택하고 싶지 않습니다. 대신에 우리는 ``부분 우주''(집합, 절 \ref{sets-section-sets-hierarchy}에서와 같이 적당히 큰 집합)를 선택하고 이 집합에 포함된 모든 스킴을 고려한다. 물론 우리는 우리가 가장 좋아하는 기지 스킴 $S$가 부분 우주에 포함되어 있는지 확인한다. 기본 범주를 선택한 후 $\tau$-피복 중 적절하게 큰 집합을 선택하여 이를 사이트로 바꿉니다. 자세한 내용은 스킴의 위상에 관한 장에 나와 있다. 선택에는 많은 자유가 있지만 결국 실제 선택은 $S$의 에탈(또는 기타) 코호몰로지에 영향을 미치지 않습니다(\cite{SGA4}에서 우주의 실제 선택은 결국 중요하지 않습니다). 더욱이, 접근하기 어려운 추기경(즉, 우주)을 사용하는 것을 좋아하는 독자는 자료가 작성되는 방식을 대신하여 그렇게 할 수 있다.

\begin{definition}
\label{etale-cohomology-definition-tau-site}
$S$를 스킴으로 설정한다.
$\tau \in \{fppf, syntomic, smooth, \etale, \linebreak[0] Zariski\}$로 놔두세요.
\begin{enumerate}
\item {\it 큰 $\tau$-사이트 의 $S$}는 사이트 중 하나이다.
$(\Sch/S)_\tau$ 위에서 설명한 대로 구성되었으며 자세한 내용은
위상, 정의
\ref{topologies-definition-big-small-fppf},
\ref{topologies-definition-big-small-syntomic},
\ref{topologies-definition-big-small-smooth},
\ref{topologies-definition-big-small-etale} 그리고
\ref{topologies-definition-big-small-Zariski}.
\item $\tau \in \{\etale, Zariski\}$인 경우
{\it 작은 $\tau$-사이트 $S$}
전체 하위 카테고리입니다$S_\tau$~의$(\Sch/S)_\tau$누구의 물건인가
~이다스킴 $T$~ 위에$S$누구의 구조사상 $T \to S$에탈이다,
각각.\ 열린 몰입. $S_\tau$의 피복은 피복이다.
$\{U_i \to U\}$ 의 $(\Sch/S)_\tau$
$U$은 $S_\tau$의 객체이다.
\end{enumerate}
\end{definition}

\noindent
사이트 $(\Sch/S)_\tau$의 기본 범주에는 합리적인
``클로저'' 성질, 즉 스킴 $T$가 주어지면 로컬로 닫힙니다.
$T$의 하위 구성표는 $(\Sch/S)_\tau$의 객체와 동형이다.
다른 폐쇄 성질은 다음과 같습니다: 스킴의 올곱 아래에서 폐쇄되었다.
셀 수 있는 분리된 결합을 취하고,
아핀 스킴이 주어지면 주어진 스킴에 대해 유한 유형 스킴을 취한다.
$\Spec(R)$ $R$의 몫을 완성하거나 현지화하거나 취할 수 있다.
범주 안에 머물면서 아이디얼에 의해
반면에, 예 임의의 분리된 결합의 경우
$(\Sch/S)_\tau$의 스킴이 당신을 그것 밖으로 데려갈 것이다.
또한 객체가 주어지면$T$~의$(\Sch/S)_\tau$존재할 것이다
$\tau$-피복 $\{T_i \to T\}_{i \in I}$ (다음과 같음)
정의 \ref{etale-cohomology-definition-tau-covering})
예에 대해 $(\Sch/S)_\tau$에서 피복이 아닌 것
스킴 $T_i$는 범주의 객체가 아니기 때문이다.
$(\Sch/S)_\tau$. 하지만 우리가 선택한 것은 사이트 $(\Sch/S)_\tau$
항상 존재하는 그런 존재야
 피복 $\{U_j \to T\}_{j \in J}$ 의 $(\Sch/S)_\tau$
피복 $\{T_i \to T\}_{i \in I}$, 참조
위상, 기본정리
\ref{topologies-lemma-fppf-induced},
\ref{topologies-lemma-syntomic-induced},
\ref{topologies-lemma-smooth-induced},
\ref{topologies-lemma-etale-induced} 그리고
\ref{topologies-lemma-zariski-induced}.
이 장에서는 이러한 문제를 대부분 무시할 것이다.

\medskip\noindent $\mathcal{F}$가 $(\Sch/S)_\tau$ 또는 $S_\tau$의 층인 경우, $$ 
H^p_\tau(U, \mathcal{F}), \text{ 특히 }
H^p_\tau(S, \mathcal{F})
$$ 코호몰로지 군을 나타냅니다. $\mathcal{F}$ 사이트의 객체 $U$에 대한 자세한 내용은 절 \ref{etale-cohomology-section-cohomology}를 참조하라. 따라서 $H^p_{fppf}(S, \mathcal{F})$, $H^p_{syntomic}(S, \mathcal{F})$, $H^p_{smooth}(S, \mathcal{F})$, $H^p_\etale(S, \mathcal{F})$ 및 $H^p_{Zar}(S, \mathcal{F})$가 있다. 마지막 두 개는 크거나 작은 에탈 또는 Zariski 사이트를 참조할 수 있으므로 잠재적으로 모호한다. 그러나 이러한 모호함은 다음 보조정리에 의해 무해한다.

\begin{lemma}
\label{etale-cohomology-lemma-compare-cohomology-big-small}
$\tau \in \{\etale, Zariski\}$로 놔두세요.
$\mathcal{F}$가 다음에 정의된 아벨 층인 경우
$(\Sch/S)_\tau$, 그런 다음
그만큼코호몰로지그룹$\mathcal{F}$~ 위에$S$에 동의하다코호몰로지
그룹$\mathcal{F}|_{S_\tau}$~ 위에$S$.
\end{lemma}

\begin{proof} 위상에 따라 정리 \ref{topologies-lemma-at-the-bottom} 및 \ref{topologies-lemma-at-the-bottom-etale}는 함자 $S_\tau \to (\Sch/S)_\tau$는 사이트, 보조정리 \ref{sites-lemma-bigger-site}의 가설을 만족한다. 따라서 보조정리는 사이트, 보조정리 \ref{sites-cohomology-lemma-cohomology-bigger-site}의 코호몰로지를 따릅니다. \end{proof}

\noindent 크고 작은 에탈에 대한 층의 범주는 $S$의 사이트 아핀으로 구성된 $(\Sch/S)_\etale$ 또는 $S_\etale$의 전체 하위 범주에만 의존하며 하나만 고려하면 된다. 그들 사이의 표준 에탈 피복 (아래에 정의됨). 이로 인해 사이트 $(\textit{Aff}/S)_\etale$ 및 $S_{affine, \etale}$가 발생한다. 위상, 정의 \ref{topologies-definition-big-small-etale}를 참조하라. 비교 결과는 위상, 보조정리 \ref{topologies-lemma-affine-big-site-etale} 및 \ref{topologies-lemma-alternative}에서 입증되었다. 다음은 이 장에서 고려할 일부 위상의 표준 피복의 정의이다.

\begin{definition}
\label{etale-cohomology-definition-standard-tau}
(보다
위상, 정의
\ref{topologies-definition-standard-fppf},
\ref{topologies-definition-standard-syntomic},
\ref{topologies-definition-standard-smooth},
\ref{topologies-definition-standard-etale} 그리고
\ref{topologies-definition-standard-Zariski}.)
\hfil\break $\tau \in \{fppf, syntomic, smooth, \etale, Zariski\}$로 놔두세요.
$T$를 아핀 스킴으로 설정한다.
{\it 표준 $\tau$-피복}은 $T$를 향하는 다음 사상족이다.
$\{f_j : U_j \to T\}_{j = 1, \ldots, m}$ 각각의 $U_j$은 유사하며,
그리고 각 $f_j$는 평평하고 유한하게 표현된다.
표준 신토믹, 표준 부드러움, 에탈, 각각.\ 하나의 몰입
표준 주체는 $T$ 및 $T = \bigcup f_j(U_j)$에서 열립니다.
\end{definition}

\begin{lemma} \label{etale-cohomology-lemma-tau-affine} $\tau \in \{fppf, syntomic, smooth, \etale, Zariski\}$로 둡니다. 아핀 스킴의 모든 $\tau$-피복은 표준 $\tau$-피복에 의해 구체화될 수 있다. \end{lemma}

\begin{proof} 위상, 기본정리 \ref{topologies-lemma-fppf-affine}, \ref{topologies-lemma-syntomic-affine}, \ref{topologies-lemma-smooth-affine}, \ref{topologies-lemma-etale-affine} 및 \ref{topologies-lemma-zariski-affine}를 참조하라. \end{proof}

\noindent 완전성을 위해 우리는 우리가 고려하고 있는 코호몰로지 군의 부분 우주 선택에 따른 불변성을 진술하고 증명한다. 아래의 보조정리 \ref{etale-cohomology-lemma-etale-topos-independent-partial-universe}에서 작은 에탈 토포스의 불변성을 증명하겠습니다. 이 보조정리에 사용된 표기법과 용어는 위상, 절 \ref{topologies-section-change-alpha}를 참조한다.

\begin{lemma}
\label{etale-cohomology-lemma-cohomology-enlarge-partial-universe}
$\tau \in \{fppf, syntomic, smooth, \etale, Zariski\}$로 놔두세요.
$S$를 스킴으로 설정한다.
$(\Sch/S)_\tau$와 $(\Sch'/S)_\tau$를 2개로 설정한다.
큰$\tau$-사이트~의$S$, 첫 번째가 두 번째에 포함되어 있다고 가정한다.
이 경우
\begin{enumerate}
\item 임의의 아벨 층 $\mathcal{F}'$ 중 $(\Sch'/S)_\tau$ 위에 정의된 것과
어떤 물건이라도$U$~의$(\Sch/S)_\tau$우리는
$$
H^p_\tau(U, \mathcal{F}'|_{(\Sch/S)_\tau}) =
H^p_\tau(U, \mathcal{F}')
$$
즉,코호몰로지~의$\mathcal{F}'$~ 위에$U$더 큰 것으로 계산사이트
더 작은 사이트로 제한된 $\mathcal{F}'$의 코호몰로지에 동의한다.
$U$ 이상.
\item 임의의 아벨 층 $\mathcal{F}$ 중 $(\Sch/S)_\tau$ 위의 것에 대하여
아벨 층 $\mathcal{F}'$ $(\Sch/S)_\tau'$에 대한 제한이 있다.
$(\Sch/S)_\tau$는 $\mathcal{F}$와 동형이다.
\end{enumerate}
\end{lemma}

\begin{proof} 위상에 따르면 보조정리 \ref{topologies-lemma-change-alpha} 포함 함자 $(\Sch/S)_\tau \to (\Sch'/S)_\tau$는 사이트, 보조정리 \ref{sites-lemma-bigger-site}의 가정을 충족한다. 이는 (2) 및 (1)가 사이트, 보조정리 \ref{sites-cohomology-lemma-cohomology-bigger-site}의 코호몰로지에서 따른다는 것을 의미한다. \end{proof}




\section{에탈 토포스} \label{etale-cohomology-section-etale-topos}

\noindent
 {\it 토포스}는 사이트에 있는 집합의 층의 범주이다.
사이트, 정의 \ref{sites-definition-topos}. 그러므로 그것은 관례이다
``에탈 토포스 의 스킴''라는 문구를 사용하여 참조
층의 범주는 스킴의 작은 에탈 사이트 위에 있다.
다음은 정식 정의이다.

\begin{definition}
\label{etale-cohomology-definition-etale-topos}
$S$를 스킴으로 설정한다.
\begin{enumerate}
\item {\it 에탈 토포스} 또는 {\it 작은 에탈 토포스}
의 $S$는 범주 $\Sh(S_\etale)$ 의 층 의 집합이다.
$S$의 작은 에탈 사이트.
\item {\it 자리스키 토포스} 또는 {\it 작은 자리스키 토포스}
의 $S$는 범주 $\Sh(S_{Zar})$ 의 층 의 집합이다.
$S$의 작은 자리스키 사이트.
\item $\tau \in \{fppf, syntomic, smooth, \etale, Zariski\}$의 경우
{\it 큰 $\tau$-토포스}는 층의 집합의 범주이다.
큰 $\tau$-토포스 $S$.
\end{enumerate}
\end{definition}

\noindent $S$의 작은 Zariski 토포스는 단순히 $S$의 기본 위상 공간에 있는 집합의 층의 범주이다. 위상을 참조하라. 보조정리 \ref{topologies-lemma-Zariski-usual}. 작은 에탈 토포스는 작은 에탈 사이트의 구성에서 이루어진 선택에 의존하지 않는 반면, 일반적으로 큰 토포스는 그러한 선택에 의존한다.

\medskip\noindent 크고 작은 에탈 토포스는 아핀으로 구성된 $(\Sch/S)_\etale$ 또는 $S_\etale$의 전체 하위 범주에만 의존하는 것으로 나타났습니다. 위상, 기본 정리 \ref{topologies-lemma-affine-big-site-etale} 및 \ref{topologies-lemma-alternative}. 다음 보조정리의 증명에서 예에 이를 사용한다.

\begin{lemma} \label{etale-cohomology-lemma-etale-topos-independent-partial-universe} $S$를 스킴으로 둡니다. $S$의 에탈 토포스는 정의 \ref{etale-cohomology-definition-tau-site}의 작은 에탈 사이트의 구성과 독립적입니다(정준적 동등성까지). \end{lemma}

\begin{proof} 두 개의 큰 에탈 사이트 $\Sch_\etale$ 및 $\Sch_\etale'$가 $S$를 포함하는 경우, $\Sh(S_\etale) \cong \Sh(S_\etale')$를 명확한 표기법으로 표시해야 한다. 위상에 따라 보조정리 \ref{topologies-lemma-contained-in} $\Sch_\etale \subset \Sch_\etale'$라고 가정할 수 있다. 집합에 의해, 보조정리 \ref{sets-lemma-what-is-in-it} $S$에 대한 모든 아핀 스킴의 에탈은 $\Sch_\etale$ 및 $\Sch_\etale'$의 객체와 동형이다. 따라서 유도된 함자 $S_{affine, \etale} \to S_{affine, \etale}'$는 등가물이다. 더욱이, 이 함자와 준역전 맵 모두 표준 에탈 피복을 표준 에탈 피복으로 변환하는 것이 분명한다. 따라서 위상의 결과는 보조정리 \ref{topologies-lemma-alternative}이다. \end{proof}





\section{코호몰로지 의 준연접층} \label{etale-cohomology-section-cohomology-quasi-coherent} %9.22.09

\noindent
우리는 간단한 보조정리로 시작한다.
정해진). 표준 피복의 Čech 복합체는 다음과 같습니다.
형식의 fpqc 피복의 Čech 복합체와 같습니다.
$\{\Spec(B) \to \Spec(A)\}$와 $A \to B$ 충실히 평탄.

\begin{lemma}
\label{etale-cohomology-lemma-cech-complex}
$\tau \in \{fppf, syntomic, smooth, \etale, Zariski\}$로 놔두세요.
$S$를 스킴으로 설정한다.
$\mathcal{F}$를 $(\Sch/S)_\tau$에서 아벨 층으로 설정하거나
$S_\tau$ $\tau = \etale$의 경우,
$\mathcal{U} = \{U_i \to U\}_{i \in I}$
이 사이트의 표준 $\tau$-피복이 된다.
$V = \coprod_{i \in I} U_i$로 놔두세요. 그 다음에
\begin{enumerate}
\item $V$는 아핀 스킴이다.
\item $\mathcal{V} = \{V \to U\}$는 fpqc 피복이다.
그리고 또한$\tau$-피복~하지 않는 한$\tau = Zariski$,
\item Čech 복합체
$\check{\mathcal{C}}^\bullet (\mathcal{U}, \mathcal{F})$ 그리고
$\check{\mathcal{C}}^\bullet (\mathcal{V}, \mathcal{F})$ 동의한다.
\end{enumerate}
\end{lemma}

\begin{proof} 표준 $\tau$-피복의 정의는 위상, 정의 \ref{topologies-definition-standard-Zariski}, \ref{topologies-definition-standard-etale}, \ref{topologies-definition-standard-smooth}, \ref{topologies-definition-standard-syntomic} 및 \ref{topologies-definition-standard-fppf}. 정의에 의해 스킴 $U_i$ 각각은 유사하고 $I$는 유한한 집합이다. 따라서 $V$는 아핀 스킴이다. $V \to U$는 단순하고 전사적이므로 $\mathcal{V}$는 fpqc 피복이다. 예 \ref{etale-cohomology-example-fpqc-coverings}를 참조하라. Zariski 사례를 제외하고 피복 $\mathcal{V}$는 $\tau$-피복이기도 한다. 위상, 정의 \ref{topologies-definition-etale-covering}, \ref{topologies-definition-smooth-covering}, \ref{topologies-definition-syntomic-covering}, 및 \ref{topologies-definition-fppf-covering}.

\medskip\noindent $\mathcal{U}$는 $\mathcal{V}$를 개선한 것이므로 Čech의 지도가 있다. 복합체 $\check{\mathcal{C}}^\bullet (\mathcal{V}, \mathcal{F}) \to
\check{\mathcal{C}}^\bullet (\mathcal{U}, \mathcal{F})$, 사이트의 코호몰로지를 참조하라. 방정식(\ref{sites-cohomology-equation-map-cech-complexes}). 다음으로, $T = \coprod_{j \in J} T_j$가 $\mathcal{F}$가 정의된 사이트에서 스킴의 분리된 결합인 경우 고정 대상 $\{T_j \to T\}_{j \in J}$을 갖는 사상 계열은 Zariski 피복이며, 그래서 \begin{equation}
\label{etale-cohomology-equation-sheaf-coprod}
\mathcal{F}(T) =
\mathcal{F}(\coprod\nolimits_{j \in J} T_j) =
\prod\nolimits_{j \in J} \mathcal{F}(T_j)
\end{equation} 에 의해 $\mathcal{F}$의 층 조건. 이는 위의 Čech 복합체 지도가 $$
V \times_U \ldots \times_U V
=
\coprod\nolimits_{i_0, \ldots i_p} U_{i_0} \times_U \ldots \times_U U_{i_p}
$$가 스킴이므로 각 등급에서 동형사상임을 의미한다. \end{proof}

\noindent 일반 전층에 대해 동등성(\ref{etale-cohomology-equation-sheaf-coprod})은 거짓이다. 층의 경우에도 사이트를 유지하지 않습니다. 왜냐하면 부산물이 피복으로 이어지지 않을 수 있고 분리되지 않을 수도 있기 때문이다. 그러나 그것은 모든 일반적인 것(적어도 우리가 공부할 모든 것)에 적용된다.

\begin{remark}
\label{etale-cohomology-remark-refinement}
보조정리 \ref{etale-cohomology-lemma-cech-complex}의 문에서 피복 $\mathcal{U}$
$\mathcal{V}$의 개선이지만 그 반대는 아닙니다. 피복
$\{V \to U\}$ 형식은 해당 항목의 초기 하위 범주를 형성하지 않습니다.
범주의 모든 피복의 $U$. 그래도 아직은 사실이다
우리는 Čech 코호몰로지 $\check H^n(U, \mathcal{F})$를 계산할 수 있다.
의 범주 반대에 대해 여극한으로 정의된다.
피복 $\mathcal{U}$ 의 $U$ Čech 코호몰로지 군 중
$\mathcal{F}$ $\mathcal{U}$에 대하여) 피복의 관점에서
$\{V \to U\}$. 우리는 정확한 보조정리를 공식화할 것입니다 (층에서만 작동합니다)
필요할 경우 여기에 추가하세요.
\end{remark}

\begin{lemma}[코호몰로지]의 지역
\label{etale-cohomology-lemma-locality-cohomology}
$\mathcal{C}$를 사이트, $\mathcal{F}$ 아벨 층으로 $\mathcal{C}$로 설정한다.
$U$ $\mathcal{C}$, $p > 0$ 정수 및 $\xi \in
H^p(U, \mathcal{F})$. 그러면 피복이 존재한다.
$\mathcal{U} = \{U_i \to U\}_{i \in I}$ 의 $U$ 에서 $\mathcal{C}$
그렇게$\xi |_{U_i} = 0$모두를 위해$i \in I$.
\end{lemma}

\begin{proof}
단사 분해 $\mathcal{F} \to \mathcal{I}^\bullet$을 선택한다. 그 다음에
$\xi$는 자전거 $\tilde{\xi} \in \mathcal{I}^p(U)$로 표현된다.
$d^p(\tilde{\xi}) = 0$로. 가정에 의해 시퀀스
$\mathcal{I}^{p - 1} \to \mathcal{I}^p \to \mathcal{I}^{p + 1}$ 정확히는
$\textit{Ab}(\mathcal{C})$, 이는 피복이 있음을 의미한다.
$\mathcal{U} = \{U_i \to U\}_{i \in I}$ 이렇게
$\tilde{\xi}|_{U_i} = d^{p - 1}(\xi_i)$ 일부에게는
$\xi_i \in \mathcal{I}^{p-1}(U_i)$. 부터
코호몰로지 클래스 $\xi|_{U_i}$는 코사이클로 표현된다.
공동경계인 $\tilde{\xi}|_{U_i}$는 사라진다.
자세한 내용은 다음을 참조하라.
코호몰로지 사이트에,
보조정리 \ref{sites-cohomology-lemma-kill-cohomology-class-on-covering}.
\end{proof}

\begin{theorem}
\label{etale-cohomology-theorem-zariski-fpqc-quasi-coherent}
$S$는 스킴이고 $\mathcal{F}$는 준일관적인 $\mathcal{O}_S$-가군라고 가정한다.
$\mathcal{C}$를 $(\Sch/S)_\tau$로 설정한다.
$\tau \in \{fppf, syntomic, smooth, \etale, Zariski\}$ 또는
$S_\etale$. 그 다음에
$$
H^p(S, \mathcal{F}) = H^p_\tau(S, \mathcal{F}^a)
$$
모든 $p \geq 0$에 대해 여기서
\begin{enumerate}
\item 왼쪽은 층의 일반적인 코호몰로지를 나타냅니다.
$\mathcal{F}$기본 위상 공간에 대해스킴 $S$, 그리고
\item 오른쪽은 코호몰로지를 나타냅니다.
아벨 층 $\mathcal{F}^a$의 (참조
명제 \ref{etale-cohomology-proposition-quasi-coherent-sheaf-fpqc})
사이트 $\mathcal{C}$에 있다.
\end{enumerate}
\end{theorem}

\begin{proof}
우리는 그것을 보여줄 것입니다
$H^p(U, f^*\mathcal{F}) = H^p_\tau(U, \mathcal{F}^a)$
어떤 대상에 대해서$f : U \to S$~의사이트 $\mathcal{C}$.
결과는 층 성질에 의한 $p = 0$에 대해 참이다.

\medskip\noindent
$U$이 유사하다고 가정한다. 그러면 우리는 그것을 증명하고 싶다.
$H^p_\tau(U, \mathcal{F}^a) = 0$모두를 위해$p > 0$. 저희는 인덕션을 이용해서$p$.
\begin{enumerate}
\item[p = 1]
$\xi \in H^1_\tau(U, \mathcal{F}^a)$을 선택한다.
보조정리 \ref{etale-cohomology-lemma-locality-cohomology}에 의해,
fpqc가 존재합니다 피복 $\mathcal{U} = \{U_i \to U\}_{i \in I}$
그렇게$\xi|_{U_i} = 0$모두를 위해$i \in I$. 정제까지
$\mathcal{U}$, $\mathcal{U}$이 표준이라고 가정할 수 있다.
$\tau$-피복. 스펙트럼 열 적용
정리 \ref{etale-cohomology-theorem-cech-ss},
$\xi$가 코호몰로지 클래스에서 나온 것을 알 수 있다.
$\check \xi \in \check H^1(\mathcal{U}, \mathcal{F}^a)$.
피복 $\mathcal{V} = \{\coprod_{i\in I} U_i \to U\}$를 고려해보세요. 에 의해
보조정리 \ref{etale-cohomology-lemma-cech-complex},
$\check H^\bullet(\mathcal{U}, \mathcal{F}^a) =
\check H^\bullet(\mathcal{V}, \mathcal{F}^a)$.
반면에 $\mathcal{V}$는 피복 형식이므로
$\{\Spec(B) \to \Spec(A)\}$ 및 $f^*\mathcal{F} = \widetilde{M}$
일부 $A$-가군 $M$의 경우 체코어 복합체가 표시된다.
$\check{\mathcal{C}}^\bullet(\mathcal{V}, \mathcal{F})$
다름 아닌 복합체 $(B/A)_\bullet \otimes_A M$이다.
이제 보조정리 \ref{etale-cohomology-lemma-descent-modules}에 의해
$H^p((B/A)_\bullet \otimes_A M) = 0$는 $p > 0$이므로 $\check \xi = 0$
그리고 $\xi = 0$.
\item[p > 1]
$\xi \in H^p_\tau(U, \mathcal{F}^a)$을 선택한다. 에 의해
보조정리 \ref{etale-cohomology-lemma-locality-cohomology},
fpqc가 존재합니다 피복 $\mathcal{U} = \{U_i \to U\}_{i \in I}$
그렇게$\xi|_{U_i} = 0$모두를 위해$i \in I$. 정제까지
$\mathcal{U}$, $\mathcal{U}$이 표준이라고 가정할 수 있다.
$\tau$-피복. 우리는 다음의 스펙트럼 열을 적용한다.
정리 \ref{etale-cohomology-theorem-cech-ss}.
교차점 $U_{i_0} \times_U \ldots \times_U U_{i_p}$을 관찰하세요.
유사하므로 귀납 가설에 따라 코호몰로지 군
$$
E_2^{p, q} = \check H^p(\mathcal{U}, \underline{H}^q(\mathcal{F}^a))
$$
모든 $0 < q < p$에 대해 사라집니다. $\xi$는 다음에서 와야 한다.
$\check \xi \in \check H^p(\mathcal{U}, \mathcal{F}^a)$. 교체
$\mathcal{U}$ 피복 $\mathcal{V}$에는 사상 하나만 포함된다.
보조정리 \ref{etale-cohomology-lemma-descent-modules}를 다시 사용하여
Čech 코호몰로지 클래스 $\check \xi$는 0이어야 한다는 것을 알 수 있다.
그러므로 $\xi = 0$.
\end{enumerate}
다음으로 $U$이 분리되어 있다고 가정한다. 개방형 아핀 선택 피복
$U = \bigcup_{i \in I} U_i$ 중 $U$이다. 가족
$\mathcal{U} = \{U_i \to U\}_{i \in I}$는 fpqc 피복이다.
그리고 모든 교차로
$U_{i_0} \times_U \ldots \times_U U_{i_p}$ 유사하다
$U$이 분리되어 있기 때문이다. 따라서 스펙트럼 열의 모든 행은
정리 \ref{etale-cohomology-theorem-cech-ss}
0번째 행을 제외하고는 0이다. 그러므로
$$
H^p_\tau(U, \mathcal{F}^a) =
\check H^p(\mathcal{U}, \mathcal{F}^a) =
\check H^p(\mathcal{U}, \mathcal{F}) = H^p(U, \mathcal{F})
$$
마지막 동등성은 표준 스킴 이론에서 비롯된다.
코호몰로지 중 스킴, 보조정리
\ref{coherent-lemma-cech-cohomology-quasi-coherent}.

\medskip\noindent 일반적인 경우는 기술적이며 (여기에 제공된 증명을 확장하려면) 스펙트럼 열 맵에 대한 논의가 필요하므로 다루지 않겠습니다. 이는 명제 \ref{descent-proposition-same-cohomology-quasi-coherent} (증명이 약간 다른 접근 방식을 취함)를 사이트, 보조정리 \ref{sites-cohomology-lemma-cohomology-of-open}에서 코호몰로지와 결합한 강하에서 따릅니다. \end{proof}

\begin{remark}
\label{etale-cohomology-remark-right-derived-global-sections}
정리 \ref{etale-cohomology-theorem-zariski-fpqc-quasi-coherent}에 대한 설명이다.
$S$는 범주 $\mathcal{C}$의 최종 객체이므로 코호몰로지
오른쪽에 있는 그룹은 단지 오른쪽 유도 함자일 뿐이다.
대역 절단 함자. 실제로 증명은 다음을 보여줍니다.
$H^p(U, f^*\mathcal{F}) = H^p_\tau(U, \mathcal{F}^a)$
어떤 대상에 대해서$f : U \to S$~의사이트 $\mathcal{C}$.
\end{remark}





\section{층} \label{etale-cohomology-section-examples-sheaves}의 예

\noindent $S$ 및 $\tau$는 절 \ref{etale-cohomology-section-big-small}와 같습니다. 우리는 이미 표현 가능한 모든 전층이 $(\Sch/S)_\tau$ 또는 $S_\tau$의 층라는 것을 확인했다. 보조정리 \ref{etale-cohomology-lemma-representable-sheaf-fpqc} 및 비고 \ref{etale-cohomology-remark-fpqc-finest}를 참조하라. 다음은 몇 가지 특별한 경우이다.

\begin{definition}
\label{etale-cohomology-definition-additive-sheaf}
사이트 $(\Sch/S)_\tau$ 또는 $S_\tau$ 중 하나에서
절 \ref{etale-cohomology-section-big-small}.
\begin{enumerate}
\item 층 $T \mapsto \Gamma(T, \mathcal{O}_T)$는 다음과 같이 표시된다.
$\mathcal{O}_S$ 또는 $\mathbf{G}_a$ 또는 $\mathbf{G}_{a, S}$
베이스 스킴을 표시하고 싶다.
\item 마찬가지로 층
$T \mapsto \Gamma(T, \mathcal{O}^*_T)$은 $\mathcal{O}_S^*$로 표시된다.
원하는 경우 $\mathbf{G}_m$ 또는 $\mathbf{G}_{m, S}$
베이스 스킴을 나타냅니다.
\item {\it 상수층} $\underline{\mathbf{Z}/n\mathbf{Z}}$
사이트는 상수 전층의 층화이다.
$U \mapsto \mathbf{Z}/n\mathbf{Z}$.
\end{enumerate}
\end{definition}

\noindent
첫 번째는 층이다.
정리 \ref{etale-cohomology-theorem-quasi-coherent}
예의 경우. 두 번째는 쉽게 볼 수 있는 첫 번째의 하위 전층이다.
층 그 자체가 되는 것이다. 세 번째는 정의에 의한 층이다.
층 각각을 표현할 수 있다는 점에 유의하세요.
스킴 $\mathbf{G}_{a, S}$에 의한 첫 번째와 두 번째
$\mathbf{G}_{m, S}$, 참조
그룹형, 절 \ref{groupoids-section-group-schemes}.
유한 에탈 그룹에 의한 세 번째 스킴 $\mathbf{Z}/n\mathbf{Z}_S$
때때로 $(\mathbf{Z}/n\mathbf{Z})_S$로 표시됨
그것은 단지$n$사본$S$부여된
$S$ 위에 명백한 그룹 스킴 구조가 있는 경우
그룹체, 예 \ref{groupoids-example-constant-group}
그리고 다음은 비고이다.

\begin{remark} \label{etale-cohomology-remark-constant-locally-constant-maps} $G$를 추상 그룹으로 둡니다. 사이트 $(\Sch/S)_\tau$ 또는 $S_\tau$ 중 절 \ref{etale-cohomology-section-big-small}에 연결된 상수 전층의 층화 $\underline{G}$ $G$ 에서 {\it Zariski 위상} 의 사이트는 이미 $$
\Gamma(U, \underline{G}) =
\{\text{자리스키 국소 상수 맵}U \to G\}
$$를 제공한다. 이 Zariski 층은 그룹으로 대표할 수 있다. 스킴 $G_S$ Groupoids에 따르면 예 \ref{groupoids-example-constant-group}. 보조정리 \ref{etale-cohomology-lemma-representable-sheaf-fpqc}에 의해 표현 가능한 모든 전층은 $\tau$-위상에 대한 층 조건도 만족하므로 자리스키 위의 층화 $\underline{G}$도 $\tau$-층화이다. \end{remark}

\begin{definition}
\label{etale-cohomology-definition-structure-sheaf}
허락하다$S$가 되다스킴. 그만큼{\it 구조층}~의$S$은층~의환
$\mathcal{O}_S$
사이트 $S_{Zar}$, $S_\etale$ 또는 $(\Sch/S)_\tau$ 중 하나
위에서 논의했다.
\end{definition}

\noindent 우리가 작업 중인 사이트에 대해 혼동이 있을 수 있는 경우 인덱스를 사용하여 이를 표시한다. 예의 경우 $\mathcal{O}_{S_\etale}$를 사용하여 $S$의 작은 에탈 사이트에 대해 작업하고 있다는 사실을 강조할 수 있다.

\begin{remark}
\label{etale-cohomology-remark-special-case-fpqc-cohomology-quasi-coherent}
위에 소개된 용어에서 특별한 경우는 다음과 같습니다.
정리 \ref{etale-cohomology-theorem-zariski-fpqc-quasi-coherent}
~이다
$$
H_{fppf}^p(X, \mathbf{G}_a) =
H_\etale^p(X, \mathbf{G}_a) =
H_{Zar}^p(X, \mathbf{G}_a) =
H^p(X, \mathcal{O}_X)
$$
모든 $p \geq 0$에 대해. 게다가, 우리는 표기법을 사용할 수 있습니다
$H^p_{fppf}(X, \mathcal{O}_X)$는 코호몰로지를 나타냅니다.
구조층 $X$의 큰 fppf 사이트에 있다.
\end{remark}




\section{Picard 군}
\label{etale-cohomology-section-picard-groups}

\noindent 다음 정리는 ``힐베르트 90''라고도 한다.

\begin{theorem} \label{etale-cohomology-theorem-picard-group} 모든 스킴 $X$에 대해 정식 ID \begin{align*}
H_{fppf}^1(X, \mathbf{G}_m) & = H^1_{syntomic}(X, \mathbf{G}_m) \\
& = H^1_{smooth}(X, \mathbf{G}_m) \\
& = H_\etale^1(X, \mathbf{G}_m) \\
& = H^1_{Zar}(X, \mathbf{G}_m) \\
& = \Pic(X) \\
& = H^1(X, \mathcal{O}_X^*)
\end{align*} \end{theorem}

\begin{proof}
$\tau$를 고려된 위상 중 하나라고 가정한다.
절 \ref{etale-cohomology-section-big-small}.
에 의해
코호몰로지 사이트, 보조정리에
\ref{sites-cohomology-lemma-h1-invertible}
우리는 그것을 본다
$H^1_\tau(X, \mathbf{G}_m) =
H^1_\tau(X, \mathcal{O}_\tau^*) =
\Pic(\mathcal{O}_\tau)$
여기서 $\mathcal{O}_\tau$는 사이트의 구조층이다.
$(\Sch/X)_\tau$. 이제 반전 가능한 $\mathcal{O}_\tau$-가군
준일관적인 $\mathcal{O}_\tau$-가군이다.
정리 \ref{etale-cohomology-theorem-quasi-coherent} 또는 더 정확한 방법으로
하강, 명제 \ref{descent-proposition-equivalence-quasi-coherent}
$\Pic(\mathcal{O}_\tau) = \Pic(X)$이 보이다.
마지막 평등은 같은 방식으로 증명된다.
\end{proof}






\section{에탈 사이트} \label{etale-cohomology-section-etale-site}

\noindent 이 시점에서 스킴의 에탈 사이트를 더 자세히 탐색하기 시작한다. 첫 번째 단계로 에탈 사상의 개념에 대해 조금 논의한다.





\section{에탈 사상}
\label{etale-cohomology-section-etale-morphism}

\noindent
자세한 내용은 다음을 참조하라.
사상, 절 \ref{morphisms-section-etale}
형식적인 정의의 경우
에탈 사상, 절
\ref{etale-section-etale-morphisms},
\ref{etale-section-structure-etale-map},
\ref{etale-section-etale-smooth},
\ref{etale-section-topological-etale},
\ref{etale-section-functorial-etale} 그리고
\ref{etale-section-properties-permanence}
에탈 사상의 흥미로운 성질에 대한 설문 조사를 위해.

\medskip\noindent
대수학을 기억하세요$A$대수적으로 닫힌 것에 대해체 $k$~이다
{\it 유한 유형이고 미분의 가군인 경우 매끄럽게}
$\Omega_{A/k}$는 차원과 동일한 순위가 없는 유한 로컬이다.
 스킴 $X$ 위의 $k$는 로컬 유한인 경우 {\it 매끄러운} 위의 $k$이다.
유형과 각 아핀 오픈은 부드러운 $k$-대수학의 스펙트럼이다.
$k$가 대수적으로 닫혀 있지 않으면 $k$-대수 $A$는 다음과 같습니다.
매끄러운$k$-대수학 경우$A \otimes_k \overline{k}$매끄럽다
$\overline{k}$-대수학. 환 지도 $A \to B$는 다음과 같은 경우 매끄러워집니다.
평평하고 유한하게 제시되며 모든 소수에 대해 $\mathfrak p \subset A$
올 환 $\kappa(\mathfrak p) \otimes_A B$는 잔여물 위에 매끄러워집니다.
체 $\kappa(\mathfrak p)$. 보다 일반적으로 스킴의 사상은 다음과 같습니다.
{\it 매끄럽게} 평면적이고 국부적으로 유한하게 표현된 경우
기하학 올은 부드럽습니다.

\medskip\noindent 이러한 사실에 대해서는 사상, 절 \ref{morphisms-section-smooth}를 참조하라. 이를 이용하여 에탈 사상을 다음과 같이 정의할 수 있다.

\begin{definition}
\label{etale-cohomology-definition-etale-morphism}
에이사상~의스킴~이다{\it 에탈}상대적으로 원활하다면차원 0.
\end{definition}

\noindent 특히, 스킴 $X \to S$의 사상은 매끄럽고 $\Omega_{X/S} = 0$이면 에탈이다.

\begin{proposition} \label{etale-cohomology-proposition-etale-morphisms} 에탈 사상에 관하여 다음이 성립한다. \begin{enumerate} \item $k$를 체라 하자. 스킴의 사상 $U \to \Spec(k)$가 에탈일 필요충분조건은 $U \cong \coprod_{i \in I} \Spec(k_i)$로 쓸 수 있고, 각 $i \in I$에 대하여 $k_i$가 $k$의 유한 분리 확대인 체인 것이다. \item $\varphi : U \to S$를 스킴의 사상이라 하자. 다음 조건들은 동치이다. \begin{enumerate} \item $\varphi$는 에탈이다. \item $\varphi$는 국소 유한 표시이고 평탄하며, 모든 올이 에탈이다. \item $\varphi$는 평탄하고 비분기이며 국소 유한 표시이다. \end{enumerate} \item 환 준동형 $A \to B$가 에탈일 필요충분조건은 $B \cong A[x_1, \ldots, x_n]/(f_1, \ldots, f_n)$로 쓸 수 있고 $\Delta = \det \left( \frac{\partial f_i}{\partial x_j} \right)$가 $B$에서 가역인 것이다. \item 에탈 사상의 밑변환은 에탈이다. \item 에탈 사상의 합성은 에탈이다. \item 에탈 사상들의 올곱과 곱은 에탈이다. \item 에탈 사상의 상대차원은 0이다. \item $Y \to X$를 에탈 사상이라 하자. $X$가 기약이면(각각 정칙이면) $Y$도 그러하다. \item 에탈 사상은 열린 사상이다. \item $X \to S$와 $Y \to S$가 에탈이면 임의의 $S$-사상 $X \to Y$도 에탈이다. \end{enumerate} \end{proposition}

\begin{proof} 우리는 이전 장에서 이러한 사실과 그 이상을 증명했다. 다음은 참조 목록입니다: (1) 사상, 보조정리 \ref{morphisms-lemma-etale-over-field}. (2) 사상, 정리 \ref{morphisms-lemma-etale-flat-etale-fibres} 및 \ref{morphisms-lemma-flat-unramified-etale}. (3) 대수학, 보조정리 \ref{algebra-lemma-etale-standard-smooth}. (4) 사상, 보조정리 \ref{morphisms-lemma-base-change-etale}. (5) 사상, 보조정리 \ref{morphisms-lemma-composition-etale}. (6) (4) 및 (5)에서 형식적으로 따릅니다. (7) 사상, 정리 \ref{morphisms-lemma-etale-locally-quasi-finite} 및 \ref{morphisms-lemma-locally-quasi-finite-rel-dimension-0}. (8) 대수학, 정리 \ref{algebra-lemma-reduced-goes-up} 및 \ref{algebra-lemma-Rk-goes-up}를 참조하라. 에탈 사상, 절 \ref{etale-section-properties-permanence}에서 이러한 종류의 더 많은 결과도 참조하라. (9) 사상, 보조정리 \ref{morphisms-lemma-fppf-open} 및 \ref{morphisms-lemma-etale-flat}를 참조하라. (10) 사상, 보조정리 \ref{morphisms-lemma-etale-permanence}를 참조하라. \end{proof}

\begin{definition} \label{etale-cohomology-definition-standard-etale} 환 지도 $A \to B$는 {\it 표준 에탈} 만약 $B \cong \left(A[t]/(f)\right)_g$라고 한다. $f, g \in A[t]$, $f$ 모닉이 있고 $\text{d}f/\text{d}t$는 $B$에서 반전 가능한다. \end{definition}

\noindent
표준 에탈 환 지도는 에탈인 것이 사실이다. 즉, 가정해보자
$B = \left(A[t]/(f)\right)_g$와 $f, g \in A[t]$, $f$ 모닉과 함께,
및 $\text{d}f/\text{d}t$는 $B$에서 반전 가능한다. 그러면 $A[t]/(f)$는 유한이다.
무료$A$-가군모닉 다항식의 차수와 동일한 순위$f$.
따라서 $B$, 이 자유 대수학의 국소화는 유한하게 제시된다.
$A$ 위에는 평평한다. $B$인 증명을 끝내려면 에탈이면 충분한다.
올 환
$$
\kappa(\mathfrak p) \otimes_A B
\cong
\kappa(\mathfrak p) \otimes_A (A[t]/(f))_g
\cong
\kappa(\mathfrak p)[t, 1/\overline{g}]/(\overline{f})
$$
유한 분리 가능한 체 확장의 유한 곱이다.
여기 $\overline{f}, \overline{g} \in \kappa(\mathfrak p)[t]$는
$f$ 및 $g$의 이미지. 허락하다
$$
\overline{f} = \overline{f}_1 \ldots \overline{f}_a
\overline{f}_{a + 1}^{e_1} \ldots \overline{f}_{a + b}^{e_b}
$$
$\overline{f}$을 쌍별 구별의 거듭제곱으로 인수분해한다.
기약 모닉 팩터 $\overline{f}_i$와 $e_1, \ldots, e_b > 1$.
가정 $\text{d}\overline{f}/\text{d}t$에 의해 반전 가능
$\kappa(\mathfrak p)[t, 1/\overline{g}]$. 그러므로 우리는
적어도 모든 $\overline{f}_i$, $i > a$는 되돌릴 수 있다. 우리는 결론을 내린다
저것
$$
\kappa(\mathfrak p)[t, 1/\overline{g}]/(\overline{f})
\cong
\prod\nolimits_{i \in I} \kappa(\mathfrak p)[t]/(\overline{f}_i)
$$
여기서 $I \subset \{1, \ldots, a\}$는 인덱스 $i$의 하위 집합이다.
$\overline{f}_i$는 $\overline{g}$를 나누지 않습니다. 게다가, 그 이미지는
$\text{d}\overline{f}/\text{d}t$ 요인에
$\kappa(\mathfrak p)[t]/(\overline{f}_i)$는 분명히
단위 시간 $\text{d}\overline{f}_i/\text{d}t$. 따라서 우리는 다음과 같이 결론을 내립니다.
$\kappa_i = \kappa(\mathfrak p)[t]/(\overline{f}_i)$는 유한한 체이다.
하나의 요소에 의해 생성된 $\kappa(\mathfrak p)$의 확장
최소 다항식은 분리 가능한다. 즉, 체 확장
$\kappa_i/\kappa(\mathfrak p)$은 원하는 대로 유한하게 분리 가능한다.

\medskip\noindent
에탈 환 지도는 지역적으로 표준인 에탈인 것으로 밝혀졌습니다.
이를 공식화하기 위해 다음 표기법을 도입한다.
환 준동형 $A \to B$가 소아이디얼 $\mathfrak q$ 중 $B$의 것에서 {\it 에탈}이라고 함은 다음 조건을 뜻한다.
$h \in B$, $h \not \in \mathfrak q$가 존재하므로 $A \to B_h$는 에탈이다.
결과는 다음과 같습니다.

\begin{theorem} \label{etale-cohomology-theorem-standard-etale} 환 지도 $A \to B$는 소수 $\mathfrak q$ 필요충분조건은에서 에탈이다. $g \in B$, $g \not \in \mathfrak q$가 존재한다. $B_g$은 $A$에 대한 표준 에탈이다. \end{theorem}

\begin{proof} 대수학, 명제 \ref{algebra-proposition-etale-locally-standard}를 참조하라. \end{proof}





\section{에탈 피복}
\label{etale-cohomology-section-etale-covering}

\noindent 정의를 기억한다.

\begin{definition} \label{etale-cohomology-definition-etale-covering} 스킴 $U$의 {\it 에탈 피복}이란 다음 조건을 만족하는 스킴 사상족 $\{\varphi_i : U_i \to U\}_{i \in I}$를 말한다. \begin{enumerate} \item 각 $\varphi_i$는 에탈 사상이다. \item $U_i$들은 $U$를 덮는다. 즉, $U = \bigcup_{i\in I}\varphi_i(U_i)$이다. \end{enumerate} \end{definition}

\begin{lemma} \label{etale-cohomology-lemma-etale-fpqc} 모든 에탈 피복은 fpqc 피복이다. \end{lemma}

\begin{proof}
(또한 참조
위상,
보조정리 \ref{topologies-lemma-zariski-etale-smooth-syntomic-fppf-fpqc}.)
$\{\varphi_i : U_i \to U\}_{i \in I}$를 에탈 피복라고 하자.
에탈 사상은 평면적이므로 피복의 요소는
목표를 커버하면 성질 fp (충실히 평탄)가 충족된다.
성질 qc(준 컴팩트)를 확인하려면 $V \subset U$를 아핀으로 둡니다.
열고 $\varphi_i^{-1}(V) = \bigcup_{j \in J_i} V_{ij}$라고 씁니다.
일부 아핀의 경우 $V_{ij} \subset U_i$이 열립니다. $\varphi_i$가 열려 있으므로
(에탈 사상이 열려 있으므로), 우리는
$V = \bigcup_{i\in I} \bigcup_{j \in J_i} \varphi_i(V_{ij})$
$V$의 열린 피복이다.
또한, $V$은 준컴팩트하므로 이 피복은 유한한
정제.
\end{proof}

\noindent 따라서 fpqc 피복에 대해 참인 모든 진술은 에탈 피복에 대해 참이다. {\it 더욱이}. 예를 들어, 에탈 사이트는 준정규적이다.

\begin{definition} \label{etale-cohomology-definition-big-etale-site} (자세한 내용은 절 \ref{etale-cohomology-section-big-small} 또는 위상, 절 \ref{topologies-section-etale}를 참조하라.) $S$를 스킴으로 둡니다. {\it 큰 에탈 사이트 위의 $S$}는 사이트 $(\Sch/S)_\etale$이다. 정의를 참조하라. \ref{etale-cohomology-definition-tau-site}. {\it 작은 에탈 사이트 위의 $S$}는 사이트 $S_\etale$이다. 정의를 참조하라. \ref{etale-cohomology-definition-tau-site}. 유사하게 {\it 큰} 및 {\it 작은 Zariski 사이트}를 $S$에 정의한다. $(\Sch/S)_{Zar}$ 및 $S_{Zar}$. \end{definition}

\noindent
대략적으로 말하면 $S$의 큰 에탈 사이트는 $S$에 대한 스킴으로 구성된다.
그리고 피복 에탈 피복. $S$의 작은 에탈 사이트를 구성한다.
밖으로스킴에탈 오버$S$~와 함께피복에탈피복.
실제로 $S_\etale$의 개체 사이에 있는 모든 사상은 에탈이다.
미덕
명제 \ref{etale-cohomology-proposition-etale-morphisms},
그러므로 그것을 확인하기 위해$\{U_i \to U\}_{i \in I}$~에$S_\etale$
피복이다. $\coprod U_i \to U$이 전사인지 확인하는 것으로 충분한다.

\medskip\noindent 작은 에탈 사이트는 큰 에탈 사이트보다 개체 수가 적으며, $S$에 있는 에탈 위상의 ``열림''만 포함한다. 큰 에탈 사이트의 전체 하위 범주이고, 그 위상은 큰 사이트의 위상에서 유도된다. 따라서 큰 에탈 사이트에서 작은 것까지의 제한 함자가 정확하고 단사를 단사로 매핑하는 것이 사실이다. 이는 다음과 같은 결과를 가져옵니다.

\begin{proposition}
\label{etale-cohomology-proposition-cohomology-restrict-small-site}
$S$를 스킴으로 하고 $\mathcal{F}$를 아벨 층으로 설정한다.
$(\Sch/S)_\etale$.
그런 다음 $\mathcal{F}|_{S_\etale}$는 $S_\etale$의 층이고
$$
H^p_\etale(S, \mathcal{F}|_{S_\etale}) =
H^p_\etale(S, \mathcal{F})
$$
모든 $p \geq 0$에 대해.
\end{proposition}

\begin{proof} 이는 보조정리 \ref{etale-cohomology-lemma-compare-cohomology-big-small}의 특별한 경우이다. \end{proof}

\noindent 절 \ref{etale-cohomology-section-big-small}에 도입된 일반 표기법에 따라 위의 코호몰로지 군에 대해 $H_\etale^p(S, \mathcal{F})$를 씁니다.





%9.24.09
\section{커머 이론} \label{etale-cohomology-section-kummer}

\noindent
$n \in \mathbf{N}$라고 하고 다음과 같이 정의된 함자 $\mu_n$를 고려한다.
$$
\begin{matrix}
\Sch^{opp} & \longrightarrow & \textit{Ab} \\
S & \longmapsto &
\mu_n(S)
=
\{t \in \Gamma(S, \mathcal{O}_S^*) \mid t^n = 1 \}.
\end{matrix}
$$
에 의해
그룹체, 예 \ref{groupoids-example-roots-of-unity}
이것은 표현 가능한 함자이고, 이를 표현하는 스킴
$\mu_n$로도 표시된다. 에 의해
보조정리 \ref{etale-cohomology-lemma-representable-sheaf-fpqc}
이 함자는 fpqc 위상에 대한 층 조건을 충족한다.
(특히, 층 조건을 만족한다.
에탈, 자리스키 등 위상).

\begin{lemma} \label{etale-cohomology-lemma-kummer-sequence} 만약 $n\in \mathcal{O}_S^*$이면 $$
0 \to
\mu_{n, S} \to
\mathbf{G}_{m, S} \xrightarrow{(\cdot)^n}
\mathbf{G}_{m, S} \to 0
$$는 $S$의 크고 작은 에탈 사이트에서 층의 짧은 완전열이다. \end{lemma}

\begin{proof}
정의에 의해 층 $\mu_{n, S}$는 지도의 핵이다.
$(\cdot)^n$. 따라서 마지막 지도가 전사임을 보여주는 것으로 충분한다.
$U$를 $S$에 대한 스킴으로 설정한다. 허락하다
$f \in \mathbf{G}_m(U) = \Gamma(U, \mathcal{O}_U^*)$.
우리는 에탈 표지를 찾을 수 있다는 것을 보여줘야 한다.
$U$제한이 있는 회원에 대하여$f$은$n$번째 힘.
두자
$$
U' =
\underline{\Spec}_U(\mathcal{O}_U[T]/(T^n-f))
\xrightarrow{\pi}
U.
$$
(보다
건설, 절 \ref{constructions-section-spec-via-glueing} 또는
\ref{constructions-section-spec}
상대 스펙트럼에 대한 논의를 위해.)
$\Spec(A) \subset U$를 아핀 개방형으로 두고 $f|_{\Spec(A)}$에 해당한다고 가정한다.
$a \in A^*$ 단위로. 그런 다음 $\pi^{-1}(\Spec(A)) = \Spec(B)$
$B = A[T]/(T^n - a)$. 환 맵 $A \to B$은 랭크 $n$가 없는 유한한 맵이다.
따라서 그것은 충실히 평탄이므로 다음과 같이 결론을 내립니다.
$\Spec(B) \to \Spec(A)$는 전사이다. 이는 모든 경우에 적용되므로
아핀 열린부분 에서 $U$ 우리는 $\pi$가 전사라고 결론을 내립니다.
또한 $n$와 $T^{n - 1}$는 $B$에서 반전이 가능하므로
$nT^{n-1} \in B^*$ 그리고 환 지도 $A \to B$는 에탈이 표준이고,
특히 에탈. 이는 $U$의 모든 아핀 개방에 대해 유지되므로
우리는 $\pi$이 에탈이라고 결론을 내립니다. 따라서
$\mathcal{U} = \{\pi : U' \to U\}$은 에탈 피복이다.
또한 $f|_{U'} = (f')^n$ 여기서 $f'$는 $T$의 클래스이다.
$\Gamma(U', \mathcal{O}_{U'}^*)$에 있으므로 $\mathcal{U}$에는 원하는 성질이 있다.
\end{proof}

\begin{remark} \label{etale-cohomology-remark-no-kummer-sequence-zariski} 보조정리 \ref{etale-cohomology-lemma-kummer-sequence}는 ``에탈''을 ``Zariski''로 바꾸면 거짓이다. 에탈 위상은 부드러운 위상보다 거칠기 때문에 위상, 보조정리 \ref{topologies-lemma-zariski-etale-smooth}를 참조하라. 이는 부드러운 위상에서도 시퀀스가 ​​정확하다는 것을 의미한다. \end{remark}

\noindent 정리 \ref{etale-cohomology-theorem-picard-group} 및 보조정리 \ref{etale-cohomology-lemma-kummer-sequence} 및 코호몰로지의 일반 성질에 의해 길고 정확한 코호몰로지 시퀀스를 얻습니다. $$
\xymatrix{
0 \ar[r] &
H_\etale^0(S, \mu_{n, S}) \ar[r] &
\Gamma(S, \mathcal{O}_S^*) \ar[r]^{(\cdot)^n} &
\Gamma(S, \mathcal{O}_S^*) \ar@(rd, ul)[rdllllr]
\\
& H_\etale^1(S, \mu_{n, S}) \ar[r] &
\Pic(S) \ar[r]^{(\cdot)^n} &
\Pic(S) \ar@(rd, ul)[rdllllr] \\
& H_\etale^2(S, \mu_{n, S}) \ar[r] &
\ldots
}
$$ 적어도 $n$가 $S$에서 반전 가능한 경우이다. $n$가 $S$에서 반전될 수 없는 경우 다음 보조정리를 적용할 수 있다.

\begin{lemma} \label{etale-cohomology-lemma-kummer-sequence-syntomic} 임의의 $n \in \mathbf{N}$에 대해 시퀀스 $$
0 \to
\mu_{n, S} \to
\mathbf{G}_{m, S} \xrightarrow{(\cdot)^n}
\mathbf{G}_{m, S} \to 0
$$는 사이트 $(\Sch/S)_{fppf}$에 있는 층의 짧은 완전열이고 $(\Sch/S)_{syntomic}$. \end{lemma}

\begin{proof}
정의에 의해 층 $\mu_{n, S}$는 지도의 핵이다.
$(\cdot)^n$. 따라서 마지막 지도가 전사임을 보여주는 것으로 충분한다.
신토믹 위상은 fppf 위상보다 약하므로 다음을 참조하라.
위상, 보조정리 \ref{topologies-lemma-zariski-etale-smooth-syntomic-fppf},
신토믹 위상에 대해 이를 증명하는 것으로 충분한다.
$U$를 $S$에 대한 스킴으로 설정한다. 허락하다
$f \in \mathbf{G}_m(U) = \Gamma(U, \mathcal{O}_U^*)$.
우리는 다음의 구문적 표지를 찾을 수 있다는 것을 보여줘야 한다.
$U$제한이 있는 회원에 대하여$f$은$n$번째 힘.
두자
$$
U' =
\underline{\Spec}_U(\mathcal{O}_U[T]/(T^n-f))
\xrightarrow{\pi}
U.
$$
(보다
건설, 절 \ref{constructions-section-spec-via-glueing} 또는
\ref{constructions-section-spec}
상대 스펙트럼에 대한 논의를 위해.)
$\Spec(A) \subset U$를 아핀 개방형으로 두고 $f|_{\Spec(A)}$에 해당한다고 가정한다.
$a \in A^*$ 단위로. 그런 다음 $\pi^{-1}(\Spec(A)) = \Spec(B)$
$B = A[T]/(T^n - a)$. 환 맵 $A \to B$은 랭크 $n$가 없는 유한한 맵이다.
따라서 그것은 충실히 평탄이므로 다음과 같이 결론을 내립니다.
$\Spec(B) \to \Spec(A)$는 전사이다. 이는 모든 경우에 적용되므로
아핀 열린부분 에서 $U$ 우리는 $\pi$가 전사라고 결론을 내립니다.
또한 $B$는 $A$에 대한 전역 상대 완전 교차점이므로
환 맵 $A \to B$은 표준 구문이다.
특히 신토믹. 이는 $U$의 모든 아핀 개방에 대해 유지되므로
우리는 $\pi$가 구문론적이라고 결론을 내립니다. 따라서
$\mathcal{U} = \{\pi : U' \to U\}$은 구문론적 피복이다.
또한 $f|_{U'} = (f')^n$ 여기서 $f'$는 $T$의 클래스이다.
$\Gamma(U', \mathcal{O}_{U'}^*)$에 있으므로 $\mathcal{U}$에는 원하는 성질이 있다.
\end{proof}

\begin{remark} \label{etale-cohomology-remark-no-kummer-sequence-smooth-etale-zariski} 보조정리 \ref{etale-cohomology-lemma-kummer-sequence-syntomic}는 매끄러운, 에탈 또는 Zariski 위상에 대해 거짓이다. \end{remark}

\noindent 정리 \ref{etale-cohomology-theorem-picard-group} 및 보조정리 \ref{etale-cohomology-lemma-kummer-sequence-syntomic} 및 코호몰로지의 일반 성질에 의해 길고 정확한 코호몰로지 시퀀스를 얻습니다. $$
\xymatrix{
0 \ar[r] &
H_{fppf}^0(S, \mu_{n, S}) \ar[r] &
\Gamma(S, \mathcal{O}_S^*) \ar[r]^{(\cdot)^n} &
\Gamma(S, \mathcal{O}_S^*) \ar@(rd, ul)[rdllllr]
\\
& H_{fppf}^1(S, \mu_{n, S}) \ar[r] &
\Pic(S) \ar[r]^{(\cdot)^n} &
\Pic(S) \ar@(rd, ul)[rdllllr] \\
& H_{fppf}^2(S, \mu_{n, S}) \ar[r] &
\ldots
}
$$ 임의의 스킴 $S$ 및 임의의 정수 $n$에 대해. 물론 신토믹 코호몰로지와 비슷한 순서가 있다.

\medskip\noindent
$n \in \mathbf{N}$로 하고 $S$는 임의의 스킴으로 둡니다.
첫 번째 코호몰로지 군을 설명하는 또 다른 더 직접적인 방법이 있다.
$\mu_n$의 값이다. 쌍을 고려하십시오
$(\mathcal{L}, \alpha)$ 여기서 $\mathcal{L}$는 $S$의 가역층이다.
그리고 $\alpha : \mathcal{L}^{\otimes n} \to \mathcal{O}_S$는 사소화이다.
~의$n$텐서 파워의$\mathcal{L}$.
$(\mathcal{L}', \alpha')$이 두 번째 쌍이라고 가정한다.
동형사상 $\varphi : (\mathcal{L}, \alpha) \to (\mathcal{L}', \alpha')$
는 동형사상 $\varphi : \mathcal{L} \to \mathcal{L}'$ 이다.
층 다이어그램은 다음과 같습니다.
$$
\xymatrix{
\mathcal{L}^{\otimes n} \ar[d]_{\varphi^{\otimes n}} \ar[r]_\alpha &
\mathcal{O}_S \ar[d]^1 \\
(\mathcal{L}')^{\otimes n} \ar[r]^{\alpha'} &
\mathcal{O}_S \\
}
$$
통근. 따라서 우리는
\begin{equation}
\label{etale-cohomology-equation-isomorphisms-pairs}
\mathit{Isom}_S((\mathcal{L}, \alpha), (\mathcal{L}', \alpha'))
=
\left\{
\begin{matrix}
\emptyset & \text{만약에} & \text{동형이 아닙니다} \\
H^0(S, \mu_{n, S})\cdot \varphi & \text{만약에} &
\varphi \text{ 동형사상 쌍}
\end{matrix}
\right.
\end{equation}
또한, 두 쌍의 $(\mathcal{L}, \alpha)$, $(\mathcal{L}', \alpha')$가 주어지면
텐서곱
$$
(\mathcal{L}, \alpha) \otimes (\mathcal{L}', \alpha')
=
(\mathcal{L} \otimes \mathcal{L}', \alpha \otimes \alpha')
$$
또 다른 쌍이다. $(\mathcal{O}_S, 1)$ 쌍은 이에 대한 ID이다.
텐서곱 연산을 수행하고 그 역은 다음과 같이 주어진다.
$$
(\mathcal{L}, \alpha)^{-1} = (\mathcal{L}^{\otimes -1}, \alpha^{\otimes -1}).
$$
따라서 동형사상 클래스 쌍의 모음은 아벨 군을 형성한다.
참고하세요
$$
(\mathcal{L}, \alpha)^{\otimes n}
=
(\mathcal{L}^{\otimes n}, \alpha^{\otimes n})
\xrightarrow{\alpha}
(\mathcal{O}_S, 1)
$$
동형사상이다.
따라서 이 그룹의 모든 요소는 $n$로 구분되는 순서를 갖습니다. 우리는 독자에게 경고합니다
이 그룹은 일반적으로{\bf~ 아니다}그만큼$n$- 비틀림$\Pic(S)$.

\begin{lemma}
\label{etale-cohomology-lemma-describe-h1-mun}
$S$를 스킴으로 설정한다. 표준 식별이 있다.
$$
H_\etale^1(S, \mu_n) =
\text{쌍 그룹 }(\mathcal{L}, \alpha)\text{ 최대 동형사상 위와 같음}
$$
$n$가 $S$에서 반전 가능한 경우. 일반적으로 우리는
$$
H_{fppf}^1(S, \mu_n) =
\text{쌍 그룹 }(\mathcal{L}, \alpha)\text{ 최대 동형사상 위와 같습니다}.
$$
fppf를 신토믹으로 대체해도 동일한 결과가 유지된다.
\end{lemma}

\begin{proof}
먼저 두 번째 동형사상을 증명한다.
$(\mathcal{L}, \alpha)$를 위와 같이 쌍으로 둡니다.
다음과 같은 아핀 오픈 피복 $S = \bigcup U_i$을 선택한다.
$\mathcal{L}|_{U_i} \cong \mathcal{O}_{U_i}$. $s_i \in \mathcal{L}(U_i)$라고 말하세요.
발전기이다. 그런 다음 $\alpha(s_i^{\otimes n}) = f_i \in \mathcal{O}_S^*(U_i)$.
$U_i = \Spec(A_i)$라고 쓰면 글로벌이 존재한다는 것을 알 수 있다.
상대 완전 교차점 $A_i \to B_i = A_i[T]/(T^n - f_i)$
그렇게$f_i$에 매핑$n$번째 전원$B_i$. 즉, 설정
$V_i = \Spec(B_i)$ 신토믹 피복을 얻습니다.
$\mathcal{V} = \{V_i \to S\}_{i \in I}$ 및 사소화
$\varphi_i : (\mathcal{L}, \alpha)|_{V_i} \to (\mathcal{O}_{V_i}, 1)$.

\medskip\noindent 이 결과(피복 $\mathcal{V}$의 존재)를 사용하여 이 쌍에 $H^1_{syntomic}(S, \mu_{n, S})$의 코호몰로지 클래스를 연결한다. 우리는 두 가지 (동등한) 구성을 제공한다.

\medskip\noindent 첫 번째 구성: čech 코호몰로지 사용. 이중 중첩 $V_i \times_S V_j$에 대해 동형사상 $$
(\mathcal{O}_{V_i \times_S V_j}, 1)
\xrightarrow{\text{pr}_0^*\varphi_i^{-1}}
(\mathcal{L}|_{V_i \times_S V_j}, \alpha|_{V_i \times_S V_j})
\xrightarrow{\text{pr}_1^*\varphi_j}
(\mathcal{O}_{V_i \times_S V_j}, 1)
$$ 쌍이 있다. (\ref{etale-cohomology-equation-isomorphisms-pairs})에 의해 이는 $\zeta_{ij} \in \mu_n(V_i \times_S V_j)$ 요소에 의해 제공된다. 우리는 이러한 $\zeta_{ij}$가 $1$-코사이클을 제공한다는 확인을 생략한다. 즉, 요소 ​​$(\zeta_{i_0i_1}) \in \check C(\mathcal{V}, \mu_n)$에 $d(\zeta_{i_0i_1}) = 0$를 제공한다. 따라서 해당 클래스는 $\check H^1(\mathcal{V}, \mu_n)$의 요소이고 정리 \ref{etale-cohomology-theorem-cech-ss}에 의해 $H^1_{syntomic}(S, \mu_{n, S})$의 코호몰로지 클래스에 매핑된다.

\medskip\noindent
두 번째 구성: 토르서를 사용한다. 전층을 고려해보세요.
$$
\mu_n(\mathcal{L}, \alpha) :
U
\longmapsto
\mathit{Isom}_U((\mathcal{O}_U, 1), (\mathcal{L}, \alpha)|_U)
$$
$(\Sch/S)_{syntomic}$에 있다.
우리는 이것을 다음의 서브프레시프로 볼 수 있다.
$\SheafHom_\mathcal{O}(\mathcal{O}, \mathcal{L})$ (내부 Hom
층, 참조
가군 사이트, 절 \ref{sites-modules-section-internal-hom}).
이 서브프리쉬프를 정의하는 조건은 지역적이므로 다음과 같습니다.
층.
(\ref{etale-cohomology-equation-isomorphisms-pairs})에 의해 이 층에는 다음과 같은 자유 동작이 있다.
층 $\mu_{n, S}$. 그러므로 우리가 확인해야 할 유일한 것은
로컬에는 절이 있다. 이는 의 존재로 인해 사실이다.
표지 $\mathcal{V}$를 평범하게 만듭니다. 그러므로 $\mu_n(\mathcal{L}, \alpha)$
$\mu_{n, S}$-토르서이고
코호몰로지 사이트, 보조정리 \ref{sites-cohomology-lemma-torsors-h1}
$H^1_{syntomic}(S, \mu_{n, S})$의 해당 요소를 얻습니다.

\medskip\noindent 좋다. 이제 다음 \begin{enumerate} \item 두 구성은 동일한 코호몰로지 클래스를 제공한다. \item 동형 쌍은 동일한 코호몰로지 클래스를 생성한다. \item $(\mathcal{L}, \alpha) \otimes (\mathcal{L}', \alpha')$의 코호몰로지 클래스는 $(\mathcal{L}, \alpha)$와 $(\mathcal{L}', \alpha')$의 코호몰로지 클래스의 합이다. \item 코호몰로지 클래스가 사소한 경우 쌍도 사소한 것이다. \item $H^1_{syntomic}(S, \mu_{n, S})$의 모든 요소는 쌍의 코호몰로지 클래스이다. \end{enumerate} (1)의 증명을 생략한다. 부분(2)은 두 번째 구성에서 명확한다. 왜냐하면 동형 토서는 동일한 코호몰로지 클래스를 제공하기 때문이다. 부분(3)은 결과로 생성된 čech 클래스가 합산되므로 첫 번째 구성부터 명확한다. 부분(4)은 토르서가 사소하기 때문에 두 번째 구성에서 명확한다. 필요충분조건은 여기에는 대역 절단이 있다. 코호몰로지, 사이트, 보조정리 \ref{sites-cohomology-lemma-trivial-torsor}를 참조하라.

\medskip\noindent
부분(5)은 다음과 같이 볼 수 있습니다(직접 증명은
선택할 만한). $\xi \in H^1_{syntomic}(S, \mu_{n, S})$라고 가정한다.
그런 다음 $\xi$는 요소에 매핑된다.
$\overline{\xi} \in H^1_{syntomic}(S, \mathbf{G}_{m, S})$
$n \overline{\xi} = 0$로. 에 의해
정리 \ref{etale-cohomology-theorem-picard-group}
$\overline{\xi}$가 가역층 $\mathcal{L}$에 해당하는 것을 알 수 있다.
$n$번째 텐서 파워는 $\mathcal{O}_S$와 동형이다.
따라서 코호몰로지를 갖는 $(\mathcal{L}, \alpha')$ 쌍이 존재한다.
클래스 $\xi'$에 동일한 이미지 $\overline{\xi'}$가 있다.
$H^1_{syntomic}(S, \mathbf{G}_{m, S})$. 그래서 보여주기엔 충분해요
$\xi - \xi'$은 쌍의 클래스이다. 건축으로, 그리고
위의 길고 정확한 코호몰로지 시퀀스를 보면
$\xi - \xi' = \partial(f)$ 일부 $f \in H^0(S, \mathcal{O}_S^*)$에 대해.
$(\mathcal{O}_S, f)$ 쌍을 고려해보세요. 검증을 생략합니다
이 쌍의 코호몰로지 클래스는 $\partial(f)$이다.
첫 번째 식별의 증명을 완료합니다(fppf를 교체함)
신토믹과 함께).

\medskip\noindent 첫 번째를 보려면 $n$가 $S$에서 반전 가능하다면 증명의 첫 번째 부분에 구성된 피복 $\mathcal{V}$는 실제로 에탈이다. 피복 (보조정리 \ref{etale-cohomology-lemma-kummer-sequence}의 증명과 비교). 증명의 나머지 부분은 Kummer 수열이 정확하다는 것을 사용하는 마지막 인수를 제외하고 위상과 독립적이다. 즉, 보조정리 \ref{etale-cohomology-lemma-kummer-sequence}을 사용한다. \end{proof}






\section{이웃, 줄기 및 지점} \label{etale-cohomology-section-stalks}

\noindent $S$의 기하점과 줄기 함자의 정확한 항목을 연결할 수 있다. $S_\etale$에 있는 층의 지도는 모든 줄기에 있는 동형사상인 경우에만 동형사상이다. 아벨 층의 복합체는 정확합니다 필요충분조건은 복합체는 줄기의 기하점이다. 이는 전체적으로 스킴 $S$의 작은 에탈 사이트에 충분한 점수가 있음을 의미한다. 또한 $S$(추상적 개념)의 작은 에탈 토포스의 임의의 점은 기하점으로 주어지는 것으로 밝혀졌습니다. 따라서 어떤 의미에서 $S$의 작은 에탈 토포스는 기하점과 이웃이라는 관점에서 이해될 수 있다.

\begin{definition}
\label{etale-cohomology-definition-geometric-point}
$S$를 스킴으로 설정한다.
\begin{enumerate}
\item {\it 기하점} $S$는 사상이다.
$\Spec(k) \to S$ 여기서 $k$는 대수적으로 닫혀 있다.
이러한 점은 일반적으로 $\overline{s}$로 표시된다.
작은 케이스 편지. 우리는 스킴을 나타내기 위해 종종 $\overline{s}$를 사용한다.
$\Spec(k)$ 뿐만 아니라 사상도 사용하며 $\kappa(\overline{s})$를 사용한다.
$k$을 나타냅니다.
\item $\overline{s}$가 $s$의 {\it 위에 있다}고 함은
$s \in S$이 $\overline{s}$의 이미지임을 나타냅니다.
\item {\it 에탈 이웃} 의 기하점 $\overline{s}$
$S$는 가환 도식이다.
$$
\xymatrix{
& U \ar[d]^\varphi \\
{\overline{s}} \ar[r]^{\overline{s}} \ar[ur]^{\bar u} & S
}
$$
여기서 $\varphi$는 스킴의 에탈 사상이다.
$(U, \overline{u}) \to (S, \overline{s})$라고 씁니다.
\item {\it 에탈 근방의 사상}
$(U, \overline{u}) \to (U', \overline{u}')$
$S$-사상 $h: U \to U'$이다.
$\overline{u}' = h \circ \overline{u}$와 같습니다.
\end{enumerate}
\end{definition}

\begin{remark} \label{etale-cohomology-remark-etale-between-etale} $U$와 $U'$는 $S$에 대해 에탈이므로, 이들 사이의 모든 $S$-사상도 에탈이다. 명제를 참조하라. \ref{etale-cohomology-proposition-etale-morphisms}. 특히 에탈 동네의 모든 사상은 에탈이다. \end{remark}

\begin{remark}
\label{etale-cohomology-remark-etale-neighbourhoods}
$S$를 스킴으로 하고 $s \in S$를 점으로 둡니다. ~ 안에
사상에 대해 자세히 알아보세요.
정의 \ref{more-morphisms-definition-etale-neighbourhood}
우리는 에탈 이웃 $(U, u) \to (S, s)$의 개념을 정의했다.
$(S, s)$. $\overline{s}$가 $S$의 기하점인 경우
$s$, 그러면 임의의 에탈 이웃 $(U, \overline{u}) \to (S, \overline{s})$
에탈 동네가 생긴다$(U, u)$~의$(S, s)$복용함으로써
$u \in U$는 $U$의 고유 지점이 된다. 즉, $\overline{u}$이다.
$u$ 위에 있다. 반대로, 에탈 이웃이 주어지면 $(U, u)$
$(S, s)$ 잔기 체 확장자 $\kappa(u)/\kappa(s)$
유한 분리 가능합니다(참조
명제 \ref{etale-cohomology-proposition-etale-morphisms})
따라서 우리는 임베딩 $\kappa(u) \subset \kappa(\overline{s})$을 찾을 수 있다.
$\kappa(s)$ 이상. 즉, 기하점을 찾을 수 있다.
$\overline{u}$ 중 $U$가 $u$ 위에 놓여 있어 $(U, \overline{u})$
은 $(S, \overline{s})$의 에탈 동네이다.
우리는 이러한 관찰을 사용하여 두 가지 유형의 유형 사이를 이동한다.
에탈 동네.
\end{remark}

\begin{lemma}
\label{etale-cohomology-lemma-cofinal-etale}
$S$를 스킴으로 두고, $\overline{s}$를 $S$의 기하점으로 둡니다.
에탈 동네의 범주가 공동 필터링된다. 더 정확하게는:
\begin{enumerate}
\item $(U_i, \overline{u}_i)_{i = 1, 2}$를 두 개의 에탈 이웃이라고 하자.
$\overline{s}$~에$S$. 그러면 세 번째 에탈 동네가 존재한다.
$(U, \overline{u})$ 및 사상
$(U, \overline{u}) \to (U_i, \overline{u}_i)$, $i = 1, 2$.
\item $h_1, h_2: (U, \overline{u}) \to (U', \overline{u}')$를 2개라고 하자.
사상 $\overline{s}$의 에탈 동네 사이. 그런 다음 존재한다.
에탈 동네 $(U'', \overline{u}'')$ 및 사상
$h : (U'', \overline{u}'') \to (U, \overline{u})$
이는 $h_1$와 $h_2$를 동일화한다. 즉,
$h_1 \circ h = h_2 \circ h$.
\end{enumerate}
\end{lemma}

\begin{proof}
부품(1)의 경우 올곱 $U = U_1 \times_S U_2$을 고려한다.
에탈 사상이므로 $U_1$과 $U_2$ 모두에 대해 에탈이다.
기본 변경 시 보존된다. 참조
명제 \ref{etale-cohomology-proposition-etale-morphisms}.
지도$\overline{s} \to U$에 의해 정의됨$(\overline{u}_1, \overline{u}_2)$
두 가지 모두에 매핑되는 에탈 이웃의 구조를 제공한다.
$U_1$ 및 $U_2$. 부품(2)에 대해 $U''$를 올곱으로 정의한다.
$$
\xymatrix{
U'' \ar[r] \ar[d] & U \ar[d]^{(h_1, h_2)} \\
U' \ar[r]^-\Delta & U' \times_S U'.
}
$$
$\overline{u}$와 $\overline{u}'$는 $S$와 $\overline{s}$에 동의하므로,
$\overline{u}'' = (\overline{u}, \overline{u}')$는 기하학적이라는 것을 알 수 있습니다
$U''$ 지점. 특히 $U'' \not = \emptyset$.
게다가 $U'$는 $S$에 대해 에탈이므로 올곱도 마찬가지이다.
$U'\times_S U'$ (참조
명제 \ref{etale-cohomology-proposition-etale-morphisms}).
따라서 수직 화살표 $(h_1, h_2)$는 에탈이다.
비고 \ref{etale-cohomology-remark-etale-between-etale} 위.
그러므로 $U''$은 염기 변화에 의해 $U'$에 대해 에탈이고, 따라서
에탈은 $S$에 걸쳐 있습니다(에탈 사상의 구성은 에탈이기 때문입니다).
따라서 $(U'', \overline{u}'')$은 문제에 대한 해결책이다.
\end{proof}

\begin{lemma} \label{etale-cohomology-lemma-geometric-lift-to-cover} $S$를 스킴으로 둡니다. $\overline{s}$를 $S$의 기하점으로 설정한다. $(U, \overline{u})$를 $\overline{s}$의 에탈 동네라 하자. $\mathcal{U} = \{\varphi_i : U_i \to U \}_{i\in I}$를 에탈 피복라고 하자. 그러면 $i \in I$와 $\overline{u}_i : \overline{s} \to U_i$가 존재하므로 $\varphi_i : (U_i, \overline{u}_i) \to (U, \overline{u})$는 에탈 동네의 사상이다. \end{lemma}

\begin{proof}
$U = \bigcup_{i\in I} \varphi_i(U_i)$로서 올곱
$\overline{s} \times_{\overline{u}, U, \varphi_i} U_i$이(가) 비어 있지 않습니다.
일부 $i$의 경우. 그런 다음 데카르트 다이어그램을 살펴보세요.
$$
\xymatrix{
\overline{s} \times_{\overline{u}, U, \varphi_i} U_i
\ar[d]^{\text{pr}_1} \ar[r]_-{\text{pr}_2} & U_i
\ar[d]^{\varphi_i} \\
\Spec(k) = \overline{s} \ar@/^1pc/[u]^\sigma
\ar[r]^-{\overline{u}} & U
}
$$
투영 $\text{pr}_1$은 에탈 사상의 기본 변경이므로
에탈이야, 봐
명제 \ref{etale-cohomology-proposition-etale-morphisms}.
따라서 $\overline{s} \times_{\overline{u}, U, \varphi_i} U_i$
는 $k$의 유한 분리 가능한 확장의 분리된 결합이다.
명제 \ref{etale-cohomology-proposition-etale-morphisms}. 여기
$\overline{s} = \Spec(k)$. 그러나 $k$은 대수적으로 닫혀 있으므로 모두
이러한 확장은 사소하며 절 $\sigma$
$\text{pr}_1$. 구성
$\text{pr}_2 \circ \sigma$는 $\overline{u}$과 호환되는 지도를 제공한다.
\end{proof}

\begin{definition}
\label{etale-cohomology-definition-stalk}
$S$를 스킴으로 설정한다.
$\mathcal{F}$를 $S_\etale$에서 전층으로 설정한다.
$\overline{s}$를 $S$의 기하점으로 설정한다.
그만큼{\it 줄기}~의$\mathcal{F}$~에$\overline{s}$~이다
$$
\mathcal{F}_{\overline{s}}
=
\colim_{(U, \overline{u})} \mathcal{F}(U)
$$
여기서 $(U, \overline{u})$는 에탈 전체에 걸쳐 실행된다.
인근 지역$\overline{s}$~에$S$.
\end{definition}

\noindent
보조정리 \ref{etale-cohomology-lemma-cofinal-etale}에 의해 이 여극한은 필터링되었다.
지수 범주, 즉 에탈 지역의 범주와 반대이다.
즉, $\mathcal{F}_{\overline{s}}$의 요소는 다음과 같습니다.
트리플 $(U, \overline{u}, \sigma)$로 생각된다.
$\sigma \in \mathcal{F}(U)$. 트리플 2개
$(U, \overline{u}, \sigma)$, $(U', \overline{u}', \sigma')$
세 번째 요소가 있는 경우 줄기의 동일한 요소를 정의한다.
에탈의 에탈 동네 $(U'', \overline{u}'')$와 사상
동네 $h : (U'', \overline{u}'') \to (U, \overline{u})$,
$h' : (U'', \overline{u}'') \to (U', \overline{u}')$ 이렇게
$h^*\sigma = (h')^*\sigma'$~에$\mathcal{F}(U'')$. 보다
범주론, 절 \ref{categories-section-directed-colimits}.

\begin{lemma}
\label{etale-cohomology-lemma-stalk-gives-point}
$S$를 스킴으로 설정한다. $\overline{s}$를 $S$의 기하점으로 설정한다.
함자를 고려해보세요.
\begin{align*}
u : S_\etale & \longrightarrow \textit{Sets}, \\
U & \longmapsto
|U_{\overline{s}}|
=
\{\overline{u} \text{이고 }(U, \overline{u})
\text{가 }\overline{s}\text{의 에탈 근방인 것}\}.
\end{align*}
여기서 $|U_{\overline{s}}|$는 기하 올의 기본 집합을 나타냅니다.
그 다음에$u$점을 정의합니다$p$~의사이트 $S_\etale$
(사이트, 정의 \ref{sites-definition-point})
및 관련 줄기 함자 $\mathcal{F} \mapsto \mathcal{F}_p$
(사이트, 방정식 \ref{sites-equation-stalk})
함자 $\mathcal{F} \mapsto \mathcal{F}_{\overline{s}}$ 입니다
위에 정의되어 있다.
\end{lemma}

\begin{proof}
증명에서
보조정리 \ref{etale-cohomology-lemma-geometric-lift-to-cover}
우리는 스킴 $U_{\overline{s}}$가 다음의 분리된 합집합이라는 것을 보았습니다.
스킴은 $\overline{s}$와 동형이다. 따라서 우리는 또한 생각할 수 있습니다
$|U_{\overline{s}}|$ 집합의 기하점의 $U$가 누워 있는 것처럼
$\overline{s}$, 즉 사상의 컬렉션으로
$\overline{u} : \overline{s} \to U$ 다이어그램에 적합
정의 \ref{etale-cohomology-definition-geometric-point}.
이것으로부터 $u(S)$는 싱글톤이고,
$u(U \times_V W) = u(U) \times_{u(V)} u(W)$
$U \to V$ 및 $W \to V$가 $S_\etale$에서 사상일 때마다.
그리고 주어진피복 $\{U_i \to U\}_{i \in I}$~에$S_\etale$
우리는 $\coprod u(U_i) \to u(U)$가 전사라는 것을 알 수 있다.
보조정리 \ref{etale-cohomology-lemma-geometric-lift-to-cover}.
따라서
사이트, 명제 \ref{sites-proposition-point-limits}
적용되므로 $p$는 사이트 $S_\etale$의 지점이다.
마지막으로, 우리의 함자 $\mathcal{F} \mapsto \mathcal{F}_{\overline{s}}$
함자와 정확히 동일한 여극한에 의해 제공된다.
$\mathcal{F} \mapsto \mathcal{F}_p$는 다음의 $p$에 연결되어 있다.
사이트, 방정식 \ref{sites-equation-stalk}
이는 최종 주장을 증명한다.
\end{proof}

\begin{remark}
\label{etale-cohomology-remark-map-stalks}
$S$를 스킴으로 두고 $\overline{s} : \Spec(k) \to S$로 설정한다.
그리고 $\overline{s}' : \Spec(k') \to S$는 2개의 기하점이어야 한다.
$S$. {\it 사상 $a : \overline{s} \to \overline{s}'$ 의 기하점}
단순히 사상 $a : \Spec(k) \to \Spec(k')$이다.
$\overline{s}' \circ a = \overline{s}$. 이러한 사상이 주어지면 우리는 다음을 얻습니다.
$\overline{s}'$의 에탈 동네의 범주에서 함자
규칙에 따라 $\overline{s}$의 에탈 동네의 범주에
$(U, \overline{u}') \mapsto (U, \overline{u}' \circ a)$. 그러므로 우리는 얻는다
표준 지도
$$
\mathcal{F}_{\overline{s}'}
=
\colim_{(U, \overline{u}')} \mathcal{F}(U)
\longrightarrow
\colim_{(U, \overline{u})} \mathcal{F}(U)
=
\mathcal{F}_{\overline{s}}
$$
범주, 보조정리 \ref{categories-lemma-functorial-colimit}에서. 사용하여
줄기 요소에 대한 설명을 삼중으로 표현하면 다음 요소가 매핑된다.
$\mathcal{F}_{\overline{s}'}$ 트리플로 표현
$(U, \overline{u}', \sigma)$를 $\mathcal{F}_{\overline{s}}$의 요소로
트리플 $(U, \overline{u}' \circ a, \sigma)$로 표현된다.
위의 함자는 분명히 동등하므로 우리는 다음과 같이 결론을 내립니다.
표준 맵은 줄기 함자의 동형사상이다.

\medskip\noindent
$a$에 해당하는 줄기의 맵이 있는지 확인하겠습니다.
올바른 방향으로. 위의 의미는 다음과 같습니다.
사이트, 정의 \ref{sites-definition-morphism-points},
저것$a$정의하다사상 $a : p \to p'$포인트 사이$p, p'$~의
그만큼사이트 $S_\etale$관련된$\overline{s}, \overline{s}'$~에 의해
보조정리 \ref{etale-cohomology-lemma-stalk-gives-point}. 보다 일반적인 사상
포인트 ($S$ 포인트의 특수화에 해당)
나중에 설명하고 동형사상이 아닌 내용은 절을 참조하라.
\ref{etale-cohomology-section-specialization}.
\end{remark}

\begin{lemma}
\label{etale-cohomology-lemma-stalk-exact}
$S$를 스킴으로 설정한다. $\overline{s}$를 $S$의 기하점으로 설정한다.
\begin{enumerate}
\item 줄기 함자
$\textit{PAb}(S_\etale) \to \textit{Ab}$,
$\mathcal{F}  \mapsto  \mathcal{F}_{\overline{s}}$
정확한다.
\item $(\mathcal{F}^\#)_{\overline{s}} = \mathcal{F}_{\overline{s}}$
누구에게나전층~의집합 $\mathcal{F}$~에$S_\etale$.
\item 함자
$\textit{Ab}(S_\etale) \to \textit{Ab}$,
$\mathcal{F} \mapsto \mathcal{F}_{\overline{s}}$이 정확한다.
\item 마찬가지로 함자
$\textit{PSh}(S_\etale) \to \textit{Sets}$ 그리고
$\Sh(S_\etale) \to \textit{Sets}$ 줄기 함자에 의해 주어진다.
$\mathcal{F} \mapsto \mathcal{F}_{\overline{s}}$ 정확합니다(참조
범주론, 정의 \ref{categories-definition-exact})
임의의 여극한으로 출퇴근한다.
\end{enumerate}
\end{lemma}

\begin{proof} 직접적인 논증으로 이를 증명하는 방법을 나타내기 전에 사이트, 절 \ref{sites-modules-section-stalks}에 대한 가군의 일반 자료에서 결과가 나온다는 점에 유의하세요. 이는 $\mathcal{F} \mapsto \mathcal{F}_{\overline{s}}$가 $S$의 작은 에탈 사이트 지점에서 오기 때문에 사실이다. 보조정리 \ref{etale-cohomology-lemma-stalk-gives-point}를 참조하라. (1), (2) 및 (3)의 직접 증명만 제공하고 (4)의 직접 증명은 생략한다.

\medskip\noindent
완전성은 $\textit{PAb}(S_\etale)$의 함자로서 형식적이다.
여극한 모든 여극한 및 유한으로 통근하도록 지시한 사실
극한. (2)에 있는 줄기의 식별은 지도를 통해 이루어집니다.
$$
\kappa :
\mathcal{F}_{\overline{s}}
\longrightarrow
(\mathcal{F}^\#)_{\overline{s}}
$$
자연적인 사상 $\mathcal{F}\to \mathcal{F}^\#$에 의해 유도된다. 참조
정리 \ref{etale-cohomology-theorem-sheafification}.
우리는 주장 이 지도가 동형사상 아벨 군임을 확인한다. 우리는 보여줄 것이다
주입성을 나타내고 전사성의 증명을 생략한다.

\medskip\noindent
$\sigma\in \mathcal{F}_{\overline{s}}$로 놔두세요.
에탈 동네가 있어요
$(U, \overline{u})\to (S, \overline{s})$ $\sigma$는 일부 이미지이다.
절 $s \in \mathcal{F}(U)$. $\kappa(\sigma) = 0$인 경우
$(\mathcal{F}^\#)_{\overline{s}}$ 그러면 에탈의 사상이 존재한다.
이웃 $(U', \overline{u}')\to (U, \overline{u})$
$s|_{U'}$는 $\mathcal{F}^\#(U')$에서 0이다. 거기에 이어
에탈이 존재합니다 피복
$\{U_i'\to U'\}_{i\in I}$ 즉, $s|_{U_i'}=0$
$\mathcal{F}(U_i')$모두를 위해$i$. 에 의해보조정리 \ref{etale-cohomology-lemma-geometric-lift-to-cover}
$i \in I$ 및 사상이 존재한다.
$\overline{u}_i': \overline{s} \to U_i'$ 이렇게
$(U_i', \overline{u}_i') \to (U', \overline{u}')\to (U, \overline{u})$
에탈 동네는 사상이다. 그러므로 $\sigma = 0$
이후 $(U_i', \overline{u}_i') \to (U, \overline{u})$
는 에탈 동네의 사상이다.
$s|_{U'_i}=0$이 있다. 이는 $\kappa$가 분사형이라는 것을 증명한다.

\medskip\noindent 함자 $\textit{Ab}(S_\etale) \to \textit{Ab}$가 정확하다는 것을 보여주기 위해 $\textit{Ab}(S_\etale)$의 짧은 완전열을 고려하라. $
0\to \mathcal{F}\to \mathcal{G}\to \mathcal H \to 0.
$ 이것은 우리에게 완전열을 제공합니다 전층 $$
0 \to \mathcal{F} \to \mathcal{G} \to \mathcal H \to
\mathcal H/^p\mathcal{G} \to 0,
$$ 여기서 $/^p$는 $\textit{PAb}(S_\etale)$의 몫을 나타냅니다. $\overline{s}$에서 줄기를 취하면 $(\mathcal H /^p\mathcal{G})_{\bar{s}} =
(\mathcal H /\mathcal{G})_{\bar{s}} = 0$를 알 수 있다. $\mathcal H/^p\mathcal{G}$의 층화가 $0$이기 때문이다. 따라서 줄기를 취하는 것은 전층에서 함자와 동일하므로 $$
0\to \mathcal{F}_{\overline{s}} \to \mathcal{G}_{\overline{s}} \to
\mathcal{H}_{\overline{s}} \to 0 = (\mathcal H/^p\mathcal{G})_{\overline{s}}
$$는 정확한다. \end{proof}

\begin{theorem}
\label{etale-cohomology-theorem-exactness-stalks}
$S$를 스킴으로 설정한다.
집합의 층의 지도 $a : \mathcal{F} \to \mathcal{G}$는 단사적이다.
(각각.\ 전사) 필요충분조건은 줄기의 지도
$a_{\overline{s}} : \mathcal{F}_{\overline{s}} \to \mathcal{G}_{\overline{s}}$
주사형이다(각각.\ 전사) 모두를 위해기하점~의$S$.
$S_\etale$의 아벨 층 시퀀스는 정확한다.
필요충분조건은 $S$의 기하점에 있는 모든 줄기에 대해 정확한다.
\end{theorem}

\begin{proof} 줄기에서 완전성의 필요성은 보조정리 \ref{etale-cohomology-lemma-stalk-exact}에서 따릅니다. 반대로, 층의 맵이 전사(각각 단사) 필요충분조건은 모든 줄기에 대해 전사(각각 단사)임을 보여주는 것으로 충분한다. 전사의 경우 이를 증명하고, 주입의 경우 증명을 생략한다.

\medskip\noindent
$\alpha : \mathcal{F} \to \mathcal{G}$를 층의 맵으로 설정한다.
그 $\mathcal{F}_{\overline{s}} \to \mathcal{G}_{\overline{s}}$
모든 기하점에 대해 전사이다. 고치다
$U \in \Ob(S_\etale)$
및 $s \in \mathcal{G}(U)$. 모든 $u \in U$에 대해 일부를 선택하세요.
$\overline{u} \to U$는 $u$와 에탈 동네 위에 누워있다.
$(V_u , \overline{v}_u) \to (U, \overline{u})$ 이렇게
$s|_{V_u} = \alpha(s_{V_u})$ 일부에게는
$s_{V_u} \in \mathcal{F}(V_u)$.
이는 $\alpha$가 전사이기 때문에 가능한다.
줄기. 그런 다음 $\{V_u \to U\}_{u \in U}$
는 $s$의 제한이 있는 에탈 피복이다.
지도 $\alpha$의 이미지에 있다.
따라서 $\alpha$는 전사이다.
사이트, 절 \ref{sites-section-sheaves-injective}.
\end{proof}

\begin{remarks}
\label{etale-cohomology-remarks-enough-points}
기하학적인 점 사이트.
\begin{enumerate}
\item 정리 \ref{etale-cohomology-theorem-exactness-stalks}는 포인트 계열이
~의$S_\etale$에 의해 주어진기하점~의$S$
(보조정리 \ref{etale-cohomology-lemma-stalk-gives-point})는 보수적이다.
사이트, 정의 \ref{sites-definition-enough-points}.
특히 $S_\etale$에는 충분한 점수가 있다.
\item $\mathcal{F}$가 큰 에탈 사이트에서 층라고 가정한다.
\label{etale-cohomology-item-stalks-big}
$S$. $T \to S$를 $S$의 큰 에탈 사이트의 객체로 두고,
$\overline{t}$를 $T$의 기하점으로 설정한다. 그런 다음 우리는 정의합니다
$\mathcal{F}_{\overline{t}}$를 줄기로
$\mathcal{F}|_{T_\etale}$ 제한 $\mathcal{F}$
$T$의 작은 에탈 사이트로. 즉, 우리는 다음과 같이 정의할 수 있다.
임의의 기하점에서 $\mathcal{F}$의 줄기
스킴 $T/S \in \Ob((\Sch/S)_\etale)$.
\item $S$의 큰 에탈 사이트도 점수가 충분한다.
이 사이트의 모든 개체 중 모든 기하점을 고려하면 다음을 참조하라.
(\ref{etale-cohomology-item-stalks-big}).
\end{enumerate}
\end{remarks}

\noindent 다음 보조정리는 첫 독해에서 건너뛰어야 한다.

\begin{lemma}
\label{etale-cohomology-lemma-points-small-etale-site}
$S$를 스킴으로 설정한다.
\begin{enumerate}
\item $p$를 작은 에탈 사이트의 점으로 하자.
$S_\etale$~의$S$에 의해 주어진함자
$u : S_\etale \to \textit{Sets}$.
그런 다음 존재한다.기하점 $\overline{s}$~의$S$그렇게
$p$는 다음과 연관된 $S_\etale$ 지점과 동형이다.
$\overline{s}$ 에
보조정리 \ref{etale-cohomology-lemma-stalk-gives-point}.
\item $p : \Sh(pt) \to \Sh(S_\etale)$를 점으로 둡니다.
$S$의 작은 에탈 토포스. 그러면 $p$는 기하점에서 나옵니다.
$S$, 즉 줄기 함자 $\mathcal{F} \mapsto \mathcal{F}_p$
에 정의된 대로 줄기 함자와 동형이다.
정의 \ref{etale-cohomology-definition-stalk}.
\end{enumerate}
\end{lemma}

\begin{proof}
에 의해
사이트, 보조정리 \ref{sites-lemma-point-site-topos}
사이트의 점과 점 사이에는 일대일 대응이 있다.
연관된 토포스이므로 (1)를 증명하는 것으로 충분한다.
에 의해
사이트, 명제 \ref{sites-proposition-point-limits}
함자 $u$에는 다음과 같은 성질이 있다.
() $u(S) = \{*\}$, (b) $u(U \times_V W) = u(U) \times_{u(V)} u(W)$ 및
(c) $\{U_i \to U\}$가 에탈 피복인 경우
$\coprod u(U_i) \to u(U)$는 전사이다.
특히, $U' \subset U$가 개방형 하위 구성표인 경우,
그런 다음 $u(U') \subset u(U)$. 게다가,
사이트, 보조정리 \ref{sites-lemma-point-site-topos}
$u(U) = p^{-1}(h_U^\#)$라고 쓸 수 있다. 즉, $u(U)$는
표현 가능한 층 $h_U$의 줄기. 만약에
$U = V \amalg W$, 그러면 $h_U = (h_V \amalg h_W)^\#$를 확인하고 다음을 얻습니다.
$u(U) = u(V) \amalg u(W)$ $p^{-1}$이 정확하기 때문이다.

\medskip\noindent $u$를 $S_{Zar}$로 제한하는 것을 고려하세요. 사이트, 예 \ref{sites-example-point-topological} 및 \ref{sites-example-point-topology}에 의해 고유한 포인트 $s \in S$가 존재하므로 $S' \subset S$ 열기에 대해 $u(S') = \{*\}$ 만약 $s \in S'$ 및 $u(S') = \emptyset$가 있다. $s \not \in S'$. $\varphi : U \to S$가 $S_\etale$의 객체인 경우 $\varphi(U) \subset S$는 열려 있고(명제 \ref{etale-cohomology-proposition-etale-morphisms} 참조) $\{U \to \varphi(U)\}$는 에탈 피복이다. 따라서 우리는 $u(U) = \emptyset \Leftrightarrow s \not \in \varphi(U)$라고 결론을 내립니다.

\medskip\noindent
고르다기하점 $\overline{s} : \overline{s} \to S$누워있다$s$, 보다
정의 \ref{etale-cohomology-definition-geometric-point}
표기법의 관례적인 남용. $\varphi : U \to S$가 객체라고 가정한다.
$S_\etale$와 $U$ 관계. $\varphi$는 분리되어 있다.
그올 $U_s$~의$\varphi$~ 위에$s$은아핀 스킴~ 위에
$\Spec(\kappa(s))$ 이는 유한 곱의 스펙트럼이다.
유한 분리 가능한 확장$k_i$~의$\kappa(s)$. 그러므로 우리는 신청할 수 있습니다
에탈 사상, 보조정리 \ref{etale-lemma-etale-etale-local-technical}
에탈 동네를 얻으려면$(V, \overline{v})$~의$(S, \overline{s})$
그렇게
$$
U \times_S V = U_1 \amalg \ldots \amalg U_n \amalg W
$$
$U_i \to V$, 동형사상 및 $W$에는 아무런 의미가 없다.
$\overline{v}$. 따라서 우리는 다음과 같이 결론을 내립니다.
$$
u(U) \times u(V) =
u(U \times_S V) =
u(U_1) \amalg \ldots \amalg u(U_n) \amalg u(W)
$$
물론 $u(U_i) = u(V)$도 있다. $V$를 조금 줄인 후에는
$V$에는 $s$ 위에 정확히 하나의 점이 있다고 가정한다. 따라서 $W$에는
$s$ 위에 있는 지점이다. 위와 같이 하면 $u(W) = \emptyset$가 된다.
그러므로 우리는 얻는다
$$
u(U) \times u(V) =
u(U_1) \amalg \ldots \amalg u(U_n) =
\coprod\nolimits_{i = 1, \ldots, n} u(V)
$$
여기서 등호는 집합 $u(V)$에 대한 맵과 호환된다.
참고하세요$u(V) \not = \emptyset$~처럼$s$의 이미지에 있습니다$V \to S$.
특히, 이 상황에서 $u(U)$는 유한하다는 것을 알 수 있다.
집합 $n$ 요소가 있는 경우, 즉 왼쪽에는 올 $u(U)$가 있다.
$u(V)$의 모든 요소와 오른쪽에는 올 카디널리티 $n$가 있다.

\medskip\noindent
극한을 고려해보세요.
$$
\lim_{(V, \overline{v})} u(V)
$$
에탈 동네 $(V, \overline{v})$의 범주에 대해
$\overline{s}$. 우리가 취할 때 동일한 가치를 얻는 것이 분명한다.
$(V, \overline{v})$의 하위 범주에 대한 극한은 $V$ 아핀을 사용한다.
이전 단락에 따라 ($V$ 역할과 $U$ 역할이 전환되어 적용됨)
이 경우 $u(V)$는 항상 비어 있지 않은 유한 집합임을 알 수 있다.
또한 극한은 공동 필터링된다.
보조정리 \ref{etale-cohomology-lemma-cofinal-etale}.
따라서
범주론, 절 \ref{categories-section-codirected-limits}
극한은 비어 있지 않습니다. 이 극한에서 요소 $x$를 선택한다.
이는 우리가 $x_{V, \overline{v}} \in u(V)$를 얻는다는 것을 의미한다.
에탈 동네마다$(V, \overline{v})$~의$(S, \overline{s})$
에탈 동네의 모든 사상에 대해
$\varphi : (V', \overline{v}') \to (V, \overline{v})$ 우리는
$u(\varphi)(x_{V', \overline{v}'}) = x_{V, \overline{v}}$.

\medskip\noindent
우리는 기능적 전단사 맵을 구성하기 위해 $x$의 선택을 사용할 것입니다
$$
c : |U_{\overline{s}}| \longrightarrow u(U)
$$
$U \in \Ob(S_\etale)$에 대해 증명을 마무리한다. 보다
보조정리 \ref{etale-cohomology-lemma-stalk-gives-point}
$|U_{\overline{s}}|$에 대한 설명은 증명이다.
먼저 주장 $U$ 아핀에 대한 맵을 구성하는 것으로 충분하다는 것을 확인했다.
이 주장의 증명을 생략한다.
추정하다$U \to S$~에$S_\etale$~와 함께$U$아핀하고하자
$\overline{u} : \overline{s} \to U$는 $|U_{\overline{s}}|$의 요소이다.
선택하세요$(V, \overline{v})$그렇게$U \times_S V$분해하다
증명의 세 번째 단락과 같습니다.
그런 다음 $(\overline{u}, \overline{v})$ 쌍은 다음의 기하점을 제공한다.
$U \times_S V$는 $\overline{v}$ 위에 놓여 있으며 다음 중 하나를 결정한다.
구성 요소$U_i$~의$U \times_S V$. 더 정확하게는 존재한다.
에이절 $\sigma : V \to U \times_S V$투영의$\text{pr}_U$
$(\overline{u}, \overline{v}) = \sigma \circ \overline{v}$와 같습니다. 집합
$c(\overline{u}) = u(\text{pr}_U)(u(\sigma)(x_{V, \overline{v}})) \in u(U)$.
이것이 $(V, \overline{v})$ 선택과 무관한지 확인해야 한다. 에 의해
보조정리 \ref{etale-cohomology-lemma-cofinal-etale}
에탈 동네의 범주가 공동 필터링된다.
그러므로 보여주기에 충분하다
에탈 동네의 사상이 주어진 것이다.
$\varphi : (V', \overline{v}') \to (V, \overline{v})$ 및 다음 중 하나를 선택한다.
절 $\sigma' : V' \to U \times_S V'$ 다음과 같은 투영
$(\overline{u}, \overline{v'}) = \sigma' \circ \overline{v}'$
$u(\sigma')(x_{V', \overline{v}'}) = u(\sigma)(x_{V, \overline{v}})$이 있다.
다이어그램을 고려하십시오
$$
\xymatrix{
V' \ar[d]^{\sigma'} \ar[r]_\varphi & V \ar[d]^\sigma \\
U \times_S V' \ar[r]^{1 \times \varphi} &
U \times_S V
}
$$
이제 이 다이어그램은 통근하는 경우가 아닐 수도 있다. 그 이유는
스킴 $V'$ 및 $V$는 연결이 아닐 수 있으므로
위의 $\sigma'$ 및 $\sigma$을 구성하는 데 사용된 분해는 다음과 같습니다.
독특하지 마십시오. 하지만 우리는 그걸 알고 있어요
$\sigma \circ \varphi \circ \overline{v}' =
(1 \times \varphi) \circ \sigma' \circ \overline{v}'$
건설로. 따라서 $U \times_S V$는 $S$에 대해 에탈이므로,
열린 동네가 있어요
$V'' \subset V'$ 의 $\overline{v'}$ 다이어그램은 다음과 같습니다.
$V''$로 제한된 경우 통근, 참조
사상, 보조정리 \ref{morphisms-lemma-value-at-one-point}.
이는 위의 다이어그램을 다음과 같이 확장할 수 있음을 의미한다.
$$
\xymatrix{
V'' \ar[r] \ar[d]^{\sigma'|_{V''}} &
V' \ar[d]^{\sigma'} \ar[r]_\varphi &
V \ar[d]^\sigma \\
U \times_S V'' \ar[r] &
U \times_S V' \ar[r]^{1 \times \varphi} &
U \times_S V
}
$$
왼쪽 정사각형과 외부 직사각형이 통근하도록 한다.
$u$는 함자이므로 이는 다음을 의미한다.
$x_{V'', \overline{v}'}$는 다음의 동일한 요소에 매핑된다.
$u(U \times_S V)$ 어떤 경로를 택하더라도
도표. 반면에 요소에 매핑된다.
$x_{V', \overline{v}'}$ 및 $x_{V, \overline{v}}$
$u(V')$ 및 $u(V)$. 이는 원하는 평등을 의미한다.
$u(\sigma')(x_{V', \overline{v}'}) = u(\sigma)(x_{V, \overline{v}})$.

\medskip\noindent 비슷한 방식으로 $c : |U_{\overline{s}}| \to u(U)$ 구성이 $U$에서 기능적임을 증명한다. 자세한 내용은 생략했다. 그리고 마지막으로 세 번째 문단의 결과를 통해 $c$ 맵은 보조정리의 증명으로 끝나는 전단사임을 알 수 있다. \end{proof}







\section{다른 위상의 점} \label{etale-cohomology-section-points-topologies}

\noindent 이 절에서는 스킴의 에탈 사이트 이외의 일부 사이트에 대한 포인트의 존재에 대해 간략하게 논의한다. 본 절에 사용된 용어는 사이트, 절 \ref{sites-section-sites-enough-points} 및 위상, 절 \ref{topologies-section-procedure} ff를 참조한다. 모든 기하학 사이트에는 충분한 점이 있다.

\begin{lemma}
\label{etale-cohomology-lemma-points-fppf}
$S$를 스킴으로 설정한다. 다음 사이트 모두 점수가 충분한다.
$S_{affine, Zar}$, $S_{Zar}$, $S_{affine, \etale}$, $S_\etale$,
$(\Sch/S)_{Zar}$, $(\textit{Aff}/S)_{Zar}$,
$(\Sch/S)_\etale$, $(\textit{Aff}/S)_\etale$,
$(\Sch/S)_{smooth}$, $(\textit{Aff}/S)_{smooth}$,
$(\Sch/S)_{syntomic}$, $(\textit{Aff}/S)_{syntomic}$,
$(\Sch/S)_{fppf}$ 및 $(\textit{Aff}/S)_{fppf}$.
\end{lemma}

\begin{proof} 각각의 큰 사이트에 대해 연관된 토포스는 사이트 $(\textit{Aff}/S)_\tau$에 의해 정의된 토포스와 동일한다. 위상, 기본 정리 \ref{topologies-lemma-affine-big-site-Zariski}를 참조하라. \ref{topologies-lemma-affine-big-site-etale}, \ref{topologies-lemma-affine-big-site-smooth}, \ref{topologies-lemma-affine-big-site-syntomic}, 그리고 \ref{topologies-lemma-affine-big-site-fppf}. 사이트 $(\textit{Aff}/S)_\tau$에 대한 결과는 Deligne의 결과 사이트, 보조정리 \ref{sites-lemma-criterion-points}에서 바로 이어집니다.

\medskip\noindent $S_{Zar}$에 대한 결과는 명확한다. $S_{affine, Zar}$에 대한 결과는 Deligne의 결과를 따릅니다. $S_\etale$에 대한 결과는 증명 의) 정리 \ref{etale-cohomology-theorem-exactness-stalks} 또는 위상에서 보조정리 \ref{topologies-lemma-alternative} 및 Deligne의 결과가 $S_{affine, \etale}$에 적용되었다. \end{proof}

\noindent
위의 보조정리는 포인트의 존재를 보장하지만 그렇지 않습니다.
이 점이 어떻게 생겼는지 알려주세요. 우리는 명시적으로 구성할 수 있다.
{\it 일부} 점은 다음과 같습니다.
$\overline{s} : \Spec(k) \to S$가 기하학이라고 가정한다.
$k$ 대수적으로 닫힌 점. 함자를 고려해보세요.
$$
u : (\Sch/S)_{fppf} \longrightarrow \textit{Sets},
\quad
u(U) = U(k) = \Mor_S(\Spec(k), U).
$$
$U \mapsto U(k)$는 다음과 같이 유한 극한으로 통근한다.
$S(k) = \{\overline{s}\}$ 그리고
$(U_1 \times_U U_2)(k) = U_1(k) \times_{U(k)} U_2(k)$.
게다가, $\{U_i \to U\}$가 fppf 피복라면,
$\coprod U_i(k) \to U(k)$는 전사이다.
에 의해
사이트, 명제 \ref{sites-proposition-point-limits}
우리는 그것을 본다$u$점을 정의합니다$p$~의$(\Sch/S)_{fppf}$~와 함께
줄기들
$$
\mathcal{F}_p = \colim_{(U, x)} \mathcal{F}(U)
$$
여기서 여극한은 평소와 같이 $U \to S$, $x \in U(k)$ 쌍 위에 있다.
하지만... 이 범주에는 초기 개체가 있다.
$(\Spec(k), \text{id})$, 따라서 우리는
$$
\mathcal{F}_p = \mathcal{F}(\Spec(k))
$$
별로 흥미롭지 않아요! 실제로 일반적으로 이러한 점은 그렇지 않습니다.
보수적인 포인트군을 형성한다. 좀 더 흥미로운 유형의 포인트
다음 비고에 설명되어 있다.

\begin{remark}
\label{etale-cohomology-remark-points-fppf-site}
\begin{reference}
이에 대해서는 \cite{Schroeer}에서 논의된다.
\end{reference}
$S = \Spec(A)$를 아핀 스킴으로 설정한다. $(p, u)$를 포인트로 삼자
사이트 $(\textit{Aff}/S)_{fppf}$, 참조
사이트, 절 \ref{sites-section-points} 및
\ref{sites-section-construct-points}. $B = \mathcal{O}_p$를 줄기로 설정한다.
$p$ 지점에서 구조층의 그것을 기억해
$$
B = \colim_{(U, x)} \mathcal{O}(U) =
\colim_{(\Spec(C), x_C)} C
$$
여기서 $x_C \in u(\Spec(C))$. 그런 일이 일어날 수 있어요
$\Spec(B)$는 $(\textit{Aff}/S)_{fppf}$의 객체이다.
그리고 다음에 매핑되는 요소 $x_B \in u(\Spec(B))$가 있다.
호환 시스템 $x_C$. 이 경우 이웃 시스템은
초기 객체가 있으며 다음과 같습니다.
$\mathcal{F}_p = \mathcal{F}(\Spec(B))$ 모든 층 $\mathcal{F}$에 대해
$(\textit{Aff}/S)_{fppf}$에 있다. 그것은 간단하다
$\mathcal{F} \mapsto \mathcal{F}(\Spec(B))$가 점을 정의하는지 확인하려면
$\Sh((\textit{Aff}/S)_{fppf})$의 다음
$B$는 모든 충실히 평탄에 대해 로컬 $A$ 대수여야 한다.
유한하게 제시된 환 지도 $B \to B'$ 절 $B' \to B$가 있다.
반대로, $A$-대수 $B$의 경우 함자
$\mathcal{F} \mapsto \mathcal{F}(\Spec(B))$는 줄기 함자이다.
한 점. 자세한 내용은 생략했다. 일반적인 내용이 무엇인지는 명확하지 않습니다.
사이트 $(\textit{Aff}/S)_{fppf}$ 처럼 보이다.
\end{remark}











\section{아벨 층} \label{etale-cohomology-section-support} 지원

\noindent 먼저 로컬 절 지원에 대해 이야기한다.

\begin{lemma}
\label{etale-cohomology-lemma-support-subsheaf-final}
$S$를 스킴으로 설정한다. $\mathcal{F}$를 최종의 하위 묶음으로 설정한다.
$S$의 에탈 토포스의 객체(참조
사이트, 예 \ref{sites-example-singleton-sheaf}).
그런 다음 고유한 개방형이 존재한다.
$W \subset S$ $\mathcal{F} = h_W$와 같습니다.
\end{lemma}

\begin{proof}
조건은 $\mathcal{F}(U)$이 싱글톤이거나
모두를 위해 비어있다$\varphi : U \to S$~에$\Ob(S_\etale)$.
특히 로컬 절은 항상 붙이다. 만약에
$\mathcal{F}(U) \not = \emptyset$, 그런 다음
$\mathcal{F}(\varphi(U)) \not = \emptyset$ 왜냐면
$\{\varphi : U \to \varphi(U)\}$은 피복이다.
그러므로 우리는 취할 수 있습니다
$W = \bigcup_{\varphi : U \to S, \mathcal{F}(U) \not = \emptyset} \varphi(U)$.
\end{proof}

\begin{lemma}
\label{etale-cohomology-lemma-zero-over-image}
$S$를 스킴으로 설정한다.
$\mathcal{F}$를 $S_\etale$에서 아벨 층으로 설정한다.
$\sigma \in \mathcal{F}(U)$를 로컬 절로 설정한다.
다음과 같은 열린 부분집합 $W \subset U$가 존재한다.
\begin{enumerate}
\item $W \subset U$는 $U$ 중에서 가장 큰 자리스키 열린 부분집합이다.
그 $\sigma|_W = 0$,
\item 모든 $\varphi : V \to U$ 중 $S_\etale$에 속하는 것에 대하여
$$
\sigma|_V = 0 \Leftrightarrow \varphi(V) \subset W,
$$
\item 모든 기하점 $\overline{u}$ 중 $U$의 것에 대하여
$$
(U, \overline{u}, \sigma) = 0\text{(}\mathcal{F}_{\overline{s}}\text{에서)}
\Leftrightarrow
\overline{u} \in W
$$
여기서 $\overline{s} = (U \to S) \circ \overline{u}$.
\end{enumerate}
\end{lemma}

\begin{proof}
에탈 위상에서 $\mathcal{F}$는 층이므로
$\mathcal{F}$ ~ $U_{Zar}$는 자리스키 위상의 $U$에 있는 층이다.
따라서 자리스키 오픈(Zariski 열린)이 존재한다.$W$데성질 (1), 보다
가군, 보조정리 \ref{modules-lemma-support-section-closed}. 허락하다
$\varphi : V \to U$는 $S_\etale$의 화살표이다. 참고하세요
$\varphi(V) \subset U$는 열린 부분집합이고
$\{V \to \varphi(V)\}$은 에탈 피복이다. 그러므로 만일
$\sigma|_V = 0$, 그런 다음 층 조건에 대해 $\mathcal{F}$에 의해 우리는
$\sigma|_{\varphi(V)} = 0$를 참조하라. 이는 (2)를 증명한다.
(3)를 증명하려면 $(U, \overline{u}, \sigma)$
$\mathcal{F}_{\overline{s}}$의 0 요소를 정의한 다음
$\overline{u} \in W$. 이는 가정이 의미하기 때문에 사실이다.
에탈 동네의 사상이 존재한다.
$(V, \overline{v}) \to (U, \overline{u})$ 이렇게
$\sigma|_V = 0$. 따라서 (2)에 의해 $V \to U$가 $W$에 매핑되는 것을 볼 수 있다.
그러므로 $\overline{u} \in W$.
\end{proof}

\noindent $S$를 스킴으로 둡니다. $s \in S$로 놔두세요. $\mathcal{F}$를 $S_\etale$에서 층으로 설정한다. 비고 \ref{etale-cohomology-remark-map-stalks}에 의해 층 $\mathcal{F}$의 줄기의 동형사상 클래스가 $s$ 위에 있는 기하점에서 잘 정의된다.

\begin{definition}
\label{etale-cohomology-definition-support}
$S$를 스킴으로 설정한다.
$\mathcal{F}$를 $S_\etale$에서 아벨 층으로 설정한다.
\begin{enumerate}
\item {\it 지지집합은 $\mathcal{F}$}의 집합이다.
포인트 $s \in S$ $\mathcal{F}_{\overline{s}} \not = 0$
(일부) 기하점 $\overline{s}$가 $s$ 위에 놓여 있다.
\item $\sigma \in \mathcal{F}(U)$를 절로 둡니다.
{\it 지지집합의 $\sigma$}는 닫힌 부분집합 $U \setminus W$이다. 여기서
$W \subset U$는 $U$ 중 $\sigma$ 중 가장 큰 열린 부분집합이다.
0으로 제한됩니다(참조
보조정리 \ref{etale-cohomology-lemma-zero-over-image}).
\end{enumerate}
\end{definition}

\noindent 일반적으로 아벨 층의 지지집합은 닫혀 있지 않습니다. 예에 대해 $S = \mathbf{A}^1_{\mathbf{C}}$라고 가정한다. $i_t : \Spec(\mathbf{C}) \to S$를 점 $t \in \mathbf{C}$의 포함으로 둡니다. 나중에 $\mathcal{F}_t = i_{t, *}(\mathbf{Z}/2\mathbf{Z})$가 아벨 층이고 지지집합이 정확히 $\{t\}$인 것을 볼 수 있다. 절 \ref{etale-cohomology-section-closed-immersions}를 참조하라. 그러면 $$
\bigoplus\nolimits_{n \in \mathbf{N}} \mathcal{F}_n
$$는 지지집합 $\{1, 2, 3, \ldots\} \subset S$가 있는 아벨 층이다. 줄기를 타고 여극한으로 통근하기 때문에 이는 사실이다. 보조정리 \ref{etale-cohomology-lemma-stalk-exact}를 참조하라. 따라서 지지집합이 있는 아벨 층의 예는 닫히지 않습니다. 다음은 층 및 절 지원에 대한 몇 가지 기본 사실이다.

\begin{lemma} \label{etale-cohomology-lemma-support-section-closed} $S$를 스킴으로 둡니다. $\mathcal{F}$를 $S_\etale$에서 아벨 층으로 설정한다. $U \in \Ob(S_\etale)$ 및 $\sigma \in \mathcal{F}(U)$로 설정한다. \begin{enumerate} \item $\sigma$의 지지집합이 $U$에서 닫혀 있다. \item $\sigma + \sigma'$의 지지집합은 $\sigma, \sigma' \in \mathcal{F}(U)$의 지원 집합에 포함되어 있다. \item $\varphi : \mathcal{F} \to \mathcal{G}$가 $S_\etale$의 아벨 층의 맵인 경우, $\varphi(\sigma)$의 지지집합은 $\sigma \in \mathcal{F}(U)$의 지지집합에 포함된다. \item $\mathcal{F}$의 지지집합은 $\mathcal{F}$의 모든 로컬 절 지원 이미지의 통합이다. \item $\mathcal{F} \to \mathcal{G}$가 전사인 경우 $\mathcal{G}$의 지지집합은 $\mathcal{F}$의 지지집합의 하위 집합이다. \item $\mathcal{F} \to \mathcal{G}$가 단사인 경우 $\mathcal{F}$의 지지집합은 $\mathcal{G}$의 지지집합의 하위 집합이다. \end{enumerate} \end{lemma}

\begin{proof} 부품(1)은 정의에 의해 유지된다. 부품 (2) 및 (3)는 $\mathcal{F}$ 및 $\mathcal{G}$ ~ $U_{Zar}$의 제한 사항에 대해 유지되므로 유지된다. 가군, 보조정리 \ref{modules-lemma-support-section-closed}를 참조하라. 부분(4)은 보조정리 \ref{etale-cohomology-lemma-zero-over-image} 부분(3)의 직접적인 결과이다. 부품 (5) 및 (6)는 다른 부품에서 이어집니다. \end{proof}

\begin{lemma} \label{etale-cohomology-lemma-support-sheaf-rings-closed} $S_\etale$의 환의 층의 지지집합이 닫혔습니다. \end{lemma}

\begin{proof} 이것은 (우리 관례에 따르면) 환이 $0$ 필요충분조건은 $1 = 0$이기 때문에 사실이다. 따라서 지지집합은 층의 환은 절 유닛의 지지집합이다. \end{proof}




\section{Hensel 환}
\label{etale-cohomology-section-henselian-ring}

\noindent Stacks 프로젝트에서 이미 여러 번 사용된 정리를 언급하는 것으로 시작한다. 이 결과에는 여러 가지 버전이 있다. 여기서 우리는 대수적 버전을 언급한다.

\begin{theorem}
\label{etale-cohomology-theorem-quasi-finite-etale-locally}
$A\to B$는 유한 유형 환 맵이고 $\mathfrak p \subset A$는 소수라고 가정한다.
아이디얼. 그러면 에탈 환 맵 $A \to A'$과 소수가 존재한다.
$\mathfrak p' \subset A'$가 $\mathfrak p$ 위에 누워서
\begin{enumerate}
\item
$\kappa(\mathfrak p) = \kappa(\mathfrak p')$,
\item
$ B \otimes_A A' = B_1\times \ldots \times B_r \times C$,
\item
$ A'\to B_i$는 유한하고 고유한 소수 $q_i\subset B_i$가 존재한다.
$\mathfrak p'$ 이상, 그리고
\item 올의 모든 기약 구성 요소
$\Spec(C \otimes_{A'} \kappa(\mathfrak p'))$ 중 $C$ 이상 $\mathfrak p'$
차원 적어도 1을 가지고 있다.
\end{enumerate}
\end{theorem}

\begin{proof} 대수학, 보조정리 \ref{algebra-lemma-etale-makes-quasi-finite-finite} 또는 \cite[Th\'eor\`eme 18.12.1]{EGA4}를 참조하라. 스킴의 사상 관련 버전에 대해서는 사상, 절 \ref{more-morphisms-section-etale-localization}에 대한 추가 정보를 참조하라. \end{proof}

\noindent 헨젤의 보조정리를 떠올려보세요. 이 보조정리에는 여러 가지 버전이 있다. 다음은 두 가지입니다: \begin{enumerate} \item[(f)] 만약 $f\in \mathbf{Z}_p[T]$ 모닉 및 $f \bmod p = g_0 h_0$ 와 함께 $gcd(g_0, h_0) = 1$ 그러면 $f$는 $f = gh$ 와 함께 $\bar g = g_0$로 간주된다. $\bar h = h_0$, \item[(r)] 만약 $f \in \mathbf{Z}_p[T]$, 모닉 $a_0 \in \mathbf{F}_p$, $\bar f(a_0) =0$ 그러나 $\bar f'(a_0) \neq 0$ 그러면 존재한다 $a \in \mathbf{Z}_p$ 와 함께 $f(a) = 0$ 그리고 $\bar a = a_0$. \end{enumerate} 두 버전 모두 사실입니다(나중에 이에 대해 살펴보겠습니다). 첫 번째 버전은 동일소 부분으로의 인수분해를 요청하고, 두 번째 버전은 극대 아이디얼 모듈로 단순근의 리프트를 요청한다. 일반 국소환에 대해 이러한 조건을 요구하는 것은 동일하며 다른 많은 조건과 동일하다는 것이 밝혀졌습니다. 우리는 루트 리프팅 성질을 헨셀식 국소환의 정의로 사용한다. 이는 종종 확인하기 가장 쉬운 방법이기 때문이다.

%10.01.09
\begin{definition} \label{etale-cohomology-definition-henselian} (대수학, 정의 \ref{algebra-definition-henselian} 참조) 국소환 $(R, \mathfrak m, \kappa)$가 호출된다. {\it heenselian} 모든 $f \in R[T]$ 모닉에 대해, 모든 $a_0 \in \kappa$에 대해 $\bar f(a_0) = 0$ 및 $\bar f'(a_0) \neq 0$에 대해 다음과 같은 $a \in R$가 존재한다. $f(a) = 0$ 및 $a \bmod \mathfrak m = a_0$. \end{definition}

\noindent 염두에 두어야 할 헨셀리안 국소환의 좋은 예는 완성된 국소환이다. (대수학, 정의 \ref{algebra-definition-complete-local-ring}) 완전한 국소환은 국소환 $(R, \mathfrak m)$이므로 $R \cong \lim_n R/\mathfrak m^n$이다. 즉, $\mathfrak m$-adic에 대해 완전하고 분리되어 있다. 위상.

\begin{theorem} \label{etale-cohomology-theorem-hensel} 완료 국소환은 헨셀리안이다. \end{theorem}

\begin{proof} 뉴턴의 방법. 대수학, 보조정리 \ref{algebra-lemma-complete-henselian}를 참조하라. \end{proof}

\begin{theorem}
\label{etale-cohomology-theorem-henselian}
$(R, \mathfrak m, \kappa)$를 국소환으로 설정한다. 다음은 동일한다.
\begin{enumerate}
\item $R$는 헨젤리안이다.
\item 임의의 $f\in R[T]$ 및 인수분해 $\bar f = g_0 h_0$에 대해
$\kappa[T]$와 $\gcd(g_0, h_0)=1$에 인수분해 $f = gh$가 존재한다.
$R[T]$와 $\bar g = g_0$ 및 $\bar h = h_0$,
\item 임의의 유한 $R$-대수학 $S$는 다음의 유한 곱과 동형이다.
국소환 $R$에 대해 유한,
\item 유한 유형 $R$-대수 $A$는 곱과 동형이다.
$A \cong A' \times C$ 여기서 $A' \cong A_1 \times \ldots \times A_r$
유한 로컬 $R$-대수와 모든 기약의 곱이다.
$C \otimes_R \kappa$의 구성요소는 차원 적어도 1를 갖고,
\item 만약 $A$는 에탈 $R$-대수이고 $\mathfrak n$는 극대 아이디얼
$A$ $\mathfrak m$ 위에 놓여 $\kappa \cong A/\mathfrak n$, 그러면 거기에
동형사상 $\varphi : A \cong R \times A'$이 존재한다.
$\varphi(\mathfrak n) = \mathfrak m \times A' \subset R \times A'$.
\end{enumerate}
\end{theorem}

\begin{proof} 이것은 대수학 결과의 하위 집합인 보조정리 \ref{algebra-lemma-characterize-henselian}이다. 위의 부분 (5)은 대수학의 부분 (8), 보조정리 \ref{algebra-lemma-characterize-henselian}에 해당하지만 약간 다르게 공식화된다. \end{proof}

\begin{lemma} \label{etale-cohomology-lemma-finite-over-henselian} $R$가 헨셀식이고 $A$가 유한 $R$-대수이면 $A$는 헨셀식 국소환의 유한 곱이다. \end{lemma}

\begin{proof} 대수학, 보조정리 \ref{algebra-lemma-finite-over-henselian}를 참조하라. \end{proof}

\begin{definition}
\label{etale-cohomology-definition-strictly-henselian}
국소환 $R$이 헨젤 환이고 그 잉여체가 분리폐체이면 이를 {\it 엄격 헨젤 환}이라 한다.
잔기 체는 분리 가능하게 닫혀 있다.
\end{definition}

\begin{example} \label{etale-cohomology-example-powerseries} $R = \mathbf{C}[[t]]$의 경우, 에탈 $R$-대수는 사소한 확장 $R \to R$와 확장 $R \to R[X, X^{-1}]/(X^n-t)$의 유한 곱이다. 후자는 $D(t) \subset \Spec(R)$ 오픈을 통해 인수이므로 어떤 에탈 피복도 피복 $\{\text{id} : \Spec(R) \to \Spec(R)\}$로 정제할 수 있다. 우리는 이것이 헨셀리안 환의 스펙트럼의 에탈 피복에 대한 다소 일반적인 사실임을 아래에서 볼 수 있다. 이는 엄격하게 Hensel 환의 스펙트럼 중 더 높은 에탈 코호몰로지가 0임을 보여줍니다. \end{example}

\begin{theorem} \label{etale-cohomology-theorem-henselization} $(R, \mathfrak m, \kappa)$는 국소환이고 $\kappa\subset\kappa^{sep}$는 분리 가능한 대수적 종결이라고 가정한다. 표준 평면 국소환 맵 $R \to R^h \to R^{sh}$가 존재한다. 여기서 \begin{enumerate} \item $R^h$, $R^{sh}$는 에탈 $R$-대수학의 여과 여극한이다. \item $R^h$는 헨셀식이고, $R^{sh}$는 엄격히 헨셀식이며, \item $\mathfrak m R^h$ (각각.\ $\mathfrak m R^{sh}$)는 다음의 극대 아이디얼이다. $R^h$ (각각.\ $R^{sh}$), \item $\kappa = R^h/\mathfrak m R^h$ 및 $\kappa^{sep} = R^{sh}/\mathfrak m R^{sh}$는 $\kappa$의 확장자이다. \end{enumerate} \end{theorem}

\begin{proof} $R^h$과 $R^{sh}$의 구조는 대수학, 정리 \ref{algebra-lemma-henselization} 및 \ref{algebra-lemma-strict-henselization}에 설명되어 있다. \end{proof}

\noindent
정리 \ref{etale-cohomology-theorem-henselization}에 구성된 환
각각 {\it 헨셀화}라고 하며
국소환 $R$의 {\it 엄격 헨젤화}라 한다. 다음을 보라.
대수학, 정의 \ref{algebra-definition-henselization}.
$R$의 성질 중 많은 부분이 (엄격한) 헨셀화에 반영된다.
대수학에 대한 자세한 내용을 참조하라.
절 \ref{more-algebra-section-permanence-henselization}.




\section{줄기 구조층} \label{etale-cohomology-section-stalks-structure-sheaf}의

\noindent 이 절에서 우리는 해당 ``보통'' 지점에서 국소환의 엄격한 헨셀화를 사용하여 기하학적 지점에서 구조층의 줄기를 식별한다.

\begin{lemma}
\label{etale-cohomology-lemma-describe-etale-local-ring}
\begin{slogan}
스킴의 구조층의 줄기
에탈 위상에는 엄격한 헨셀화가 적용된다.
\end{slogan}
$S$를 스킴으로 설정한다.
허락하다$\overline{s}$가 되다기하점~의$S$누워있다$s \in S$.
$\kappa = \kappa(s)$를 허용하고
$\kappa \subset \kappa^{sep} \subset \kappa(\overline{s})$ 표시
분리 가능한 대수적 종결$\kappa$~에$\kappa(\overline{s})$.
그런 다음 정식 식별이 있다.
$$
(\mathcal{O}_{S, s})^{sh}
\cong
(\mathcal{O}_S)_{\overline{s}}
$$
여기서 왼쪽은 국소환의 엄격한 헨젤화이다.
$\mathcal{O}_{S, s}$에 설명된 대로
정리 \ref{etale-cohomology-theorem-henselization}
오른쪽은 구조층의 줄기이다.
$\mathcal{O}_S$ 에 $S_\etale$ 에
기하점 $\overline{s}$.
\end{lemma}

\begin{proof}
$\Spec(A) \subset S$를 $s$의 유사 이웃이라고 가정한다.
$\mathfrak p \subset A$를 $s$에 대응하는 소 아이디얼로 둡니다.
이러한 선택을 통해 표준 동형사상이 있다.
$\mathcal{O}_{S, s} = A_{\mathfrak p}$ 및 $\kappa(s) = \kappa(\mathfrak p)$.
따라서 우리는
$\kappa(\mathfrak p) \subset \kappa^{sep} \subset \kappa(\overline{s})$.
그것을 기억해
$$
(\mathcal{O}_S)_{\overline{s}} =
\colim_{(U, \overline{u})} \mathcal{O}(U)
$$
여기서 극한은 $(S, \overline{s})$의 에탈 이웃 위에 있다.
에탈 동네 $(U, \overline{u})$에 의해 동종 시스템이 제공된다.
$U$는 유사하고 $U \to S$는 $\Spec(A)$를 통해 인수가 된다.
즉, 우리는
$$
(\mathcal{O}_S)_{\overline{s}} = \colim_{(B, \mathfrak q, \phi)} B
$$
여기서 여극한 이다 위의 에탈 $A$-algebras $B$ 소수가 부여됨
$\mathfrak q$가 $\mathfrak p$ 위에 누워 있고
$\kappa(\mathfrak p)$-대수학 지도
$\phi : \kappa(\mathfrak q) \to \kappa(\overline{s})$.
$\kappa(\mathfrak q)$는 유한하게 분리 가능하므로 주의하라.
$\kappa(\mathfrak p)$ $\phi$의 이미지는 $\kappa^{sep}$에 포함되어 있다.
이러한 번역을 통해 보조정리의 결과는 동일한다.
결과에
대수학, 보조정리 \ref{algebra-lemma-strict-henselization-different}.
\end{proof}

\begin{definition}
\label{etale-cohomology-definition-etale-local-rings}
$S$를 스킴으로 설정한다. $\overline{s}$를 $S$의 기하점으로 설정한다.
$s \in S$ 지점 위에 놓여 있다.
\begin{enumerate}
\item $S$의 $\overline{s}$에서의 {\it 에탈 국소환}
은줄기~의구조층 $\mathcal{O}_S$~에$S_\etale$
$\overline{s}$에 있다. 우리는 때때로 이것을
{\it $\mathcal{O}_{S, s}$} 상대의 엄격한 헨젤화
기하점 $\overline{s}$로.
사용된 표기법: $\mathcal{O}_{S, \overline{s}}^{sh}$.
\item {\it $\mathcal{O}_{S, s}$}의 헨젤화는 다음과 같습니다.
헨젤화국소환~의$S$~에$s$. 보다
대수학, 정의 \ref{algebra-definition-henselization},
그리고
정리 \ref{etale-cohomology-theorem-henselization}.
표기법: $\mathcal{O}_{S, s}^h$.
\item $S$의 $\overline{s}$에서의 {\it 엄격 헨젤화}는
스킴 $\Spec(\mathcal{O}_{S, \overline{s}}^{sh})$이다.
\item $S$의 $s$에서의 {\it 헨젤화}는 스킴
$\Spec(\mathcal{O}_{S, s}^h)$.
\end{enumerate}
\end{definition}

\noindent
$f : T \to S$를 스킴의 사상으로 설정한다. $\overline{t}$
가 되다기하점~의$T$이미지 포함$\overline{s}$~에$S$.
$t \in T$ 및 $s \in S$를 이미지로 설정한다.
그런 다음 표준 가환 도식을 얻습니다.
$$
\xymatrix{
\Spec(\mathcal{O}^{sh}_{T, \overline{t}}) \ar[r] \ar[d] &
\Spec(\mathcal{O}^h_{T, t}) \ar[r] \ar[d] &
T \ar[d]^f \\
\Spec(\mathcal{O}^{sh}_{S, \overline{s}}) \ar[r] &
\Spec(\mathcal{O}^h_{S, s}) \ar[r] &
S
}
$$
$T$ 및 $S$의 헨젤화 및 엄격한 헨젤화. 당신은 할 수 있습니다
$t$ 및 $s$의 유사 이웃을 선택하고 다음을 사용하여 이를 증명한다.
다음과 같이 주어진 (엄격한) 헨젤화의 기능
대수학, 기본형 \ref{algebra-lemma-henselian-functorial-improve} 및
\ref{algebra-lemma-strictly-henselian-functorial-improve}.

\begin{lemma}
\label{etale-cohomology-lemma-describe-henselization}
$S$를 스킴으로 설정한다. $s \in S$로 놔두세요. 그럼 우리는
$$
\mathcal{O}_{S, s}^h =
\colim_{(U, u)} \mathcal{O}(U)
$$
여기서 여극한은 필터링된 범주 위에 있다.
에탈 지역$(U, u)$~의$(S, s)$그렇게
$\kappa(s) = \kappa(u)$.
\end{lemma}

\begin{proof} 이 보조정리는 사상, 보조정리 \ref{more-morphisms-lemma-describe-henselization}에 대한 추가 정보의 사본이다. \end{proof}

\begin{remark} \label{etale-cohomology-remark-henselization-Noetherian} $S$를 스킴으로 둡니다. $s \in S$로 놔두세요. $S$가 로컬로 뇌터인 경우 $\mathcal{O}_{S, s}^h$도 뇌터이며 동일한 완비화를 갖습니다. $$
\widehat{\mathcal{O}_{S, s}} \cong \widehat{\mathcal{O}_{S, s}^h}.
$$ 특히, $\mathcal{O}_{S, s} \subset
\mathcal{O}_{S, s}^h \subset
\widehat{\mathcal{O}_{S, s}}$이다. $\mathcal{O}_{S, s}$의 헨셀화는 일반적으로 완비화보다 훨씬 작으며 성질의 많은 부분을 상속한다. 예의 경우 $\mathcal{O}_{S, s}$가 감소하면 $\mathcal{O}_{S, s}^h$도 감소하지만 완비화의 경우 일반적으로 그렇지 않습니다. 여기에 향후 참조를 삽입하라. \end{remark}

\begin{lemma} \label{etale-cohomology-lemma-etale-site-locally-ringed} $S$를 스킴으로 둡니다. 사이트 $S_\etale$에 구조층 $\mathcal{O}_S$가 부여된 작은 에탈은 로컬로 울리는 사이트이다. 가군에서 사이트, 정의를 참조하라. \ref{sites-modules-definition-locally-ringed}. \end{lemma}

\begin{proof} 이는 줄기 $(\mathcal{O}_S)_{\overline{s}} = \mathcal{O}^{sh}_{S, \overline{s}}$가 로컬이고 $S_\etale$에 충분한 포인트가 있기 때문에 다음과 같습니다. 보조정리 \ref{etale-cohomology-lemma-describe-etale-local-ring}, 정리 \ref{etale-cohomology-theorem-exactness-stalks} 및 설명을 참조하라 \ref{etale-cohomology-remarks-enough-points}. 사이트의 가군, 보조정리 \ref{sites-modules-lemma-locally-ringed-stalk} 및 \ref{sites-modules-lemma-ringed-stalk-not-zero}를 참조하라. 이는 작은 에탈 사이트가 로컬로 링이 발생함을 의미한다. \end{proof}






\section{작은 에탈의 기능성 토포스} \label{etale-cohomology-section-functoriality}

\noindent 지금까지 우리는 에탈 사이트의 기능, 즉 스킴의 사상이 주어졌을 때 무슨 일이 일어나는지에 대해 아직 논의하지 않았습니다. 정확한 공식 토론은 위상, 절 \ref{topologies-section-etale}에서 찾을 수 있다. 이번과 다음 절에서는 이 자료를 작은 에탈 사이트의 설정에서 구체적으로 간략하게 논의한다.

\medskip\noindent
$f : X \to Y$를 스킴의 사상으로 설정한다. 우리는 함자를 얻습니다.
\begin{equation}
\label{etale-cohomology-equation-functorial}
u : Y_\etale \longrightarrow X_\etale, \quad
V/Y \longmapsto X \times_Y V/X.
\end{equation}
이 함자에는 다음과 같은 중요한 성질이 있다.
\begin{enumerate}
\item $u(\text{최종 목적}) = \text{최종 목적}$,
\item $u$는 올곱을 유지한다.
\item $\{V_j \to V\}$가 $Y_\etale$의 피복인 경우
$\{u(V_j) \to u(V)\}$은 $X_\etale$의 피복이다.
\end{enumerate}
각각 확인하기 쉽습니다(생략). 결과적으로 우리는 무엇을 얻습니까?
이라고 불린다{\it 사상~의사이트}
$$
f_{small} : X_\etale \longrightarrow Y_\etale,
$$
보다
사이트, 정의 \ref{sites-definition-morphism-sites}
그리고
사이트, 명제 \ref{sites-proposition-get-morphism}.
추상적인 개념을 자세히 알 필요는 없다.
에탈 층 및 에탈 코호몰로지로 작업하려면
일반적으로 함자
$f_{small, *}$(직상) 및 $f_{small}^{-1}$(당김)
에탈 층에 대해 알아보고, 간단한 성질에 대해 알아보세요.
다음 절에서 성질에 대해 논의하겠지만,
때로는 증명을 위해 좀 더 추상적인 자료를 참조하기도 한다.
그것은 종종 그것을 증명하는 자연스러운 환경이다.


\section{직상}
\label{etale-cohomology-section-direct-image}

\noindent 전층의 직상을 정의해 보자.

\begin{definition}
\label{etale-cohomology-definition-direct-image-presheaf}
$f: X\to Y$를 스킴의 사상으로 설정한다.
$\mathcal{F} $를 집합의 전층을 $X_\etale$로 설정한다.
{\it 직상} 또는 {\it 직상} $\mathcal{F}$
($f$ 아래)는
$$
f_*\mathcal{F} : Y_\etale^{opp} \longrightarrow \textit{Sets}, \quad
(V/Y) \longmapsto \mathcal{F}(X \times_Y V/X).
$$
다른 것과 구별하기 위해 $f_* = f_{small, *}$를 쓰는 경우도 있다.
직상 함자 (예: 일반적인 Zariski 직상 또는 $f_{big, *}$).
\end{definition}

\noindent 에탈 사상의 기본 변화가 다시 에탈이므로 이는 잘 정의된 에탈 전층이다. 좀 더 범주적으로 말하면 $f_*\mathcal{F}$는 함자 $\mathcal{F} \circ u$의 구성이라고 할 수 있다. 여기서 $u$는 식(\ref{etale-cohomology-equation-functorial})과 같습니다. 이는 구성이 전층 $\mathcal{F}$에서 기능적이라는 것을 명확하게 하며 따라서 함자 $$
f_* = f_{small, *} :
\textit{PSh}(X_\etale)
\longrightarrow
\textit{PSh}(Y_\etale)
$$를 얻습니다. $\mathcal{F}$가 아벨 군의 전층인 경우 $f_*\mathcal{F}$는 아벨 군의 전층이기도 하며 이전과 같이 $$
f_* = f_{small, *} :
\textit{PAb}(X_\etale)
\longrightarrow
\textit{PAb}(Y_\etale)
$$를 얻습니다(즉, 정확히 동일한 규칙으로 정의됨).

\begin{remark}
\label{etale-cohomology-remark-direct-image-sheaf}
층의 직상이 층임을 주장한다.
즉, 만약$\{V_j \to V\}$에탈이다피복~에$Y_\etale$
그러면 $\{X \times_Y V_j \to X \times_Y V\}$는 에탈 피복이다.
$X_\etale$. 따라서 $\mathcal{F}$에 대한 층 조건
$\{X \times_Y V_i \to X \times_Y V\}$로
는 $f_*\mathcal{F}$에 대한 층 조건과 동일한다.
$\{V_i \to V\}$. 따라서 $\mathcal{F}$가 층인 경우
$f_*\mathcal{F}$.
\end{remark}

\begin{definition}
\label{etale-cohomology-definition-direct-image-sheaf}
$f: X\to Y$를 스킴의 사상으로 설정한다.
$\mathcal{F} $를 집합의 층을 $X_\etale$로 설정한다.
{\it 직상} 또는 {\it 직상} $\mathcal{F}$
($f$ 아래)는
$$
f_*\mathcal{F} : Y_\etale^{opp} \longrightarrow \textit{Sets}, \quad
(V/Y) \longmapsto \mathcal{F}(X \times_Y V/X)
$$
이는 층이다.
비고 \ref{etale-cohomology-remark-direct-image-sheaf}.
다른 것과 구별하기 위해 $f_* = f_{small, *}$를 쓰는 경우도 있다.
직상 함자 (예: 일반적인 Zariski 직상 또는 $f_{big, *}$).
\end{definition}

\noindent 위와 똑같은 논의가 적용되며 함자 $$
f_* = f_{small, *} :
\Sh(X_\etale)
\longrightarrow
\Sh(Y_\etale)
$$ 및 $$
f_* = f_{small, *} :
\textit{Ab}(X_\etale)
\longrightarrow
\textit{Ab}(Y_\etale)
$$를 {\it 직상}라고 다시 얻습니다.

\medskip\noindent 아벨 층의 함자 $f_*$는 왼쪽 완전이다. (호몰로지, 절 \ref{homology-section-functors} 아벨 범주 사이의 함자가 왼쪽 완전이 된다는 것이 무엇을 의미하는지 참조하라.) 즉, $0 \to \mathcal{F}_1 \to \mathcal{F}_2 \to \mathcal{F}_3$가 정확하다면 $X_\etale$, 그러면 모든 $U/X \in \Ob(X_\etale)$에 대해 아벨 군 $0 \to \mathcal{F}_1(U) \to \mathcal{F}_2(U) \to \mathcal{F}_3(U)$의 순서가 정확한다. 따라서 모든 $V/Y \in \Ob(Y_\etale)$에 대해 아벨 군 $0 \to f_*\mathcal{F}_1(V) \to f_*\mathcal{F}_2(V) \to f_*\mathcal{F}_3(V)$의 시퀀스는 정확한다. 이는 $U = X \times_Y V$의 이전 시퀀스이기 때문이다.

\begin{definition}
\label{etale-cohomology-definition-higher-direct-images}
$f: X \to Y$를 스킴의 사상으로 설정한다.
오른쪽 유도 함자 $\{R^pf_*\}_{p \geq 1}$
$f_* : \textit{Ab}(X_\etale) \to \textit{Ab}(Y_\etale)$
{\it 고차 직상}라고 한다.
\end{definition}

\noindent 고차 직상 및 해당 유도 범주 변형은 (여기에 향후 참조 삽입)에서 더 자세히 설명된다.



\section{역상}
\label{etale-cohomology-section-inverse-image}

\noindent 이 절에서는 작은 에탈 사이트에 대한 층의 철수에 대해 간략하게 논의한다. 이것의 정확한 구성은 위상, 절 \ref{topologies-section-etale}에 있다.

\begin{definition}
\label{etale-cohomology-definition-inverse-image}
$f: X\to Y$를 스킴의 사상으로 설정한다. {\it 역상} 또는
{\it 당김}\footnote{당김에는 $f^{-1}$ 표기법을 사용한다.
층의 집합 또는 층의 아벨 군, 그리고 $f^*$를 위해 예약한다.
링이 있는 사이트/토포스의 사상을 통해 가군층의 당김.}
함자는 함자이다.
$$
f^{-1} = f_{small}^{-1} :
\Sh(Y_\etale)
\longrightarrow
\Sh(X_\etale)
$$
그리고
$$
f^{-1} = f_{small}^{-1} :
\textit{Ab}(Y_\etale)
\longrightarrow
\textit{Ab}(X_\etale)
$$
왼쪽 수반부터 $f_* = f_{small, *}$까지이다. 따라서
$f^{-1}$의 특징은 다음과 같습니다.
$$
\Hom_{{\Sh(X_\etale)}} (f^{-1}\mathcal{G}, \mathcal{F})
=
\Hom_{\Sh(Y_\etale)} (\mathcal{G}, f_*\mathcal{F})
$$
기능적으로 모든 $\mathcal{F} \in \Sh(X_\etale)$에 대해
$\mathcal{G} \in \Sh(Y_\etale)$. 우리도 비슷하게
$$
\Hom_{{\textit{Ab}(X_\etale)}} (f^{-1}\mathcal{G}, \mathcal{F})
=
\Hom_{\textit{Ab}(Y_\etale)} (\mathcal{G}, f_*\mathcal{F})
$$
$\mathcal{F} \in \textit{Ab}(X_\etale)$ 및
$\mathcal{G} \in \textit{Ab}(Y_\etale)$.
\end{definition}

\noindent
그러한 수반이 존재한다는 것은 사소한 일이 아닙니다.
반면에 이는 상당히 일반적인 환경에 존재한다.
비고 \ref{etale-cohomology-remark-functoriality-general}
아래에. 일반 기계에서는 $f^{-1}\mathcal{G}$
전층에 연결된 층이다.
\begin{equation}
\label{etale-cohomology-equation-pullback}
U/X
\longmapsto
\colim_{U \to X \times_Y V} \mathcal{G}(V/Y)
\end{equation}
여기서 여극한은 범주 쌍 위에 있다.
$(V/Y, \varphi : U/X \to X \times_Y V/X)$.
이것을 적용하려면
사이트, 명제 \ref{sites-proposition-get-morphism}
방정식의 함자 $u$ (\ref{etale-cohomology-equation-functorial})
$u_s = (u_p\ )^\#$ 설명을 사용하세요.
사이트, 절 \ref{sites-section-continuous-functors} 및
\ref{sites-section-functoriality-PSh}.
우리는 때때로 당김에 대해 이 공식을 사용할 것이다.
기본 성질 중 일부를 증명하기 위해.

\begin{lemma} \label{etale-cohomology-lemma-stalk-pullback} $f : X \to Y$를 스킴의 사상으로 둡니다. \begin{enumerate} \item 함자 $f^{-1} : \textit{Ab}(Y_\etale) \to \textit{Ab}(X_\etale)$가 정확한다. \item 함자 $f^{-1} : \Sh(Y_\etale) \to \Sh(X_\etale)$는 정확한다. 즉, 유한 극한 및 여극한으로 통근한다. 범주, 정의 \ref{categories-definition-exact}를 참조하라. \item $\overline{x} \to X$를 기하점으로 둡니다. $\mathcal{G}$를 $Y_\etale$에서 층으로 설정한다. 그런 다음 표준 식별 $$
(f^{-1}\mathcal{G})_{\overline{x}} = \mathcal{G}_{\overline{y}}.
$$이 있으며 여기서 $\overline{y} = f \circ \overline{x}$이다. \item 모든 $V \to Y$ 에탈에 대해 $f^{-1}h_V = h_{X \times_Y V}$가 있다. \end{enumerate} \end{lemma}

\begin{proof}
집합의 층에 있는 $f^{-1}$의 완전성은 다음의 결과이다.
사이트, 명제 \ref{sites-proposition-get-morphism}
우리에게 적용함자 $u$방정식 (\ref{etale-cohomology-equation-functorial}).
실제로 철회의 완전성은 정의의 일부이다.
사상 또는 토포스(원하는 경우 사이트). 따라서 우리는 (2) 보유를 확인한다.
부분(1) 주어진 이후아벨 층 $\mathcal{G}$~에
$Y_\etale$
$f^{-1}\mathcal{F}$의 집합의 기본 층은 동일한다.
~처럼$f^{-1}$근본적인 것의층~의집합~의$\mathcal{F}$, 보다
사이트, 절 \ref{sites-section-sheaves-algebraic-structures}.
또한보십시오
가군 사이트, 보조정리 \ref{sites-modules-lemma-flat-pullback-exact}.
문헌에서 (1) 및 (2)는 때때로 (3)를 통해 추론된다.
정리 \ref{etale-cohomology-theorem-exactness-stalks}.

\medskip\noindent
부분(3)은 하락의 줄기에 대한 일반적인 사실이다.
사이트, 보조정리 \ref{sites-lemma-point-morphism-sites}.
우리는 또한 증명할 것입니다 (3) 바로 다음과 같습니다. 참고로
보조정리 \ref{etale-cohomology-lemma-stalk-exact}
줄기를 타고 층화로 통근한다.
이제 $f^{-1}\mathcal{G}$가 층라는 점을 기억하세요.
전층에 연관된
$$
U \longrightarrow \colim_{U \to X \times_Y V} \mathcal{G}(V),
$$
방정식(\ref{etale-cohomology-equation-pullback})을 참조하라.
따라서 우리는
\begin{align*}
(f^{-1}\mathcal{G})_{\overline{x}}
& = \colim_{(U, \overline{u})} f^{-1}\mathcal{G}(U) \\
& = \colim_{(U, \overline{u})}
\colim_{a : U \to X \times_Y V} \mathcal{G}(V) \\
& = \colim_{(V, \overline{v})} \mathcal{G}(V) \\
& = \mathcal{G}_{\overline{y}}
\end{align*}
세 번째 평등에서는 $(U, \overline{u})$ 쌍과 지도
$a : U \to X \times_Y V$은 $(V, a \circ \overline{u})$ 쌍에 해당한다.

\medskip\noindent
부분(4)은 여극한을 식별하여 비슷한 방식으로 증명할 수 있다.
$f^{-1}h_V$을 정의한다. 아니면 당신은 사용할 수 있습니다
요네다의 보조정리 (범주, 보조정리 \ref{categories-lemma-yoneda})
그리고 기능적 평등
$$
\Mor_{\Sh(X_\etale)}(f^{-1}h_V, \mathcal{F}) =
\Mor_{\Sh(Y_\etale)}(h_V, f_*\mathcal{F}) =
f_*\mathcal{F}(V) = \mathcal{F}(X \times_Y V)
$$
표현 가능한 전층은 층라는 사실과 결합된다. 또한보십시오
사이트, 보조정리 \ref{sites-lemma-pullback-representable-sheaf}
완전히 일반적인 결과를 얻으려면.
\end{proof}

\noindent
함자 $(f_*, f^{-1})$ 쌍은 작은 에탈 토포스의 사상을 정의한다.
$$
f_{small} :
\Sh(X_\etale)
\longrightarrow
\Sh(Y_\etale)
$$
층의 코호몰로지에 대한 많은 일반 사항은 토포스에 적용되며
토포스의 사상. 우리는 결과가 언제 나오는지 지적하려고 노력할 것이다.
일반적인 경우와 에탈 토포스에 특정한 경우이다.

\begin{remark}
\label{etale-cohomology-remark-functoriality-general}
보다 일반적으로 $\mathcal{C}_1, \mathcal{C}_2$를 사이트로 두고,
최종 개체와 올곱이 있다고 가정한다. 허락하다
$u: \mathcal{C}_2 \to \mathcal{C}_1$는 다음을 만족하는 함자이어야 한다.
\begin{enumerate}
\item $\{V_i \to V\}$가 $\mathcal{C}_2$의 피복인 경우
$\{u(V_i) \to u(V)\}$는 $\mathcal{C}_1$의 피복입니다(우리는
$u$는 {\it 연속})이고,
\item $u$는 유한 극한으로 통근합니다(즉, $u$는 왼쪽 완전이다. 즉,
$u$는 올곱 및 최종 객체를 보존한다.
\end{enumerate}
그러면 정의할 수 있다.
$f_*: \Sh(\mathcal{C}_1) \to \Sh(\mathcal{C}_2)$
$ f_* \mathcal{F}(V) = \mathcal{F}(u(V))$로.
게다가 완전 함자 $f^{-1}$가 존재하는데,
왼쪽 수반 ~ $f_*$이다.
사이트, 정의 \ref{sites-definition-morphism-sites} 및
명제 \ref{sites-proposition-get-morphism}.
경고: 단순히 $u$가 연속적이라고 요구하는 것만으로는 충분하지 않습니다.
그리고 토포스의 사상을 얻기 위해 올곱으로 출퇴근한다.
\end{remark}




\section{큰 기능 토포스} \label{etale-cohomology-section-functoriality-big-topoi}

\noindent 사상 의 스킴 $f : X \to Y$가 주어지면 $f$와 연관된 토포스의 사상의 전체 호스트가 있다. 위상, 절을 참조하라. \ref{topologies-section-change-topologies} 목록이다. 아마도 가장 많이 사용되는 것은 사상 의 토포스 $$
f_{big} = f_{big, \tau} :
\Sh((\Sch/X)_\tau)
\longrightarrow
\Sh((\Sch/Y)_\tau)
$$ 여기서 $\tau \in \{Zariski, \etale, smooth, syntomic, fppf\}$일 것이다. 이는 각각 최종 객체인 올곱 및 피복을 보존하는 연속적인 함자 $$
(\Sch/Y)_\tau \longrightarrow (\Sch/X)_\tau, \quad
V/Y \longmapsto X \times_Y V/X
$$에 해당하므로 사이트 $$
f_{big} : (\Sch/X)_\tau \longrightarrow (\Sch/Y)_\tau.
$$의 사상을 정의한다. 위상을 참조하라. 절 \ref{topologies-section-zariski}, \ref{topologies-section-etale}, \ref{topologies-section-smooth}, \ref{topologies-section-syntomic}, \ref{topologies-section-fppf}. 특히 $f_{big}$에 따른 직상은 $$
(f_{big, *}\mathcal{F})(V/Y) = \mathcal{F}(X \times_Y V/X)
$$ 규칙에 의해 제공된다. 토포스의 사상은 설명하기 매우 쉬운 반전 이미지 함자 $f_{big}^{-1}$를 갖는 것으로 나타났습니다. 즉, $$
(f_{big}^{-1}\mathcal{G})(U/X) = \mathcal{G}(U/Y)
$$가 있다. 여기서 $U/Y$의 구조 사상은 구조 사상 $U \to X$와 $f$의 구성이다. 위상, 기본정리 \ref{topologies-lemma-morphism-big}를 참조하라. \ref{topologies-lemma-morphism-big-etale}, \ref{topologies-lemma-morphism-big-smooth}, \ref{topologies-lemma-morphism-big-syntomic}, 그리고 \ref{topologies-lemma-morphism-big-fppf}.







\section{기능성과 가군층} \label{etale-cohomology-section-morphisms-modules}

\noindent
이 절에서 우리는 설명된 자료 중 일부를 재구성할 것이다.
하강, 절 \ref{descent-section-quasi-coherent-sheaves},
\ref{descent-section-quasi-coherent-cohomology} 그리고
\ref{descent-section-quasi-coherent-sheaves-bis}
에탈 위상 설정에서. $f : X \to Y$를 사상으로 설정한다.
스킴. 우리는 위에서 보았습니다.
절 \ref{etale-cohomology-section-functoriality}, \ref{etale-cohomology-section-direct-image} 및
\ref{etale-cohomology-section-inverse-image}
이는 작은 에탈 사이트의 사상 $f_{small}$를 유발한다. ~ 안에
하강, 비고 \ref{descent-remark-change-topologies-ringed}
우리는 $f$도 자연 지도를 유도한다는 것을 확인했다.
$$
f_{small}^\sharp :
\mathcal{O}_{Y_\etale}
\longrightarrow
f_{small, *}\mathcal{O}_{X_\etale}
$$
층 중 환 중
$Y_\etale$ $(f_{small}, f_{small}^\sharp)$
링링된 사이트의 사상이다. 보다
가군 사이트, 정의 \ref{sites-modules-definition-ringed-site}
링이 있는 사이트의 사상의 정의에 대해.
여기서 $f_{small}^\sharp$는 다음과 같이 정의된다는 점을 기억해 보자.
호환되는 지도 시스템
$$
\text{pr}_V^\sharp : \mathcal{O}(V) \longrightarrow \mathcal{O}(X \times_Y V)
$$
$V$의 경우 $Y_\etale$의 개체에 따라 달라집니다.

\medskip\noindent 이 구성은 스킴의 사상 구성과 호환되는 것이 분명한다. 보다 정확하게는 $f : X \to Y$ 및 $g : Y \to Z$가 스킴의 사상인 경우 링 토포스의 사상로서 $$
(g_{small}, g_{small}^\sharp)
\circ
(f_{small}, f_{small}^\sharp)
=
((g \circ f)_{small}, (g \circ f)_{small}^\sharp)
$$를 갖습니다. 또한, 가군 위의 사이트, 정의 \ref{sites-modules-definition-pushforward}에 의해 사상 $f : X \to Y$ 스킴이 주어지면 잘 정의된 당김과 직상을 얻을 수 있음을 알 수 있다. 함자 \begin{align*}
f_{small}^* :
\textit{Mod}(\mathcal{O}_{Y_\etale})
\longrightarrow
\textit{Mod}(\mathcal{O}_{X_\etale}), \\
f_{small, *} :
\textit{Mod}(\mathcal{O}_{X_\etale})
\longrightarrow
\textit{Mod}(\mathcal{O}_{Y_\etale})
\end{align*} 일반적인 방식으로 인접해 있다. $g : Y \to Z$가 스킴의 또 다른 사상인 경우, 작곡에 대해 말한 내용 때문에 $(g \circ f)_{small}^* = f_{small}^* \circ g_{small}^*$ 및 $(g \circ f)_{small, *} = g_{small, *} \circ f_{small, *}$가 있다.

\medskip\noindent
범주 사이에는 꽤 많은 차이가 있다.
모두의$\mathcal{O}_X$ 가군~에$X$그리고범주모두 사이
$\mathcal{O}_{X_\etale}$-가군 $X_\etale$에 있다. 그러나
결과
하강, 절 \ref{descent-section-quasi-coherent-sheaves},
\ref{descent-section-quasi-coherent-cohomology} 그리고
\ref{descent-section-quasi-coherent-sheaves-bis}
준일관성을 고려하는 것 사이에는 큰 차이가 없다고 말해주세요.
$S$의 가군 및 $S_\etale$의 준연접 가군.
(우리는 이미 이것을 보았습니다.
정리 \ref{etale-cohomology-theorem-quasi-coherent}
예의 경우)
특히, $f : X \to Y$가 스킴 중 사상인 경우
당김 함자 $f_{small}^*$ 및 $f^*$ 일치
준연접층, 참조
하강,
명제 \ref{descent-proposition-equivalence-quasi-coherent-functorial}.
또한 $f$이 제공되는 직상의 경우에도 마찬가지이다.
준컴팩트 및 준분리형, 참조
하강, 보조정리 \ref{descent-lemma-higher-direct-images-small-etale}.

\medskip\noindent 큰 사이트에서 구조층의 기능에 대한 몇 마디이다. $f : X \to Y$를 스킴의 사상으로 설정한다. 위상 $\tau \in \{Zariski,\linebreak[0]
\etale,\linebreak[0] smooth,\linebreak[0] syntomic,\linebreak[0]
fppf\}$ 중 하나를 선택한다. 그러면 사상 $f_{big} : (\Sch/X)_\tau \to (\Sch/Y)_\tau$는 지도 $$
f_{big}^\sharp : \mathcal{O}_Y \longrightarrow f_{big, *}\mathcal{O}_X
$$에 의해 고리 모양의 사이트의 사상이 된다. 하강, 비고 \ref{descent-remark-change-topologies-ringed}를 참조하라. 실제로는 앞서 설명한 작은 사이트의 경우와 동일한 구성으로 주어진다.












\section{위상 비교} \label{etale-cohomology-section-compare-topologies}

\noindent 이 절에서는 다양한 위상과 관련하여 층을 비교할 때 어떤 일이 발생하는지 연구하기 시작한다.

\begin{lemma}
\label{etale-cohomology-lemma-where-sections-are-equal}
$S$를 스킴으로 설정한다. $\mathcal{F}$를 $S_\etale$에서 집합의 층으로 설정한다.
$s, t \in \mathcal{F}(S)$로 놔두세요. 그런 다음 열린 $W \subset S$가 존재한다.
다음과 같은 특징이 있다. 성질: 사상 $f : T \to S$
$W$ 필요충분조건은 $s|_T = t|_T$를 통한 요인 (제한 사항은
$f_{small}$)에 의해 철회된다.
\end{lemma}

\begin{proof}
$U \in \Ob(S_\etale)$에 빈 집합을 할당하는 전층을 고려하세요.
만약 $s|_U \not = t|_U$ 및 싱글톤 그 밖에는. 이것이 분명하다
$\Sh(S_\etale)$의 최종 객체의 하위 묶음이다. 에 의해
보조정리 \ref{etale-cohomology-lemma-support-subsheaf-final}
이 전층을 나타내는 열린 $W \subset S$를 찾습니다.
에 대한기하점 $\overline{x}$~의$S$우리는 그것을 본다$\overline{x} \in W$
필요충분조건은그만큼줄기~의$s$그리고$t$~에$\overline{x}$동의하다.
줄기에 대한 설명을 통해
보조정리 \ref{etale-cohomology-lemma-stalk-pullback}
$W$에 원하는 성질이 있음을 알 수 있다.
\end{proof}

\begin{lemma} \label{etale-cohomology-lemma-describe-pullback} $S$를 스킴으로 둡니다. $\tau \in \{Zariski, \etale\}$로 놔두세요. 위상의 사상 $$
\pi_S : (\Sch/S)_\tau \longrightarrow S_\tau
$$, 보조정리 \ref{topologies-lemma-at-the-bottom} 또는 \ref{topologies-lemma-at-the-bottom-etale}를 고려하세요. $\mathcal{F}$를 $S_\tau$에서 층으로 설정한다. 그런 다음 $\pi_S^{-1}\mathcal{F}$는 $$
(\pi_S^{-1}\mathcal{F})(T) = \Gamma(T_\tau, f_{small}^{-1}\mathcal{F})
$$ 규칙에 의해 제공된다. 여기서 $f : T \to S$이다. 또한, $\pi_S^{-1}\mathcal{F}$는 fpqc 피복에 대해 층 조건을 만족한다. \end{lemma}

\begin{proof} 사상 $i_f : \Sh(T_\tau) \to \Sh(\Sch/S)_\tau)$가 $\pi_S \circ i_f = f_{small}$와 같이 사상 $T_\tau \to S_\tau$인 것을 확인한다. 위상, 기본정리 \ref{topologies-lemma-put-in-T}, \ref{topologies-lemma-morphism-big-small}를 참조하라. \ref{topologies-lemma-put-in-T-etale} 및 \ref{topologies-lemma-morphism-big-small-etale}. 당김은 전이적이므로 $i_f^{-1} \pi_S^{-1}\mathcal{F} = f_{small}^{-1}\mathcal{F}$가 원하는 대로 표시된다.

\medskip\noindent
$\{g_i : T_i \to T\}_{i \in I}$를 fpqc 피복으로 설정한다.
최종 문장은 다음을 의미합니다: 층 $\mathcal{G}$
$T_\tau$ 및 주어진 절
$s_i \in \Gamma(T_i, g_{i, small}^{-1}\mathcal{G})$ 그의 철회
에게$T_i \times_T T_j$동의해, 독특한 게 있어절 $s$~의$\mathcal{G}$
$T$에 대해 $T_i$에 대한 당김이 $s_i$와 일치한다.

\medskip\noindent
$V \to T$를 $T_\tau$의 객체로 두고, $t \in \mathcal{G}(V)$로 둡니다.
모든 $i$에 대해 가장 큰 공개 $W_i \subset T_i \times_T V$가 있다.
$s_i$ 및 $t$의 하락이 하락의 절과 일치하도록
$\mathcal{G}$ ~ $W_i \subset T_i \times_T V$ 중, 참조
보조정리 \ref{etale-cohomology-lemma-where-sections-are-equal}.
$s_i$ 및 $s_j$가 $T_i \times_T T_j$에 동의하기 때문에
$W_i$ 및 $W_j$ 당김이 동일한 오픈 오버
$T_i \times_T T_j \times_T V$. 에 의해
하강, 보조정리 \ref{descent-lemma-open-fpqc-covering}
우리는 열린 곳을 찾았어요$W \subset V$누구의역상에게$T_i \times_T V$
$W_i$을(를) 복구한다.

\medskip\noindent
$g_{i, small}^{-1}\mathcal{G}$의 구성으로 존재한다.
 $\tau$-피복 $\{T_{ij} \to T_i\}_{j \in J_i}$, 각 $j$에 대해
열린 몰입 또는 에탈 사상 $V_{ij} \to T$, 절
$t_{ij} \in \mathcal{G}(V_{ij})$ 및 가환 도식
$$
\xymatrix{
T_{ij} \ar[r] \ar[d] & V_{ij} \ar[d] \\
T_i \ar[r] & T
}
$$
$s_i|_{T_{ij}}$은 $t_{ij}$의 철회이다. 다시 말해서,
피복 $\{T_i \to T\}$를 $\{T_{ij} \to T\}$로 교체한 후
우리는 $T_i \to V_i \to T$ 인수분해가 있다고 가정할 수 있다.
$V_i \in \Ob(T_\tau)$ 및 절 $t_i \in \mathcal{G}(V_i)$
뒤로 당기다$s_i$~ 위에$T_i$.
이전 단락의 결과로 $W_i \subset V_i$가 열립니다.
그렇게$t_i|_{W_i}$모두 '동의'$s_j$~ 위에$T_j \times_T W_i$.
$T_i \to V_i$는 $W_i$을 통해 고려된다.
따라서 $\{W_i \to T\}$는 $\tau$-피복이고 보조정리는 증명되었다.
\end{proof}

\begin{lemma} \label{etale-cohomology-lemma-sections-upstairs} $S$를 스킴으로 둡니다. $f : T \to S$를 사상으로 하여 \begin{enumerate} \item $f$가 평평하고 준컴팩트하며 \item $f$의 기하학적 올은 다음과 같습니다. 연결. \end{enumerate} $\mathcal{F}$를 $S_\etale$에서 층으로 설정한다. 그런 다음 $\Gamma(S, \mathcal{F}) = \Gamma(T, f^{-1}_{small}\mathcal{F})$. \end{lemma}

\begin{proof} 표준 맵 $\Gamma(S, \mathcal{F}) \to \Gamma(T, f_{small}^{-1}\mathcal{F})$이 있다. $f$는 전사이므로(올은 연결이므로) 이 지도는 단사임을 알 수 있다.

\medskip\noindent
지도가 전사임을 보여주기 위해,
$\alpha \in \Gamma(T, f_{small}^{-1}\mathcal{F})$.
$\{T \to S\}$는 fpqc 피복이므로 다음을 사용할 수 있다.
보조정리 \ref{etale-cohomology-lemma-describe-pullback} 그걸 보면 충분하다는 걸 증명할 수 있어요
$\alpha$는 $T \times_S T$에 대해 동일한 절로 되돌립니다.
두 가지 예측. $\overline{s} \to S$를 기하점으로 설정한다.
$(T \times_S T)_{\overline{s}}$에 대해 계약이 유지됨을 입증하는 것으로 충분한다.
$T \times_S T$의 모든 기하점은 다음 중 하나에 포함되어 있다.
이러한 기하학적인 올. 즉, 우리는 그것을 보여 주려고 노력하고 있습니다
$\alpha|_{T_{\overline{s}}}$는 동일한 절로 되돌아갑니다.
$$
(T \times_S T)_{\overline{s}} =
T_{\overline{s}} \times_{\overline{s}} T_{\overline{s}}
$$
$T_{\overline{s}}$에 대한 두 가지 투영에 의해.
그러나 $\mathcal{F}|_{T_{\overline{s}}}$는
$\mathcal{F}|_{\overline{s}}$의 철회는 상수층이다.
값 $\mathcal{F}_{\overline{s}}$. 이후 $T_{\overline{s}}$
가정에 의해 연결이고, 상수층의 모든 절은 상수이다.
따라서 $\alpha|_{T_{\overline{s}}}$는 요소에 해당한다.
$\mathcal{F}_{\overline{s}}$. 따라서 두 가지 철수는
$(T \times_S T)_{\overline{s}}$ 둘 다
이 동일한 요소에 해당하고 결론을 내립니다.
\end{proof}

\noindent 다음은 사상이 플랫하다고 가정하지 않는 보조정리 \ref{etale-cohomology-lemma-sections-upstairs}의 버전이다.

\begin{lemma} \label{etale-cohomology-lemma-sections-upstairs-submersive} $S$를 스킴으로 둡니다. $f : X \to S$를 사상으로 하여 \begin{enumerate} \item $f$가 잠수형이고 \item $f$의 기하학적인 올은 다음과 같습니다. 연결. \end{enumerate} $\mathcal{F}$를 $S_\etale$에서 층으로 설정한다. 그런 다음 $\Gamma(S, \mathcal{F}) = \Gamma(X, f^{-1}_{small}\mathcal{F})$. \end{lemma}

\begin{proof} 표준 맵 $\Gamma(S, \mathcal{F}) \to \Gamma(X, f_{small}^{-1}\mathcal{F})$이 있다. $f$는 전사이므로(올은 연결이므로) 이 지도는 단사임을 알 수 있다.

\medskip\noindent
지도가 전사임을 보여주기 위해,
$\tau \in \Gamma(X, f_{small}^{-1}\mathcal{F})$.
찾아보면 충분하다
에탈 피복 $\{U_i \to S\}$ 그리고
절 $\sigma_i \in \mathcal{F}(U_i)$
$\sigma_i$이 $\tau|_{X \times_S U_i}$로 되돌아갑니다.
즉, 위에 표시된 주입성은 다음을 보장한다.
$\sigma_i$ 및 $\sigma_j$는 동일하게 제한된다.
$\mathcal{F}$의 $U_i \times_S U_j$에 대한 절.
따라서 우리는 고유한 절 $\sigma \in \mathcal{F}(S)$를 얻습니다.
이는 다음으로 제한된다.$\sigma_i$~ 위에$U_i$.
그러면 $\sigma$에서 $X$로의 철회는 $\tau$이다.
왜냐하면 이것이 현지에서는 사실이기 때문이다.

\medskip\noindent
$\overline{x}$를 $X$ 이미지 $\overline{s}$의 기하점으로 설정한다.
$S$에서. 줄기에서 $\tau$의 이미지를 고려하세요.
$$
(f_{small}^{-1}\mathcal{F})_{\overline{x}} = \mathcal{F}_{\overline{s}}
$$
보조정리 \ref{etale-cohomology-lemma-stalk-pullback}를 참조하라.
에탈 동네를 찾을 수 있어요$U \to S$~의$\overline{s}$
그리고 이 이미지에 매핑되는 절 $\sigma \in \mathcal{F}(U)$
줄기에서. 따라서 $S$을 $U$로, $X$를 $X \times_S U$로 바꾼 후
우리는 거기에 종료가 있다고 가정할 수 있습니다절 $\sigma$~의$\mathcal{F}$~ 위에$S$
$(f_{small}^{-1}\mathcal{F})_{\overline{x}}$에 있는 이미지는
$\tau$와 같습니다.

\medskip\noindent
보조정리 \ref{etale-cohomology-lemma-where-sections-are-equal}에 의해
다음과 같은 최대 개방 $W \subset X$이 존재한다.
$f_{small}^{-1}\sigma$ 및 $\tau$는 $W$에 동의한다.
$W$의 형성은 추가 철수로 통근한다.
$\mathcal{F}$이(가) 뒤로 물러나는 것을 관찰하세요.
기하학 올 $X_{\overline{s}}$은 당김이다.
$\mathcal{F}_{\overline{s}}$의 층으로 표시됨
$\overline{s}$ 에 의해 $X_{\overline{s}} \to \overline{s}$.
따라서 $\tau$ 및 $\sigma$는 절을 제공한다.
값이 $\mathcal{F}_{\overline{s}}$인 상수층
$X_{\overline{s}}$에 대해 한 가지 점에서는 동의한다. 부터
$X_{\overline{s}}$는 가정에 의해 연결이므로 결론을 내립니다.
$W$에는 $X_s$이 포함되어 있다. 같은 주장이 다른데
기하학 올은 $W$가 충족하는 모든 올을 포함한다는 것을 보여줍니다.
$f$는 잠수적이므로 $W$는 반대라고 결론을 내립니다.
열린 동네의 이미지$s$~에$S$.
이로써 증명이 완료된다.
\end{proof}

\begin{lemma}
\label{etale-cohomology-lemma-sections-base-field-extension}
$K/k$는 $k$와 체의 확장으로 분리 가능한다.
대수적으로 닫혀있다. $S$를 $k$에 대한 스킴으로 설정한다. 표시하다
$p : S_K = S \times_{\Spec(k)} \Spec(K) \to S$ 투영이다.
$\mathcal{F}$를 $S_\etale$에서 층으로 설정한다.
그런 다음 $\Gamma(S, \mathcal{F}) = \Gamma(S_K, p^{-1}_{small}\mathcal{F})$.
\end{lemma}

\begin{proof}
보조정리 \ref{etale-cohomology-lemma-sections-upstairs}에서 이어집니다. 즉, 분명하다.
$p$은 기본 변경 사항과 같이 평평하고 준 컴팩트한다.
$\Spec(K) \to \Spec(k)$. 반면에 $\overline{s} : \Spec(L) \to S$
는기하점, 그런 다음올~의$p$~ 위에$\overline{s}$
는 $K \otimes_k L$의 스펙트럼이다. 이는 기약이므로 연결이다.
대수학, 보조정리 \ref{algebra-lemma-separably-closed-irreducible}.
\end{proof}









\section{사상} \label{etale-cohomology-section-morphisms} 복구 중

\noindent 이 절에서 우리는 스킴의 로컬 링 작은 에탈 토포스에 연결되는 규칙이 적절한 의미에서 완전히 충실하다는 것을 증명한다. 정리 \ref{etale-cohomology-theorem-fully-faithful}을 참조하라.

\begin{lemma}
\label{etale-cohomology-lemma-morphism-locally-ringed}
$f : X \to Y$를 스킴의 사상으로 설정한다.
울린 사이트 $(f_{small}, f_{small}^\sharp)$의 사상
$f$에 연결된 것은 로컬로 울리는 사이트의 사상이다.
가군 사이트에,
정의 \ref{sites-modules-definition-morphism-locally-ringed-topoi}.
\end{lemma}

\begin{proof}
우리가 다음을 보았으므로 이 주장은 의미가 있다는 점에 유의하라.
$(X_\etale, \mathcal{O}_{X_\etale})$ 그리고
$(Y_\etale, \mathcal{O}_{Y_\etale})$
로컬로 벨이 울립니다 사이트. 참조
보조정리 \ref{etale-cohomology-lemma-etale-site-locally-ringed}.
게다가, 우리는 $X_\etale$에 충분한 포인트가 있다는 것을 알고 있다.
정리 \ref{etale-cohomology-theorem-exactness-stalks}
그리고
비고 \ref{etale-cohomology-remarks-enough-points}.
따라서 $(f_{small}, f_{small}^\sharp)$을 증명하는 것으로 충분한다.
다음의 조건 (3)을 충족한다.
가군 사이트에,
보조정리 \ref{sites-modules-lemma-locally-ringed-morphism}.
이것을 보려면 요점이 필요한다.$p$~의$X_\etale$. 에 의해
보조정리 \ref{etale-cohomology-lemma-points-small-etale-site}
$p$에 해당기하점 $\overline{x}$~의$X$.
에 의해
보조정리 \ref{etale-cohomology-lemma-stalk-pullback}
$q = f_{small} \circ p$ 지점은
기하점 $\overline{y} = f \circ \overline{x}$의 $Y$.
따라서 우리가 증명해야 할 주장은 유도된 지도가
줄기
$$
(\mathcal{O}_Y)_{\overline{y}} \longrightarrow (\mathcal{O}_X)_{\overline{x}}
$$
국소환 지도이다. $a \in (\mathcal{O}_Y)_{\overline{y}}$
최대 요소에 매핑되는 왼쪽 요소이다.
오른쪽의 아이디얼. $a$이 동등 클래스라고 가정한다.
트리플의$(V, \overline{v}, a)$~와 함께$V \to Y$ 에탈,
$\overline{v} : \overline{x} \to V$ 오버 $Y$ 및 $a \in \mathcal{O}(V)$.
이는 동등한 클래스에 매핑된다.
$(X \times_Y V, \overline{x} \times \overline{v}, \text{pr}_V^\sharp(a))$
국소환 $(\mathcal{O}_X)_{\overline{x}}$에서. 그러나 분명한 것은
극대 아이디얼에 있다는 것은 $\text{pr}_V^\sharp(a)$를 뒤로 당기는 것을 의미한다.
$\kappa(\overline{x})$ 요소에 0을 제공한다. 그러므로 뒤로 당기는 것도
$a$ ~ $\kappa(\overline{x})$은 0이다. 이는 $a$가
극대 아이디얼 중 $(\mathcal{O}_Y)_{\overline{y}}$.
\end{proof}

\begin{lemma}
\label{etale-cohomology-lemma-2-morphism}
$X$, $Y$를 스킴으로 설정한다. $f : X \to Y$를 스킴의 사상으로 설정한다.
허락하다$t$가 되다$2$-사상~에서$(f_{small}, f_{small}^\sharp)$그 자체로, 참조
가군 사이트에,
정의 \ref{sites-modules-definition-2-morphism-ringed-topoi}.
그런 다음 $t = \text{id}$.
\end{lemma}

\begin{proof}
즉, $t : f^{-1}_{small} \to f^{-1}_{small}$
다음과 같은 함자의 변환이다.
$$
\xymatrix{
f_{small}^{-1}\mathcal{O}_Y
\ar[rd]_{f_{small}^\sharp}  & &
f_{small}^{-1}\mathcal{O}_Y \ar[ll]^t \ar[ld]^{f_{small}^\sharp} \\
& \mathcal{O}_X
}
$$
교환 가능한다. $V \to Y$가 $V$ 아핀과 함께 에탈이라고 가정한다. 에 의해
사상, 보조정리 \ref{morphisms-lemma-quasi-affine-finite-type-over-S}
우리는 몰입을 선택할 수 있습니다$i : V \to \mathbf{A}^n_Y$~ 위에$Y$.
층 측면에서 이는 $i$가 주입을 유도함을 의미한다.
층의 $h_i : h_V \to \prod_{j = 1, \ldots, n} \mathcal{O}_Y$.
베이스 변경$i'$~의$i$에게$X$몰입이다
(스킴, 보조정리 \ref{schemes-lemma-base-change-immersion}).
따라서 $i' : X \times_Y V \to \mathbf{A}^n_X$는 몰입이다.
차례로 다음을 의미한다.
$h_{i'} : h_{X \times_Y V} \to \prod_{j = 1, \ldots, n} \mathcal{O}_X$
층의 주입이다.
식별 $f_{small}^{-1}h_V = h_{X \times_Y V}$을 통해
보조정리 \ref{etale-cohomology-lemma-stalk-pullback}
지도 $h_{i'}$는 다음과 같습니다.
$$
\xymatrix{
f_{small}^{-1}h_V \ar[r]^-{f^{-1}h_i} &
\prod_{j = 1, \ldots, n} f_{small}^{-1}\mathcal{O}_Y
\ar[r]^{\prod f^\sharp} &
\prod_{j = 1, \ldots, n} \mathcal{O}_X
}
$$
(확인 생략). 즉 지도는
$t : f_{small}^{-1}h_V \to f_{small}^{-1}h_V$
가환 도식에 적합
$$
\xymatrix{
f_{small}^{-1}h_V \ar[r]^-{f^{-1}h_i} \ar[d]^t &
\prod_{j = 1, \ldots, n} f_{small}^{-1}\mathcal{O}_Y
\ar[r]^-{\prod f^\sharp} \ar[d]^{\prod t} &
\prod_{j = 1, \ldots, n} \mathcal{O}_X \ar[d]^{\text{id}}\\
f_{small}^{-1}h_V \ar[r]^-{f^{-1}h_i} &
\prod_{j = 1, \ldots, n} f_{small}^{-1}\mathcal{O}_Y
\ar[r]^-{\prod f^\sharp} &
\prod_{j = 1, \ldots, n} \mathcal{O}_X
}
$$
오른쪽 사각형의 교환성은 $t$에 대한 가정에 의해 유지된다.
위에서 설명했다.
가로 화살표의 구성은 주입식이므로
위의 논의를 통해 우리는 왼쪽 수직 화살표가
정체성 지도이기도 한다. 집합 중 모든 층
$Y_\etale$는 층의 (거대한) 부산물로부터의 전사를 인정한다.
$h_V$ 형식과 $V$ 아핀(결합
위상, 보조정리 \ref{topologies-lemma-alternative}
~와 함께
사이트, 보조정리 \ref{sites-lemma-sheaf-coequalizer-representable}).
따라서 우리는 $t : f_{small}^{-1} \to f_{small}^{-1}$
원하는대로 신원 변환이다.
\end{proof}

\begin{lemma}
\label{etale-cohomology-lemma-faithful}
$X$, $Y$를 스킴으로 설정한다.
스킴 중 임의의 두 사상 $a, b : X \to Y$
$2$-동형사상이 존재하는 경우
$(a_{small}, a_{small}^\sharp) \cong (b_{small}, b_{small}^\sharp)$
링형 토포스의 $2$-범주는 동일한다.
\end{lemma}

\begin{proof} 이 부분은 다소 혼란스럽기 때문에 신중하게 논의해 보자. $t : a_{small}^{-1} \to b_{small}^{-1}$를 $2$-동형사상으로 설정한다. 열린 $V \subset Y$을 고려하라. $h_V$는 최종 층 $*$의 하위 묶음이다. 따라서 $a_{small}^{-1}h_V = h_{a^{-1}(V)}$와 $b_{small}^{-1}h_V = h_{b^{-1}(V)}$는 모두 최종 층의 서브시브이다. 따라서 동형사상 $$
t : a_{small}^{-1}h_V = h_{a^{-1}(V)} \to b_{small}^{-1}h_V = h_{b^{-1}(V)}
$$는 신원이어야 하고, $a^{-1}(V) = b^{-1}(V)$여야 한다. $a$ 및 $b$는 기본 위상 공간에서 동일한다. 다음으로, 절 $f \in \mathcal{O}_Y(V)$를 선택하세요. 이는 집합 $f : h_V \to \mathcal{O}_Y$의 층 맵에 의해 결정되고 결정된다. 이것을 뒤로 당기고 $t$를 적용하여 가환 도식 $$
\xymatrix{
h_{b^{-1}(V)} \ar@{=}[r] &
b_{small}^{-1}h_V \ar[d]^{b_{small}^{-1}(f)} & &
a_{small}^{-1}h_V \ar[d]^{a_{small}^{-1}(f)} \ar[ll]^t &
h_{a^{-1}(V)} \ar@{=}[l]
\\
& b_{small}^{-1}\mathcal{O}_Y
\ar[rd]_{b^\sharp}  & &
a_{small}^{-1}\mathcal{O}_Y \ar[ll]^t \ar[ld]^{a^\sharp} \\
& & \mathcal{O}_X
}
$$를 얻습니다. 여기서 삼각형은 $2$-동형사상에서 가군의 정의로 교환 가능한다. 사이트, 절 \ref{sites-modules-section-2-category}. 위에서 우리는 상단 수평 화살표의 구성이 $a^{-1}(V) = b^{-1}(V)$ 항등에서 비롯된다는 것을 확인했다. 따라서 다이어그램의 교환성은 $a_{small}^\sharp(f) = b_{small}^\sharp(f)$가 $\mathcal{O}_X(a^{-1}(V)) = \mathcal{O}_X(b^{-1}(V))$에 있음을 알려줍니다. 이는 모든 열린 $V$ 및 모든 $f \in \mathcal{O}_Y(V)$에 대해 적용되므로 $a = b$는 스킴의 사상으로 결론을 내립니다. \end{proof}

\begin{lemma}
\label{etale-cohomology-lemma-morphism-ringed-etale-topoi-affines}
$X$, $Y$를 아핀 스킴으로 설정한다.
허락하다
$$
(g, g^\#) :
(\Sh(X_\etale), \mathcal{O}_X)
\longrightarrow
(\Sh(Y_\etale), \mathcal{O}_Y)
$$
로컬 링 토포스의 사상이어야 한다. 그런 다음 존재한다.
스킴 $f : X \to Y$의 고유한 사상
$(g, g^\#)$는 $2$와 $(f_{small}, f_{small}^\sharp)$ 동형이다.
보다
가군 사이트에,
정의 \ref{sites-modules-definition-2-morphism-ringed-topoi}.
\end{lemma}

\begin{proof}
이 증명에서는 구조층에 대해 $\mathcal{O}_X$를 씁니다.
작은 에탈 사이트 $X_\etale$의 경우와 마찬가지로
$\mathcal{O}_Y$. $Y = \Spec(B)$ 및 $X = \Spec(A)$라고 말한다. 부터
$B = \Gamma(Y_\etale, \mathcal{O}_Y)$,
$A = \Gamma(X_\etale, \mathcal{O}_X)$
$g^\sharp$가 환 맵 $\varphi : B \to A$을 유도하는 것을 볼 수 있다.
$f = \Spec(\varphi) : X \to Y$를 해당 사상으로 설정한다.
아핀 스킴. $f$이 작업을 수행하는 것을 보여드리겠습니다.

\medskip\noindent
$V \to Y$를 $Y$에 대한 아핀 스킴의 에탈로 둡니다. 따라서 우리는 다음과 같이 쓸 수 있다.
$V = \Spec(C)$ 와 $C$ 에탈 $B$-대수학. 우리는 쓸 수 있다
$$
C = B[x_1, \ldots, x_n]/(P_1, \ldots, P_n)
$$
$P_i$ 다음과 같은 다항식 사용: $\Delta = \det(\partial P_i/ \partial x_j)$
$C$에서 반전 가능한다. 예를 참조하라.
대수학, 보조정리 \ref{algebra-lemma-etale-standard-smooth}.
만약에$T$는스킴~ 위에$Y$, 그런 다음$T$-가치 있는 점$V$에 의해 주어진다
$n$ 절~의$\Gamma(T, \mathcal{O}_T)$다항식을 만족하는 것은
방정식 $P_1 = 0, \ldots, P_n = 0$. 즉, 층 $h_V$
위의 $Y_\etale$은 두 맵의 이퀄라이저이다.
$$
\xymatrix{
\prod\nolimits_{i = 1, \ldots, n} \mathcal{O}_Y
\ar@<1ex>[r]^a \ar@<-1ex>[r]_b &
\prod\nolimits_{j = 1, \ldots, n} \mathcal{O}_Y
}
$$
여기서 $b(h_1, \ldots, h_n) = 0$ 및
$a(h_1, \ldots, h_n) =
(P_1(h_1, \ldots, h_n), \ldots, P_n(h_1, \ldots, h_n))$.
$g^{-1}$는 정확하므로 우리는
다음 솔리드 가환 도식도 등화자 도식이다.
$$
\xymatrix{
g^{-1}h_V \ar[r] \ar@{..>}[d] &
\prod\nolimits_{i = 1, \ldots, n} g^{-1}\mathcal{O}_Y
\ar@<1ex>[r]^{g^{-1}a} \ar@<-1ex>[r]_{g^{-1}b} \ar[d]^{\prod g^\sharp} &
\prod\nolimits_{j = 1, \ldots, n} g^{-1}\mathcal{O}_Y \ar[d]^{\prod g^\sharp}\\
h_{X \times_Y V} \ar[r] &
\prod\nolimits_{i = 1, \ldots, n} \mathcal{O}_X
\ar@<1ex>[r]^{a'} \ar@<-1ex>[r]_{b'} &
\prod\nolimits_{j = 1, \ldots, n} \mathcal{O}_X  \\
}
$$
여기서 $b'$는 제로 맵이고 $a'$는 다음과 같이 정의된 맵이다.
동일한 이미지를 통해 $P'_i = \varphi(P_i) \in A[x_1, \ldots, x_n]$
규칙
$a'(h_1, \ldots, h_n) =
(P'_1(h_1, \ldots, h_n), \ldots, P'_n(h_1, \ldots, h_n))$.
$a$에 의해 정의되었다. 다이어그램의 교환성은 다음과 같습니다.
대역 절단에 $\varphi = g^\sharp$이 있다는 사실. 낮은
행은 등화자 도식 이기도 한다.
$X \times_Y V$는 아핀 스킴이기 때문에 이전에
$\Spec(A \otimes_B C)$ 그리고
$A \otimes_B C = A[x_1, \ldots, x_n]/(P'_1, \ldots, P'_n)$.
따라서 우리는 독특한 점선 화살표를 얻습니다.
$g^{-1}h_V \to h_{X \times_Y V}$ 다이어그램에 적합

\medskip\noindent
우리는 주장 층 $g^{-1}h_V \to h_{X \times_Y V}$의 지도를
동형사상이다. $X$의 작은 에탈 사이트에는 점수가 충분하므로
(정리 \ref{etale-cohomology-theorem-exactness-stalks})
줄기에서 이를 증명하면 충분한다. 따라서 $\overline{x}$를
$X$의 기하점, 그리고 $p$의 연관 포인트를 나타냅니다.
$X$의 작은 에탈 토포스. 집합 $q = g \circ p$. 이것이 포인트이다
$Y$의 작은 에탈 토포스. 에 의해
보조정리 \ref{etale-cohomology-lemma-points-small-etale-site}
$q$는 다음의 기하점 $\overline{y}$에 해당한다.
$Y$. 줄기의 지도를 생각해 보세요.
$$
(g^\sharp)_p :
(\mathcal{O}_Y)_{\overline{y}} =
\mathcal{O}_{Y, q} =
(g^{-1}\mathcal{O}_Y)_p
\longrightarrow
\mathcal{O}_{X, p} =
(\mathcal{O}_X)_{\overline{x}}
$$
$(g, g^\sharp)$는 {\it 국소} 환 달린 토포스의 사상이므로
$(g^\sharp)_p$는 국소환 엄밀한 헨셀의 동형이다.
국소환. 위의 큰 가환 도식에 국소화를 적용하고
대수학, 보조정리 \ref{algebra-lemma-strictly-henselian-solutions}
우리는 $(g^{-1}h_V)_p \to (h_{X \times_Y V})_p$가 동형사상라고 결론을 내립니다.
원하는대로.

\medskip\noindent 주장 동형사상 $g^{-1}h_V \to h_{X \times_Y V}$는 기능적이다. 즉, $V_1 \to V_2$는 $Y$에 대한 아핀 스킴의 에탈의 사상라고 가정하자. $V_i = \Spec(C_i)$를 $$
C_i = B[x_{i, 1}, \ldots, x_{i, n_i}]/(P_{i, 1}, \ldots, P_{i, n_i})
$$로 작성한다. 사상 $V_1 \to V_2$는 $B$ 대수 맵 $C_2 \to C_1$에 의해 제공되며 이는 차례로 일부 다항식 $Q_j \in B[x_{1, 1}, \ldots, x_{1, n_1}]$에 의해 제공된다. $j = 1, \ldots, n_2$. 그러면 층 $$
\xymatrix{
h_{V_1} \ar[d] \ar[r] & \prod_{i = 1, \ldots, n_1} \mathcal{O}_Y
\ar[d]^{Q_1, \ldots, Q_{n_2}}\\
h_{V_2} \ar[r] & \prod_{i = 1, \ldots, n_2} \mathcal{O}_Y
}
$$의 다이어그램이 교환 가능하다는 것을 보여주는 것은 쉬운 일이며, $X_\etale$로 돌아가서 솔리드 가환 도식 $$
\xymatrix{
g^{-1}h_{V_1} \ar@{..>}[dd] \ar[rrd] \ar[r] &
\prod_{i = 1, \ldots, n_1} g^{-1}\mathcal{O}_Y
\ar[dd]^{g^\sharp}
\ar[rrd]^{Q_1, \ldots, Q_{n_2}} \\
& & g^{-1}h_{V_2} \ar@{..>}[dd] \ar[r] &
\prod_{i = 1, \ldots, n_2} g^{-1}\mathcal{O}_Y
\ar[dd]^{g^\sharp} \\
h_{X \times_Y V_1} \ar[r] \ar[rrd] &
\prod\nolimits_{i = 1, \ldots, n_1} \mathcal{O}_X
\ar[rrd]^{Q'_1, \ldots, Q'_{n_2}} \\
& & h_{X \times_Y V_2} \ar[r] &
\prod\nolimits_{i = 1, \ldots, n_2} \mathcal{O}_X
}
$$를 얻습니다. 여기서 $Q'_j \in A[x_{1, 1}, \ldots, x_{1, n_1}]$는 다음을 통해 $Q_j$의 이미지이다. $\varphi$. 점선 화살표가 존재하므로 두 개의 사각형을 통근하게 하고 수평 화살표는 단사이므로 전체 다이어그램이 통근하는 것을 볼 수 있다. 이는 기능성을 증명합니다(또한 $g^{-1}h_V \to h_{X \times_Y V}$의 구성은 표시 선택과 무관하지만 엄밀히 말하면 이를 보여줄 필요는 없습니다).

\medskip\noindent
이 시점에서 우리는 $f_{small, *} \cong g_*$을 보여줄 수 있다.
즉, $\mathcal{F}$를 $X_\etale$의 층으로 설정한다. 모든
$V \in \Ob(X_\etale)$ 우리는 유사하다
\begin{align*}
(g_*\mathcal{F})(V)
& =
\Mor_{\Sh(Y_\etale)}(h_V, g_*\mathcal{F}) \\
& =
\Mor_{\Sh(X_\etale)}(g^{-1}h_V, \mathcal{F}) \\
& =
\Mor_{\Sh(X_\etale)}(h_{X \times_Y V}, \mathcal{F}) \\
& =
\mathcal{F}(X \times_Y V) \\
& =
f_{small, *}\mathcal{F}(V)
\end{align*}
세 번째 평등에서는 동형사상을 사용한다.
$g^{-1}h_V \cong h_{X \times_Y V}$ 위에 구성되었다. 이 동형사상
$\mathcal{F}$에서는 분명히 기능적이며 $V$에서는 기능적이다.
동형사상 $g^{-1}h_V \cong h_{X \times_Y V}$는 기능적이다.
이제 $Y_\etale$의 모든 층은 제한 사항에 따라 결정된다.
아핀 스킴 하위 카테고리로
(위상, 보조정리 \ref{topologies-lemma-alternative}),
따라서 우리는 함자 $f_{small, *} \cong g_*$의 동형사상을 얻습니다.
원하는대로.

\medskip\noindent
마지막으로 동형사상을 통해 이를 확인해야 한다.
위의 $f_{small, *} \cong g_*$, 지도 $f_{small}^\sharp$ 및
$g^\sharp$ 동의한다. 구조적으로 이것은 이미 다음의 경우에 해당된다.
$\mathcal{O}_Y$의 대역 절단, 즉 $B$의 요소에 대한 것이다.
우리는 결과를 확인하기 만하면된다.
절 위의 아핀 $V$ 에탈 위의 $Y$ (에 의해
위상, 보조정리 \ref{topologies-lemma-alternative}
다시). 글쓰기
$V = \Spec(C)$, $C = B[x_i]/(P_j)$ 이전과 마찬가지로 충분한다.
좌표 함수 $x_i$가 매핑되어 있는지 확인하려면
같은절~의$\mathcal{O}_X$~ 위에$X \times_Y V$.
그리고 이것이 바로 다이어그램이 의미하는 바이다.
$$
\xymatrix{
g^{-1}h_V \ar[r] \ar@{..>}[d] &
\prod\nolimits_{i = 1, \ldots, n} g^{-1}\mathcal{O}_Y
\ar[d]^{\prod g^\sharp} \\
h_{X \times_Y V} \ar[r] &
\prod\nolimits_{i = 1, \ldots, n} \mathcal{O}_X
}
$$
통근. 따라서 보조정리가 증명되었다.
\end{proof}

\noindent 일반 스킴 버전이다.

\begin{theorem} \label{etale-cohomology-theorem-fully-faithful} $X$, $Y$를 스킴으로 둡니다. $$
(g, g^\#) :
(\Sh(X_\etale), \mathcal{O}_X)
\longrightarrow
(\Sh(Y_\etale), \mathcal{O}_Y)
$$를 로컬 링 토포스의 사상으로 설정한다. 그러면 스킴 $f : X \to Y$의 고유한 사상이 존재하므로 $(g, g^\#)$는 $(f_{small}, f_{small}^\sharp)$와 동형이다. 즉, $$
\Sch \longrightarrow \textit{국소환 달린 토포스},
\quad
X \longrightarrow (X_\etale, \mathcal{O}_X)
$$ 구성은 완전히 충실합니다(오른쪽의 사상부터 $2$-동형사상까지). \end{theorem}

\begin{proof} 보조정리 \ref{etale-cohomology-lemma-morphism-ringed-etale-topoi-affines}의 증명 인수를 전역 설정으로 신중하게 조정하여 정리를 증명할 수 있다. 그러나 보조정리 \ref{etale-cohomology-lemma-2-morphism}의 결과가 제공하는 견고성으로 인해 보조정리의 결과를 전역 사상으로 묶을 수 있는 방법을 표시하고 싶다. 불행히도 이것은 약간 지저분한다.

\medskip\noindent
$Y$이 유사할 때 존재를 증명해 보자. 이 경우 다음 중 하나를 선택하세요.
개방형 피복 $X = \bigcup U_i$. 각 $i$에 대해 포함
사상 $j_i : U_i \to X$는 로컬 링 토포스의 사상을 유도한다.
$(j_{i, small}, j_{i, small}^\sharp) :
(\Sh(U_{i, \etale}), \mathcal{O}_{U_i})
\to
(\Sh(X_\etale), \mathcal{O}_X)$
~에 의해
보조정리 \ref{etale-cohomology-lemma-morphism-locally-ringed}.
이것을 $(g, g^\sharp)$로 구성하여 사상을 얻을 수 있다.
로컬로 울린 토포스
$$
(g, g^\sharp) \circ (j_{i, small}, j_{i, small}^\sharp) :
(\Sh(U_{i, \etale}), \mathcal{O}_{U_i})
\to
(\Sh(Y_\etale), \mathcal{O}_Y)
$$
보다
가군 사이트에,
보조정리 \ref{sites-modules-lemma-composition-morphisms-locally-ringed-topoi}.
에 의해
보조정리 \ref{etale-cohomology-lemma-morphism-ringed-etale-topoi-affines}
스킴 $f_i : U_i \to Y$의 고유한 사상이 존재한다.
그리고 $2$-동형사상
$$
t_i :
(f_{i, small}, f_{i, small}^\sharp)
\longrightarrow
(g, g^\sharp) \circ (j_{i, small}, j_{i, small}^\sharp).
$$
집합 $U_{i, i'} = U_i \cap U_{i'}$, 그리고 $j_{i, i'} : U_{i, i'} \to U_i$를 나타냅니다.
포함 사상. 우리가 가지고 있기 때문에
$j_i \circ j_{i, i'} = j_{i'} \circ j_{i', i}$
우리는 그것을 본다
\begin{align*}
(g, g^\sharp) \circ
(j_{i, small}, j_{i, small}^\sharp) \circ
(j_{i, i', small}, j_{i, i', small}^\sharp)
= \\
(g, g^\sharp) \circ
(j_{i', small}, j_{i', small}^\sharp) \circ
(j_{i', i, small}, j_{i', i, small}^\sharp)
\end{align*}
따라서 고유성에 의해 (참조
보조정리 \ref{etale-cohomology-lemma-faithful})
우리는 다음과 같이 결론을 내립니다.
$f_i \circ j_{i, i'} = f_{i'} \circ j_{i', i}$ 즉,
사상 의 스킴 $f_i = f \circ j_i$ 은
스킴 $f : X \to Y$의 전역 사상. 다이어그램을 고려하십시오
$2$-동형사상 (여기서 ${}^\sharp$ 구성요소를 삭제하여
표기법)
$$
\xymatrix{
g \circ j_{i, small} \circ j_{i, i', small}
\ar[rr]^{t_i \star \text{id}_{j_{i, i', small}}}
\ar@{=}[d] & &
f_{small} \circ j_{i, small} \circ j_{i, i', small} \ar@{=}[d] \\
g \circ j_{i', small} \circ j_{i', i, small}
\ar[rr]^{t_{i'} \star \text{id}_{j_{i', i, small}}} & &
f_{small} \circ j_{i', small} \circ j_{i', i, small}
}
$$
$\star$ 표기는 가로 구성을 나타냅니다.
범주론, 정의 \ref{categories-definition-2-category}
일반적으로 그리고
사이트, 절 \ref{sites-section-2-category}
우리의 특별한 경우에. 의 결과로
보조정리 \ref{etale-cohomology-lemma-2-morphism}
이 다이어그램은 통근한다. 따라서 층 $\mathcal{G}$
$Y_\etale$에서 동형사상
$t_i : f_{small}^{-1}\mathcal{G}|_{U_i} \to g^{-1}\mathcal{G}|_{U_i}$
$U_{i, i'}$에 동의하고 전역 동형사상을 얻습니다.
$t : f_{small}^{-1}\mathcal{G} \to g^{-1}\mathcal{G}$.
이 동형사상은 $\mathcal{G}$에서 기능적임이 분명한다.
$f_{small}^\sharp$ 및 $g^\sharp$ 지도와 호환된다.
(이 지도와 로컬로 호환되기 때문입니다)
이는 $Y$가 유사할 경우 정리를 증명한다.

\medskip\noindent
일반적인 경우에는 $V \subset Y$를 아핀 개방형으로 둡니다.
그 다음에$h_V$결승전의 하위단이다층 $*$~에$Y_\etale$.
$g$가 정확하므로 $g^{-1}h_V$가 최종 항목의 하위 뭉치라는 것을 알 수 있다.
층 - $X_\etale$. 따라서
보조정리 \ref{etale-cohomology-lemma-support-subsheaf-final}
공개 하위 계획이 존재합니다$W \subset X$그렇게$g^{-1}h_V = h_W$. 에 의해
가군 사이트에,
보조정리 \ref{sites-modules-lemma-localize-morphism-locally-ringed-topoi}
로컬로 울린 사상 중 가환 도식이 존재한다.
토포스
$$
\xymatrix{
(\Sh(W_\etale), \mathcal{O}_W) \ar[r] \ar[d]_{g'} &
(\Sh(X_\etale), \mathcal{O}_X) \ar[d]^g \\
(\Sh(V_\etale), \mathcal{O}_V) \ar[r] &
(\Sh(Y_\etale), \mathcal{O}_Y)
}
$$
여기서 가로 화살표는 국소화 사상이다.
(사상 $V \to Y$ 및 $W \to X$ 포함으로 인해 발생함)
여기서 $g'$는 $g$에서 유도된다. 앞선 결과로
단락에서 스킴 $f' : W \to V$의 사상을 얻습니다.
$2$-동형사상
$t : (f'_{small}, (f'_{small})^\sharp) \to (g', (g')^\sharp)$.
정확히 이전과 마찬가지로 사상 $f'$ (다양한 아핀 개방의 경우 $V \subset Y$)
고유성에 의한 중복에 동의하므로 사상 $f : X \to Y$을 얻습니다.
또한, $2$-동형사상 $t$는
보조정리 \ref{etale-cohomology-lemma-2-morphism}
다시 한 번 전역 $2$-동형사상을 얻습니다.
$(f_{small}, (f_{small})^\sharp) \to (g, (g)^\sharp)$.
원하는대로. 일부 세부 사항은 생략되었다.
\end{proof}










\section{밀고 당기기} \label{etale-cohomology-section-monomorphisms}

\noindent
$f : X \to Y$를 스킴의 사상으로 설정한다.
다음에서 고려할 조건 목록은 다음과 같습니다.
\begin{enumerate}
\item[()] 모든 에탈 사상 $U \to X$ 및 $u \in U$가 존재한다.
에탈 사상 $V \to Y$ 및 분리된 결합 분해
$X \times_Y V = W \amalg W'$ 및 사상 $h : W \to U$ 위의 $X$
~와 함께$u$의 이미지에서$h$.
\item[(B)] 매 $V \to Y$ 에탈, 그리고 매 에탈 피복
$\{U_i \to X \times_Y V\}$ 에탈이 존재합니다 피복
$\{V_j \to V\}$ 각 $j$에 대해
$X \times_Y V_j = \coprod W_{ij}$ 여기서 $W_{ij} \to X \times_Y V$
요인을 통해$U_i \to X \times_Y V$어떤 사람들에게는$i$.
\item[(C)] 모든 $U \to X$ 에탈에는 $V \to Y$ 에탈이 존재한다.
그리고 전사사상 $X \times_Y V \to U$~ 위에$X$.
\end{enumerate}
이들 성질 각각은 다음과 같은 의미를 갖고 있는 것으로 밝혀졌습니다.
함자 $f_{small, *}$의 동작. 우리는 이것을 할 것이다
다음 몇 절에 나올 것이다.



\section{성질 ()} \label{etale-cohomology-section-A}

\noindent 성질 ()의 정의에 대해서는 절 \ref{etale-cohomology-section-monomorphisms}를 참조하라.

\begin{lemma}
\label{etale-cohomology-lemma-property-A-implies}
$f : X \to Y$를 스킴의 사상으로 설정한다.
()를 가정해보자.
\begin{enumerate}
\item
$f_{small, *} :
\textit{Ab}(X_\etale)
\to
\textit{Ab}(Y_\etale)$
주사와 추측을 반영한다.
\item $f_{small}^{-1}f_{small, *}\mathcal{F} \to \mathcal{F}$
는 어떤 것에 대해서도 전사적이다아벨 층 $\mathcal{F}$~에$X_\etale$,
\item
$f_{small, *} :
\textit{Ab}(X_\etale)
\to
\textit{Ab}(Y_\etale)$
충실하다.
\end{enumerate}
\end{lemma}

\begin{proof}
$\mathcal{F}$를 $X_\etale$에서 아벨 층으로 설정한다.
$U$를 $X_\etale$의 객체로 둡니다. 가정으로 우리는
피복 $\{W_i \to U\}$ 에서 $X_\etale$ 각 $W_i$는 다음과 같습니다.
일부 개체에 대한 $X \times_Y V_i$의 개방형 및 폐쇄형 하위 구성표
$V_i$ 중 $Y_\etale$이다. 층 조건은 다음을 보여줍니다.
$$
\mathcal{F}(U) \subset \prod \mathcal{F}(W_i)
$$
그리고 $\mathcal{F}(W_i)$는 다음의 직접적인 합이다.
$\mathcal{F}(X \times_Y V_i) = f_{small, *}\mathcal{F}(V_i)$.
따라서 $f_{small, *}$이 주사를 반영하는 것이 분명한다.

\medskip\noindent
다음으로, $a : \mathcal{G} \to \mathcal{F}$가 다음의 맵이라고 가정한다.
아벨 층 $f_{small, *}a$는 전사이다. 허락하다
$s \in \mathcal{F}(U)$ 위와 같이 $U$를 사용한다. $W_i$, $V_i$를 다음과 같이 사용한다.
위에서 우리는 $s|_{W_i}$가 지역적으로 에탈임을 보여주는 것으로 충분하다는 것을 알 수 있다.
의 이미지절~의$\mathcal{G}$아래에$a$. 부터$\mathcal{F}(W_i)$
는 $\mathcal{F}(X \times_Y V_i)$의 직접적인 합이다.
$V \in \Ob(Y_\etale)$
모든 요소 $s \in \mathcal{F}(X \times_Y V)$
의 절 이미지는 $X \times_Y V$에 로컬로 에탈되어 있다.
$\mathcal{G}$ 아래 $a$. 부터
$\mathcal{F}(X \times_Y V) = f_{small, *}\mathcal{F}(V)$
가정을 통해 피복 $\{V_j \to V\}$가 존재한다는 것을 알 수 있다.
$s$는 다음의 이미지이다.
$s_j \in f_{small, *}\mathcal{G}(V_j) = \mathcal{G}(X \times_Y V_j)$.
이는 $f_{small, *}$가 추측을 반영한다는 것을 증명한다.

\medskip\noindent 부품(2), (3)는 부품(1)에서 형식적으로 따릅니다. 사이트, 보조정리의 가군을 참조하라. \ref{sites-modules-lemma-reflect-surjections}. \end{proof}

\begin{lemma} \label{etale-cohomology-lemma-locally-quasi-finite-A} $f : X \to Y$를 스킴의 분리된 국부적 준유한 사상으로 둡니다. 그러면 위의 성질 ()가 유지된다. \end{lemma}

\begin{proof}
$U \to X$를 에탈 사상 및 $u \in U$로 가정한다.
기하학적 설명()은 $U$ 및 $Y$인 경우로 직접적으로 축소된다.
아핀 스킴이다. $x \in X$ 및 $y \in Y$를 나타냅니다.
$u$의 이미지. $X \to Y$는 지역적으로 준유한이고 $U \to X$는
국소적으로 준유한(참조
사상, 보조정리 \ref{morphisms-lemma-etale-locally-quasi-finite})
우리는 $U \to Y$이 지역적으로 준 유한하다는 것을 알 수 있습니다 (참조
사상, 보조정리 \ref{morphisms-lemma-composition-quasi-finite}).
또한 $X \to Y$와 $U \to Y$는 모두 분리되어 있다. 따라서
사상, 보조정리에 대한 추가 정보
\ref{more-morphisms-lemma-etale-splits-off-quasi-finite-part-technical-variant}
사상 둘 다에 적용된다. 이는 우리가 에탈 지역을 선택할 수 있음을 의미한다.
$(V, v) \to (Y, y)$ 이렇게
$$
X \times_Y V = W \amalg R, \quad
U \times_Y V = W' \amalg R'
$$
그리고 포인트 $w \in W$, $w' \in W'$는 다음과 같습니다.
\begin{enumerate}
\item $W$, $R$는 $X \times_Y V$에서 열려 있고 닫혀 있다.
\item $W'$, $R'$는 $U \times_Y V$에서 열려 있고 닫혀 있다.
\item $W \to V$ 및 $W' \to V$는 유한한다.
\item $w$, $w'$는 $v$에 매핑된다.
\item $\kappa(v) \subset \kappa(w)$ 및 $\kappa(v) \subset \kappa(w')$
완전히 분리될 수 없으며,
\item $W$ 또는 $W'$의 다른 지점은 $v$에 매핑되지 않습니다.
\end{enumerate}
여기 가환 도식이 있다.
$$
\xymatrix{
U \ar[d] & U \times_Y V \ar[l] \ar[d] & W' \amalg R' \ar[d] \ar[l] \\
X \ar[d] & X \times_Y V \ar[l] \ar[d] & W \amalg R \ar[l] \\
Y & V \ar[l]
}
$$
$V$를 축소한 후 $W'$가 $W$에 매핑된다고 가정할 수 있다.
그냥 이미지를 삭제하세요역상~의$R$~에$W'$; 이것은
닫힌 집합 ($W' \to V$은 유한하므로) $v$을 포함하지 않습니다.
그 다음에$W' \to W$둘 다 유한하기 때문에$W \to V$그리고$W' \to V$유한한다.
따라서 $W' \to W$는 유한한 에탈이고,
올~ 위에$w$~와 함께$\kappa(w) = \kappa(w')$. 따라서$W' \to W$은
동형사상열린 동네에서$W^\circ$~의$w$, 보다
에탈 사상, 보조정리 \ref{etale-lemma-finite-etale-one-point}.
$W \to V$는 유한하므로 $W \setminus W^\circ$의 이미지는 닫힌 이미지이다.
하위 집합$T$~의$V$포함하지 않음$v$. 그래서 교체 후$V$~에 의해
$V \setminus T$ $W' \to W$는 동형사상라고 가정할 수 있다.
이제 분해 $X \times_Y V = W \amalg R$ 및 사상
$W \to U$ 원하는 대로 되었고 우리가 승리했다.
\end{proof}

\begin{lemma} \label{etale-cohomology-lemma-integral-A} $f : X \to Y$를 정역 사상 스킴으로 설정한다. 그러면 성질 ()가 유지된다. \end{lemma}

\begin{proof}
$U \to X$를 에탈, $u \in U$를 점이라고 하자.
우리는 분리된 조합 분해인 $V \to Y$ 에탈을 찾아야 한다.
$X \times_Y V = W \amalg W'$ 및 $X$-사상 $W \to U$
이미지에 $u$이 있다. $U$ 및 $Y$를 축소하고 가정할 수 있다.
$U$와 $Y$는 유사한다. 이 경우에는 $X$도 유사한다.
정역 사상은 정의와 유사한다. $Y = \Spec(A)$를 쓰고,
$X = \Spec(B)$ 및 $U = \Spec(C)$. 그러면 $A \to B$는
정역 환 지도이고, $B \to C$는 에탈 환 지도이다. 에 의해
대수학, 보조정리 \ref{algebra-lemma-etale}
유한 $A$-하대수학 $B' \subset B$과 에탈 환을 찾을 수 있다.
$B' \to C'$을 $C = B \otimes_{B'} C'$로 매핑한다. 따라서 질문
에탈 사상으로 감소한다.
$U' = \Spec(C') \to X' = \Spec(B')$
유한 사상 $X' \to Y$에 대해. 이 경우 결과는 다음과 같습니다.
보조정리 \ref{etale-cohomology-lemma-locally-quasi-finite-A}.
\end{proof}

\begin{lemma} \label{etale-cohomology-lemma-when-push-pull-surjective} $f : X \to Y$를 스킴의 사상으로 둡니다. 작은 에탈 토포스의 연관된 사상을 $f_{small} :
\Sh(X_\etale)
\to
\Sh(Y_\etale)$로 표시한다. 다음 \begin{enumerate} \item $f$ 중 적어도 하나가 정역이거나 \item $f$가 분리되고 국부적으로 준유한이라고 가정한다. \end{enumerate} 그러면 함자 $f_{small, *} :
\textit{Ab}(X_\etale)
\to
\textit{Ab}(Y_\etale)$는 다음과 같습니다. 성질 \begin{enumerate} \item 지도는 $f_{small}^{-1}f_{small, *}\mathcal{F} \to \mathcal{F}$ 항상 전사적이다. \item $f_{small, *}$는 충실하며, \item $f_{small, *}$는 주사와 추측을 반영한다. \end{enumerate} \end{lemma}

\begin{proof} 기본정리 \ref{etale-cohomology-lemma-locally-quasi-finite-A}, \ref{etale-cohomology-lemma-integral-A} 및 \ref{etale-cohomology-lemma-property-A-implies}를 결합한다. \end{proof}



\section{성질 (B)} \label{etale-cohomology-section-B}

\noindent 성질 (B)의 정의에 대해서는 절 \ref{etale-cohomology-section-monomorphisms}를 참조하라.

\begin{lemma} \label{etale-cohomology-lemma-property-B-implies} $f : X \to Y$를 스킴의 사상으로 둡니다. (B)가 성립한다고 가정한다. 그런 다음 함자 $f_{small, *} :
\Sh(X_\etale)
\to
\Sh(Y_\etale)$는 전사를 전사로 변환한다. \end{lemma}

\begin{proof} 이는 사이트, 보조정리 \ref{sites-lemma-weaker}에 이어집니다. \end{proof}

\begin{lemma} \label{etale-cohomology-lemma-simplify-B} $f : X \to Y$를 스킴의 사상으로 둡니다. \begin{enumerate} \item $V \to Y$는 스킴의 에탈 사상이고, \item $\{U_i \to X \times_Y V\}$는 에탈 피복이고, \item $v \in V$는 점이다. \end{enumerate} 이러한 데이터에 대해 에탈 이웃 $(V', v') \to (V, v)$, 분리된 결합 분해 $X \times_Y V' = \coprod W'_i$, 사상 $W'_i \to U_i$가 $X \times_Y V$ 위에 존재한다고 가정한다. 그러면 성질 (B)가 유지된다. \end{lemma}

\begin{proof} 생략한다. \end{proof}

\begin{lemma} \label{etale-cohomology-lemma-finite-B} $f : X \to Y$를 스킴의 유한 사상으로 설정한다. 그러면 성질 (B)가 유지된다. \end{lemma}

\begin{proof} $V \to Y$ 에탈, $\{U_i \to X \times_Y V\}$ 에탈 피복 및 $v \in V$를 고려하세요. $V' \to V$를 찾아서 보조정리 \ref{etale-cohomology-lemma-simplify-B}처럼 분해하고 매핑해야 한다. $V$와 $Y$를 축소할 수 있으므로 $V$와 $Y$는 유사하다고 가정할 수 있다. $X$는 $Y$에 대해 유한하므로 이는 또한 $X$가 유사하다는 것을 의미한다. 증명 동안 우리는 $(V, v)$를 에탈 이웃 $(V', v')$로 대체하고 이에 따라 피복 $\{U_i \to X \times_Y V\}$를 $\{V' \times_V U_i \to X \times_Y V'\}$로 대체할 수 있다.

\medskip\noindent
$X \times_Y V \to V$는 유한하므로 유한하게 존재한다.
많은(쌍별로 구별되는) 포인트 $x_1, \ldots, x_n \in X \times_Y V$
$v$에 매핑된다. 우리는 신청할 수 있습니다
사상, 보조정리에 대한 추가 정보
\ref{more-morphisms-lemma-etale-splits-off-quasi-finite-part-technical-variant}
$X \times_Y V \to V$에 포인트 $x_1, \ldots, x_n$가 놓여 있다.
$v$ 그리고 에탈 동네 $(V', v') \to (V, v)$를 찾습니다.
그렇게
$$
X \times_Y V' = R \amalg \coprod T_a
$$
$T_a \to V'$ 유한하고 정확히 하나의 점 $p_a$이 $v'$ 위에 놓여 있음
게다가 $\kappa(v') \subset \kappa(p_a)$는 완전히 분리될 수 없으며,
$R \to V'$에는 $v'$ 위에 빈 올이 있다.
$X \to Y$은 유한하므로 $R \to V'$도 유한한다. 따라서 이후
축소 $V'$ $R = \emptyset$라고 가정할 수 있다. 따라서 우리는
$X \times_Y V = X_1 \amalg \ldots \amalg X_n$를 다음과 같이 가정한다.
정확히 한 점 $x_l \in X_l$이 $v$ 위에 놓여 있고 게다가
$\kappa(v) \subset \kappa(x_l)$ 완전히 분리될 수 없다. 참고하세요
성질은 에탈 동네의 정제된 상태로 보존된다.
$(V, v)$.

\medskip\noindent
각각에 대해$l$선택하다$i_l$그리고 포인트$u_l \in U_{i_l}$매핑$x_l$.
이제 유한 사상에 대해 성질 ()를 적용한다.
$X \times_Y V \to V$ 그리고 에탈
사상 $U_{i_l} \to X \times_Y V$ 및 점 $u_l$.
이는 다음에 의해 허용된다.
보조정리 \ref{etale-cohomology-lemma-integral-A}
이는 에탈 이웃 $(V', v') \to (V, v)$을 생성한다.
그리고 분해
$$
X \times_Y V' = W_l \amalg R_l
$$
및 $X$-사상 $a_l : W_l \to U_{i_l}$ 이미지에 $u_{i_l}$가 포함되어 있다.
여기 사진이 있습니다:
$$
\xymatrix{
& & & U_{i_l} \ar[d] & \\
W_l \ar[rrru] \ar[r] & W_l \amalg R_l \ar@{=}[r] &
X \times_Y V' \ar[r] \ar[d] &
X \times_Y V \ar[r] \ar[d] & X \ar[d] \\
& & V' \ar[r] & V \ar[r] & Y
}
$$
$(V, v)$을 $(V', v')$로 바꾼 후 우리는 각각의 결론을 내립니다.
$x_l$는 열린 이웃과 닫힌 이웃 $W_l$에 포함되어 있다.
포함 사상 $W_l \to X \times_Y V$ 요소를 통해
$U_i \to X \times_Y V$ 일부 $i$에 대해. $W_l$를 $W_l \cap X_l$로 교체
우리는 열린 것과 닫힌 것 집합이 분리되어 있다는 것을 알 수 있다.
그 $\{x_1, \ldots, x_n\} \subset W_1 \cup \ldots \cup W_n$.
$X \times_Y V \to V$는 유한하므로 $V$를 축소하고 다음과 같이 가정할 수 있다.
$X \times_Y V = W_1 \amalg \ldots \amalg W_n$ 원하는 대로.
\end{proof}

\begin{lemma} \label{etale-cohomology-lemma-integral-B} $f : X \to Y$를 정역 사상 스킴으로 설정한다. 그러면 성질 (B)가 유지된다. \end{lemma}

\begin{proof}
$V \to Y$ 에탈, $\{U_i \to X \times_Y V\}$ 에탈 피복을 고려하고,
$v \in V$. 우리는 다음과 같이 $V' \to V$와 분해 및 맵을 찾아야 한다.
보조정리 \ref{etale-cohomology-lemma-simplify-B}.
$V$와 $Y$를 축소할 수 있으므로 $V$와 $Y$는 유사하다고 가정할 수 있다.
$X$는 $Y$에 대해 정역이므로 이는 $X$ 및
$X \times_Y V$는 유사한다. 피복을 개선할 수 있다.
$\{U_i \to X \times_Y V\}$, 따라서 우리는 다음과 같이 가정할 수 있다.
$\{U_i \to X \times_Y V\}_{i = 1, \ldots, n}$은 표준 에탈 피복이다.
$Y = \Spec(A)$, $X = \Spec(B)$,
$V = \Spec(C)$ 및 $U_i = \Spec(B_i)$.
그러면 $A \to B$는 정역 환 지도이고 $B \otimes_A C \to B_i$는
에탈 환 지도. 에 의해
대수학, 보조정리 \ref{algebra-lemma-etale}
유한 $A$-하대수학 $B' \subset B$과 에탈 환을 찾을 수 있다.
지도$B' \otimes_A C \to B'_i$~을 위한$i = 1, \ldots, n$
$B_i = B \otimes_{B'} B'_i$와 같습니다. 따라서 질문
에탈 피복으로 감소한다.
$\{\Spec(B'_i) \to X' \times_Y V\}_{i = 1, \ldots, n}$
~와 함께$X' = \Spec(B')$유한한 이상$Y$.
이 경우 결과는 다음과 같습니다.
보조정리 \ref{etale-cohomology-lemma-finite-B}.
\end{proof}

\begin{lemma}
\label{etale-cohomology-lemma-what-integral}
$f : X \to Y$를 스킴의 사상으로 설정한다.
$f$가 정역라고 가정합니다(예 유한의 경우).
그 다음에
\begin{enumerate}
\item $f_{small, *}$ 전사를 전사로 변환합니다(층
집합 및 아벨 층),
\item $f_{small}^{-1}f_{small, *}\mathcal{F} \to \mathcal{F}$
는 어떤 것에 대해서도 전사적이다아벨 층 $\mathcal{F}$~에$X_\etale$,
\item
$f_{small, *} :
\textit{Ab}(X_\etale)
\to
\textit{Ab}(Y_\etale)$
충실하고 주사와 추측을 반영하며,
\item
$f_{small, *} :
\textit{Ab}(X_\etale)
\to
\textit{Ab}(Y_\etale)$
정확한다.
\end{enumerate}
\end{lemma}

\begin{proof} 부품 (2), (3)는 보조정리 \ref{etale-cohomology-lemma-when-push-pull-surjective}에서 확인했다. 부분(1)은 기본정리 \ref{etale-cohomology-lemma-integral-B} 및 \ref{etale-cohomology-lemma-property-B-implies}에서 이어집니다. 부분(4)은 부분(1)의 결과이다. 사이트, 보조정리 \ref{sites-modules-lemma-exactness}의 가군을 참조하라. \end{proof}








\section{성질 (C)} \label{etale-cohomology-section-C}

\noindent 성질 (C)의 정의에 대해서는 절 \ref{etale-cohomology-section-monomorphisms}를 참조하라.

\begin{lemma} \label{etale-cohomology-lemma-property-C-implies} $f : X \to Y$를 스킴의 사상으로 둡니다. (C)가 성립한다고 가정한다. 그러면 함자 $f_{small, *} :
\Sh(X_\etale)
\to
\Sh(Y_\etale)$는 주입과 전사를 반영한다. \end{lemma}

\begin{proof} 사이트, 보조정리 \ref{sites-lemma-cover-from-below}에 따릅니다. 성질 (C)는 함자 $Y_\etale \to X_\etale$, $V \mapsto X \times_Y V$가 사이트, 보조정리의 가정을 만족한다는 것을 의미한다는 검증을 생략한다. \ref{sites-lemma-cover-from-below}. \end{proof}

\begin{remark} \label{etale-cohomology-remark-property-C-strong} 성질 (C)는 $f : X \to Y$가 열린 몰입인 경우 유지된다. 즉, $U \in \Ob(X_\etale)$이면 $U$를 $Y_\etale$와 $U \times_Y X = U$의 객체로도 볼 수 있다. 따라서 성질 (C)는 $f_{small, *}$가 정확하다는 것을 의미하지 않습니다. 이는 열린 몰입(일반적으로)의 경우가 아니기 때문이다. \end{remark}

\begin{lemma} \label{etale-cohomology-lemma-property-C-closed-implies} $f : X \to Y$를 스킴의 사상으로 둡니다. 모든 $V \to Y$ 에탈에 대해 \begin{enumerate} \item $X \times_Y V \to V$에는 성질 (C)가 있고 \item $X \times_Y V \to V$는 닫혀 있다고 가정한다. \end{enumerate} 그러면 함자 $Y_\etale \to X_\etale$, $V \mapsto X \times_Y V$는 거의 공연속이다. 사이트, 정의 \ref{sites-definition-almost-cocontinuous}를 참조하라. \end{lemma}

\begin{proof}
$V \to Y$를 $Y_\etale$의 객체로 두고
$\{U_i \to X \times_Y V\}_{i \in I}$는 $X_\etale$의 피복이다.
가정 (1) 각 $i$에 대해 에탈 사상을 찾을 수 있다.
$h_i : V_i \to V$ 및 전사 사상 $X \times_Y V_i \to U_i$
$X \times_Y V$ 이상. $\bigcup h_i(V_i) \subset V$는
열린집합 닫힌 집합 $Z = \Im(X \times_Y V \to V)$을 포함한다.
$h_0 : V_0 = V \setminus Z \to V$를 열린 몰입으로 설정한다.
$\{V_i \to V\}_{i \in I \cup \{0\}}$은(는)
에탈 피복 각 $i \in I \cup \{0\}$에 대해 우리는
$V_i \times_Y X = \emptyset$ (즉, $i = 0$인 경우) 또는
$V_i \times_Y X \to V \times_Y X$는 $U_i \to X \times_Y V$를 통해 요인이 된다.
($i \not = 0$인 경우). 따라서 함자 $Y_\etale \to X_\etale$
거의 연속적이다.
\end{proof}

\begin{lemma} \label{etale-cohomology-lemma-integral-homeo-onto-image-C} $f : X \to Y$를 정역 사상의 스킴으로 정의한다. 이는 $X$의 동형을 닫힌 부분집합으로 정의한다. $Y$. 그러면 성질 (C)가 유지된다. \end{lemma}

\begin{proof}
$g : U \to X$를 에탈 사상라고 하자. 우리는 물건을 찾아야 해요
$V \to Y$~의$Y_\etale$그리고 전사사상 $X \times_Y V \to U$
$X$ 이상. 모든 $u \in U$에 대해 객체를 찾을 수 있다고 가정한다.
$V_u \to Y$ 중 $Y_\etale$ 및 사상 $h_u : X \times_Y V_u \to U$
$X$를 $u \in \Im(h_u)$로 초과한다. 그러면 $V = \coprod V_u$를 취하면 된다.
그리고 $h = \coprod h_u$ 그러면 우리가 승리한다. 따라서 포인트가 주어졌습니다.
$u \in U$ 위와 같이 $(V_u, h_u)$ 쌍을 찾습니다. 이를 위해 우리는
$U$을 축소하고 $U$이 유사하다고 가정한다. 이 경우
$g : U \to X$는 지역적으로 준유한적이다. 허락하다
$g^{-1}(g(\{u\})) = \{u, u_2, \ldots, u_n\}$. 없기 때문에
특수화 $u_i \leadsto u$ $U$을 유사 이웃으로 대체할 수 있다.
그래서 $g^{-1}(g(\{u\})) = \{u\}$.

\medskip\noindent
이미지 $g(U) \subset X$가 열려 있다.
따라서 $f(g(U))$는 $Y$에서 로컬로 닫혀 있다. 다음과 같은 열린 $V \subset Y$를 선택하세요.
그 $f(g(U)) = f(X) \cap V$. $g$는 다음을 통해 고려된다.
$X \times_Y V$ 결과 $\{U \to X \times_Y V\}$는 에탈이다.
피복. $f$에는 성질 (B)가 있으므로 다음을 참조하라.
보조정리 \ref{etale-cohomology-lemma-integral-B},
우리는 에탈 피복 $\{V_j \to V\}$가 존재한다는 것을 알 수 있다.
$X \times_Y V_j \to X \times_Y V$는 $U$을 통해 인수분해된다.
이는 $V' = \coprod V_j$가 $Y$에 대해 에탈이고
사상 $h : X \times_Y V' \to U$ 그 이미지
$g(U)$에 대한 추측이다. $u$는 올의 유일한 지점이므로
$h$의 이미지를 갖고 있으면 우리가 승리한다.
\end{proof}

\noindent 다음의 보조정리를 중간정착소로 생각하시기 바란다.\footnote{중간역은 긴 여행 중에 잠시 멈춰 식사를 하고 쉬어가는 곳이다.} $f_{small, *}$에 대한 궁극적인 진실을 향한 여정에 정역 보편적으로 분사되는 사상.

\begin{lemma} \label{etale-cohomology-lemma-integral-universally-injective} $f : X \to Y$를 스킴의 사상으로 둡니다. $f$는 보편적으로 단사적이며 정역(예 및 닫힌 몰입의 경우)라고 가정한다. 그런 다음 \begin{enumerate} \item $f_{small, *} :
\Sh(X_\etale)
\to
\Sh(Y_\etale)$는 주입 및 주입을 반영하고, \item $f_{small, *} :
\Sh(X_\etale)
\to
\Sh(Y_\etale)$는 푸시아웃 및 동일화 장치를 사용하여 통근합니다(더 일반적으로는 유한 연결 여극한). \item $f_{small, *}$는 전사를 전사로 변환합니다(집합의 층 및 아벨 층), \item 지도 $f_{small}^{-1}f_{small, *}\mathcal{F} \to \mathcal{F}$는 전사이다. 층 (집합 또는 아벨 군) $\mathcal{F}$ 위의 $X_\etale$, \item 함자 $f_{small, *}$는 충실합니다 (위의 층 중 집합 및 아벨 층)에서 \item $f_{small, *} :
\textit{Ab}(X_\etale)
\to
\textit{Ab}(Y_\etale)$는 완전하다. \item 함자 $Y_\etale \to X_\etale$, $V \mapsto X \times_Y V$는 거의 여연속이다. \end{enumerate} \end{lemma}

\begin{proof}
에 의해
기본형 \ref{etale-cohomology-lemma-integral-A},
\ref{etale-cohomology-lemma-integral-B} 그리고
\ref{etale-cohomology-lemma-integral-homeo-onto-image-C}
우리는 사상 $f$에 성질 (), (B), (C)가 있다는 것을 알고 있다.
게다가,
보조정리 \ref{etale-cohomology-lemma-property-C-closed-implies}
우리는 함자 $Y_\etale \to X_\etale$가
거의 연속적이다. 이제 우리는
\begin{enumerate}
\item 성질 (C)는 다음과 같이 (1)를 의미한다.
보조정리 \ref{etale-cohomology-lemma-property-C-implies},
\item 거의 연속적인 의미는 (2)
사이트, 보조정리 \ref{sites-lemma-morphism-of-sites-almost-cocontinuous},
\item 성질 (B)는 (3)를 다음과 같이 의미한다.
보조정리 \ref{etale-cohomology-lemma-property-B-implies}.
\end{enumerate}
성질 (4), (5) 및 (6)는 공식적으로 처음 세 개를 따릅니다.
사이트, 보조정리 \ref{sites-lemma-exactness-properties}
그리고
가군 사이트, 보조정리 \ref{sites-modules-lemma-exactness}.
성질 (7) 위에서 본 것이다.
\end{proof}




\section{작은 에탈의 위상적 불변성 사이트} \label{etale-cohomology-section-topological-invariance}

\noindent
다음 정리에서 작은 에탈 사이트가 위상수학임을 보여줍니다.
다음 의미에서 불변: $f : X \to Y$가 스킴의 사상인 경우
이것이 보편적 동형이라면 $X_\etale \cong Y_\etale$
사이트로. 이는 다음의 결과를 향상시킵니다.
에탈 사상, 정리 \ref{etale-theorem-remarkable-equivalence}.
먼저 사상에 대한 결과를 증명한 다음 결과를 설명한다.
범주의 경우.

\begin{theorem} \label{etale-cohomology-theorem-etale-topological} $X$ 및 $Y$를 밑수 스킴 $S$ 위에 두 개의 스킴으로 둡니다. $S' \to S$를 보편적인 동형이라고 하자. $X'$ (각각.\ $Y'$)를 $S'$로 변경한다. $X$가 $S$에 대해 에탈인 경우 맵 $$
\Mor_S(Y, X) \longrightarrow \Mor_{S'}(Y', X')
$$는 전단사적 맵이다. \end{theorem}

\begin{proof} $Y \to S$을 통해 염기를 변경한 후에는 $Y = S$라고 가정할 수 있다. 따라서 우리는 $X$와 $Y$ 둘 다 $S$에 대해 에탈이라고 가정할 수 있고 가정한다. 즉, 정리는 기본 변경 함자가 스킴의 범주에서 $S$에 대한 범주에 완전히 충실한 함자임을 나타냅니다. 스킴 $S'$에 대한 에탈.

\medskip\noindent 잊어버린 함자 \begin{equation}
\label{etale-cohomology-equation-descent-etale-forget}
\begin{matrix}
\text{강하 데이터 }(X', \varphi')\text{(}S'/S\text{에 관한)} \\
\text{단, }X'\text{는 }S'\text{ 위에서 에탈}
\end{matrix}
\longrightarrow
\text{스킴 }X'\text{(}S'\text{ 위에서 에탈)}
\end{equation} 우리는 주장 이 함자는 등가물이다. 한편, 함자 \begin{equation}
\label{etale-cohomology-equation-descent-etale}
\text{스킴 }X\text{(}S\text{ 위에서 에탈)} \longrightarrow
\begin{matrix}
\text{강하 데이터 }(X', \varphi')\text{(}S'/S\text{에 관한)} \\
\text{단, }X'\text{는 }S'\text{ 위에서 에탈}
\end{matrix}
\end{equation}는 에탈 사상, 보조정리에 의해 완전히 충실한다. \ref{etale-lemma-fully-faithful-cases}. 따라서 주장은 정리를 의미한다.

\medskip\noindent
주장의 증명.
보편적인 동형은 다음과 같다는 것을 기억하세요.
정역, 보편적으로 단사, 전사 사상, 참조
사상, 보조정리 \ref{morphisms-lemma-universal-homeomorphism}.
특히 대각선 $\Delta : S' \to S' \times_S S'$이 두꺼워지고 있다.
사상, 보조정리 \ref{morphisms-lemma-universally-injective} 기준.
따라서 에탈 사상, 정리
\ref{etale-theorem-etale-topological}
우리는 주어진 $X' \to S'$ 에탈에 고유한 동형사상이 있음을 알 수 있다.
$$
\varphi' : X' \times_S S' \to S' \times_S X'
$$
의 스킴의 에탈 위의 $S' \times_S S'$ 아래로 물러남
$\Delta$ ~ $\text{id} : X' \to X'$, $S'$.
이후 $S' \to S' \times_S S' \times_S S'$
두껍게 하는 것이기도 합니다(전단사적이며 닫힌 몰입).
우리는 $(X', \varphi')$가 $S'/S$에 상대적인 하강 기준점이라고 결론을 내립니다.
$\varphi'$ 구성의 정식 성격은 다음과 같습니다.
$S'$에 대해 스킴의 에탈 사이에서 사상과 호환된다.
즉, 준역(quasi-inverse)을 얻습니다.
함자의 $X' \mapsto (X', \varphi')$
(\ref{etale-cohomology-equation-descent-etale-forget}). 이는 주장과
정리의 증명을 완료한다.
\end{proof}

\begin{theorem} \label{etale-cohomology-theorem-topological-invariance} \begin{reference} \cite[IV Theorem 18.1.2]{EGA} \end{reference}
$f : X \to Y$를 스킴의 사상으로 둡니다. $f$는 정역, 보편적으로 단사 및 전사라고 가정합니다(즉, $f$는 보편적 동형사상이다. 사상, 보조정리 \ref{morphisms-lemma-universal-homeomorphism} 참조). 함자 $$
V \longmapsto V_X = X \times_Y V
$$는 범주의 동치 $$
\{
\text{스킴 }V\text{(}Y\text{ 위에서 에탈)}
\}
\leftrightarrow
\{
\text{스킴 }U\text{(}X\text{ 위에서 에탈)}
\}
$$ \end{theorem}를 정의한다.

\noindent 우리는 두 가지 증거를 제시한다. 첫 번째는 전사 정역 사상에 상대적인 준 컴팩트, 분리, 에탈 사상에 대한 하강의 유효성을 사용한다. 두 번째는 이 장의 앞부분에서 논의한 성질 (), (B), (C)의 자료를 사용한다.

\begin{proof}[첫 번째 증명] 정리 \ref{etale-cohomology-theorem-etale-topological}에 의해 함자가 완전히 충실하다는 것을 알 수 있다. 함자가 본질적으로 전사임을 보여주는 것이 남아 있다. $U \to X$를 스킴의 에탈 사상라고 하자.

\medskip\noindent
$U$와 $Y$가 유사하면 결과가 유지된다고 가정한다.
이 경우, 우리는 개방형 아핀 피복을 선택한다.
$U = \bigcup U_i$ 각 $U_i$가 매핑되도록
$Y$의 아핀 오픈으로. 가정 (아핀 케이스)에 의해 우리는
에탈 사상 $V_i \to Y$를 찾아 $X \times_Y V_i \cong U_i$
$X$에 대해 스킴으로 표시된다. $V_{i, i'} \subset V_i$
기본 위상 공간이 있는 개방형 하위 체계여야 한다.
$U_i \cap U_{i'}$에 해당한다. 왜냐하면 우리는 동형사상
$$
X \times_Y V_{i, i'} \cong U_i \cap U_{i'} \cong X \times_Y V_{i', i}
$$
스킴 $X$에 대해 우리는 우리가
동형사상 획득
$\theta_{i, i'} : V_{i, i'} \to V_{i', i}$의 스킴 위의 $Y$.
우리는 이러한 동형사상이 다음을 만족하는지 확인을 생략한다.
스킴, 절 \ref{schemes-section-glueing-schemes}의 조건.
스킴, 보조정리 \ref{schemes-lemma-glue-schemes} 적용
우리는 다음과 같이 스킴 $V \to Y$를 얻습니다.
접착스킴 $V_i$신분증을 따라$\theta_{i, i'}$.
$V \to Y$는 에탈이고 $X \times_Y V \cong U$임이 분명한다.
건설로.

\medskip\noindent
따라서 $U$과 $Y$이 유사할 경우에는 보조정리를 표시하는 것으로 충분한다.
정리 \ref{etale-cohomology-theorem-etale-topological}의 증명에서 이를 기억하세요.
우리는 $U$가 고유한 하강 데이터와 함께 제공된다는 것을 보여주었습니다.
$(U, \varphi)$는 $X/Y$에 상대적이다. 에 의해
에탈 사상, 명제 \ref{etale-proposition-effective}
($U \to X$이 준컴팩트하고 분리되어 있기 때문에 적용된다.
아핀 케이스로의 축소에 의한 에탈도 마찬가지입니다)
에탈 사상 $V \to Y$이 존재한다.
$X \times_Y V \cong U$ 및 증명이 완료되었다.
\end{proof}

\begin{proof}[두 번째 증명] 정리 \ref{etale-cohomology-theorem-etale-topological}에 의해 함자가 완전히 충실하다는 것을 알 수 있다. 함자가 본질적으로 전사임을 보여주는 것이 남아 있다. $U \to X$를 스킴의 에탈 사상라고 하자.

\medskip\noindent
$U$와 $Y$가 유사하면 결과가 유지된다고 가정한다.
이 경우, 우리는 개방형 아핀 피복을 선택한다.
$U = \bigcup U_i$ 각 $U_i$가 매핑되도록
$Y$의 아핀 오픈으로. 가정 (아핀 케이스)에 의해 우리는
에탈 사상 $V_i \to Y$를 찾아 $X \times_Y V_i \cong U_i$
$X$에 대해 스킴으로 표시된다. $V_{i, i'} \subset V_i$
기본 위상 공간이 있는 개방형 하위 체계여야 한다.
$U_i \cap U_{i'}$에 해당한다. 왜냐하면 우리는 동형사상
$$
X \times_Y V_{i, i'} \cong U_i \cap U_{i'} \cong X \times_Y V_{i', i}
$$
스킴 $X$에 대해 우리는 우리가
동형사상 획득
$\theta_{i, i'} : V_{i, i'} \to V_{i', i}$의 스킴 위의 $Y$.
우리는 이러한 동형사상이 다음을 만족하는지 확인을 생략한다.
스킴, 절 \ref{schemes-section-glueing-schemes}의 조건.
스킴, 보조정리 \ref{schemes-lemma-glue-schemes} 적용
우리는 다음과 같이 스킴 $V \to Y$를 얻습니다.
접착스킴 $V_i$신분증을 따라$\theta_{i, i'}$.
$V \to Y$는 에탈이고 $X \times_Y V \cong U$임이 분명한다.
건설로.

\medskip\noindent 따라서 함자 \begin{equation}
\label{etale-cohomology-equation-affine-etale}
\{
\text{아핀 스킴 }V\text{(}Y\text{ 위에서 에탈)}
\}
\leftrightarrow
\{
\text{아핀 스킴 }U\text{(}X\text{ 위에서 에탈)}
\}
\end{equation}는 $X$와 $Y$가 유사할 때 본질적으로 전사이다.

\medskip\noindent
$U \to X$를 $X$에 대한 아핀 스킴의 에탈로 둡니다.
우리는 $V \to Y$ 에탈 (및 아핀)을 찾아야 $X \times_Y V$
와 동형이다$U$~ 위에$X$. 에탈사상아핀의
보편적으로 제한된 올을 가지고 있다.
사상,
기본형 \ref{morphisms-lemma-etale-locally-quasi-finite} 및
\ref{morphisms-lemma-locally-quasi-finite-qc-source-universally-bounded}.
따라서 우리는 올의 차수를 제한하는 정수 $n$에 대한 귀납법을 수행할 수 있다.
$U \to X$. 보다
사상, 보조정리 \ref{morphisms-lemma-etale-universally-bounded}
에탈 사상의 경우 이 정수에 대한 설명이다.
$n = 1$인 경우 $U \to X$는 열린 몰입입니다(참조:
에탈 사상, 정리 \ref{etale-theorem-etale-radicial-open}),
결과는 분명한다. $n > 1$를 가정한다.

\medskip\noindent
에 의해
보조정리 \ref{etale-cohomology-lemma-integral-homeo-onto-image-C}
스킴 $W \to Y$의 에탈 사상이 있고
전사사상 $W_X \to U$~ 위에$X$.
$U$은 준컴팩트하므로 $W$을 다음의 분리된 합집합으로 대체할 수 있다.
$W$의 유한한 수의 아핀 오픈, 따라서 $W$라고 가정할 수 있다.
역시 친해요. 여기 다이어그램이 있다.
$$
\xymatrix{
U \ar[d] & U \times_Y W \ar[l] \ar[d] & W_X \amalg R \ar@{=}[l]\\
X \ar[d] & W_X \ar[l] \ar[d] \\
Y & W \ar[l]
}
$$
분리된 결합 분해는 구성에 의해 발생한다.
아핀 스킴 $U \times_Y W \to W_X$의 에탈 사상에는 절이 있다.
자, 이제 사상 $R \to X \times_Y W$가 에탈이라는 것을 알 수 있다.
사상 아핀 스킴의 올 차수가 보편적으로 제한됨
$n - 1$로. 따라서 유도 가정에 의해 스킴이 존재한다.
$V' \to W$ 에탈은 $R \cong W_X \times_W V'$과 같습니다.
취득$V'' = W \amalg V'$우리는스킴 $V''$에탈 오버$W$누구의
$W_X$에 대한 기본 변경은 $U \times_Y W$와 동형이다.
$X \times_Y W$ 이상.

\medskip\noindent
이 시점에서 우리는 강하를 사용하여 다음을 찾을 수 있다.$V$~ 위에$Y$누구의 기지
로 변경$X$와 동형이다$U$~ 위에$X$. 즉, 완전히
함자 (\ref{etale-cohomology-equation-affine-etale})의 충실도
보편적 동형론에 해당
$X \times_Y (W \times_Y W) \to (W \times_Y W)$
고유한 동형사상 $\varphi : V'' \times_Y W \to W \times_Y V''$가 존재한다.
$X \times_Y (W \times_Y W)$에 대한 기본 변경이 표준이다.
하강 기준점$U \times_Y W$~ 위에$X \times_Y W$. 특히
$\varphi$는 코사이클 조건을 만족한다. 따라서
하강, 보조정리 \ref{descent-lemma-affine}
$\varphi$가 효과적이라는 것을 알 수 있습니다(위의 모든 스킴은 유사하다는 점을 기억하세요).
따라서 우리는 $V \to Y$ 및 동형사상 $V'' \cong W \times_Y V$를 얻습니다.
$W \times_Y V/W/Y$의 표준 하강 데이터가 일치하도록
$\varphi$로. $V \to Y$는 에탈이다.
하강, 보조정리 \ref{descent-lemma-descending-property-etale}.
게다가, 다음에서 오는 동형사상 $V_X \cong U$가 있다.
동형사상 내림차순
$$
V_X \times_X W_X =
X \times_Y V \times_Y W =
(X \times_Y W) \times_W (W \times_Y V) \cong
W_X  \times_W V'' \cong U \times_Y W
$$
우리는 건설을 통해 가지고 있다. 일부 세부 사항은 생략되었다.
\end{proof}

\begin{remark} \label{etale-cohomology-remark-affine-inside-equivalence} 정리 \ref{etale-cohomology-theorem-topological-invariance}의 상황에서 $V \mapsto V_X$는 사상 $V \to Y$와 $V$ 사이의 동등성을 유도하는 것도 사실이다. 아핀과 그 에탈 사상 $U \to X$는 $U$ 아핀이다. 극한, 명제 \ref{limits-proposition-affine}의 예에 대한 내용이다. \end{remark}

\begin{proposition}[에탈의 위상적 불변성 코호몰로지] \label{etale-cohomology-proposition-topological-invariance} $X_0 \to X$를 스킴의 보편적 동형이라고 하자(예에 대해 닫힌 몰입은 nilpotent에 의해 정의됨) 층 중 아이디얼). 그런 다음 \begin{enumerate} \item 에탈 사이트 $X_\etale$ 및 $(X_0)_\etale$는 동형이고, \item 에탈 토포스 $\Sh(X_\etale)$ 및 $\Sh((X_0)_\etale)$는 동등하며 \item $H^q_\etale(X, \mathcal{F}) = H^q_\etale(X_0, \mathcal{F}|_{X_0})$는 모든 $q$와 아벨 층 $\mathcal{F}$에 대해 $X_\etale$에 적용된다. \end{enumerate} \end{proposition}

\begin{proof} 범주의 동치 $X_\etale \to (X_0)_\etale$는 정리 \ref{etale-cohomology-theorem-topological-invariance}에 의해 제공된다. 이 등가 하에서 에탈 피복에 해당하는 증명을 생략한다. 따라서 (1)가 성립한다. 부분 (2) 및 (3)는 공식적으로 (1)를 따릅니다. \end{proof}






\section{닫힌 몰입 및 직상} \label{etale-cohomology-section-closed-immersions}

\noindent 명제 \ref{etale-cohomology-proposition-closed-immersion-pushforward}를 올바른 일반성으로 진술하고 증명하기 전에 닫힌 몰입에 대해 간략하게 진술하고 증명한다. 즉, 앞의 주장 중 일부는 닫힌 몰입의 경우 따라가기가 훨씬 더 쉽기 때문에 여기서는 단순화된 형식으로 반복한다.

\medskip\noindent 나머지 부분에서는 절 $i : Z \to X$는 닫힌 몰입이다. 함자 $$
\Sch/X \longrightarrow \Sch/Z, \quad
U \longmapsto U_Z = Z \times_X U
$$는 표시된 대로 $U \mapsto U_Z$로 표시된다. 닫힌 몰입은 임의의 기본 변경에도 유지되므로 스킴 $U_Z$는 $U$의 폐쇄형 하위 체계이다.

\begin{lemma}
\label{etale-cohomology-lemma-closed-immersion-almost-full}
$i : Z \to X$를 스킴의 닫힌 몰입으로 설정한다.
$U, U'$를 $X$에 대해 스킴의 에탈로 둡니다. $h : U_Z \to U'_Z$
$Z$보다 사상이어야 한다. 그런 다음 다이어그램이 존재한다.
$$
\xymatrix{
U & W \ar[l]_a \ar[r]^b & U'
}
$$
$X_\etale$에서 $a_Z : W_Z \to U_Z$
동형사상 및 $h = b_Z \circ (a_Z)^{-1}$이다.
\end{lemma}

\begin{proof}
스킴 $M = U \times_X U'$를 고려해보세요. 그래프 $\Gamma_h \subset M_Z$
$h$이(가) 열려 있다. 이는 예에 해당된다. $\Gamma_h$는 다음의 이미지이기 때문이다.
절 의 에탈 사상 $\text{pr}_{1, Z} : M_Z \to U_Z$, 참조
에탈 사상, 명제 \ref{etale-proposition-properties-sections}.
따라서 다음과 교차하는 개방형 하위 체계 $W \subset M$가 존재한다.
닫힌 부분집합 $M_Z$는 $\Gamma_h$이다. 집합 $a = \text{pr}_1|_W$
및 $b = \text{pr}_2|_W$.
\end{proof}

\begin{lemma} \label{etale-cohomology-lemma-closed-immersion-almost-essentially-surjective} $i : Z \to X$를 스킴의 닫힌 몰입으로 둡니다. $V \to Z$를 스킴의 에탈 사상라고 하자. 에탈 사상 $U_i \to X$와 사상 $U_{i, Z} \to V$가 존재하므로 $\{U_{i, Z} \to V\}$는 $V$의 자리스키 피복이다. \end{lemma}

\begin{proof}
스킴으로 구성된 $V$의 자리스키 피복만 찾으면 되므로
$U_Z$ 형식과 $U$ 에탈이 $X$에 걸쳐 있는 경우 Zariski는 $X$에 지역화할 수 있다.
및 $V$. 따라서 $X$ 및 $V$ 유사하다고 가정할 수 있다. 아핀의 경우 이는 다음과 같습니다.
대수학, 보조정리 \ref{algebra-lemma-lift-etale}.
\end{proof}

\noindent $\overline{x} : \Spec(k) \to X$가 $X$의 기하점인 경우 $\overline{x}$는 닫힌 하위 체계 $Z$ 또는 $Z_{\overline{x}} = \emptyset$를 통해 (고유하게) 인수분해된다. $\overline{x}$가 $Z$를 통해 고려된다면 $\overline{x}$는 $Z$의 기하점라고 말하고(그렇기 때문에) ``$\overline{x} \in Z$'' 표기법을 사용하여 이를 나타냅니다.

\begin{lemma} \label{etale-cohomology-lemma-stalk-pushforward-closed-immersion} $i : Z \to X$를 스킴의 닫힌 몰입으로 둡니다. $\mathcal{G}$를 $Z_\etale$에서 집합의 층으로 설정한다. $\overline{x}$를 $X$의 기하점으로 설정한다. 그런 다음 $$
(i_{small, *}\mathcal{G})_{\overline{x}} =
\left\{
\begin{matrix}
* & \text{if} & \overline{x} \not \in Z \\
\mathcal{G}_{\overline{x}} & \text{if} & \overline{x} \in Z
\end{matrix}
\right.
$$ 여기서 $*$는 싱글톤 집합을 나타냅니다. \end{lemma}

\begin{proof}
$U = X \setminus Z$로 놔두세요.
$i_{small, *}\mathcal{G}|_{U_\etale} = *$이 최종이다.
의 개체범주에탈의층~에$U$, 즉층
이는 싱글톤 집합을 $U$의 각 스킴의 에탈에 연결한다.
이것은 $(i_{small, *}\mathcal{G})_{\overline{x}}$의 값을 설명한다.
만약 $\overline{x} \not \in Z$.

\medskip\noindent
다음으로 $\overline{x} \in Z$라고 가정한다. 참고하세요
$$
(i_{small, *}\mathcal{G})_{\overline{x}}
=
\colim_{(U, \overline{u})} \mathcal{G}(U_Z)
$$
그리고 다른 한편으로는
$$
\mathcal{G}_{\overline{x}}
=
\colim_{(V, \overline{v})} \mathcal{G}(V).
$$
$\mathcal{C}_1 = \{(U, \overline{u})\}^{opp}$를 반대라고 하자.
범주에탈 지역 중$\overline{x}$~에$X$, 그리고
$\mathcal{C}_2 = \{(V, \overline{v})\}^{opp}$ 반대가 되다
범주에탈 지역 중$\overline{x}$~에$Z$. 표준 지도
$$
\mathcal{G}_{\overline{x}}
\longrightarrow
(i_{small, *}\mathcal{G})_{\overline{x}}
$$
함자 $F : \mathcal{C}_1 \to \mathcal{C}_2$에 해당한다.
$F(U, \overline{u}) = (U_Z, \overline{x})$. 지금
기본형 \ref{etale-cohomology-lemma-closed-immersion-almost-essentially-surjective} 및
\ref{etale-cohomology-lemma-closed-immersion-almost-full}
$\mathcal{C}_1$가 $\mathcal{C}_2$에서 동일함을 암시한다.
범주론, 정의 \ref{categories-definition-cofinal}.
따라서 표시된 화살표는 동형사상이다.
범주론, 보조정리 \ref{categories-lemma-cofinal}.
\end{proof}

\begin{proposition}
\label{etale-cohomology-proposition-closed-immersion-pushforward}
$i : Z \to X$를 스킴의 닫힌 몰입으로 설정한다.
\begin{enumerate}
\item 함자
$$
i_{small, *} :
\Sh(Z_\etale)
\longrightarrow
\Sh(X_\etale)
$$
완전히 충실하며 그 필수 이미지는 집합의 층이다.
$\mathcal{F}$ $X_\etale$의 $X \setminus Z$에 대한 제한은 다음과 같습니다.
$*$와 동형이고,
\item 함자
$$
i_{small, *} :
\textit{Ab}(Z_\etale)
\longrightarrow
\textit{Ab}(X_\etale)
$$
완전히 충실하며 그 핵심 이미지는 아벨 층이다.
$X_\etale$, 지지집합이 $Z$에 포함되어 있다.
\end{enumerate}
두 경우 모두 $i_{small}^{-1}$는 함자의 왼쪽 역수이다.
$i_{small, *}$.
\end{proposition}

\begin{proof}
집합 중 층의 경우를 살펴보자.
누구에게나층 $\mathcal{G}$~에$Z$그만큼사상
$i_{small}^{-1}i_{small, *}\mathcal{G} \to \mathcal{G}$
는 동형사상이다.
보조정리 \ref{etale-cohomology-lemma-stalk-pushforward-closed-immersion}
(그리고
정리 \ref{etale-cohomology-theorem-exactness-stalks}).
이는 공식적으로 $i_{small, *}$이 완전히 충실하다는 것을 의미한다.
사이트, 보조정리 \ref{sites-lemma-exactness-properties}.
$i_{small, *}\mathcal{G}|_{U_\etale} \cong *$임이 분명한다.
여기서 $U = X \setminus Z$. 반대로 $\mathcal{F}$라고 가정해 보자.
$X$에 있는 집합의 층이므로 $\mathcal{F}|_{U_\etale} \cong *$이다.
부속물 매핑을 고려하세요.
$$
\mathcal{F} \longrightarrow i_{small, *}i_{small}^{-1}\mathcal{F}
$$
결합
기본형 \ref{etale-cohomology-lemma-stalk-pushforward-closed-immersion} 및
\ref{etale-cohomology-lemma-stalk-pullback}
동형사상임을 알 수 있다. 이것으로 (1)의 증명이 완료된다.
(2)의 증명이 동일한다.
\end{proof}





\section{정역 보편적 단사 사상}
\label{etale-cohomology-section-integral-universally-injective}

\noindent 명제 \ref{etale-cohomology-proposition-closed-immersion-pushforward}의 일반 버전이다.

\begin{proposition}
\label{etale-cohomology-proposition-integral-universally-injective-pushforward}
$f : X \to Y$를 정역인 스킴의 사상으로 설정한다.
보편적으로 주사제이다.
\begin{enumerate}
\item 함자
$$
f_{small, *} :
\Sh(X_\etale)
\longrightarrow
\Sh(Y_\etale)
$$
완전히 충실하며 그 필수 이미지는 집합의 층이다.
$\mathcal{F}$ $Y_\etale$의 $Y \setminus f(X)$에 대한 제한은 다음과 같습니다.
$*$와 동형이고,
\item 함자
$$
f_{small, *} :
\textit{Ab}(X_\etale)
\longrightarrow
\textit{Ab}(Y_\etale)
$$
완전히 충실하며 그 핵심 이미지는 아벨 층이다.
$Y_\etale$, 지지집합이 $f(X)$에 포함되어 있다.
\end{enumerate}
두 경우 모두 $f_{small}^{-1}$는 함자의 왼쪽 역수이다.
$f_{small, *}$.
\end{proposition}

\begin{proof} $f$을 $$
\xymatrix{
X \ar[r]^h & Z \ar[r]^i & Y
}
$$로 인수분해할 수 있다. 여기서 $h$는 정역이고 보편적으로 단사와 전사이고 $i : Z \to Y$는 닫힌 몰입이다. 명제 \ref{etale-cohomology-proposition-closed-immersion-pushforward}를 $i$에 적용하고 정리 \ref{etale-cohomology-theorem-topological-invariance}를 $h$에 적용한다. \end{proof}









\section{큰 사이트 및 직상} \label{etale-cohomology-section-big}

\noindent 이 절에서는 스킴의 사상의 특정 유형에 대한 $f_{big, *}$에 대한 몇 가지 기술적 결과를 증명한다.

\begin{lemma}
\label{etale-cohomology-lemma-monomorphism-big-push-pull}
$\tau \in \{Zariski, \etale, smooth, syntomic, fppf\}$로 놔두세요.
$f : X \to Y$를 스킴의 단사 사상으로 설정한다.
그런 다음 표준 맵
$f_{big}^{-1}f_{big, *}\mathcal{F} \to \mathcal{F}$
는 층 $\mathcal{F}$에 대한 동형사상이다.
$(\Sch/X)_\tau$.
\end{lemma}

\begin{proof} 이 경우 함자 $(\Sch/X)_\tau \to (\Sch/Y)_\tau$는 연속적이고 불연속적이며 완전히 충실한다. 따라서 결과는 사이트, 보조정리 \ref{sites-lemma-back-and-forth}의 결과이다. \end{proof}

\begin{remark} \label{etale-cohomology-remark-push-pull-shriek} 보조정리 \ref{etale-cohomology-lemma-monomorphism-big-push-pull}의 상황에서 표준 맵 $\mathcal{F} \to f_{big}^{-1}f_{big!}\mathcal{F}$은 집합의 모든 층에 대해 동형사상이다. $\mathcal{F}$는 $(\Sch/X)_\tau$에 있다. 증명도 마찬가지다. 이는 아벨 군의 층에도 적용된다. 그러나 아벨 군의 층에 대한 함자 $f_{big!}$는 사이트, 절 \ref{sites-modules-section-exactness-lower-shriek}의 가군에 정의되어 있으며 일반적으로 다음과 다릅니다. 집합의 층에 $f_{big!}$. 아벨 군의 층에 대한 결과는 사이트, 보조정리 \ref{sites-modules-lemma-back-and-forth}의 가군에서 따릅니다. \end{remark}

\begin{lemma} \label{etale-cohomology-lemma-closed-immersion-cover-from-below} $f : X \to Y$를 스킴의 닫힌 몰입으로 둡니다. $U \to X$를 신토믹(각각.\ 매끄러운, 각각.\ 에탈) 사상라고 하자. 그런 다음 신토믹 (각각.\ 매끄러운, 각각.\ 에탈) 사상 $V_i \to Y$ 및 사상 $V_i \times_Y X \to U$가 존재하므로 $\{V_i \times_Y X \to U\}$은 Zariski이다. 피복 중 $U$. \end{lemma}

\begin{proof}
$\tau = syntomic$일 때 보조정리를 증명해 보자.
질문은 $U$에 대한 로컬이다. 따라서 우리는 $U$가 다음과 같다고 가정할 수 있다.
$Y$의 아핀으로 매핑되는 아핀 스킴.
따라서 우리는 다음 사례를 증명하는 것으로 축소한다.
$Y = \Spec(A)$, $X = \Spec(A/I)$ 및
$U = \Spec(\overline{B})$, 여기서
$A/I \to \overline{B}$ 신토믹 환 맵이어야 한다.
대수학으로, 보조정리 \ref{algebra-lemma-lift-syntomic}
$\overline{g}_i \in \overline{B}$ 요소를 찾을 수 있다.
그렇게
$\overline{B}_{\overline{g}_i} = A_i/IA_i$ 특정 신토믹 환 맵의 경우
$A \to A_i$.
이는 신토믹 사례에서 보조정리를 증명한다.
부드러운 케이스의 증명은 다음을 사용하는 것을 제외하면 동일한다.
대수학, 보조정리 \ref{algebra-lemma-lift-smooth}.
에탈의 경우 사용
대수학, 보조정리 \ref{algebra-lemma-lift-etale}.
\end{proof}

\begin{lemma}
\label{etale-cohomology-lemma-prepare-closed-immersion-almost-cocontinuous}
$f : X \to Y$를 스킴의 닫힌 몰입으로 설정한다.
$\{U_i \to X\}$를 신토믹(각각.\ 매끄러운, 각각.\ 에탈) 피복으로 가정한다.
신토믹이 존재한다 (각각.\ 매끄러운, 각각.\ 에탈) 피복 $\{V_j \to Y\}$
각 $j$에 대해 $V_j \times_Y X = \emptyset$ 또는
사상 $V_j \times_Y X \to X$요인을 통해$U_i$어떤 사람들에게는$i$.
\end{lemma}

\begin{proof} 각 $i$에 대해 신토믹 (각각.\ 매끄러운, 각각.\ 에탈) 사상 $g_{ij} : V_{ij} \to Y$ 및 사상을 선택할 수 있다. $V_{ij} \times_Y X \to U_i$ 오버 $X$, $\{V_{ij} \times_Y X \to U_i\}$는 Zariski 피복이다. 보조정리 \ref{etale-cohomology-lemma-closed-immersion-cover-from-below}을 참조하라. 이는 특히 $\bigcup_{ij} g_{ij}(V_{ij})$에 닫힌 부분집합 $f(X)$가 포함되어 있음을 의미한다. 따라서 신토믹 (각각.\ 매끄러운, 각각.\ 에탈)은 $g_{ij}$를 열린 몰입 $Y \setminus f(X) \to Y$와 함께 매핑하여 원하는 신토믹 (각각.\ 매끄러운, 각각.\ 에탈) 피복 의 $Y$. \end{proof}

\begin{lemma}
\label{etale-cohomology-lemma-closed-immersion-almost-cocontinuous}
$f : X \to Y$를 스킴의 닫힌 몰입으로 설정한다.
$\tau \in \{syntomic, smooth, \etale\}$로 놔두세요.
함자 $V \mapsto X \times_Y V$는 거의
연속 함자 (참조
사이트, 정의 \ref{sites-definition-almost-cocontinuous})
$(\Sch/Y)_\tau \to (\Sch/X)_\tau$ 사이
큰 $\tau$ 사이트.
\end{lemma}

\begin{proof} 우리는 다음을 보여주어야 합니다: 사상 $V \to Y$ 및 모든 신토믹(각각.\ 매끄러움, 각각.\ 에탈) 피복 $\{U_i \to X \times_Y V\}$가 존재한다. (각각.\ 부드러움, 반응\ 에탈) 피복 $\{V_j \to V\}$ 각 $j$에 대해 $X \times_Y V_j$는 비어 있거나 $X \times_Y V_j \to Z \times_Y V$는 다음 중 하나를 통해 요인이 된다. $U_i$. 위의 보조정리 \ref{etale-cohomology-lemma-prepare-closed-immersion-almost-cocontinuous}를 닫힌 몰입 $X \times_Y V \to V$에 적용한 것이다. \end{proof}

\begin{lemma} \label{etale-cohomology-lemma-closed-immersion-pushforward-exact} $f : X \to Y$를 스킴의 닫힌 몰입으로 둡니다. $\tau \in \{syntomic, smooth, \etale\}$로 놔두세요. \begin{enumerate} \item 직상 $f_{big, *} :
\Sh((\Sch/X)_\tau)
\to
\Sh((\Sch/Y)_\tau)$는 코이퀄라이저 및 푸시아웃을 사용하여 통근한다. \item 직상 $f_{big, *} :
\textit{Ab}((\Sch/X)_\tau)
\to
\textit{Ab}((\Sch/Y)_\tau)$가 정확한다. \end{enumerate} \end{lemma}

\begin{proof}
이것은 다음과 같습니다
사이트, 보조정리 \ref{sites-lemma-morphism-of-sites-almost-cocontinuous},
가군 사이트에,
보조정리 \ref{sites-modules-lemma-morphism-ringed-sites-almost-cocontinuous},
그리고
보조정리 \ref{etale-cohomology-lemma-closed-immersion-almost-cocontinuous}
위에.
\end{proof}

\begin{remark}
\label{etale-cohomology-remark-fppf-closed-immersion-not-closed}
보조정리 \ref{etale-cohomology-lemma-closed-immersion-pushforward-exact} 경우 $\tau = fppf$
누락되었다. 그 이유는 환 $A$, 아이디얼 $I$ 및
충실히 평탄, 유한하게 제시된 환 지도 $A/I \to \overline{B}$, 저기
{\it 모든} 플랫이 유한하게 제시된 환을 찾을 수 있다고 생각할 이유가 없다.
$A \to B$를 $B/IB \not = 0$로 매핑하여 $A/I \to B/IB$가 다음을 고려하도록 한다.
$\overline{B}$. 따라서 증명
보조정리 \ref{etale-cohomology-lemma-closed-immersion-almost-cocontinuous}
fppf 위상에서는 작동하지 않습니다.
실제로는 다음이 거짓일 가능성이 높습니다.
$f_{big, *} : \textit{Ab}((\Sch/X)_{fppf})
\to \textit{Ab}((\Sch/Y)_{fppf})$
$f$가 닫힌 몰입일 때 정확한다.
예를 알고 계시다면 이메일을 보내주세요.
\href{mailto:stacks.project@gmail.com}{stacks.project@gmail.com}.
\end{remark}












\section{완전성 낮은 비명소리} \label{etale-cohomology-section-exactness-lower-shriek}

\noindent 이는 다음과 같은 기술적 결과일 뿐이다. 함자 $f_{big!}$는 일반적으로 코호몰로지와 컴팩트한 지지집합과 아무런 관련이 없다.

\begin{lemma} \label{etale-cohomology-lemma-exactness-lower-shriek} $\tau \in \{Zariski, \etale, smooth, syntomic, fppf\}$로 둡니다. $f : X \to Y$를 스킴의 사상으로 설정한다. $$
f_{big} :
\Sh((\Sch/X)_\tau)
\longrightarrow
\Sh((\Sch/Y)_\tau)
$$를 위상에서와 같이 토포스의 해당 사상으로 설정한다. 보조정리 \ref{topologies-lemma-morphism-big}, \ref{topologies-lemma-morphism-big-etale}, \ref{topologies-lemma-morphism-big-smooth}, \ref{topologies-lemma-morphism-big-syntomic} \ref{topologies-lemma-morphism-big-fppf}. \begin{enumerate} \item 함자 $f_{big}^{-1} : \textit{Ab}((\Sch/Y)_\tau) \to \textit{Ab}((\Sch/X)_\tau)$에는 정확한 왼쪽 수반 $$
f_{big!} : \textit{Ab}((\Sch/X)_\tau) \to \textit{Ab}((\Sch/Y)_\tau)
$$이 있다. \item 함자 $f_{big}^* :
\textit{Mod}((\Sch/Y)_\tau, \mathcal{O})
\to
\textit{Mod}((\Sch/X)_\tau, \mathcal{O})$에는 왼쪽 수반 $$
f_{big!} :
\textit{Mod}((\Sch/X)_\tau, \mathcal{O})
\to
\textit{Mod}((\Sch/Y)_\tau, \mathcal{O})
$$가 있는데 이는 정확한다. \end{enumerate} 게다가 두 함자 $f_{big!}$는 아벨 군의 기본 층에 동의한다. \end{lemma}

\begin{proof}
$f_{big}$는 토포스의 사상이다.
연속 및 동시연속 함자
$u : (\Sch/X)_\tau \to (\Sch/Y)_\tau$, $U/X \mapsto U/Y$.
게다가 $f_{big}^{-1}\mathcal{O} = \mathcal{O}$도 있다.
따라서 $f_{big!}$의 존재는 다음과 같습니다.
가군 사이트, 보조정리 \ref{sites-modules-lemma-g-shriek-adjoint},
각기
가군 사이트, 보조정리 \ref{sites-modules-lemma-lower-shriek-modules}.
$U$가 $(\Sch/X)_\tau$의 객체인 경우 함자
$u$는 범주의 동치를 유도한다.
$$
u' :
(\Sch/X)_\tau/U
\longrightarrow
(\Sch/Y)_\tau/U
$$
화살표의 양쪽이 $(\Sch/U)_\tau$와 같기 때문이다.
따라서 기본 아벨 층에 대한 $f_{big!}$의 합의는 다음과 같습니다.
의 토론에서 다음과 같습니다.
가군 사이트, 비고 \ref{sites-modules-remark-when-shriek-equal}.
$f_{big!}$의 완전성은 다음과 같습니다.
가군 사이트, 보조정리 \ref{sites-modules-lemma-exactness-lower-shriek}
올곱 및 이퀄라이저로 통근하는 위의 함자 $u$과 같습니다.
\end{proof}

\noindent 다음으로, 대수 스택과 연관된 다른 사이트의 가군층을 비교할 때 나중에 유용할 기술적인 보조정리를 증명한다.

\begin{lemma}
\label{etale-cohomology-lemma-compare-structure-sheaves}
$X$를 스킴으로 설정한다. 허락하다
$\tau \in \{Zariski, \etale, smooth, syntomic, fppf\}$.
$\mathcal{C}_1 \subset \mathcal{C}_2 \subset (\Sch/X)_\tau$를 꽉 차게 하세요.
다음 성질이 포함된 하위 카테고리:
\begin{enumerate}
\item 대상 $U/X$ 중 $\mathcal{C}_t$에 속하는 것에 대하여,
\begin{enumerate}
\item $\{U_i \to U\}$가 $(\Sch/X)_\tau$의 피복인 경우
$U_i/X$는 $\mathcal{C}_t$의 객체이고,
\item $U \times \mathbf{A}^1/X$는 $\mathcal{C}_t$의 객체이다.
\end{enumerate}
\item $X/X$는 $\mathcal{C}_t$의 객체이다.
\end{enumerate}
$\mathcal{C}_t$에 사이트의 구조를 부여한다. 그 구조는 피복이다.
정확히 피복 $\{U_i \to U\}$ $(\Sch/X)_\tau$
$U \in \Ob(\mathcal{C}_t)$. 그 다음에
\begin{enumerate}
\item[()] 함자 $\mathcal{C}_1 \to \mathcal{C}_2$
완전히 충실하고 연속적이며 동시 연속적이다.
\end{enumerate}
$g : \Sh(\mathcal{C}_1) \to \Sh(\mathcal{C}_2)$ 해당을 나타냅니다.
토포스의 사상. $\mathcal{O}_t$ $\mathcal{O}$의 제한을 나타냅니다.
$\mathcal{C}_t$로 변경된다. $g_!$의 함자를 나타냅니다.
가군 사이트, 정의 \ref{sites-modules-definition-g-shriek}.
\begin{enumerate}
\item[(b)] 표준 맵 $g_!\mathcal{O}_1 \to \mathcal{O}_2$
동형사상이다.
\end{enumerate}
\end{lemma}

\begin{proof}
주장 ()는 정의에서 바로 나옵니다.
이 증명 모든 스킴은 $X$ 위에 스킴이고 모든 사상은
스킴은 $X$에 대한 스킴의 사상이다. $g^{-1}$는
제한에 의해 주어지므로 객체에 대해$U$~의$\mathcal{C}_1$
$\mathcal{O}_1(U) = \mathcal{O}_2(U) = \mathcal{O}(U)$이 있다.
$g_!\mathcal{O}_1$는 전층에 연결된 층라는 점을 기억하세요.
$g_{p!}\mathcal{O}_1$에 연관되는$V$~에$\mathcal{C}_2$그룹
$$
\colim_{V \to U} \mathcal{O}(U)
$$
여기서 $U$는 $\mathcal{C}_1$의 개체 위로 실행되고 여극한은 다음과 같습니다.
아벨 군의 범주에서 촬영되었다. 아래에서는 자주 사용하겠습니다.
만약에
$$
V \to U \to U'
$$
사상이고 $U, U' \in \Ob(\mathcal{C}_1)$이다.
$f' \in \mathcal{O}(U')$가 $f \in \mathcal{O}(U)$로 제한되는 경우,
그런 다음 $(V \to U, f)$ 및 $(V \to U', f')$는
여극한. 또한 $g_!\mathcal{O}_1 \to \mathcal{O}_2$는 요소를 매핑한다.
$(V \to U, f)$는 단순히 $f$을 $V$로 철회하는 것이다.

\medskip\noindent
생존성. $V$를 스킴으로 하고 $h \in \mathcal{O}(V)$로 둡니다.
그런 다음 우리는사상 $V \to X \times \mathbf{A}^1$에 의해 유발$h$
및 구조 사상 $V \to X$. 글쓰기
$\mathbf{A}^1 = \Spec(\mathbf{Z}[x])$ 요소가 보이다.
$x \in \mathcal{O}(X \times \mathbf{A}^1)$ 당기기
$h$로 돌아갑니다. $X \times \mathbf{A}^1$는 $\mathcal{C}_1$의 객체이므로
(1)(b) 및 (2) 가정을 통해 원하는 전사성을 얻습니다.

\medskip\noindent 주입성. $V$를 스킴으로 설정한다. $s = \sum_{i = 1, \ldots, n} (V \to U_i, f_i)$를 위에 표시된 여극한의 요소로 가정한다. 모든 $i$에 대해 사상 $f_i : U_i \to X \times \mathbf{A}^1$를 사용하여 $(V \to U_i, f_i)$가 여극한의 동일한 요소를 $(f_i : V \to X \times \mathbf{A}^1, x)$로 정의하는지 확인할 수 있다. 그런 다음 $$
f_1 \times \ldots \times f_n : V \to X \times \mathbf{A}^n
$$를 고려하면 $s$가 여극한과 $$
\sum\nolimits_{i = 1, \ldots, n}
(f_1 \times \ldots \times f_n : V \to X \times \mathbf{A}^n, x_i) =
(f_1 \times \ldots \times f_n : V \to X \times \mathbf{A}^n,
x_1 + \ldots + x_n)
$$에서 동일하다는 것을 알 수 있다. 이제 $x_1 + \ldots + x_n$가 $V$에서 0으로 제한되면 다음을 통해 $f_1 \times \ldots \times f_n$ 요소를 볼 수 있다. $X \times \mathbf{A}^{n - 1} = V(x_1 + \ldots + x_n)$. 따라서 $s$는 여극한에서 0과 동일하다는 것을 알 수 있다. \end{proof}











%9.29.09
\section{에탈 코호몰로지}
\label{etale-cohomology-section-etale-cohomology}

\noindent 다음 절에서 우리는 에탈 코호몰로지에 대한 몇 가지 기본 결과를 증명한다. 여기에 에탈 코호몰로지에 대해서도 유지되는 위상 공간의 코호몰로지에 대해 우리가 알고 있는 것의 예가 있다.

\begin{lemma}[에탈을 위한 메이어-비에토리스 코호몰로지]
\label{etale-cohomology-lemma-mayer-vietoris}
$X$를 스킴으로 설정한다. $X = U \cup V$가 다음과 같다고 가정한다.
두 사람의 결합이 열립니다. 누구에게나아벨 층 $\mathcal{F}$~에$X_\etale$
길고 정확한 코호몰로지 시퀀스가 ​​존재합니다
$$
\begin{matrix}
0 \to
H^0_\etale(X, \mathcal{F}) \to
H^0_\etale(U, \mathcal{F}) \oplus H^0_\etale(V, \mathcal{F}) \to
H^0_\etale(U \cap V, \mathcal{F}) \phantom{\to \ldots} \\
\phantom{0} \to H^1_\etale(X, \mathcal{F}) \to
H^1_\etale(U, \mathcal{F}) \oplus H^1_\etale(V, \mathcal{F}) \to
H^1_\etale(U \cap V, \mathcal{F}) \to \ldots
\end{matrix}
$$
이 긴 완전열은 $\mathcal{F}$에서 기능적이다.
\end{lemma}

\begin{proof}
$\mathcal{I}$가 단사 아벨 층인 경우,
$$
0 \to \mathcal{I}(X) \to \mathcal{I}(U) \oplus \mathcal{I}(V) \to
\mathcal{I}(U \cap V) \to 0
$$
정확한다. 이는 $\mathcal{I}$과 같이 첫 번째와 중간 지점에 해당된다.
층이다. 오른쪽이 사실이다. 왜냐하면
$\mathcal{I}(U) \to \mathcal{I}(U \cap V)$는 전사이다
코호몰로지 사이트, 보조정리에
\ref{sites-cohomology-lemma-restriction-along-monomorphism-surjective}.
이를 증명하는 또 다른 방법은 여핵
$\mathcal{I}(U) \oplus \mathcal{I}(V) \to
\mathcal{I}(U \cap V)$는 첫 번째 čech 코호몰로지 군이다.
피복에 대하여 $\mathcal{I}$의
$X = U \cup V$ 사라지는
기본정리 \ref{etale-cohomology-lemma-hom-injective} 및 \ref{etale-cohomology-lemma-forget-injectives}.
따라서 $\mathcal{F} \to \mathcal{I}^\bullet$가 분사라면
해상도, 그럼
$$
0 \to \mathcal{I}^\bullet(X) \to
\mathcal{I}^\bullet(U) \oplus \mathcal{I}^\bullet(V) \to
\mathcal{I}^\bullet(U \cap V) \to 0
$$
복합체의 짧은 완전열이고 관련 긴
정확한 코호몰로지 시퀀스는 보조정리 명령문의 시퀀스이다.
\end{proof}

\begin{lemma}[상대 메이어-비에토리스]
\label{etale-cohomology-lemma-relative-mayer-vietoris}
$f : X \to Y$를 스킴의 사상으로 설정한다. $X = U \cup V$
두 개의 개방형 하위 계획을 결합한 것이다. 표시하다
$a = f|_U : U \to Y$, $b = f|_V : V \to Y$ 및
$c = f|_{U \cap V} : U \cap V \to Y$.
모든아벨 층 $\mathcal{F}$~에$X_\etale$
긴 완전열이 존재합니다
$$
0 \to
f_*\mathcal{F} \to
a_*(\mathcal{F}|_U) \oplus b_*(\mathcal{F}|_V) \to
c_*(\mathcal{F}|_{U \cap V}) \to
R^1f_*\mathcal{F} \to \ldots
$$
$Y_\etale$에 있다.
이 긴 완전열은 $\mathcal{F}$에서 기능적이다.
\end{lemma}

\begin{proof}
$\mathcal{F} \to \mathcal{I}^\bullet$를 단사 분해로 설정한다.
~의$\mathcal{F}$~에$X_\etale$. 우리주장우리는
복합체의 짧은 완전열을 얻습니다.
$$
0 \to
f_*\mathcal{I}^\bullet \to
a_*\mathcal{I}^\bullet|_U \oplus b_*\mathcal{I}^\bullet|_V \to
c_*\mathcal{I}^\bullet|_{U \cap V} \to
0.
$$
즉, 어떤 경우에도$W$~에$Y_\etale$, 그리고 어떤 경우에도$n \geq 0$그만큼
$W$에 대한 절 그룹의 해당 시퀀스
$$
0 \to
\mathcal{I}^n(W \times_Y X) \to
\mathcal{I}^n(W \times_Y U)
\oplus \mathcal{I}^n(W \times_Y V) \to
\mathcal{I}^n(W \times_Y (U \cap V)) \to
0
$$
보조정리 \ref{etale-cohomology-lemma-mayer-vietoris}의 증명에서 정확히 짧은 것으로 나타났습니다.
보조정리는 코호몰로지 층을 취하여 다음과 같은 사실을 사용한다.
$\mathcal{I}^\bullet|_U$는 $\mathcal{F}|_U$의 단사 분해이다.
$\mathcal{I}^\bullet|_V$, $\mathcal{I}^\bullet|_{U \cap V}$에 대해서도 마찬가지이다.
\end{proof}







\section{여극한}
\label{etale-cohomology-section-colimit}

\noindent $(\mathcal{F}_i, \varphi_{ii'})$가 사이트 $\mathcal{C}$에서 층의 다이어그램인 경우 해당 여극한(층의 범주에서)는 층화의 전층 $U \mapsto \colim_i \mathcal{F}_i(U)$. 사이트, 보조정리 \ref{sites-lemma-colimit-sheaves}를 참조하라. 시스템이 지시되면 $U$는 준컴팩트 객체에 의해 피복의 동종 시스템을 갖는 $\mathcal{C}$의 준컴팩트 객체이고, 그러면 $\mathcal{F}(U) = \colim \mathcal{F}_i(U)$, 사이트, 보조정리를 참조하라. \ref{sites-lemma-directed-colimits-sections}. 아벨 층의 여극한의 상위 코호몰로지 군을 처리하는 결과는 사이트, 보조정리 \ref{sites-cohomology-lemma-colim-works-over-collection}의 코호몰로지를 참조하라.

\medskip\noindent 코호몰로지 위의 사이트, 보조정리 \ref{sites-cohomology-lemma-colimit}에서 이 결과를 사이트의 역 시스템에서 층 시스템으로 일반화한다. 다음은 스킴의 역체계에 살고 있는 에탈 층 체계의 경우에 해당하는 개념이다.

\begin{definition}
\label{etale-cohomology-definition-inverse-system-sheaves}
$I$를 선주문된 집합으로 설정한다. $(X_i, f_{i'i})$를 역으로 하자
$I$에 대한 스킴 시스템이다.
층의 {\it 시스템 $(\mathcal{F}_i, \varphi_{i'i})$
$(X_i, f_{i'i})$}에 대해 다음과 같이 제공된다.
\begin{enumerate}
\item 층 $\mathcal{F}_i$ 중 $(X_i)_\etale$ 위의 것을 각 $i \in I$에 대하여 하나씩,
\item~을 위한$i' \geq i$지도
$\varphi_{i'i} : f_{i'i}^{-1}\mathcal{F}_i \to \mathcal{F}_{i'}$
$(X_{i'})_\etale$의 층 중
\end{enumerate}
$\varphi_{i''i} = \varphi_{i''i'} \circ f_{i'' i'}^{-1}\varphi_{i'i}$
$i'' \geq i' \geq i$마다.
\end{definition}

\noindent
정의 \ref{etale-cohomology-definition-inverse-system-sheaves}의 상황에서,
$I$는 방향이 있는 집합이고 전이는 사상 $f_{i'i}$ 유사하다고 가정한다.
허락하다$X = \lim X_i$될극한에서범주~의스킴, 보다
극한, 절 \ref{limits-section-limits}.
$f_i : X \to X_i$ 투영 사상을 표시하고 지도를 고려한다.
$$
f_i^{-1}\mathcal{F}_i = f_{i'}^{-1}f_{i'i}^{-1}\mathcal{F}_i
\xrightarrow{f_{i'}^{-1}\varphi_{i'i}}
f_{i'}^{-1}\mathcal{F}_{i'}
$$
이는 $f_i^{-1}\mathcal{F}_i$를 $X_\etale$에서 층 시스템으로 전환한다.
$I$ 이상 (이것을 확인하는 것은 좋은 연습입니다).
우리는 종종 동형사상이 있는지 알고 싶어한다.
$$
H^q_\etale(X, \colim f_i^{-1}\mathcal{F}_i) =
\colim H^q_\etale(X_i, \mathcal{F}_i)
$$
$X_i$이 준컴팩트하고 준분리된 경우 이는 사실임이 밝혀집니다.
모든 $i$에 대해서는 정리 \ref{etale-cohomology-theorem-colimit}를 참조하라.

\begin{lemma} \label{etale-cohomology-lemma-colimit-affine-sites} $I$를 집합으로 지정한다. $(X_i, f_{i'i})$는 아핀 전이 사상을 갖는 $I$에 대한 스킴의 역 시스템이라고 가정한다. $X = \lim_{i \in I} X_i$로 놔두세요. 위상과 같은 표기법을 사용하면 보조정리 \ref{topologies-lemma-alternative} $$
X_{affine, \etale} = \colim (X_i)_{affine, \etale}
$$는 사이트, 보조정리 \ref{sites-lemma-colimit-sites}의 의미에서 사이트로 표시된다. \end{lemma}

\begin{proof}
$X$과 $X_i$가 준컴팩트하고
모든 $i$에 대해 준분리됨(이는 다음의 모든 경우에 해당됨)
관심). 이 경우 어떤 대상의
$X_{affine, \etale}$, 각각 \ $(X_i)_{affine, \etale}$는
$X$에 대한 유한한 표현. 더욱이,
$X$에 대한 유한 표현의 스킴의 범주는
여극한의 범주의 스킴의 유한 표현
$X_i$ 이상, 참조
극한, 보조정리 \ref{limits-lemma-descend-finite-presentation}.
아핀 객체 에탈의 하위 범주에 대해서도 마찬가지이다.
$X$ 에 의해 극한, 보조정리
\ref{limits-lemma-limit-affine} 및 \ref{limits-lemma-descend-etale}.
마지막으로, $\{U^j \to U\}$가 $X_{affine, \etale}$의 피복인 경우
그리고 $U_i^j \to U_i$가 사상 의 아핀 스킴인 경우 에탈이 끝났습니다.
$X_i$의 $X$에 대한 기본 변경 사항은 $U^j \to U$이다. 그러면 다음을 알 수 있다.
$\{U^j_i \to U_i\}$의 일부 $X_{i'}$로의 기본 변경은 다음과 같습니다.
충분히 큰 $i'$에 대한 피복, 참조
극한, 보조정리 \ref{limits-lemma-descend-surjective}.

\medskip\noindent
일반적인 경우에는 $U$를 $X_{affine, \etale}$의 객체로 둡니다.
그러면 $U \to X$는 에탈이고 분리됩니다 ($U$가 분리됨에 따라)
그러나 일반적으로 준컴팩트하지는 않습니다. 그래도 $U \to X$는 로컬에 있다.
유한한 제시의 결과로
극한, 보조정리 \ref{limits-lemma-descend-finite-presentation-variant}
$i$, 준컴팩트하고 준분리된 스킴 $U_i$가 있고,
로컬로 유한하게 표현되는 사상 $U_i \to X_i$
$X$에 대한 기본 변경 사항은 $U \to X$이다. 그러면 $U = \lim_{i' \geq i} U_{i'}$
여기서 $U_{i'} = U_i \times_{X_i} X_{i'}$.
$i$를 증가시킨 후 $U_i$가 유사하다고 가정할 수 있다.
극한, 보조정리 \ref{limits-lemma-limit-affine}.
그것을 확인하려면$U_i \to X_i$에탈이다$i$충분히 크고,
유한 아핀 열기 선택 피복 $U_i = U_{i, 1} \cup \ldots \cup U_{i, m}$
$U_{i, j} \to U_i \to X_i$는 아핀 개방형으로 매핑된다.
$W_{i, j} \subset X_i$. 그럼 신청해볼까요
극한, 보조정리 \ref{limits-lemma-descend-etale}
$U_{i, j} \to W_{i, j}$이 에탈인지 확인하려면
$i$를 증가시킨 후.
이런 식으로 우리는 함자
$\colim (X_i)_{affine, \etale} \to X_{affine, \etale}$
본질적으로 전사적이다. 온전한 신실하심이 따르느니라
이미 사용된 것에서 직접
극한, 보조정리 \ref{limits-lemma-descend-finite-presentation-variant}.
피복에 대한 진술은 정확히 동일하게 증명된다.
증명의 첫 번째 단락에서와 같은 방식으로 수행된다.
\end{proof}

\noindent 위를 사용하여 여극한 및 코호몰로지에 대해 다음과 같은 일반적인 결과를 얻습니다.

\begin{theorem}
\label{etale-cohomology-theorem-colimit}
$X = \lim_{i \in I} X_i$를 스킴의 유도 시스템의 극한으로 설정한다.
아핀 전이 사상 $f_{i'i} : X_{i'} \to X_i$. 우리는 가정한다
$X_i$은 모든 $i \in I$에 대해 준간소하고 준분리되어 있다.
$(\mathcal{F}_i, \varphi_{i'i})$를 아벨 층의 시스템으로 설정한다.
$(X_i, f_{i'i})$에 있다. $f_i : X \to X_i$ 투영을 표시하고 집합
$\mathcal{F} = \colim f_i^{-1}\mathcal{F}_i$. 그 다음에
$$
\colim_{i\in I} H_\etale^p(X_i, \mathcal{F}_i) = H_\etale^p(X, \mathcal{F}).
$$
모든 $p \geq 0$에 대해.
\end{theorem}

\begin{proof}
위상별, 보조정리 \ref{topologies-lemma-alternative}
우리는 코호몰로지를 계산할 수 있습니다
~의$\mathcal{F}$~에$X_{affine, \etale}$.
따라서 다음의 조합에 의한 결과는
보조정리 \ref{etale-cohomology-lemma-colimit-affine-sites}
그리고
코호몰로지 사이트, 보조정리 \ref{sites-cohomology-lemma-colimit}.
\end{proof}

\noindent 다음 두 결과는 위의 정리의 특별한 경우이다.

\begin{lemma}
\label{etale-cohomology-lemma-colimit}
$X$를 준간소화하고 준분리된 스킴라고 하자. $I$
집합 지시를 받습니다. $(\mathcal{F}_i, \varphi_{ij})$를 시스템으로 설정
$X_\etale$의 $I$에 대한 아벨 층의. 그 다음에
$$
\colim_{i\in I} H_\etale^p(X, \mathcal{F}_i) = H_\etale^p(X,
\colim_{i\in I} \mathcal{F}_i).
$$
\end{lemma}

\begin{proof}
이는 정리 \ref{etale-cohomology-theorem-colimit}의 특별한 경우이다.
또한 직접 증명을 스케치한다.
우리는 모두를 위해 그것을 증명합니다$X$동시에 유도에 의해$p$.
\begin{enumerate}
\item
준컴팩트 및 준분리형 스킴 $X$ 및 에탈 피복
$\mathcal{U}$ 중 $X$, 개선 사항이 있음을 보여줍니다.
$\mathcal{V} = \{V_j \to X\}_{j\in J}$ $J$ 유한 및 각 $V_j$
준컴팩트하고 준분리되어 모든 것이
$V_{j_0} \times_X \ldots \times_X V_{j_p}$도 있다.
준컴팩트하고 준분리형이다.
\item
이전 단계와 범주의 여극한의 정의를 사용하여
층, 정리가 $p = 0$ 및 모든 $X$에 대해 유지됨을 보여줍니다.
\item
코호몰로지의 지역성 사용
(보조정리 \ref{etale-cohomology-lemma-locality-cohomology}),
Čech-투-코호몰로지 스펙트럼 열
(정리 \ref{etale-cohomology-theorem-cech-ss}) 그리고 유도 사실은
가설은 모두에 적용된다
$V_{j_0} \times_X \ldots \times_X V_{j_p}$
위 상황에서 유도 단계 $p \to p + 1$를 증명하라.
\end{enumerate}
\end{proof}

\begin{lemma}
\label{etale-cohomology-lemma-directed-colimit-cohomology}
$A$는 환, $(I, \leq)$는 방향이 있는 집합 및 $(B_i, \varphi_{ij})$는
$A$-대수 체계. 집합 $B = \colim_{i\in I} B_i$. $X \to \Spec(A)$
스킴의 준컴팩트하고 준분리 사상이어야 한다. 허락하다
$\mathcal{F}$ 아벨 층 - $X_\etale$.
$Y_i = X \times_{\Spec(A)} \Spec(B_i)$를 나타냅니다.
$Y = X \times_{\Spec(A)} \Spec(B)$,
$\mathcal{G}_i = (Y_i \to X)^{-1}\mathcal{F}$ 그리고
$\mathcal{G} = (Y \to X)^{-1}\mathcal{F}$. 그 다음에
$$
H_\etale^p(Y, \mathcal{G}) =
\colim_{i\in I} H_\etale^p (Y_i, \mathcal{G}_i).
$$
\end{lemma}

\begin{proof} 이는 정리 \ref{etale-cohomology-theorem-colimit}의 특별한 경우이다. 또한 다음과 같이 직접 증명의 개요를 설명한다. \begin{enumerate} \item $V \to Y$ 에탈이 $V$ 준밀착형이고 준분리형이면 $i\in I$과 $V_i \to Y_i$가 존재하며 $V = V_i \times_{Y_i} Y$와 같습니다. 고려된 모든 스킴이 아핀인 경우 이는 다음 대수 진술에 해당합니다: $B = \colim B_i$ 및 $B \to C$가 에탈이면 $i \in I$ 및 $B_i \to C_i$ 에탈이 존재하여 $C \cong B \otimes_{B_i} C_i$이 된다. 이는 대수학 보조정리 \ref{algebra-lemma-etale}에서 증명된다. \item (1) 상황에서 $\mathcal{G}(V) = \colim_{i' \geq i} \mathcal{G}_{i'}(V_{i'})$를 보여줍니다. 여기서 $V_{i'}$는 $V_i$에서 $Y_{i'}$로의 기본 변경이다. \item (1)에 의해, 우리는 모든 에탈 피복 $\mathcal{V} = \{V_j \to Y\}_{j\in J}$에 대해 $J$ 유한하고 $V_j$s 준컴팩트하고 준분리되어 있음을 알 수 있다. $i \in I$ 및 에탈 피복 $\mathcal{V}_i = \{V_{ij} \to Y_i\}_{j \in J}$는 $\mathcal{V} \cong \mathcal{V}_i \times_{Y_i} Y$이다. \item (2) 및 (3)가 $$
\check H^*(\mathcal{V}, \mathcal{G})=
\colim_{i\in I} \check H^*(\mathcal{V}_i, \mathcal{G}_i).
$$ \item 체흐-코호몰로지 스펙트럼 열(정리 \ref{etale-cohomology-theorem-cech-ss})을 영리하게 사용한다. \end{enumerate} \end{proof}

\begin{lemma} \label{etale-cohomology-lemma-higher-direct-images} $f: X\to Y$를 스킴 및 $\mathcal{F}\in
\textit{Ab}(X_\etale)$의 사상으로 설정한다. 그러면 $R^pf_*\mathcal{F}$는 전층 $$
(V \to Y) \longmapsto H_\etale^p(X \times_Y V, \mathcal{F}|_{X \times_Y V}).
$$에 연결된 층이다. 보다 일반적으로 $K \in D(X_\etale)$에 대해 $R^pf_*K$는 전층에 연결된 층이다. $$
(V \to Y) \longmapsto H_\etale^p(X \times_Y V, K|_{X \times_Y V}).
$$ \end{lemma}

\begin{proof} 이 보조정리는 위상 공간에 유효하며 이 경우의 증명도 같습니다. 층의 경우 사이트, 보조정리 \ref{sites-cohomology-lemma-higher-direct-images}에서 코호몰로지를 참조하고 사이트, 보조정리에서 코호몰로지를 참조하라. \ref{sites-cohomology-lemma-sheafification-cohomology}는 아벨 층의 복합체인 경우이다. \end{proof}

\begin{lemma}
\label{etale-cohomology-lemma-relative-colimit}
$S$를 스킴으로 설정한다. $X = \lim_{i \in I} X_i$를 극한으로 설정한다.
아핀 전이 사상이 있는 $S$에 대한 스킴의 방향 시스템
$f_{i'i} : X_{i'} \to X_i$. 사상 구조를 가정한다.
$g_i : X_i \to S$ 및 $g : X \to S$은 준 컴팩트하고 준 분리되어 있다.
$(\mathcal{F}_i, \varphi_{i'i})$를 아벨 층의 시스템으로 설정한다.
$(X_i, f_{i'i})$에 있다. $f_i : X \to X_i$ 투영을 표시하고 집합
$\mathcal{F} = \colim f_i^{-1}\mathcal{F}_i$. 그 다음에
$$
\colim_{i\in I} R^p g_{i, *} \mathcal{F}_i = R^p g_* \mathcal{F}
$$
모든 $p \geq 0$에 대해.
\end{lemma}

\begin{proof} 보조정리 \ref{etale-cohomology-lemma-higher-direct-images}) $R^p g_{i, *} \mathcal{F}_i$는 전층 $U \mapsto H^p_\etale(U \times_S X_i, \mathcal{F}_i)$와 연관된 층이고 $R^pg_*\mathcal{F}$에 대해서도 유사하다는 점을 상기하라. 또한, 층 시스템의 여극한은 전층 레벨의 여극한의 층화이다. $S_\etale$의 모든 객체는 준 컴팩트하고 준분리된 객체(예: 아핀 스킴)에 의해 피복을 갖습니다. 게다가, $U$가 준컴팩트하고 준분리된 객체라면, $$
\colim H^p_\etale(U \times_S X_i, \mathcal{F}_i) =
H^p_\etale(U \times_S X, \mathcal{F})
$$ 에 의해 정리 \ref{etale-cohomology-theorem-colimit}가 된다. 따라서 보조정리가 이어집니다. \end{proof}

\begin{lemma}
\label{etale-cohomology-lemma-relative-colimit-general}
$I$를 집합으로 지정한다. $g_i : X_i \to S_i$를 다음의 역계라고 하자.
사상의 스킴 위의 $I$. $g_i$이 준컴팩트하고
준분리 및 $i' \geq i$ 전이 사상
$f_{i'i} : X_{i'} \to X_i$와 $h_{i'i} : S_{i'} \to S_i$는 유사한다.
$g : X \to S$를 사상 $g_i$의 극한으로 설정한다.
극한, 절 \ref{limits-section-limits}.
$f_i : X \to X_i$ 및 $h_i : S \to S_i$ 투영을 나타냅니다.
$(\mathcal{F}_i, \varphi_{i'i})$를 층의 시스템으로 설정한다.
$(X_i, f_{i'i})$에 있다. 집합 $\mathcal{F} = \colim f_i^{-1}\mathcal{F}_i$. 그 다음에
$$
R^p g_* \mathcal{F} =
\colim_{i \in I} h_i^{-1}R^p g_{i, *} \mathcal{F}_i
$$
모든 $p \geq 0$에 대해.
\end{lemma}

\begin{proof} 보조정리의 지도는 어떻게 구성되나요? $i' \geq i$에 대해 가환 도식 $$
\xymatrix{
X \ar[r]_{f_{i'}} \ar[d]_g &
X_{i'} \ar[r]_{f_{i'i}} \ar[d]_{g_{i'}} &
X_i \ar[d]^{g_i} \\
S \ar[r]^{h_{i'}} &
S_{i'} \ar[r]^{h_{i'i}} &
S_i
}
$$가 있다. 기본 변경 맵 $h_{i'i}^{-1}Rg_{i, *}\mathcal{F}_i \to Rg_{i', *}f_{i'i}^{-1}\mathcal{F}_i$ (코호몰로지 위의 사이트, 보조정리 \ref{sites-cohomology-lemma-base-change-map-flat-case} 또는 비고 \ref{sites-cohomology-remark-base-change}) 지도 $Rg_{i', *}\varphi_{i'i}$를 사용하여 $\psi_{i'i} : h_{i' i}^{-1} R^p g_{i, *} \mathcal{F}_i \to
R^pg_{i', *} \mathcal{F}_{i'}$를 얻습니다. 마찬가지로 다이어그램의 왼쪽 사각형을 사용하여 맵 $\psi_i : h_i^{-1}R^pg_{i, *}\mathcal{F}_i \to R^pg_*\mathcal{F}$을 얻습니다. $h_{i'}^{-1}\psi_{i'i}$ 및 $\psi_i$ 맵은 보조정리의 문에 사용된 맵이다. 이를 이해하려면 $\psi_{i''i} = \psi_{i''i'} \circ h_{i''i'}^{-1}\psi_{i'i}$ 및 $\psi_{i'} \circ h_{i'}^{-1}\psi_{i'i} = \psi_i$를 확인해야 한다. 이는 사이트, 비고 \ref{sites-cohomology-remark-compose-base-change-horizontal}의 코호몰로지에서 이어집니다.

\medskip\noindent
증명 평등의. 먼저 증명을 사용하여
차원 이동\footnote{이 방법을 사용할 수도 있습니다
보조정리.}에서 맵을 생성한다. 누구에게나
$U$아핀 및 에탈 오버$X$~에 의해정리 \ref{etale-cohomology-theorem-colimit}
우리는
$$
g_*\mathcal{F}(U) =
H^0(U \times_S X, \mathcal{F}) =
\colim H^0(U_i \times_{S_i} X_i, \mathcal{F}_i) =
\colim g_{i, *}\mathcal{F}_i(U_i)
$$
여기서 여극한은 $i$보다 커서 다음과 같습니다.
존재한다$i$그리고$U_i$아핀 에탈 오버$S_i$
누구의 기본 변화는$U$~ 위에$S$(보다
보조정리 \ref{etale-cohomology-lemma-colimit-affine-sites}).
오른쪽은 다음과 같습니다.
$(\colim h_i^{-1}g_{i, *}\mathcal{F}_i)(U)$ 에 의해
사이트, 보조정리 \ref{sites-lemma-colimit}.
이는 $p = 0$에 대한 보조정리를 증명한다.
$(\mathcal{G}_i, \varphi_{i'i})$이 다음과 같은 시스템인 경우
$\mathcal{G} = \colim f_i^{-1}\mathcal{G}_i$
$\mathcal{G}_i$는 $X_i$의 단사 아벨 층이다.
모든 $i$에 대해, 모든 $U$ 아핀 및 $X$에 대한 에탈에 대해
정리 \ref{etale-cohomology-theorem-colimit} 우리는
$$
H^p(U \times_S X, \mathcal{G}) =
\colim H^p(U_i \times_{S_i} X_i, \mathcal{G}_i) = 0
$$
$p > 0$ (이전과 동일 여극한). 그러므로 $R^pg_*\mathcal{G} = 0$
그런 시스템에 대해 $p > 0$에 대한 결과를 얻습니다.
일반적으로 우리는 짧은 완전열 시스템을 선택할 수 있다.
$$
0 \to (\mathcal{F}_i, \varphi_{i'i}) \to
(\mathcal{G}_i, \varphi_{i'i}) \to
(\mathcal{Q}_i, \varphi_{i'i}) \to 0
$$
여기서 $(\mathcal{G}_i, \varphi_{i'i})$는 위와 같습니다.
코호몰로지 사이트, 보조정리에
\ref{sites-cohomology-lemma-colim-sites-injective}.
귀납법에 의해 보조정리는 $p - 1$에 대해 유지되고 위의 내용에 의해 우리는 다음을 얻습니다.
$p$ 및 $(\mathcal{G}_i, \varphi_{i'i})$의 경우 소멸이다.
따라서 $p$에 대한 결과는 다음과 같습니다.
그리고 긴 완전열에 의한 $(\mathcal{F}_i, \varphi_{i'i})$
코호몰로지.

\medskip\noindent
두 번째 증명.
$S_{affine, \etale} = \colim (S_i)_{affine, \etale}$를 기억하세요.
보조정리 \ref{etale-cohomology-lemma-colimit-affine-sites}. 따라서 $U$가 다음의 객체인 경우
$S_{affine, \etale}$, 그러면 $U = U_i \times_{S_i} S$라고 쓸 수 있다.
어떤 사람들에게는$i$그리고 일부$U_i$~에$(S_i)_{affine, \etale}$그리고
$$
(\colim_{i \in I} h_i^{-1}R^p g_{i, *} \mathcal{F}_i)(U) =
\colim_{i' \geq i} (R^p g_{i', *}\mathcal{F}_{i'})(U_i \times_{S_i} S_{i'})
$$
사이트, 보조정리 \ref{sites-lemma-colimit} 및 구성에 의해
위에서 설명한 시스템의 전환 맵이다. 부터
$R^pg_{i', *}\mathcal{F}_{i'}$는 전층에 연결된 층이다.
$U_{i'} \mapsto H^p(U_{i'} \times_{S_{i'}} X_{i'}, \mathcal{F}_{i'})$
그리고 $R^pg_*\mathcal{F}$는 전층과 연관된 층이기 때문에
$U \mapsto H^p(U \times_S X, \mathcal{F})$
(보조정리 \ref{etale-cohomology-lemma-higher-direct-images})
우리는 정식 가환 도식을 얻습니다.
$$
\xymatrix{
\colim_{i' \geq i}
H^p(U_i \times_{S_i} X_{i'}, \mathcal{F}_{i'}) \ar[r] \ar[d] &
\colim_{i' \geq i}
(R^p g_{i', *}\mathcal{F}_{i'})(U_i \times_{S_i} S_{i'}) \ar[d] \\
H^p(U \times_S X, \mathcal{F}) \ar[r] &
R^pg_*\mathcal{F}(U)
}
$$
왼쪽 수직 화살표가 동형사상인지 확인하세요.
정리 \ref{etale-cohomology-theorem-colimit} 기준. 우리는 그것을 보여주려고 노력하고 있어요
오른쪽 수직 화살표는 동형사상이다. 그러나 우리는
이 화살의 출처와 목표는 이미 알고 있습니다
$S_{affine, \etale}$에 층이 있다. 그러므로 충분하다
보이자: (1) 대상의 요소는 로컬에서 가져옵니다.
소스의 요소 및 (2) 소스의 요소
대상에서 0으로 매핑되는 것은 로컬로 사라집니다.
부분(1)은 위의 내용과 바로 이어지며 다음과 같은 사실이 나옵니다.
아래쪽 수평 화살표는 전층 지도에서 나옵니다.
이는 층화 다음에 동형사상이 된다.
부품(2)의 경우 $\xi \in  \colim_{i' \geq i}
(R^p g_{i', *}\mathcal{F}_{i'})(U_i \times_{S_i} S_{i'})$
핵에 있다. $i' \geq i$을 선택하고
$\xi_{i'} \in (R^p g_{i', *}\mathcal{F}_{i'})(U_i \times_{S_i} S_{i'})$
$\xi$을 나타냅니다.
표준을 선택하세요 에탈 피복
$\{U_{i', k} \to U_i \times_{S_i} S_{i'}\}_{k = 1, \ldots, m}$
$\xi_{i'}|_{U_{i', k}}$는 다음에서 유래한다.
$\xi_{i', k} \in H^p(U_{i', k} \times_{S_{i'}} X_{i'}, \mathcal{F}_{i'})$.
$\xi$이 로컬에서 죽는다는 것을 증명하는 것으로 충분하므로,
$U$을 에탈 멤버로 대체할 수 있다.
피복 $\{U_{i', k} \times_{S_{i'}} S \to U = U_i \times_{S_i} S\}$.
이 교체 후에는 $\xi$이 다음 이미지임을 알 수 있다.
그룹의 요소 $\xi'$
$\colim_{i' \geq i} H^p(U_i \times_{S_i} X_{i'}, \mathcal{F}_{i'})$
위의 다이어그램에서. $\xi'$는 $R^pg_*\mathcal{F}(U)$에서 0으로 매핑되므로
또 다른 교체를 수행하고 $\xi'$ 맵을 가정할 수 있다.
$H^p(U \times_S X, \mathcal{F})$에서 0으로 설정한다.
그러나 왼쪽 수직 화살표는 동형사상이므로
그런 다음 원하는 대로 $\xi' = 0$ 결론을 내립니다. 따라서 $\xi = 0$이다.
\end{proof}

\begin{lemma} \label{etale-cohomology-lemma-linus-hamann} $X = \lim_{i \in I} X_i$를 아핀 전이 사상 $f_{i'i}$ 및 투영 사상을 사용하여 스킴의 방향성 극한으로 설정한다. $f_i : X \to X_i$. $\mathcal{F}$를 $X_\etale$에서 층으로 설정한다. 그런 다음 \begin{enumerate} \item 표준 맵 $\varphi_{i'i} : f_{i'i}^{-1}f_{i, *}\mathcal{F} \to f_{i', *}\mathcal{F}$이 있으며 $(f_{i, *}\mathcal{F}, \varphi_{i'i})$는 $(X_i, f_{i'i})$에서 정의 \ref{etale-cohomology-definition-inverse-system-sheaves}의 층 시스템이다. \item $\mathcal{F} = \colim f_i^{-1}f_{i, *}\mathcal{F}$. \end{enumerate} \end{lemma}

\begin{proof} 위상을 통해 보조정리 \ref{topologies-lemma-alternative} 및 보조정리 \ref{etale-cohomology-lemma-colimit-affine-sites} 이는 사이트, 보조정리 \ref{sites-lemma-colimit-push-pull}의 특별한 경우이다. \end{proof}

\begin{lemma} \label{etale-cohomology-lemma-compute-strangely} $I$를 집합으로 지정한다. $g_i : X_i \to S_i$를 $I$에 대한 스킴의 사상의 역 시스템이라고 가정한다. $g_i$는 준컴팩트하고 준분리되어 있고 $i' \geq i$의 경우 전이 사상 $X_{i'} \to X_i$ 및 $S_{i'} \to S_i$는 유사하다고 가정한다. $g : X \to S$를 사상 $g_i$의 극한으로 설정한다. 극한, 절 \ref{limits-section-limits}를 참조하라. $f_i : X \to X_i$ 및 $h_i : S \to S_i$ 투영을 나타냅니다. $\mathcal{F}$를 $X$에서 아벨 층으로 설정한다. 그런 다음 $$
R^pg_*\mathcal{F} = \colim_{i \in I} h_i^{-1}R^pg_{i, *}(f_{i, *}\mathcal{F})
$$ \end{lemma}가 있다.

\begin{proof} 기본정리 \ref{etale-cohomology-lemma-relative-colimit-general} 및 \ref{etale-cohomology-lemma-linus-hamann}의 형식적인 조합이다. \end{proof}











\section{여극한 및 복합체} \label{etale-cohomology-section-colimit-complexes}

\noindent 이 절에서는 복합체 시스템의 코호몰로지를 다양한 설정으로 사용하는 방법에 대해 논의하고, 절 \ref{etale-cohomology-section-colimit}에서 시작된 층에 대한 논의를 계속한다. 우리는 꼭 필요한 경우가 아니라면 독자 여러분이 이 절을 읽지 말 것을 강력히 권고한다.

\begin{lemma}
\label{etale-cohomology-lemma-colimit-variant-complexes}
$X = \lim_{i \in I} X_i$를 아핀 전이 사상
$f_{i'i} : X_{i'} \to X_i$를 갖는 스킴 유향계의 극한이라 하자.
$X_i$가 모든 $i \in I$에 대해 준콤팩트하고 준분리라고 가정하자.
$\mathcal{F}_i^\bullet$를 $X_{i, \etale}$ 위의 아벨 층의 복합체라 하자.
$\varphi_{i'i} : f_{i'i}^{-1}\mathcal{F}_i^\bullet \to
\mathcal{F}_{i'}^\bullet$를 $X_{i, \etale}$ 위의 복합체 사상이라 하고,
$\varphi_{i''i} = \varphi_{i''i'} \circ f_{i'' i'}^{-1}\varphi_{i'i}$가
$i'' \geq i' \geq i$일 때마다 성립한다고 하자. 어떤 정수 $a$가 존재하여
$\mathcal{F}_i^n = 0$이 $n < a$ 및 모든 $i \in I$에 대해 성립한다고 가정하자. 그러면
$$
H^p_\etale(X, \colim f_i^{-1}\mathcal{F}_i^\bullet) =
\colim H^p_\etale(X_i, \mathcal{F}^\bullet_i)
$$
여기서 $f_i : X \to X_i$는 투영이다.
\end{lemma}

\begin{proof}
이는 정리 \ref{etale-cohomology-theorem-colimit}의 결과이다. 집합
$\mathcal{F}^\bullet = \colim f_i^{-1}\mathcal{F}_i^\bullet$.
정리는 다음을 알려줍니다.
$$
\colim_{i \in I} H_\etale^p(X_i, \mathcal{F}_i^n) =
H_\etale^p(X, \mathcal{F}^n)
$$
모든 $n, p \in \mathbf{Z}$에 대해. 스펙트럼 열을 사용해보자
$$
E_{1, i}^{s, t} = H_\etale^t(X_i, \mathcal{F}_i^s) \Rightarrow
H_\etale^{s + t}(X_i, \mathcal{F}_i^\bullet)
$$
그리고
$$
E_1^{s, t} = H_\etale^t(X, \mathcal{F}^s) \Rightarrow
H_\etale^{s + t}(X, \mathcal{F}^\bullet)
$$
유도 범주, 보조정리 \ref{derived-lemma-two-ss-complex-functor}.
부터$\mathcal{F}_i^n = 0$~을 위한$n < a$(와 함께$a$독립하다$i$)
우리는 고정된 유한 수의 용어 $E_{1, i}^{s, t}$만 있음을 알 수 있다.
($i$와 무관) 및 $E_1^{s, t}$에 기여
$H^q_\etale(X_i, \mathcal{F}_i^\bullet)$ 그리고
$H^q_\etale(X, \mathcal{F}^\bullet)$ 및 $E_1^{s, t} = \colim E_{i, i}^{s, t}$.
이는 우리가 원하는 것을 의미한다. 일부 세부 사항은 생략되었다.
(``어리석은'' 잘림을 사용하는 또 다른 주장이 있다.
스펙트럼 열 사용을 방지하는 복합체.)
\end{proof}

\begin{lemma}
\label{etale-cohomology-lemma-direct-sum-bounded-below-cohomology}
$X$를 준컴팩트하고 준분리된 스킴라고 하자. 허락하다
$K_i \in D(X_\etale)$, $i \in I$는 객체 계열이다.
주어진 것으로 가정$a \in \mathbf{Z}$그렇게$H^n(K_i) = 0$~을 위한$n < a$
및 $i \in I$. 그러면 $R\Gamma(X, \bigoplus_i K_i) =
\bigoplus_i R\Gamma(X, K_i)$.
\end{lemma}

\begin{proof}
$H^p(X, \bigoplus_i K_i) =
\bigoplus_i H^p(X, K_i)$가 모든 $p \in \mathbf{Z}$에 대해
성립함을 보이면 충분하다. 복합체 $\mathcal{F}_i^\bullet$로 $K_i$를 나타내되,
$\mathcal{F}_i^n = 0$이 $n < a$에 대해 성립하도록 선택하자.
단사 대상, 보조정리 \ref{injectives-lemma-derived-products}에 의해
복합체 $\mathcal{F}_i^\bullet$의 직합은 대상 $\bigoplus K_i$를 나타낸다.
$\bigoplus \mathcal{F}^\bullet$은 유한 직합들의 여과 여극한이므로,
결론은 보조정리 \ref{etale-cohomology-lemma-colimit-variant-complexes}에서 따른다.
\end{proof}

\begin{lemma}
\label{etale-cohomology-lemma-colimit-complexes}
$S$를 스킴으로 설정한다. $X = \lim_{i \in I} X_i$를 극한으로 설정한다.
아핀 전이 사상이 있는 $S$에 대한 스킴의 방향 시스템
$f_{i'i} : X_{i'} \to X_i$. 우리는 $X_i$가 준컴팩트하고
모든 $i \in I$에 대해 준분리된다. $K \in D^+(S_\etale)$로 놔두세요. 그 다음에
$$
\colim_{i \in I} H_\etale^p(X_i, K|_{X_i}) = H_\etale^p(X, K|_X).
$$
모든 $p \in \mathbf{Z}$에 대해 $K|_{X_i}$ 및 $K|_X$
$K$에서 $X_i$ 및 $X$로의 철회이다.
\end{lemma}

\begin{proof}
우리는 대표할 수 있습니다$K$아래의 경계로복합체 $\mathcal{G}^\bullet$
~의아벨 층~에$S_\etale$. 말하다$\mathcal{G}^n = 0$~을 위한$n < a$.
하락세를 $\mathcal{F}^\bullet_i$ 및 $\mathcal{F}^\bullet$로 표시
이 복합체 중 $X_i$ 및 $X$이다. 이 복합체는 다음을 나타냅니다.
객체 $K|_{X_i}$ 및 $K|_X$ 그리고 우리는
$\mathcal{F}^\bullet = \colim f_i^{-1}\mathcal{F}_i^\bullet$
텀별로. 따라서 보조정리는 다음과 같습니다.
보조정리 \ref{etale-cohomology-lemma-colimit-variant-complexes}.
\end{proof}

\begin{lemma}
\label{etale-cohomology-lemma-relative-colimit-general-complexes}
$I$, $g_i : X_i \to S_i$, $g : X \to S$, $f_i$, $g_i$, $h_i$는 다음과 같습니다.
보조정리 \ref{etale-cohomology-lemma-relative-colimit-general}.
$0 \in I$ 및 $K_0 \in D^+(X_{0, \etale})$로 설정한다.
$i \geq 0$에 대해 $K_i$는 $K_0$에서 $X_i$로의 철회를 나타냅니다.
$K$는 $K_0$에서 $X$로의 철회를 나타냅니다. 그 다음에
$$
R^pg_*K = \colim_{i \geq 0} h_i^{-1}R^pg_{i, *}K_i
$$
모든 $p \in \mathbf{Z}$에 대해.
\end{lemma}

\begin{proof}
정수 $p_0 \in \mathbf{Z}$를 수정한다.
허락하다$a$다음과 같은 정수가 되어야 한다.$H^j(K_0) = 0$~을 위한$j < a$.
우리는 모든 $p \leq p_0$에 대해 공식이 성립함을 증명하겠습니다.
$a$에서 유도를 내림차순으로 수행한다. 만약에
$a > p_0$, 그러면 왼쪽과 오른쪽이
소멸에 대한 공식은 $p \leq p_0$에 대해 0이다.
유도 범주, 보조정리 \ref{derived-lemma-negative-vanishing}.
$a \leq p_0$를 가정한다. 구별된 삼각형을 고려해보세요
$$
H^a(K_0)[-a] \to K_0 \to \tau_{\geq a + 1}K_0
$$
이 구별된 삼각형을 $X_i$ 및 $X$로 되돌립니다.
$K_i$ 및 $K$에 대해 호환 가능한 구별 삼각형을 제공한다.
$p \leq p_0$에 대해 가환 도식을 고려한다.
$$
\xymatrix{
\colim_{i \geq 0} h_i^{-1}R^{p - 1}g_{i, *}(\tau_{\geq a + 1}K_i)
\ar[r]_-\alpha \ar[d] &
R^{p - 1}g_*(\tau_{\geq a + 1}K) \ar[d] \\
\colim_{i \geq 0} h_i^{-1}R^pg_{i, *}(H^a(K_i)[-a])
\ar[r]_-\beta \ar[d] &
R^pg_*(H^a(K)[-a]) \ar[d] \\
\colim_{i \geq 0} h_i^{-1}R^pg_{i, *}K_i
\ar[r]_-\gamma \ar[d] &
R^pg_*K \ar[d] \\
\colim_{i \geq 0} R^pg_{i, *}\tau_{\geq a + 1}K_i
\ar[r]_-\delta \ar[d] &
R^pg_*\tau_{\geq a + 1}K \ar[d] \\
\colim_{i \geq 0} R^{p + 1}g_{i, *}(H^a(K_i)[-a])
\ar[r]^-\epsilon &
R^{p + 1}g_*(H^a(K)[-a])
}
$$
정확한 열이 있다. 화살표 $\beta$ 및 $\epsilon$는 동형사상이다.
보조정리 \ref{etale-cohomology-lemma-relative-colimit-general}.
화살표 $\alpha$ 및 $\delta$는 동형사상이다.
유도 가설에 의해.
따라서 $\gamma$는 원하는 대로 동형사상이다.
\end{proof}

\begin{lemma}
\label{etale-cohomology-lemma-relative-colimit-general-really-complexes}
$I$, $g_i : X_i \to S_i$, $g : X \to S$, $f_{ii'}$, $f_i$, $g_i$, $h_i$를
보조정리 \ref{etale-cohomology-lemma-relative-colimit-general}에서와 같이 두자.
$\mathcal{F}_i^\bullet$를 $X_{i, \etale}$ 위의 아벨 층의 복합체라 하자.
$\varphi_{i'i} : f_{i'i}^{-1}\mathcal{F}_i^\bullet \to
\mathcal{F}_{i'}^\bullet$를 $X_{i, \etale}$ 위의 복합체 사상이라 하고,
$\varphi_{i''i} = \varphi_{i''i'} \circ f_{i'' i'}^{-1}\varphi_{i'i}$가
$i'' \geq i' \geq i$일 때마다 성립한다고 하자. 어떤 정수 $a$가 존재하여
$\mathcal{F}_i^n = 0$이 $n < a$ 및 모든 $i \in I$에 대해 성립한다고 가정하자. 그러면
$$
R^pg_*(\colim f_i^{-1}\mathcal{F}_i^\bullet) =
\colim_{i \geq 0} h_i^{-1}R^pg_{i, *}\mathcal{F}_i^\bullet
$$
모든 $p \in \mathbf{Z}$에 대해 성립한다.
\end{lemma}

\begin{proof}
이는 보조정리 \ref{etale-cohomology-lemma-relative-colimit-general}의 결과이다. 집합
$\mathcal{F}^\bullet = \colim f_i^{-1}\mathcal{F}_i^\bullet$.
보조정리는 다음을 알려줍니다.
$$
\colim_{i \in I} h_i^{-1}R^pg_{i, *}\mathcal{F}_i^n = R^pg_*\mathcal{F}^n
$$
모든 $n, p \in \mathbf{Z}$에 대해. 스펙트럼 열을 사용해보자
$$
E_{1, i}^{s, t} = R^tg_{i, *}\mathcal{F}_i^s \Rightarrow
R^{s + t}g_{i, *}\mathcal{F}_i^\bullet
$$
그리고
$$
E_1^{s, t} = R^tg_*\mathcal{F}^s \Rightarrow
R^{s + t}g_*\mathcal{F}^\bullet
$$
유도 범주, 보조정리 \ref{derived-lemma-two-ss-complex-functor}.
부터$\mathcal{F}_i^n = 0$~을 위한$n < a$(와 함께$a$독립하다$i$)
우리는 고정된 유한 수의 용어 $E_{1, i}^{s, t}$만 있음을 알 수 있다.
($i$와 무관) 및 $E_1^{s, t}$가 기여하고
$E_1^{s, t} = \colim E_{i, i}^{s, t}$.
이는 우리가 원하는 것을 의미한다. 일부 세부 사항은 생략되었다.
(``어리석은'' 잘림을 사용하는 또 다른 주장이 있다.
스펙트럼 열 사용을 방지하는 복합체.)
\end{proof}

\begin{lemma}
\label{etale-cohomology-lemma-direct-sum-bounded-below-pushforward}
$f : X \to Y$를 준간소하고 준분리된 사상으로 가정한다.
스킴. $K_i \in D(X_\etale)$, $i \in I$를 객체 계열로 설정한다.
주어진 것으로 가정$a \in \mathbf{Z}$그렇게$H^n(K_i) = 0$~을 위한$n < a$
및 $i \in I$. 그런 다음 $Rf_*(\bigoplus_i K_i) = \bigoplus_i Rf_*K_i$.
\end{lemma}

\begin{proof}
$R^pf_*(\bigoplus_i K_i) =
\bigoplus_i R^pf_*K_i$가 모든 $p \in \mathbf{Z}$에 대해
성립함을 보이면 충분하다. 복합체 $\mathcal{F}_i^\bullet$로 $K_i$를 나타내되,
$\mathcal{F}_i^n = 0$이 $n < a$에 대해 성립하도록 선택하자.
단사 대상, 보조정리 \ref{injectives-lemma-derived-products}에 의해
복합체 $\mathcal{F}_i^\bullet$의 직합은 대상 $\bigoplus K_i$를 나타낸다.
$\bigoplus \mathcal{F}^\bullet$은 유한 직합들의 여과 여극한이므로,
결론은 보조정리 \ref{etale-cohomology-lemma-relative-colimit-general-really-complexes}에서 따른다.
\end{proof}














\section{줄기 의 고차 직상} \label{etale-cohomology-section-stalks-direct-image}

\noindent 고차 직상의 줄기는 종종 다음과 같이 계산될 수 있다.

\begin{theorem}
\label{etale-cohomology-theorem-higher-direct-images}
$f: X \to S$를 스킴의 준컴팩트하고 준분리 사상으로 두고,
$\mathcal{F}$ 아벨 층 $X_\etale$ 및 $\overline{s}$
기하점의 $S$이 $s \in S$ 위에 놓여 있다. 그 다음에
$$
\left(R^nf_* \mathcal{F}\right)_{\overline{s}} =
H_\etale^n( X \times_S \Spec(\mathcal{O}_{S, \overline{s}}^{sh}),
p^{-1}\mathcal{F})
$$
여기서 $p : X \times_S \Spec(\mathcal{O}_{S, \overline{s}}^{sh}) \to X$
투영이다. $K \in D^+(X_\etale)$ 및 $n \in \mathbf{Z}$의 경우
우리는
$$
\left(R^nf_*K\right)_{\overline{s}} =
H_\etale^n(X \times_S \Spec(\mathcal{O}_{S, \overline{s}}^{sh}), p^{-1}K)
$$
사실, 우리는
$$
\left(Rf_*K\right)_{\overline{s}}
=
R\Gamma_\etale(X \times_S \Spec(\mathcal{O}_{S, \overline{s}}^{sh}), p^{-1}K)
$$
$D^+(\textit{Ab})$에서.
\end{theorem}

\begin{proof}
$\mathcal{I}$를 $\overline{s}$의 에탈 동네의 범주라고 하자.
$S$에 있다. 보조정리 \ref{etale-cohomology-lemma-higher-direct-images}에 의해
우리는
$$
(R^nf_*\mathcal{F})_{\overline{s}} =
\colim_{(V, \overline{v}) \in \mathcal{I}^{opp}}
H_\etale^n(X \times_S V, \mathcal{F}|_{X \times_S V}).
$$
$\mathcal{I}$을 초기 하위 카테고리로 대체할 수 있다.
$\overline{s}$의 아핀 에탈 이웃. 그것을 관찰하십시오
$$
\Spec(\mathcal{O}_{S, \overline{s}}^{sh}) =
\lim_{(V, \overline{v}) \in \mathcal{I}} V
$$
보조정리 \ref{etale-cohomology-lemma-describe-etale-local-ring} 및
극한, 보조정리
\ref{limits-lemma-directed-inverse-system-affine-schemes-has-limit}.
올곱 극한으로 통근하므로 우리는 또한 다음을 얻습니다.
$$
X \times_S \Spec(\mathcal{O}_{S, \overline{s}}^{sh}) =
\lim_{(V, \overline{v}) \in \mathcal{I}} X \times_S V
$$
보조정리 \ref{etale-cohomology-lemma-directed-colimit-cohomology}로 결론을 내립니다.
두 번째 변형의 경우 다음을 사용하여 동일한 인수를 사용한다.
보조정리 \ref{etale-cohomology-lemma-colimit-complexes} 대신
보조정리 \ref{etale-cohomology-lemma-directed-colimit-cohomology}.

\medskip\noindent
마지막 진술이 참인지 확인하려면 지도를 생성하는 것으로 충분한다.
$\left(Rf_*K\right)_{\overline{s}} \to
R\Gamma_\etale(X \times_S \Spec(\mathcal{O}_{S, \overline{s}}^{sh}), p^{-1}K)$
코호몰로지에서 동형성을 구현하는 $D^+(\textit{Ab})$에서
학위 그룹$n$무엇보다도$n$. 이렇게 하려면 다음을 선택하세요.
분사형 아벨리안의 복합체 $\mathcal{J}^\bullet$ 아래로 경계됨
층~에$X_\etale$대표하는$K$. 그만큼복합체
$f_*\mathcal{J}^\bullet$는 $Rf_*K$을 나타냅니다. 따라서 복합체
$$
(f_*\mathcal{J}^\bullet)_{\overline{s}} =
\colim_{(V, \overline{v}) \in \mathcal{I}^{opp}}
(f_*\mathcal{J}^\bullet)(V)
$$
$(Rf_*K)_{\overline{s}}$을 나타냅니다. 각 $V$에 대해 지도가 있다.
$$
(f_*\mathcal{J}^\bullet)(V) =
\Gamma(X \times_S V, \mathcal{J}^\bullet)
\longrightarrow
\Gamma(X \times_S \Spec(\mathcal{O}_{S, \overline{s}}^{sh}),
p^{-1}\mathcal{J}^\bullet)
$$
대상 복합체는 다음을 나타냅니다.
$R\Gamma_\etale(X \times_S \Spec(\mathcal{O}_{S, \overline{s}}^{sh}), p^{-1}K)$
$D^+(\textit{Ab})$에서.
이 지도의 여극한을 사용하여 결과를 얻습니다.
\end{proof}

\begin{remark}
\label{etale-cohomology-remark-stalk-fibre}
$f : X \to S$를 스킴의 사상으로 설정한다.
$K \in D(X_\etale)$로 놔두세요.
$\overline{s}$를 $S$의 기하점으로 설정한다.
항상 표준 지도가 있다.
$$
(Rf_*K)_{\overline{s}}
\longrightarrow
R\Gamma(X \times_S \Spec(\mathcal{O}_{S, \overline{s}}^{sh}), p^{-1}K)
\longrightarrow
R\Gamma(X_{\overline{s}}, K|_{X_{\overline{s}}})
$$
여기서 $p : X \times_S \Spec(\mathcal{O}_{S, \overline{s}}^{sh}) \to X$
투영이다. 즉, 가환 도식을 고려하라.
$$
\xymatrix{
X_{\overline{s}} \ar[r] \ar[d]^{f_{\overline{s}}} &
X \times_S \Spec(\mathcal{O}_{S, \overline{s}}^{sh}) \ar[r]_-p \ar[d]^{f'} &
X \ar[d]^f \\
\overline{s} \ar[r]^-i &
\Spec(\mathcal{O}_{S, \overline{s}}^{sh}) \ar[r]^-j &
S
}
$$
기본 변경 맵이 있다.
$$
i^{-1}Rf'_*(p^{-1}K) \to Rf_{\overline{s}, *}(K|_{X_{\overline{s}}})
\quad\text{및}\quad
j^{-1}Rf_*K \to Rf'_*(p^{-1}K)
$$
(코호몰로지 사이트, 비고 \ref{sites-cohomology-remark-base-change})
이 다이어그램의 두 사각형에 대해.
대역 절단을 사용하여 원하는 지도를 얻습니다.
사이트, 비고에서 코호몰로지에 의해
\ref{sites-cohomology-remark-compose-base-change-horizontal}
이 두 맵의 구성은 다음과 같습니다.
평소(기본 변경) 지도
$(Rf_*K)_{\overline{s}} \to R\Gamma(X_{\overline{s}}, K|_{X_{\overline{s}}})$.
\end{remark}






\section{르레이 스펙트럼 열} \label{etale-cohomology-section-leray}

\begin{lemma} \label{etale-cohomology-lemma-prepare-leray} $f: X \to Y$를 사상이라 하고 $\mathcal{I}$를 $\textit{Ab}(X_\etale)$의 단사 대상으로 두자. $V \in \Ob(Y_\etale)$라 하자. 그러면 \begin{enumerate} \item 임의의 피복 $\mathcal{V} = \{V_j\to V\}_{j \in J}$에 대하여 $\check H^p(\mathcal{V}, f_*\mathcal{I}) = 0$이 모든 $p > 0$에 대해 성립한다. \item $f_*\mathcal{I}$는 함자 $\Gamma(V, -)$에 대하여 비순환적이다. \item $g : Y \to Z$이면 $f_*\mathcal{I}$는 $g_*$에 대하여 비순환적이다. \end{enumerate} \end{lemma}

\begin{proof}
$\check{\mathcal{C}}^\bullet(\mathcal{V}, f_*\mathcal{I}) =
\check{\mathcal{C}}^\bullet(\mathcal{V} \times_Y X, \mathcal{I})$
소멸 코호몰로지 군이 보조정리 \ref{etale-cohomology-lemma-hom-injective}만큼 높습니다.
이는 (1)를 증명한다. 두 번째 문은 다음과 같은 층으로 이어집니다.
소멸 모든 피복에 대한 상위 Čech 코호몰로지 군에는 소멸이 있다.
더 높은 코호몰로지 군. 이것은
Čech-to-코호몰로지 스펙트럼 열, 그러나 참조
코호몰로지 사이트, 보조정리 \ref{sites-cohomology-lemma-cech-vanish-collection}
자세한 내용과 보다 정확하고 일반적인 설명을 보려면
부분 (3)은 (2)의 결과이며 다음에 대한 설명이다.
$R^pg_*$~에보조정리 \ref{etale-cohomology-lemma-higher-direct-images}.
\end{proof}

\noindent Grothendieck의 형식주의 스펙트럼 열을 사용하면 다음과 같습니다.

\begin{proposition}[Leray 스펙트럼 열] \label{etale-cohomology-proposition-leray} $f: X \to Y$는 스킴의 사상이고 $\mathcal{F}$는 에탈 층이다. $X$. 그런 다음 스펙트럼 열 $$
E_2^{p, q} = H_\etale^p(Y, R^qf_*\mathcal{F}) \Rightarrow
H_\etale^{p+q}(X, \mathcal{F}).
$$ \end{proposition}가 있다.

\begin{proof} 보조정리 \ref{etale-cohomology-lemma-prepare-leray} 및 유도 범주, 절 \ref{derived-section-composition-right-derived-functors}를 참조하라. \end{proof}









\section{소멸 유한 고차 직상} \label{etale-cohomology-section-vanishing-finite-morphism}

\noindent 다음 목표는 스킴의 유한 사상의 고차 직상이 사라지는 것을 증명하는 것이다.

\begin{lemma}
\label{etale-cohomology-lemma-vanishing-etale-cohomology-strictly-henselian}
$R$는 엄밀히 말하면 헨셀식 국소환이라고 가정한다. 집합 $S = \Spec(R)$ 그리고
$\overline{s}$이 닫힌 지점이 된다. 그러면 글로벌
절 함자
$\Gamma(S, -) : \textit{Ab}(S_\etale) \to \textit{Ab}$
정확한다. 사실 우리는 $\Gamma(S, \mathcal{F}) = \mathcal{F}_{\overline{s}}$
집합 $\mathcal{F}$ 중 층에 대해. 특히
$$
\forall p\geq 1, \quad H_\etale^p(S, \mathcal{F})=0
$$
모든 $\mathcal{F}\in \textit{Ab}(S_\etale)$에 대해.
\end{lemma}

\begin{proof}
$\Gamma(S, \mathcal{F}) = \mathcal{F}_{\overline{s}}$를 보여주면
그러면 $\Gamma(S, -)$는 줄기 함자가 정확하듯이 정확한다.
$(U, \overline{u})$를 $\overline{s}$의 에탈 이웃이라고 하자.
아핀 오픈 동네를 선택하세요$\Spec(A)$~의$\overline{u}$~에$U$.
그러면 $R \to A$는 에탈이고 $\kappa(\overline{s}) = \kappa(\overline{u})$이다.
정리 \ref{etale-cohomology-theorem-henselian}에 의해 $A \cong R \times A'$를 알 수 있다.
지도와 호환되는 $R$-대수로
$\kappa(\overline{s}) = \kappa(\overline{u})$.
따라서 우리는 절을 얻습니다.
$$
\xymatrix{
\Spec(A) \ar[r] & U \ar[d]\\
& S \ar[ul]
}
$$
$\overline{s}$의 에탈 이웃 시스템에서는 다음과 같습니다.
신원 맵 $(S, \overline{s}) \to (S, \overline{s})$은 공동 최종이다.
그러므로 $\Gamma(S, \mathcal{F}) = \mathcal{F}_{\overline{s}}$.
보조정리의 마지막 문장은 더 높은 파생어와 같습니다.
완전 함자의 함자는 0이다. 참조
유도 범주, 보조정리 \ref{derived-lemma-right-derived-exact-functor}.
\end{proof}

\begin{proposition}
\label{etale-cohomology-proposition-finite-higher-direct-image-zero}
$f : X \to Y$를 스킴의 유한한 사상으로 둡니다.
\begin{enumerate}
\item 모든 기하점 $\overline{y} : \Spec(k) \to Y$에 대해 우리는
$$
(f_*\mathcal{F})_{\overline{y}} =
\prod\nolimits_{\overline{x} : \Spec(k) \to X,\ f(\overline{x}) =
\overline{y}} \mathcal{F}_{\overline{x}}.
$$
~을 위한$\mathcal{F}$~에$\Sh(X_\etale)$그리고
$$
(f_*\mathcal{F})_{\overline{y}} =
\bigoplus\nolimits_{\overline{x} : \Spec(k) \to X,\ f(\overline{x}) =
\overline{y}} \mathcal{F}_{\overline{x}}.
$$
~을 위한$\mathcal{F}$~에$\textit{Ab}(X_\etale)$.
\item 모든 $q \geq 1$에 대해 $R^q f_*\mathcal{F} = 0$가 있다.
~을 위한$\mathcal{F}$~에$\textit{Ab}(X_\etale)$.
\end{enumerate}
\end{proposition}

\begin{proof}
$X_{\overline{y}}^{sh}$는 올곱을 나타냅니다.
$X \times_Y \Spec(\mathcal{O}_{Y, \overline{y}}^{sh})$.
정리 \ref{etale-cohomology-theorem-higher-direct-images}에 의해
그만큼줄기~의$R^qf_*\mathcal{F}$~에$\overline{y}$에 의해 계산됩니다
$H_\etale^q(X_{\overline{y}}^{sh}, \mathcal{F})$.
$f$는 유한하므로 $X_{\bar y}^{sh}$는 유한한다.
$\Spec(\mathcal{O}_{Y, \overline{y}}^{sh})$, 따라서
$X_{\bar y}^{sh} = \Spec(A)$ 일부 환 $A$
$\mathcal{O}_{Y, \bar y}^{sh}$에 대해 유한한다.
후자는 엄격히 헨셀적이기 때문에,
보조정리 \ref{etale-cohomology-lemma-finite-over-henselian}
$A$는 헨셀식 국소환의 유한 곱임을 의미한다.
$A = A_1 \times \ldots \times A_r$. 잔여물 체 이후
$\mathcal{O}_{Y, \overline{y}}^{sh}$은 분리 가능하게 닫혀 있다.
각 $A_i$에 대해서도 마찬가지이다. 따라서 $A_i$은 엄밀히 말하면 헨셀식이다.
이는 $X_{\overline{y}}^{sh} = \coprod_{i = 1}^r \Spec(A_i)$을 의미한다.
소멸의
보조정리 \ref{etale-cohomology-lemma-vanishing-etale-cohomology-strictly-henselian}
다음을 암시한다$(R^qf_*\mathcal{F})_{\overline{y}} = 0$~을 위한$q > 0$
이는 (2) 에 의해정리 \ref{etale-cohomology-theorem-exactness-stalks}.
부분(1)은
보조정리 \ref{etale-cohomology-lemma-vanishing-etale-cohomology-strictly-henselian}.
\end{proof}

\begin{lemma}
\label{etale-cohomology-lemma-finite-pushforward-commutes-with-base-change}
데카르트 정사각형을 고려해보세요
$$
\xymatrix{
X' \ar[r]_{g'} \ar[d]_{f'} & X \ar[d]^f \\
Y' \ar[r]^g & Y
}
$$
스킴과 $f$의 유한한 사상.
누구에게나층~의집합 $\mathcal{F}$~에$X_\etale$우리는
$f'_*(g')^{-1}\mathcal{F} = g^{-1}f_*\mathcal{F}$.
\end{lemma}

\begin{proof}
일반적으로 당김 맵이 있다.
$g^{-1}f_*\mathcal{F} \to f'_*(g')^{-1}\mathcal{F}$, 참조
사이트, 절 \ref{sites-section-pullback}.
줄기 (정리 \ref{etale-cohomology-theorem-exactness-stalks})를 확인하면 충분한다.
$\overline{y}' : \Spec(k) \to Y'$를 기하점으로 설정한다.
우리는
\begin{align*}
(f'_*(g')^{-1}\mathcal{F})_{\overline{y}'}
& =
\prod\nolimits_{\overline{x}' : \Spec(k) \to X',\ f' \circ \overline{x}' =
\overline{y}'}
((g')^{-1}\mathcal{F})_{\overline{x}'} \\
& =
\prod\nolimits_{\overline{x}' : \Spec(k) \to X',\ f' \circ \overline{x}' =
\overline{y}'} \mathcal{F}_{g' \circ \overline{x}'} \\
& =
\prod\nolimits_{\overline{x} : \Spec(k) \to X,\ f \circ \overline{x} =
g \circ \overline{y}'} \mathcal{F}_{\overline{x}} \\
& =
(f_*\mathcal{F})_{g \circ \overline{y}'} \\
& =
(g^{-1}f_*\mathcal{F})_{\overline{y}'}
\end{align*}
최초의 평등
명제 \ref{etale-cohomology-proposition-finite-higher-direct-image-zero}.
두 번째 평등
보조정리 \ref{etale-cohomology-lemma-stalk-pullback}.
세 번째 등식은 다이어그램이 데카르트 정사각형이기 때문에 유지된다.
그래서 지도
$$
\{\overline{x}' : \Spec(k) \to X',\ f' \circ \overline{x}' =
\overline{y}'\}
\longrightarrow
\{\overline{x} : \Spec(k) \to X,\ f \circ \overline{x} =
g \circ \overline{y}'\}
$$
$\overline{x}'$을 $g' \circ \overline{x}'$로 보내는 것은 전단사이다.
네 번째 평등
명제 \ref{etale-cohomology-proposition-finite-higher-direct-image-zero}.
다섯 번째 평등
보조정리 \ref{etale-cohomology-lemma-stalk-pullback}.
\end{proof}

\begin{lemma}
\label{etale-cohomology-lemma-integral-pushforward-commutes-with-base-change}
데카르트 정사각형을 고려해보세요
$$
\xymatrix{
X' \ar[r]_{g'} \ar[d]_{f'} & X \ar[d]^f \\
Y' \ar[r]^g & Y
}
$$
스킴과 $f$ 및 정역 사상.
누구에게나층~의집합 $\mathcal{F}$~에$X_\etale$우리는
$f'_*(g')^{-1}\mathcal{F} = g^{-1}f_*\mathcal{F}$.
\end{lemma}

\begin{proof}
질문은 $Y$에 대한 로컬이므로 $Y$이 유사하다고 가정할 수 있다.
그럼 우리는 쓸 수 있습니다$X = \lim X_i$~와 함께$f_i : X_i \to Y$한정된
(아핀 경우에는 쉽지만,
극한, 보조정리 \ref{limits-lemma-integral-limit-finite-and-finite-presentation}
참고용). $p_{i'i} : X_{i'} \to X_i$ 표시
전환 사상 및 $p_i : X \to X_i$ 투영.
$\mathcal{F}_i = p_{i, *}\mathcal{F}$ 설정을 얻습니다.
보조정리 \ref{etale-cohomology-lemma-linus-hamann}에서
시스템 $(\mathcal{F}_i, \varphi_{i'i})$
$\mathcal{F} = \colim p_i^{-1}\mathcal{F}_i$로.
우리는 $f_*\mathcal{F} = \colim f_{i, *}\mathcal{F}_i$를 얻습니다.
보조정리 \ref{etale-cohomology-lemma-relative-colimit}에서.
집합 $X'_i = Y' \times_Y X_i$ 투영 $f'_i$ 및 $g'_i$가 포함된다.
그런 다음 $X' = \lim X'_i$를 극한으로 극한으로 통근한다.
$p'_i : X' \to X'_i$ 투영을 나타냅니다. 우리는
\begin{align*}
g^{-1}f_*\mathcal{F}
& =
g^{-1} \colim f_{i, *}\mathcal{F}_i \\
& =
\colim g^{-1}f_{i, *}\mathcal{F}_i \\
& =
\colim f'_{i, *}(g'_i)^{-1}\mathcal{F}_i \\
& =
f'_*(\colim (p'_i)^{-1}(g'_i)^{-1}\mathcal{F}_i) \\
& =
f'_*(\colim (g')^{-1}p_i^{-1}\mathcal{F}_i) \\
& =
f'_*(g')^{-1} \colim p_i^{-1}\mathcal{F}_i \\
& =
f'_*(g')^{-1}\mathcal{F}
\end{align*}
원하는대로. 첫 번째 평등에 대해서는 위를 참조하라.
두 번째 사용의 경우 당김은 여극한으로 통근한다.
유한 사례를 세 번째로 사용하려면 다음을 참조하라.
보조정리 \ref{etale-cohomology-lemma-finite-pushforward-commutes-with-base-change}.
네 번째에는 보조정리 \ref{etale-cohomology-lemma-relative-colimit}를 사용한다.
다섯 번째로는 $g'_i \circ p'_i = p_i \circ g'$을 사용하세요.
여섯 번째 사용의 경우 당김은 여극한으로 통근한다.
일곱 번째에는 $\mathcal{F} = \colim p_i^{-1}\mathcal{F}_i$을 사용한다.
\end{proof}

\noindent 다음 보조정리는 에탈 층과 유한 전사 사상을 다루는 동종계의 경우이다. 고유 기저변환 정리를 증명하면 이 결과를 상당히 일반화할 것이다.

\begin{lemma}
\label{etale-cohomology-lemma-cohomological-descent-finite}
$f : X \to Y$를 스킴의 전사 유한 사상이라고 하자.
집합 $f_n : X_n \to Y$는 $(n + 1)$배 올곱과 같습니다.
~의$X$~ 위에$Y$. 을 위한$\mathcal{F} \in \textit{Ab}(Y_\etale)$ 집합
$\mathcal{F}_n = f_{n, *}f_n^{-1}\mathcal{F}$. 정확한 내용이 있습니다
순서
$$
0 \to \mathcal{F} \to \mathcal{F}_0 \to \mathcal{F}_1 \to
\mathcal{F}_2 \to \ldots
$$
$Y_\etale$에 있다. 게다가 스펙트럼 열
$$
E_1^{p, q} = H^q_\etale(X_p, f_p^{-1}\mathcal{F})
$$
$H^{p + q}(Y_\etale, \mathcal{F})$로 수렴된다.
이 스펙트럼 열은 $\mathcal{F}$에서 기능적이다.
\end{lemma}

\begin{proof}
보조정리의 첫 번째 진술을 증명하면 스펙트럼을 얻습니다.
$E_1^{p, q} = H^q_\etale(Y, \mathcal{F}_p)$ 수렴하는 시퀀스
$H^{p + q}(Y_\etale, \mathcal{F})$로, 참조
유도 범주, 보조정리 \ref{derived-lemma-two-ss-complex-functor}.
반면에, 이후
$R^if_{p, *}f_p^{-1}\mathcal{F} = 0$ 에 대하여 $i > 0$
(명제 \ref{etale-cohomology-proposition-finite-higher-direct-image-zero})
우리는 얻는다
$$
H^q_\etale(X_p, f_p^{-1}\mathcal{F}) =
H^q_\etale(Y, f_{p, *}f_p^{-1} \mathcal{F}) =
H^q_\etale(Y, \mathcal{F}_p)
$$
명제 \ref{etale-cohomology-proposition-leray} 기준
보조정리의 스펙트럼 열을 얻습니다.

\medskip\noindent
보조정리의 첫 번째 진술을 증명하려면 다음을 관찰하라.
$X_n$는 $Y$에 대해 단순 스킴을 형성한다.
단순, 예 \ref{simplicial-example-fibre-products-simplicial-object}.
또한 각 예측에 대해 다음을 관찰하라.
$d_j : X_{n + 1} \to X_n$ 지도가 있습니다
$d_j^{-1} f_n^{-1}\mathcal{F} \to f_{n + 1}^{-1}\mathcal{F}$.
이 지도는 지도를 유도한다.
$$
\delta_j : \mathcal{F}_n \to \mathcal{F}_{n + 1}
$$
$j = 0, \ldots, n + 1$의 경우. 우리는 이 지도들의 교대합을 사용한다.
미분 $\mathcal{F}_n \to \mathcal{F}_{n + 1}$을 정의한다.
마찬가지로, 정식 확장 $\mathcal{F} \to \mathcal{F}_0$이 있다.
즉 이것은 단지 표준 맵 $\mathcal{F} \to f_*f^{-1}\mathcal{F}$이다.
이 층 시퀀스가 ​​정확한 복합체인지 확인하려면 충분한다.
기하점 (정리 \ref{etale-cohomology-theorem-exactness-stalks})에서 줄기를 확인한다.
따라서 $\overline{y} : \Spec(k) \to Y$를 기하점으로 둡니다. 허락하다
$E = \{\overline{x} : \Spec(k) \to X \mid f(\overline{x}) = \overline{y}\}$.
그러면 $E$는 비어 있지 않은 유한한 집합이며 다음을 알 수 있다.
$$
(\mathcal{F}_n)_{\overline{y}} =
\bigoplus\nolimits_{e \in E^{n + 1}} \mathcal{F}_{\overline{y}}
$$
명제 \ref{etale-cohomology-proposition-finite-higher-direct-image-zero} 기준
및 보조정리 \ref{etale-cohomology-lemma-stalk-pullback}.
따라서 우리는 아벨 군 $M$ 시퀀스가 ​​주어지는 것을 확인해야 한다.
$$
0 \to M \to \bigoplus\nolimits_{e \in E} M \to
\bigoplus\nolimits_{e \in E^2} M \to \ldots
$$
정확한다. 여기서 첫 번째 지도는 대각선 지도이고 지도는
$\bigoplus_{e \in E^{n + 1}} M  \to \bigoplus_{e \in E^{n + 2}} M$
는 $(n + 2)$에 의해 유도된 맵의 교번 합이다.
투영 $E^{n + 2} \to E^{n + 1}$. 이건 바로 보여드릴 수 있어요
또는 단순, 보조정리를 적용하여 추론한다.
\ref{simplicial-lemma-fibre-products-simplicial-object-w-section}
지도 $E \to \{*\}$로 이동한다.
\end{proof}

\begin{remark} \label{etale-cohomology-remark-cohomological-descent-finite} 보조정리 \ref{etale-cohomology-lemma-cohomological-descent-finite}의 상황에서 $\mathcal{G}$가 집합의 층인 경우 $Y_\etale$에 대해 $$
\Gamma(Y, \mathcal{G}) =
\text{평형 장치}(
\xymatrix{
\Gamma(X_0, f_0^{-1}\mathcal{G})
\ar@<1ex>[r] \ar@<-1ex>[r] &
\Gamma(X_1, f_1^{-1}\mathcal{G})
}
)
$$ 이는 층 $\mathcal{G}$가 두 맵 $f_{0, *}f_0^{-1}\mathcal{G} \to f_{1, *}f_1^{-1}\mathcal{G}$의 이퀄라이저임을 보여줌으로써 정확히 동일한 방식으로 증명된다. \end{remark}









%10.06.09
\section{줄기} \label{etale-cohomology-section-galois-action-stalks}에 대한 갈루아 조치

\noindent
이 절에서 우리는 잔기의 절대 갈루아 그룹의 작용을 정의합니다
체한 점의$s$~의$S$에줄기 함자언제든지기하점거짓말하는
$s$ 이상.

\medskip\noindent
줄기에 대한 갈루아의 행동.
$S$를 스킴으로 설정한다.
$\overline{s}$를 $S$의 기하점으로 설정한다.
$\sigma \in \text{Aut}(\kappa(\overline{s})/\kappa(s))$로 놔두세요.
다음 작업을 정의한다.$\sigma$에줄기 $\mathcal{F}_{\overline{s}}$
층 $\mathcal{F}$의 다음과 같습니다.
\begin{equation}
\label{etale-cohomology-equation-galois-action}
\begin{matrix}
\mathcal{F}_{\overline{s}} &
\longrightarrow &
\mathcal{F}_{\overline{s}} \\
(U, \overline{u}, t) &
\longmapsto &
(U, \overline{u} \circ \Spec(\sigma), t).
\end{matrix}
\end{equation}
여기서는 줄기 요소에 대한 설명을 트리플 측면에서 사용한다.
다음 논의에서처럼
정의 \ref{etale-cohomology-definition-stalk}.
이것은 왼쪽 동작이다.
$\sigma_i \in \text{Aut}(\kappa(\overline{s})/\kappa(s))$
그 다음에
\begin{align*}
\sigma_1 \cdot (\sigma_2 \cdot (U, \overline{u}, t))
& =
\sigma_1 \cdot (U, \overline{u} \circ \Spec(\sigma_2), t) \\
& =
(U, \overline{u} \circ \Spec(\sigma_2) \circ \Spec(\sigma_1), t) \\
& =
(U, \overline{u} \circ \Spec(\sigma_1 \circ \sigma_2), t) \\
& =
(\sigma_1 \circ \sigma_2) \cdot (U, \overline{u}, t)
\end{align*}
이 작업은 층 $\mathcal{F}$에서 기능적이라는 것이 분명한다.
우리는 다음을 직접 참조하여 이 작업을 정의할 수 있다는 점에 주목한다.
비고 \ref{etale-cohomology-remark-map-stalks}.

\begin{definition}
\label{etale-cohomology-definition-algebraic-geometric-point}
$S$를 스킴으로 설정한다.
허락하다$\overline{s}$가 되다기하점요점 위에 누워있다$s$~의$S$.
$\kappa(s) \subset \kappa(s)^{sep} \subset \kappa(\overline{s})$
대수적으로 $\kappa(s)$의 분리 가능한 대수적 종결을 나타냅니다.
닫혔습니다 체 $\kappa(\overline{s})$.
\begin{enumerate}
\item 이 상황에서 {\it 절대 갈루아 그룹} $\kappa(s)$
$\text{Gal}(\kappa(s)^{sep}/\kappa(s))$이다. 표기하는 경우도 있습니다
$\text{Gal}_{\kappa(s)}$.
\item 기하점 $\overline{s}$를 호출한다.
{\it 대수} 만약 $\kappa(s) \subset \kappa(\overline{s})$
$\kappa(s)$의 대수적 종결.
\end{enumerate}
\end{definition}

\begin{example} \label{etale-cohomology-example-stupid} 기하점 $\Spec(\mathbf{C}) \to \Spec(\mathbf{Q})$는 대수가 아닙니다. \end{example}

\noindent
$\kappa(s) \subset \kappa(s)^{sep} \subset \kappa(\overline{s})$
정의와 같습니다. $\kappa(\overline{s})$은 대수적으로
지도를 닫았어
$$
\text{Aut}(\kappa(\overline{s})/\kappa(s))
\longrightarrow
\text{Gal}(\kappa(s)^{sep}/\kappa(s)) = \text{Gal}_{\kappa(s)}
$$
전사적이다. $(U, \overline{u})$가 다음과 같다고 가정한다.
에탈 동네 $\overline{s}$, 그리고 $\overline{u}$이 놓여 있다고 말함
요점$u$~의$U$. 부터$U \to S$는 에탈, 잔여물체확대
$\kappa(u)/\kappa(s)$은 유한 분리 가능한다.
이는 다음을 의미한다.
\begin{enumerate}
\item 만약 $\sigma \in \text{Aut}(\kappa(\overline{s})/\kappa(s)^{sep})$
그러면 $\sigma$는 $\mathcal{F}_{\overline{s}}$에 대해 사소한 역할을 한다.
\item 보다 정확하게는
$\text{Aut}(\kappa(\overline{s})/\kappa(s))$
절대 갈루아 그룹의 행동에 의해 결정되고 결정됩니다
$\text{Gal}_{\kappa(s)}$ - $\mathcal{F}_{\overline{s}}$.
\item $(U, \overline{u}, t)$가 원소 $\xi$를 나타낼 때,
$\mathcal{F}_{\overline{s}}$ 모든 요소
$\text{Gal}(\kappa(s)^{sep}/K)$는 사소하게 작동한다.
$\kappa(s) \subset K \subset \kappa(s)^{sep}$는 다음의 이미지이다.
$\overline{u}^\sharp : \kappa(u) \to \kappa(\overline{s})$.
\end{enumerate}
전체적으로 우리는 $\mathcal{F}_{\overline{s}}$가
$\text{Gal}_{\kappa(s)}$-집합 (참조
기본 그룹, 정의 \ref{pione-definition-G-set-continuous}).
따라서 줄기 함자를 함자로 생각할 수 있다.
$$
\Sh(S_\etale) \longrightarrow
\text{Gal}_{\kappa(s)}\textit{-Sets},
\quad
\mathcal{F} \longmapsto \mathcal{F}_{\overline{s}}
$$
이제부터 우리는 줄기 함자에 대해 보통 이런 식으로 생각하게 된다.

\begin{theorem} \label{etale-cohomology-theorem-equivalence-sheaves-point} $S = \Spec(K)$를 $K$와 체로 둡니다. $\overline{s}$를 $S$의 기하점으로 설정한다. $G = \text{Gal}_{\kappa(s)}$는 절대 갈루아 그룹을 나타냅니다. 줄기를 복용하면 범주의 동치 $$
\Sh(S_\etale) \longrightarrow G\textit{-Sets},
\quad
\mathcal{F} \longmapsto \mathcal{F}_{\overline{s}}.
$$ \end{theorem}가 유발된다.

\begin{proof}
이 함자의 역함수를 구성해 보자. ~ 안에
기본 그룹, 보조정리 \ref{pione-lemma-sheaves-point}
우리는 $G$-집합 $M$가 주어지면 에탈 사상이 존재한다는 것을 보았습니다.
$X \to \Spec(K)$
$\Mor_K(\Spec(K^{sep}), X)$는 다음과 같습니다.
$M$와 $G$-집합과 동형이다. 층을 고려해보세요.
$\mathcal{F}$에 의해 정의된 $\Spec(K)_\etale$의
규칙 $U \mapsto \Mor_K(U, X)$. 에탈은 층 이다.
위상은 하위 표준이다. 그러면 우리는 그것을 본다
$\mathcal{F}_{\overline{s}} = \Mor_K(\Spec(K^{sep}), X) = M$
$G$-집합(자세한 내용은 생략). 이는 함자의 역을 제공하고
우리가 승리한다.
\end{proof}

\begin{remark}
\label{etale-cohomology-remark-every-sheaf-representable}
결론을 표현하는 또 다른 방법
정리 \ref{etale-cohomology-theorem-equivalence-sheaves-point} 및
기본 그룹, 보조정리 \ref{pione-lemma-sheaves-point}
즉, $\Spec(K)_\etale$의 모든 층을 표현할 수 있다는 뜻이다.
에 의해 스킴 $X$ 에탈 위의 $\Spec(K)$.
이는 모든 층이 다음과 같은 의미로 표현 가능하다는 의미는 아닙니다.
사이트, 정의 \ref{sites-definition-representable-sheaf}.
그 이유는 $\Spec(K)_\etale$의 구성에서
우리는 $\Spec(K)$ 대신 스킴의 에탈 중 충분히 큰 집합을 선택했다.
$\Spec(K)_\etale$의 층은 적절한 클래스를 형성한다.
\end{remark}

\begin{lemma} \label{etale-cohomology-lemma-global-sections-point} 정리 \ref{etale-cohomology-theorem-equivalence-sheaves-point}와 같은 가정 및 표기법이다. 기능적 전단사가 있습니다 $$
\Gamma(S, \mathcal{F}) = (\mathcal{F}_{\overline{s}})^G
$$ \end{lemma}

\begin{proof}
우리는 공식적인 논증과 다음의 결과를 사용하여 이를 증명할 수 있다.
정리 \ref{etale-cohomology-theorem-equivalence-sheaves-point}
다음과 같이. 층 $\mathcal{F}$에 해당하는 경우
$G$-집합 $M = \mathcal{F}_{\overline{s}}$ 우리는
\begin{eqnarray*}
\Gamma(S, \mathcal{F}) & = &
\Mor_{\Sh(S_\etale)}(h_{\Spec(K)},  \mathcal{F})
\\
& = & \Mor_{G\textit{-Sets}}(\{*\}, M) \\
& = & M^G
\end{eqnarray*}
여기서 첫 번째 식별에 대해 설명한다.
사이트, 절 \ref{sites-section-presheaves} 및
\ref{sites-section-representable-sheaves},
두 번째 결과
정리 \ref{etale-cohomology-theorem-equivalence-sheaves-point}
세 번째는 분명한다. 또한 증명\footnote{에 대해 직접 제공할 것이다.
저 밖에는 의심하는 토마스 가족이 있다.}.

\medskip\noindent
$t \in \Gamma(S, \mathcal{F})$이 대역 절단라고 가정한다.
그런 다음 트리플 $(S, \overline{s}, t)$은 다음 요소를 정의한다.
$\mathcal{F}_{\overline{s}}$ 이는 다음과 같이 명확하게 변하지 않습니다.
$G$의 동작. 반대로 $(U, \overline{u}, t)$라고 가정해 보자.
불변인 $\mathcal{F}_{\overline{s}}$의 요소를 정의한다.
그런 다음 $U$를 축소하고 일부에 대해 $U = \Spec(L)$를 가정할 수 있다.
유한 분리 가능 체 $K$의 확장, 참조
명제 \ref{etale-cohomology-proposition-etale-morphisms}.
이 경우 지도는 $\mathcal{F}(U) \to \mathcal{F}_{\overline{s}}$
에탈 지역의 사상에 대해 단사형이다.
$(U', \overline{u}') \to (U, \overline{u})$ 제한 사상
$\mathcal{F}(U) \to \mathcal{F}(U')$는 $U' \to U$이므로 분사형이다.
$S_\etale$의 피복이다.
$L$를 조금 확대한 후 $K \subset L$가 유한하다고 가정할 수 있다.
갈루아 확장. 이 시점에서 우리는 그것을 사용합니다
$$
\Spec(L) \times_{\Spec(K)} \Spec(L)
=
\coprod\nolimits_{\sigma \in \text{Gal}(L/K)} \Spec(L)
$$
지도는 $\Spec(L) \to \Spec(L \otimes_K L)$
환 지도 $a \otimes b \mapsto a\sigma(b)$에서 왔습니다. 그러므로 우리는
$(U, \overline{u}, t)$이 변하지 않는 조건을 확인하세요.
$G$ 아래의 모든 것은 $t \in \mathcal{F}(\Spec(L))$을 의미한다.
의 동일한 요소에 매핑된다.
$\mathcal{F}(\Spec(L) \times_{\Spec(K)} \Spec(L))$
투영에 의한 제한을 통해(이것은 언급된 주입성을 사용한다.
위에; 자세한 내용은 생략) 따라서 $\mathcal{F}$의 층 조건
에탈 피복 $\{\Spec(L) \to \Spec(K)\}$ 킥을 위해
에서 $t$는 $\mathcal{F}$의 고유한 절에서 온다고 결론을 내립니다.
$\Spec(K)$ 이상.
\end{proof}

\begin{remark} \label{etale-cohomology-remark-stalk-pullback} $S$는 스킴이고 $\overline{s} : \Spec(k) \to S$는 $S$의 기하점이다. 정의는 $k$가 대수적으로 닫혀 있음을 의미한다. 특히 $k$의 절대 갈루아 그룹은 사소한다. 따라서 정리 \ref{etale-cohomology-theorem-equivalence-sheaves-point}에 의해 $\Spec(k)_\etale$의 층의 범주는 집합의 범주와 동일한다. $\Spec(k)$에 절을 더하면 동등성이 부여된다. 이는 마침내 줄기 함자의 대안인 정의를 제공한다. 즉, 함자 $$
\Sh(S_\etale) \longrightarrow \textit{Sets}, \quad
\mathcal{F} \longmapsto \mathcal{F}_{\overline{s}}
$$는 함자 $$
\Sh(S_\etale)
\longrightarrow
\Sh(\Spec(k)_\etale) = \textit{Sets},
\quad
\mathcal{F} \longmapsto \overline{s}^*\mathcal{F}
$$와 동형이다. 이를 엄격하게 증명하려면 보조정리 \ref{etale-cohomology-lemma-stalk-pullback} 부분(3)을 다음과 같이 사용할 수 있다. $f = \overline{s}$. 게다가 이렇게 말하면 보조정리 \ref{etale-cohomology-lemma-stalk-pullback} 부분(3)의 일반적인 경우는 하락 기능에 따른 것이다. \end{remark}






\section{군 코호몰로지}
\label{etale-cohomology-section-group-cohomology}

\noindent
다음에서 $H^i(G, M)$라고 쓰면 $G$를 의미한다.
은 위상 그룹이고 $M$ 이산 $G$-가군
연속적인 $G$-동작이고 $H^i(G, -)$는 $i$번째 오른쪽 파생이다.
함자 범주 $\text{Mod}_G$ $G$-가군에 대한 내용은 다음을 참조하라.
정의 \ref{etale-cohomology-definition-G-module-continuous} 및
\ref{etale-cohomology-definition-galois-cohomology}. 여기에는 다음이 포함됩니다
추상 그룹 $G$의 경우, 이는 단순히 $G$이 표시됨을 의미한다.
이산 위상을 갖는 위상 그룹으로.

\medskip\noindent 가군에 불연속적인 위상이 있는 경우 $H^i_{cont}(G, M)$ 표기법을 사용하여 \cite{Tate}에 도입된 연속적인 코호몰로지 군을 나타냅니다. 절 \ref{etale-cohomology-section-continuous-group-cohomology}.

\begin{definition}
\label{etale-cohomology-definition-G-module-continuous}
$G$를 위상 그룹으로 가정한다.
\begin{enumerate}
\item {\it $G$-가군}, 때로는 {\it 이산이라고도 함 $G$-가군},
은아벨 군 $M$왼쪽 행동이 부여됨$a : G \times M \to M$
그룹 동형에 의해 다음과 같이$a$는 연속적일 때$M$주어진다
이산 위상.
\item {\it 사상 의 $G$-가군} $f : M \to N$는
$G$-$M$에서 $N$까지의 등가 동형이다.
\item $G$-가군의 범주는 $\text{Mod}_G$로 표시된다.
\end{enumerate}
$R$를 환으로 설정한다.
\begin{enumerate}
\item {\it $R\text{-}G$-가군}는 $R$-가군 $M$이다.
좌익 행동$a : G \times M \to M$~에 의해$R$- 선형 맵은 다음과 같습니다.$a$
$M$에 이산 위상이 주어지면 연속이다.
\item {\it 사상 의 $R\text{-}G$-가군} $f : M \to N$는
$G$-등변 $R$-가군은 $M$에서 $N$로 매핑된다.
\item $R\text{-}G$-가군의 범주는 $\text{Mod}_{R, G}$로 표시된다.
\end{enumerate}
\end{definition}

\noindent
$a : G \times M \to M$ 연속인 조건은 동일한다.
조건을 사용하면 $x \in M$의 안정 장치가 $G$에 열려 있다.
$G$가 추상 그룹인 경우 이는 다음의 개념에 해당한다.
아벨 군에는 $G$-액션이 부여되어 $G$에 다음이 부여된다.
이산 위상. $\text{Mod}_{\mathbf{Z}, G} = \text{Mod}_G$을 확인하세요.

\medskip\noindent
범주 $\text{Mod}_G$에는 충분한 분사가 있다.
분사, 보조정리 \ref{injectives-lemma-G-modules}.
왼쪽 완전 함자를 고려하라.
$$
\text{Mod}_G \longrightarrow \textit{Ab},
\quad
M \longmapsto M^G =
\{x \in M \mid g \cdot x = x\ \forall g \in G\}
$$
$M^G = H^0(G, M)$를 표시할 때도 있고, 다음과 같이 쓰기도 한다.
$M^G = \Gamma_G(M)$. 이 함자는 완전한 권리 유도 함자를 갖고 있다.
$R\Gamma_G(M)$ 및 $i$번째 오른쪽 유도 함자
$R^i\Gamma_G(M) = H^i(G, M)$ 모든 $i \geq 0$에 대해.

\medskip\noindent $H^0(G, -) : \text{Mod}_{R, G} \to \text{Mod}_R$에도 동일한 구성이 적용된다. 보조정리 \ref{etale-cohomology-lemma-modules-abelian}에서 이것이 기본 $G$-가군의 코호몰로지와 일치함을 볼 수 있다.

\begin{definition}
\label{etale-cohomology-definition-galois-cohomology}
$G$를 위상 그룹으로 가정한다. $M$를 이산 $G$-가군으로 가정한다.
지속적인 $G$ 동작으로. 즉, $M$는 객체이다.
범주 $\text{Mod}_G$에 소개되었다.
정의 \ref{etale-cohomology-definition-G-module-continuous}.
\begin{enumerate}
\item $H^i(G, M)$, 즉 $H^0(G, M)$의 오른쪽 유도함자를
범주 $\text{Mod}_G$는
{\it 연속 군 코호몰로지 군} $M$.
\item $G$가 이산 위상을 부여받은 추상 그룹인 경우
그런 다음$H^i(G, M)$이라고 불린다{\it 군 코호몰로지여러 떼}~의$M$.
\item $G$가 갈루아 그룹인 경우 $H^i(G, M)$ 그룹이 호출된다.
{\it Galois 코호몰로지 군} $M$.
\item $G$가 체 $K$의 절대 갈루아 그룹인 경우 그룹은
$H^i(G, M)$는 {\it Galois 코호몰로지 $K$ 그룹이라고도 한다.
$M$}에서 계수를 사용한다. 이런 경우 가끔 글을 씁니다.
$H^i(K, M)$대신에$H^i(G, M)$.
\end{enumerate}
\end{definition}

\begin{lemma} \label{etale-cohomology-lemma-modules-abelian} $G$를 위상 그룹으로 둡니다. $R$를 환으로 설정한다. 모든 $i \geq 0$에 대해 수직 화살표가 망각된 함자인 다이어그램 $$
\xymatrix{
\text{Mod}_{R, G} \ar[rr]_{H^i(G, -)} \ar[d] & &
\text{Mod}_R \ar[d] \\
\text{Mod}_G \ar[rr]^{H^i(G, -)} & &
\textit{Ab}
}
$$는 교환 가능한다. \end{lemma}

\begin{proof}
건망증이 있는 함자 $F : \text{Mod}_{R, G} \to \text{Mod}_G$을 표시해 보자.
그런 다음 $F$에는 왼쪽 수반 $H : \text{Mod}_G \to \text{Mod}_{R, G}$가 있다.
$H(M) = M \otimes_\mathbf{Z} R$에 의해 주어진다. 모든 대상이
$\text{Mod}_G$는 가군 형식의 직합의 몫이다.
$\mathbf{Z}[G/U]$ 여기서 $U \subset G$는 열린 하위 그룹이다.
여기서 $\mathbf{Z}[G/U]$는 $G$-가군을 의미한다.
유한 $\mathbf{Z}$-선형 조합
권리의$U$일치 클래스$G$왼쪽으로 부여$G$-행동.
따라서 $\text{Mod}_G$에서 복합체 위의 모든 경계는 준동형이다.
$\text{Mod}_G$의 복합체 위에 있는 경계로
용어는 단순한다. $\mathbf{Z}$-가군
(유도 범주, 보조정리 \ref{derived-lemma-subcategory-left-resolution}).
따라서 $LH$이 $D^-(\text{Mod}_G)$에 존재하고 다음과 같이 계산된다는 것이 분명한다.
평가하다$H$무엇이든복합체조건이 균일하지 않은 사람$\mathbf{Z}$-가군;
이것은 다음과 같습니다
유도 범주, 보조정리 \ref{derived-lemma-subcategory-left-acyclics} 및
명제 \ref{derived-proposition-enough-acyclics}.
유도 범주, 보조정리로 결론을 내립니다.
\ref{derived-lemma-pre-derived-adjoint-functors}
저것
$$
\text{Ext}^i(\mathbf{Z}, F(M)) = \text{Ext}^i(R, M)
$$
~을 위한$M$~에$\text{Mod}_{R, G}$.
$H^0(G, -) = \Hom(\mathbf{Z}, -)$가 켜져 있는지 확인하세요.
$\text{Mod}_G$ 여기서 $\mathbf{Z}$는 $G$-가군을 나타냅니다.
사소한 행동으로. 따라서
$H^i(G, -) = \text{Ext}^i(\mathbf{Z}, -)$ - $\text{Mod}_G$.
마찬가지로 $H^i(G, -) = \text{Ext}^i(R, -)$가 켜져 있다.
$\text{Mod}_{R, G}$. 우리가 보는 모든 것을 결합하면 보조정리가 참이라는 것을 알 수 있다.
\end{proof}

\begin{lemma}
\label{etale-cohomology-lemma-ext-modules-hom}
$G$를 위상 그룹으로 가정한다. $R$를 환으로 설정한다.
$M$, $N$를 $R\text{-}G$-가군으로 설정한다. $M$가 유한 투영인 경우
$R$-가군으로, 그런 다음
$\text{Ext}^i(M, N) = H^i(G, M^\vee \otimes_R N)$ (기호용
증명)를 참조하라.
\end{lemma}

\begin{proof}
가군 $M^\vee = \Hom_R(M, R)$는 반대의 속성을 부여받았습니다.
$G$의 동작. 즉, $(g \cdot \lambda)(m) = \lambda(g^{-1} \cdot m)$
$g \in G$, $\lambda \in M^\vee$, $m \in M$의 경우. $G$의 동작
$M^\vee \otimes_R N$는 대각선이다. 즉, 다음과 같이 지정된다.
$g \cdot (\lambda \otimes n) = g \cdot \lambda \otimes g \cdot n$.
세 번째 $R\text{-}G$-가군 $E$에 대해서는
$\Hom(E, M^\vee \otimes_R N) = \Hom(M \otimes_R E, N)$.
즉, $R$-가군 수준에서는 다음과 같습니다.
대수학, 기본형 \ref{algebra-lemma-hom-from-tensor-product} 및
\ref{algebra-lemma-evaluation-map-iso-finite-projective}
$G$-액션의 정의는 다음과 같이 선택된다.
$R\text{-}G$-가군에 대해서는 참로 유지된다. 그것은 다음과 같습니다
$M^\vee \otimes_R N$는 단사 $R\text{-}G$-가군이다.
$N$가 단사 $R\text{-}G$-가군인 경우. 그러므로 만일
$N \to N^\bullet$는 단사 분해이다.
$M^\vee \otimes_R N \to M^\vee \otimes_R N^\bullet$
단사 분해이다. 그 다음에
$$
\Hom(M, N^\bullet) = \Hom(R, M^\vee \otimes_R N^\bullet) =
(M^\vee \otimes_R N^\bullet)^G
$$
왼쪽은 $\text{Ext}^i(M, N)$를 계산하고 오른쪽은
손 측은 $H^i(G, M^\vee \otimes_R N)$를 계산하며 증명이 완료된다.
\end{proof}

\begin{lemma} \label{etale-cohomology-lemma-finite-dim-group-cohomology} $G$를 위상 그룹으로 둡니다. $k$를 체로 설정한다. $V$를 $k\text{-}G$-가군으로 설정한다. $G$가 위상적으로 유한하게 생성되고 $\dim_k(V) < \infty$이면 $\dim_k H^1(G, V) < \infty$이다. \end{lemma}

\begin{proof}
$g_1, \ldots, g_r \in G$를 위상적으로 $G$를 생성하는 요소라고 하면,
즉, 이는 $g_1, \ldots, g_r$에 의해 생성된 하위 그룹이 조밀임을 의미한다.
보조정리 \ref{etale-cohomology-lemma-ext-modules-hom}에 의해
$H^1(G, V)$는 확장의 $k$-벡터공간이다.
$$
0 \to V \to E \to k \to 0
$$
$k\text{-}G$-가군. $e \in E$를 $1 \in k$에 매핑하는 것을 선택한다.
쓰다
$$
g_i \cdot e = v_i + e
$$
일부 $v_i \in V$의 경우. 이는 $g_i \cdot 1 = 1$이기 때문에 가능한다.
우리는 주장 요소 목록 $v_1, \ldots, v_r \in V$
확장 $E$의 동형사상 클래스를 결정한다.
이것을 증명하면 보조정리는 다음과 같습니다.
Ext 벡터공간은 $k$-벡터의 부분몫과 동형이다.
공간 $V^{\oplus r}$; 일부 세부 사항은 생략되었다.
$E$은 범주의 객체이므로
정의 \ref{etale-cohomology-definition-G-module-continuous}
우리는 다음과 같은 개방형 하위 그룹 $U$이 있다는 것을 알고 있다.
$u \cdot e = e$모두를 위해$u \in U$.
이제 $g \in G$ 중 하나를 선택하세요. 그러면 $gU$에는 $w$라는 단어가 포함되어 있다.
요소 $g_1, \ldots, g_r$.
$gu = w$라고 말하세요. $w \cdot e$ 요소는 다음과 같이 결정된다.
$v_1, \ldots, v_r$, $g \cdot e = (gu) \cdot e = w \cdot e$
너무이다.
\end{proof}

\begin{lemma} \label{etale-cohomology-lemma-profinite-group-cohomology-torsion} $G$를 유한위상군이라 하자. 그러면 \begin{enumerate} \item $H^i(G, M)$은 $i > 0$ 및 모든 $G$-가군 $M$에 대하여 꼬임군이다. \item $H^i(G, M) = 0$이다. 여기서 $M$은 $\mathbf{Q}$-벡터공간이다. \end{enumerate} \end{lemma}

\begin{proof}
증명 (1). 차원 이동을 통해 충분하다는 것을 알 수 있다.
$H^1(G, M)$가 모든 $G$-가군 $M$에 대해 비틀림임을 보여줍니다.
선택하세요완전열 $0 \to M \to I \to N \to 0$~와 함께$I$
$G$-가군의 범주의 단사 대상.
그러면 $H^1(G, M)$의 모든 요소는 요소의 이미지이다.
$y \in N^G$. $x \in I$를 $y$에 매핑하는 것을 선택한다.
안정제$U \subset G$~의$x$열려 있으므로
유한 인덱스 $r$가 있다. $g_1, \ldots, g_r \in G$를 시스템으로 설정
$G/U$의 대표자 수이다. 그러면 $\sum g_i(x)$는 불변이다.
요소$I$이는 다음과 같이 매핑된다.$ry$. 따라서$r$요소를 죽인다
$H^1(G, M)$부터 시작했다. 부분(2)은 다음과 같습니다.
$H^i(G, M)$는 $\mathbf{Q}$-벡터공간이자 비틀림이다.
\end{proof}





\section{테이트의 연속 코호몰로지} \label{etale-cohomology-section-continuous-group-cohomology}

\noindent
Tate의 연속형 코호몰로지 (\cite{Tate})는 다음의 복합체로 정의된다.
연속적인 불균일한 코체인. $M$이 다음과 같을 때 이를 정의할 수 있다.
연속적인 $G$ 동작이 부여된 임의의 위상 아벨 군.
즉, 복합체를 고려한다.
$$
C^\bullet_{cont}(G, M) :
M \to \text{지도}_{cont}(G, M) \to
\text{지도}_{cont}(G \times G, M) \to \ldots
$$
여기서 경계 맵은 규칙에 따라 $n \geq 1$에 대해 정의된다.
\begin{align*}
\text{d}(f)(g_1, \ldots, g_{n + 1})
& = g_1(f(g_2, \ldots, g_{n + 1})) \\
&
+ \sum\nolimits_{j = 1, \ldots, n}
(-1)^jf(g_1, \ldots, g_jg_{j + 1}, \ldots, g_{n + 1}) \\
&
+ (-1)^{n + 1}f(g_1, \ldots, g_n)
\end{align*}
$n = 0$의 경우 $m \in M$을 지도 $g \mapsto g(m) - m$로 보냅니다. 우리는 정의합니다
$$
H^i_{cont}(G, M) = H^i(C^\bullet_{cont}(G, M))
$$
복합체의 항은 $G$의 연속 맵을 포함하므로
$G$의 자체 곱을 위상 가군 $M$로 변환하면 명확하지 않습니다.
이는 위상 가군의 짧은 완전열을
길고 정확한 코호몰로지 시퀀스. 또 다른 어려움은 범주
위상학적 아벨 군은 아벨 범주가 아닙니다!

\medskip\noindent
그러나 이산 $G$-가군의 짧은 완전열은 다음을 제공한다.
연속적인 코체인의 복합체 중 짧은 완전열로 상승
따라서 길고 정확한 코호몰로지 연속 연속 시퀀스
코호몰로지 군 $H^i_{cont}(G, -)$.
따라서 범주 $\text{Mod}_G$에
정의 \ref{etale-cohomology-definition-G-module-continuous} 함자
$H^i_{cont}(G, M)$는 정의된 대로 동동 $\delta$-함자를 형성한다.
호몰로지, 절 \ref{homology-section-cohomological-delta-functor}.
코호몰로지 $H^i(G, M)$ 이후
정의 \ref{etale-cohomology-definition-galois-cohomology}
범용 $\delta$-함자이다.
(유도 범주, 보조정리 \ref{derived-lemma-right-derived-delta-functor})
우리는 표준지도를 얻습니다
$$
H^i(G, M) \longrightarrow H^i_{cont}(G, M)
$$
$M \in \text{Mod}_G$의 경우. 이 지도는 다음과 같은 것으로 알려져 있다.
동형사상 $G$가 추상 그룹인 경우(예: $G$에는
이산 위상) 또는 $G$가 유한 그룹인 경우
(여기에 향후 참조를 삽입하십시오).
이 지도를 표시하는 예가 동형사상이 아니라는 것을 알고 계시다면
위상 그룹 $G$ 및 $M \in \Ob(\text{Mod}_G)$
이메일을 보내주세요
\href{mailto:stacks.project@gmail.com}{stacks.project@gmail.com}.




\section{코호몰로지 점} \label{etale-cohomology-section-cohomology-point}

\noindent 이전 절에서의 논의의 결과로 우리는 체의 스펙트럼의 에탈 코호몰로지와 Galois 코호몰로지의 등가성을 얻습니다.

\begin{lemma} \label{etale-cohomology-lemma-equivalence-abelian-sheaves-point} $S = \Spec(K)$를 $K$와 체로 둡니다. $\overline{s}$를 $S$의 기하점으로 설정한다. $G = \text{Gal}_{\kappa(s)}$는 절대 갈루아 그룹을 나타냅니다. 줄기 함자는 범주의 동치 $$
\textit{Ab}(S_\etale) \longrightarrow \text{Mod}_G,
\quad
\mathcal{F} \longmapsto \mathcal{F}_{\overline{s}}.
$$ \end{lemma}를 유도한다.

\begin{proof} 정리 \ref{etale-cohomology-theorem-equivalence-sheaves-point}에서 집합과 $G$-집합의 층 사이의 동등성을 확인했다. 현재 보조정리는 아벨 층이 교환 그룹 법칙을 부여받은 집합의 층이고, $G$-가군은 단지 $G$-집합 교환 그룹 법칙이 부여된다. \end{proof}

\begin{lemma}
\label{etale-cohomology-lemma-compare-cohomology-point}
다음과 같은 표기법과 가정
보조정리 \ref{etale-cohomology-lemma-equivalence-abelian-sheaves-point}.
$\mathcal{F}$를 $\Spec(K)_\etale$에서 아벨 층으로 설정한다.
이는 $G$-가군 $M$에 해당한다.
그 다음에
\begin{enumerate}
\item $D(\textit{Ab})$에는 정식 동형사상이 있다.
$R\Gamma(S, \mathcal{F}) = R\Gamma_G(M)$,
\item $H_\etale^0(S, \mathcal{F}) = M^G$ 및
\item $H_\etale^q(S, \mathcal{F}) = H^q(G, M)$.
\end{enumerate}
\end{lemma}

\begin{proof} 보조정리 \ref{etale-cohomology-lemma-equivalence-abelian-sheaves-point}를 보조정리 \ref{etale-cohomology-lemma-global-sections-point}와 결합한다. \end{proof}

\begin{example}
\label{etale-cohomology-example-sheaves-point}
층 - $\Spec(K)_\etale$.
$G = \text{Gal}(K^{sep}/K)$를 $K$의 절대 갈루아군이라고 가정한다.
\begin{enumerate}
\item 상수층 $\underline{\mathbf{Z}/n\mathbf{Z}}$는 다음에 해당한다.
가군 $\mathbf{Z}/n\mathbf{Z}$ 간단한 $G$ 동작을 사용하여
\item 층 $\mathbf{G}_m|_{\Spec(K)_\etale}$는 다음에 해당한다.
$(K^{sep})^*$와 $G$-액션,
\item 층 $\mathbf{G}_a|_{\Spec(K^{sep})}$는 다음에 해당한다.
$(K^{sep}, +)$와 $G$-액션
\item 층 $\mu_n|_{\Spec(K^{sep})}$는 다음에 해당한다.
$\mu_n(K^{sep})$를 $G$-액션으로 사용한다.
\end{enumerate}
에 의해
비고 \ref{etale-cohomology-remark-special-case-fpqc-cohomology-quasi-coherent}
그리고
정리 \ref{etale-cohomology-theorem-picard-group}
코호몰로지 군에 대해 다음과 같은 식별 정보가 있다.
\begin{align*}
H_\etale^0(S_\etale, \mathbf{G}_m) & =
\Gamma(S, \mathcal{O}_S^*) \\
H_\etale^1(S_\etale, \mathbf{G}_m) & =
H_{Zar}^1(S, \mathcal{O}_S^*) = \Pic(S) \\
H_\etale^i(S_\etale, \mathbf{G}_a) & =
H_{Zar}^i(S, \mathcal{O}_S)
\end{align*}
또한, 누구에게나준연접층 $\mathcal{F}$~에$S_\etale$우리는
$$
H^i(S_\etale, \mathcal{F}) = H_{Zar}^i(S, \mathcal{F}),
$$
보다
정리 \ref{etale-cohomology-theorem-zariski-fpqc-quasi-coherent}.
특히 이는 다음과 같은 등식 시퀀스를 제공한다.
$$
0 =
\Pic(\Spec(K)) =
H_\etale^1(\Spec(K)_\etale, \mathbf{G}_m) =
H^1(G, (K^{sep})^*)
$$
이는 다름 아닌 힐베르트의 90 정리이다. 마찬가지로 $i \geq 1$의 경우
$$
0 = H^i(\Spec(K), \mathcal{O})
= H_\etale^i(\Spec(K)_\etale, \mathbf{G}_a)
= H^i(G, K^{sep})
$$
여기서 $K^{sep}$는 $K^{sep}$를 갈루아 가군으로 나타냅니다.
단체법으로. 이런 식으로 우리는 지금까지 수행한 작업을 고려할 수 있다.
Galois 코호몰로지 군을 계산하는 복잡한 방법이다.
\end{example}

\noindent 다음 결과는 호기심이므로 첫 독해에서는 건너뛰어야 한다.

\begin{lemma}
\label{etale-cohomology-lemma-all-modules-quasi-coherent}
$R$를 차원 $0$의 국소환으로 설정한다. $S = \Spec(R)$로 놔두세요.
그러면 매$\mathcal{O}_S$-가군~에$S_\etale$준일관적이다.
\end{lemma}

\begin{proof}
허락하다$\mathcal{F}$가 되다$\mathcal{O}_S$-가군~에$S_\etale$.
우리는 $\mathcal{F}$가 다음에 의해 결정된다는 것을 보여야 한다.
$R$-가군 $M = \Gamma(S, \mathcal{F})$.
더 정확하게 말하면, $\pi : X \to S$가 에탈이라면 우리는
$\Gamma(X, \mathcal{F}) = \Gamma(X, \pi^*\widetilde{M})$을 보여줍니다.

\medskip\noindent $\mathfrak m \subset R$를 극대 아이디얼로 두고 $\kappa$를 나머지 체로 둡니다. 대수학에 따르면 보조정리 \ref{algebra-lemma-local-dimension-zero-henselian} 국소환 $R$는 헨셀식이다. $X \to S$가 에탈인 경우 $X$의 기본 위상 공간은 사상, 보조정리 \ref{morphisms-lemma-etale-over-field}에 의해 이산되므로 $X$는 각각 하나의 점을 갖는 아핀 스킴의 분리된 합집합이다. 게다가, $X = \Spec(A)$가 아핀이고 하나의 점을 갖는다면, $R \to A$는 대수학, 보조정리 \ref{algebra-lemma-mop-up}에 의해 유한 에탈이다. 이 경우에는 $\Gamma(X, \mathcal{F}) = M \otimes_R A$임을 보여줘야 한다.

\medskip\noindent
함자 $A \mapsto A/\mathfrak m A$는 다음과 같은 동등성을 정의한다.
유한 에탈 $R$-대수의 범주
유한 분리 가능한 $\kappa$-대수의 범주를 사용하여
대수학, 보조정리 \ref{algebra-lemma-henselian-cat-finite-etale}.
먼저 $A/\mathfrak m A$ 경우를 고려해 보자.
Galois 확장은 다음과 같습니다.$\kappa$갈루아 그룹과 함께$G$.
각 $\sigma \in G$에 대해 $\sigma : A \to A$는
대응하는 자기형성$A$~ 위에$R$.
$N = \Gamma(X, \mathcal{F})$로 놔두세요.
그러면 $\Spec(\sigma) : X \to X$는 $S$에 대한 자기동형이다.
따라서 이에 의한 당김은 맵 $\sigma : N \to N$을 정의한다.
이는 $\sigma$ 선형 지도입니다: $\sigma(an) = \sigma(a) \sigma(n)$
$a \in A$ 및 $n \in N$의 경우.
준연접 가군에 갈루아 하강을 적용하겠습니다.
$\widetilde{N}$ $X$에 동형사상 권한을 부여함
에 대한 작업에서 발생$\sigma$~에$N$. 하강을 참조하라.보조정리
\ref{descent-lemma-galois-descent-more-general}.
이 보조정리는 동형사상 $N = N^G \otimes_R A$이 있음을 알려줍니다.
반면, 층 성질에 의해 $N^G = M$인 것이 분명하다.
$\mathcal{F}$의 경우. 따라서 필수 동형사상이 유지된다.

\medskip\noindent 일반적인 경우($A$ 로컬 및 유한 에탈 위의 $R$)는 갈루아 경우에서 다음과 같이 추론된다. $A \to B$ 유한 에탈을 선택하여 $B$가 $\kappa$ 위에 잔차 체 갈루아가 있는 로컬이 되도록 선택한다. $G = \text{Aut}(B/R) = \text{Gal}(\kappa_B/\kappa)$로 놔두세요. $H \subset G$를 갈루아 확장 $\kappa_B/\kappa_A$에 대응하는 갈루아 그룹이라고 가정한다. 그러면 위와 같이 $\Gamma(X, \mathcal{F}) = \Gamma(\Spec(B), \mathcal{F})^H$이 표시된다. Galois 확장(두 번 사용)의 결과로 원하는 대로 $$
\Gamma(X, \mathcal{F}) = (M \otimes_R B)^H = M \otimes_R A
$$을 얻습니다. \end{proof}








\section{코호몰로지 곡선} \label{etale-cohomology-section-cohomology-curves}

\noindent 다음 작업은 꼬임 계수를 사용하여 대수적으로 닫힌 체에 대한 부드러운 곡선의 에탈 코호몰로지를 계산하고, 특히 적어도 3 정도에서 사라지는 것을 보여주는 것이다. 이를 증명하기 위해 일반 지점에서 코호몰로지를 계산할 것이며 이는 일부 Galois 코호몰로지에 해당한다.





\section{Brauer 군}
\label{etale-cohomology-section-brauer-groups}

\noindent Brauer 군 의 체 은 유한 중심 단순 대수를 사용하여 정의된다. 이 절에서는 Brauer 군에 대한 관련 사실을 검토한다. 대부분의 내용은 Brauer 군, 절 \ref{brauer-section-introduction} 장에서 논의된다. 다른 참고 자료는 \cite{SerreCorpsLocaux}, \cite{SerreGaloisCohomology} 또는 \cite{Weil}를 참조하라.

\begin{theorem}
\label{etale-cohomology-theorem-central-simple-algebra}
$K$를 체로 설정한다. 단위, 결합(반드시 가환적일 필요는 없음)
$K$-대수 $A$ 다음은 동일한다.
\begin{enumerate}
\item $A$는 유한 중심 단순 $K$-대수이다.
\item $A$는 유한차원 $K$-벡터공간이고, $K$는 $A$의 중심이고,
그리고 $A$에는 중요한 양면 아이디얼이 없다.
\item $d \geq 1$가 존재한다.
$A \otimes_K \bar K \cong \text{Mat}(d \times d, \bar K)$,
\item $d \geq 1$가 존재한다.
$A \otimes_K K^{sep} \cong \text{Mat}(d \times d, K^{sep})$,
\item $d \geq 1$ 및 유한한 갈루아 확장 $K'/K$이 존재한다.
그렇게
$A \otimes_K K' \cong \text{Mat}(d \times d, K')$,
\item $n \geq 1$ 및 유한한 중앙 스큐 체 $D$가 존재한다.
$K$ 이상 $A \cong \text{Mat}(n \times n, D)$.
\end{enumerate}
정수 $d$를 $A$의 {\it 차수}라 한다.
\end{theorem}

\begin{proof} Brauer 군, 보조정리 \ref{brauer-lemma-finite-central-simple-algebra}의 사본이다. \end{proof}

\begin{lemma} \label{etale-cohomology-lemma-brauer-inverse} $A$를 $K$에 대한 유한 중심 단순 대수라고 가정한다. 그러면 $$
\begin{matrix}
A \otimes_K A^{opp} & \longrightarrow & \text{End}_K(A) \\
\ a \otimes a' & \longmapsto & (x \mapsto a x a')
\end{matrix}
$$는 $K$에 대한 대수학의 동형사상이다. \end{lemma}

\begin{proof} Brauer 군, 보조정리 \ref{brauer-lemma-inverse}를 참조하라. \end{proof}

\begin{definition}
\label{etale-cohomology-definition-brauer-equivalent}
두 개의 유한 중심 단순 대수$A_1$그리고$A_2$~ 위에$K$호출된다
{\it 유사} 또는 {\it 동등} $m, n \geq 1$
$\text{Mat}(n \times n, A_1) \cong \text{Mat}(m \times m, A_2)$. We write $A_1 \sim A_2$.
\end{definition}

\noindent Brauer 군, 보조정리 \ref{brauer-lemma-similar}에 의해 동치관계가 된다.

\begin{definition}
\label{etale-cohomology-definition-brauer-group}
허락하다$K$가 되다체. 그만큼{\it Brauer 군}~의$K$은집합 $\text{Br} (K)$
$K$에 대한 유한 중심 단순 대수의 유사성 클래스,
텐서곱에 의해 유도된 군법칙($K$ 이상). $A$ 클래스
$\text{Br}(K)$은 $[A]$로 표시된다. 중립 요소는
$[K] = [\text{Mat}(d \times d, K)]$ 모든 $d \geq 1$에 대해.
\end{definition}

\noindent
앞의 보조정리는 역이 존재하고 $-[A] = [A^{opp}]$를 의미한다.
체의 Brauer 군은 항상 비틀림이다.
실제로 $[A]$에는 $\deg(A)$를 나누는 순서가 있음을 알 수 있다.
유한 중심 단순 대수 $A$ (참조
보조정리 \ref{etale-cohomology-lemma-annihilated-by-degree}).
일반적으로 Brauer 군은 예에 대해 유한하게 생성되지 않습니다.
아르키메데스가 아닌 로컬 체의 Brauer 군은 $\mathbf{Q}/\mathbf{Z}$이다.
$\mathbf{C}(x, y)$의 Brauer 군은 셀 수 없다.

\begin{lemma}
\label{etale-cohomology-lemma-central-simple-algebra-pgln}
\begin{slogan}
중앙 단순 대수는 PGL의 Galois 코호몰로지로 분류된다.
\end{slogan}
$K$를 체로 두고 $K^{sep}$를 분리 가능한 대수적 종결자로 둡니다.
그런 다음 동형사상 클래스의 중앙 단순 대수학의 집합
$d$ 위의 $K$은 비아벨식 코호몰로지와 전단사를 사용한다.
$H^1(\text{Gal}(K^{sep}/K), \text{PGL}_d(K^{sep}))$.
\end{lemma}

\begin{proof}[증명의 스케치.] Skolem-Noether 정리 (Brauer 군, 정리 \ref{brauer-theorem-skolem-noether} 참조)는 모든 체에 대해 다음을 의미한다. $L$ 그룹 $\text{Aut}_{L\text{-대수학}}(\text{Mat}_d(L))$는 $\text{PGL}_d(L)$와 같습니다. 정리 \ref{etale-cohomology-theorem-central-simple-algebra}에 의해, $d$ 정도의 중앙 단순 대수는 $K$-대수 $\text{Mat}_d(K)$의 형태에 해당한다는 것을 알 수 있다. 결합하여 우리는 $d$ 정도의 동형사상 클래스 클래스가 $H^1(\text{Gal}(K^{sep}/K), \text{PGL}_d(K^{sep}))$의 요소에 해당한다는 것을 알 수 있다. 비틀기에 대한 자세한 내용은 예 \cite{SilvermanEllipticCurves}를 참조하라. \end{proof}

\noindent
만약에$A$유한 중심 단순 대수학이다$d$이상체 $K$,
우리는 $\xi_A$ 해당 코호몰로지 클래스를 나타냅니다.
$H^1(\text{Gal}(K^{sep}/K), \text{PGL}_d(K^{sep}))$.
짧은 완전열을 고려해보세요.
$$
1 \to (K^{sep})^* \to \text{GL}_d(K^{sep}) \to \text{PGL}_d(K^{sep}) \to 1,
$$
이는 길고 정확한 코호몰로지 시퀀스(최대 2 수준)를 생성한다.
공동경계 지도
$$
\delta_d :
H ^1(\text{Gal}(K^{sep}/K), \text{PGL}_d(K^{sep}))
\longrightarrow
H^2(\text{Gal}(K^{sep}/K), (K^{sep})^*).
$$
명시적으로 이는 다음과 같이 제공됩니다: $\xi$가 코호몰로지 클래스인 경우
1-코사이클 $(g_\sigma)$로 표시되는 경우 $\delta_d(\xi)$는
2-코사이클 클래스
\begin{equation}
\label{etale-cohomology-equation-two-cocycle}
(\sigma, \tau)
\longmapsto
\tilde g_\sigma^{-1} \tilde g_{\sigma \tau} \sigma(\tilde g_\tau^{-1})
\in (K^{sep})^*
\end{equation}
여기서 $\tilde g_\sigma \in \text{GL}_d(K^{sep})$는 $g_\sigma$의 리프트이다.
이를 사용하여 지도를 명시적으로 만들 수 있다.
$$
\delta : \text{Br}(K) \longrightarrow H^2(\text{Gal}(K^{sep}/K), (K^{sep})^*),
\quad
[A] \longmapsto \delta_{\deg A} (\xi_A)
$$
다음과 같이. 추정하다$A$학위가 있다$d$~ 위에$K$. 선택하세요동형사상
$\varphi : \text{Mat}_d(K^{sep}) \to A \otimes_K K^{sep}$. 을 위한
$\sigma \in \text{Gal}(K^{sep}/K)$ 요소를 선택하세요
$\tilde g_\sigma \in \text{GL}_d(K^{sep})$ 이렇게
$\varphi^{-1} \circ \sigma(\varphi)$는 지도와 같습니다.
$x \mapsto \tilde g_\sigma x \tilde g_\sigma^{-1}$. $H^2$의 클래스
는 두 개의 공주기(\ref{etale-cohomology-equation-two-cocycle})로 정의된다.

\begin{theorem} \label{etale-cohomology-theorem-brauer-delta} $K$를 분리 가능한 대수적 종결부 $K^{sep}$가 있는 체로 둡니다. 위에서 정의한 맵 $\delta : \text{Br}(K) \to H^2(\text{Gal}(K^{sep}/K), (K^{sep})^*)$는 그룹 동형사상이다. \end{theorem}

\begin{proof}[증명]의 스케치 $\delta$가 그룹 동형사상을 정의한다는 것, 즉 $\delta(A \otimes_K B) = \delta(A) + \delta(B)$를 정의하기 위해 공자전거를 사용하여 직접 계산한다.

\medskip\noindent
$\delta$의 주입성. 아벨의 경우($d = 1$)에서는
신분증
$$
H^1(\text{Gal}(K^{sep}/K), \text{GL}_d(K^{sep})) =
H_\etale^1(\Spec(K), \text{GL}_d(\mathcal{O}))
$$
후자는 fpqc 하강에 의해 사소한 것이다. 만약에 이것이 사실이었다면
아벨이 아닌 경우, 이는 $\delta$의 주입성을 쉽게 암시한다. (보다
\cite{SGA4.5}.) 오히려 이를 증명하기 위해 $\delta([A])$을 다음과 같이 재해석할 수 있다.
존재의 방해$K$-벡터공간 $V$왼쪽으로$A$-가군
$\dim_K V = \deg A$와 같은 구조이다. $V$이 존재하는 경우,
$A \cong \text{End}_K(V)$이 있다.

\medskip\noindent
전사성을 위해 다음 중 하나를 선택하세요.
코호몰로지 클래스 $\xi \in H^2(\text{Gal}(K^{sep}/K), (K^{sep})^*)$,
그러면 유한한 갈루아 확장 $K^{sep}/K'/K$이 존재한다.
$\xi$는 일부 이미지이다.
$\xi' \in H^2(\text{Gal}(K'|K), (K')^*)$. 그럼 쓰세요
명시적인 중앙 단순 대수학 아래로$K$데이터를 사용하여$K', \xi'$.
\end{proof}










\section{스킴} \label{etale-cohomology-section-brauer-scheme}의 Brauer 군

\noindent
$S$를 스킴으로 설정한다. $\mathcal{O}_S$-대수
$\mathcal{A}$는 지역적으로 에탈인 경우 {\it Azumaya}라고 한다.
대수학, 즉 에탈 피복이 존재하는 경우
$\mathcal{U} = \{ \varphi_i : U_i \to S\}_{i \in I}$ 이렇게
$\varphi_i^*\mathcal{A} \cong \text{Mat}_{d_i}(\mathcal{O}_{U_i})$
일부 $d_i \geq 1$의 경우. 그런 두 가지
$\mathcal{A}$ 및 $\mathcal{B}$는 {\it 등가물}이라고 한다.
유한 로컬 자유 $\mathcal{O}_S$-가군 $\mathcal{F}$ 및 $\mathcal{G}$
$s \in S$마다 양수 순위를 가지므로
$$
\mathcal{A} \otimes_{\mathcal{O}_S}
\SheafHom_{\mathcal{O}_S}(\mathcal{F}, \mathcal{F})
\cong
\mathcal{B} \otimes_{\mathcal{O}_S}
\SheafHom_{\mathcal{O}_S}(\mathcal{G}, \mathcal{G})
$$
$\mathcal{O}_S$-대수와 같습니다. {\it Brauer 군}
$S$는 Azumaya의 동등 클래스의 집합 $\text{Br}(S)$이다.
$\mathcal{O}_S$-텐서곱에 의해 유도된 연산을 사용하는 대수(이상
$\mathcal{O}_S$).

\begin{lemma}
\label{etale-cohomology-lemma-end-unique-up-to-invertible}
$S$를 스킴으로 설정한다. $\mathcal{F}$ 및 $\mathcal{G}$를 로컬로 유한하게 만듭니다.
양수 순위의 $\mathcal{O}_S$-가군 중 층을 무료로 제공한다. 만약 거기에
동형사상이 존재한다.
$\SheafHom_{\mathcal{O}_S}(\mathcal{F}, \mathcal{F}) \cong
\SheafHom_{\mathcal{O}_S}(\mathcal{G}, \mathcal{G})$ /
$\mathcal{O}_S$-대수학이면 가역층이 존재한다.
$\mathcal{L}$ $S$에 대해
$\mathcal{F} \otimes_{\mathcal{O}_S} \mathcal{L} \cong \mathcal{G}$
그리고 이 동형사상은 주어진 동형사상을 유도한다.
내형성 대수학.
\end{lemma}

\begin{proof}
동형사상 문제를 해결하세요.
$\SheafHom_{\mathcal{O}_S}(\mathcal{F}, \mathcal{F}) \to
\SheafHom_{\mathcal{O}_S}(\mathcal{G}, \mathcal{G})$.
층 $\mathcal{L} \subset \SheafHom(\mathcal{F}, \mathcal{G})$를 고려하라.
로컬 동형사상에 의해 $\mathcal{O}_S$-가군으로 생성됨
$\varphi : \mathcal{F} \to \mathcal{G}$ 다음과 같은 활용이 가능한다.
$\varphi$는 내형성 대수의 주어진 동형사상이다.
로컬 계산($\mathcal{F}$ 및 $\mathcal{G}$인 경우로 축소)
유한 자유이고 $S$는 유사함)은 $\mathcal{L}$이 가역적임을 보여줍니다.
또 다른 로컬 계산은 평가 맵이
$$
\mathcal{F} \otimes_{\mathcal{O}_S} \mathcal{L} \longrightarrow \mathcal{G}
$$
동형사상이다.
\end{proof}

\noindent 다음 보조정리의 증명에 주어진 인수는 \cite{Saltman-torsion}에서 확인할 수 있다.

\begin{lemma}
\label{etale-cohomology-lemma-annihilated-by-degree}
\begin{reference}
\cite{Saltman-torsion}에서 가져온 인수이다.
\end{reference}
$S$를 스킴으로 설정한다. $\mathcal{A}$를 Azumaya 대수라고 하자.
지역적으로는 계급이 없음$d^2$~ 위에$S$. 그럼 수업은
~의$\mathcal{A}$에서Brauer 군~의$S$에 의해 소멸된다$d$.
\end{lemma}

\begin{proof}
에탈 피복 $\{U_i \to S\}$를 선택하고 동형사상을 선택한다.
$\mathcal{A}|_{U_i} \to \SheafHom(\mathcal{F}_i, \mathcal{F}_i)$
일부 지역 무료 $\mathcal{O}_{U_i}$-가군 $\mathcal{F}_i$
순위 $d$. ($\mathcal{F}_i$는 무료라고 가정할 수 있다.)
구성
$$
p_i : \mathcal{F}_i^{\otimes d} \to
\wedge^d(\mathcal{F}_i) \to \mathcal{F}_i^{\otimes d}
$$
첫 번째 화살표는 일반적인 투영이고 두 번째 화살표는
$\mathcal{F}_i$의 상단 외부 전력의 동형사상
변환하는 $\mathcal{F}_i^{\otimes d}$의 절 서브모듈
대칭 그룹의 동작에 따라 기호 문자에 따라
$d$ 글자로. 그러면 $p_i^2 = d! p_i$이고 $p_i$의 순위는 $1$이다.
주어진 동형사상을 사용하여
$\mathcal{A}|_{U_i} \to \SheafHom(\mathcal{F}_i, \mathcal{F}_i)$
그리고 정식 동형사상
$$
\SheafHom(\mathcal{F}_i, \mathcal{F}_i)^{\otimes d} =
\SheafHom(\mathcal{F}_i^{\otimes d}, \mathcal{F}_i^{\otimes d})
$$
$p_i$을 $\mathcal{A}^{\otimes d}$의 절로 생각할 수 있다.
$U_i$ 이상. 우리는 주장 $p_i|_{U_i \times_S U_j} = p_j|_{U_i \times_S U_j}$
$\mathcal{A}^{\otimes d}$의 절로. 즉, 신청
보조정리 \ref{etale-cohomology-lemma-end-unique-up-to-invertible}
우리는 가역층 $\mathcal{L}_{ij}$ 및 표준 동형사상을 얻습니다.
$$
\mathcal{F}_i|_{U_i \times_S U_j} \otimes \mathcal{L}_{ij}
\longrightarrow
\mathcal{F}_j|_{U_i \times_S U_j}.
$$
이 동형사상을 사용하면 $p_i$가 $p_j$에 매핑되는 것을 볼 수 있다.
$\mathcal{A}^{\otimes d}$는 $S_\etale$의 층이므로
(명제 \ref{etale-cohomology-proposition-quasi-coherent-sheaf-fpqc}) 우리는 표준을 찾습니다
대역 절단 $p \in \Gamma(S, \mathcal{A}^{\otimes d})$. 로컬 계산
그것을 보여줍니다
$$
\mathcal{H} =
\Im(\mathcal{A}^{\otimes d} \to \mathcal{A}^{\otimes d}, f \mapsto fp)
$$
랭크 $d^d$ 의 로컬 프리 가군 이고 그 (왼쪽) 곱셈입니다
$\mathcal{A}^{\otimes d}$에 의해 동형사상을 유도한다.
$\mathcal{A}^{\otimes d} \to \SheafHom(\mathcal{H}, \mathcal{H})$.
즉, $\mathcal{A}^{\otimes d}$는 사소한 요소이다.
원하는 대로 $S$의 Brauer 군을 선택한다.
\end{proof}

\noindent
이 설정에서 동형사상 $\delta$의 아날로그
정리 \ref{etale-cohomology-theorem-brauer-delta}
지도이다
$$
\delta_S: \text{Br}(S) \to H_\etale^2(S, \mathbf{G}_m).
$$
$\delta_S$가 단사라는 것은 사실이다. $S$이 준 컴팩트하거나
연결이면 $\text{Br}(S)$는 비틀림 그룹이므로 이 경우
$\delta_S$의 상은 $S$의 {\it 코호몰로지적 Brauer 군}에 포함된다.
$$
\text{Br}'(S) := H_\etale^2(S, \mathbf{G}_m)_\text{torsion}.
$$
따라서 $S$가 준컴팩트 또는 연결인 경우 $\text{Br}(S)
\subset \text{Br}'(S)$. 이것이 항상 평등한 것은 아닙니다.
분리되지 않은 단일 표면$S$$\text{Br}(S) \subset
\text{Br}'(S)$ 이다 strict inclusion. 만약 $S$는 준투사법이다.
$\text{Br}(S) = \text{Br}'(S)$. 그러나 이것이 적용되는지 여부는 알려지지 않았습니다.
예를 들어 $\mathbf{C}$에 대한 부드러운 적절한 다양체이다.




\section{아르틴-슈라이어 수열} \label{etale-cohomology-section-artin-schreier}

\noindent $p$를 소수라고 하자. $S$를 표수 $p$에서 스킴으로 설정한다. {\it Artin-Schreier} 시퀀스는 짧은 완전열 $$
0 \longrightarrow \underline{\mathbf{Z}/p\mathbf{Z}}_S \longrightarrow
\mathbf{G}_{a, S} \xrightarrow{F-1} \mathbf{G}_{a, S} \longrightarrow 0
$$이다. 여기서 $F - 1$는 맵 $x \mapsto x^p - x$이다.

\begin{lemma}
\label{etale-cohomology-lemma-vanishing-affine-char-p-p}
$p$를 소수로 둡니다. $S$를 표수 $p$의 스킴으로 설정한다.
\begin{enumerate}
\item 만약 $S$이 유사하다면,
$H_\etale^q(S, \underline{\mathbf{Z}/p\mathbf{Z}}) = 0$ 모두에 대해
$q \geq 2$.
\item $S$가 다음의 준컴팩트하고 준분리된 스킴인 경우
차원 $d$, 그런 다음 $H_\etale^q(S, \underline{\mathbf{Z}/p\mathbf{Z}}) = 0$
모든 $q \geq 2 + d$에 대해.
\end{enumerate}
\end{lemma}

\begin{proof} 구조층의 에탈 코호몰로지는 기본 위상 공간(정리 \ref{etale-cohomology-theorem-zariski-fpqc-quasi-coherent})의 코호몰로지와 같습니다. 첫 번째 진술은 Artin-Schreier 완전열 및 소멸 의 코호몰로지 의 구조층 위의 아핀 스킴 (코호몰로지 의 스킴, 보조정리 \ref{coherent-lemma-quasi-coherent-affine-cohomology-zero}). 두 번째 진술은 코호몰로지, 명제 \ref{cohomology-proposition-cohomological-dimension-spectral}의 소멸에서 동일한 인수와 $S$가 스펙트럼 공간(성질, 보조정리이라는 사실이 뒤따릅니다. \ref{properties-lemma-quasi-compact-quasi-separated-spectral}). \end{proof}

\begin{lemma}
\label{etale-cohomology-lemma-F-1}
$k$를 표수 $p > 0$의 대수적으로 닫힌 체로 설정한다.
$V$를 유한 차원 $k$-벡터공간이라고 가정한다. $F : V \to V$
frobenius 선형 맵, 즉 다음과 같은 추가 맵이어야 한다.
$F(\lambda v) = \lambda^p F(v)$모두를 위해$\lambda \in k$그리고$v \in V$.
그런 다음 $F - 1 : V \to V$는 핵 유한 차원의 전사이다.
$\mathbf{F}_p$-벡터공간의 차원 $\leq \dim_k(V)$.
\end{lemma}

\begin{proof} $F = 0$이면 명령문이 유지된다. $V$를 $F$-안정적인 하위 벡터 공간으로 필터링하여 등급이 매겨진 각 조각에 대해 진술이 유지되면 $(V, F)$에 대해 유지된다. 이 두 가지 설명을 결합하면 $F$의 핵이 0이라고 가정할 수 있다.

\medskip\noindent
기준을 선택하세요$v_1, \ldots, v_n$~의$V$그리고 쓰다
$F(v_i) = \sum a_{ij} v_j$. $v = \sum \lambda_i v_i$를 확인하세요.
핵 필요충분조건은 $\sum \lambda_i^p a_{ij} v_j = 0$에 있다.
$k$는 대수적으로 닫혀 있으므로 이는 행렬 $(a_{ij})$을 의미한다.
반전 가능한다. $(b_{ij})$를 그 반대라고 하자. 그런 다음 $F - 1$를 확인하려면
전사라는 것은 $w = \sum \mu_i v_i \in V$를 선택하고 풀려고 노력하는 것이다.
$$
(F - 1)(\sum \lambda_iv_i) =
\sum \lambda_i^p a_{ij} v_j - \sum \lambda_j v_j = \sum \mu_j v_j
$$
이는 다음과 같습니다.
$$
\sum \lambda_j^p v_j - \sum b_{ij} \lambda_i v_j = \sum b_{ij} \mu_i v_j
$$
다시 말해서
$$
\lambda_j^p - \sum b_{ij} \lambda_i = \sum b_{ij} \mu_i,
\quad j = 1, \ldots, \dim(V).
$$
대수학
$$
A = k[x_1, \ldots, x_n]/
(x_j^p - \sum b_{ij} x_i - \sum b_{ij} \mu_i)
$$
$k$에 대해 표준 평활이다.
(대수학, 정의 \ref{algebra-definition-standard-smooth})
왜냐하면 행렬 $(b_{ij})$은 가역적이고 편도함수는
$x_j^p$는 0이다. 의 한 기저는 $A$ 위의 $k$ 이다 집합 의 단항식들
$x_1^{e_1} \ldots x_n^{e_n}$를 $e_i < p$로 사용하므로 $\dim_k(A) = p^n$이다.
$k$는 대수적으로 닫혀 있으므로 $\Spec(A)$는 정확히
$p^n$ 포인트. $F - 1$는 전사이고 모든 올
$p^n$ 포인트가 있다. 즉, $F - 1$의 핵은 $p^n$ 요소로 구성된 그룹이다.
\end{proof}

\begin{lemma} \label{etale-cohomology-lemma-F-2} 보조정리 \ref{etale-cohomology-lemma-F-1}의 상황에서 $k'/k$를 대수적으로 닫힌 체의 확장으로 둡니다. 집합 $V' = V \otimes_k k'$ 및 $F' : V' \to V'$는 $F$에 의해 유도된 Frobenius 선형 맵이다. 그런 다음 $\Ker(F - 1)$는 $\Ker(F' - 1)$에 동형으로 매핑된다. \end{lemma}

\begin{proof}
$v_1, \ldots, v_n$를 $V$의 기초로 두고 $F(v_i) = \sum a_{ij}v_j$라고 말한다.
그러면 $v = \sum \lambda_i v_i$는 $\Ker(F - 1)$ 필요충분조건은에 있다.
$\sum \lambda_i^pa_{ij} = \lambda_j$~을 위한$j = 1, \ldots, n$.
보조정리 \ref{etale-cohomology-lemma-F-1} 이 방정식에 의해
폐쇄된 하위 계획 $Z$을 잘라냅니다.
$\mathbf{A}^n_k = \Spec(k[\lambda_1, \ldots, \lambda_n])$
$\Spec(k)$에 대해 유한한다. 그런 다음 $Z(k) = Z(k')$ 결론을 내립니다.
\end{proof}

\begin{lemma} \label{etale-cohomology-lemma-top-cohomology-coherent} $X$를 체 $k$에 걸쳐 유한 유형의 분리된 스킴으로 둡니다. $\mathcal{F}$는 $\mathcal{O}_X$-가군의 일관된 층라고 가정한다. 그런 다음 $\dim_k H^d(X, \mathcal{F}) < \infty$ 여기서 $d = \dim(X)$이다. \end{lemma}

\begin{proof}
$d$에 대한 귀납법을 통해 이를 증명하겠습니다. $d = 0$ 사례는 다음과 같은 이유로 유지된다.
이 경우 $X$는 유한 차원 $k$-대수학 $A$의 스펙트럼이다.
(다양체, 보조정리 \ref{varieties-lemma-algebraic-scheme-dim-0})
그리고 모든 일관된 층 $\mathcal{F}$는 유한 $A$-가군에 해당한다.
$M = H^0(X, \mathcal{F})$에는 $\dim_k M < \infty$이 있다.

\medskip\noindent
$d > 0$를 가정하고 결과는 분리된 스킴에 대해 표시되었다.
차원 $< d$의 유한 유형이다. 스킴 $X$는 뇌터이다. 고려하다
그만큼성질 $\mathcal{P}$일관된층~에$X$규칙에 의해 정의됨
$$
\mathcal{P}(\mathcal{F}) \Leftrightarrow
\dim_k H^d(X, \mathcal{F}) < \infty
$$
우리는 다음의 결과를 사용할 것이다.
코호몰로지 중 스킴, 보조정리 \ref{coherent-lemma-property-initial}
$\mathcal{P}$가 $X$의 모든 일관성 있는 층에 대해 성립함을 증명한다.

\medskip\noindent $$
0 \to \mathcal{F}_1 \to \mathcal{F} \to \mathcal{F}_2 \to 0
$$를 $X$에서 일관된 층의 짧은 완전열로 설정한다. 코호몰로지 $$
H^d(X, \mathcal{F}_1) \to
H^d(X, \mathcal{F}) \to
H^d(X, \mathcal{F}_2)
$$의 긴 완전열을 고려하라. 따라서 $\mathcal{P}$가 $\mathcal{F}_1$ 및 $\mathcal{F}_2$에 대해 유지되면 $\mathcal{F}$에 대해 유지된다.

\medskip\noindent $Z \subset X$를 정역 폐쇄형 하위 체계라고 가정한다. $\mathcal{I}$를 $Z$에서 아이디얼의 일관된 층으로 설정한다. 증명을 마치려면 $H^d(X, i_*\mathcal{I}) = H^d(Z, \mathcal{I})$가 유한 차원임을 보여야 한다. $\dim(Z) < d$인 경우 코호몰로지 군이 0이 되기 때문에 결과가 유지됩니다(코호몰로지, 명제 \ref{cohomology-proposition-vanishing-Noetherian}). 이런 식으로 우리는 다음 단락에서 논의되는 상황으로 축소된다.

\medskip\noindent $X$는 차원 $d$의 다양체이고 $\mathcal{F} = \mathcal{I}$는 일관성 있는 아이디얼라고 가정한다. 층. 이 경우 짧은 완전열 $$
0 \to \mathcal{I} \to \mathcal{O}_X \to i_*\mathcal{O}_Z \to 0
$$가 있다. 여기서 $i : Z \to X$는 $\mathcal{I}$에 의해 정의된 폐쇄형 하위 체계이다. 귀납 가설에 의해 우리는 $H^{d - 1}(Z, \mathcal{O}_Z) = H^{d - 1}(X, i_*\mathcal{O}_Z)$가 유한 차원이라는 것을 알 수 있다. 따라서 구조층에 대한 결과를 증명하는 것으로 충분하다는 것을 알 수 있다.

\medskip\noindent
Chow의 보조정리를 적용할 수 있다.
(코호몰로지 중 스킴, 보조정리 \ref{coherent-lemma-chow-Noetherian})
사상 $X \to \Spec(k)$로. 따라서 우리는 다이어그램을 얻습니다.
$$
\xymatrix{
X \ar[rd]_g & X' \ar[d]^-{g'} \ar[l]^\pi \ar[r]_i & \mathbf{P}^n_k \ar[dl] \\
& \Spec(k) &
}
$$
Chow의 보조정리의 진술과 같습니다. 또한 $U \subset X$를
$\pi^{-1}(U) \to U$이 동형사상이 되도록 조밀 개방형 하위 구성표.
$X'$도 다양체라고 가정할 수 있다.
코호몰로지 중 스킴, 비고 \ref{coherent-remark-chow-Noetherian}.
사상 $i' = (i, \pi) : X' \to \mathbf{P}^n_X$는
 닫힌 몰입 (앞의 인용문헌). 따라서
$$
\mathcal{L} = i^*\mathcal{O}_{\mathbf{P}^n_k}(1) \cong
(i')^*\mathcal{O}_{\mathbf{P}^n_X}(1)
$$
$\pi$-상대적으로 충분함(예의 경우)
사상, 보조정리 \ref{morphisms-lemma-characterize-ample-on-finite-type}).
따라서 스킴의 코호몰로지에 의해, 보조정리 \ref{coherent-lemma-kill-by-twisting}
$n \geq 0$이 존재한다.
$R^p\pi_*\mathcal{L}^{\otimes n} = 0$모두를 위해$p > 0$.
두자 $\mathcal{G} = \pi_*\mathcal{L}^{\otimes n}$.
0이 아닌 것을 선택하세요대역 절단 $s$~의$\mathcal{L}^{\otimes n}$.
$\mathcal{G} = \pi_*\mathcal{L}^{\otimes n}$ 이후, 절 $s$
$\mathcal{G}$의 절에 해당한다. 즉, 지도
$\mathcal{O}_X \to \mathcal{G}$.
$s|_U \not = 0$는 $X'$이므로 다양체 및 $\mathcal{L}$이다.
반전 가능, $\mathcal{O}_X|_U \to \mathcal{G}|_U$
0이 아닙니다. $\mathcal{G}|_U = \mathcal{L}^{\otimes n}|_{\pi^{-1}(U)}$
가역적이라는 결론을 내리면 짧은 완전열
$$
0 \to \mathcal{O}_X \to \mathcal{G} \to \mathcal{Q} \to 0
$$
여기서 $\mathcal{Q}$는 일관되고 적절한 방식으로 지원된다.
$X$의 하위 체계가 폐쇄되었다. 귀납법을 사용하기 전처럼 논쟁하기
가설, 우리는 그것을 본다
$\dim H^d(X, \mathcal{G}) < \infty$을 증명하는 것으로 충분한다.

\medskip\noindent Leray 스펙트럼 열 (코호몰로지, 보조정리 \ref{cohomology-lemma-apply-Leray})에 따르면 $H^d(X, \mathcal{G}) = H^d(X', \mathcal{L}^{\otimes n})$이 표시된다. $\overline{X}' \subset \mathbf{P}^n_k$를 $X'$의 종료로 설정한다. 그러면 $\overline{X}'$는 차원 $d$의 $k$에 대한 사영 다양체이고 $X' \subset \overline{X}'$는 조밀 개방형이다. 가역층 $\mathcal{L}^{\otimes n}$는 $\mathcal{O}_{\overline{X}'}(n)$에서 $X'$까지의 제한 사항이다. 코호몰로지, 명제 \ref{cohomology-proposition-cohomological-dimension-spectral}에 의해 지도 $$
H^d(\overline{X}', \mathcal{O}_{\overline{X}'}(n))
\longrightarrow
H^d(X', \mathcal{L}^{\otimes n})
$$는 전사이다. 왼쪽의 코호몰로지 군은 스킴의 코호몰로지에 의해 차원으로 유한하므로 보조정리 \ref{coherent-lemma-coherent-projective} 증명이 완성된다. \end{proof}

\begin{lemma}
\label{etale-cohomology-lemma-vanishing-variety-char-p-p}
$X$를 대수적으로 닫힌 식에 대해 유한 유형으로 분리하겠습니다.
체 $k$의 표수 $p > 0$. 그 다음에
$H_\etale^q(X, \underline{\mathbf{Z}/p\mathbf{Z}}) = 0$
$q \geq dim(X) + 1$.
\end{lemma}

\begin{proof} $d = \dim(X)$로 둡니다. 보조정리 \ref{etale-cohomology-lemma-vanishing-affine-char-p-p}에 설정된 소멸에 의해 $H_\etale^{d + 1}(X, \underline{\mathbf{Z}/p\mathbf{Z}}) = 0$임을 나타내면 충분하다. 보조정리 \ref{etale-cohomology-lemma-top-cohomology-coherent}에 의해 $H^d(X, \mathcal{O}_X)$가 유한 차원 $k$-벡터공간임을 알 수 있다. 따라서 Artin-Schreier 시퀀스와 관련된 긴 정확한 코호몰로지 시퀀스는 $$
H^d(X, \mathcal{O}_X) \xrightarrow{F - 1}
H^d(X, \mathcal{O}_X) \to H^{d + 1}_\etale(X, \mathbf{Z}/p\mathbf{Z}) \to 0
$$로 끝납니다. 에 의해 보조정리 \ref{etale-cohomology-lemma-F-1} 이 시퀀스의 맵 $F - 1$는 전사이다. 이는 보조정리를 증명한다. \end{proof}

\begin{lemma}
\label{etale-cohomology-lemma-finiteness-proper-variety-char-p-p}
대수적으로 닫힌 방정식에 대해 $X$를 적절한 스킴으로 설정한다.
체 $k$의 표수 $p > 0$. 그 다음에
\begin{enumerate}
\item $H_\etale^q(X, \underline{\mathbf{Z}/p\mathbf{Z}})$
유한하다$\mathbf{Z}/p\mathbf{Z}$-가군모두를 위해$q$, 그리고
\item $H^q_\etale(X, \underline{\mathbf{Z}/p\mathbf{Z}}) \to
H^q_\etale(X_{k'}, \underline{\mathbf{Z}/p\mathbf{Z}})$
$k'/k$가 대수적으로 확장된 경우 동형사상이다.
체를 닫았습니다.
\end{enumerate}
\end{lemma}

\begin{proof}
스킴의 코호몰로지에 의해, 보조정리
\ref{coherent-lemma-proper-over-affine-cohomology-finite}
그리고 코호몰로지의 비교
정리 \ref{etale-cohomology-theorem-zariski-fpqc-quasi-coherent} 코호몰로지 군
$H^q_\etale(X, \mathbf{G}_a) = H^q(X, \mathcal{O}_X)$는
유한차원 $k$-벡터공간. 따라서
보조정리 \ref{etale-cohomology-lemma-F-1} 긴 정확한 코호몰로지 시퀀스
Artin-Schreier 서열과 연관되어 다음과 같이 분할된다.
짧은 완전열들
$$
0 \to H_\etale^q(X, \underline{\mathbf{Z}/p\mathbf{Z}}) \to
H^q(X, \mathcal{O}_X) \xrightarrow{F - 1} H^q(X, \mathcal{O}_X) \to 0
$$
그리고 코호몰로지 군의 $\mathbf{F}_p$-차원
$H_\etale^q(X, \underline{\mathbf{Z}/p\mathbf{Z}})$는 다음보다 작거나 같습니다.
$k$-차원의 벡터공간 $H^q(X, \mathcal{O}_X)$.
이것은 첫 번째 진술을 증명한다. 두 번째 진술은 다음과 같습니다
보조정리 \ref{etale-cohomology-lemma-F-2}에서
$H^q(X, \mathcal{O}_X) \otimes_k k' \to H^q(X_{k'}, \mathcal{O}_{X_{k'}})$
평평한 베이스 변경에 의한 동형사상이다.
(코호몰로지 중 스킴,
보조정리 \ref{coherent-lemma-flat-base-change-cohomology}).
\end{proof}








\section{국소상수층}
\label{etale-cohomology-section-locally-constant}

\noindent 이 절은 에탈 사이트에 대한 사이트, 절 \ref{sites-modules-section-locally-constant}의 가군과 유사한다.

\begin{definition}
\label{etale-cohomology-definition-finite-locally-constant}
$X$를 스킴으로 설정한다.
$\mathcal{F}$를 $X_\etale$에서 집합의 층으로 설정한다.
\begin{enumerate}
\item $E$를 집합으로 둡니다. $\mathcal{F}$는
{\it 상수층 값 $E$}, $\mathcal{F}$인 경우
전층 $U \mapsto E$의 층화.
표기법: $\underline{E}_X$ 또는 $\underline{E}$.
\item $\mathcal{F}$는 {\it 상수층}라고 말한다.
(1)에서와 같이 층과 동형이다.
\item $\mathcal{F}$는 {\it 로컬 상수}라고 말한다.
피복 $\{U_i \to X\}$ 즉, $\mathcal{F}|_{U_i}$는 상수층이다.
\item $\mathcal{F}$는 {\it 유한 지역 상수}라고 말한다.
는 지역적으로 상수이고 상수 값은 유한 집합이다.
\end{enumerate}
$\mathcal{F}$를 $X_\etale$에서 아벨 군의 층으로 설정한다.
\begin{enumerate}
\item $A$를 아벨 군으로 설정한다.
우리는 말한다$\mathcal{F}$은{\it 상수층가치있는$A$}만약에
$\mathcal{F}$은 전층 $U \mapsto A$의 층화이다.
표기법: $\underline{A}_X$ 또는 $\underline{A}$.
\item $\mathcal{F}$는 동형인 경우 {\it 상수층}라고 말한다.
(1)에서와 같이 아벨 층에서 층으로.
\item $\mathcal{F}$는 {\it 로컬 상수}라고 말한다.
피복 $\{U_i \to X\}$ 즉, $\mathcal{F}|_{U_i}$는 상수층이다.
\item $\mathcal{F}$는 {\it 유한 지역 상수}라고 말한다.
는 지역적으로 상수이고 상수 값은 유한 아벨 군이다.
\end{enumerate}
$\Lambda$를 환으로 설정한다. $\mathcal{F}$를 $\Lambda$-가군의 층으로 설정한다.
$X_\etale$에 있다.
\begin{enumerate}
\item $M$를 $\Lambda$-가군라고 가정한다.
우리는 말한다$\mathcal{F}$은{\it 상수층가치있는$M$}만약에
$\mathcal{F}$은 전층 $U \mapsto M$의 층화이다.
표기법: $\underline{M}_X$ 또는 $\underline{M}$.
\item $\mathcal{F}$는 동형인 경우 {\it 상수층}라고 말한다.
로서층~의$\Lambda$-가군에층에서와 같이 (1).
\item $\mathcal{F}$는 {\it 로컬 상수}라고 말한다.
피복 $\{U_i \to X\}$ 즉, $\mathcal{F}|_{U_i}$는 상수층이다.
\end{enumerate}
\end{definition}

\begin{lemma}
\label{etale-cohomology-lemma-pullback-locally-constant}
$f : X \to Y$를 스킴의 사상으로 설정한다. $\mathcal{G}$인 경우
국소상수층 또는 집합, 아벨 군 또는 $\Lambda$-가군
$Y_\etale$에서는 $f^{-1}\mathcal{G}$도 마찬가지이다.
$X_\etale$에 있다.
\end{lemma}

\begin{proof} 토포스의 사상에 대해 유지된다. 사이트, 보조정리 \ref{sites-modules-lemma-pullback-locally-constant}의 가군을 참조하라. \end{proof}

\begin{lemma} \label{etale-cohomology-lemma-pushforward-locally-constant} $f : X \to Y$를 스킴의 유한 에탈 사상라고 하자. $\mathcal{F}$가 집합의 (유한) 국소상수층인 경우, 아벨 군의 (유한) 국소상수층 또는 (유한 유형) 국소상수층의 $\Lambda$-가군 $X_\etale$, $f_*\mathcal{F}$의 $Y_\etale$도 마찬가지이다. \end{lemma}

\begin{proof} $f_*$의 구성은 에탈 국소화로 통근한다. 유한 에탈 사상은 동형사상의 분리된 합집합과 국소적으로 동형이다. 에탈 사상, 보조정리 \ref{etale-lemma-finite-etale-etale-local}를 참조하라. 따라서 보조정리는 $\mathcal{F}_i$, $i = 1, \ldots, n$가 집합의 (유한) 국소상수층인 경우 $\prod_{i = 1, \ldots, n} \mathcal{F}_i$도 마찬가지라고 말한다. 이것은 분명한다. 아벨 군 및 가군의 층도 마찬가지이다. \end{proof}

\begin{lemma}
\label{etale-cohomology-lemma-characterize-finite-locally-constant}
허락하다$X$가 되다스킴그리고$\mathcal{F}$에이층~의집합~에$X_\etale$.
그러면 다음은 동일한다.
\begin{enumerate}
\item $\mathcal{F}$는 유한한 국부적으로 상수이며,
\item $\mathcal{F} = h_U$ 일부 유한 에탈 사상 $U \to X$.
\end{enumerate}
\end{lemma}

\begin{proof}
유한 에탈 사상은 분리된 합집합과 국소적으로 동형이다.
동형사상의 내용을 참조하라.
에탈 사상, 보조정리 \ref{etale-lemma-finite-etale-etale-local}.
따라서 (2)는 (1)를 의미한다. 반대로, $\mathcal{F}$가 지역적으로 유한한 경우
상수이면 다음과 같은 에탈 피복 $\{X_i \to X\}$가 존재한다.
$\mathcal{F}|_{X_i}$는 $U_i \to X_i$ 유한 에탈로 표현 가능한다.
증명 에서와 똑같이 주장한다.
하강, 보조정리 \ref{descent-lemma-descent-data-sheaves}
우리는 스킴 $(U_i, \varphi_{ij})$에 대한 하강 데이터를 얻습니다.
$\{X_i \to X\}$ (자세한 내용은 생략) 이 하강 기준점은 다음에 효과적이다.
예 하강으로, 보조정리 \ref{descent-lemma-affine}
그리고 스킴 $U \to X$의 결과 사상은 유한 에탈이다.
강하 기준, 보조정리 \ref{descent-lemma-descending-property-finite} 및
\ref{descent-lemma-descending-property-etale}.
\end{proof}

\begin{lemma}
\label{etale-cohomology-lemma-morphism-locally-constant}
$X$를 스킴으로 설정한다.
\begin{enumerate}
\item $\varphi : \mathcal{F} \to \mathcal{G}$를 지도로 설정
국소상수층의 집합의 $X_\etale$.
$\mathcal{F}$가 국소적으로 유한한 상수인 경우
에탈 피복 $\{U_i \to X\}$
$\varphi|_{U_i}$는 다음과 연관된 상수 층의 맵이다.
집합의 지도.
\item $\varphi : \mathcal{F} \to \mathcal{G}$를 지도로 설정
국소상수층의 아벨 군의 $X_\etale$.
$\mathcal{F}$가 국소적으로 유한한 상수이면 에탈이 존재한다.
피복 $\{U_i \to X\}$ $\varphi|_{U_i}$는 다음의 맵이다.
아벨 군의 지도에 연결된 상수 아벨 층.
\item $\Lambda$를 환으로 둡니다.
$\varphi : \mathcal{F} \to \mathcal{G}$를 지도로 설정
~의국소상수층~의$\Lambda$-가군~에$X_\etale$.
$\mathcal{F}$가 유한 유형이면 에탈 피복이 존재한다.
$\{U_i \to X\}$ $\varphi|_{U_i}$는 상수의 맵이다.
층~의$\Lambda$-가군지도와 관련된$\Lambda$-가군.
\end{enumerate}
\end{lemma}

\begin{proof} 이는 모든 사이트에 적용된다. 사이트, 보조정리 \ref{sites-modules-lemma-morphism-locally-constant}의 가군을 참조하라. \end{proof}

\begin{lemma}
\label{etale-cohomology-lemma-kernel-finite-locally-constant}
$X$를 스킴으로 설정한다.
\begin{enumerate}
\item 범주의 유한 국소상수층 집합
$\Sh(X_\etale)$ 내부의 유한 극한 및 여극한 아래에서 닫혀 있다.
\item 유한 지역 상수 아벨 층의 범주는
$\textit{Ab}(X_\etale)$의 약한 Serre 하위 범주.
\item $\Lambda$를 뇌터 환으로 둡니다. 범주의
유한 유형, $\Lambda$-가군의 국소상수층
$X_\etale$은 약한 Serre 하위 범주이다.
$\textit{Mod}(X_\etale, \Lambda)$.
\end{enumerate}
\end{lemma}

\begin{proof} 이는 모든 사이트에 적용된다. 사이트, 보조정리 \ref{sites-modules-lemma-kernel-finite-locally-constant}의 가군을 참조하라. \end{proof}

\begin{lemma}
\label{etale-cohomology-lemma-tensor-product-locally-constant}
$X$를 스킴으로 설정한다. $\Lambda$를 환으로 설정한다.
$\Lambda$-가군의 두 국소상수층의 텐서곱
$X_\etale$에서 $\Lambda$-가군의 국소상수층이다.
\end{lemma}

\begin{proof} 이는 모든 사이트에 적용된다. 사이트, 보조정리 \ref{sites-modules-lemma-tensor-product-locally-constant}의 가군을 참조하라. \end{proof}

\begin{lemma}
\label{etale-cohomology-lemma-connected-locally-constant}
$X$를 연결 스킴으로 설정한다. $\Lambda$를 환으로 두고
$\mathcal{F}$는 $\Lambda$-가군의 국소상수층이다.
그런 다음 $\Lambda$-가군 $M$ 및 에탈 피복이 존재한다.
$\{U_i \to X\}$ $\mathcal{F}|_{U_i} \cong \underline{M}|_{U_i}$와 같습니다.
\end{lemma}

\begin{proof}
에탈을 선택하세요 피복
$\{U_i \to X\}$ $\mathcal{F}|_{U_i}$가 일정하다고 가정해 보자.
$\mathcal{F}|_{U_i} \cong \underline{M_i}_{U_i}$.
그것을 관찰하십시오$U_i \times_X U_j$비어 있는 경우$M_i$동형이 아니다
$M_j$로 변경된다.
각각의 $\Lambda$-가군 $M$에 대해 $I_M = \{i \in I \mid M_i \cong M\}$를 허용한다.
에탈 사상이 열려 있음을 알 수 있다.
$U_M = \bigcup_{i \in I_M} \Im(U_i \to X)$
$X$의 열린 부분집합이다. 그러면 $X = \coprod U_M$는 분리된
$X$의 피복을 엽니다. $X$는 연결이므로 $U_M$ 하나만 비어 있지 않습니다.
그리고 보조정리가 이어집니다.
\end{proof}






\section{국소상수층 및 기본군} \label{etale-cohomology-section-pione}

\noindent 어떤 경우에는 국소상수층을 스킴의 기본군과 연관시킬 수 있다.

\begin{lemma} \label{etale-cohomology-lemma-locally-constant-on-connected} $X$를 연결 스킴으로 둡니다. $\overline{x}$를 $X$의 기하점으로 설정한다. \begin{enumerate} \item 범주의 동치 $$
\left\{
\begin{matrix}
\text{유한 국소상수인}\\
\text{집합의 층으로서 }X_\etale\text{ 위의 것}
\end{matrix}
\right\}
\longleftrightarrow
\left\{
\begin{matrix}
\text{유한 }\pi_1(X, \overline{x})\text{-집합}
\end{matrix}
\right\}
$$ \item 범주의 동치 $$
\left\{
\begin{matrix}
\text{유한 국소상수인}\\
\text{아벨 군의 층으로서 }X_\etale\text{ 위의 것}
\end{matrix}
\right\}
\longleftrightarrow
\left\{
\begin{matrix}
\text{유한 }\pi_1(X, \overline{x})\text{-가군}
\end{matrix}
\right\}
$$ \item $\Lambda$가 있다. 유한 환. 범주의 동치 $$
\left\{
\begin{matrix}
\text{유한형 국소상수인}\\
\Lambda\text{-가군층으로서 }X_\etale\text{ 위의 것}
\end{matrix}
\right\}
\longleftrightarrow
\left\{
\begin{matrix}
\text{유한 }\pi_1(X, \overline{x})\text{-가군이고,}\\
\text{가환하는 }\Lambda\text{-가군 구조를 갖는 것}
\end{matrix}
\right\}
$$ \end{enumerate} \end{lemma}가 있다.

\begin{proof}
우리는 $\pi_1(X, \overline{x})$가 유한한 것을 관찰한다.
위상 그룹, 기본 그룹, 정의 참조
\ref{pione-definition-fundamental-group}.
왼손 범주는 다음과 같이 정의된다.
절 \ref{etale-cohomology-section-locally-constant}.
오른손에 사용된 표기법 범주는
기본 그룹, 정의 \ref{pione-definition-G-set-continuous}
집합 및
정의 \ref{etale-cohomology-definition-G-module-continuous} - 아벨 군.
이것은 표기법을 설명한다.

\medskip\noindent 주장(1)은 보조정리 \ref{etale-cohomology-lemma-characterize-finite-locally-constant} 및 기본 그룹, 정리 \ref{pione-theorem-fundamental-group}에서 따릅니다. 부분 (2) 및 (3)는 기본 (층 의) 집합에 추가 구조를 부여하여 이로부터 즉시 이어집니다. 예의 경우, $X_\etale$에 있는 아벨 군의 유한 국소상수층은 집합 $\mathcal{F}$의 유한 국소상수층과 맵 $+ : \mathcal{F} \times \mathcal{F} \to \mathcal{F}$를 만족하는 것과 동일한다. 일반적인 공리. (1)의 등가는 제품을 제품으로 보내고 따라서 $+$을 해당 유한 $\pi_1(X, \overline{x})$-집합의 추가 항목으로 보냅니다. $\pi_1(X, \overline{x})$-가군은 호환되는 아벨 군 구조를 갖는 $\pi_1(X, \overline{x})$-집합과 동일하므로 (2)를 얻습니다. 부분(3)도 똑같은 방식으로 증명된다. \end{proof}

\begin{lemma} \label{etale-cohomology-lemma-locally-constant-on-connected-geom-unibranch} $X$를 기약, 기하학적으로 단일 가지 스킴라고 가정한다. $\overline{x}$를 $X$의 기하점으로 설정한다. $\Lambda$를 환으로 설정한다. 범주의 동치 $$
\left\{
\begin{matrix}
\text{유한형 국소상수인}\\
\Lambda\text{-가군층으로서 }X_\etale\text{ 위의 것}
\end{matrix}
\right\}
\longleftrightarrow
\left\{
\begin{matrix}
\text{유한 생성 }\Lambda\text{-가군 }M\text{이고 연속적인}\\
\pi_1(X, \overline{x})\text{-작용을 갖는 것}
\end{matrix}
\right\}
$$ \end{lemma}가 있다.

\begin{proof} 보조정리 \ref{etale-cohomology-lemma-locally-constant-on-connected}에 제공된 증명은 유한 $\Lambda$-가군 $M$가 유한 기본 집합으로 작동하지 않을 수 있다.

\medskip\noindent $\nu : X^\nu \to X$를 정규화 사상으로 둡니다. 사상, 보조정리 \ref{morphisms-lemma-normalization-geom-unibranch-univ-homeo}에 의해 이것은 보편적인 동형이다. 기본 그룹에 따라 명제 \ref{pione-proposition-universal-homeomorphism} 이는 동형사상 $\pi_1(X^\nu, \overline{x}) \to \pi_1(X, \overline{x})$를 유도하고 정리 \ref{etale-cohomology-theorem-topological-invariance}에 의해 유한 유형, 로컬 상수 사이에서 범주의 동등성을 얻습니다. $\Lambda$-가군 $X_\etale$ 및 $X^\nu_\etale$에 있다. 이는 $X$가 정역 정규 스킴인 경우로 축소된다.

\medskip\noindent
$X$는 정역 정규 스킴라고 가정한다. $\eta \in X$를 일반적인 지점으로 둡니다.
$\overline{\eta}$는 $\eta$ 위에 놓여 있는 기하점라고 가정한다.
기본 그룹별, 명제 \ref{pione-proposition-normal}
계속해서 추측을 하다
$$
\text{Gal}(\kappa(\eta)^{sep}/\kappa(\eta)) = \pi_1(\eta, \overline{\eta})
\longrightarrow
\pi_1(X, \overline{\eta})
$$
핵은 기본 그룹, 보조정리에 설명되어 있다.
\ref{pione-lemma-normal-pione-quotient-inertia}.
$\mathcal{F}$를 유한 유형으로 두고, $\Lambda$-가군의 국소상수층
$X_\etale$에 있다. $M = \mathcal{F}_{\overline{\eta}}$
될줄기~의$\mathcal{F}$~에$\overline{\eta}$.
$\text{Gal}(\kappa(\eta)^{sep}/\kappa(\eta))$의 연속 동작을 얻습니다.
$M$에 의해 절 \ref{etale-cohomology-section-galois-action-stalks}에 의해.
우리의 목표는 이러한 행동이 다음을 통해 영향을 미친다는 것을 보여주는 것이다.
표시되는 추측.
$\mathcal{F}$는 유한 유형이므로 $M$는 유한 $\Lambda$-가군이다.
$\mathcal{F}$는 지역적으로 일정하므로 모든 $x \in X$에 대해
$\mathcal{F}$를 $\Spec(\mathcal{O}_{X, x}^{sh})$로 제한
일정한다. 따라서 $\text{Gal}(K^{sep}/K_x^{sh})$의 동작은 다음과 같습니다.
(기본 그룹과 같은 표기법 사용, 보조정리
\ref{pione-lemma-normal-pione-quotient-inertia})
$M$에서는 사소한 일이다. 우리는 원하는 대로 인수분해를 했다고 결론을 내렸습니다.

\medskip\noindent
반면에 유한 $\Lambda$-가군 $M$
$\pi_1(X, \overline{\eta})$의 연속 동작으로.
우리는 $\mathcal{F}$를 다음과 같이 구성할 것이다.
$M \cong \mathcal{F}_{\overline{\eta}}$
$\Lambda[\pi_1(X, \overline{\eta})]$-가군.
생성기 $m_1, \ldots, m_r \in M$를 선택한다.
의 행동 이후$\pi_1(X, \overline{\eta})$~에$M$~이다
연속, 각 $i$에 대해 열린 하위 그룹이 존재한다.
$N_i$ 유한 그룹 $\pi_1(X, \overline{\eta})$의
모든 $\gamma \in H_i$는 $m_i$을 수정한다. 우리는 다음과 같이 결론을 내립니다.
열린 하위 그룹 $H = \bigcap_{i = 1, \ldots, r} H_i$의 모든 요소
$M$의 모든 요소를 ​​수정한다. $H$를 축소한 후 $H$라고 가정할 수 있다.
$\pi_1(X, \overline{\eta})$의 열린 정규 하위 그룹이다.
집합 $G = \pi_1(X, \overline{\eta})/H$. $f : Y \to X$를
상응하는 갈루아 유한 에탈 $G$-피복.
$f_*\underline{\mathbf{Z}}$을 층으로 볼 수 있다.
$\mathbf{Z}[G]$-가군 $X_\etale$에 있다. 그럼 그냥 가져가자
$$
\mathcal{F} =
f_*\underline{\mathbf{Z}} \otimes_{\underline{\mathbf{Z}[G]}} \underline{M}
$$
계산은 독자에게 맡깁니다.
$\mathcal{F}_{\overline{\eta}}$.
우리는 또한 이 건설에 대한 검증을 생략한다.
이전 단락의 구성과 반대이다.
\end{proof}

\begin{remark}
\label{etale-cohomology-remark-functorial-locally-constant-on-connected}
기본정리 \ref{etale-cohomology-lemma-locally-constant-on-connected}와
\ref{etale-cohomology-lemma-locally-constant-on-connected-geom-unibranch}
당김과 호환된다. 예에 대해 $f : Y \to X$라고 가정한다.
연결 스킴의 사상이다. $\overline{y}$를 기하학적이라고 하자
$Y$ 및 집합 $\overline{x} = f(\overline{y})$의 지점이다.
그런 다음 다이어그램
$$
\xymatrix{
\text{집합의 유한 국소상수층으로서 }Y_\etale\text{ 위의 것}
\ar[r] &
\text{유한 }\pi_1(Y, \overline{y})\text{-집합} \\
\text{집합의 유한 국소상수층으로서 }X_\etale\text{ 위의 것}
\ar[r] \ar[u]_{f^{-1}} &
\text{유한 }\pi_1(X, \overline{x})\text{-집합} \ar[u]
}
$$
가환적이며 오른쪽의 수직 화살표가 나타납니다.
연속 동형으로부터
$\pi_1(Y, \overline{y}) \to \pi_1(X, \overline{x})$
$f$에 의해 유도되었다. 이는 다음에서 바로 이어집니다.
가환 도식
기본 그룹, 정리 \ref{pione-theorem-fundamental-group}.
다른 경우에도 비슷한 결과가 나타납니다.
\end{remark}










\section{추적 방법} \label{etale-cohomology-section-trace-method}

\noindent 이 절에 대한 참조는 \cite[Expos\'e IX, \S 5]{SGA4}이다. 이곳의 재료는 아래 보조정리 \ref{etale-cohomology-lemma-vanishing-easier}의 증명에 사용된다.

\medskip\noindent
$f : Y \to X$를 스킴의 에탈 사상이라고 하자. 거기
시퀀스이다
$$
f_!, f^{-1}, f_*
$$
수반 함자 사이
$\textit{Ab}(X_\etale)$ 및 $\textit{Ab}(Y_\etale)$.
함자 $f_!$는 절 \ref{etale-cohomology-section-extension-by-zero}에서 논의된다.
부속 맵 $\text{id} \to f_* f^{-1}$은 {\it 제한}이라고 한다.
부속 맵 $f_! f^{-1} \to \text{id}$은 종종
{\it 트레이스 맵}라고 한다. $f$가 유한 에탈이면 $f_* = f_!$
(보조정리 \ref{etale-cohomology-lemma-shriek-equals-star-finite-etale}) 및
이것을 지도 $f_*f^{-1} \to \text{id}$로 볼 수 있다.

\begin{definition} \label{etale-cohomology-definition-trace-map} $f : Y \to X$를 스킴의 유한 에탈 사상라고 하자. 위와 명시적으로 아래에 설명된 지도 $f_* f^{-1} \to \text{id}$를 {\it 추적}이라고 한다. \end{definition}

\noindent
$f : Y \to X$를 스킴의 유한 에탈 사상라고 하자. 트레이스 맵은
다음 두 가지 성질이 특징이다.
\begin{enumerate}
\item $X$에 에탈 국소화를 타고 출퇴근하고
\item 만약 $Y = \coprod_{i = 1}^d X$이면 트레이스 맵은 다음과 같습니다.
합계 맵 $f_*f^{-1} \mathcal{F} = \mathcal{F}^{\oplus d} \to \mathcal{F}$.
\end{enumerate}
에탈 사상, 보조정리 \ref{etale-lemma-finite-etale-etale-local}
모든 유한 에탈 사상 $f : Y \to X$는 $X$에서 로컬로 에탈이다.
(에 주어진 형태의2) 일부 정수의 경우$d \geq 0$. 그러므로 우리는
주어진 특성화를 사용하여 추적 맵을 정의할 수 있다. 특히
우리는 $f_!$의 존재와 합의에 대해 알 필요가 없다.
트레이스 맵을 구성하려면 $f_!$을 $f_*$로 사용하세요.
이 설명은 $f$가 일정한 차수 $d$를 갖는 경우 다음을 보여줍니다.
구성
$$
\mathcal{F} \xrightarrow{res}
f_* f^{-1} \mathcal{F} \xrightarrow{trace}
\mathcal{F}
$$
는 $d$의 곱이다. ``m\'추적 방법''
다음 관찰은 다음과 같습니다: 만약 $\mathcal{F}$
$X_\etale$의 아벨 층이므로 $d$을 곱하면 된다.
$\mathcal{F}$에서 동형사상이면 지도가
$$
H^n_\etale(X, \mathcal{F}) \longrightarrow H^n_\etale(Y, f^{-1}\mathcal{F})
$$
주입식이다. 즉, 우리는
$$
H^n_\etale(Y, f^{-1}\mathcal{F}) = H^n_\etale(X, f_*f^{-1}\mathcal{F})
$$
고차 직상의 소멸에 의해
(명제 \ref{etale-cohomology-proposition-finite-higher-direct-image-zero})
그리고 르레이 스펙트럼 열
(명제 \ref{etale-cohomology-proposition-leray}).
따라서 우리는 지도를 고려할 수 있다
$$
H^n_\etale(X, \mathcal{F}) \to
H^n_\etale(Y, f^{-1}\mathcal{F})= H^n_\etale(X, f_*f^{-1}\mathcal{F})
\xrightarrow{trace}
H^n_\etale(X, \mathcal{F})
$$
구성은 동형사상 (가정 위의 $\mathcal{F}$
및 $f$). 특히, 만약
$H_\etale^q(Y, f^{-1}\mathcal{F}) = 0$ 그러면
$H_\etale^q(X, \mathcal{F}) = 0$도 그렇습니다.
실제로 $d$를 곱하면
동형사상 $H_\etale^q(X, \mathcal{F})$에 대해 다음을 고려한다.
$H_\etale^q(Y, f^{-1}\mathcal{F})= 0$.

\medskip\noindent 이것은 종종 다음과 결합된다.

\begin{lemma}
\label{etale-cohomology-lemma-pullback-filtered}
$S$를 연결 스킴으로 설정한다. $\ell$를 소수라고 하자. 허락하다
$\mathcal{F}$ 유한 유형이어야 하며, 국소상수층 의
$\mathbf{F}_\ell$-벡터공간 $S_\etale$에 있다.
그러면 유한한 에탈 사상이 존재한다.
$f : T \to S$는 $\ell$에 소수인 $f^{-1}\mathcal{F}$와 같습니다.
연속된 몫이 다음과 같은 유한한 필터링을 가집니다.
$\underline{\mathbf{Z}/\ell\mathbf{Z}}_T$.
\end{lemma}

\begin{proof}
선택하세요기하점 $\overline{s}$~의$S$.
보조정리 \ref{etale-cohomology-lemma-locally-constant-on-connected}의 동등성을 통해
층 $\mathcal{F}$는 유한 차원에 해당한다.
$\mathbf{F}_\ell$-벡터공간 $V$ 연속
$\pi_1(S, \overline{s})$-액션.
$G \subset \text{Aut}(V)$를 동형의 이미지로 둡니다.
$\rho : \pi_1(S, \overline{s}) \to \text{Aut}(V)$ 동작을 제공한다.
$G$가 유한하다는 점에 유의하세요.
전사 연속 동형
$\overline{\rho} : \pi_1(S, \overline{s}) \to G$
Galois 객체 $Y \to S$에 해당한다.
$\textit{F\'Et}_S$ 자동형성 그룹 $G = \text{Aut}(Y/S)$ 포함, 참조
기본 그룹, 절 \ref{pione-section-finite-etale-under-galois}.
$H \subset G$를 $\ell$-Sylow 하위 그룹으로 설정한다.
우리는 주장 $T = Y/H \to S$ 작동한다. 즉, $\overline{t} \in T$
$\overline{s}$보다 기하점이어야 한다. 이미지
$\pi_1(T, \overline{t}) \to \pi_1(S, \overline{s})$
functorial에서 다음과 같이 $(\overline{\rho})^{-1}(H)$이다.
기본 그룹의 성격. 따라서 $\pi_1(T, \overline{t})$의 동작은 다음과 같습니다.
~에$V$에 해당하는$f^{-1}\mathcal{F}$통해
지도 $\pi_1(T, \overline{t}) \to H$, 참조
비고 \ref{etale-cohomology-remark-functorial-locally-constant-on-connected}. 처럼
$H$는 유한한 $\ell$ 그룹이며, 기약 구성 요소는
표현 $\rho|_{\pi_1(T, \overline{t})}$
각각의 순위는 $1$입니다(이것은 간단한 보조정리이다.
유한 그룹의 표현 이론; 여기에 향후 참조를 삽입하십시오).
동등성을 통해
보조정리 \ref{etale-cohomology-lemma-locally-constant-on-connected}
이는 $f^{-1}\mathcal{F}$이 다음의 연속적인 확장임을 의미한다.
값 $\underline{\mathbf{Z}/\ell\mathbf{Z}}_T$의 상수 층.
또한 $T = Y/H \to S$의 차수는 $\ell$의 소수이다.
지수와 같으니까.$H$~에$G$.
\end{proof}

\begin{lemma}
\label{etale-cohomology-lemma-action-l-group}
$\Lambda$를 뇌터 환으로 설정한다.
$\ell$를 소수, $n \geq 1$로 설정한다.
$H$를 유한 $\ell$-그룹으로 둡니다.
$M$를 $\Lambda[H]$-가군이 $\ell^n$에 의해 소멸되는 유한한 것으로 가정한다.
그런 다음 유한한 필터링이 있다.
$0 = M_0 \subset M_1 \subset \ldots \subset M_t = M$
$\Lambda[H]$-하위 모듈
$H$는 모든 사용자에 대해 $M_{i + 1}/M_i$에 대해 사소하게 작동한다.
$i = 0, \ldots, t - 1$.
\end{lemma}

\begin{proof}
생략. 힌트: 증강 아이디얼 $\mathfrak m$
비가환적 환 $\mathbf{Z}/\ell^n\mathbf{Z}[H]$
무능하다.
\end{proof}

\begin{lemma}
\label{etale-cohomology-lemma-pullback-filtered-modules}
$S$를 기약, 기하학적으로 단일 분기 스킴라고 가정한다.
$\ell$를 소수, $n \geq 1$로 설정한다. $\Lambda$
뇌터 환이어야 한다. $\mathcal{F}$를 유한 유형으로 두고,
국소상수층~의$\Lambda$-가군~에$S_\etale$
$\ell^n$에 의해 소멸된다. 그런 다음 존재한다.
유한 에탈 사상 $f : T \to S$ 소수 ~ $\ell$
$f^{-1}\mathcal{F}$에는 유한한 필터링이 있다.
연속된 몫은 $\underline{M}_T$ 형식이다.
일부 유한 $\Lambda$-가군 $M$의 경우.
\end{lemma}

\begin{proof}
선택하세요기하점 $\overline{s}$~의$S$.
동등성을 통해
보조정리 \ref{etale-cohomology-lemma-locally-constant-on-connected-geom-unibranch}
층 $\mathcal{F}$는 유한 $\Lambda$-가군에 해당한다.
$M$ 연속적인 $\pi_1(S, \overline{s})$ 동작을 사용한다.
$G \subset \text{Aut}(V)$를 동형의 이미지로 둡니다.
$\rho : \pi_1(S, \overline{s}) \to \text{Aut}(M)$ 동작을 제공한다.
$G$는 유한한 $M$는 유한한 $\Lambda$-가군이므로 유한한다.
(보조정리 \ref{etale-cohomology-lemma-locally-constant-on-connected-geom-unibranch}의 증명 참조).
전사 연속 동형
$\overline{\rho} : \pi_1(S, \overline{s}) \to G$
Galois 객체 $Y \to S$에 해당한다.
$\textit{F\'Et}_S$ 자동형성 그룹 $G = \text{Aut}(Y/S)$ 포함, 참조
기본 그룹, 절 \ref{pione-section-finite-etale-under-galois}.
$H \subset G$를 $\ell$-Sylow 하위 그룹으로 설정한다.
우리는 주장 $T = Y/H \to S$ 작동한다. 즉, $\overline{t} \in T$
$\overline{s}$보다 기하점이어야 한다. 이미지
$\pi_1(T, \overline{t}) \to \pi_1(S, \overline{s})$
functorial에서 다음과 같이 $(\overline{\rho})^{-1}(H)$이다.
기본 그룹의 성격. 따라서 $\pi_1(T, \overline{t})$의 동작은 다음과 같습니다.
~에$M$에 해당하는$f^{-1}\mathcal{F}$통해
지도 $\pi_1(T, \overline{t}) \to H$, 참조
비고 \ref{etale-cohomology-remark-functorial-locally-constant-on-connected}.
$0 = M_0 \subset M_1 \subset \ldots \subset M_t = M$
보조정리 \ref{etale-cohomology-lemma-action-l-group}와 같습니다.
이는 여과를 유도한다.
$0 = \mathcal{F}_0 \subset \mathcal{F}_1 \subset \ldots
\subset \mathcal{F}_t = f^{-1}\mathcal{F}$
연속된 몫은 다음과 같이 일정한다.
값 $M_{i + 1}/M_i$.
마지막으로, $T = Y/H \to S$의 차수는 $\ell$의 소수이다.
지수와 같으니까.$H$~에$G$.
\end{proof}






\section{Galois 코호몰로지}
\label{etale-cohomology-section-galois-cohomology}

\noindent
이 절에서는 Galois 코호몰로지에 대한 결과를 증명한다.
(명제 \ref{etale-cohomology-proposition-serre-galois})
에탈 코호몰로지 과 트릭을 사용하여
절 \ref{etale-cohomology-section-trace-method}.
이를 통해 더 높은 에탈 코호몰로지 군의 소멸을 증명할 수 있다.
체의 스펙트럼에 걸쳐 있다.

\begin{lemma} \label{etale-cohomology-lemma-nonvanishing-inherited} $\ell$는 소수이고 $n$는 정수 $> 0$라고 가정한다. $S$를 준컴팩트하고 준분리된 스킴라고 하자. $X = \lim_{i \in I} X_i$를 $S$-스킴의 방향 시스템의 극한라고 가정한다. 각 $X_i \to S$는 $\ell$에 상대적으로 소수인 일정한 차수의 유한 에탈이다. 다음은 동일합니다: \begin{enumerate} \item $\ell$-거듭제곱 꼬임층 $\mathcal{G}$ 위의 $S$ 존재하므로 $H_\etale^n(S, \mathcal{G}) \neq 0$ 및 \item 존재 $\ell$-전원 꼬임층 $\mathcal{F}$를 $X$에서 $H_\etale^n(X, \mathcal{F}) \neq 0$로 전환한다. \end{enumerate} 사실, $\mathcal{G}$가 주어지면 $\mathcal{F} = g^{-1}\mathcal{F}$를 취할 수 있고, $\mathcal{F}$가 주어지면 $\mathcal{G} = g_*\mathcal{F}$를 취할 수 있다. \end{lemma}

\begin{proof}
$g : X \to S$ 및 $g_i : X_i \to S$는 구조 사상을 나타냅니다.
$\ell$-전원 꼬임층 $\mathcal{G}$ $S$ 수정
$H^n_\etale(S, \mathcal{G}) \not = 0$로.
$\mathcal{G}_i = g_i^{-1}\mathcal{G}$에 의해 주어진 시스템
정리 \ref{etale-cohomology-theorem-colimit}의 조건을 만족함
$g^{-1}\mathcal{G}$에 의해 주어진 여극한 층으로. 이것은 말한다
우리는:
$$
\colim_{i\in I} H^n_\etale(X_i, g_i^{-1}\mathcal{G}) =
H^n_\etale(X, \mathcal{G})
$$
$g_i$가 소수인 에탈 사상이기 때문에
$\ell$에 ``la m\'ethode de la Trace''를 적용하면 다음을 찾을 수 있다.
지도
$$
H^n_\etale(S, \mathcal{G}) \to H^n_\etale(X_i, g_i^{-1}\mathcal{G})
$$
모두 주입적입니다(그리고 전환 맵과 호환됩니다).
절 \ref{etale-cohomology-section-trace-method}를 참조하라. 따라서 여극한은 0이 아닙니다. 즉,
$H^n(X,g^{-1}\mathcal{G}) \neq 0$, 다음과 같이 원하는 결과를 제공한다.
$\mathcal{F} = g^{-1}\mathcal{G}$.

\medskip\noindent
반대로, 주어진다고 가정하자.$\ell$-힘꼬임층 $\mathcal{F}$~에$X$
$H^n_\etale(X, \mathcal{F}) \not = 0$로. $g_i$ 이후
유한하다 사상 고차 직상 사라지다
(명제 \ref{etale-cohomology-proposition-finite-higher-direct-image-zero}).
그런 다음 보조정리 \ref{etale-cohomology-lemma-relative-colimit}를 적용하여
$g$에 대해서도 동일한 결론을 내릴 수 있다.
고차 직상의 소멸은 다음을 알려줍니다.
$H^n_\etale(X, \mathcal{F}) = H^n(S, g_*\mathcal{F}) \neq 0$
작성자: Leray (명제 \ref{etale-cohomology-proposition-leray})
$\mathcal{G} = g_*\mathcal{F}$로 우리가 원하는 것을 제공한다.
\end{proof}

\begin{lemma}
\label{etale-cohomology-lemma-reduce-to-l-group}
$\ell$는 소수이고 $n$는 정수 $> 0$라고 가정한다.
$K$를 $G = Gal(K^{sep}/K)$를 사용하여 체로 만들고
$H \subset G$최대한의 프로가 되세요$\ell$하위 그룹$L/K$
해당 체 확장자가 된다. 그 다음에
$H^n_\etale(\Spec(K), \mathcal{F}) = 0$ 모두에 대해
$\ell$-동력 비틀림 $\mathcal{F}$ 필요충분조건은
$H^n_\etale(\Spec(L), \underline{\mathbf{Z}/\ell\mathbf{Z}}) = 0$.
\end{lemma}

\begin{proof}
$L = \bigcup L_i$을 $K$에 대한 유한 하위 확장의 합집합으로 작성한다.
$H$을 선택하면 $[L_i : K]$가 $\ell$의 소수라는 의미이다.
따라서 $\Spec(L) = \lim_{i \in I} \Spec(L_i)$ 와 같이
보조정리 \ref{etale-cohomology-lemma-nonvanishing-inherited}.
따라서 $K$을 $L$로 대체하고 다음과 같이 가정할 수 있다.
절대 갈루아 그룹$G$~의$K$는
한정된 프로-$\ell$ 그룹.

\medskip\noindent
$H^n(\Spec(K), \underline{\mathbf{Z}/\ell\mathbf{Z}}) = 0$를 가정한다.
$\mathcal{F}$를 $\ell$-전력 꼬임층으로 $\Spec(K)_\etale$로 설정한다.
$H^n_\etale(\Spec(K), \mathcal{F}) = 0$을 보여드리겠습니다.
다음에 명시된 서신을 통해
보조정리 \ref{etale-cohomology-lemma-equivalence-abelian-sheaves-point}
우리의 층 $\mathcal{F}$는 $\ell$-파워 비틀림에 해당한다.
$G$-가군 $M$. 요소 $x_1, \ldots, x_m \in M$의 유한한 집합
연속성을 통해 열린 하위 그룹 $U$에 의해 수정되어야 한다.
$M'$를 $x_1, \ldots, x_m$의 궤도에 걸쳐 있는 가군으로 설정한다.
이것은 유한 아벨 $\ell$-그룹이다.
각각의 $x_i$은 $\ell$의 힘에 의해 살해된다.
그리고 궤도는 유한한다. $M$는 여과 여극한이므로
이러한 하위 모듈 $M'$, $\mathcal{F}$이 필터링된 것을 알 수 있다.
해당 서브시브 $\mathcal{F}' \subset \mathcal{F}$의 여극한.
이 여극한에 정리 \ref{etale-cohomology-theorem-colimit}를 적용하면
$\mathcal{F}$이 유한한 국소상수층인 경우.

\medskip\noindent
$M$를 연속 동작을 갖는 유한 아벨 $\ell$ 그룹으로 설정한다.
유한한 프로-$\ell$ 그룹 $G$의. 그런 다음 $G$-불변이 있다.
여과법
$$
0 = M_0 \subset M_1 \subset \ldots \subset M_r = M
$$
$M_{i + 1}/M_i \cong \mathbf{Z}/\ell \mathbf{Z}$와 같이
사소한 $G$-동작(이것은 표현에 대한 간단한 보조정리이다.
유한 그룹 이론; 여기에 향후 참조를 삽입하십시오).
따라서 해당 층 $\mathcal{F}$에는 필터링이 있다.
$$
0 = \mathcal{F}_0 \subset \mathcal{F}_1 \subset \ldots \subset
\mathcal{F}_r = \mathcal{F}
$$
다음과 같은 연속적인 몫을 갖는
$\underline{\mathbf{Z}/\ell \mathbf{Z}}$.
따라서 귀납법과 긴 정확한 코호몰로지에 의해
우리는 결론을 내립니다.
\end{proof}

\begin{lemma}
\label{etale-cohomology-lemma-reduce-to-l-group-higher}
$\ell$는 소수이고 $n$는 정수 $> 0$라고 가정한다.
$K$를 $G = Gal(K^{sep}/K)$를 사용하여 체로 만들고
$H \subset G$ 최대 친$\ell$ 하위 그룹이 된다.
$L/K$는 해당 체 확장자이다.
그 다음에$H^q_\etale(\Spec(K),\mathcal{F}) = 0$~을 위한$q \geq n$그리고 모두
$\ell$-꼬임층 $\mathcal{F}$ 다음과 같은 경우에만
$H^n_\etale(\Spec(L), \underline{\mathbf{Z}/\ell\mathbf{Z}}) = 0$.
\end{lemma}

\begin{proof}
순방향은 간단하므로 역방향만 증명하면 된다.
$q$에 대한 유도를 진행한다. $q = n$의 경우는
보조정리 \ref{etale-cohomology-lemma-reduce-to-l-group}. 이제 보자
$\mathcal{F}$가 되다$\ell$-힘꼬임층~에$\Spec(K)$.
$f : \Spec(K^{sep}) \rightarrow \Spec(K)$
기하점을 포함해야 한다.
그런 다음 완전열을 고려하라.
$$
0 \rightarrow \mathcal{F} \xrightarrow{res}
f_* f^{-1} \mathcal{F} \rightarrow f_* f^{-1} \mathcal{F}/\mathcal{F} 
\rightarrow 0
$$
$K^{sep}$는 유한의 여과 여극한으로 쓸 수 있다.
분리 가능한 확장자. 따라서 $f$
유한 에탈 사상의 방향성 시스템의 극한이다.
증명에서 볼 수 있듯이
보조정리 \ref{etale-cohomology-lemma-nonvanishing-inherited}, $f$에 다음이 있다고 결론 내립니다.
소멸 고차 직상. 따라서 우리는 더 높은 코호몰로지를 표현할 수 있다.
$f_* f^{-1} \mathcal{F}$는 기하점의 더 높은 코호몰로지이다.
분명히 사라집니다. 따라서 여기에 있는 모든 것이 여전히 $\ell$-비틀림이므로 다음을 사용할 수 있다.
길고 정확한 코호몰로지 시퀀스와 관련된 귀납적 가설
$q + 1$에 대한 결과를 결론짓습니다.
\end{proof}

\begin{proposition}
\label{etale-cohomology-proposition-serre-galois}
\begin{reference}
\cite[Chapter II, Section 3, Proposition 5]{SerreGaloisCohomology}
\end{reference}
$K$는 분리 가능한 대수적 종결부 $K^{sep}$를 갖는 체라고 가정한다.
유한 확장에 대해 다음을 가정한다.$K'$~의$K$우리는
$\text{Br}(K') = 0$. 그 다음에
\begin{enumerate}
\item $H^q(\text{Gal}(K^{sep}/K), (K^{sep})^*) = 0$
모든 $q \geq 1$에 대해
\item $H^q(\text{Gal}(K^{sep}/K), M) = 0$
비틀림 $\text{Gal}(K^{sep}/K)$-가군 $M$ 및 $q \geq 2$에 대해
\end{enumerate}
\end{proposition}

\begin{proof}
두자 $p = \text{char}(K)$.
보조정리 \ref{etale-cohomology-lemma-compare-cohomology-point}에 의해,
정리 \ref{etale-cohomology-theorem-brauer-delta},
그리고 예 \ref{etale-cohomology-example-sheaves-point}
명제는 다음을 보여주는 것과 동일한다.
$H^2(\Spec(K'),\mathbf{G}_m|_{\Spec(K')_\etale}) = 0$
모든 유한 확장 $K'/K$에 대해 다음을 수행한다.
\begin{itemize}
\item $H^q(\Spec(K),\mathbf{G}_m|_{\Spec(K)_\etale}) = 0$
모든 $q \geq 1$에 대해
\item $H^q(\Spec(K),\mathcal{F}) = 0$
모든 꼬임층 $\mathcal{F}$ 및 $q \geq 2$에 대해.
\end{itemize}
우리는 두 번째 부분을 먼저 증명한다.
$\mathcal{F}$는 꼬임층이므로 다음을 사용할 수 있다.
$\ell$-1차 분해 및 호환성
코호몰로지 여극한 (즉, 직합, 정리 \ref{etale-cohomology-theorem-colimit} 참조)
$H^q(\Spec(K),\mathcal{F}) = 0$, $q \geq 2$ 표시로 줄이다.
모두를 위해$\ell$-힘꼬임층모든 소수에 대해$\ell$. 이것
각 소수를 개별적으로 분석할 수 있다.

\medskip\noindent
$\ell \neq p$라고 가정한다. 확장자 $K'/K$
Kummer 수열을 고려하십시오 (보조정리 \ref{etale-cohomology-lemma-kummer-sequence})
$$
0 \to
\mu_{\ell, \Spec{K'}} \to
\mathbf{G}_{m, \Spec{K'}} \xrightarrow{(\cdot)^{\ell}}
\mathbf{G}_{m, \Spec{K'}} \to 0
$$
이후 $H^q(\Spec{K'},\mathbf{G}_m|_{\Spec(K')_\etale}) = 0$
$q = 2$의 경우 가정 및 $q = 1$의 경우
정리 \ref{etale-cohomology-theorem-picard-group}와 결합됨
$\Pic(K) = (0)$. 따라서, 길고 정확한 코호몰로지
시퀀스에서 $H^2(\Spec{K'}, \mu_\ell) = 0$에 대해 결론을 내릴 수 있다.
분리 가능한 모든 $K'/K$. 이제 $H$를 최대 프로-$\ell$ 하위 그룹으로 지정한다.
$K$의 절대 갈루아 그룹에 속하며 $L$를
해당 확장자. 유한 확장의 여극한으로 $L$를 쓸 수 있다.
정리 \ref{etale-cohomology-theorem-colimit}를 이 여극한에 적용하면
$H^2(\Spec(L), \mu_\ell) = 0$.
이제 $\mu_\ell$는 상수층여야 한다. 그렇지 않은 경우에는 다음을 의미한다.
상대적으로 소인 정도의 갈루아 확장이 존재한다.
$\ell$ 의 $L$ 이는 정의 의 $L$에 의해 참이 아닙니다(즉, 확장자는
$\ell$번째 통일근을 $L$)에 연결함으로써 얻을 수 있다.
따라서 보조정리 \ref{etale-cohomology-lemma-reduce-to-l-group-higher}를 통해,
$\ell \neq p$에 대한 결과를 결론짓습니다.

\medskip\noindent
이제 $\ell = p$라고 가정해 보자. 우리는 다음을 고려한다.
아르틴-슈리어 완전열 (절 \ref{etale-cohomology-section-artin-schreier})
$$
0 \longrightarrow \underline{\mathbf{Z}/p\mathbf{Z}}_{\Spec{K}} \longrightarrow
\mathbf{G}_{a, \Spec{K}} \xrightarrow{F-1} \mathbf{G}_{a, \Spec{K}} 
\longrightarrow 0
$$
여기서 $F - 1$는 지도 $x \mapsto x^p - x$이다. 그렇다면 더 높다는 점에 유의하세요.
$\mathbf{G}_{a, \Spec{K}}$ 중 코호몰로지가 사라집니다.
비고 \ref{etale-cohomology-remark-special-case-fpqc-cohomology-quasi-coherent} 그리고
소멸 코호몰로지의 구조층의 아핀 스킴
(코호몰로지 중 스킴, 보조정리
\ref{coherent-lemma-quasi-coherent-affine-cohomology-zero}).
이는 다음의 모든 체에 적용될 수 있다.
표수 $p$. 특히, 체 확장자 $L$에 적용할 수 있다.
최대 프로-$p$ 하위 그룹 $H$에 의해 정의된다. 이를 통해 우리는 결론을 내릴 수 있다.
$H^n(\Spec{L}, \underline{\mathbf{Z}/p\mathbf{Z}}_{\Spec{L}}) = 0$
$n \geq 2$에 대해 $\ell = p$에 대한 결과는 다음과 같습니다.
보조정리 \ref{etale-cohomology-lemma-reduce-to-l-group-higher}.

\medskip\noindent
증명을 끝내려면 여전히 다음을 보여줘야 한다.
$H^q(\text{Gal}(K^{sep}/K), (K^{sep})^*) = 0$모두를 위해$q \geq 1$.
집합 $G = \text{Gal}(K^{sep}/K)$ 및 집합 $M = (K^{sep})^*$
$G$-가군으로 간주된다. 우리는 이미 (위에서) 그것을 보여주었습니다.
$H^1(G, M) = 0$ 및 $H^2(G, M) = 0$. 완전열을 고려해보세요.
$$
0 \to A \to M \to M \otimes \mathbf{Q} \to B \to 0
$$
$G$-가군. 위에서 우리는 $H^i(G, A) = 0$
그리고$H^i(G, B) = 0$~을 위한$i > 1$~부터$A$그리고$B$~이다
비틀림 $G$-가군. 에 의해
보조정리 \ref{etale-cohomology-lemma-profinite-group-cohomology-torsion}
우리는$H^i(G, M \otimes \mathbf{Q}) = 0$~을 위한$i > 0$.
이것이 다음을 의미한다는 것을 보는 것은 즐거운 운동이다.
$H^i(G, M) = 0$는 $i \geq 3$에도 적용된다.
\end{proof}

%10.08.09

\begin{definition}
\label{etale-cohomology-definition-Cr}
 체 $K$는 {\it $C_r$}라고 한다.
모든 $0 < d^r < n$과 모든 동차다항식 $f \in K[T_1, \ldots, T_n]$ 중 차수가 $d$인 것에 대하여, $\alpha = (\alpha_1, \ldots, \alpha_n)$가 존재하여 $\alpha_i \in K$가 모두 영은 아니고 $f(\alpha) = 0$을 만족한다는 뜻이다.
그러한 $\alpha$를 $f$의 {\it 비자명해}라 한다.
\end{definition}

\begin{example} \label{etale-cohomology-example-algebraically-closed-field-Cr} 대수적으로 닫힌 체는 $C_r$이다. \end{example}

\noindent 실제로 다음과 같은 간단한 보조정리가 있다.

\begin{lemma}
\label{etale-cohomology-lemma-algebraically-closed-find-solutions}
$k$는 대수적으로 닫힌 체라고 가정한다. 허락하다
$f_1, \ldots, f_s \in k[T_1, \ldots, T_n]$
차수가 $d_1, \ldots, d_s$이고 $d_i > 0$인 동차다항식들이라 하자. $s < n$이면 $f_1 = \ldots = f_s = 0$은
공통의 비자명해를 갖는다.
\end{lemma}

\begin{proof} 이는 차원 이론을 따르며, 예에 대해 다양체, 보조정리 \ref{varieties-lemma-intersection-in-affine-space}가 $s - 1$번 적용되는 형식이다. \end{proof}

\noindent 다음 결과는 $C_1$ 체의 Brauer 군을 계산한다.

\begin{theorem} \label{etale-cohomology-theorem-C1-brauer-group-zero} $K$를 $C_1$ 체로 둡니다. 그런 다음 $\text{Br}(K) = 0$. \end{theorem}

\begin{proof}
허락하다$D$유한 차원 분할 대수학이 되십시오$K$센터와 함께$K$. 우리
그걸 봤어
$$
D \otimes_K K^{sep} \cong \text{Mat}_d(K^{sep})
$$
내부 동형사상까지 고유한다. 따라서 행렬식 $\det : \text{Mat}_d(K^{sep}) \to K^{sep}$는 갈루아입니다 불변이고 다음으로 내려간다.
동차차 $d$ 맵
$$
\det = N_\text{red} : D \longrightarrow K
$$
{\it 감소된 표준}이라고 한다. $K$는 $C_1$이므로 $d > 1$이면
0이 아닌 값이 존재합니다$x \in D$~와 함께$N_\text{red}(x) = 0$. 이는 분명히 암시한다.
$x$는 되돌릴 수 없으며 이는 모순이다. 그러므로 $\text{Br}(K) = 0$.
\end{proof}

\begin{definition} \label{etale-cohomology-definition-variety} $k$를 체로 둡니다. {\it 다양체}는 $k$에 대해 유한 유형의 정역 스킴으로 구분된다. {\it 곡선}은 차원 $1$의 다양체이다. \end{definition}

\begin{theorem}[Tsen의 정리] \label{etale-cohomology-theorem-tsen} 대수적으로 닫힌 체에 대한 차원 $r$의 다양체의 함수 체 $k$는 $C_r$이다. \end{theorem}

\begin{proof} 투영 공간의 경우 체 $k(x_1, \ldots, x_r)$가 $C_r$(연습)임을 직접적으로 보여줄 수 있다.

\medskip\noindent
일반적인 경우. 일반성을 잃지 않으면서 $X$가 투영적이라고 가정할 수 있다.
$f \in k(X)[T_1, \ldots, T_n]_d$를 $0 < d^r < n$로 설정한다.
$f$의 계수가 $\Gamma(X, \mathcal{O}_X(H))$에 있다고 가정한다.
충분한 $H \subset X$에 대해. 허락하다
$\mathbf{\alpha} = (\alpha_1, \ldots, \alpha_n)$이고 $\alpha_i \in \Gamma(X, \mathcal{O}_X(eH))$라 하자. 그러면 $f(\mathbf{\alpha}) \in \Gamma(X, \mathcal{O}_X((de + 1)H))$이다. 방정식계 $f(\mathbf{\alpha}) =0$을 생각하자. 그러면 점근적 Riemann--Roch 정리에 의해
(다양체, 명제 \ref{varieties-proposition-asymptotic-riemann-roch})
$c > 0$이 존재한다.
\begin{itemize}
\item 변수의 개수는
$n\dim_k \Gamma(X, \mathcal{O}_X(eH)) \sim n e^r c$ 그리고
\item 방정식의 개수는 다음과 같습니다.
$\dim_k \Gamma(X, \mathcal{O}_X((de + 1)H)) \sim (de + 1)^r c$.
\end{itemize}
$n > d^r$이므로 방정식보다 변수가 더 많습니다. 방정식은 다음과 같습니다
동질적이므로 다음과 같은 솔루션이 있다.
보조정리 \ref{etale-cohomology-lemma-algebraically-closed-find-solutions}.
\end{proof}

\begin{lemma} \label{etale-cohomology-lemma-curve-brauer-zero} $C$는 대수적으로 닫힌 체 $k$에 대한 곡선이라고 가정한다. 그러면 $C$의 체 기능의 Brauer 군은 0입니다: $\text{Br}(k(C)) = 0$. \end{lemma}

\begin{proof} 이는 Tsen의 정리, 정리 \ref{etale-cohomology-theorem-tsen} 및 정리 \ref{etale-cohomology-theorem-C1-brauer-group-zero}에서 분명한다. \end{proof}

\begin{lemma}
\label{etale-cohomology-lemma-cohomology-Gm-function-field-curve}
$k$는 대수적으로 닫힌 체 및 $K/k$ 체 확장이라고 가정한다.
초월 등급 1. 그런 다음 모든 $q \geq 1$에 대해
$H_\etale^q(\Spec(K), \mathbf{G}_m) = 0$.
\end{lemma}

\begin{proof}
그것을 기억해
$H_\etale^q(\Spec(K), \mathbf{G}_m) = H^q(\text{Gal}(K^{sep}/K), (K^{sep})^*)$
보조정리 \ref{etale-cohomology-lemma-compare-cohomology-point} 기준.
따라서 명제 \ref{etale-cohomology-proposition-serre-galois}
$K'/K$가 유한한 체 확장인 경우,
$\text{Br}(K') = 0$. 이제 $K' = \colim K''$를 관찰하세요. 여기서 $K''$가 실행된다.
유한하게 생성된 하위 확장에 대해$k$에 포함된$K'$~의
초월도 $1$.
우리를 감소시키는 $\text{Br}(K') = \colim \text{Br}(K'')$에 주목하세요.
유한하게 생성된 체 초월의 확장 $K''/k$
정도 $1$. 이러한 체는 $k$에 대한 곡선의 체 함수이다.
따라서 사소한 Brauer 군이 있다.
보조정리 \ref{etale-cohomology-lemma-curve-brauer-zero}.
\end{proof}






\section{곱셈 그룹의 경우 더 높은 소멸 } \label{etale-cohomology-section-higher-Gm}

\noindent 이 절에서는 대수적으로 닫힌 체 $k$와 $X$를 $k$ 위에 고정한다. $i_x : x \hookrightarrow X$는 $X$의 닫힌 점 포함을 나타내고 $j : \eta \hookrightarrow X$는 일반 점의 포함을 나타냅니다. 또한 $X_0$는 $X$의 폐점의 집합을 나타냅니다.

\begin{theorem}[기본적인 완전열] \label{etale-cohomology-theorem-fundamental-exact-sequence} $X$ $$
0 \longrightarrow
\mathbf{G}_{m, X} \longrightarrow
j_* \mathbf{G}_{m, \eta} \longrightarrow
\bigoplus\nolimits_{x \in X_0} {i_x}_* \underline{\mathbf{Z}}
\longrightarrow 0.
$$ \end{theorem}에 에탈 층의 짧은 완전열이 있다.

\begin{proof} $\varphi : U \to X$를 에탈 사상으로 하자. 그런 다음 에탈 사상 (명제 \ref{etale-cohomology-proposition-etale-morphisms}), $U = \coprod_i U_i$의 성질에 의해 각 $U_i$는 $X$에 매핑되는 부드러운 곡선이다. 위의 $U$에 대한 시퀀스는 각 $U_i$에 대한 해당 시퀀스의 곱이므로 $U$가 연결, 즉 기약인 경우를 처리하면 충분한다. 이 경우 잘 알려진 완전열 $$
1 \longrightarrow
\Gamma(U, \mathcal{O}_U^*) \longrightarrow
k(U)^* \longrightarrow \bigoplus\nolimits_{y \in U_0} \mathbf{Z}_y.
$$ 이는 시퀀스 $$
0 \longrightarrow \Gamma(U, \mathcal{O}_U^*) \longrightarrow
\Gamma(\eta \times_X U, \mathcal{O}_{\eta \times_X U}^*) \longrightarrow
\bigoplus\nolimits_{x \in X_0}
\Gamma(x \times_X U, \underline{\mathbf{Z}})
$$에 해당하며 전개 정의는 시퀀스 $$
0 \longrightarrow \mathbf{G}_m(U) \longrightarrow
j_* \mathbf{G}_{m, \eta}(U) \longrightarrow
\left(\bigoplus\nolimits_{x \in X_0} {i_x}_* \underline{\mathbf{Z}}\right)(U).
$$에 지나지 않습니다. 이는 기본 완전열의 맵을 정의하고 마지막 단계를 제외하고는 정확함을 보여줍니다. 전사성을 확인하기 위해 $U$가 비특이 곡선이고 $D$가 $U$의 제수인 경우 피복 $\{U_j \to U\}$의 $U$가 존재한다는 점을 상기해 보자. $D|_{U_j} = \text{div}(f_j)$ 일부 $f_j \in k(U)^*$에 대해. \end{proof}

\begin{lemma} \label{etale-cohomology-lemma-higher-direct-jstar-Gm} 모든 $q \geq 1$, $R^q j_*\mathbf{G}_{m, \eta} = 0$에 대해. \end{lemma}

\begin{proof}
우리는 모든 것에 대해 $(R^q j_*\mathbf{G}_{m, \eta})_{\bar x} = 0$를 보여줘야 한다.
기하점 $\bar x$의 $X$.

\medskip\noindent
가정$\bar x$폐쇄된 지점 위에 위치$x$~의$X$. 허락하다$\Spec(A)$
친밀하고 개방적인 이웃이 되다$x$~에$X$, 그리고$K$분수체
$A$. 그 다음에
$$
\Spec(\mathcal{O}^{sh}_{X, \bar x}) \times_X \eta =
\Spec(\mathcal{O}^{sh}_{X, \bar x} \otimes_A K).
$$
환 $\mathcal{O}^{sh}_{X, \bar x} \otimes_A K$는 다음의 국소화이다.
이산적 평가 환 $\mathcal{O}^{sh}_{X, \bar x}$이므로 다음 중 하나이다.
$\mathcal{O}^{sh}_{X, \bar x}$ 다시 또는 그 분수 체
$K^{sh}_{\bar x}$. 하지만 일부 로컬 균일화 프로그램은 반전되기 때문에
후자이세요. 따라서
$$
(R^q j_*\mathbf{G}_{m, \eta})_{(X, \bar x)} = H_\etale^q(\Spec
K^{sh}_{\bar x}, \mathbf{G}_m).
$$
이제 $\mathcal{O}^{sh}_{X, \bar x} =
\colim_{(U, \bar u) \to \bar x} \mathcal{O}(U) = \colim_{A \subset B} B$
여기서 $A \to B$는 에탈이므로 $K^{sh}_{\bar x}$는 대수이다.
$K = k(X)$ 확장을 적용할 수 있다.
보조정리 \ref{etale-cohomology-lemma-cohomology-Gm-function-field-curve}
소멸을 얻으려면

\medskip\noindent
가정$\bar x = \bar \eta$일반적인 지점 위에 위치$\eta$~의$X$(안에
사실, 이 경우는 불필요합니다). 그 다음에
$\mathcal{O}^{sh}_{X, \bar \eta} = \kappa(\eta)^{sep}$ 따라서
\begin{eqnarray*}
(R^q j_*\mathbf{G}_{m, \eta})_{\bar \eta}
& = &
H_\etale^q(\Spec(\kappa(\eta)^{sep}) \times_X \eta, \mathbf{G}_m) \\
& = & H_\etale^q (\Spec(\kappa(\eta)^{sep}), \mathbf{G}_m) \\
& = & 0 \ \ \text{~을 위한} q \geq 1
\end{eqnarray*}
대응하는 갈루아 그룹은 사소하기 때문이다.
\end{proof}

\begin{lemma} \label{etale-cohomology-lemma-cohomology-jstar-Gm} 모든 $p \geq 1$, $H_\etale^p(X, j_*\mathbf{G}_{m, \eta}) = 0$에 대해. \end{lemma}

\begin{proof}
Leray 스펙트럼 열은 다음과 같습니다.
$$
E_2^{p, q} = H_\etale^p(X, R^qj_*\mathbf{G}_{m, \eta}) \Rightarrow
H_\etale^{p+q}(\eta, \mathbf{G}_{m, \eta}),
$$
$p + q \geq 1$에 대해 오른쪽이 사라집니다.
보조정리 \ref{etale-cohomology-lemma-cohomology-Gm-function-field-curve}.
$q \geq 1$에 대해 왼쪽이 사라집니다.
보조정리 \ref{etale-cohomology-lemma-higher-direct-jstar-Gm}. 따라서 우리는 원하는 소멸을 얻습니다.
\end{proof}

\begin{lemma} \label{etale-cohomology-lemma-cohomology-istar-Z} 모든 $q \geq 1$, $H_\etale^q(X, \bigoplus_{x \in X_0} {i_x}_*
\underline{\mathbf{Z}}) = 0$에 대해. \end{lemma}

\begin{proof}
$X$ 준소형 및 준분리형의 경우 코호몰로지는 여극한으로 통근한다.
따라서 $H_\etale^q(X, {i_x}_* \underline{\mathbf{Z}})$. 그러나 그러면 inclusion $i_x$의 폐점은 다음과 같습니다.
유한해서$R^p {i_x}_* \underline{\mathbf{Z}} = 0$모두를 위해$p \geq 1$~에 의해
명제 \ref{etale-cohomology-proposition-finite-higher-direct-image-zero}.
Leray 스펙트럼 열을 적용하면 다음을 알 수 있다.
$H_\etale^q(X, {i_x}_* \underline{\mathbf{Z}}) =
H_\etale^q(x, \underline{\mathbf{Z}})$.
마지막으로, $x$는
대수적으로 닫힌 체, $x$의 더 높은 코호몰로지가 모두 사라집니다.
\end{proof}

\noindent 이 일련의 정리를 마치면 다음과 같은 결과를 얻습니다.

\begin{theorem} \label{etale-cohomology-theorem-vanishing-cohomology-Gm-curve} $X$가 대수적으로 닫힌 체에 대한 부드러운 곡선이라고 가정한다. 그러면 $$
H_\etale^q(X, \mathbf{G}_m) = 0 \ \ \text{모두를 위해} q \geq 2.
$$ \end{theorem}

\begin{proof} 위의 논의를 참조하라. \end{proof}

\noindent 코호몰로지 긴 완전열 $$
0 \to
H_\etale^0(X, \mathbf{G}_m) \to
H_\etale^0(X, j_*\mathbf{G}_{m\eta}) \to
H_\etale^0(X, \bigoplus {i_x}_*\underline{\mathbf{Z}}) \to
H_\etale^1(X, \mathbf{G}_m) \to 0
$$도 얻습니다. 비록 이것이 친숙한 $$
0 \to H_{Zar}^0(X, \mathcal{O}_X^*) \to k(X)^* \to \text{Div}(X)
\to \Pic(X) \to 0.
$$이지만





\section{Picard 군 곡선} \label{etale-cohomology-section-pic-curves}

\noindent 다음 단계는 Kummer 수열을 사용하여 유한 계수를 갖는 곡선의 코호몰로지 군에 대한 일부 정보를 추론하는 것이다. 긴 완전열에서 소멸을 얻기 위해 Picard 군에 대한 몇 가지 사실을 검토한다.

\medskip\noindent
$X$는 대수적으로 닫힌 체 $k$에 대한 부드러운 사영 곡선이라고 가정한다.
$g = \dim_k H^1(X, \mathcal{O}_X)$를 $X$의 속으로 둡니다.
짧은 완전열이 존재한다.
$$
0 \to \Pic^0(X) \to \Pic(X) \xrightarrow{\deg} \mathbf{Z} \to 0.
$$
아벨 군 $\Pic^0(X)$는 다음으로 식별할 수 있다.
$\Pic^0(X) = \underline{\Picardfunctor}^0_{X/k}(k)$ 즉,
아벨 다양체의 $k$ 값 포인트
$\underline{\Picardfunctor}^0_{X/k}$~ 위에$k$~의차원 $g$.
결과적으로 $n \in k^*$이면
$\Pic^0(X)[n] \cong (\mathbf{Z}/n\mathbf{Z})^{2g}$
아벨 군으로. 보다
곡선의 피카드 스킴, 절 \ref{pic-section-picard-curve}
그리고
그룹형, 절 \ref{groupoids-section-abelian-varieties}.
이 핵심 사실, 즉 비틀림에 대한 설명이다.
대수학적으로 부드러운 사영 곡선의 Picard 군
닫힌 체에는 기본 증명이 없는 것 같습니다.

\begin{lemma}
\label{etale-cohomology-lemma-cohomology-smooth-projective-curve}
$X$를 $g$ 속의 부드러운 사영 곡선으로 정의한다.
대수적으로 닫힌 체 $k$ 그리고 $n \geq 1$가 $k$에서 반전 가능하도록 둡니다.
그런 다음 정식 식별이 있다.
$$
H_\etale^q(X, \mu_n) =
\left\{
\begin{matrix}
\mu_n(k) & \text{만약에}q = 0, \\
\Pic^0(X)[n] & \text{만약에}q = 1, \\
\mathbf{Z}/n\mathbf{Z} & \text{만약에}q = 2, \\
0 & \text{만약에}q \geq 3.
\end{matrix}
\right.
$$
$\mu_n \cong \underline{\mathbf{Z}/n\mathbf{Z}}$ 이후 이는 다음과 같습니다.
(비표준) 식별
$$
H_\etale^q(X, \underline{\mathbf{Z}/n\mathbf{Z}}) \cong
\left\{
\begin{matrix}
\mathbf{Z}/n\mathbf{Z} & \text{만약에}q = 0, \\
(\mathbf{Z}/n\mathbf{Z})^{2g} & \text{만약에}q = 1, \\
\mathbf{Z}/n\mathbf{Z} & \text{만약에}q = 2, \\
0 & \text{만약에}q \geq 3.
\end{matrix}
\right.
$$
\end{lemma}

\begin{proof}
정리 \ref{etale-cohomology-theorem-picard-group} 및
\ref{etale-cohomology-theorem-vanishing-cohomology-Gm-curve}
에탈을 결정하다코호몰로지~의$\mathbf{G}_m$~에$X$
$X$의 Picard 군에 관하여.
Kummer 수열 $0\to \mu_{n, X} \to \mathbf{G}_{m, X}
\to \mathbf{G}_{m, X}\to 0$ (보조정리 \ref{etale-cohomology-lemma-kummer-sequence})
그런 다음 길고 정확한 코호몰로지 시퀀스를 제공한다.
$$
\xymatrix{
0 \ar[r] & \mu_n(k) \ar[r] &
k^* \ar[r]^{(\cdot)^n} &
k^* \ar@(rd, ul)[rdllllr] \\
& H_\etale^1(X, \mu_n) \ar[r] &
\Pic(X) \ar[r]^{(\cdot)^n} &
\Pic(X) \ar@(rd, ul)[rdllllr] \\
& H_\etale^2(X, \mu_n) \ar[r] & 0 \ar[r] & 0 \ldots
}
$$
$n$번째 파워 맵 $k^* \to k^*$은 대수적으로 $k$이므로 전사이다.
닫은. 따라서 지도의 핵 및 여핵을 계산해야 한다.
$n : \Pic(X) \to \Pic(X)$. 가환 도식을 고려해보세요.
정확한 행
$$
\xymatrix{
0 \ar[r] &
\Pic^0(X) \ar[r] \ar@{>>}[d]^{(\cdot)^n} &
\Pic(X) \ar[r]_-\deg \ar[d]^{(\cdot)^n} &
\mathbf{Z} \ar[r] \ar@{^{(}->}[d]^n & 0 \\
0 \ar[r] &
\Pic^0(X) \ar[r] &
\Pic(X) \ar[r]^-\deg &
\mathbf{Z} \ar[r] & 0
}
$$
$\Pic^0(X)$ 그룹은 $k$ 포인트이다.
그룹 스킴 $\underline{\Picardfunctor}^0_{X/k}$, 참조
곡선의 Picard 스킴, 보조정리 \ref{pic-lemma-picard-pieces}.
동일한 보조정리는 $\underline{\Picardfunctor}^0_{X/k}$를 알려줍니다.
는$g$-차원 아벨리안다양체~ 위에$k$정의된 대로
그룹형, 정의 \ref{groupoids-definition-abelian-variety}.
따라서 왼쪽 수직 맵은 다음과 같이 전사적이다.
그룹형, 명제 \ref{groupoids-proposition-review-abelian-varieties}.
뱀 보조정리를 적용하면 명시된 대로 정식 식별이 제공된다.
보조정리에서.

\medskip\noindent 보조정리의 비정규 식별을 얻으려면 $n : \Pic^0(X) \to \Pic^0(X)$의 핵이 $(\mathbf{Z}/n\mathbf{Z})^{\oplus 2g}$과 동형임을 보여야 한다. 이것은 또한 Groupoids, 명제 \ref{groupoids-proposition-review-abelian-varieties}의 일부이다. \end{proof}

\begin{lemma} \label{etale-cohomology-lemma-pullback-on-h2-curve} $\pi : X \to Y$는 대수적으로 닫힌 체 $k$에 대한 부드러운 투영 곡선의 비상수 사상이고 $n \geq 1$는 반전 가능한다. $k$. 지도 $$
\pi^* : H^2_\etale(Y, \mu_n) \longrightarrow H^2_\etale(X, \mu_n)
$$는 $\pi$의 차수를 곱하여 주어집니다. \end{lemma}

\begin{proof}
우리가 두 가지를 모두 식별했으므로 진술이 의미가 있음을 관찰하라.
코호몰로지 군 $ H^2_\etale(Y, \mu_n)$ 및 $H^2_\etale(X, \mu_n)$
$\mathbf{Z}/n\mathbf{Z}$ 포함
보조정리 \ref{etale-cohomology-lemma-cohomology-smooth-projective-curve}.
실제로 $\mathcal{L}$가 차수 $1$의 라인 묶음인 경우
$Y$ 클래스 $[\mathcal{L}] \in H^1_\etale(Y, \mathbf{G}_m)$에서,
그러면 $[\mathcal{L}]$의 공동 경계는 다음의 생성기이다.
$H^2_\etale(Y, \mu_n)$. 여기서 공동경계는 공동경계이다.
Kummer와 관련된 코호몰로지의 긴 완전열 중
순서. 따라서 보조정리의 결과는 다음과 같습니다.
라인 묶음의 정도$\pi^*\mathcal{L}$~에$X$~이다$\deg(\pi)$.
일부 세부 사항은 생략되었다.
\end{proof}

\begin{lemma} \label{etale-cohomology-lemma-vanishing-cohomology-mu-smooth-curve} $X$를 대수적으로 닫힌 체 $k$ 및 $n \in k^*$에 대한 아핀 매끄러운 곡선으로 가정한다. $X \subset \overline{X}$를 부드러운 투영 압축(다양체, 비고 \ref{varieties-remark-smooth-projective-compactification})이라고 한다. $g$를 $\overline{X}$의 속으로 두고 $r$를 $\overline{X} \setminus X$의 포인트 수로 둡니다. 그런 다음 \begin{enumerate} \item $H_\etale^0(X, \mu_n) = \mu_n(k)$; \item $H_\etale^1(X, \mu_n) \cong (\mathbf{Z}/n\mathbf{Z})^{2g+r-1}$ 및 \item $H_\etale^q(X, \mu_n) = 0$ 모든 $q \geq 2$에 대해. \end{enumerate} \end{lemma}

\begin{proof}
$X = \overline{X} - \{ x_1, \ldots, x_r\}$를 씁니다. 그 다음에
$\Pic(X) = \Pic(\overline{X})/ R$, 여기서 $R$는 다음에 의해 생성된 하위 그룹이다.
$\mathcal{O}_{\overline{X}}(x_i)$, $1 \leq i \leq r$.
$r \geq 1$ 이후로 $\Pic^0(\overline{X}) \to \Pic(X)$
전사이므로 $\Pic(X)$는 나누어질 수 있습니다(증명의 논의 참조).
보조정리 \ref{etale-cohomology-lemma-cohomology-smooth-projective-curve}).
Kummer 수열을 적용하면 (1) 및 (3)를 얻습니다. (2)에 대해 우리는 다음을 사용한다.
\begin{align*}
H_\etale^1(X, \mu_n)
& =
\{(\mathcal L, \alpha)\}/\cong \\
& =
(\{(\bar{\mathcal L}, D, \bar \alpha)\}/\cong)/\tilde{R}
\end{align*}
첫 번째 등호 및 표기법은 보조정리 \ref{etale-cohomology-lemma-describe-h1-mun}를 참조하라.
두 번째 줄의 트리플 $(\bar{\mathcal L}, D, \bar \alpha)$
인버터로 구성가군 $\bar{\mathcal L}$~에$\overline{X}$
학위$0$, 제수$D$~에$\overline{X}$학위$0$지원됨
$\left\{x_1, \ldots, x_r\right\}$ 및 동형사상
$\bar{\alpha}: \bar{\mathcal L}^{\otimes n} \cong \mathcal{O}_{\bar{X}}(D)$.
트리플의 동형사상은 이전 논의에서와 같이 정의된다.
보조정리 \ref{etale-cohomology-lemma-describe-h1-mun}.
또한 $\tilde{R}$는 다음 형식의 트리플로 생성된 하위 그룹이다.
$(\mathcal{O}_{\overline{X}}(D'), n D', 1^{\otimes n})$
여기서 $D'$는 $\left\{x_1, \ldots, x_r\right\}$에서 지원되며 등급은 0이다.
위의 논의는 다음을 통해 두 번째 등호가 있음을 보여줍니다.
트리플 $(\bar{\mathcal L},\ D,\ \bar \alpha)$을 보내는 지도
$(\bar{\mathcal L}|_X, \bar \alpha|_X)$ 쌍으로. 따라서,
우리는 완전열을 얻습니다.
$$
0 \longrightarrow
H_\etale^1(\overline{X}, \mu_n) \longrightarrow
H_\etale^1(X, \mu_n) \longrightarrow
\bigoplus_{i = 1}^r \mathbf{Z}/n\mathbf{Z}
\xrightarrow{\ \sum\ }
\mathbf{Z}/n\mathbf{Z} \longrightarrow 0
$$
중간 맵이 트리플 클래스를 보내는 곳
$(\bar{ \mathcal L}, D, \bar \alpha)$ 와 $D = \sum_{i = 1}^r a_i (x_i)$
$r$-튜플 $(a_i)_{i = 1}^r$로. 이제 사용해도 충분합니다
보조정리 \ref{etale-cohomology-lemma-cohomology-smooth-projective-curve}
순위를 계산한다.
\end{proof}

\begin{remark}
\label{etale-cohomology-remark-natural-proof}
이전 것을 증명하는 '자연스러운' 방법따름정리소비세를 내는 것이다$X$$\bar X$. 이것은 가능한다. 우리는 그 이론을 개발하지 않았을 뿐이다.
\end{remark}

\begin{remark}
\label{etale-cohomology-remark-normalize-H1-Gm}
$k$는 대수적으로 닫힌 체라고 가정한다. $n$를 정수 소수로 설정한다.
$k$의 표수. 그것을 기억해
$$
\mathbf{G}_{m, k} = \mathbf{A}^1_k \setminus \{0\} =
\mathbf{P}^1_k \setminus \{0, \infty\}
$$
우리는 주장 표준 동형사상이 있다.
$$
H^1_\etale(\mathbf{G}_{m, k}, \mu_n) = \mathbf{Z}/n\mathbf{Z}
$$
이것은 무엇을 의미합니까? 이는 $1_k$ 요소가 있음을 의미한다.
$H^1_\etale(\mathbf{G}_{m, k}, \mu_n)$ 그런 식으로
사상 $\Spec(k') \to \Spec(k)$ 당김 맵마다
에탈 코호몰로지 지도 $\mathbf{G}_{m, k'} \to \mathbf{G}_{m, k}$
$1_k$을 $1_{k'}$에 매핑한다. (특히 이 요소는
$k$의 모든 자동형에 따라 고정된다.) 이를 확인하려면 다음을 고려하세요.
$\mu_{n, \mathbf{Z}}$-토르서
$\mathbf{G}_{m, \mathbf{Z}} \to \mathbf{G}_{m, \mathbf{Z}}$,
$x \mapsto x^n$. 토르서를 식별함으로써
먼저 코호몰로지, 이는 표준 요소 $1_k$를 제공하기 위해 뒤로 당겨집니다.
뒤돌아 보면 표준 식별이 있음을 알 수 있다.
$$
H^1_\etale(\mathbf{G}_{m, k}, \mathbf{Z}/n\mathbf{Z}) =
\Hom(\mu_n(k), \mathbf{Z}/n\mathbf{Z}),
$$
즉, 이 동형사상은 다음의 지도와 호환된다.
대수적으로 닫힌 체, 특히 다음과 관련하여
$k$의 자동형.
\end{remark}










\section{0으로 확장} \label{etale-cohomology-section-extension-by-zero}

\noindent 사이트, 절 \ref{sites-modules-section-localize}의 가군에 있는 일반 자료를 사용하면 다음 정의를 만들 수 있다.

\begin{definition} \label{etale-cohomology-definition-extension-zero} $j : U \to X$를 스킴의 에탈 사상으로 둡니다. \begin{enumerate} \item 제한사항 함자 $j^{-1} : \Sh(X_\etale) \to \Sh(U_\etale)$에는 왼쪽 수반 $j_!^{Sh} : \Sh(U_\etale) \to \Sh(X_\etale)$이 있다. \item 제한 함자 $j^{-1} : \textit{Ab}(X_\etale) \to \textit{Ab}(U_\etale)$에는 $j_! : \textit{Ab}(U_\etale) \to \textit{Ab}(X_\etale)$로 표시되고 {\it 확장자}라고 불리는 왼쪽 수반이 있다. \item $\Lambda$를 환으로 둡니다. 제한 함자 $j^{-1} : \textit{Mod}(X_\etale, \Lambda) \to
\textit{Mod}(U_\etale, \Lambda)$에는 $j_! : \textit{Mod}(U_\etale, \Lambda) \to
\textit{Mod}(X_\etale, \Lambda)$로 표시되고 {\it 확장자}라고 불리는 왼쪽 수반이 있다. \end{enumerate} \end{definition}

\noindent
$\mathcal{F}$가 $X_\etale$의 아벨 층인 경우
$j_!\mathcal{F} \not = j_!^{Sh}\mathcal{F}$ 일반적으로. 반면에
$j_!$~을 위한층~의$\Lambda$-가군동의하다$j_!$기초에
아벨 층
(가군 사이트, 비고 \ref{sites-modules-remark-localize-shriek-equal}).
함자 $j_!$는 함수형 동형사상을 특징으로 한다.
$$
\Hom_X(j_!\mathcal{F}, \mathcal{G}) = \Hom_U(\mathcal{F}, j^{-1}\mathcal{G})
$$
모든 $\mathcal{F} \in \textit{Ab}(U_\etale)$에 대해 그리고
$\mathcal{G} \in \textit{Ab}(X_\etale)$. 마찬가지로
층 중 $\Lambda$-가군.

\medskip\noindent 정의 \ref{etale-cohomology-definition-extension-zero}의 함자를 더 명시적으로 설명하려면 $j^{-1}$는 함자 $U_\etale \to X_\etale$를 통한 제한일 뿐이라는 점을 기억하세요. 즉, $j^{-1}\mathcal{G}(U') = \mathcal{G}(U')$는 $U'$에 대해 $U$에 대해 에탈이다. 반면, $\mathcal{F} \in \textit{Ab}(U_\etale)$의 경우 전층 \begin{equation}
\label{etale-cohomology-equation-j-p-shriek}
j_{p!}\mathcal{F} : X_\etale \longrightarrow \textit{Ab},
\quad
V \longmapsto \bigoplus\nolimits_{V \to U} \mathcal{F}(V \to U)
\end{equation} 그러면 $j_!\mathcal{F}$는 층화이다. $j_{p!}\mathcal{F}$. 이는 사이트, 보조정리 \ref{sites-modules-lemma-extension-by-zero}의 가군에서 입증되었다. 더 일반적으로는 사이트, 절 \ref{sites-modules-section-localize} 및 \ref{sites-modules-section-exactness-lower-shriek}에 대한 가군의 토론을 참조하라.

\begin{exercise}
\label{etale-cohomology-exercise-jshriek-direct}
함자 $j_!$이 다음의 층화로 정의되었음을 직접 증명
함자 $j_{p!}$에 주어진
(\ref{etale-cohomology-equation-j-p-shriek})는 왼쪽 수반 ~ $j^{-1}$이다.
\end{exercise}

\begin{proposition}
\label{etale-cohomology-proposition-describe-jshriek}
$j : U \to X$를 스킴의 에탈 사상이라고 하자.
허락하다$\mathcal{F}$~에$\textit{Ab}(U_\etale)$.
$\overline{x} : \Spec(k) \to X$가 $X$의 기하점인 경우
$$
(j_!\mathcal{F})_{\overline{x}} =
\bigoplus\nolimits_{\overline{u} : \Spec(k) \to U,\ j(\overline{u}) =
\overline{x}} \mathcal{F}_{\bar{u}}.
$$
특히, $j_!$은 완전 함자이다.
\end{proposition}

\begin{proof} 완전성 의 $j_!$는 매우 일반적이다. 가군 위의 사이트, 보조정리 \ref{sites-modules-lemma-extension-by-zero-exact}를 참조하라. 물론 줄기의 설명에도 따릅니다. 줄기에 대한 공식은 사이트, 보조정리 \ref{sites-modules-lemma-stalk-j-shriek}의 가군과 기하점에 대한 작은 에탈 사이트의 설명을 따릅니다. 보조정리 \ref{etale-cohomology-lemma-points-small-etale-site}.

\medskip\noindent
나중에 사용할 수 있도록 동형사상
\begin{align*}
(j_!\mathcal{F})_{\overline{x}}
& =
(j_{p!}\mathcal{F})_{\overline{x}} \\
& =
\colim_{(V, \overline{v})} j_{p!}\mathcal{F}(V) \\
& =
\colim_{(V, \overline{v})}
\bigoplus\nolimits_{\varphi : V \to U}
\mathcal{F}(V \xrightarrow{\varphi} U) \\
& \to
\bigoplus\nolimits_{\overline{u} : \Spec(k) \to U,\ j(\overline{u}) =
\overline{x}} \mathcal{F}_{\bar{u}}.
\end{align*}
사이트, 보조정리 \ref{sites-modules-lemma-stalk-j-shriek}에 가군에 구축됨
보낸다$(V, \overline{v}, \varphi, s)$~의 수업에$s$에서줄기
~의$\mathcal{F}$~에$\overline{u} = \varphi(\overline{v})$.
\end{proof}

\begin{lemma}
\label{etale-cohomology-lemma-jshriek-open}
$j : U \to X$를 스킴의 열린 몰입으로 설정한다. 누구에게나
아벨 층 $\mathcal{F}$ $U_\etale$의 부속 매핑
$j^{-1}j_*\mathcal{F} \to \mathcal{F}$ 그리고
$\mathcal{F} \to j^{-1}j_!\mathcal{F}$는 동형사상이다.
실제로 $j_!\mathcal{F}$는 $X_\etale$의 고유한 아벨 층이다.
$U$에 대한 제한은 $\mathcal{F}$이고 줄기는
$X \setminus U$의 기하점은 0이다.
\end{lemma}

\begin{proof}
우리는 독자가 다음을 통해 첫 번째 진술을 증명하도록 권장한다.
그러나 여기서는 그것이 매우 일반적인 정의의 특별한 경우라는 점만 사용한다.
가군 사이트, 보조정리 \ref{sites-modules-lemma-restrict-back}.
두 번째 명령문에서는 $\mathcal{G}$가 아벨 층인 경우를 관찰하세요.
~에$X_\etale$누구의 제한$U$~이다$\mathcal{F}$, 그러면 우리는 얻는다
지도 $j_!\mathcal{F} \to \mathcal{G}$와 인접해 있다. 이 지도
그러면 $U$의 기하점의 줄기에 있는 동형사상은 다음과 같습니다.
명제 \ref{etale-cohomology-proposition-describe-jshriek}.
따라서 $\mathcal{G}$에 소멸 줄기가 기하점에 있는 경우
$X \setminus U$의 경우 $j_!\mathcal{F} \to \mathcal{G}$는
동형사상 에 의해 정리 \ref{etale-cohomology-theorem-exactness-stalks}.
\end{proof}

\begin{lemma}[기본 변경으로 통근 제로 연장] \label{etale-cohomology-lemma-shriek-base-change} $f: Y \to X$를 스킴의 사상으로 둡니다. $j: V \to X$를 에탈 사상라고 하자. 올곱 $$
\xymatrix{
V' = Y \times_X V \ar[d]_{f'} \ar[r]_-{j'} & Y \ar[d]^f \\
V \ar[r]^j & X
}
$$를 고려하면 아벨 층과 가군층에 $j'_! f'^{-1} = f^{-1} j_!$가 있다. \end{lemma}

\begin{proof} 이는 $j'_! f'^{-1}$가 왼쪽 수반 ~ $f'_* (j')^{-1}$이고 $f^{-1} j_!$가 왼쪽 수반 ~ $j^{-1}f_*$이기 때문에 이는 사실이다. 또한 $f'_* (j')^{-1} = j^{-1}f_*$는 $f_*$가 에탈 국소화로 출퇴근하기 때문이다. 실제로 보조정리는 사이트의 사상 설정에서 매우 일반적으로 유지된다. 사이트, 보조정리 \ref{sites-modules-lemma-localize-morphism-ringed-sites}의 가군을 참조하라. \end{proof}

\begin{lemma} \label{etale-cohomology-lemma-shriek-into-star-separated-etale} $j : U \to X$를 분리하여 에탈이라고 한다. 그런 다음 아벨 층에 기능적 분사 맵 $j_!\mathcal{F} \to j_*\mathcal{F}$이 있고 $\Lambda$-가군의 층이 있다. \end{lemma}

\begin{proof}
아벨 층의 경우에서 이를 증명한다. 표준지도를 만들어보자
$$
j_{p!}\mathcal{F} \to j_*\mathcal{F}
$$
모든 아벨 층에 대해 $X_\etale$에 대한 아벨 전층의
$\mathcal{F}$ 위의 $U_\etale$ 여기서 $j_{p!}$는 다음과 같습니다.
(\ref{etale-cohomology-equation-j-p-shriek}). 이 지도의 층화는
원하는 맵 $j_!\mathcal{F} \to j_*\mathcal{F}$이 된다.
$V \to X$에 대해 양쪽을 평가하면 에탈을 얻습니다.
$$
j_{p!}\mathcal{F}(V) =
\bigoplus\nolimits_{\varphi : V \to U} \mathcal{F}(V \xrightarrow{\varphi} U)
\quad\text{이고}\quad
j_*\mathcal{F}(V) = \mathcal{F}(V \times_X U)
$$
각 $\varphi$에 대해 열려 있고 닫힌 몰입이 있다.
$$
\Gamma_\varphi = (1, \varphi) : V \longrightarrow V \times_X U
$$
$U$ 이상. 스킴의 에탈 사이에 사상이므로 열려 있다.
$U$ 위에 있고 스킴으로 분리된 절이므로 닫혀 있다.
$V$ (스킴, 보조정리 \ref{schemes-lemma-section-immersion}) 이상이다.
따라서 절 $s_\varphi \in \mathcal{F}(V \xrightarrow{\varphi} U)$
독특한 존재가 있다절 $s'_\varphi$~에$\mathcal{F}(V \times_X U)$
다시 끌어당기는$s_\varphi$~에 의해$\Gamma_\varphi$
이는 $\Gamma_\varphi$ 이미지의 보수에서 0으로 제한된다.

\medskip\noindent
우리의 지도가 단사적이라는 것을 보여주기 위해 다음과 같이 가정하자
$\sum_{i = 1, \ldots, n} s_{\varphi_i}$는 다음의 요소이다.
위 수식에서 $j_{p!}\mathcal{F}(V)$는 0에 매핑된다.
$j_*\mathcal{F}(V)$에서. 우리의 임무는 그것을 보여주는 것입니다
$\sum_{i = 1, \ldots, n} s_{\varphi_i}$ 0으로 제한된다.
$V$의 에탈 피복 멤버에 대해.
모든 페어별 이퀄라이저 살펴보기
($V$에서 열려 있고 닫혀 있음)의 사상
$\varphi_i : V \to U$ 그리고 $V$에서 로컬로 작업하고 있다.
사상의 이미지를 가정할 수 있다.
$\Gamma_{\varphi_1}, \ldots, \Gamma_{\varphi_n}$
쌍으로 분리되어 있다. 우리의 가정은
$\sum_{i = 1, \ldots, n} s'_{\varphi_i} = 0$
그런 다음 즉시 $s'_{\varphi_i} = 0$라고 결론을 내립니다.
각 $i$에 대해 (이러한 지지대의 불일치로 인해)
절), 어디서$s_{\varphi_i} = 0$모두를 위해$i$원하는대로.
\end{proof}

\begin{lemma}
\label{etale-cohomology-lemma-shriek-equals-star-finite-etale}
$j : U \to X$를 유한하고 에탈이라고 하자. 그러면 지도
$j_! \to j_*$ 의 보조정리 \ref{etale-cohomology-lemma-shriek-into-star-separated-etale}
동형사상이다.
$\Lambda$-가군의 아벨 층 및 층에 있다.
\end{lemma}

\begin{proof} $j_!\mathcal{F} \to j_*\mathcal{F}$이 동형사상 에탈인지 $X$에서 로컬로 확인하면 충분한다. 따라서 $U \to X$는 동형사상의 유한 분리 합집합이라고 가정할 수 있다. 에탈 사상, 보조정리 \ref{etale-lemma-finite-etale-etale-local}를 참조하라. 이 경우에는 증명을 생략한다. \end{proof}

\begin{lemma}
\label{etale-cohomology-lemma-ses-associated-to-open}
$X$를 스킴으로 설정한다. $Z \subset X$를 폐쇄형 하위 구성으로 두고
$U \subset X$ 보완이 된다. $i : Z \to X$ 및 $j : U \to X$를 나타냅니다.
포함사상. 모든아벨 층 $\mathcal{F}$~에$X_\etale$
정식 짧은 완전열이 있다.
$$
0 \to j_!j^{-1}\mathcal{F} \to \mathcal{F} \to i_*i^{-1}\mathcal{F} \to 0
$$
$X_\etale$에 있다.
\end{lemma}

\begin{proof}
함자의 인접성 성질로 지도를 얻습니다.
관련된. 에 대한기하점 $\overline{x}$~에$X$우리는 둘 중 하나를 가지고 있습니다
$\overline{x} \in U$ 이 경우 왼쪽 지도는
줄기의 동형사상이고 $i_*i^{-1}\mathcal{F}$의 줄기이다.
0이거나 $\overline{x} \in Z$인 경우 오른쪽 지도는
줄기의 동형사상이고 $j_!j^{-1}\mathcal{F}$의 줄기이다.
0이다. 여기서는 줄기의 설명을 사용했다.
보조정리 \ref{etale-cohomology-lemma-stalk-pushforward-closed-immersion} 및
명제 \ref{etale-cohomology-proposition-describe-jshriek}.
\end{proof}

\begin{lemma} \label{etale-cohomology-lemma-compatible-shriek-push-finite} 스킴 $$
\xymatrix{
U \ar[d]_g \ar[r]_{j'} & X \ar[d]^f \\
V \ar[r]^j & Y
}
$$의 데카르트 다이어그램을 고려한다. 여기서 $f$는 유한하고 $j$는 열린 몰입이다. 그런 다음 $f_* \circ j'_! = j_! \circ g_*$를 함자 $\textit{Ab}(U_\etale) \to \textit{Ab}(Y_\etale)$로 지정한다. \end{lemma}

\begin{proof}
$\mathcal{F}$를 $\textit{Ab}(U_\etale)$의 객체로 둡니다.
$\overline{y}$를 $Y$의 기하점으로 설정한다.
열린 $V$에서. 그 다음에
$$
(f_*j'_!\mathcal{F})_{\overline{y}} =
\bigoplus\nolimits_{\overline{x},\ f(\overline{x}) = \overline{y}}
(j'_!\mathcal{F})_{\overline{x}} = 0
$$
명제 \ref{etale-cohomology-proposition-finite-higher-direct-image-zero} 및
왜냐하면줄기~의$j'_!\mathcal{F}$~에$\overline{x} \not \in U$
보조정리 \ref{etale-cohomology-lemma-jshriek-open}에 의해 0이 된다. 반면에 우리는
$$
j^{-1}f_*j'_!\mathcal{F} =
g_*(j')^{-1}j'_!\mathcal{F} =
g_*\mathcal{F}
$$
기본 정리 \ref{etale-cohomology-lemma-finite-pushforward-commutes-with-base-change} 및
\ref{etale-cohomology-lemma-jshriek-open}.
따라서 특성화를 통해
$j_!$ 보조정리 \ref{etale-cohomology-lemma-jshriek-open}에서 우리는
$f_*j'_!\mathcal{F} = j_!g_*\mathcal{F}$.
이 신분증에 대한 확인을 생략한다.
$\mathcal{F}$에서 기능적이다.
\end{proof}







\section{구성가능층}
\label{etale-cohomology-section-constructible}

\noindent $X$를 스킴으로 둡니다. {\it 구성 가능한 로컬 폐쇄형 서브스킴} 의 $X$는 로컬로 폐쇄형 서브스킴 $T \subset X$이므로 $T$의 기본 위상 공간은 다음의 구성 가능한 하위 집합이다. $X$. $T, T' \subset X$가 동일한 기본 위상 공간을 가진 로컬로 닫힌 서브스킴인 경우 $T_\etale \cong T'_\etale$는 에탈의 위상 불변 사이트 (정리 \ref{etale-cohomology-theorem-topological-invariance})이다. 따라서 다음 정의에서는 로컬로 닫힌 하위 계획이 축소되었다고 가정할 수 있다.

\begin{definition}
\label{etale-cohomology-definition-constructible}
$X$를 스킴으로 설정한다.
\begin{enumerate}
\item $X_\etale$에 있는 집합의 층은 {\it 구성 가능한다.}
모든 아핀 개방 $U \subset X$에 대해 유한 분해가 존재하는 경우
$U$을 구성 가능한 로컬 폐쇄 하위 구성 $U = \coprod_i U_i$으로
$\mathcal{F}|_{U_i}$는 모든 $i$에 대해 유한한 지역적으로 상수이다.
\item $X_\etale$에 있는 아벨 군의 층은 {\it 구성 가능한다.}
모든 아핀 개방 $U \subset X$에 대해 유한 분해가 존재하는 경우
$U$을 구성 가능한 로컬 폐쇄 하위 구성 $U = \coprod_i U_i$으로
$\mathcal{F}|_{U_i}$는 모든 $i$에 대해 유한한 지역적으로 상수이다.
\item $\Lambda$를 뇌터 환으로 둡니다. $\Lambda$-가군의 층
위의 $X_\etale$는 {\it 구성 가능}이다.
$U \subset X$ 유한 분해가 존재합니다
$U$을 건설 가능한 로컬 폐쇄 하위 계획으로
$U = \coprod_i U_i$ 이렇게
$\mathcal{F}|_{U_i}$은 유한 유형이며 모든 $i$에 대해 지역적으로 상수이다.
\end{enumerate}
\end{definition}

\noindent 이것이 허용되는 정의인 것 같습니다. 스킴의 에탈 사이트를 넘어서는 일반화에 더 쉽게 적합한 대안은 $h_U$, $j_{U!}\mathbf{Z}/n\mathbf{Z}$ 및 $j_{U!}\underline{\Lambda}$로 시작하여 구성가능층을 정의하는 것이었다. 여기서 $U$ $X_\etale$의 모든 준컴팩트 및 준분리 개체에 걸쳐 실행된 다음, 이들을 포함하고 유한 극한 및 여극한 아래에 닫혀 있는 $\Sh(X_\etale)$, $\textit{Ab}(X_\etale)$ 및 $\textit{Mod}(X_\etale, \underline{\Lambda})$의 가장 작은 전체 하위 범주를 취한다. 보조정리 \ref{etale-cohomology-lemma-constructible-abelian} 및 기본정리 \ref{etale-cohomology-lemma-category-constructible-sets}, \ref{etale-cohomology-lemma-category-constructible-abelian} 및 \ref{etale-cohomology-lemma-category-constructible-modules}에서 $X$가 준간소하고 준분리된 경우 동일한 범주를 생성한다. 그러나 일반적으로 이는 동일한 범주를 생성하지 않습니다.

\medskip\noindent 로컬로 닫힌 하위 체계에 의한 스킴의 분리된 통합 분해 $U = \coprod U_i$를 {\it 파티션}이라고 한다. $U$ (위상, 절 \ref{topology-section-stratifications}와 비교).

\begin{lemma}
\label{etale-cohomology-lemma-constructible-quasi-compact-quasi-separated}
$X$를 준간소화하고 준분리된 스킴라고 하자. $\mathcal{F}$
$X_\etale$에서 집합의 층이다. 다음은 동일한다.
\begin{enumerate}
\item $\mathcal{F}$는 구성 가능한다.
\item 다음과 같은 열린 피복 $X = \bigcup U_i$가 존재한다.
$\mathcal{F}|_{U_i}$은 구성 가능하며,
\item 구성 가능한 파티션 $X = \bigcup X_i$이 존재한다.
$\mathcal{F}|_{X_i}$이 유한하도록 로컬로 닫힌 하위 체계
지역적으로 일정한다.
\end{enumerate}
아벨 층 및 층에 대해서도 유사한 설명이 적용된다.
$\Lambda$-가군만약에$\Lambda$~이다뇌터.
\end{lemma}

\begin{proof} (1)가 (2)를 의미하는 것은 분명한다.

\medskip\noindent
(2)를 가정한다. 모든 $x \in X$에 대해 $i$와 아핀 오픈을 찾을 수 있다.
이웃$V_x \subset U_i$~의$x$. 그러므로 우리는 유한함수를 찾을 수 있다
아핀 열린부분 피복 $X = \bigcup V_j$ 각 $j$에 대해
국소적으로 닫혀 있는 유한 분해 $V_j = \coprod V_{j, k}$가 존재한다.
$\mathcal{F}|_{V_{j, k}}$가 유한한 구성 가능한 부분 집합
지역적으로 일정한다. 위상, 보조정리에 의해
\ref{topology-lemma-quasi-compact-open-immersion-constructible-image}
각 $V_{j, k}$은 $X$의 하위 집합으로 구성 가능한다.
위상, 보조정리에 의해
\ref{topology-lemma-constructible-partition-refined-by-stratification}
우리는 구성 가능한 유한 계층화 $X = \coprod X_l$를 찾을 수 있다.
국지적으로 닫힌 지층은 각각
$V_{j, k}$는 $X_l$의 합집합이다. 따라서 (3)가 성립한다.

\medskip\noindent (3)가 성립한다고 가정한다. $U \subset X$를 아핀 개방형으로 둡니다. 그러면 $U \cap X_i$는 $U$의 로컬에서 구성 가능한 닫힌 부분집합입니다(예의 경우 성질, 보조정리 \ref{properties-lemma-locally-constructible}). $U = \coprod U \cap X_i$는 다음의 파티션이다. $U$ 정의 \ref{etale-cohomology-definition-constructible}와 같습니다. 따라서 (1)가 성립한다. \end{proof}

\begin{lemma}
\label{etale-cohomology-lemma-constructible-constructible}
$X$를 준간소화하고 준분리된 스킴라고 하자.
$\mathcal{F}$를 집합, 아벨 군의 층라고 가정한다.
$\Lambda$-가군 ($\Lambda$ 뇌터) - $X_\etale$.
건설 가능한 로컬 폐쇄형 하위 계획이 있는 경우 $T_i \subset X$
() $X = \bigcup T_j$ 및 (b) $\mathcal{F}|_{T_j}$는
구성 가능하면 $\mathcal{F}$도 구성 가능한다.
\end{lemma}

\begin{proof}
먼저, $X$이 준컴팩트하므로 피복이 유한하다고 가정할 수 있다.
스펙트럼에서 위상
(위상, 보조정리 \ref{topology-lemma-constructible-hausdorff-quasi-compact}
그리고
성질, 보조정리
\ref{properties-lemma-quasi-compact-quasi-separated-spectral}).
각 $T_i$은 준컴팩트하고 준분리되어 있음을 확인하세요.
스킴 그 자체로($X$에서 구성 가능하기 때문이다. 자세한 내용은 생략됨)
따라서 우리는 유한 파티션 $T_i = \coprod T_{i, j}$을 찾을 수 있다.
$T_i$의 국부적으로 폐쇄된 건설 가능한 부분
$\mathcal{F}|_{T_{i, j}}$는 유한한 국소 상수이다.
(보조정리 \ref{etale-cohomology-lemma-constructible-quasi-compact-quasi-separated}).
위상, 보조정리 \ref{topology-lemma-constructible-in-constructible}에 의해
$T_{i, j}$는 $X$의 구성 가능한 로컬 폐쇄형 하위 체계라는 것을 알 수 있다.
그런 다음 위상, 보조정리를 적용할 수 있다.
\ref{topology-lemma-constructible-partition-refined-by-stratification}
$X = \bigcup T_{i, j}$로 $X$의 원하는 파티션을 찾습니다.
\end{proof}

\begin{lemma} \label{etale-cohomology-lemma-constructible-local} $X$를 스킴으로 둡니다. 집합, 아벨 군, $\Lambda$-가군($\Lambda$ 뇌터)의 층 구성 가능성을 확인하는 것은 Zariski 로컬에서 수행할 수 있다. $X$. \end{lemma}

\begin{proof} 이 진술은 $X = \bigcup U_i$가 개방형 피복이고 $\mathcal{F}|_{U_i}$가 구성 가능하다면 $\mathcal{F}$도 구성 가능함을 의미한다. $U \subset X$가 개방형이면 $U = \bigcup U \cap U_i$와 $\mathcal{F}|_{U \cap U_i}$가 구성 가능합니다(구성가능층을 개방형으로 제한하는 것은 구성 가능합니다). 보조정리 \ref{etale-cohomology-lemma-constructible-quasi-compact-quasi-separated}로부터 $\mathcal{F}|_U$이 구성 가능한다. 즉, $U$의 적합한 파티션이 존재한다. \end{proof}

\begin{lemma}
\label{etale-cohomology-lemma-pullback-constructible}
$f : X \to Y$를 스킴의 사상으로 설정한다. $\mathcal{F}$인 경우
구성가능층 또는 집합, 아벨 군 또는 $\Lambda$-가군
($\Lambda$ 뇌터) $Y_\etale$에 대해 동일
에 대한 사실이다$f^{-1}\mathcal{F}$~에$X_\etale$.
\end{lemma}

\begin{proof}
보조정리 \ref{etale-cohomology-lemma-constructible-local}에 의해 이는 다음과 같은 경우로 줄어듭니다.
여기서 $X$ 및 $Y$는 유사한다. 에 의해
보조정리 \ref{etale-cohomology-lemma-constructible-quasi-compact-quasi-separated}
구성 가능한 $X$의 유한 분할을 찾는 것으로 충분한다.
$f^{-1}\mathcal{F}$이 로컬로 유한하도록 로컬로 닫힌 하위 구성
그들 각각에 대해 일정한다.
그것을 찾으려면 다음에 적응된 $Y$의 파티션을 뒤로 당기면 된다.
$\mathcal{F}$ 그리고 사용
보조정리 \ref{etale-cohomology-lemma-pullback-locally-constant}.
\end{proof}

\begin{lemma}
\label{etale-cohomology-lemma-constructible-abelian}
$X$를 스킴으로 설정한다.
\begin{enumerate}
\item 범주의 구성가능층의 집합
$\Sh(X_\etale)$ 내부의 유한 극한 및 여극한 아래에서 닫혀 있다.
\item 건설 가능한 아벨 층의 범주는
$\textit{Ab}(X_\etale)$의 약한 Serre 하위 범주.
\item $\Lambda$를 뇌터 환으로 둡니다. 범주의
구성가능층 중 $\Lambda$-가군 의
$X_\etale$은 약한 Serre 하위 범주이다.
$\textit{Mod}(X_\etale, \Lambda)$.
\end{enumerate}
\end{lemma}

\begin{proof}
(3)를 증명한다. 의 기준을 사용하겠습니다.
호몰로지, 보조정리 \ref{homology-lemma-characterize-weak-serre-subcategory}.
$\varphi : \mathcal{F} \to \mathcal{G}$
$\Lambda$-가군의 구성가능층 맵이다.
우리는 $\mathcal{K} = \Ker(\varphi)$와
$\mathcal{Q} = \Coker(\varphi)$은 구성 가능한다.
마찬가지로,
$0 \to \mathcal{F} \to \mathcal{E} \to \mathcal{G} \to 0$
짧은 완전열 의 층 의 $\Lambda$-가군 이다.
$\mathcal{F}$, $\mathcal{G}$로 구성 가능한다. 우리는 보여줘야 해
$\mathcal{E}$는 구성 가능한다.
두 경우 모두 $X$를 멤버로 바꿀 수 있다.
아핀 오픈 피복. 그러므로 우리는 $X$가 유사하다고 가정할 수 있다.
그런 다음 $X$를 유한 요소로 대체할 수 있다.
구성 가능한 로컬 폐쇄형 하위 계획에 의한 $X$ 분할
$\mathcal{F}$ 및 $\mathcal{G}$는 유한 유형이고
지역적으로 일정한다. 따라서 우리는 신청할 수 있습니다
보조정리 \ref{etale-cohomology-lemma-kernel-finite-locally-constant} 결론을 내립니다.

\medskip\noindent (1)와 (2)의 증명은 매우 유사하므로 생략한다. \end{proof}

\begin{lemma}
\label{etale-cohomology-lemma-support-constructible}
$X$를 준간소화하고 준분리된 스킴라고 하자.
\begin{enumerate}
\item $\mathcal{F} \to \mathcal{G}$를 건설 가능한 맵으로 설정한다.
$X_\etale$의 집합의 층. 그런 다음 $x \in X$ 점의 집합
여기서 $\mathcal{F}_{\overline{x}} \to \mathcal{G}_{\overline{x}}$
전사, 각각.\ 단사, 각각.\ 는 주어진 지도와 동형이다.
집합은 $X$에서 구성 가능한다.
\item $\mathcal{F}$를 $X_\etale$에서 구성 가능한 아벨 층으로 설정한다.
$\mathcal{F}$의 지지집합은 구성 가능한다.
\item $\Lambda$를 뇌터 환으로 둡니다.
허락하다$\mathcal{F}$가 되다구성가능층~의$\Lambda$-가군~에$X_\etale$.
$\mathcal{F}$의 지지집합은 구성 가능한다.
\end{enumerate}
\end{lemma}

\begin{proof}
증명 (1).
$X = \coprod X_i$를 국부적으로 폐쇄된 시공 가능으로 $X$의 파티션으로 설정한다.
$\mathcal{F}$ 및 $\mathcal{G}$가 모두 다음과 같은 하위 구성표
부품에 대한 유한 지역 상수(사용
보조정리 \ref{etale-cohomology-lemma-constructible-quasi-compact-quasi-separated}
$\mathcal{F}$ 및 $\mathcal{G}$ 모두에 대해 공통을 선택한다.
정제). 그런 다음 보조정리 \ref{etale-cohomology-lemma-morphism-locally-constant}를 적용한다.
지도를 각 부분으로 제한한다.

\medskip\noindent (2) 및 (3)의 증명이 생략된다. \end{proof}

\noindent 다음 보조정리는 나중에 매우 유용할 것이다. 구성가능층의 범주는 일종의 약한 ``뇌터'' 성질을 가지고 있다고 대략적으로 말한다.

\begin{lemma} \label{etale-cohomology-lemma-colimit-constructible} $X$를 준컴팩트하고 준분리된 스킴라고 하자. $\mathcal{F} = \colim_{i \in I} \mathcal{F}_i$를 집합, 아벨 층 또는 가군층의 층의 필터링된 여극한으로 가정한다. \begin{enumerate} \item $\mathcal{F}$ 및 $\mathcal{F}_i$가 집합의 구성가능층인 경우 Ind-대상 $\mathcal{F}_i$는 본질적으로 값 $\mathcal{F}$로 일정한다. \item $\mathcal{F}$ 및 $\mathcal{F}_i$가 아벨 군의 구성가능층인 경우 Ind-대상 $\mathcal{F}_i$는 본질적으로 값 $\mathcal{F}$로 일정한다. \item $\Lambda$를 뇌터 환으로 둡니다. $\mathcal{F}$ 및 $\mathcal{F}_i$가 $\Lambda$-가군의 구성가능층인 경우 Ind-대상 $\mathcal{F}_i$는 본질적으로 값 $\mathcal{F}$로 일정한다. \end{enumerate} \end{lemma}

\begin{proof}
증명 (1). 더 이상 언급하지 않고 유한 극한을 사용하겠습니다.
및 구성가능층의 여극한은 구성 가능한다.
(보조정리 \ref{etale-cohomology-lemma-kernel-finite-locally-constant}).
각 $i$에 대해 $T_i \subset X$를 점 $x \in X$의 집합으로 둡니다.
여기서 $\mathcal{F}_{i, \overline{x}} \to \mathcal{F}_{\overline{x}}$
전사적이지 않습니다. 왜냐하면 $\mathcal{F}_i$와 $\mathcal{F}$는
구성 가능 $T_i$은 $X$의 구성 가능한 하위 집합이다.
(보조정리 \ref{etale-cohomology-lemma-support-constructible}).
$\mathcal{F}$의 줄기는 유한하므로
그리고 $\mathcal{F} = \colim_{i \in I} \mathcal{F}_i$ 이후 우리는
모두를 위한 것$x \in X$우리는$x \not \in T_i$~을 위한$i$충분히 크다.
$X$는 성질에 의한 스펙트럼 공간이므로, 보조정리
\ref{properties-lemma-quasi-compact-quasi-separated-spectral}
$X$의 구성 가능한 위상은 다음과 같이 준컴팩트한다.
위상, 보조정리 \ref{topology-lemma-constructible-hausdorff-quasi-compact}.
따라서$T_i = \emptyset$~을 위한$i$충분히 크다. 따라서
$\mathcal{F}_i \to \mathcal{F}$는 충분히 큰 $i$의 전사이다.
이제 $\mathcal{F}_i \to \mathcal{F}$가 모든 $i$에 대해 전사라고 가정한다.
$i \in I$을 선택한다. $i' \geq i$의 경우 $S_{i'} \subset X$를 나타냅니다. 집합
포인트 $x$는
$\Im(\mathcal{F}_{i, \overline{x}} \to \mathcal{F}_{\overline{x}})$
의 요소 수와 같지 않습니다.
$\Im(\mathcal{F}_{i, \overline{x}} \to \mathcal{F}_{i', \overline{x}})$.
$\mathcal{F}_i$, $\mathcal{F}_{i'}$ 및 $\mathcal{F}$가
구성 가능 $S_{i'}$은 $X$의 구성 가능한 하위 집합이다.
(세부 사항 생략; 힌트: 보조정리 \ref{etale-cohomology-lemma-support-constructible}를 사용하세요).
$\mathcal{F}_i$ 및 $\mathcal{F}$의 줄기 이후
유한하고 이후 $\mathcal{F} = \colim_{i' \geq i} \mathcal{F}_{i'}$
우리는 모두 그것을 본다$x \in X$우리는$x \not \in S_{i'}$~을 위한$i'$
충분히 크다. 위와 동일한 주장을 통해 우리는 다음과 같은 큰 $i'$를 찾을 수 있다.
그 $S_{i'} = \emptyset$. 따라서 $\mathcal{F}_i \to \mathcal{F}_{i'}$
원하는 대로 $\mathcal{F}$을 통해 인수를 적용한다.

\medskip\noindent 증명 (2). 구성 가능한 아벨 층은 집합의 구성 가능한 층임을 관찰하세요. 따라서 (2) 사례는 (1)에서 이어집니다.

\medskip\noindent
증명 (3). 우리는 더 이상 언급하지 않고 다음의 범주를 사용할 것이다.
$\Lambda$-가군의 구성가능층은 아벨적이다.
(보조정리 \ref{etale-cohomology-lemma-kernel-finite-locally-constant}).
각 $i$에 대해 $\mathcal{Q}_i$를 지도의 여핵으로 설정한다.
$\mathcal{F}_i \to \mathcal{F}$. 그만큼지지집합 $T_i$~의$\mathcal{Q}_i$
의 구성 가능한 하위 집합이다.$X$~처럼$\mathcal{Q}_i$건설 가능하다
(보조정리 \ref{etale-cohomology-lemma-support-constructible}).
$\mathcal{F}$의 줄기는 유한하므로 $\Lambda$-가군
그리고 $\mathcal{F} = \colim_{i \in I} \mathcal{F}_i$ 이후 우리는
모두를 위한 것$x \in X$우리는$x \not \in T_i$~을 위한$i$충분히 크다.
$X$는 성질에 의한 스펙트럼 공간이므로, 보조정리
\ref{properties-lemma-quasi-compact-quasi-separated-spectral}
$X$의 구성 가능한 위상은 다음과 같이 준컴팩트한다.
위상, 보조정리 \ref{topology-lemma-constructible-hausdorff-quasi-compact}.
따라서$T_i = \emptyset$~을 위한$i$충분히 크다. 이것이 첫번째로 증명된다
역설. 두 번째로 이제 다음과 같이 가정한다.
$\mathcal{F}_i \to \mathcal{F}$은 모든 $i$에 대한 전사이다.
$i \in I$을 선택한다. $i' \geq i$의 경우 $\mathcal{K}_{i'}$를 나타냅니다.
이미지$\Ker(\mathcal{F}_i \to \mathcal{F})$~에$\mathcal{F}_{i'}$.
그만큼지지집합 $S_{i'}$~의$\mathcal{K}_{i'}$
의 구성 가능한 하위 집합이다.$X$~처럼$\mathcal{K}_{i'}$시공 가능한다.
$\Ker(\mathcal{F}_i \to \mathcal{F})$의 줄기 이후
유한 $\Lambda$-가군이고 이후
$\mathcal{F} = \colim_{i' \geq i} \mathcal{F}_{i'}$ 알겠습니다
모두를 위한 것$x \in X$우리는$x \not \in S_{i'}$~을 위한$i'$충분히 크다.
위와 동일한 주장을 통해 우리는 다음과 같은 큰 $i'$를 찾을 수 있다.
그 $S_{i'} = \emptyset$. 따라서 $\mathcal{F}_i \to \mathcal{F}_{i'}$
원하는 대로 $\mathcal{F}$을 통해 인수를 적용한다.
\end{proof}

\begin{lemma}
\label{etale-cohomology-lemma-tensor-product-constructible}
$X$를 스킴으로 설정한다. $\Lambda$를 뇌터 환으로 설정한다.
$\Lambda$-가군의 두 구성가능층의 텐서곱
$X_\etale$에서 $\Lambda$-가군의 구성가능층이다.
\end{lemma}

\begin{proof} 질문은 즉시 $X$이 연관성이 있는 경우로 축소된다. 건설 가능한 국지적으로 닫힌 지층이 있는 $X$의 두 파티션은 동일한 유형의 공통 세분화를 가지며 당김은 텐서곱으로 통근하므로 보조정리 \ref{etale-cohomology-lemma-tensor-product-locally-constant}로 줄이다. \end{proof}

\begin{lemma}
\label{etale-cohomology-lemma-tensor-constructible}
$\Lambda \to \Lambda'$를 뇌터 환의 동형이라고 하자.
$X$를 스킴으로 설정한다. $\mathcal{F}$를 구성 가능한 것으로 둡니다.
층~의$\Lambda$-가군~에$X_\etale$. 그 다음에
$\mathcal{F} \otimes_{\underline{\Lambda}} \underline{\Lambda'}$
$\Lambda'$-가군의 구성가능층이다.
\end{lemma}

\begin{proof} 생략한다. 힌트: 지역적으로 유사하게 동일한 계층화를 사용할 수 있다. \end{proof}










\section{사상} \label{etale-cohomology-section-stratify-morphisms}의 보조 기본정리

\noindent 구성가능층의 기능성 성질을 증명하는 데 유용한 일부 기본형.

\begin{lemma}
\label{etale-cohomology-lemma-etale-stratified-finite}
$U \to X$를 준간밀하고 준분리된 에탈 사상으로 하자.
스킴 (예의 경우 뇌터 스킴의 에탈 사상). 그럼 거기
로컬로 닫힌 구성 가능한 파티션 $X = \coprod_i X_i$이 존재한다.
$X_i \times_X U \to X_i$가 모든 $i$에 대해 유한 에탈이 되도록 하는 하위 계획이다.
\end{lemma}

\begin{proof} $U \to X$가 분리된 경우 사상, 보조정리 \ref{more-morphisms-lemma-stratify-flat-fp-qf}에 대한 자세한 내용이다. 일반적으로 $X$는 유사하다고 가정할 수 있다. 유한 아핀 개방형 피복 $U = \bigcup U_j$를 선택한다. 모든 사상 $U_j \to X$ 및 $U_j \cap U_{j'} \to X$에 이전 사례를 적용하고 결과 파티션의 공통 세분화 $X = \coprod X_i$를 선택한다. 파티션을 더욱 구체화한 후 $X_i$도 유사하다고 가정할 수 있다. $i$ 및 집합 $V = U \times_X X_i$를 수정하세요. 사상 $V_j = U_j \times_X X_i \to X_i$ 및 $V_{jj'} = (U_j \cap U_{j'}) \times_X X_i \to X_i$는 유한 에탈이다. 따라서 $V_j$ 및 $V_{jj'}$는 아핀 스킴이고 $V_{jj'} \subset V_j$는 닫혀 있으면서도 열려 있습니다($V_{jj'} \to X_i$가 적절하므로 사상, 보조정리 \ref{morphisms-lemma-image-proper-scheme-closed}가 적용됩니다). 그러면 $V = \bigcup V_j$는 $\mathcal{O}(V_j) \to \mathcal{O}(V_{jj'})$가 전사이기 때문에 분리된다. 스킴, 보조정리 \ref{schemes-lemma-characterize-separated}를 참조하라. 따라서 이전 사례는 $V \to X_i$에 적용되며 필요한 경우 파티션을 더 세분화할 수 있습니다(실제로는 그렇지 않지만 필요하지는 않습니다). \end{proof}

\noindent 뇌터의 경우 앞의 보조정리를 뇌터 귀납법으로 증명할 수 있고 다음은 재미있는 보조정리를 증명할 수 있다.

\begin{lemma}
\label{etale-cohomology-lemma-generically-finite}
$f: X \to Y$를 준컴팩트한 스킴의 사상으로 두고,
준분리되고 국부적으로 유한 유형이다. $\eta$이 일반 포인트인 경우
~의기약구성 요소$Y$그렇게$f^{-1}(\eta)$그렇다면 유한하다
열려 있는 곳이 있다$V \subset Y$포함하는$\eta$그렇게
$f^{-1}(V) \to V$은 유한한다.
\end{lemma}

\begin{proof} 이것은 사상, 보조정리 \ref{morphisms-lemma-generically-finite}이다. \end{proof}

\noindent 다음 보조정리의 서술문은 조금 더 강화될 수 있다.

\begin{lemma}
\label{etale-cohomology-lemma-decompose-quasi-finite-morphism}
$f : Y \to X$을 준유한이고 유한하게 표현한다고 가정한다.
아핀 스킴의 사상.
\begin{enumerate}
\item 아핀 스킴 $X' \to X$의 전사 사상이 있고
폐쇄형 하위 계획 $Z' \subset Y' = X' \times_X Y$
\begin{enumerate}
\item $Z' \subset Y'$는 두꺼워지고,
\item $Z' \to X'$는 유한 에탈 사상이다.
\end{enumerate}
\item 유한한 파티션 $X = \coprod X_i$이 존재한다.
국부적으로 닫혀 있고, 구성 가능하며, 아핀 지층, 국부적으로 전사 유한
무료 사상 $X'_i \to X_i$
$Y'_i = X'_i \times_X Y \to X'_i$는 다음과 동형이다.
$\coprod_{j = 1}^{n_i} (X'_i)_{red} \to (X'_i)_{red}$ 일부 $n_i$에 대해.
\end{enumerate}
\end{lemma}

\begin{proof} $X' = \coprod X'_i$를 설정하면 (2)가 (1)를 의미함을 알 수 있다. $X = \Spec(A)$ 및 $Y = \Spec(B)$를 씁니다. $A$를 유한 유형 $\mathbf{Z}$-대수 $A_i$의 여과 여극한으로 씁니다. $B$는 유한 표현의 $A$-대수이므로 $0 \in I$과 유한 유형 환 맵 $A_0 \to B_0$가 $B = \colim B_i$와 $B_i = A_i \otimes_{A_0} B_0$가 있음을 알 수 있다. 대수학을 참조하라. 보조정리 \ref{algebra-lemma-colimit-category-fp-algebras}. $i$가 충분히 큰 경우 $A_i \to B_i$는 준 유한함을 알 수 있다. 극한, 보조정리 \ref{limits-lemma-descend-quasi-finite}를 참조하라. 따라서 우리는 $\mathbf{Z}$에 대한 유한 유형 대수의 경우로 축소하고, 특히 뇌터 경우로 축소한다. (자세한 내용은 생략한다.)

\medskip\noindent
$X$ 및 $Y$ 뇌터를 가정한다. 이 경우 현지 폐쇄
$X$의 하위 집합을 구성할 수 있다. 보조정리 \ref{etale-cohomology-lemma-generically-finite}에 의해
그리고 뇌터 유도는 다음과 같습니다.
유한한 파티션이 있다$X = \coprod X_i$~의$X$
$Y \times_X X_i \to X_i$이 유한하도록 국부적으로 닫힌 지층에 의해.
우리는 이 분할을 개선하여 아핀 지층을 얻을 수 있다.
따라서 $X$을 $X' = \coprod X_i$로 바꾼 후 다음과 같이 가정할 수 있다.
$Y \to X$은 유한한다.

\medskip\noindent
$X$ 및 $Y$ 뇌터 및 $Y \to X$ 유한하다고 가정한다.
염기 변경 후 전사에 의해 (2)를 증명할 수 있다고 가정하면,
편평하고 준 유한 사상 $U \to X$. 따라서 우리는 파티션을 가지고 있습니다
$U = \coprod U_i$ 및 유한 로컬 자유 사상 $U'_i \to U_i$
$U'_i \times_X Y \to U'_i$는 다음과 동형이다.
$\coprod_{j = 1}^{n_i} (U'_i)_{red} \to (U'_i)_{red}$ 일부 $n_i$에 대해.
그러면 이전 단락의 논증을 통해 다음을 찾을 수 있다.
다음과 같이 로컬로 닫힌 아핀 계층이 있는 파티션 $X = \coprod X_j$
$X_j \times_X U_i \to X_j$은 모든 $i, j$에 대해 유한한다. 에 의해
사상, 보조정리 \ref{morphisms-lemma-finite-flat}
각 $X_j \times_X U_i \to X_j$는 로컬에서 유한하게 사용 가능한다.
따라서 $X_j \times_X U'_i \to X_j$는 유한한 지역적 자유이다.
(사상, 보조정리 \ref{morphisms-lemma-composition-finite-locally-free}).
$X = \coprod X_j$ 및 $X_j' = \coprod_i X_j \times_X U'_i$
$Y \to X$에 대한 솔루션이다. 그러므로 증명하면 충분하다
전사 편평 준유한 이후의 결과(뇌터 경우)
베이스 변경.

\medskip\noindent
사상, 보조정리 \ref{morphisms-lemma-massage-finite} 적용
$Y$는 폐쇄형 하위 체계라고 가정할 수 있다.
아핀 스킴 $Z$ 이는 (집합 이론적으로는 유한 합집합)
$Z = \bigcup_{i \in I} Z_i$ 동형적으로 매핑된 폐쇄형 하위 구성표
$X$로 변경된다. 이 경우 우리는 $X = \coprod X_j$의 유한한 파티션을 찾을 것이다.
작동하는 아핀 국소 폐쇄 지층(즉, $X'_j = X_j$)을 사용한다.
집합 $T_i = Y \cap Z_i$. 이것은 $X$의 폐쇄형 하위 체계이다.
$X$는 뇌터이므로 $X = \coprod X_j$의 유한 분할을 찾을 수 있다.
로컬로 폐쇄된 하위 계획을 사용하여 각
$X_j \times_X T_i$는 (집합 이론적으로) 계층 $X_j \times_X Z_i$의 결합이다.
$X$를 $X_j$로 바꾸면 $I = I_1 \amalg I_2$라고 가정할 수 있다.
~와 함께$Z_i \subset Y$~을 위한$i \in I_1$그리고$Z_i \cap Y = \emptyset$~을 위한
$i \in I_2$. $Z$를 $\bigcup_{i \in I_1} Z_i$로 바꾸면
$Y = Z$라고 가정할 수 있다.
마지막으로 $X$를 파티션 멤버로 다시 바꿀 수 있다.
위와 같이 모든 $i, i' \subset I$ 교차로에 대해
$Z_i \cap Z_{i'}$는 비어 있거나 (집합 이론적으로) 동일한다.
$Z_i$ 및 $Z_{i'}$로 변경된다. 이는 이론적으로 $Y$가 (집합이라는 것을 의미합니다)
우리가 보여주고 싶었던 것은 $Z_i$의 분리된 합집합과 같습니다.
\end{proof}








\section{구성가능층} \label{etale-cohomology-section-more-constructible}에 대해 자세히 알아보기

\noindent
$\Lambda$를 뇌터 환으로 설정한다. $X$를 스킴으로 설정한다.
우리는 종종 $X_\etale$을 고리 모양의 사이트로 간주한다.
환 $\underline{\Lambda}$의 층. 아벨 층의 경우
우리는 종종 $\Lambda = \mathbf{Z}/n\mathbf{Z}$를 사용하여 적절한
정수 $n$.

\begin{lemma}
\label{etale-cohomology-lemma-jshriek-constructible}
$j : U \to X$를 준컴팩트의 에탈 사상으로 두고
준분리된 스킴.
\begin{enumerate}
\item 층 $h_U$는 집합의 구성가능층이다.
\item 층 $j_!\underline{M}$는 구성 가능한 아벨 층이다.
유한 아벨 군 $M$의 경우.
\item $\Lambda$가 뇌터 환이고 $M$가 유한 $\Lambda$-가군인 경우,
그러면 $j_!\underline{M}$는 $\Lambda$-가군의 구성가능층이다.
$X_\etale$에 있다.
\end{enumerate}
\end{lemma}

\begin{proof}
보조정리 \ref{etale-cohomology-lemma-etale-stratified-finite}에 의해 파티션이 있습니다
$\coprod_i X_i$ $\pi_i : j^{-1}(X_i) \to X_i$는 유한 에탈이다.
$h_U$에서 $X_i$로의 제한은 $h_{j^{-1}(X_i)}$이며 이는 유한한다.
보조정리 \ref{etale-cohomology-lemma-characterize-finite-locally-constant}에 의해 지역적으로 상수이다.
(2) 및 (3) 사례의 경우
$$
j_!(\underline{M})|_{X_i} =
\pi_{i!}(\underline{M}) =
\pi_{i*}(\underline{M})
$$
기본 정리 \ref{etale-cohomology-lemma-shriek-base-change} 및
\ref{etale-cohomology-lemma-shriek-equals-star-finite-etale}.
따라서 $\pi : Y \to X$ 유한 에탈에 대해 보조정리를 표시하면 충분한다.
이것은 보조정리 \ref{etale-cohomology-lemma-pushforward-locally-constant}이다.
\end{proof}

\begin{lemma}
\label{etale-cohomology-lemma-torsion-colimit-constructible}
$X$를 준간소화하고 준분리된 스킴라고 하자.
\begin{enumerate}
\item $\mathcal{F}$를 $X_\etale$에서 집합의 층으로 설정한다.
그러면 $\mathcal{F}$는 구성 가능한 여과 여극한이다.
집합의 층.
\item $\mathcal{F}$를 $X_\etale$의 비틀림 아벨 층이라고 가정한다.
그러면 $\mathcal{F}$는 구성 가능한 아벨 층의 여과 여극한이다.
\item $\Lambda$는 뇌터 환 및 $\mathcal{F}$는 층
$\Lambda$-가군의 $X_\etale$. 그 다음에
$\mathcal{F}$는 구성가능층의 여과 여극한이다.
$\Lambda$-가군.
\end{enumerate}
\end{lemma}

\begin{proof} $\mathcal{B}$를 $X_\etale$의 준컴팩트하고 준분리된 객체의 집합이라고 가정한다. 가군에 의해 사이트, 보조정리 \ref{sites-modules-lemma-module-filtered-colimit-constructibles} 집합의 모든 층은 여과 여극한 형식의 층이다. $$
\text{코이퀄라이저}\left(
\xymatrix{
\coprod\nolimits_{j = 1, \ldots, m} h_{V_j}
\ar@<1ex>[r] \ar@<-1ex>[r] &
\coprod\nolimits_{i = 1, \ldots, n} h_{U_i}
}
\right)
$$는 $V_j$ 및 $U_i$ $X_\etale$의 준간밀하고 준분리된 객체이다. 보조정리 \ref{etale-cohomology-lemma-jshriek-constructible} 및 \ref{etale-cohomology-lemma-constructible-abelian}에 의해 이러한 공등화기는 구성 가능한다. 이는 (1)를 증명한다.

\medskip\noindent $\Lambda$를 뇌터 환으로 둡니다. 가군에 의해 사이트, 보조정리 \ref{sites-modules-lemma-module-filtered-colimit-constructibles} $\Lambda$-가군 $\mathcal{F}$는 여과 여극한의 가군이다. $$
\Coker\left(
\bigoplus\nolimits_{j = 1, \ldots, m} j_{V_j!}\underline{\Lambda}_{V_j}
\longrightarrow
\bigoplus\nolimits_{i = 1, \ldots, n} j_{U_i!}\underline{\Lambda}_{U_i}
\right)
$$를 $V_j$ 및 $U_i$ $X_\etale$의 준간밀하고 준분리된 객체로 구성한다. 보조정리 \ref{etale-cohomology-lemma-jshriek-constructible} 및 \ref{etale-cohomology-lemma-constructible-abelian}에 의해 이러한 코커널은 구성 가능한다. 이는 (3)를 증명한다.

\medskip\noindent 증명 (2). 먼저 $\mathcal{F} = \bigcup \mathcal{F}[n]$를 작성한다. 여기서 $\mathcal{F}[n]$는 $n$ 비틀림 서브쉬프이다. 그러면 $\mathcal{F}[n]$를 $\mathbf{Z}/n\mathbf{Z}$-가군의 층으로 보고 (3)를 적용할 수 있다. \end{proof}

\begin{lemma}
\label{etale-cohomology-lemma-check-constructible}
$f : X \to Y$를 준 컴팩트의 대사 사상으로 두고
준분리된 스킴.
\begin{enumerate}
\item $\mathcal{F}$를 $Y_\etale$에서 집합의 층으로 설정한다. 그러면 $\mathcal{F}$
건설 가능합니다 필요충분조건은 $f^{-1}\mathcal{F}$ 건설 가능한다.
\item $\mathcal{F}$를 $Y_\etale$에서 아벨 층으로 설정한다. 그러면 $\mathcal{F}$
건설 가능합니다 필요충분조건은 $f^{-1}\mathcal{F}$ 건설 가능한다.
\item $\Lambda$를 뇌터 환으로 둡니다.
허락하다$\mathcal{F}$이다층~의$\Lambda$-가군~에$Y_\etale$.
그런 다음 $\mathcal{F}$는 구성 가능한다. 필요충분조건은 $f^{-1}\mathcal{F}$
시공 가능한다.
\end{enumerate}
\end{lemma}

\begin{proof}
보조정리 \ref{etale-cohomology-lemma-pullback-constructible}에서 한 가지 의미가 나옵니다.
반대로, $f^{-1}\mathcal{F}$는 구성 가능하다고 가정한다.
$\mathcal{F} = \colim \mathcal{F}_i$를 다음과 같이 씁니다.
여과 여극한의 구성가능층 (집합, 아벨 군 또는 가군)
보조정리 \ref{etale-cohomology-lemma-torsion-colimit-constructible}를 사용한다.
$f^{-1}$는 왼쪽 수반이므로 여극한으로 통근한다.
(범주, 보조정리 \ref{categories-lemma-adjoint-exact}) 그리고 우리는
$f^{-1}\mathcal{F} = \colim f^{-1}\mathcal{F}_i$.
보조정리 \ref{etale-cohomology-lemma-colimit-constructible}에 의해 우리는
$f^{-1}\mathcal{F}_i \to f^{-1}\mathcal{F}$
충분히 큰 모든 $i$에 대해 전사이다.
$f$는 전사이므로 (줄기를 보면 다음과 같이 결론을 내립니다.
보조정리 \ref{etale-cohomology-lemma-stalk-pullback} 및
정리 \ref{etale-cohomology-theorem-exactness-stalks})
$\mathcal{F}_i \to \mathcal{F}$는 모든 $i$에 대해 충분히 큰 전사이다.
따라서 $\mathcal{F}$은 구성가능층 $\mathcal{G}$의 몫이다.
주장을 다시 한 번 적용해 보면
$\mathcal{G} \times_\mathcal{F} \mathcal{G}$ 또는
$\mathcal{G} \to \mathcal{F}$의 핵
우리는 $f^{-1}$가 정확하고 범주를 사용하여 결론을 내립니다.
구성가능층 (집합, 아벨 군 또는 가군)는
내부의 유한 (co)극한 또는 (co)커널 아래에 보존됨
$\Sh(Y_\etale)$, $\Sh(X_\etale)$, $\textit{Ab}(Y_\etale)$,
$\textit{Ab}(X_\etale)$, $\textit{Mod}(Y_\etale, \Lambda)$ 및
$\textit{Mod}(X_\etale, \Lambda)$, 참조
보조정리 \ref{etale-cohomology-lemma-constructible-abelian}.
\end{proof}

\begin{lemma}
\label{etale-cohomology-lemma-pushforward-constructible}
$f : X \to Y$를 스킴의 유한 에탈 사상라고 하자. $\Lambda$를
뇌터 환. $\mathcal{F}$가 집합의 구성가능층인 경우,
구성가능층 중 아벨 군 또는 구성가능층 중
$\Lambda$-가군 $X_\etale$의 경우에도 마찬가지이다.
$f_*\mathcal{F}$ - $Y_\etale$.
\end{lemma}

\begin{proof} 보조정리 \ref{etale-cohomology-lemma-constructible-local}에 의해 이 Zariski를 $Y$에서 로컬로 확인하는 것으로 충분하며 보조정리 \ref{etale-cohomology-lemma-check-constructible}에 의해 $Y$를 에탈 덮개로 대체할 수 있습니다(구성 $f_*$ 에탈 국소화)로 통근한다. 유한 에탈 사상은 동형사상의 분리된 합집합에 에탈 국소적으로 동형이다. 에탈 사상, 보조정리 \ref{etale-lemma-finite-etale-etale-local}을 참조하라. 따라서 집합의 층의 경우, 보조정리는 $\mathcal{F}_i$, $i = 1, \ldots, n$가 집합의 구성가능층인 경우 $\prod_{i = 1, \ldots, n} \mathcal{F}_i$도 마찬가지임을 의미한다. 이것은 분명한다. 아벨 군 및 가군의 층도 마찬가지이다. \end{proof}

\begin{lemma}
\label{etale-cohomology-lemma-category-constructible-sets}
$X$를 준간소화하고 준분리된 스킴라고 하자. 범주의
집합의 구성가능층은 $\Sh(X_\etale)$의 전체 하위 범주이다.
보조 이퀄라이저인 층 $\mathcal{F}$로 구성됨
$$
\xymatrix{
\mathcal{F}_1
\ar@<1ex>[r] \ar@<-1ex>[r]
&
\mathcal{F}_0 \ar[r]
&
\mathcal{F}}
$$
$\mathcal{F}_i$, $i = 0, 1$는 다음의 층의 유한한 부산물이다.
형태$h_U$~와 함께$U$준컴팩트하고 준분리형
$X_\etale$의 객체이다.
\end{lemma}

\begin{proof}
보조정리 \ref{etale-cohomology-lemma-torsion-colimit-constructible}의 증명에서
우리는 이 형식의 층이 구성 가능하다는 것을 확인했다.
반대로, 모든 구성가능층에 대해
집합 $\mathcal{F}$ 전사 $\mathcal{F}_0 \to \mathcal{F}$를 찾을 수 있다.
보조정리에서와 같이 $\mathcal{F}_0$을 사용한다. 그런 다음 우리는 추측을 찾습니다.
$\mathcal{F}_1 \to \mathcal{F}_0 \times_\mathcal{F} \mathcal{F}_0$
왜냐하면 후자는 보조정리 \ref{etale-cohomology-lemma-constructible-abelian}로 구성 가능하기 때문이다.

\medskip\noindent
위상, 보조정리에 의해
\ref{topology-lemma-constructible-partition-refined-by-stratification}
우리는 유한한 계층화를 선택할 수 있습니다
$X = \coprod_{i \in I} X_i$ $\mathcal{F}$는 지역적으로 유한한다.
각 지층에서 일정한다. 그 결과를 귀납법으로 증명하겠습니다.
$I$의 카디널리티이다. $i \in I$를 최소 요소로 둡니다.
$I$의 부분 순서이다. 그런 다음 $X_i \subset X$가 닫힙니다.
인덕션으로는 준컴팩트형과 준분리형이 유한하게 많이 존재한다.
사물$U_\alpha$~의$(X \setminus X_i)_\etale$그리고 전사
지도 $\coprod h_{U_\alpha} \to \mathcal{F}|_{X \setminus X_i}$.
이것들이 지도를 결정한다
$$
\coprod h_{U_\alpha} \to \mathcal{F}
$$
이는 $X \setminus X_i$로 제한한 후 전사이다. 에 의해
보조정리 \ref{etale-cohomology-lemma-characterize-finite-locally-constant}
우리는 그것을 본다$\mathcal{F}|_{X_i} = h_V$어떤 사람들에게는스킴 $V$유한한 에탈
$X_i$ 이상. $\overline{v}$를 $V$의 기하점으로 설정한다.
$\overline{x} \in X_i$ 이상. 우리는 $\overline{v}$를 다음과 같이 생각할 수 있다.
줄기 $\mathcal{F}_{\overline{x}} = V_{\overline{x}}$의 요소이다.
따라서 우리는 에탈 이웃 $(U, \overline{u})$을 찾을 수 있다.
$\overline{x}$ 및 절 $s \in \mathcal{F}(U)$의 줄기
$\overline{x}$는 $\overline{v}$을 제공한다. $s$를 지도로 생각하기
$s : h_U \to \mathcal{F}$, $X_i$로 제한하여 사상을 얻습니다.
$s|_{X_i} : U \times_X X_i \to V$ $X_i$에 대해 $\overline{u}$를 매핑함
$\overline{v}$로 변경된다. $V$는 준컴팩트이므로(닫힌 공간에 대해 유한함)
하위 계획$X_i$준컴팩트의스킴 $X$) 유한수
$s^{(1)}, \ldots, s^{(m)}$ 중 절 중 $\mathcal{F}$ 이상
$U^{(1)}, \ldots, U^{(m)}$는 공동으로 결정한다.
전사 지도
$$
\coprod s^{(j)}|_{X_i} : \coprod U^{(j)} \times_X X_i \longrightarrow V
$$
그러면 우리는 추측을 얻습니다.
$$
\coprod h_{U_\alpha} \amalg \coprod h_{U^{(j)}} \to \mathcal{F}
$$
원하는대로.
\end{proof}

\begin{lemma}
\label{etale-cohomology-lemma-category-constructible-modules}
$X$를 준간소화하고 준분리된 스킴라고 하자. $\Lambda$
뇌터 환이어야 한다. 구성가능층의 범주
$\Lambda$-가군은 정확히 범주 형식의 가군이다.
$$
\Coker\left(
\bigoplus\nolimits_{j = 1, \ldots, m} j_{V_j!}\underline{\Lambda}_{V_j}
\longrightarrow
\bigoplus\nolimits_{i = 1, \ldots, n} j_{U_i!}\underline{\Lambda}_{U_i}
\right)
$$
$V_j$ 및 $U_i$ 준컴팩트 및 준분리 객체 사용
$X_\etale$. 실제로 $U_i$ 및 $V_j$ 유사성을 가정할 수도 있다.
\end{lemma}

\begin{proof}
보조정리 \ref{etale-cohomology-lemma-torsion-colimit-constructible}의 증명에서
우리는 이 형식의 가군이 구성 가능하다는 것을 보았습니다. 이후
구성 가능한 가군의 범주는 아벨적이다.
(보조정리 \ref{etale-cohomology-lemma-constructible-abelian})
구성 가능한 가군 $\mathcal{F}$가 주어지면 이를 증명하는 것으로 충분한다.
추측이있다
$$
\bigoplus\nolimits_{i = 1, \ldots, n} j_{U_i!}\underline{\Lambda}_{U_i}
\longrightarrow \mathcal{F}
$$
일부 아핀 객체의 경우$U_i$~에$X_\etale$. 에 의해
가군 사이트, 보조정리에
\ref{sites-modules-lemma-module-filtered-colimit-constructibles}
추측이있다
$$
\Psi :
\bigoplus\nolimits_{i \in I} j_{U_i!}\underline{\Lambda}_{U_i}
\longrightarrow
\mathcal{F}
$$
$U_i$ 아핀과 직합이 무한할 가능성이 있는 경우
인덱스 집합 $I$. 모든 유한 부분집합 $I' \subset I$ 집합
$$
T_{I'} = \text{Supp}(\Coker(
\bigoplus\nolimits_{i \in I'} j_{U_i!}\underline{\Lambda}_{U_i}
\longrightarrow \mathcal{F}))
$$
구성가능층의 바로 정의에 의해 집합 $T_{I'}$
는 $X$의 구성 가능한 하위 집합이다. 우리는 $T_{I'} = \emptyset$를 보여주고 싶다.
일부 $I'$의 경우. 모든 줄기 $\mathcal{F}_{\overline{x}}$는
유한 유형 $\Lambda$-가군 그리고 $\Psi$는 전사이므로
모든$x \in X$있다$I'$그렇게$x \not \in T_{I'}$.
즉, 우리는
$\emptyset = \bigcap_{I' \subset I\text{한정된}} T_{I'}$. 부터
$X$는 성질, 보조정리에 의한 스펙트럼 공간이다.
\ref{properties-lemma-quasi-compact-quasi-separated-spectral}
$X$의 구성 가능한 위상은 다음과 같이 준컴팩트한다.
위상, 보조정리 \ref{topology-lemma-constructible-hausdorff-quasi-compact}.
따라서 $T_{I'} = \emptyset$ 일부 $I' \subset I$ 유한
원하는대로.
\end{proof}

\begin{lemma}
\label{etale-cohomology-lemma-category-constructible-abelian}
$X$를 준간소화하고 준분리된 스킴라고 하자. 범주의
구성가능 아벨 층 은 정확히 아벨리안의 범주 입니다
층 형식의
$$
\Coker\left(
\bigoplus\nolimits_{j = 1, \ldots, m}
j_{V_j!}\underline{\mathbf{Z}/m_j\mathbf{Z}}_{V_j}
\longrightarrow
\bigoplus\nolimits_{i = 1, \ldots, n}
j_{U_i!}\underline{\mathbf{Z}/n_i\mathbf{Z}}_{U_i}
\right)
$$
$V_j$ 및 $U_i$ 준컴팩트 및 준분리 객체 사용
$X_\etale$ 및 $m_j$, $n_i$ 양의 정수.
실제로 $U_i$ 및 $V_j$ 유사성을 가정할 수도 있다.
\end{lemma}

\begin{proof}
이는 보조정리 \ref{etale-cohomology-lemma-category-constructible-modules}에서 따온 것이다.
$\Lambda = \mathbf{Z}/n\mathbf{Z}$로 적용됨
그리고 $X$은 준컴팩트하기 때문에 모든 구성 가능한
아벨 층은 어떤 양의 정수 $n$에 의해 소멸됩니다(자세한 내용은 생략됨).
\end{proof}

\begin{lemma}
\label{etale-cohomology-lemma-constructible-is-compact}
$X$를 준간소화하고 준분리된 스킴라고 하자. $\Lambda$를
뇌터 환. $\mathcal{F}$를 집합의 구성가능층으로 설정한다.
그룹 또는$\Lambda$-가군~에$X_\etale$. 허락하다
$\mathcal{G} = \colim \mathcal{G}_i$는 여과 여극한의 층
집합, 아벨 군 또는 $\Lambda$-가군. 그 다음에
$$
\Mor(\mathcal{F}, \mathcal{G}) = \colim \Mor(\mathcal{F}, \mathcal{G}_i)
$$
범주의 층의 집합, 아벨 군 또는 $\Lambda$-가군의
$X_\etale$.
\end{lemma}

\begin{proof} 집합의 층의 경우이다. 보조정리 \ref{etale-cohomology-lemma-category-constructible-sets}에 의해 $h_U$에 대한 보조정리를 증명하는 것으로 충분한다. 여기서 $U$는 $X_\etale$의 준컴팩트하고 준분리된 개체이다. $\Mor(h_U, \mathcal{G}) = \mathcal{G}(U)$를 기억하세요. 따라서 결과는 사이트, 보조정리 \ref{sites-lemma-directed-colimits-sections}의 결과이다.

\medskip\noindent
아벨 층 또는 가군층의 경우 결과는 다음과 같습니다.
같은 방법으로 이용해서
기본형 \ref{etale-cohomology-lemma-category-constructible-abelian} 및
\ref{etale-cohomology-lemma-category-constructible-modules}.
아벨 층의 경우 다음을 추가한다.
$\Mor(j_{U!}\underline{\mathbf{Z}/n\mathbf{Z}}, \mathcal{G})$
는$n$- 비틀림 요소$\mathcal{G}(U)$.
\end{proof}

\begin{lemma}
\label{etale-cohomology-lemma-finite-pushforward-constructible}
$f : X \to Y$를 스킴의 유한하고 유한하게 제시된 사상라고 가정한다.
$\Lambda$를 뇌터 환으로 설정한다. $\mathcal{F}$가 구성 가능한 경우
층의 집합, 아벨 군 또는 $\Lambda$-가군의 $X_\etale$,
그러면 $f_*\mathcal{F}$도 마찬가지이다.
\end{lemma}

\begin{proof} $X$과 $Y$가 보조정리 \ref{etale-cohomology-lemma-constructible-local}에 의해 유사할 때 이를 증명하면 충분한다. 기본정리 \ref{etale-cohomology-lemma-finite-pushforward-commutes-with-base-change} 및 \ref{etale-cohomology-lemma-check-constructible}에 따라 $X$에 대한 아핀 스킴 전사로 기본 변경이 가능한다. 보조정리 \ref{etale-cohomology-lemma-decompose-quasi-finite-morphism}에 의해 이는 유한 에탈 사상의 경우로 축소됩니다(두꺼워지면 에탈 토포스와 심지어 작은 에탈 사이트와 동등하게 되기 때문이다. 정리를 참조하라. \ref{etale-cohomology-theorem-topological-invariance}). 유한 에탈 경우는 보조정리 \ref{etale-cohomology-lemma-pushforward-constructible}이다. \end{proof}

\begin{lemma}
\label{etale-cohomology-lemma-category-is-colimit}
$X = \lim_{i \in I} X_i$를 극한으로 지정한다.
아핀 전이 사상이 있는 스킴 시스템.
$X_i$은 준컴팩트하고 준분리되어 있다고 가정한다.
모든 $i \in I$에 대해.
\begin{enumerate}
\item 범주의 구성가능층의 집합의 $X_\etale$
은 여극한 의 범주 의 구성가능층 의 집합
$(X_i)_\etale$에 있다.
\item $X_\etale$에 건설 가능한 아벨 층의 범주
건설 가능한 아벨 층의 범주 중 여극한이다.
$(X_i)_\etale$에 있다.
\item $\Lambda$를 뇌터 환으로 둡니다. 건설 가능한 범주
층~의$\Lambda$-가군~에$X_\etale$은여극한~의
범주의 구성가능층의 $\Lambda$-가군의 $(X_i)_\etale$.
\end{enumerate}
\end{lemma}

\begin{proof} 증명 (1). 투영 지도를 $f_i : X \to X_i$로 표시한다. 증명에는 ``신실한'', ``완전히 충실한'', ``본질적으로 전사적인''에 해당하는 3 부분이 있다.

\medskip\noindent
충실한. $0 \in I$을 선택하고 $\mathcal{F}_0$, $\mathcal{G}_0$를
구성가능층 - $X_0$. 가정해보자
$a, b : \mathcal{F}_0 \to \mathcal{G}_0$는 다음과 같은 지도이다.
$f_0^{-1}a = f_0^{-1}b$. $E \subset X_0$를 집합으로 설정한다.
포인트$x \in X_0$그렇게$a_{\overline{x}} = b_{\overline{x}}$.
보조정리 \ref{etale-cohomology-lemma-support-constructible} 하위 집합에 의해
$E \subset X_0$은 구성 가능한다. 가정 $X \to X_0$는 $E$에 매핑된다.
극한, 보조정리 \ref{limits-lemma-limit-contained-in-constructible}에 의해
우리는$i \geq 0$그렇게$X_i \to X_0$에 매핑$E$.
그러므로 $f_{i0}^{-1}a = f_{i0}^{-1}b$.

\medskip\noindent
완전히 충실한다. $0 \in I$을 선택하고 $\mathcal{F}_0$, $\mathcal{G}_0$를
구성가능층 - $X_0$. 가정해보자
$a : f_0^{-1}\mathcal{F}_0 \to f_0^{-1}\mathcal{G}_0$는 지도이다.
우리는 주장 $i$와 지도가 있다.
$a_i : f_{i0}^{-1}\mathcal{F}_0 \to f_{i0}^{-1}\mathcal{G}_0$
다시 끌어당기는$a$~에$X$.
보조정리 \ref{etale-cohomology-lemma-category-constructible-sets}에 의해
$\mathcal{F}_0$를 대체할 수 있다.
준컴팩트(quasi-compact)로 표현되는 층의 유한 부산물에 의해
및 $(X_0)_\etale$의 준분리된 객체.
따라서 우리는 다음을 보여야 합니다: $U_0 \to X_0$가 그러한 객체인 경우
$(X_0)_\etale$의 다음
$$
f_0^{-1}\mathcal{G}(U) = \colim_{i \geq 0} f_{i0}^{-1}\mathcal{G}(U_i)
$$
여기서 $U = X \times_{X_0} U_0$ 및 $U_i = X_i \times_{X_0} U_0$이다.
이는 정리 \ref{etale-cohomology-theorem-colimit}의 특별한 경우이다.

\medskip\noindent
본질적으로 전사적이다. 우리는 건설 가능한 모든 것을 보여줘야 한다. $\mathcal{F}$
위의 $X$은 일부 구성 가능한 경우 $f_i^{-1}\mathcal{F}$와 동형이다.
$\mathcal{F}_i$ - $X_i$. 신청
보조정리 \ref{etale-cohomology-lemma-category-constructible-sets}
이전 두 단락의 결과를 사용하여 다음을 알 수 있다.
일부 준소형의 경우 $h_U$에 대해 이를 증명하는 것으로 충분한다.
그리고 준분리된 물체$U$~의$X_\etale$.
이 경우 $U$가 다음의 기본 변경임을 보여야 한다.
준컴팩트하고 준분리형스킴에탈 오버$X_i$어떤 사람들에게는$i$.
이것은 다음과 같습니다
극한, 기본형 \ref{limits-lemma-descend-finite-presentation} 및
\ref{limits-lemma-descend-etale}.

\medskip\noindent
증명 (3). 주장은 의 주장과 매우 ​​유사한다.
층 중 집합이지만 다음을 사용하고 있다.
보조정리 \ref{etale-cohomology-lemma-category-constructible-modules}
대신에
보조정리 \ref{etale-cohomology-lemma-category-constructible-sets}. 자세한 내용은 생략했다.
부분 (2)은 부분 (3)에서 이어집니다. 왜냐하면 모든 구성 가능한 아벨리안
층 준소형 스킴에 대한 구성가능층은 다음과 같습니다.
$\mathbf{Z}/n\mathbf{Z}$-가군 일부 $n$에 대해.
\end{proof}

\begin{lemma}
\label{etale-cohomology-lemma-category-loc-cst-is-colimit}
$X = \lim_{i \in I} X_i$를 극한으로 지정한다.
아핀 전이 사상이 있는 스킴 시스템.
$X_i$은 준컴팩트하고 준분리되어 있다고 가정한다.
모든 $i \in I$에 대해.
\begin{enumerate}
\item $X_\etale$의 유한 국소상수층의 범주
유한 국소상수층의 범주의 여극한이다.
$(X_i)_\etale$에 있다.
\item $X_\etale$에 대한 유한 지역 상수 아벨 층의 범주
유한 지역 상수 아벨 층의 범주의 여극한이다.
$(X_i)_\etale$에 있다.
\item $\Lambda$를 뇌터 환으로 둡니다. 유한형의 범주,
국소상수층~의$\Lambda$-가군~에$X_\etale$
유한 유형의 범주 중 여극한이며, 지역적으로 상수이다.
층~의$\Lambda$-가군~에$(X_i)_\etale$.
\end{enumerate}
\end{lemma}

\begin{proof}
보조정리 \ref{etale-cohomology-lemma-category-is-colimit} 각 경우에 함자
온전히 충실한다. 마찬가지로 보조정리, 끝까지 보여드려야 한다.
(1)의 경우 증명은 다음과 같습니다. 구성가능층이 주어지면
$\mathcal{F}_i$ $X_i$에서 $\mathcal{F}$에서 $X$로의 철회는
유한 지역 상수, $i' \geq i$가 존재한다.
철수$\mathcal{F}_{i'}$~의$\mathcal{F}_i$에게$X_{i'}$
유한 지역적으로 상수이다. 가정에 의해 에탈 피복이 존재한다.
$\mathcal{U} = \{U_j \to X\}_{j \in J}$
그렇게$\mathcal{F}|_{U_j} \cong \underline{S_j}$일부 유한한 경우집합 $S_j$.
$U_j$는 모든 $j \in J$에 대해 유사하다고 가정할 수 있다.
$X$는 준컴팩트하므로 $J$는 유한하다고 가정할 수 있다.
보조정리 \ref{etale-cohomology-lemma-colimit-affine-sites}로 $i' \geq i$를 찾을 수 있다.
그리고 에탈 피복 $\mathcal{U}_{i'} = \{U_{i', j} \to X_{i'}\}_{j \in J}$
$X$에 대한 기본 변경 사항은 $\mathcal{U}$이다. 그 다음에
$\mathcal{F}_{i'}|_{U_{i', j}}$ 및 $\underline{S_j}$는
구성가능층 $(U_{i', j})_\etale$에 대한 철회
$U_j$는 동형이다. 따라서 $i'$를 증가시킨 후에 우리는 다음을 얻습니다.
$\mathcal{F}_{i'}|_{U_{i', j}}$ 및 $\underline{S_j}$
동형이다. 따라서 $\mathcal{F}_{i'}$는 유한한 국소적 상수이다.
(2) 및 (3)의 경우 증명은 완전히 동일한다.
\end{proof}

\begin{lemma}
\label{etale-cohomology-lemma-irreducible-subsheaf-constant-zero}
$X$를 일반 포인트 $\eta$가 있는 기약 스킴으로 설정한다.
\begin{enumerate}
\item $S' \subset S$를 집합의 포함으로 간주한다. 우리가 가지고 있다면
$\underline{S'} \subset \mathcal{G} \subset \underline{S}$
$\Sh(X_\etale)$ 및 $S' = \mathcal{G}_{\overline{\eta}}$에서
$\mathcal{G} = \underline{S'}$.
\item $A' \subset A$를 아벨 군의 포함으로 간주한다. 우리가 가지고 있다면
$\underline{A'} \subset \mathcal{G} \subset \underline{A}$
$\textit{Ab}(X_\etale)$ 및 $A' = \mathcal{G}_{\overline{\eta}}$에서
$\mathcal{G} = \underline{A'}$.
\item $M' \subset M$를 환 $\Lambda$ 위에 가군을 포함한다고 가정한다.
$\underline{M'} \subset \mathcal{G} \subset \underline{M}$
$\textit{Mod}(X_\etale, \underline{\Lambda})$
그리고 $M' = \mathcal{G}_{\overline{\eta}}$, 그리고
$\mathcal{G} = \underline{M'}$.
\end{enumerate}
\end{lemma}

\begin{proof} 모든 에탈 사상 $U \to X$에 대해 $U \not = \emptyset$ 지점 $\eta$이 이미지에 있기 때문에 이는 사실이다. \end{proof}

\begin{lemma} \label{etale-cohomology-lemma-push-constant-sheaf-from-open} $X$를 정역 정규 스킴 함수 체 $K$로 설정한다. $E$를 집합으로 둡니다. \begin{enumerate} \item $g : \Spec(K) \to X$를 일반 포인트의 포함으로 둡니다. 그런 다음 $g_*\underline{E} = \underline{E}$. \item $j : U \to X$는 비어 있지 않은 열린 항목을 포함한다고 가정한다. 그런 다음 $j_*\underline{E} = \underline{E}$. \end{enumerate} \end{lemma}

\begin{proof} 증명 (1). $x \in X$를 점으로 둡니다. $\mathcal{O}^{sh}_{X, \overline{x}}$를 $\mathcal{O}_{X, x}$의 엄격한 헨젤화로 가정한다. 대수학에 대한 자세한 내용을 보면 보조정리 \ref{more-algebra-lemma-henselization-normal} $\mathcal{O}^{sh}_{X, \overline{x}}$가 정규 도메인임을 알 수 있다. 따라서 $\Spec(K) \times_X \Spec(\mathcal{O}^{sh}_{X, \overline{x}})$는 기약이다. 줄기 $(g_*\underline{E})_{\overline{x}}$는 $E$와 같습니다. 정리 \ref{etale-cohomology-theorem-higher-direct-images}를 참조하라.

\medskip\noindent
증명 (2). $g$는 $j$를 통해 요소가 되므로 지도가 있다.
$j_*\underline{E} \to g_*\underline{E}$. 이 지도는 주입식이므로
스킴 $V$ 에탈 $X$ 집합 $\Spec(K) \times_X V$에 대해
$U \times_X V$의 조밀이다. 반면에 지도가 있다.
$\underline{E} \to j_*\underline{E}$ 그리고 결론을 내립니다.
\end{proof}

\begin{lemma}
\label{etale-cohomology-lemma-zero-in-generic-point}
$X$를 준간소화하고 준분리된 스킴라고 하자. 허락하다
$\eta \in X$은 $X$의 기약 구성 요소의 일반적인 지점이다.
\begin{enumerate}
\item $\mathcal{F}$를 $X_\etale$에서 비틀림 아벨 층으로 설정한다.
줄기 $\mathcal{F}_{\overline{\eta}}$는 0이다.
그러면 $\mathcal{F} = \colim \mathcal{F}_i$는 다음의 여과 여극한이다.
아벨 층 $\mathcal{F}_i$ 각 $i$에 대해 구성 가능
$\mathcal{F}_i$의 지지집합이 포함되어 있다.
$\eta$을 포함하지 않는 폐쇄형 하위 체계에서.
\item $\Lambda$는 뇌터 환 및 $\mathcal{F}$는 층
~의$\Lambda$-가군~에$X_\etale$누구의줄기
$\mathcal{F}_{\overline{\eta}}$은 0이다. 그 다음에
$\mathcal{F} = \colim \mathcal{F}_i$
는 구성가능층의 여과 여극한이다.
$\Lambda$-가군 $\mathcal{F}_i$ 각 $i$에 대해
$\mathcal{F}_i$의 지지집합이 폐쇄된 하위 구성표에 포함되어 있다.
$\eta$을 포함하지 않습니다.
\end{enumerate}
\end{lemma}

\begin{proof}
증명 (1). $\mathcal{F} = \colim_{i \in I} \mathcal{F}_i$라고 쓸 수 있다.
$\mathcal{F}_i$ 구성 가능한 아벨리안 사용
보조정리 \ref{etale-cohomology-lemma-torsion-colimit-constructible}.
$i \in I$을 선택한다. $\mathcal{F}|_\eta$는 가정에 의해 0이므로
$i'(i) \geq i$가 있는지 확인하세요.
$\mathcal{F}_i|_\eta \to \mathcal{F}_{i'(i)}|_\eta$은 0이다. 참조
보조정리 \ref{etale-cohomology-lemma-colimit-constructible}.
그런 다음 $\mathcal{G}_i = \Im(\mathcal{F}_i \to \mathcal{F}_{i'(i)})$
구성 가능한 아벨 층 (보조정리 \ref{etale-cohomology-lemma-constructible-abelian})
$\eta$의 줄기가 0이다.
따라서지지집합 $E_i$~의$\mathcal{G}_i$건설 가능하다
하위 집합$X$포함하지 않음$\eta$. 부터
$\eta$는 기약 구성 요소의 일반적인 지점이다.
$X$, 우리는 $\eta \not \in Z_i = \overline{E_i}$를 통해
위상, 보조정리 \ref{topology-lemma-generic-point-in-constructible}.
새로운 감독 정의집합 $I'$을 사용하여집합 $I$~와 함께
규칙에 의해 정의된 순서
$i_1$는 $i_2$ 필요충분조건은 $i_1 \geq i'(i_2)$보다 크거나 같습니다.
그런 다음 층 $\mathcal{G}_i$는 $I'$를 통해 시스템을 형성한다.
여극한 $\mathcal{F}$ 및 증명이 완료되었다.

\medskip\noindent (2)의 경우 증명은 완전히 동일하므로 생략한다. \end{proof}








\section{구성가능층 뇌터 스킴} \label{etale-cohomology-section-noetherian-constructible}

\noindent
$X$가 뇌터 스킴인 경우 로컬에서 닫힌 부분집합은
로컬에서 시공 가능 닫힌 부분집합
(위상, 보조정리 \ref{topology-lemma-constructible-Noetherian-space}).
따라서아벨 층 $\mathcal{F}$~에$X_\etale$건설 가능하다
필요충분조건은 유한한 파티션이 존재한다. $X = \coprod X_i$
$\mathcal{F}|_{X_i}$는 유한한 지역적으로 상수이다.
(관례적으로 위상 공간의 분할은 지역적으로
닫힌 부품, 위상, 절 \ref{topology-section-stratifications} 참조)
즉, 형용사 '구성 가능'을 생략할 수 있다.
정의 \ref{etale-cohomology-definition-constructible}. 사실, 범주
구성가능층 중
뇌터 스킴에는 추가 성질이 있다.
이 절의 카탈로그이다.

\begin{proposition}
\label{etale-cohomology-proposition-constructible-over-noetherian}
$X$를 뇌터 스킴으로 설정한다. $\Lambda$를 뇌터 환으로 설정한다.
\begin{enumerate}
\item 집합의 구성가능층의 하위 또는 몫 층
시공 가능한다.
\item $X_\etale$에 건설 가능한 아벨 층의 범주는
(강함) $\textit{Ab}(X_\etale)$의 하위 카테고리이다. 특히,
구성 가능한 아벨 층의 모든 하위 및 몫 층
위의 $X_\etale$은 구성 가능한다.
\item 구성가능층의 $\Lambda$-가군의 범주
위의 $X_\etale$은 (강력한) Serre 하위 카테고리이다.
$\textit{Mod}(X_\etale, \Lambda)$. 특히, 모든 하위 모듈
및 $\Lambda$-가군의 구성가능층의 몫 가군
위의 $X_\etale$은 구성 가능한다.
\end{enumerate}
\end{proposition}

\begin{proof}
증명 (1). $\mathcal{G} \subset \mathcal{F}$를 $\mathcal{F}$로 설정
$X_\etale$에 있는 집합의 구성가능층. $\eta \in X$를 일반으로 설정
$X$의 기약 구성 요소 지점이다. 뇌터 유도로
열린 동네를 찾는 것만으로도 충분합니다$U$~의$\eta$그렇게
$\mathcal{G}|_U$는 지역적으로 상수이다. 이를 위해 $X$를 대체할 수 있다.
$\eta$의 에탈 동네로.
따라서 $\mathcal{F}$는 상수이고 $X$는 기약라고 가정할 수 있다.

\medskip\noindent
말하다$\mathcal{F} = \underline{S}$일부 유한한 경우집합 $S$.
그런 다음 $S' = \mathcal{G}_{\overline{\eta}} \subset S$
$S' = \{s_1, \ldots, s_t\}$라고 말하세요.
에탈 지역을 선택하세요$(U, \overline{u})$~의$\overline{\eta}$
및 절 $\sigma_1, \ldots, \sigma_t \in \mathcal{G}(U)$에 매핑된다.
$s_i$~에$\mathcal{G}_{\overline{\eta}} \subset S$.
$\sigma_i$는 요소에 매핑되므로
$s_i \in S' \subset S = \Gamma(X, \mathcal{F})$
$\sigma_i$에서 $U \times_X U$로의 두 가지 철회가 있음을 알 수 있다.
$\mathcal{G}$의 절과 같습니다. 층 조건에 의해
$\mathcal{G}$에 대해 $\sigma_i$가 절에서 온다는 것을 알 수 있다.
~의$\mathcal{G}$공개된 곳에서$\Im(U \to X)$~의$X$.
축소 $X$ 가정할 수 있다.
$\underline{S'} \subset \mathcal{G} \subset \underline{S}$.
그러면 우리는 $\underline{S'} = \mathcal{G}$를 볼 수 있다.
보조정리 \ref{etale-cohomology-lemma-irreducible-subsheaf-constant-zero}.

\medskip\noindent
$\mathcal{F} \to \mathcal{Q}$를 $\mathcal{F}$로 변환한다.
$X_\etale$에 있는 집합의 구성가능층. 그러면 집합
$\mathcal{G} = \mathcal{F} \times_\mathcal{Q} \mathcal{F}$.
증명의 첫 번째 부분에서 $\mathcal{G}$가
$\mathcal{F} \times \mathcal{F}$의 하위 묶음으로 구성 가능한다.
이는 결국 $\mathcal{Q}$가 구성 가능함을 의미한다.
보조정리 \ref{etale-cohomology-lemma-constructible-abelian}.

\medskip\noindent 증명 (3). 우리는 가군의 구성가능층이 약한 Serre 하위 범주를 형성한다는 것을 이미 알고 있다. 보조정리 \ref{etale-cohomology-lemma-constructible-abelian}를 참조하라. 따라서 하위 모듈에 대한 설명을 표시하는 것으로 충분한다.

\medskip\noindent
$\mathcal{G} \subset \mathcal{F}$를
구성가능층~의$\Lambda$-가군~에$X_\etale$. 허락하다$\eta \in X$
$X$의 기약 구성 요소의 일반적인 지점이 된다. 뇌터 유도로
열린 동네를 찾는 것만으로도 충분합니다$U$~의$\eta$그렇게
$\mathcal{G}|_U$는 지역적으로 상수이다. 이를 위해 $X$를 대체할 수 있다.
$\eta$의 에탈 동네로. 따라서 우리는 $\mathcal{F}$라고 가정할 수 있다.
는 상수이고 $X$는 기약이다.

\medskip\noindent
말하다$\mathcal{F} = \underline{M}$일부 유한한 경우$\Lambda$-가군 $M$.
그런 다음 $M' = \mathcal{G}_{\overline{\eta}} \subset M$. 무한히 골라라
많은 요소 $s_1, \ldots, s_t$가 $M'$을 $\Lambda$-가군으로 생성한다.
($\Lambda$은 뇌터이고 $M$는 유한하기 때문에 가능한다.)
에탈 지역을 선택하세요$(U, \overline{u})$~의$\overline{\eta}$
및 절 $\sigma_1, \ldots, \sigma_t \in \mathcal{G}(U)$에 매핑된다.
$s_i$~에$\mathcal{G}_{\overline{\eta}} \subset M$.
$\sigma_i$는 요소에 매핑되므로
$s_i \in M' \subset M = \Gamma(X, \mathcal{F})$
$\sigma_i$에서 $U \times_X U$로의 두 가지 철회가 있음을 알 수 있다.
$\mathcal{G}$의 절과 같습니다. 층 조건에 의해
$\mathcal{G}$에 대해 $\sigma_i$가 절에서 온다는 것을 알 수 있다.
~의$\mathcal{G}$공개된 곳에서$\Im(U \to X)$~의$X$.
축소 $X$ 가정할 수 있다.
$\underline{M'} \subset \mathcal{G} \subset \underline{M}$.
그러면 우리는 $\underline{M'} = \mathcal{G}$를 볼 수 있다.
보조정리 \ref{etale-cohomology-lemma-irreducible-subsheaf-constant-zero}.

\medskip\noindent 증명 (2). 이는 (3)에서 일반적인 방식으로 이어집니다. 자세한 내용은 생략했다. \end{proof}

\noindent 다음 보조정리는 $X$에 있는 구성가능층의 아벨 범주의 모든 객체가 ``뇌터''임을 알려줍니다. 즉,.c.c.\ 하위 객체용.

\begin{lemma}
\label{etale-cohomology-lemma-constructible-over-noetherian-noetherian}
$X$를 뇌터 스킴으로 설정한다. $\Lambda$를 뇌터 환으로 설정한다.
포함사항을 고려하세요
$$
\mathcal{F}_1 \subset \mathcal{F}_2 \subset \mathcal{F}_3 \subset \ldots
\subset \mathcal{F}
$$
범주의 층, 집합, 아벨 군 또는 $\Lambda$-가군.
$\mathcal{F}$이 구성 가능하다면 일부 $n$에 대해서는
$\mathcal{F}_n = \mathcal{F}_{n + 1} = \mathcal{F}_{n + 2} = \ldots$이 있다.
\end{lemma}

\begin{proof} 명제 \ref{etale-cohomology-proposition-constructible-over-noetherian}에 의해 $\mathcal{F}_i$와 $\colim \mathcal{F}_i$가 구성 가능하다는 것을 알 수 있다. 그런 다음 보조정리는 보조정리 \ref{etale-cohomology-lemma-colimit-constructible}에서 이어집니다. \end{proof}

\begin{lemma}
\label{etale-cohomology-lemma-constructible-maps-into-constant}
$X$를 뇌터 스킴으로 설정한다.
\begin{enumerate}
\item $\mathcal{F}$를 $X_\etale$에서 집합의 구성가능층으로 설정한다.
층의 주입형 맵이 존재한다.
$$
\mathcal{F} \longrightarrow
\prod\nolimits_{i = 1, \ldots, n} f_{i, *}\underline{E_i}
$$
여기서 $f_i : Y_i \to X$는 유한 사상이고 $E_i$는 유한 집합이다.
\item $\mathcal{F}$를 $X_\etale$에서 구성 가능한 아벨 층으로 설정한다.
아벨 층의 주입형 맵이 존재한다.
$$
\mathcal{F} \longrightarrow
\bigoplus\nolimits_{i = 1, \ldots, n} f_{i, *}\underline{M_i}
$$
여기서 $f_i : Y_i \to X$는 유한한 사상이고
$M_i$는 유한한 아벨 군이다.
\item $\Lambda$를 뇌터 환으로 둡니다.
허락하다$\mathcal{F}$가 되다구성가능층~의$\Lambda$-가군~에$X_\etale$.
가군층의 주입형 맵이 존재한다.
$$
\mathcal{F} \longrightarrow
\bigoplus\nolimits_{i = 1, \ldots, n} f_{i, *}\underline{M_i}
$$
여기서 $f_i : Y_i \to X$는 유한한 사상이고
$M_i$는 유한한 $\Lambda$-가군이다.
\end{enumerate}
게다가, 우리는 각각의 $Y_i$가 기약이고 축소되어 매핑된다고 가정할 수 있다.
기약 및 축소된 폐쇄형 하위 계획 $Z_i \subset X$
$Y_i \to Z_i$는 $Z_i$의 비어 있지 않은 개방에 대해 유한 에탈이다.
\end{lemma}

\begin{proof}
증명 (1). 왜냐하면 우리는 다음에 대한 오름차순 체인 조건을 가지고 있기 때문이다.
$\mathcal{F}$의 서브시브
(보조정리 \ref{etale-cohomology-lemma-constructible-over-noetherian-noetherian}), 그것은
모든 포인트 $x \in X$에 대해 우리가
$\varphi : \mathcal{F} \to f_*\underline{E}$ 지도를 찾을 수 있다.
$f : Y \to X$는 유한하고 $E$는 유한한 집합이다.
$\varphi_{\overline{x}} : \mathcal{F}_{\overline{x}} \to
(f_*S)_{\overline{x}}$는 단사형이다.
(이 인수는 다음과 같이 $X$ 파티션을 선택하여 피할 수 있다.
보조정리 \ref{etale-cohomology-lemma-constructible-quasi-compact-quasi-separated}
그리고 각 기약 구성 요소에 대해 $Y_i \to X$를 구성한다.
각 부분의.)
$Z \subset X$를 유도된 환원된 스킴 구조라고 하자
(스킴, 정의 \ref{schemes-definition-reduced-induced-scheme})
$\overline{\{x\}}$에 있다.
$\mathcal{F}$은 구성 가능하므로 유한한 분리 가능이 있다.
확장자 $K/\kappa(x)$
$\mathcal{F}|_{\Spec(K)}$는 값이 $E$인 상수층이다.
일부 유한 집합 $E$의 경우. $Y \to Z$를 정규화라고 하자.
~의$Z$~에$\Spec(K)$.
사상, 보조정리 \ref{morphisms-lemma-normal-normalization}에 의해
$Y$는 정규 정역 스킴라는 것을 알 수 있다.
$K/\kappa(x)$는 유한 확장이므로 $K$가 함수임이 분명한다.
체 중 $Y$. $g : \Spec(K) \to Y$ 포함을 나타냅니다.
지도 $\mathcal{F}|_{\Spec(K)} \to \underline{E}$가 인접해 있다.
지도로 $\mathcal{F}|_Y \to g_*\underline{E} = \underline{E}$
(보조정리 \ref{etale-cohomology-lemma-push-constant-sheaf-from-open}).
이것은 차례로 지도에 인접해 있다.
$\varphi : \mathcal{F} \to f_*\underline{E}$.
기하점에서 $\varphi$의 줄기를 관찰하세요.
$\overline{x}$에서의 사상은 전사이다. 실제로 올림 $\overline{y} \in Y$를
$\overline{x}$ 위에서 취할 수 있고, 가환 도식
$$
\xymatrix{
\mathcal{F}_{\overline{x}} \ar@{=}[r] \ar[d] &
(\mathcal{F}|_Y)_{\overline{y}} \ar@{=}[d] \\
(f_*\underline{E})_{\overline{x}} \ar[r] &
\underline{E}_{\overline{y}}
}
$$
주입성을 증명한다. 그러나 우리는 아직 끝나지 않았습니다.
사상 $f : Y \to Z$는 정역이지만 일반적으로 그렇지 않습니다.
유한\footnote{$X$가 나가타 스킴인 경우, 예의 경우 유한
체 위에 입력하면 $Y \to Z$는 유한한다.}.

\medskip\noindent
이전 단락의 마지막 문장에 명시된 문제를 해결하려면,
우리는 쓴다$Y = \lim_{i \in I} Y_i$~와 함께$Y_i$ 기약, 정역, 그리고
$Z$에 대해 유한한다. 즉, 성질, 보조정리를 적용한다.
\ref{properties-lemma-integral-algebra-directed-colimit-finite}
$f_*\mathcal{O}_Y$는 $\mathcal{O}_Z$-대수학의 층으로 간주된다.
함자 $\underline{\Spec}_Z$를 적용한다.
그런 다음 $f_*\underline{E} = \colim f_{i, *}\underline{E}$
보조정리 \ref{etale-cohomology-lemma-relative-colimit} 기준.
보조정리 \ref{etale-cohomology-lemma-constructible-is-compact} 지도로
$\mathcal{F} \to f_*\underline{E}$
요인을 통해$f_{i, *}\underline{E}$어떤 사람들에게는$i$.
$Y_i \to Z$는 정역 스킴의 유한 사상이므로
그리고 체 확장 기능 이후
이 사상에 의해 유도된 는 유한 분리가능하며, 우리는
사상은 $Z$의 비어 있지 않은 개방에 대해 유한 에탈입니다(사용
대수학, 보조정리 \ref{algebra-lemma-smooth-at-generic-point}; 자세한 내용은 생략)
이것으로 (1)의 증명이 완료된다.

\medskip\noindent (2) 및 (3)의 증명은 (1)의 증명과 동일한다. \end{proof}

\noindent 다음 보조정리에서는 매우 일반적인 설명을 뇌터 사례로 줄이기 위해 표준 트릭을 사용한다.

\begin{lemma}
\label{etale-cohomology-lemma-constructible-maps-into-constant-general}
\begin{reference}
\cite[Exposee IX, Proposition 2.14]{SGA4}
\end{reference}
$X$를 준간소화하고 준분리된 스킴라고 하자.
\begin{enumerate}
\item $\mathcal{F}$를 $X_\etale$에서 집합의 구성가능층으로 설정한다.
층의 주입형 맵이 존재한다.
$$
\mathcal{F} \longrightarrow
\prod\nolimits_{i = 1, \ldots, n} f_{i, *}\underline{E_i}
$$
여기서 $f_i : Y_i \to X$는 유한하고 유한하게 표시된다. 사상
$E_i$는 유한한 집합이다.
\item $\mathcal{F}$를 $X_\etale$에서 구성 가능한 아벨 층으로 설정한다.
아벨 층의 주입형 맵이 존재한다.
$$
\mathcal{F} \longrightarrow
\bigoplus\nolimits_{i = 1, \ldots, n} f_{i, *}\underline{M_i}
$$
여기서 $f_i : Y_i \to X$는 유한하고 유한하게 표시된다. 사상
$M_i$는 유한한 아벨 군이다.
\item $\Lambda$를 뇌터 환으로 둡니다.
허락하다$\mathcal{F}$가 되다구성가능층~의$\Lambda$-가군~에$X_\etale$.
가군층의 주입형 맵이 존재한다.
$$
\mathcal{F} \longrightarrow
\bigoplus\nolimits_{i = 1, \ldots, n} f_{i, *}\underline{M_i}
$$
여기서 $f_i : Y_i \to X$는 유한하고 유한하게 표시된다. 사상
$M_i$는 유한한 $\Lambda$-가군이다.
\end{enumerate}
\end{lemma}

\begin{proof}
이 보조정리를 절대 뇌터로 뇌터 사례로 줄이겠습니다.
근사. 즉,
극한, 명제 \ref{limits-proposition-approximate}
우리는 쓸 수 있어요$X = \lim_{t \in T} X_t$각각$X_t$유한 유형의
$\Spec(\mathbf{Z})$ 및 아핀 전이 사상이 있다. 에 의해
보조정리 \ref{etale-cohomology-lemma-category-is-colimit}
구성가능층의 범주 (집합, 아벨 군 또는
$\Lambda$-가군) 에$X_\etale$은여극한해당하는
$X_t$에 대한 범주. 따라서 우리의 구성가능층 $\mathcal{F}$
비슷한 구성가능층 $\mathcal{F}_t$의 당김이다.
~ 위에$X_t$어떤 사람들에게는$t$. 그런 다음 우리는뇌터사례
(보조정리 \ref{etale-cohomology-lemma-constructible-maps-into-constant})
주사를 찾으러
$$
\mathcal{F}_t \longrightarrow
\prod\nolimits_{i = 1, \ldots, n} f_{i, *}\underline{E_i}
\quad\text{or}\quad
\mathcal{F}_t \longrightarrow
\bigoplus\nolimits_{i = 1, \ldots, n} f_{i, *}\underline{M_i}
$$
~ 위에$X_t$일부 유한한 경우사상 $f_i : Y_i \to X_t$.
$X_t$는 뇌터이므로 사상 $f_i$는 유한하게 표현된다.
철수가 정확하고 $f_{i, *}$ 형성 이후 통근
베이스 변경으로
(보조정리 \ref{etale-cohomology-lemma-finite-pushforward-commutes-with-base-change}) 결론을 내립니다.
\end{proof}

\begin{lemma}
\label{etale-cohomology-lemma-support-in-subset}
$X$를 뇌터 스킴으로 설정한다. $E \subset X$를 닫힌 하위 집합으로 둡니다.
특수화 아래에 있다.
\begin{enumerate}
\item $\mathcal{F}$를 $X_\etale$에서 비틀림 아벨 층으로 설정한다.
지지집합이 $E$에 포함되어 있다. 그런 다음 $\mathcal{F} = \colim \mathcal{F}_i$
건설 가능한 아벨 층 $\mathcal{F}_i$의 여과 여극한이다.
각 $i$에 대해 $\mathcal{F}_i$의 지지집합이 다음에 포함된다.
$E$에 포함된 닫힌 부분집합이다.
\item $\Lambda$는 뇌터 환 및 $\mathcal{F}$는 층
~의$\Lambda$-가군~에$X_\etale$누구의지지집합에 포함되어 있습니다$E$.
그런 다음 $\mathcal{F} = \colim \mathcal{F}_i$
는 구성가능층의 여과 여극한이다.
$\Lambda$-가군 $\mathcal{F}_i$ 각 $i$에 대해
$\mathcal{F}_i$의 지지집합이 닫힌 부분집합에 포함되어 있다.
$E$에 포함되어 있다.
\end{enumerate}
\end{lemma}

\begin{proof}
증명 (1). $\mathcal{F} = \colim_{i \in I} \mathcal{F}_i$라고 쓸 수 있다.
$\mathcal{F}_i$ 구성 가능한 아벨리안 사용
보조정리 \ref{etale-cohomology-lemma-torsion-colimit-constructible}.
명제 \ref{etale-cohomology-proposition-constructible-over-noetherian}에 의해
이미지 $\mathcal{F}'_i \subset \mathcal{F}$
지도 $\mathcal{F}_i \to \mathcal{F}$는 건설 가능한다.
따라서 $\mathcal{F} = \colim \mathcal{F}'_i$ 및 지지집합
$\mathcal{F}'_i$이 $E$에 포함되어 있다.
$\mathcal{F}'_i$의 지지집합은 구성 가능하므로
(구성가능층의 정의에 의해), 우리는
클로저도 $E$에 포함되어 있다. 예를 참조하라.
위상, 보조정리
\ref{topology-lemma-constructible-stable-specialization-closed}.

\medskip\noindent (2)의 경우 증명은 완전히 동일하므로 생략한다. \end{proof}








\section{특수화 및 에탈 층} \label{etale-cohomology-section-specialization}

\noindent
위상 그림: $X$를 위상 공간으로 두고 $x' \leadsto x$를
$X$에 있는 포인트의 특수화이다. 그런 다음 $x$의 모든 열린 동네
포함$x'$. 따라서 누구에게나층 $\mathcal{F}$~에$X$있다
{\it 특수화 지도}
$$
sp : \mathcal{F}_x \longrightarrow \mathcal{F}_{x'}
$$
줄기 쌍 $(U, s)$의 동등 클래스를 보내는 중
$\mathcal{F}_x$를 $(U, s)$ 쌍의 동등 클래스로
$\mathcal{F}_{x'}$; 층, 절 \ref{sheaves-section-stalks} 참조
쌍의 동등 클래스 측면에서 줄기에 대한 설명이다.
물론 이 지도는 $\mathcal{F}$에서 기능적이다. 즉, $sp$
함자의 변환이다.

\medskip\noindent
에탈 위상의 층에 대해 우리는 이 구조를 흉내낼 수 있다.
\cite[Exposee VIII, 7.7, page 397]{SGA4}. 이를 위해 우리는
스킴 $S$, 기하점 $\overline{s}$ $S$, 그리고 기하학적
가리키다$\overline{t}$~의$\Spec(\mathcal{O}^{sh}_{S, \overline{s}})$.
누구에게나층 $\mathcal{F}$~에$S_\etale$우리는
{\it 특수화 지도}
$$
sp : \mathcal{F}_{\overline{s}} \longrightarrow \mathcal{F}_{\overline{t}}
$$
여기서 우리는 언어를 남용했다: 글을 쓰는 대신
$\mathcal{F}_{\overline{t}}$ $\mathcal{F}_{p(\overline{t})}$라고 써야 한다.
여기서 $p : \Spec(\mathcal{O}^{sh}_{S, \overline{s}}) \to S$는 표준이다.
사상. 그것을 기억해
$$
\mathcal{F}_{\overline{s}} = \colim_{(U, \overline{u})} \mathcal{F}(U)
$$
여기서 여극한은 모든 에탈 동네 $(U, \overline{u})$에 걸쳐 있다.
$(S, \overline{s})$의 경우, 절 \ref{etale-cohomology-section-stalks}를 참조하라. 부터
$\mathcal{O}^{sh}_{S, \overline{s}}$는 줄기이다.
구조층, 모든 에탈 동네에 대해 찾습니다.
$(U, \overline{u})$의 $(S, \overline{s})$ 표준 맵
$\mathcal{O}_{U, u} \to \mathcal{O}^{sh}_{S, \overline{s}}$.
따라서 우리는 고유 인수분해를 얻습니다.
$$
\Spec(\mathcal{O}^{sh}_{S, \overline{s}}) \to U \to S
$$
만약에$\overline{v}$의 이미지를 나타냅니다.$\overline{t}$~에$U$, 그러면 우리는 그것을 볼 수 있습니다
$(U, \overline{v})$은 $(S, \overline{t})$의 에탈 동네이다.
이 구성은 에탈 지역의 범주에서 함자를 정의한다.
$(S, \overline{s})$에서 에탈 동네의 범주까지
$(S, \overline{t})$. 따라서 우리는 지도를 정의할 수 있다.
$sp : \mathcal{F}_{\overline{s}} \to \mathcal{F}_{\overline{t}}$
동등한 클래스를 보내서
$(U, \overline{u}, \sigma)$ 여기서 $\sigma \in \mathcal{F}(U)$
$(U, \overline{v}, \sigma)$의 동등 클래스로

\medskip\noindent $K \in D(S_\etale)$로 둡니다. 위와 같이 $\overline{s}$ 및 $\overline{t}$를 사용하면 {\it 특수화 맵} $$
sp : K_{\overline{s}} \longrightarrow K_{\overline{t}}
\quad\text{~에}\quad D(\textit{Ab})
$$이 있다. 즉, $K$가 복합체 $\mathcal{F}^\bullet$ 의 아벨 층, 그러면 우리는 위에서 구성된 층에 대한 특수화 맵에 의해 용어별로 주어진 $$
K_{\overline{s}} = \mathcal{F}^\bullet_{\overline{s}}
\longrightarrow
\mathcal{F}^\bullet_{\overline{t}} = K_{\overline{t}}
$$ 맵을 간단히 정의한다. 이는 줄기 함자의 완전성에 의해 $K$를 나타내는 복합체의 선택과 무관합니다(즉, 복합체의 줄기를 취하는 것은 유도 범주).

\medskip\noindent 분명히 구성은 층 $\mathcal{F}$ $S_\etale$에서 기능적이다. 줄기 함자를 토포스 $\overline{s}, \overline{t} : \textit{Sets} \to \Sh(S_\etale)$의 사상으로 생각하면 $sp$를 $2$-사상으로 생각할 수 있다. 토포스의 $$
\xymatrix{
\textit{Sets}
\rrtwocell^{\overline{t}}_{\overline{s}}{\ sp}
&
&
\Sh(S_\etale)
}
$$.

\begin{remark}[sp의 대체 설명]
\label{etale-cohomology-remark-alternative-sp}
$S$, $\overline{s}$, $\overline{t}$를 위와 같이 해보자.
특수화 지도를 설명하는 또 다른 방법은
$$
\mathcal{F}_{\overline{s}} =
\Gamma(\Spec(\mathcal{O}^{sh}_{S, \overline{s}}), p^{-1}\mathcal{F})
\quad\text{이고 }\quad
\mathcal{F}_{\overline{t}} = \Gamma(\overline{t},
\overline{t}^{-1}p^{-1}\mathcal{F})
$$
첫 번째 동일성은 정리 \ref{etale-cohomology-theorem-higher-direct-images}에서 따릅니다.
$\text{id}_S : S \to S$에 적용되고 두 번째 동일성은 다음과 같습니다.
보조정리 \ref{etale-cohomology-lemma-stalk-pullback}에서. 그러면 $sp$를 생각해 볼 수 있다.
지도로
$$
sp :
\mathcal{F}_{\overline{s}} =
\Gamma(\Spec(\mathcal{O}^{sh}_{S, \overline{s}}), p^{-1}\mathcal{F})
\xrightarrow{\text{에 의한 당김 }\overline{t}}
\Gamma(\overline{t}, \overline{t}^{-1}p^{-1}\mathcal{F}) =
\mathcal{F}_{\overline{t}}
$$
\end{remark}

\begin{remark}[sp에 대한 또 다른 설명]
\label{etale-cohomology-remark-another-sp}
$S$, $\overline{s}$, $\overline{t}$를 위와 같이 해보자.
또 다른 대안은 고유한 사상을 사용하는 것이다.
$$
c :
\Spec(\mathcal{O}^{sh}_{S, \overline{t}})
\longrightarrow
\Spec(\mathcal{O}^{sh}_{S, \overline{s}})
$$
주어진 사상과 호환되는 $S$ 이상
$\overline{t} \to \Spec(\mathcal{O}^{sh}_{S, \overline{s}})$
그리고 사상
$\overline{t} \to \Spec(\mathcal{O}^{sh}_{t, \overline{t}})$.
표시된 화살표의 고유성과 존재 여부
대수학, 보조정리 \ref{algebra-lemma-map-into-henselian-colimit}
$\mathcal{O}_{S, s}$, $\mathcal{O}^{sh}_{S, \overline{t}}$에 적용됨
$\mathcal{O}^{sh}_{S, \overline{s}} \to \kappa(\overline{t})$.
우리는 얻는다
$$
sp :
\mathcal{F}_{\overline{s}} =
\Gamma(\Spec(\mathcal{O}^{sh}_{S, \overline{s}}), \mathcal{F})
\xrightarrow{\text{에 의한 당김 }c}
\Gamma(\Spec(\mathcal{O}^{sh}_{S, \overline{t}}), \mathcal{F}) =
\mathcal{F}_{\overline{t}}
$$
(명백한 표기법에 따라).
실제로 이 절차는 객체에도 적용된다.$K$~에$D(S_\etale)$:
$K$에 대한 특수화 지도는 지도와 같습니다.
$$
sp :
K_{\overline{s}} =
R\Gamma(\Spec(\mathcal{O}^{sh}_{S, \overline{s}}), K)
\xrightarrow{\text{에 의한 당김 }c}
R\Gamma(\Spec(\mathcal{O}^{sh}_{S, \overline{t}}), K) =
K_{\overline{t}}
$$
등호는 대역 절단을 대신하여 유효한다.
엄격한 헨셀식 스킴
$\Spec(\mathcal{O}^{sh}_{S, \overline{s}})$ 그리고
$\Spec(\mathcal{O}^{sh}_{S, \overline{t}})$
정확합니다(그리고 $\overline{s}$ 및 $\overline{t}$에서 줄기를 취하는 것과 동일)
따라서 $K$가 제한되지 않을 수 있다는 사실과 관련된 미묘함은 없다.
생기다.
\end{remark}

\begin{remark}[들어올리기 특수화]
\label{etale-cohomology-remark-can-lift}
$S$를 스킴으로 두고 $t \leadsto s$를 특수화로 두십시오.
$S$에 점을 찍습니다. 기하점 $\overline{t}$ 및 $\overline{s}$를 선택한다.
$t$와 $s$ 위에 놓여 있다. $t$는
$\Spec(\mathcal{O}_{S, s})$ 에 의해
스킴, 보조정리 \ref{schemes-lemma-specialize-points}
그리고 이후 $\mathcal{O}_{S, s} \to \mathcal{O}^{sh}_{S, \overline{s}}$
충실히 평탄이므로 점을 찾을 수 있다.
$t' \in \Spec(\mathcal{O}^{sh}_{S, \overline{s}})$를 $t$에 매핑한다.
$\Spec(\mathcal{O}^{sh}_{S, \overline{s}})$는 다음의 극한이다.
스킴의 에탈은 $S$에 대해 $\kappa(t')/\kappa(t)$를 참조하라.
분리 가능한 대수적 확장입니다(물론 일반적으로 유한하지는 않습니다).
$\kappa(\overline{t})$는 대수적으로 닫혀 있으므로 다음이 가능한다.
임베딩 선택 $\kappa(t') \to \kappa(\overline{t})$
$\kappa(t)$의 확장자로. 이 선택은 우리에게
가환 도식
$$
\xymatrix{
\overline{t} \ar[d] \ar[r] &
\Spec(\mathcal{O}^{sh}_{S, \overline{s}}) \ar[d] &
\overline{s} \ar[l] \ar[d] \\
t \ar[r] & S & s \ar[l]
}
$$
포인트 및 기하점. 따라서 $t \leadsto s$
우리는 항상 $\overline{t}$를 기하점으로 ``리프트''할 수 있다.
엄격한 헨젤화의$S$~에$\overline{s}$
위와 같이 특수화 지도를 얻습니다.
\end{remark}

\begin{lemma}
\label{etale-cohomology-lemma-specialization-map-pullback}
$g : S' \to S$를 스킴의 사상으로 설정한다. $\mathcal{F}$를 층으로 설정한다.
$S_\etale$에 있다. $\overline{s}'$를 $S'$의 기하점으로 두고,
$\overline{t}'$는 다음의 기하점이다.
$\Spec(\mathcal{O}^{sh}_{S', \overline{s}'})$. 표시하다
$\overline{s} = g(\overline{s}')$ 및 $\overline{t} = h(\overline{t}')$
여기서 $h : \Spec(\mathcal{O}^{sh}_{S', \overline{s}'}) \to
\Spec(\mathcal{O}^{sh}_{S, \overline{s}})$는 정식 사상이다.
누구에게나층 $\mathcal{F}$~에$S_\etale$그만큼특수화지도
$$
sp :
(g^{-1}\mathcal{F})_{\overline{s}'}
\longrightarrow
(g^{-1}\mathcal{F})_{\overline{t}'}
$$
특수화 지도와 같습니다.
$sp : \mathcal{F}_{\overline{s}} \to \mathcal{F}_{\overline{t}}$
신분증을 통해
$(g^{-1}\mathcal{F})_{\overline{s}'} = \mathcal{F}_{\overline{s}}$ 그리고
$(g^{-1}\mathcal{F})_{\overline{t}'} = \mathcal{F}_{\overline{t}}$
보조정리 \ref{etale-cohomology-lemma-stalk-pullback}.
\end{lemma}

\begin{proof} 생략한다. \end{proof}

\begin{lemma} \label{etale-cohomology-lemma-characterize-locally-constant} $S$를 스킴으로 설정하여 $S$의 모든 준 컴팩트 개방에는 유한한 수의 기약 구성 요소가 있습니다(예의 경우 $S$에는 뇌터 기본 위상 공간이 있거나 $S$가 로컬로 뇌터인 경우). $\mathcal{F}$를 $S_\etale$에서 집합의 층으로 설정한다. 다음은 등가이다. \begin{enumerate} \item $\mathcal{F}$는 유한 지역 상수이고 \item $\mathcal{F}$의 모든 줄기는 유한 집합 및 모든 특수화 맵이다. $sp : \mathcal{F}_{\overline{s}} \to \mathcal{F}_{\overline{t}}$는 전단사이다. \end{enumerate} \end{lemma}

\begin{proof}
(2)를 가정한다. $\overline{s}$를 $S$의 기하점으로 놔두세요
$s \in S$. (1)를 증명하려면 에탈 이웃을 찾아야 한다.
$(U, \overline{u})$ $(S, \overline{s})$의 $\mathcal{F}|_U$
일정한다. 우리는 $S$가 유사하다고 가정할 수도 있고 실제로 가정한다.

\medskip\noindent
$\mathcal{F}_{\overline{s}}$는 유한하므로 다음을 선택할 수 있다.
$(U, \overline{u})$, $n \geq 0$ 및 쌍별 개별 요소
$\sigma_1, \ldots, \sigma_n \in \mathcal{F}(U)$
$\{\sigma_1, \ldots, \sigma_n\} \subset \mathcal{F}(U)$ 매핑되도록
지도를 통해 $\mathcal{F}_{\overline{s}}$에 전단적으로
$\mathcal{F}(U) \to \mathcal{F}_{\overline{s}}$.
지도를 고려해보세요
$$
\varphi : \underline{\{1, \ldots, n\}} \longrightarrow \mathcal{F}|_U
$$
~에$U_\etale$에 의해 정의됨$\sigma_1, \ldots, \sigma_n$.
이 지도는 $\overline{u}$에서 줄기에 대한 전단사이다.
건설로. 하위 집합을 고려해 보자.
$$
E = \{u' \in U \mid \varphi_{\overline{u}'}\text{는 전단사}\} \subset U
$$
여기서 $\overline{u}'$는 $U$ 중 $u'$ 위에 있는 기하점이다.
(조건은 비고 \ref{etale-cohomology-remark-map-stalks}의 선택과 무관한다.)
이미지$u \in U$~의$\overline{u}$에 있습니다$E$.
$\mathcal{F}$에 대한 특수화 지도의 가정에 의해,
비고 \ref{etale-cohomology-remark-can-lift} 및
보조정리 \ref{etale-cohomology-lemma-specialization-map-pullback}
$E$가 특수화 아래에 닫혀 있는 것을 볼 수 있다.
그리고 위상 공간에서의 일반화 $U$.

\medskip\noindent $U$를 축소한 후 $U$도 유사하다고 가정할 수 있다. 강하를 통해 보조정리 \ref{descent-lemma-locally-finite-nr-irred-local-fppf} $U$에는 유한한 수의 기약 구성 요소가 있음을 알 수 있다. $u$을 통과하지 않는 기약 구성 요소를 제거한 후 $U$의 모든 기약 구성 요소가 $u$를 통과한다고 가정할 수 있다. $U$는 냉정한 위상 공간이므로 $E = U$를 따르며 $\varphi$는 정리 \ref{etale-cohomology-theorem-exactness-stalks}에 의해 동형사상라고 결론을 내립니다. 따라서 (1)는 다음과 같습니다.

\medskip\noindent (1)가 (2)을 의미하는 증명을 생략한다. \end{proof}

\begin{lemma}
\label{etale-cohomology-lemma-characterize-locally-constant-module}
$S$를 스킴으로 하여 $S$의 모든 준소형 개방이
유한한 수의 기약 구성 요소(예의 경우 $S$에
뇌터 기본 위상 공간 또는 $S$가 로컬로 뇌터인 경우.
$\Lambda$를 뇌터 환으로 설정한다.
허락하다$\mathcal{F}$가 되다층~의$\Lambda$-가군~에$S_\etale$.
다음은 동일한다.
\begin{enumerate}
\item $\mathcal{F}$는 유한 유형이고, 국소상수층
$\Lambda$-가군 및
\item $\mathcal{F}$의 모든 줄기는 유한 $\Lambda$-가군이고
모든 특수화 지도
$sp : \mathcal{F}_{\overline{s}} \to \mathcal{F}_{\overline{t}}$
전단적이다.
\end{enumerate}
\end{lemma}

\begin{proof}
이 보조정리의 증명은 다음의 증명과 동일한다.
보조정리 \ref{etale-cohomology-lemma-characterize-locally-constant}.
(2)를 가정한다. $\overline{s}$를 $S$의 기하점으로 놔두세요
$s \in S$. (1)를 증명하려면 에탈 이웃을 찾아야 한다.
$(U, \overline{u})$ $(S, \overline{s})$의 $\mathcal{F}|_U$
일정한다. 우리는 $S$가 유사하다고 가정할 수도 있고 실제로 가정한다.

\medskip\noindent
$M = \mathcal{F}_{\overline{s}}$는 유한한 $\Lambda$-가군이므로
그리고 $\Lambda$는 뇌터이므로 프레젠테이션을 선택할 수 있다.
$$
\Lambda^{\oplus m} \xrightarrow{A} \Lambda^{\oplus n} \to M \to 0
$$
일부 매트릭스의 경우$A = (a_{ji})$~와 함께계수~에$\Lambda$.
$(U, \overline{u})$ 및 요소를 선택할 수 있다.
$\sigma_1, \ldots, \sigma_n \in \mathcal{F}(U)$
그렇게$\sum a_{ji}\sigma_i = 0$~에$\mathcal{F}(U)$그리고
그런 이미지는$\sigma_i$~에$\mathcal{F}_{\overline{s}} = M$
의 표준 기본 요소의 이미지이다.
$\Lambda^n$ 위에 주어진 $M$의 프리젠테이션에서.
지도를 고려해보세요
$$
\varphi : \underline{M} \longrightarrow \mathcal{F}|_U
$$
~에$U_\etale$에 의해 정의됨$\sigma_1, \ldots, \sigma_n$.
이 지도는 $\overline{u}$에서 줄기에 대한 전단사이다.
건설로. 하위 집합을 고려해 보자.
$$
E = \{u' \in U \mid \varphi_{\overline{u}'}\text{는 전단사}\} \subset U
$$
여기서 $\overline{u}'$는 $U$ 중 $u'$ 위에 있는 기하점이다.
(조건은 비고 \ref{etale-cohomology-remark-map-stalks}의 선택과 무관한다.)
이미지$u \in U$~의$\overline{u}$에 있습니다$E$.
$\mathcal{F}$에 대한 특수화 지도의 가정에 의해,
비고 \ref{etale-cohomology-remark-can-lift} 및
보조정리 \ref{etale-cohomology-lemma-specialization-map-pullback}
$E$가 특수화 아래에 닫혀 있는 것을 볼 수 있다.
그리고 위상 공간에서의 일반화 $U$.

\medskip\noindent $U$를 축소한 후 $U$도 유사하다고 가정할 수 있다. 강하를 통해 보조정리 \ref{descent-lemma-locally-finite-nr-irred-local-fppf} $U$에는 유한한 수의 기약 구성 요소가 있음을 알 수 있다. $u$을 통과하지 않는 기약 구성 요소를 제거한 후 $U$의 모든 기약 구성 요소가 $u$를 통과한다고 가정할 수 있다. $U$는 냉정한 위상 공간이므로 $E = U$를 따르며 $\varphi$는 정리 \ref{etale-cohomology-theorem-exactness-stalks}에 의해 동형사상라고 결론을 내립니다. 따라서 (1)는 다음과 같습니다.

\medskip\noindent (1)가 (2)을 의미하는 증명을 생략한다. \end{proof}

\begin{lemma} \label{etale-cohomology-lemma-specialization-map-pushforward} $f : X \to S$를 스킴의 준컴팩트 및 준분리 사상으로 설정한다. $K \in D^+(X_\etale)$로 놔두세요. $\overline{s}$를 $S$의 기하점으로 두고 $\overline{t}$를 $\Spec(\mathcal{O}^{sh}_{S, \overline{s}})$의 기하점으로 둡니다. $$
\xymatrix{
(Rf_*K)_{\overline{s}} \ar[r]_{sp} \ar@{=}[d] &
(Rf_*K)_{\overline{t}} \ar@{=}[d] \\
R\Gamma(X \times_S \Spec(\mathcal{O}^{sh}_{S, \overline{s}}), K)
\ar[r] &
R\Gamma(X \times_S \Spec(\mathcal{O}^{sh}_{S, \overline{t}}), K)
}
$$ 아래 수평 화살표가 사상 $\text{id}_X \times c$에 의한 후퇴로 발생하는 교환 다이어그램이 있다. 여기서 $c : \Spec(\mathcal{O}^{sh}_{S, \overline{t}}) \to
\Spec(\mathcal{O}^{sh}_{S, \overline{s}})$는 사상에 도입된 비고 \ref{etale-cohomology-remark-another-sp}이다. 수직 화살표는 정리 \ref{etale-cohomology-theorem-higher-direct-images}로 지정된다. \end{lemma}

\begin{proof}
설명에 바로 이어집니다
~의$sp$~에비고 \ref{etale-cohomology-remark-another-sp}.
\end{proof}

\begin{remark}
\label{etale-cohomology-remark-specialization-map-and-fibres}
$f : X \to S$를 스킴의 사상으로 설정한다. $K \in D(X_\etale)$로 놔두세요.
$\overline{s}$를 $S$의 기하점으로 두고 $\overline{t}$로 설정한다.
$\Spec(\mathcal{O}^{sh}_{S, \overline{s}})$의 기하점이어야 한다.
$c$를 비고 \ref{etale-cohomology-remark-another-sp}와 같이 설정한다.
우리는 항상 가환 도식을 만들 수 있다.
$$
\xymatrix{
(Rf_*K)_{\overline{s}} \ar[r] \ar[d]_{sp} &
R\Gamma(X \times_S \Spec(\mathcal{O}_{S, \overline{s}}^{sh}), K)
\ar[r] \ar[d]_{(\text{id}_X \times c)^{-1}} &
R\Gamma(X_{\overline{s}}, K) \\
(Rf_*K)_{\overline{t}} \ar[r] &
R\Gamma(X \times_S \Spec(\mathcal{O}_{S, \overline{t}}^{sh}), K) \ar[r] &
R\Gamma(X_{\overline{t}}, K)
}
$$
여기서 수평 화살표는 비고 \ref{etale-cohomology-remark-stalk-fibre}의 화살표이다.
일반적으로 오른쪽에는 수직 지도가 없다.
올에 대한 $K$의 동족성 사이가 이에 적합한다.
보조정리 \ref{etale-cohomology-lemma-specialization-map-pushforward}의 경우에도 마찬가지이다.
\end{remark}















\section{복합체 구성 가능한 코호몰로지} \label{etale-cohomology-section-Dc}

\noindent
$\Lambda$를 환으로 설정한다. $D(X_\etale, \Lambda)$를 유도 범주로 표시
~의층~의$\Lambda$-가군~에$X_\etale$.
$D^b(X_\etale, \Lambda)$ (각각 $D^+$, $D^-$)로 표시한다.
경계가 있는(각각 위, 아래) 복합체의 전체 하위 범주
$D(X_\etale, \Lambda)$.

\begin{definition} \label{etale-cohomology-definition-c} $X$를 스킴으로 둡니다. $\Lambda$를 뇌터 환으로 설정한다. {\it $D_c(X_\etale, \Lambda)$}는 복합체의 $D(X_\etale, \Lambda)$의 전체 하위 범주를 나타냅니다. 그 하위 범주는 코호몰로지 층이다. 구성가능층 중 $\Lambda$-가군. \end{definition}

\noindent 이 정의는 보조정리 \ref{etale-cohomology-lemma-constructible-abelian} 및 유도 범주, 절 \ref{derived-section-triangulated-sub}에 의해 의미가 있다. 따라서 $D_c(X_\etale, \Lambda)$는 $D(X_\etale, \Lambda)$의 완전하고 포화된 삼각형 하위 범주임을 알 수 있다.

\begin{lemma} \label{etale-cohomology-lemma-restrict-and-shriek-from-etale-c} $\Lambda$를 뇌터 환으로 둡니다. 만약 $j : U \to X$는 스킴의 에탈 사상이고, \begin{enumerate} \item $K|_U \in D_c(U_\etale, \Lambda)$는 $K \in D_c(X_\etale, \Lambda)$이고, \item $j_!M \in D_c(X_\etale, \Lambda)$ 만약 $M \in D_c(U_\etale, \Lambda)$이고 사상 $j$는 준간소하고 준분리되어 있다. \end{enumerate} \end{lemma}

\begin{proof} 첫 번째 주장은 명확한다. 두 번째는 $j_!$가 정확하고 보조정리 \ref{etale-cohomology-lemma-jshriek-constructible}라는 사실에서 비롯된다. \end{proof}

\begin{lemma} \label{etale-cohomology-lemma-pullback-c} $\Lambda$를 뇌터 환으로 둡니다. $f : X \to Y$를 스킴의 사상으로 설정한다. $K \in D_c(Y_\etale, \Lambda)$이면 $Lf^*K \in D_c(X_\etale, \Lambda)$이다. \end{lemma}

\begin{proof} 이는 $f^{-1} = f^*$가 정확하고 보조정리 \ref{etale-cohomology-lemma-pullback-constructible}이므로 다음과 같습니다. \end{proof}

\begin{lemma}
\label{etale-cohomology-lemma-one-constructible}
$X$를 준간소화하고 준분리된 스킴라고 하자.
$\Lambda$를 뇌터 환으로 설정한다. $K \in D(X_\etale, \Lambda)$
그리고 $b \in \mathbf{Z}$는 $H^b(K)$이 구성 가능하도록 한다.
그러면 유한한 직합인 층 $\mathcal{F}$가 존재한다.
$j_{U!}\underline{\Lambda}$와 $U \in \Ob(X_\etale)$ 아핀 및
지도$\mathcal{F}[-b] \to K$~에$D(X_\etale, \Lambda)$
전사 $\mathcal{F} \to H^b(K)$를 유도한다.
\end{lemma}

\begin{proof}
대표하다$K$에 의해복합체 $\mathcal{K}^\bullet$~의층~의
$\Lambda$-가군. 추측을 고려해보세요
$$
\Ker(\mathcal{K}^b \to \mathcal{K}^{b + 1})
\longrightarrow
H^b(K)
$$
사이트, 보조정리에서 가군에 의해
\ref{sites-modules-lemma-module-quotient-direct-sum}
우리는 추측을 선택할 수 있습니다
$\bigoplus_{i \in I} j_{U_i!} \underline{\Lambda} \to
\Ker(\mathcal{K}^b \to \mathcal{K}^{b + 1})$
$U_i$ 아핀을 사용한다. $I' \subset I$ 유한의 경우 다음을 나타냅니다.
$\mathcal{H}_{I'} \subset H^b(K)$ 이미지
$\bigoplus_{i \in I'} j_{U_i!} \underline{\Lambda}$. 에 의해
보조정리 \ref{etale-cohomology-lemma-colimit-constructible}
우리는 그것을 본다$\mathcal{H}_{I'} = H^b(K)$어떤 사람들에게는$I' \subset I$한정된.
보조정리는 다음과 같습니다.
$\mathcal{F} = \bigoplus_{i \in I'} j_{U_i!} \underline{\Lambda}$.
\end{proof}

\begin{lemma} \label{etale-cohomology-lemma-bounded-above-c} $X$를 준컴팩트하고 준분리된 스킴라고 하자. $\Lambda$를 뇌터 환으로 설정한다. $K \in D^-(X_\etale, \Lambda)$로 놔두세요. 그런 다음 다음은 등가이다. \begin{enumerate} \item $K$는 $D_c(X_\etale, \Lambda)$에 있으며, \item $K$는 유한한 직합 위에 제한된 복합체로 표현될 수 있다. $j_{U!}\underline{\Lambda}$와 $U \in \Ob(X_\etale)$ 아핀, \item $K$는 $\Lambda$-가군의 평면 구성가능층의 경계 위의 복합체로 표현될 수 있다. \end{enumerate} \end{lemma}

\begin{proof}
(2)는 (3)를 의미하고 (3)는 (1)를 의미한다는 것이 분명한다.
$K$이 $D_c^-(X_\etale, \Lambda)$에 있다고 가정한다.
말하다$H^i(K) = 0$~을 위한$i > b$. 인덕션으로$a$
우리는 복합체 $\mathcal{F}^a \to \ldots \to \mathcal{F}^b$를 구성할 것이다.
각 $\mathcal{F}^i$는 유한한 직합이다.
$j_{U!}\underline{\Lambda}$ 및 $U \in \Ob(X_\etale)$ 아핀
그리고 동형사상을 유도하는 지도 $\mathcal{F}^\bullet \to K$
$H^i(\mathcal{F}^\bullet) \to H^i(K)$~을 위한$i > a$그리고 추측
$H^a(\mathcal{F}^\bullet) \to H^a(K)$.
$a = b$의 경우 보조정리 \ref{etale-cohomology-lemma-one-constructible}로 수행할 수 있다.
그러한 기준이 주어지면 구별되는 삼각형을 선택하십시오
$$
\mathcal{F}^\bullet \to K \to L \to \mathcal{F}^\bullet[1]
$$
그러면 우리는 그것을 본다$H^i(L) = 0$~을 위한$i \geq a$. 선택하다
$\mathcal{F}^{a - 1}[-a +1] \to L$ 에서와 같이
보조정리 \ref{etale-cohomology-lemma-one-constructible}. 구성
$\mathcal{F}^{a - 1}[-a +1] \to L \to \mathcal{F}^\bullet$
지도 $\mathcal{F}^{a - 1} \to \mathcal{F}^a$에 해당한다.
$\mathcal{F}^a \to \mathcal{F}^{a + 1}$의 구성이
0이다. TR4에 의해 지도를 얻습니다.
$$
(\mathcal{F}^{a - 1} \to \ldots \to \mathcal{F}^b) \to K
$$
$D(X_\etale, \Lambda)$에서. 이것으로 유도 단계가 완료되고
보조정리의 증명.
\end{proof}

\begin{lemma} \label{etale-cohomology-lemma-tensor-c} $X$를 스킴으로 둡니다. $\Lambda$를 뇌터 환으로 설정한다. $K, L \in D_c^-(X_\etale, \Lambda)$로 놔두세요. 그러면 $K \otimes_\Lambda^\mathbf{L} L$는 $D_c^-(X_\etale, \Lambda)$에 있다. \end{lemma}

\begin{proof} 이는 기본정리 \ref{etale-cohomology-lemma-bounded-above-c} 및 \ref{etale-cohomology-lemma-tensor-product-constructible}에서 따릅니다. \end{proof}




\section{구성 가능한 Tor 유한 코호몰로지} \label{etale-cohomology-section-ctf}

\noindent
$X$는 스킴이고 $\Lambda$는 뇌터 환이다. 자주 사용되는
유도 범주 $D_c(X_\etale, \Lambda)$ 정의의 하위 범주
절 \ref{etale-cohomology-section-Dc}는 객체로 구성된 전체 하위 범주이다.
(국지적으로) 유한한 tor 차원을 가집니다. 다음은 정식 정의이다.

\begin{definition} \label{etale-cohomology-definition-ctf} $X$를 스킴으로 둡니다. $\Lambda$를 뇌터 환으로 설정한다. {\it $D_{ctf}(X_\etale, \Lambda)$}는 국부적으로 유한한 tor 차원을 갖는 객체로 구성된 $D_c(X_\etale, \Lambda)$의 전체 하위 범주를 나타냅니다. \end{definition}

\noindent
이는 엄격하게 완전하고 포화된 삼각형 하위 범주이다.
$D_c(X_\etale, \Lambda)$ 및 $D(X_\etale, \Lambda)$. 우리의 관례에 따라 다음을 참조하라.
코호몰로지 사이트, 정의 \ref{sites-cohomology-definition-tor-amplitude},
우리는 그것을 본다
$$
D_{ctf}(X_\etale, \Lambda) \subset D^b_c(X_\etale, \Lambda)
\subset D^b(X_\etale, \Lambda)
$$
$X$가 준컴팩트인 경우. 사물에 대해 생각하는 좋은 방법
$D_{ctf}(X_\etale, \Lambda)$는 보조정리 \ref{etale-cohomology-lemma-when-ctf}에 주어집니다.

\begin{remark} \label{etale-cohomology-remark-different} 유도 범주 $D_{ctf}(X_\etale, \Lambda)$의 객체는 어떤 의미에서 $D(\mathcal{O}_X)$의 완벽한 객체보다 전역 성질이 더 좋다. 즉, $\mathcal{O}_X$-가군의 복합체는 유한한 로컬 자유 $\mathcal{O}_X$-가군의 유한 복합체에 대해 로컬적으로 준동형인 것이 발생할 수 있다. 복합체 로컬 무료 $\mathcal{O}_X$-가군. 다음 보조정리는 뇌터 스킴의 $D_{ctf}$에 대해서는 이러한 일이 발생하지 않음을 보여줍니다. \end{remark}

\begin{lemma} \label{etale-cohomology-lemma-when-ctf} $\Lambda$를 뇌터 환으로 둡니다. $X$를 준컴팩트하고 준분리된 스킴라고 하자. $K \in D(X_\etale, \Lambda)$로 놔두세요. 다음은 \begin{enumerate} \item $K \in D_{ctf}(X_\etale, \Lambda)$와 동일하며, \item $K$는 다음의 시공 가능한 평면 층의 유한 복합체로 표현될 수 있다. $\Lambda$-가군. \end{enumerate} 사실, $K$가 $[a, b]$에 토르 진폭을 갖고 있다면 $K$를 복합체 $\mathcal{F}^a \to \ldots \to \mathcal{F}^b$와 $\mathcal{F}^p$ 건설 가능한 평면 층으로 나타낼 수 있다. $\Lambda$-가군. \end{lemma}

\begin{proof} $\Lambda$-가군의 건설 가능한 평면 층의 유한 복합체가 유한한 토르 차원을 갖는 것이 분명한다. $D_c(X_\etale, \Lambda)$의 대상이라는 것도 분명하다. 따라서 (2)는 (1)를 의미함을 알 수 있다.

\medskip\noindent
(1)를 가정한다. 다음과 같이 $a, b \in \mathbf{Z}$를 선택한다.
$H^i(K \otimes_\Lambda^\mathbf{L} \mathcal{G}) = 0$ 만약
$i \not \in [a, b]$ $\Lambda$-가군 $\mathcal{G}$의 모든 층에 대해.
$b - a$에 대한 귀납법을 통해 최종 주장이 성립함을 증명하겠습니다. 만약에
$a = b$, 그러면 $K = H^a(K)[-a]$는 평평한 구성가능층이다.
결과는 유지된다. 다음으로 $b > a$를 가정한다. $K$를 나타냅니다.
복합체 $\mathcal{K}^\bullet$의 층의 $\Lambda$-가군에 의해.
추측을 고려해보세요
$$
\Ker(\mathcal{K}^b \to \mathcal{K}^{b + 1})
\longrightarrow
H^b(K)
$$
보조정리 \ref{etale-cohomology-lemma-category-constructible-modules}에 의해
우리는 무한히 많은 것을 찾을 수 있습니다아핀 스킴 $U_i$에탈 오버$X$그리고
추측 $\bigoplus j_{U_i!}\underline{\Lambda}_{U_i} \to H^b(K)$.
$U_i$을 표준 에탈 피복 $\{U_{ij} \to U_i\}$로 교체한 후
우리는 이 추측이 지도에 나타난다고 가정할 수 있다.
$\mathcal{F} = \bigoplus j_{U_i!}\underline{\Lambda}_{U_i} \to
\Ker(\mathcal{K}^b \to \mathcal{K}^{b + 1})$.
이 지도는 구별되는 삼각형을 결정합니다
$$
\mathcal{F}[-b] \to K \to L \to \mathcal{F}[-b + 1]
$$
$D(X_\etale, \Lambda)$에서. $D_{ctf}(X_\etale, \Lambda)$는 삼각측량이므로
하위 범주에는 $L$도 포함되어 있다. 실제로 $L$에는
$[a, b - 1]$의 진폭은 $\mathcal{F}$에 투영된다.
$H^b(K)$ (자세한 내용은 생략) 귀납 가설을 통해 우리는 다음을 찾을 수 있다.
유한 복합체 $\mathcal{F}^a \to \ldots \to \mathcal{F}^{b - 1}$
아파트의구성가능층~의$\Lambda$-가군대표하는$L$.
지도 $L \to \mathcal{F}[-b + 1]$는 지도에 해당한다.
$\mathcal{F}^b \to \mathcal{F}$ 이미지를 소멸시킵니다.
$\mathcal{F}^{b - 1} \to \mathcal{F}^b$. 그러면 다음과 같다
공리 TR에서3저것$K$로 표현된다.복합체
$$
\mathcal{F}^a \to \ldots \to \mathcal{F}^{b - 1} \to \mathcal{F}^b
$$
그러면 증명이 완료된다.
\end{proof}

\begin{remark}
\label{etale-cohomology-remark-projective-each-degree}
$\Lambda$를 뇌터 환으로 설정한다. $X$를 스킴으로 설정한다.
제한된복합체 $K^\bullet$건설 가능한 아파트$\Lambda$-가군
$X_\etale$에
각 줄기 $K^p_{\overline{x}}$는 유한 사영 $\Lambda$-가군이다.
따라서 복합체의 줄기는 $\Lambda$-가군의 복합체로 완벽한다.
\end{remark}

\begin{lemma} \label{etale-cohomology-lemma-restrict-and-shriek-from-etale-ctf} $\Lambda$를 뇌터 환으로 둡니다. 만약 $j : U \to X$는 스킴의 에탈 사상이고, \begin{enumerate} \item $K|_U \in D_{ctf}(U_\etale, \Lambda)$는 $K \in D_{ctf}(X_\etale, \Lambda)$이고, \item $j_!M \in D_{ctf}(X_\etale, \Lambda)$ 만약 $M \in D_{ctf}(U_\etale, \Lambda)$이고 사상 $j$는 준간소하고 준분리되어 있다. \end{enumerate} \end{lemma}

\begin{proof}
아마도 이것을 보조정리 증명하는 가장 쉬운 방법은 다음과 같이 줄이는 것이다.
$X$가 유사하고 보조정리 \ref{etale-cohomology-lemma-when-ctf}를 적용하는 경우
유한 복합체에 대한 진술로 번역하려면
$\Lambda$-가군의 플랫 구성가능층
여기서 결과는 보조정리 \ref{etale-cohomology-lemma-jshriek-constructible}이다.
\end{proof}

\begin{lemma} \label{etale-cohomology-lemma-pullback-ctf} $\Lambda$를 뇌터 환으로 둡니다. $f : X \to Y$를 스킴의 사상으로 설정한다. $K \in D_{ctf}(Y_\etale, \Lambda)$이면 $Lf^*K \in D_{ctf}(X_\etale, \Lambda)$이다. \end{lemma}

\begin{proof} 보조정리 \ref{etale-cohomology-lemma-when-ctf}를 적용하여 $\Lambda$-가군의 플랫 구성가능층의 유한 복합체에 대한 질문으로 줄이다. 그런 다음 $f^{-1} = f^*$가 정확하고 보조정리 \ref{etale-cohomology-lemma-pullback-constructible}라는 명령문이 따릅니다. \end{proof}

\begin{lemma}
\label{etale-cohomology-lemma-connected-ctf-locally-constant}
$X$를 연결 스킴으로 설정한다. $\Lambda$를 뇌터 환으로 설정한다.
$K \in D_{ctf}(X_\etale, \Lambda)$이 지역적으로 상수 코호몰로지 층을 갖는다고 가정한다.
그러면 유한 사영 $\Lambda$-가군의 유한 복합체가 존재한다.
$M^\bullet$ 및 에탈 피복 $\{U_i \to X\}$
$K|_{U_i} \cong \underline{M^\bullet}|_{U_i}$~에$D(U_{i, \etale}, \Lambda)$.
\end{lemma}

\begin{proof}
에탈 피복 $\{U_i \to X\}$을 선택하여 $K|_{U_i}$
상수이다. $K|_{U_i} \cong \underline{M_i^\bullet}_{U_i}$라고 말한다.
유한 $\Lambda$-가군 $M_i^\bullet$의 유한 복합체에 대해.
사이트, 보조정리의 코호몰로지를 참조하라.
\ref{sites-cohomology-lemma-locally-constant}.
그것을 관찰하십시오$U_i \times_X U_j$비어 있는 경우$M_i^\bullet$
와 동형이 아니다$M_j^\bullet$~에$D(\Lambda)$.
$\Lambda$-가군 $M^\bullet$의 각 복합체에 대해
$I_{M^\bullet} =
\{i \in I \mid M_i^\bullet \cong M^\bullet\text{(}D(\Lambda)\text{에서)}\}$.
에탈 사상이 열려 있음을 알 수 있다.
$U_{M^\bullet} = \bigcup_{i \in I_{M^\bullet}} \Im(U_i \to X)$
$X$의 열린 부분집합이다. 그러면 $X = \coprod U_{M^\bullet}$는 분리된
$X$의 피복을 엽니다. $X$는 연결 단 하나의 $U_{M^\bullet}$이므로
비어 있지 않습니다. $K$가 $D_{ctf}(X_\etale, \Lambda)$에 있으므로 $M^\bullet$가 표시된다.
$\Lambda$-가군의 완벽한 복합체이다.
대수학에 대한 자세한 내용은 보조정리 \ref{more-algebra-lemma-perfect}이다.
그러므로 우리는 $M^\bullet$가 유한 사영의 유한 복합체라고 가정할 수 있다.
$\Lambda$-가군.
\end{proof}








\section{꼬임층}
\label{etale-cohomology-section-torsion}

\noindent 비틀림 아벨 층과 그 에탈 코호몰로지에 대한 간략한 절이다. $\mathcal{C}$를 사이트로 설정한다. 우리는 사이트, 보조정리 \ref{sites-cohomology-lemma-torsion}의 코호몰로지에서 코호몰로지 층이 꼬임층인 $D(\mathcal{C})$의 모든 개체를 다음으로 나타낼 수 있음을 보여주었습니다. 복합체 그 용어는 모두 비틀림이다.

\begin{lemma}
\label{etale-cohomology-lemma-torsion-cohomology}
$X$를 준간소화하고 준분리된 스킴라고 하자.
\begin{enumerate}
\item $\mathcal{F}$가 $X_\etale$의 비틀림 아벨 층인 경우
$H^n_\etale(X, \mathcal{F})$는 모든 $n$에 대한 비틀림 아벨 군이다.
\item $K$ 에서 $D^+(X_\etale)$에 비틀림 코호몰로지 층이 있는 경우
$H^n_\etale(X, K)$는 모든 $n$에 대한 비틀림 아벨 군이다.
\end{enumerate}
\end{lemma}

\begin{proof} (1)를 증명하기 위해 $\mathcal{F} = \bigcup \mathcal{F}[n]$라고 씁니다. 여기서 $\mathcal{F}[d]$는 $d$ 비틀림 서브쉬프이다. 보조정리 \ref{etale-cohomology-lemma-colimit}에 의해 $H^n_\etale(X, \mathcal{F}) = \colim H^n_\etale(X, \mathcal{F}[d])$이 된다. 이는 1)가 $H^n_\etale(X, \mathcal{F}[d])$가 $d$에 의해 소멸됨을 증명한다.

\medskip\noindent (2)를 증명하기 위해 스펙트럼 열 $E_2^{p, q} = H^p_\etale(X, H^q(K))$를 사용하여 $H^n_\etale(X, K)$ (유도 범주, 보조정리로 수렴할 수 있다. \ref{derived-lemma-two-ss-complex-functor}) 및 층에 대한 결과이다. \end{proof}

\begin{lemma}
\label{etale-cohomology-lemma-torsion-direct-image}
$f : X \to Y$를 준컴팩트하고 준분리된 것으로 하자.
스킴의 사상.
\begin{enumerate}
\item $\mathcal{F}$가 $X_\etale$의 비틀림 아벨 층인 경우
$R^nf_*\mathcal{F}$비틀림이다아벨 층~에$Y_\etale$모두를 위해$n$.
\item $K$ 에서 $D^+(X_\etale)$에 비틀림 코호몰로지 층이 있는 경우
$Rf_*K$는 $D^+(Y_\etale)$의 객체이며, 그 객체는 코호몰로지 층이다.
비틀림 아벨 층.
\end{enumerate}
\end{lemma}

\begin{proof}
증명 (1). $R^nf_*\mathcal{F}$는 층과 연관되어 있음을 기억하세요.
전층 $V \mapsto H^n_\etale(X \times_Y V, \mathcal{F})$로
$Y_\etale$에 있다. 사이트의 코호몰로지를 참조하라.
보조정리 \ref{sites-cohomology-lemma-higher-direct-images}.
$V$ 아핀을 선택하면 $X \times_Y V$는 다음과 같습니다.
$f$이므로 준컴팩트하고 준분리되어 있다.
보조정리 \ref{etale-cohomology-lemma-torsion-cohomology}를 적용하여 확인할 수 있다.
$H^n_\etale(X \times_Y V, \mathcal{F})$는 비틀림이다.

\medskip\noindent
증명 (2). $R^nf_*K$는 층과 연관되어 있음을 기억하세요.
전층 $V \mapsto H^n_\etale(X \times_Y V, K)$로
$Y_\etale$에 있다. 사이트의 코호몰로지를 참조하라.
보조정리 \ref{sites-cohomology-lemma-unbounded-describe-higher-direct-images}.
$V$ 아핀을 선택하면 $X \times_Y V$는 다음과 같습니다.
$f$이므로 준컴팩트하고 준분리되어 있다.
보조정리 \ref{etale-cohomology-lemma-torsion-cohomology}를 적용하여 확인할 수 있다.
$H^n_\etale(X \times_Y V, K)$는 비틀림이다.
\end{proof}









\section{지지집합을 가진 코호몰로지 닫힌 하위 체계} \label{etale-cohomology-section-cohomology-support}

\noindent
$X$를 스킴으로 두고 $Z \subset X$를 폐쇄형 하위 체계로 둡니다.
$\mathcal{F}$를 $X_\etale$에서 아벨 층으로 설정한다. 우리는
$$
\Gamma_Z(X, \mathcal{F}) =
\{s \in \mathcal{F}(X) \mid \text{Supp}(s) \subset Z\}
$$
$Z$ (정의 \ref{etale-cohomology-definition-support})에서 지지집합과 절이 된다.
이것은 일반적으로 정확하지 않은 왼쪽 완전 함자이다.
따라서 우리는 유도 함자를 얻습니다.
$$
R\Gamma_Z(X, -) : D(X_\etale) \longrightarrow D(\textit{Ab})
$$
및 $Z$에 지지집합이 있는 코호몰로지 군은 다음과 같이 정의된다.
$H^q_Z(X, \mathcal{F}) = R^q\Gamma_Z(X, \mathcal{F})$.

\medskip\noindent
$\mathcal{I}$를 $X_\etale$의 단사 아벨 층으로 설정한다.
$U = X \setminus Z$로 놔두세요.
그런 다음 제한 사상 $\mathcal{I}(X) \to \mathcal{I}(U)$는 전사이다.
(코호몰로지 위의 사이트, 보조정리
\ref{sites-cohomology-lemma-restriction-along-monomorphism-surjective})
핵 $\Gamma_Z(X, \mathcal{I})$로. 그것은 즉시 다음과 같습니다
$K \in D(X_\etale)$에는 구별되는 삼각형이 있다.
$$
R\Gamma_Z(X, K) \to R\Gamma(X, K) \to R\Gamma(U, K) \to R\Gamma_Z(X, K)[1]
$$
$D(\textit{Ab})$에서. 결과적으로 우리는 긴 정확한 코호몰로지를 얻습니다.
순서
$$
\ldots \to H^i_Z(X, K) \to H^i(X, K) \to H^i(U, K) \to
H^{i + 1}_Z(X, K) \to \ldots
$$
누구에게나$K$~에$D(X_\etale)$.

\medskip\noindent
대한아벨 층 $\mathcal{F}$~에$X_\etale$우리는 고려할 수 있다
{\it 절의 하위 단은 $Z$}의 지지집합으로 표시된다.
$\mathcal{H}_Z(\mathcal{F})$, 규칙에 의해 정의됨
$$
\mathcal{H}_Z(\mathcal{F})(U) =
\{s \in \mathcal{F}(U) \mid \text{Supp}(s) \subset U \times_X Z\}
$$
여기서는 정의 \ref{etale-cohomology-definition-support}의 절의 지지집합을 사용한다.
동등성을 사용하여
명제 \ref{etale-cohomology-proposition-closed-immersion-pushforward}
$\mathcal{H}_Z(\mathcal{F})$를 아벨 층으로 볼 수 있다.
$Z_\etale$. 따라서 우리는 함자를 얻습니다.
$$
\textit{Ab}(X_\etale) \longrightarrow \textit{Ab}(Z_\etale),\quad
\mathcal{F} \longmapsto \mathcal{H}_Z(\mathcal{F})
$$
이는 왼쪽 완전이지만 일반적으로 정확하지는 않습니다.

\begin{lemma} \label{etale-cohomology-lemma-sections-with-support-acyclic} $i : Z \to X$를 스킴의 닫힌 몰입으로 둡니다. $\mathcal{I}$를 $X_\etale$의 단사 아벨 층으로 둡니다. 그러면 $\mathcal{H}_Z(\mathcal{I})$는 $Z_\etale$의 단사 아벨 층이다. \end{lemma}

\begin{proof}
어떤 경우에도 이를 관찰하세요.아벨 층 $\mathcal{G}$~에$Z_\etale$
우리는
$$
\Hom_Z(\mathcal{G}, \mathcal{H}_Z(\mathcal{F})) =
\Hom_X(i_*\mathcal{G}, \mathcal{F})
$$
결국 $i_*\mathcal{G}$의 모든 절에는 $Z$에 지지집합이 있기 때문이다.
$i_*$는 정확하기 때문에 (절 \ref{etale-cohomology-section-closed-immersions})
$\mathcal{I}$는 $X_\etale$에 대한 단사형이다.
$\mathcal{H}_Z(\mathcal{I})$는 $Z_\etale$에 대한 분사형이다.
\end{proof}

\noindent $$
R\mathcal{H}_Z : D(X_\etale) \longrightarrow D(Z_\etale)
$$는 유도 함자를 나타냅니다. 집합 $\mathcal{H}^q_Z(\mathcal{F}) = R^q\mathcal{H}_Z(\mathcal{F})$하여 $\mathcal{H}^0_Z(\mathcal{F}) = \mathcal{H}_Z(\mathcal{F})$한다. 위의 보조정리에는 Grothendieck 스펙트럼 열 $$
E_2^{p, q} = H^p(Z, \mathcal{H}^q_Z(\mathcal{F}))
\Rightarrow H^{p + q}_Z(X, \mathcal{F})
$$가 있다.

\begin{lemma}
\label{etale-cohomology-lemma-cohomology-with-support-sheaf-on-support}
$i : Z \to X$를 스킴의 닫힌 몰입으로 설정한다.
$\mathcal{G}$를 $Z_\etale$의 단사 아벨 층으로 설정한다.
그 다음에$\mathcal{H}^p_Z(i_*\mathcal{G}) = 0$~을 위한$p > 0$.
\end{lemma}

\begin{proof} 이는 함자 $i_*$가 정확하고 단사 아벨 층을 단사 아벨 층으로 변환하기 때문에 사실입니다 (코호몰로지 위의 사이트, 보조정리 \ref{sites-cohomology-lemma-pushforward-injective-flat}). \end{proof}

\begin{lemma}
\label{etale-cohomology-lemma-cohomology-with-support-triangle}
$i : Z \to X$를 스킴의 닫힌 몰입으로 설정한다.
$j : U \to X$를 $Z$의 보수를 포함한다고 하자.
$\mathcal{F}$를 $X_\etale$에서 아벨 층으로 설정한다.
뚜렷한 삼각형이 있다
$$
i_*R\mathcal{H}_Z(\mathcal{F}) \to \mathcal{F} \to Rj_*(\mathcal{F}|_U) \to
i_*R\mathcal{H}_Z(\mathcal{F})[1]
$$
$D(X_\etale)$에서. 그러면 완전열이 생성된다.
$$
0 \to i_*\mathcal{H}_Z(\mathcal{F}) \to \mathcal{F} \to
j_*(\mathcal{F}|_U) \to i_*\mathcal{H}^1_Z(\mathcal{F}) \to 0
$$
그리고 동형사상
$R^pj_*(\mathcal{F}|_U) \cong i_*\mathcal{H}^{p + 1}_Z(\mathcal{F})$
$p \geq 1$의 경우.
\end{lemma}

\begin{proof} 구별된 삼각형을 얻으려면 단사 분해 $\mathcal{F} \to \mathcal{I}^\bullet$를 선택한다. 그런 다음 위의 논의를 통해 복합체 $$
0 \to
i_*\mathcal{H}_Z(\mathcal{I}^\bullet) \to \mathcal{I}^\bullet
\to j_*(\mathcal{I}^\bullet|_U) \to 0
$$의 짧고 정확한 시퀀스를 얻습니다. 따라서 유도 범주, 절 \ref{derived-section-canonical-delta-functor}에 의해 구별되는 삼각형이다. \end{proof}

\noindent $X$를 스킴으로 두고 $Z \subset X$를 폐쇄형 하위 체계로 둡니다. $D_Z(X_\etale)$는 복합체로 구성된 $D(X_\etale)$의 완전 포화 삼각분할 하위 범주로 표시되며, 코호몰로지 층은 $Z$에서 지원된다. $D_Z(X_\etale)$는 $X$의 기본 닫힌 부분집합에만 의존한다는 점에 유의하세요.

\begin{lemma}
\label{etale-cohomology-lemma-complexes-with-support-on-closed}
$i : Z \to X$를 스킴의 닫힌 몰입으로 설정한다.
지도 $Ri_{small, *} = i_{small, *} : D(Z_\etale) \to D(X_\etale)$
준역원으로 등가성 $D(Z_\etale) \to D_Z(X_\etale)$을 유도한다.
$$
i_{small}^{-1}|_{D_Z(X_\etale)} = R\mathcal{H}_Z|_{D_Z(X_\etale)}
$$
\end{lemma}

\begin{proof} $i_{small}^{-1}$ 및 $i_{small, *}$는 완전 함자의 인접 쌍이므로 $i_{small}^{-1}i_{small, *}$는 아벨 층의 식별 함자와 동형이다. 명제 \ref{etale-cohomology-proposition-closed-immersion-pushforward} 및 보조정리 \ref{etale-cohomology-lemma-stalk-pullback}를 참조하라. 따라서 $i_{small, *} : D(Z_\etale) \to D_Z(X_\etale)$은 완전히 충실하고 $i_{small}^{-1}$는 왼쪽 역원을 결정한다. 한편, $K$이 $D_Z(X_\etale)$의 객체라고 가정하고, 부속 맵 $K \to i_{small, *}i_{small}^{-1}K$을 고려해보자. $i_{small, *}$ 및 $i_{small}^{-1}$의 완전성을 사용하면 코호몰로지 층에 부속 맵 $H^n(K) \to i_{small, *}i_{small}^{-1}H^n(K)$이 유도된다. 이러한 코호몰로지 층은 $Z$에서 지원되므로 이러한 부속 맵은 동형사상이며 $D(Z_\etale) \to D_Z(X_\etale)$는 동등하다고 결론을 내립니다.

\medskip\noindent
증명을 끝내려면 $R\mathcal{H}_Z(K) = i_{small}^{-1}K$를 보여야 한다.
$K$가 $D_Z(X_\etale)$의 객체인 경우. 이를 위해 우리는 그것을 사용할 수 있습니다
$K = i_{small, *}i_{small}^{-1}K$
우리가 방금 증명했듯이 이것이 사실이다. 그러면 우리는
K-주사 대표자 $\mathcal{I}^\bullet$를 선택할 수 있다.
$i_{small}^{-1}K$.
$i_{small, *}$는 오른쪽 수반에서 완전 함자이므로
$i_{small}^{-1}$,
복합체 $i_{small, *}\mathcal{I}^\bullet$는 K-단사형이다.
(유도 범주, 보조정리 \ref{derived-lemma-adjoint-preserve-K-injectives}).
$R\mathcal{H}_Z(K)$는 다음과 같이 계산된다.
$\mathcal{H}_Z(i_{small, *}\mathcal{I}^\bullet) = \mathcal{I}^\bullet$
원하는대로.
\end{proof}

\begin{lemma}
\label{etale-cohomology-lemma-cohomology-with-support-quasi-coherent}
$X$를 스킴으로 설정한다. $Z \subset X$를 폐쇄형 하위 구성이라고 가정한다.
$\mathcal{F}$를 준일관적 $\mathcal{O}_X$-가군으로 설정한다.
그리고 $\mathcal{F}^a$ 관련 준연접층을 표시한다.
$X$의 작은 에탈 사이트에
(명제 \ref{etale-cohomology-proposition-quasi-coherent-sheaf-fpqc}). 그 다음에
\begin{enumerate}
\item $H^q_Z(X, \mathcal{F})$는 $H^q_Z(X_\etale, \mathcal{F}^a)$에 동의한다.
\item $Z$의 보수가 $X$에서 소급 압축되면
$i_*\mathcal{H}^q_Z(\mathcal{F}^a)$는 다음의 준연접층이다.
$\mathcal{O}_X$-가군은 $(i_*\mathcal{H}^q_Z(\mathcal{F}))^a$과 같습니다.
\end{enumerate}
\end{lemma}

\begin{proof} $j : U \to X$를 $Z$의 보수 포함이라고 하자. 코호몰로지 군의 문 (1)은 코호몰로지에 대한 긴 완전열에서 지원 및 $H^q(X_\etale, \mathcal{F}^a) = H^q(X, \mathcal{F})$ 및 $H^q(U_\etale, \mathcal{F}^a) = H^q(U, \mathcal{F})$에 대한 동의를 따릅니다. 정리 \ref{etale-cohomology-theorem-zariski-fpqc-quasi-coherent}를 참조하라. $j : U \to X$가 준컴팩트 사상인 경우, 즉 $U \subset X$가 레트로컴팩트인 경우 $R^qj_*$는 준연접층을 준연접층(코호몰로지로 변환한다. 스킴, 보조정리 \ref{coherent-lemma-quasi-coherence-higher-direct-images}) 에탈 사이트 (하강, 보조정리 \ref{descent-lemma-higher-direct-images-small-etale})에서 관련 층을 타고 통근한다. 보조정리 \ref{etale-cohomology-lemma-cohomology-with-support-triangle}를 적용하여 결론을 내립니다. \end{proof}







\section{스킴 엄밀 Hensel 국소환} \label{etale-cohomology-section-strictly-henselian-local-rings}

\noindent 이 절에서 우리는 국소환이 엄격하게 헨젤리안인 스킴의 에탈 코호몰로지에 대한 몇 가지 결과를 수집한다. 예에 대해 보조정리 \ref{etale-cohomology-lemma-vanishing-etale-cohomology-strictly-henselian}의 재미있는 일반화는 다음과 같습니다.

\begin{lemma}
\label{etale-cohomology-lemma-local-rings-strictly-henselian}
$S$는 모두 국소환이 엄격하게 헨젤리안인 스킴라고 가정한다.
그렇다면 누구에게나아벨 층 $\mathcal{F}$~에$S_\etale$우리는
$H^i(S_\etale, \mathcal{F}) = H^i(S_{Zar}, \mathcal{F})$.
\end{lemma}

\begin{proof} $\epsilon : S_\etale \to S_{Zar}$를 함자를 포함하여 주어지는 사이트의 사상으로 둡니다. 자리스키 층 $R^p\epsilon_*\mathcal{F}$는 전층 $U \mapsto H^p_\etale(U, \mathcal{F})$와 연관된 층이다. 따라서 $x \in X$의 줄기는 $\colim H^p_\etale(U, \mathcal{F}) =
H^p_\etale(\Spec(\mathcal{O}_{X, x}), \mathcal{G}_x)$이다. 여기서 $\mathcal{G}_x$는 $\mathcal{F}$에서 $\Spec(\mathcal{O}_{X, x})$로의 철회를 나타냅니다. 보조정리 \ref{etale-cohomology-lemma-directed-colimit-cohomology}를 참조하라. 따라서 $R^p\epsilon_*\mathcal{F}$의 고차 직상은 보조정리 \ref{etale-cohomology-lemma-vanishing-etale-cohomology-strictly-henselian}에 의해 0이고 Leray 스펙트럼 열로 결론을 내립니다. \end{proof}

\begin{lemma}
\label{etale-cohomology-lemma-gabber-criterion}
$R$는 모두 국소환이 엄격하게 헨젤리안인 환라고 가정한다.
$\mathcal{F}$를 $\Spec(R)_\etale$에서 층으로 설정한다.
모든 $f, g \in R$에 대해 핵
$$
H^1_\etale(D(f + g), \mathcal{F})
\longrightarrow
H^1_\etale(D(f(f + g)), \mathcal{F}) \oplus
H^1_\etale(D(g(f + g)), \mathcal{F})
$$
0이다. 그 다음에$H^q_\etale(\Spec(R), \mathcal{F}) = 0$~을 위한$q > 0$.
\end{lemma}

\begin{proof}
보조정리 \ref{etale-cohomology-lemma-local-rings-strictly-henselian}에 의해 우리는
$\mathcal{F}$의 에탈 코호몰로지는 Zariski 코호몰로지에 동의한다.
$\Spec(R)$을 열 때마다. $i$에 대한 귀납법으로 증명하겠습니다.
명령문: $h \in R$에 대해 $H^q_\etale(D(h), \mathcal{F}) = 0$이 있다.
$1 \leq q \leq i$의 경우. 기본 사례 $i = 0$는 간단한다. $i \geq 1$를 가정한다.

\medskip\noindent
허락하다$\xi \in H^q_\etale(D(h), \mathcal{F})$어떤 사람들에게는$1 \leq q \leq i$
및 $h \in R$. $q < i$이면 귀납법이 완료되었으므로 $q = i$라고 가정한다.
$R$을 $R_h$로 바꾼 후 다음과 같이 가정할 수 있다.
$\xi \in H^i_\etale(\Spec(R), \mathcal{F})$; 일부 세부 사항은 생략되었다.
$I \subset R$를 $f \in R$ 요소의 집합으로 설정한다.
$\xi|_{D(f)} = 0$. $\xi$는 Zariski 지역적으로 사소하기 때문에 다음과 같습니다.
그것은 모든 소수에 대해$\mathfrak p$~의$R$존재한다$f \in I$
$f \not \in \mathfrak p$로. 따라서 $I$이
아이디얼, $1 \in I$ 그러면 끝이다. 그것은 분명하다
$f \in I$, $r \in R$는 $rf \in I$를 의미한다. 따라서 우리는 $f, g \in I$
그리고 우리는 $f + g \in I$를 보여줍니다. $q = i = 1$인 경우 이는 정확히
가정의 보조정리! $i = 1$에 대한 결과이다. 을 위한
$q = i > 1$, 참고하세요
$$
D(f + g) = D(f(f + g)) \cup D(g(f + g))
$$
Mayer-Vietoris 작성 (코호몰로지, 보조정리 \ref{cohomology-lemma-mayer-vietoris}
이는 $\Spec(R)$의 개방형 하위 구성에서 에탈 코호몰로지로 적용된다.
자리스키 코호몰로지) 완전열
$$
\xymatrix{
H^{i - 1}_\etale(D(fg(f + g)), \mathcal{F}) \ar[d] \\
H^i_\etale(D(f + g), \mathcal{F}) \ar[d] \\
H^i_\etale(D(f(f + g)), \mathcal{F}) \oplus
H^i_\etale(D(g(f + g)), \mathcal{F})
}
$$
첫 번째 그룹은 유도에 의해 0이 된다.
\end{proof}

\begin{lemma}
\label{etale-cohomology-lemma-affine-only-closed-points}
$S$를 다음과 같이 아핀 스킴으로 설정한다.
(1) 모든 점은 닫혀 있고 (2) 모든 나머지 체는 분리 가능한다.
대수적으로 닫혀있다. 그 다음에
누구에게나아벨 층 $\mathcal{F}$~에$S_\etale$우리는
$H^i(S_\etale, \mathcal{F}) = 0$~을 위한$i > 0$.
\end{lemma}

\begin{proof} 조건 (1)는 $S$의 기본 위상 공간이 유한하다는 것을 의미한다. 대수학, 보조정리 \ref{algebra-lemma-ring-with-only-minimal-primes}를 참조하라. 따라서 위상 공간 $S$(즉, Zariski 코호몰로지)에서 아벨 층의 상위 코호몰로지 군은 사소한 것이다. 코호몰로지, 보조정리 \ref{cohomology-lemma-vanishing-for-profinite}를 참조하라. 국소환은 대수학, 보조정리 \ref{algebra-lemma-local-dimension-zero-henselian}에 의해 엄격하게 헨셀식이다. 따라서 $S$의 에탈 코호몰로지는 Zariski에 의해 코호몰로지에 의해 보조정리 \ref{etale-cohomology-lemma-local-rings-strictly-henselian}로 계산되고 증명이 완료된다. \end{proof}

\noindent 절대정수적으로 닫힌 환의 스펙트럼은 스킴의 예이며, 그 모두는 국소환이 엄밀하게 헨젤리안이다. 대수학에 대한 추가 정보, 보조정리를 참조하라. \ref{more-algebra-lemma-absolutely-integrally-closed-strictly-henselian}. 분리가능하게 닫힌 분수 체를 갖는 정규 도메인은 다음 보조정리에 설명된 바와 같이 훨씬 더 강한 성질을 갖는 것으로 나타났습니다.

\begin{lemma} \label{etale-cohomology-lemma-normal-scheme-with-alg-closed-function-field} $X$는 분리 가능하게 닫힌 기능 체가 있는 정역 정규 스킴으로 설정한다. \begin{enumerate} \item 분리된 에탈 사상 $U \to X$는 열린 몰입의 분리된 합집합이다. \item $X$의 모든 국소환은 엄격히 헨셀식이다. \end{enumerate} \end{lemma}

\begin{proof}
$R$는 분수 체가 대수적으로 분리 가능한 정규 도메인이라고 가정한다.
닫은. $R \to A$를 에탈 환 맵이라고 하자. 그 다음에
$A \otimes_R K$는 $K$-대수와 마찬가지로 유한 곱이다.
$\prod_{i = 1, \ldots, n} K$의 $K$ 사본. $e_i$, $i = 1, \ldots, n$
$A \otimes_R K$의 해당 멱등원이 된다. $A$는 정규이므로
(대수학, 보조정리 \ref{algebra-lemma-normal-goes-up})
멱등원 $e_i$은 $A$에 있다.
(대수학, 보조정리 \ref{algebra-lemma-normal-ring-integrally-closed}).
따라서 $A = \prod Ae_i$이고 $A \otimes_R K = K$라고 가정할 수 있다.
$A \subset A \otimes_R K = K$ 이후 ($R \to A$의 평탄도에 따라
$R \subset K$) 이후 $A$은 도메인이라고 결론을 내립니다.
같은 주장으로 우리는 다음과 같이 결론을 내린다.
$A \otimes_R A \subset (A \otimes_R A) \otimes_R K = K$.
지도 $A \otimes_R A \to A$는 다음과 같습니다.
주사뿐만 아니라 전사. 따라서 $R \to A$는
열린 몰입
사상, 보조정리 \ref{morphisms-lemma-universally-injective}
그리고
에탈 사상, 정리 \ref{etale-theorem-etale-radicial-open}.

\medskip\noindent
$f : U \to X$를 분리된 에탈 사상이라고 하자. $\eta \in X$
일반적인 지점이 되고 $f^{-1}(\{\eta\}) = \{\xi_i\}_{i \in I}$로 둡니다.
이전 단락의 결과는 다음과 같습니다.
모든 아핀 개방에 대해$U' \subset U$누구의 이미지인가$X$에 포함되어 있습니다
우리는 $U' = \coprod_{i \in I} U'_i$를 가지고 있다. 여기서 $U'_i$
은집합~의$U'$어느 것특수화~의$\xi_i$.
또한, 사상 $U'_i \to X$는 열린 몰입이다.
$U_i = \overline{\{\xi_i\}}$는 열려 있고 닫혀 있다.
$U$의 하위 구성표 및 $U_i \to X$가 소스에 로컬로 있음
동형사상. 사상에 의해,
보조정리 \ref{morphisms-lemma-distinct-local-rings}
$U_i \to X$이 분리되어 있다는 사실은 다음을 의미한다.
$U_i \to X$는 단사적이므로 $U_i \to X$라고 결론을 내립니다.
열린 몰입, 즉 (1)가 성립한다.

\medskip\noindent
부분(2)은 부분(1)과 엄격한
헨젤화$\mathcal{O}_{X, x}$으로국소환~에$\overline{x}$
$X$ (보조정리 \ref{etale-cohomology-lemma-describe-etale-local-ring})의 에탈 사이트에 있다.
직접적으로 증명할 수도 있다.
기본 그룹, 보조정리
\ref{pione-lemma-normal-local-domain-separablly-closed-fraction-field}.
\end{proof}

\begin{lemma}
\label{etale-cohomology-lemma-Rf-star-zero-normal-with-alg-closed-function-field}
$f : X \to Y$를 스킴의 사상으로 설정한다. 여기서 $X$는 정역이다.
정규 스킴 분리 가능하게 닫힌 기능 체가 있다. 그 다음에
$R^qf_*\underline{M} = 0$~을 위한$q > 0$그리고 어떤아벨 군 $M$.
\end{lemma}

\begin{proof}
$R^qf_*\underline{M}$는 층과 연관되어 있음을 기억하세요.
에전층 $V \mapsto H^q_\etale(V \times_Y X, M)$~에$Y_\etale$, 보다
보조정리 \ref{etale-cohomology-lemma-higher-direct-images}.
$V$가 유사하면 $V \times_Y X \to X$가 분리되어 에탈이 된다.
따라서 $V \times_Y X = \coprod U_i$는 개방형의 분리된 합집합이다.
하위 계획$U_i$~의$X$, 보다
보조정리 \ref{etale-cohomology-lemma-normal-scheme-with-alg-closed-function-field}.
보조정리 \ref{etale-cohomology-lemma-local-rings-strictly-henselian}에 의해
$H^q_\etale(U_i, M)$는 다음과 같습니다.
$H^q_{Zar}(U_i, M)$. 이것은 사라집니다
코호몰로지, 보조정리 \ref{cohomology-lemma-irreducible-constant-cohomology-zero}.
\end{proof}

\begin{lemma} \label{etale-cohomology-lemma-closed-of-affine-normal-scheme-with-alg-closed-function-field} $X$를 분리 가능하게 닫힌 함수 체가 있는 아핀 정역 정규 스킴으로 설정한다. $Z \subset X$를 폐쇄형 하위 구성이라고 가정한다. $V \to Z$는 $V$ 아핀을 갖는 에탈 사상라고 하자. 그러면 $V$는 $Z$의 개방형 하위 구성표의 유한 분리 결합이다. $V \to Z$가 전사이고 유한 에탈이면 $V \to Z$에는 절이 있다. \end{lemma}

\begin{proof} 대수학으로 보조정리 \ref{algebra-lemma-lift-etale} $V$를 아핀 스킴 $U$ 에탈로 $X$로 끌어올릴 수 있다. 보조정리 \ref{etale-cohomology-lemma-normal-scheme-with-alg-closed-function-field}를 $U \to X$에 적용하여 첫 번째 명령문을 가져옵니다.

\medskip\noindent
마지막 진술은 첫 번째 진술의 결과이다.
$V = \coprod_{i = 1, \ldots, n} V_i$를 유한하다고 하자
개방형으로 분해
$V_i \to Z$ 및 열린 몰입을 사용하여 하위 계획을 닫았습니다.
$V \to Z$는 유한하므로 $V_i \to Z$도 닫혀 있음을 알 수 있다.
$U_i \subset Z$를 이미지로 둡니다. 그런 다음 분해가 있다.
개방형 및 폐쇄형 하위 계획으로
$$
Z =
\coprod\nolimits_{(A, B)}
\bigcap\nolimits_{i \in A} U_i \cap
\bigcap\nolimits_{i \in B} U_i^c
$$
여기서 분리된 합집합은 $\{1, \ldots, n\} = A \amalg B$ 위에 있다.
여기서 $A$에는 하나 이상의 요소가 있다.
각 지층은 단일 $U_i$에 포함되어 있으며
절을 찾습니다.
\end{proof}

\begin{lemma} \label{etale-cohomology-lemma-gabber-for-h1-absolutely-algebraically-closed} $X$는 분리 가능하게 닫힌 기능 체가 있는 정규 정역 아핀 스킴으로 설정한다. $Z \subset X$를 폐쇄형 하위 구성이라고 가정한다. 임의의 유한 아벨 군 $M$에 대해 $H^1_\etale(Z, \underline{M}) = 0$가 있다. \end{lemma}

\begin{proof}
사이트에 코호몰로지로, 보조정리 \ref{sites-cohomology-lemma-torsors-h1}
$H^1_\etale(Z, \underline{M})$의 요소는
$\underline{M}$-$\mathcal{F}$를 $Z_\etale$에 저장한다.
그러한 토서는 분명히 유한한 국소상수층이다.
따라서 $\mathcal{F}$는 스킴 $V$ 유한으로 표현 가능한다.
$Z$, 보조정리 \ref{etale-cohomology-lemma-characterize-finite-locally-constant}에 대한 에탈이다.
물론 $V \to Z$는 토르서가 지역적으로 사소하기 때문에 전사이다.
$V \to Z$에는 절이 있으므로
보조정리 \ref{etale-cohomology-lemma-closed-of-affine-normal-scheme-with-alg-closed-function-field}
우리는 끝났습니다.
\end{proof}

\begin{lemma}
\label{etale-cohomology-lemma-gabber-for-absolutely-algebraically-closed}
$X$를 분리 가능하게 닫힌 정규 정역 아핀 스킴으로 설정한다.
기능 체. $Z \subset X$를 폐쇄형 하위 구성이라고 가정한다.
유한한 아벨 군 $M$에 대해 우리는
$H^q_\etale(Z, \underline{M}) = 0$~을 위한$q \geq 1$.
\end{lemma}

\begin{proof} $X = \Spec(R)$ 및 $Z = \Spec(R')$를 작성하여 환 $R \to R'$의 전사를 갖도록 한다. $R'$의 모든 국소환은 보조정리 \ref{etale-cohomology-lemma-normal-scheme-with-alg-closed-function-field} 및 대수학, 보조정리 \ref{algebra-lemma-quotient-strict-henselization}에 의해 엄격하게 헨셀식이다. 또한 모든 $f' \in R'$에 대해 $R_f \to R'_{f'}$라는 전사가 있음을 알 수 있다. 여기서 $f \in R$는 $f'$의 리프트이다. $R_f$는 분리 가능하게 닫힌 분수 체를 갖는 정규 도메인이므로 $H^1_\etale(D(f'), \underline{M}) = 0$는 보조정리 \ref{etale-cohomology-lemma-gabber-for-h1-absolutely-algebraically-closed}로 표시된다. 따라서 보조정리 \ref{etale-cohomology-lemma-gabber-criterion}를 $Z = \Spec(R')$에 적용하여 결론을 내릴 수 있다. \end{proof}

\begin{lemma} \label{etale-cohomology-lemma-integral-cover-trivial-cohomology} $X$를 아핀 스킴으로 둡니다. \begin{enumerate} \item 정역 전사 사상 $X' \to X$가 존재하므로 모든 닫힌 하위 체계 $Z' \subset X'$, 모든 유한 아벨 군 $M$, 그리고 $q \geq 1$마다 $H^q_\etale(Z', \underline{M}) = 0$가 있다. \item 모든 닫힌 하위 체계 $Z \subset X$, 유한 아벨 군 $M$, $q \geq 1$ 및 $\xi \in H^q_\etale(Z, \underline{M})$에 대해 유한 전사 사상이 존재한다. $X' \to X$는 $\xi$가 $H^q_\etale(X' \times_X Z, \underline{M})$에서 0으로 돌아가도록 유한하게 표현된다. \end{enumerate} \end{lemma}

\begin{proof}
쓰다$X = \Spec(A)$. 쓰다$A = \mathbf{Z}[x_i]/J$어떤 사람들에게는아이디얼 $J$.
$R$를 대수학에서 $\mathbf{Z}[x_i]$의 정역 클로저로 설정한다.
$\mathbf{Z}[x_i]$의 분수 체를 닫습니다. 허락하다
$A' = R/JR$ 및 집합 $X' = \Spec(A')$. 이는 (1)와 같이 예를 제공한다.
보조정리 \ref{etale-cohomology-lemma-gabber-for-absolutely-algebraically-closed}.

\medskip\noindent
증명 (2). $X' \to X$를 우리가 찾은 정역 전사 사상라고 합시다.
위에. 확실히 $\xi$는 0으로 매핑된다.
$H^q_\etale(X' \times_X Z, \underline{M})$. $X'$를 다음과 같이 쓸 수 있다.
극한 $X' = \lim X'_i$ 스킴 유한 및 유한 표현
$X$ 이상; 이것은 현재 아핀 사례에서는 하기 쉽지만
보다 일반적인 극한, 보조정리의 특별한 경우이다.
\ref{limits-lemma-integral-limit-finite-and-finite-presentation}.
보조정리 \ref{etale-cohomology-lemma-directed-colimit-cohomology}에 의해
$\xi$가 $H^q_\etale(X'_i \times_X Z, \underline{M})$에서 0으로 매핑되는 것을 볼 수 있다.
일부 $i$의 경우 충분히 큽니다.
\end{proof}








\section{절대정수적으로 닫힌 소멸}
\label{etale-cohomology-section-aic-vanishing}

\noindent
환 $R$는 절대정수적으로 닫힌라고 말한 것을 기억하세요.
$R$ 이상의 모든 모닉 다항식이 $R$에 근을 갖는 경우
(대수학에 대한 자세한 내용, 정의
\ref{more-algebra-definition-absolutely-integrally-closed}).
이 절에서 우리는 의 에탈 코호몰로지를 증명한다.
유한 비틀림 그룹에서 $\Spec(R)$와 계수
양의 각도로 사라집니다. (명제 \ref{etale-cohomology-proposition-aic-vanishing})
이로써 이전 버전이 약간 개선되었다.
보조정리 \ref{etale-cohomology-lemma-gabber-for-absolutely-algebraically-closed}.
독자는 이 절을 건너뛰는 것이 좋다.

\begin{lemma}
\label{etale-cohomology-lemma-find-extension}
$A$를 환으로 설정한다. $a, b \in A$ 이렇게 하자
$aA + bA = A$ 및 $a \bmod bA$는 화합의 근원이다.
그러면 단일 유전자 확장 $A \subset B$이 존재한다.
그리고 요소$y \in B$그렇게$u = a - by$단위이다.
\end{lemma}

\begin{proof}
$a^n \equiv 1 \bmod bA$라고 말하세요. 특히 $a^i$는 모듈로 단위이다.
$b^mA$모두를 위해$i, m \geq 1$. 우리주장존재한다
$a_1, \ldots, a_n \in A$ 이렇게
$$
1 = a^n + a_1 a^{n - 1}b + a_2 a^{n - 2}b^2 + \ldots + a_n b^n
$$
즉, $1 - a^n \in bA$부터 $a_1 \in A$ 요소를 찾을 수 있다.
성질 단위를 사용하여 $1 - a^n - a_1 a^{n - 1} b \in b^2A$
$a^{n - 1}$ 모듈로 $bA$이다. 다음으로 $a_2 \in A$ 요소를 찾을 수 있다.
$1 - a^n - a_1 a^{n - 1} b - a_2 a^{n - 2} b^2 \in b^3A$와 같습니다.
등. 결국 우리는 $a_1, \ldots, a_{n - 1} \in A$를 발견했다.
$1 - (a^n + a_1 a^{n - 1}b + a_2 a^{n - 2}b^2 +
\ldots + a_{n - 1} ab^{n - 1}) \in b^nA$. 이를 통해 우리는 다음을 찾을 수 있다.
$a_n \in A$ 표시된 동등성이 유지된다.

\medskip\noindent 위와 같이 $a_1, \ldots, a_n$를 사용하면 주장 $$
B = A[y]/(y^n + a_1 y^{n - 1} + a_2 y^{n - 2} + \ldots + a_n)
$$ 설정이 작동한다. 즉, $\mathfrak q \subset B$가 $\mathfrak p \subset A$ 위에 있는 소 아이디얼라고 가정한다. 모순을 얻으려면 $u = a - by$가 $\mathfrak q$에 있다고 가정한다. $b \in \mathfrak p$인 경우 $a \not \in \mathfrak p$는 $aA + bA = A$이므로 $u$는 $\mathfrak q$에 없다. 따라서 $b \not \in \mathfrak p$, 즉 $b \not \in \mathfrak q$라고 가정할 수 있다. 이는 $y \bmod \mathfrak q$가 $a/b \bmod \mathfrak q$과 동일함을 의미한다. 그러나 그러면 $$
0 = y^n + a_1 y^{n - 1} + a_2 y^{n - 2} + \ldots + a_n =
b^{-n}(a^n + a_1 a^{n - 1}b + a_2 a^{n - 2}b^2 + \ldots + a_nb^n) =
b^{-n}
$$ 모순이 발생한다. 이것으로 증명이 완료된다. \end{proof}

\noindent
증명을 설명하려면 스킴 그룹을 소개해야 한다.
소수 $\ell$를 수정하세요. 허락하다
$$
A = \mathbf{Z}[\zeta] =
\mathbf{Z}[x]/(x^{\ell - 1} + x^{\ell - 2} + \ldots + 1)
$$
즉, $A$는 다음에 의해 생성된 $\mathbf{Z}$의 단일 유전자 확장이다.
기본 $\ell$번째 통일근 $\zeta$. 우리는 집합
$$
\pi = \zeta - 1
$$
계산(생략)은 $\ell$이 $\pi^{\ell - 1}$로 나누어진다는 것을 보여줍니다.
$A$에서. $A$에 대한 첫 번째 그룹 스킴은
$$
G = \Spec(A[s, \frac{1}{\pi s + 1}])
$$
공곱셈에 의해 주어진 그룹 법칙
$$
\mu :
A[s, \frac{1}{\pi s + 1}]
\longrightarrow
A[s, \frac{1}{\pi s + 1}] \otimes_A A[s, \frac{1}{\pi s + 1}],\quad
s \longmapsto \pi s \otimes s + s \otimes 1 + 1 \otimes s
$$
이 선택으로 우리는
$$
\mu(\pi s + 1) = (\pi s + 1) \otimes (\pi s + 1)
$$
따라서 우리는 실제로 표시된 대로 $A$ 대수 맵을 갖게 된다. 우리는
이것이 실제로 그룹법을 정의한다는 것을 확인한다. 우리의 두 번째 그룹
스킴 초과 $A$는
$$
H = \Spec(A[t, \frac{1}{\pi^\ell t + 1}])
$$
공곱셈에 의해 주어진 그룹 법칙
$$
\mu : A[t, \frac{1}{\pi^\ell t + 1}]
\longrightarrow
A[t, \frac{1}{\pi^\ell t + 1}] \otimes_A A[t, \frac{1}{\pi^\ell t + 1}],\quad
t \longmapsto \pi^\ell t \otimes t + t \otimes 1 + 1 \otimes t
$$
이전과 동일한 검증을 통해 이것이 그룹 법칙을 정의한다는 것을 알 수 있다.
다음으로, 우리는 다항식을 관찰합니다
$$
\Phi(s) = \frac{(\pi s + 1)^\ell - 1}{\pi^\ell}
$$
$A[s]$에 있고 $\ell$ 등급이고 $s$에 모닉이 있다. 즉, 계수
~의$s^i$~을 위한$0 < i < \ell$는 다음과 같다${\ell \choose i}\pi^{i - \ell}$
그리고 그 이후로$\pi^{\ell - 1}$나누다$\ell$~에$A$이것은의 요소입니다$A$.
우리는 환 지도를 얻습니다.
$$
A[t, \frac{1}{\pi^\ell t + 1}]
\longrightarrow
A[s, \frac{1}{\pi s + 1}],\quad
t \longmapsto \Phi(s)
$$
독자가 쉽게 확인할 수 있는 것은 공곱셈과 호환된다.
따라서 우리는 스킴 그룹의 사상을 얻습니다.
$$
f : G \to H
$$
특히 다음 보조정리는 이 사상이 충실히 있음을 보여줍니다.
flat(사실 우리는 그것이 유한한 에탈 전사라는 것을 알게 될 것입니다).

\begin{lemma} \label{etale-cohomology-lemma-monogenic-one} $$
A[s, \frac{1}{\pi s + 1}] =
\left(A[t, \frac{1}{\pi^\ell t + 1}]\right)[s]/(\Phi(s) - t)
$$ 특히, $G$의 Hopf 대수는 $H$의 Hopf 대수를 단일하게 확장한 것이다. \end{lemma}

\begin{proof} 위의 논의와 $\Phi(s)$의 모양을 따릅니다. 특히, $\Phi(s) = t$를 사용하면 $\frac{1}{\pi^\ell t + 1}$ 요소가 $\frac{1}{(\pi s + 1)^\ell}$ 요소가 된다. \end{proof}

\noindent
다음으로 $f$의 핵을 계산해 보자. $H$의 출처는
에 의해 주어진$t = 0$~에$H$우리는핵~의$f$주어진다
$\Phi(s) = 0$로. 이제 $A$ 값 포인트를 관찰하세요.
$\sigma_0, \ldots, \sigma_{\ell - 1}$
$G$에 의해 제공됨
$$
\sigma_i :
s = \frac{\zeta^i - 1}{\pi} = \frac{\zeta^i - 1}{\zeta - 1} =
\zeta^{i - 1} + \zeta^{i - 2} + \ldots + 1,\quad
i = 0, 1, \ldots, \ell - 1
$$
확실히 $\Ker(f)$에 포함되어 있다. 게다가 이것들은 모두 쌍을 이룬다.
{\bf 모든} 올 $G \to \Spec(A)$에서 구별된다. 또한 독자는
$\sigma_i +_G \sigma_j = \sigma_{i + j \bmod \ell}$을 계산한다.
따라서 스킴 그룹의 닫힌 몰입을 찾습니다.
$$
\underline{\mathbf{Z}/\ell \mathbf{Z}}_A \longrightarrow \Ker(f)
$$
$i$을 $\sigma_i$에게 보냅니다. 그러나 공사로 인해 $\Ker(f)$
유한한 플랫 오버$\Spec(A)$학위$\ell$. 그러므로 우리는 결론을 내린다
이 지도는 동형사상이다. 대체로 우리는 다음과 같이 결론을 내립니다.
짧은 완전열이 있다.
$A$
\begin{equation}
\label{etale-cohomology-equation-ses}
0 \to
\underline{\mathbf{Z}/\ell \mathbf{Z}}_A
\to G
\to H
\to 0
\end{equation}에 대한 그룹 스킴.

\begin{lemma} \label{etale-cohomology-lemma-lift-points-H-to-G} $R$를 $A$ 대수로 정의한다. 이는 절대정수적으로 닫힌이다. 그러면 $G(R) \to H(R)$는 전사이다. \end{lemma}

\begin{proof}
$h \in H(R)$를 $A$-대수 맵에 대응시키십시오.
$A[t, \frac{1}{\pi^\ell t + 1}] \to R$ $t$을 $a \in A$에게 보냅니다.
$\Phi(s)$는 모닉이므로 $b \in A$를 찾을 수 있다.
$\Phi(b) = a$로. 보조정리 \ref{etale-cohomology-lemma-monogenic-one}에 의해
$s$를 $b$로 전송하여 고유한 $A$ 대수 맵을 얻습니다.
$A[s, \frac{1}{\pi s + 1}] \to R$ 호환
위의 지도 $A[t, \frac{1}{\pi^\ell t + 1}] \to R$를 사용한다.
이는 차례로 $g \in G(R)$ 요소에 해당한다.
$h \in H(R)$에 매핑된다.
\end{proof}

\begin{lemma}
\label{etale-cohomology-lemma-interpolate}
$R$를 절대정수적으로 닫힌인 $A$ 대수라고 가정한다.
$I, J \subset R$를 $I + J = R$로 아이디얼로 설정한다.
$g \in G(R)$ 같은 것이 존재합니다
$g \bmod I = \sigma_0$ 및 $g \bmod J = \sigma_1$이다.
\end{lemma}

\begin{proof}
선택하다$x \in I$그렇게$x \equiv 1 \bmod J$.
$I$을 $xR$로, $J$를 $(x - 1)R$로 바꿀 수 있다.
그런 다음 우리는 다음과 같은 $s \in R$를 찾고 있다.
\begin{enumerate}
\item $1 + \pi s$는 단위이며,
\item $s \equiv 0 \bmod xR$ 및
\item $s \equiv 1 \bmod (x - 1)R$.
\end{enumerate}
마지막 두 조건은 다음과 같습니다.$s = x + x(x - 1)y$어떤 사람들에게는$y \in R$.
첫 번째 조건은 $1 + \pi s = 1 + \pi x + \pi x (x - 1) y$라고 말한다.
$R$ 단위여야 한다.
그러나 $1 + \pi x$ 및 $\pi x (x - 1)$는
단위아이디얼~의$R$그리고 그거$1 + \pi x$은$\ell$번째 루트$1$
모듈로 $\pi x (x - 1)$\footnote{$1 + \pi x$가 합동이므로
$1$ 모듈로 $\pi$, $1$ 모듈로 $x$와 합동, 그리고 다음과 합동
$1 + \pi = \zeta$ 모듈로 $x - 1$ 그리고 우리는
$(\pi) \cap (x) \cap (x - 1) = (\pi x (x - 1))$ 에서 $A[x]$.}.
따라서 우리는 보조정리 \ref{etale-cohomology-lemma-find-extension}로 승리하고
$R$는 절대정수적으로 닫힌라는 사실이다.
\end{proof}

\begin{proposition}
\label{etale-cohomology-proposition-aic-vanishing}
$R$를 절대정수적으로 닫힌 환으로 설정한다.
$M$를 유한 아벨 군이라고 가정한다.
그 다음에$H^i_\etale(\Spec(R), \underline{M}) = 0$~을 위한$i > 0$.
\end{proposition}

\begin{proof} 임의의 유한 아벨 군에는 하위 몫이 소수 순서로 순환하는 유한 필터링이 있으므로 $M = \mathbf{Z}/\ell\mathbf{Z}$라고 가정할 수 있다. 여기서 $\ell$는 소수이다.

\medskip\noindent $R$의 모든 국소환이 엄격하게 헨셀식임을 확인하세요. 대수학에 대한 추가 정보, 보조정리 \ref{more-algebra-lemma-absolutely-integrally-closed-strictly-henselian}를 참조하라. 또한, $R$의 모든 국소화는 대수학에 관한 추가 사항, 보조정리 \ref{more-algebra-lemma-absolutely-integrally-closed-quotient-localization}에 의해 절대정수적으로 닫힌이기도 한다. 따라서 보조정리 \ref{etale-cohomology-lemma-gabber-criterion}는 $$
H^1_\etale(D(f + g), \mathbf{Z}/\ell\mathbf{Z})
\longrightarrow
H^1_\etale(D(f(f + g)), \mathbf{Z}/\ell\mathbf{Z}) \oplus
H^1_\etale(D(g(f + g)), \mathbf{Z}/\ell\mathbf{Z})
$$의 핵이 모든 $f, g \in R$에 대해 0임을 보여주는 것으로 충분하다고 알려줍니다. $R$를 $R_{f + g}$로 바꾼 후 다음과 같이 주장으로 줄이다. 주어진 $\xi \in H^1_\etale(\Spec(R), \mathbf{Z}/\ell\mathbf{Z})$ 및 아핀 개방 피복 $\Spec(R) = U \cup V$로 $\xi|_U$ 및 $\xi|_V$는 다음과 같습니다. 사소한 다음에는 $\xi = 0$이다.

\medskip\noindent $A = \mathbf{Z}[\zeta]$를 위와 같이 둡니다. $\mathbf{Z} \subset A$는 단일 유전자이기 때문에 환 맵 $A \to R$을 찾을 수 있다. 이제부터 $R$를 $A$ 대수로 생각하고 $\Spec(R)$를 $\Spec(A)$에 대한 스킴으로 생각한다. 짧은 완전열 (\ref{etale-cohomology-equation-ses})를 $\Spec(R)$로 기본 변경하고 에탈 코호몰로지를 취하면 $$
G(R) \to H(R) \to
H^1_\etale(\Spec(R), \mathbf{Z}/\ell\mathbf{Z}) \to
H^1_\etale(\Spec(R), G)
$$를 얻습니다. 증명의 나머지 부분에서 이 점을 명심하라.

\medskip\noindent
$\tau \in \Gamma(U \cap V, \mathbf{Z}/\ell\mathbf{Z})$
Mayer-Vietoris 시퀀스의 경계를 갖는 절이어야 한다.
(보조정리 \ref{etale-cohomology-lemma-mayer-vietoris})는 $\xi$를 제공한다.
$i = 0, 1, \ldots, \ell - 1$의 경우 $A_i \subset U \cap V$를 허용한다.
개방형이고 닫힌 부분집합이다. 여기서 $\tau$에는 $i \bmod \ell$ 값이 있다.
따라서 우리는 유한 분리된 결합 분해를 가집니다.
$$
U \cap V = A_0 \amalg \ldots \amalg A_{\ell - 1}
$$
$\tau$는 각 $A_i$에서 일정한다. 을 위한
$i = 0, 1, \ldots, \ell - 1$
$\tau_i \in H^0(U \cap V, \mathbf{Z}/\ell\mathbf{Z})$을 나타냄
와 같은 요소$1$~에$A_i$그리고
같음$0$~에$A_j$~을 위한$j \not = i$.
그런 다음 $\tau$는 $\tau_i$\footnote{모듈로의 배수의 합이다.
$\tau = \sum i \tau_i$.} 계산 오류가 있다.
따라서 코호몰로지 클래스에 해당하는 것을 보여주는 것으로 충분한다.
$\tau_i$는 사소한 일이다. 이는 우리를 다음과 같은 경우로 축소시킵니다.
$\tau$는 두 개의 고유한 값, 즉 $1$ 및 $0$만 사용한다.

\medskip\noindent
$\tau$는 $1$ 및 $0$ 값만 사용한다고 가정한다. 쓰다
$$
U \cap V = A \amalg B
$$
여기서 $A$는 $\tau = 0$이고 $B$는 궤적이다.
여기서 $\tau = 1$. 그러면 $A$와 $B$는 서로 분리되어 있다. 닫힌 부분집합.
$\overline{A}$ 및 $\overline{B}$는 $A$ 및 $B$의 종료를 나타냅니다.
$\Spec(R)$. 그러면 ``바나나''가 있습니다: 즉
우리는
$$
\overline{A} \cap \overline{B} = Z_1 \amalg Z_2
$$
$Z_1 \subset U$ 및 $Z_2 \subset V$ 분리된 닫힌 부분집합을 사용한다.
집합 $T_1 = \Spec(R) \setminus V$ 및 $T_2 = \Spec(R) \setminus U$.
$Z_1 \subset T_1 \subset U$, $Z_2 \subset T_2 \subset V$ 및
$T_1 \cap T_2 = \emptyset$. 위상적으로 우리는 다음과 같이 쓸 수 있다.
$$
\Spec(R) = \overline{A} \cup \overline{B} \cup T_1 \cup T_2
$$
이를 시각화하기 위해 그림을 그리는 것이 좋다.
$\xi$이 0임을 증명하기 위해 $R$을 다음과 같이 바꿀 수 있다.
감소합니다(명제 \ref{etale-cohomology-proposition-topological-invariance}).
아래에서는 $A$, $\overline{A}$, $B$, $\overline{B}$, $T_1$, $T_2$를 생각한다.
$\Spec(R)$의 축소된 폐쇄형 하위 계획으로. 다음으로 $Z_1$의 스킴 구조체와 같이
그리고 $Z_2$ 우리는
$$
Z_1 = \overline{A} \cap (\overline{B} \cup T_1)
\quad\text{및}\quad
Z_2 = \overline{A} \cap (\overline{B} \cup T_2)
$$
(스킴 사상, 정의에서와 같은 이론적 합집합 및 교차점
\ref{morphisms-definition-scheme-theoretic-intersection-union}).

\medskip\noindent
표시하다$X$그만큼$G$-토르소 오버$\Spec(R)$이미지에 해당하는
~의$\xi$~에$H^1(\Spec(R), G)$. 만약에$X$사소한 일이라면$\xi$
$h \in H(R)$ 요소에서 옵니다(코호몰로지의 완전열 참조).
위에). 그러나 보조정리 \ref{etale-cohomology-lemma-lift-points-H-to-G}
$h$ 요소가 $G(R)$ 요소로 상승하고 결론을 내립니다.
$\xi = 0$ 원하는 대로. 따라서 우리의 목표는 $X$이 사소하다는 것을 증명하는 것이다.

\medskip\noindent
임베딩 $\mathbf{Z}/\ell \mathbf{Z} \to G(R)$
$i \bmod \ell$을 $\sigma_i \in G(R)$에게 보냅니다. $\overline{A}$를 확인하세요.
는 절대정수적으로 닫힌 환의 스펙트럼입니다(즉
$R$)의 몫. 에 의해
보조정리 \ref{etale-cohomology-lemma-interpolate} 우리는 $g \in G(\overline{A})$를 찾을 수 있습니다
$g|_{\overline{A} \cap Z_1} = \sigma_0$ 그리고
$g|_{\overline{A} \cap Z_2} = \sigma_1$ (스킴 이론적으로는).
그러면 우리는 정의할 수 있다.
\begin{enumerate}
\item $g_1 \in G(U)$이는$g$~에$\overline{A} \cap U$, 어느
~이다$\sigma_0$~에$\overline{B} \cap U$, 그리고$\sigma_0$~에$T_1$, 그리고
\item $g_2 \in G(V)$이는$g$~에$\overline{A} \cap V$, 어느
~이다$\sigma_1$~에$\overline{B} \cap V$, 그리고$\sigma_1$~에$T_2$.
\end{enumerate}
즉, 찾기 위해서는$g_1$에서와 같이 (1) 우리는절
$\sigma_0$ 의 $\Omega = (\overline{B} \cup T_1) \cap U$
의 제한에절 $g$~에$\Omega' = \overline{A} \cap U$.
$U = \Omega \cup \Omega'$ (스킴 이론적으로는)
왜냐하면 $U$는 감소하고 $\Omega \cap \Omega' = Z_1$ (이론적으로는 스킴)
$Z_1$을 선택한다. 따라서
사상, 보조정리 \ref{morphisms-lemma-scheme-theoretic-union}
$U$는 $\Omega$ 및 $\Omega'$의 푸시아웃이다.
$Z_1$을 따라. 따라서 $g_1$을 찾을 수 있다.
마찬가지로 존재에 대해서도$g_2$안에 (2).
그럼 우리는
$$
\tau = g_2|_{A \cup B} - g_1|_{A \cup B} \quad(\text{집단법칙에 추가})
$$
그리고 우리는 $X$가 사소하여 증명을 마무리한다는 것을 알 수 있다.
\end{proof}




















\section{고유 기저변환} \label{etale-cohomology-section-gabber-affine-proper}의 아핀 아날로그

\noindent 이 절에서는 Ofer Gabber의 결과를 논의한다. \cite{gabber-affine-proper}을 참조하라. 이는 Roland Huber에 의해서도 증명되었다. \cite{Huber-henselian}를 참조하라. 우리는 절 \ref{etale-cohomology-section-strictly-henselian-local-rings}에서 Gabber의 증명에 필요한 작업 중 일부를 이미 수행했다.

\begin{lemma}
\label{etale-cohomology-lemma-efface-cohomology-on-closed-by-finite-cover}
$X$를 아핀 스킴으로 설정한다. $\mathcal{F}$를 비틀림 아벨 층이라고 합시다.
$X_\etale$에 있다. $Z \subset X$를 폐쇄형 하위 구성이라고 가정한다. 허락하다
$\xi \in H^q_\etale(Z, \mathcal{F}|_Z)$ 일부 $q > 0$에 대해.
그러면 주입 맵 $\mathcal{F} \to \mathcal{F}'$이 존재한다.
$X_\etale$의 비틀림 아벨 층
이미지$\xi$~에$H^q_\etale(Z, \mathcal{F}'|_Z)$0이다.
\end{lemma}

\begin{proof}
기본정리 \ref{etale-cohomology-lemma-torsion-colimit-constructible} 및 \ref{etale-cohomology-lemma-colimit}
우리는 지도를 찾을 수 있어요$\mathcal{G} \to \mathcal{F}$~와 함께$\mathcal{G}$
건설 가능한아벨 층그리고$\xi$요소에서 오는$\zeta$~의
$H^q_\etale(Z, \mathcal{G}|_Z)$. 주입 지도를 찾을 수 있다고 가정하자.
$\mathcal{G} \to \mathcal{G}'$ 비틀림 아벨 층 $X_\etale$
그런 이미지를$\zeta$~에$H^q_\etale(Z, \mathcal{G}'|_Z)$0이다.
그런 다음 $\mathcal{F}'$을 푸시아웃으로 사용할 수 있다.
$$
\mathcal{F}' = \mathcal{G}' \amalg_{\mathcal{G}} \mathcal{F}
$$
따라서 보조정리의 결론을 얻는다. (실제로 $Z$로의 제한은 완전하므로
유한 극한 및 여극한과 가환하고, 왼쪽 수반이므로 임의의 여극한과도 가환한다.)
따라서 우리는 $\mathcal{F}$이 구성 가능하다고 가정할 수 있다.

\medskip\noindent
$\mathcal{F}$이 구성 가능하다고 가정한다. 에 의해
보조정리 \ref{etale-cohomology-lemma-constructible-maps-into-constant-general}
$\mathcal{F}$ 일 때 결과를 증명하는 것으로 충분한다.
$f_*\underline{M}$ 형식이다. 여기서 $M$는 유한 아벨 군이다.
그리고 $f : Y \to X$는 유한한 표현의 유한한 사상이다.
(층은 여전히 다음과 같이 구성 가능한다.
보조정리 \ref{etale-cohomology-lemma-finite-pushforward-constructible}
하지만 우리는 이것이 필요하지 않습니다).
$f_*$의 형성으로 인해 어떠한 베이스 변화에도 통근이 가능한다.
(보조정리 \ref{etale-cohomology-lemma-finite-pushforward-commutes-with-base-change})
$f_*\underline{M}$에서 $Z$로의 제한은 다음과 같습니다.
$\underline{M}$를 통해 직상하는 것과 같습니다.
$Y \times_X Z \to Z$. 르레이의 스펙트럼 열
(명제 \ref{etale-cohomology-proposition-leray})
그리고 고차 직상의 소멸
(명제 \ref{etale-cohomology-proposition-finite-higher-direct-image-zero}),
우리는 발견
$$
H^q_\etale(Z, f_*\underline{M}|_Z) = H^q_\etale(Y \times_X Z, \underline{M}).
$$
보조정리 \ref{etale-cohomology-lemma-integral-cover-trivial-cohomology}에 의해
유한 표현의 유한 전사 사상 $Y' \to Y$를 찾을 수 있습니다
$\xi$는 $H^q(Y' \times_X Z, \underline{M})$에서 0으로 매핑된다.
$f' : Y' \to X$ 구성 $Y' \to Y \to X$ 표시 주장
지도
$$
f_*\underline{M} \longrightarrow f'_*\underline{M}
$$
위에서 말한 내용으로 증명을 마무리하는 단사이다.
원하는 주입성을 보려면 줄기를 보면 된다. 즉,
$\overline{x} : \Spec(k) \to X$가 기하점인 경우
$$
(f_*\underline{M})_{\overline{x}} =
\bigoplus\nolimits_{f(\overline{y}) = \overline{x}} M
$$
명제 \ref{etale-cohomology-proposition-finite-higher-direct-image-zero} 기준
다른 층에 대해서도 마찬가지이다.
$Y' \to Y$는 전사적이고 유한하므로 우리는 다음을 알 수 있다.
기하점 리프팅 $\overline{x}$에 대한 유도 지도는 다음과 같습니다.
전사도 하고 결론을 내린다.
\end{proof}

\noindent 위의 보조정리는 Gabber 결과에서 상위 코호몰로지 군을 처리한다. 다음 보조정리는 대역 절단을 처리하는 데 사용된다.

\begin{lemma}
\label{etale-cohomology-lemma-gabber-h0}
$X$를 준간소화하고 준분리된 스킴라고 하자.
$i : Z \to X$를 닫힌 몰입으로 설정한다. 가정
\begin{enumerate}
\item 임의의 층 $\mathcal{F}$ 중 $X_{Zar}$ 위의 것에 대하여 사상
$\Gamma(X, \mathcal{F}) \to \Gamma(Z, i^{-1}\mathcal{F})$
전단적이며,
\item 유한 사상 $X' \to X$ 가정 (1)에 대해
$Z \times_X X' \to X'$의 경우.
\end{enumerate}
그렇다면 누구에게나층 $\mathcal{F}$~에$X_\etale$우리는
$\Gamma(X, \mathcal{F}) = \Gamma(Z, i^{-1}_{small}\mathcal{F})$.
\end{lemma}

\begin{proof}
$\mathcal{F}$를 $X_\etale$에서 층으로 설정한다. 정식이 있습니다
(베이스 변경) 지도
$$
i^{-1}(\mathcal{F}|_{X_{Zar}})
\longrightarrow
(i_{small}^{-1}\mathcal{F})|_{Z_{Zar}}
$$
$Z_{Zar}$의 층. 다음을 보면 이 지도가 단정적이라는 것을 알 수 있다.
줄기에 있다. $z \in Z$ 왼쪽의 줄기
은줄기~의$\mathcal{F}|_{X_{Zar}}$~에$z$. 그만큼줄기오른쪽에
손 쪽은 모든 초등학교 에탈 동네에 대한 여극한이다.
$(U, u) \to (X, z)$ $U \times_X Z \to Z$에는 절이 포함된다.
$z$ 동네. 에탈 사상이 오픈되면서,
$U \to X$열린 동네다$U_0$~의$z$~에$X$. 지도
$\mathcal{F}(U_0) \to \mathcal{F}(U)$는 층에 의해 단사적이다.
조건~을 위한$\mathcal{F}$에탈에 관해서피복 $U \to U_0$.
모든 $U$ 및 $U_0$에 대해 여극한을 취하여 줄기에 대한 주입성을 얻습니다.

\medskip\noindent
이것과 가정 (1)로부터 지도는
$\Gamma(X, \mathcal{F}) \to \Gamma(Z, i^{-1}_{small}\mathcal{F})$
주입식이다. 에 의해 (2) 같은 것은 모든 것에 해당됩니다$X'$유한한 이상$X$.

\medskip\noindent
전사성을 증명하기 전에 몇 가지 표기법을 소개하겠습니다. 을 위한
어떤 물건이라도$U$~의$X_\etale$그리고$s \in \mathcal{F}(U)$우리
$i^{-1}s \in
(i_{small}^{-1}\mathcal{F})(U \times_X Z)$을 나타냄
철수(보다 정확하게는 지도 아래의 $s$ 이미지이다.
사이트, 보조정리 \ref{sites-lemma-recover}의 지도와 결합됨
전층 당김에서 층 당김으로) 이 공사는
$X_\etale$의 사상 및 (유한)에 의한 당김과 호환 가능
사상 $X' \to X$ 명확하게 하는 것은 독자의 몫이다.

\medskip\noindent
$t \in \Gamma(Z, i^{-1}_{small}\mathcal{F})$로 놔두세요. 건설로
$i^{-1}_{small}\mathcal{F}$ 에탈이 존재합니다 피복
$\{V_j \to Z\}_{j \in J}$, 에탈 사상 $U_j \to X$, 절
$s_j \in \mathcal{F}(U_j)$ 및 사상 $V_j \to U_j$ $X$
$t|_{V_j}$는 $i^{-1}s_j$의 후퇴이다.
$V_j \to U_j \times_X Z$. $V_j \to U_j \times_X Z$는 에탈이므로,
이미지는 공개 하위 구성표이다. $U_j \times_X Z \subset U_j$는
폐쇄형인 경우 $U_j$를 개방형 하위 체계로 대체한 후 다음과 같이 가정할 수 있다.
$V_j \to U_j \times_X Z$는 전사이다. 그런 다음 $\{V_j \to U_j \times_X Z\}$
는 에탈 피복이고 우리는 $t|_{U_j \times_X Z} = i^{-1}s_j$라고 결론을 내립니다.
모든 비어 있지 않은 닫힌 하위 구성표 $T \subset X$가 $Z$를 충족하는지 확인하세요.
가정 (1)에 의해 층 $(T \to X)_*\underline{\mathbf{Z}}$에 적용됨
예의 경우. 따라서 $\coprod U_j \to X$는 전사임을 알 수 있다.
사상, 보조정리에 대한 추가 정보
\ref{more-morphisms-lemma-there-is-a-scheme-integral-over}
유한 전사 사상 $X' \to X$를 찾을 수 있다.
$X' \to X$ Zariski는 $\coprod U_j \to X$을 통해 국소적으로 인수분해한다.
$X' = \bigcup_{\alpha \in A} W_\alpha$가 열려 있는 피복라고 가정해 보세요.
그리고 $j : A \to J$는 지도이므로 사상이 존재한다.
$g_\alpha : W_\alpha \to U_{j(\alpha)}$~ 위에$X$. 표시하다
$\mathcal{F}'$ $\mathcal{F}$을 $X'_\etale$로 철회한다.
$s'_\alpha \in \mathcal{F}'(W_\alpha)$ 하락을 나타냄
$s_{j(\alpha)}$ 에 의해 $g_\alpha$.
집합 $Z' = X' \times_X Z$. $i' : Z' \to X'$ 표시
$i$의 기본 변경 사항이다. 그러면 $(i'_{small})^{-1}\mathcal{F}' =
(Z' \to Z)_{small}^{-1}i_{small}^{-1}\mathcal{F}$. 이 신분증을 통해
$t' \in
\Gamma(Z', (i')^{-1}_{small}\mathcal{F}')$
$t$을 $Z'$로 철회한다.
그러면 $t'|_{W_\alpha \times_X Z}$는 $(i')^{-1}s'_\alpha$와 같습니다.
건설. 우리는 $t'$가 다음의 절라고 결론을 내립니다.
서브 다발
$$
(i')^{-1}(\mathcal{F}'|_{X'_{Zar}})
\subset
(i'_{small})^{-1}\mathcal{F}')|_{Z'_{Zar}}
$$
가정 (2)에 의해 절 $t'$는 절에서 비롯된다.
$s' \in \Gamma(X', \mathcal{F}'|_{X'_{Zar}}) = \Gamma(X', \mathcal{F}')$.
두 번째 문단에서 증명된 주입성에 의해,
우리는 $s'$의 두 가지 후퇴가 다음과 같다고 결론을 내렸습니다.
$X' \times_X X'$ 동일합니다(결국 이는
$t'$에서 $Z' \times_Z Z'$)로의 철회. 그러므로 우리는 결론을 내린다
$s'$는 $\mathcal{F}$의 절에서 $X$에 의해 생성된다.
비고 \ref{etale-cohomology-remark-cohomological-descent-finite}.
\end{proof}

\begin{lemma}
\label{etale-cohomology-lemma-connected-topological}
$Z \subset X$를 위상 공간 $X$의 닫힌 부분집합으로 가정한다.
추정하다
\begin{enumerate}
\item $X$는 스펙트럼 공간이다.
(위상, 정의 \ref{topology-definition-spectral-space}) 및
\item~을 위한$x \in X$교차로$Z \cap \overline{\{x\}}$
연결입니다(특히 비어 있지 않음).
\end{enumerate}
$Z = Z_1 \amalg Z_2$와 $Z_i$가 $Z$에서 닫힌 경우,
그러면 분해 $X = X_1 \amalg X_2$가 존재한다.
$X_i$는 $X$ 및 $Z_i = Z \cap X_i$에서 닫혔습니다.
\end{lemma}

\begin{proof}
$Z_i$이 준컴팩트하다는 점에 유의하세요. 따라서 포인트의 집합
$W_i$ $Z_i$에 특화된 것은 건설 가능한 위상에서 닫힙니다.
위상, 보조정리 \ref{topology-lemma-make-spectral-space} 기준.
가정 (2)는 $X = W_1 \amalg W_2$을 의미한다.
$x \in \overline{W_1}$로 놔두세요. 위상, 보조정리에 의해
\ref{topology-lemma-constructible-stable-specialization-closed} 부분 (1)
존재한다특수화 $x_1 \leadsto x$~와 함께$x_1 \in W_1$.
따라서 $\overline{\{x\}} \subset \overline{\{x_1\}}$ 그리고 우리는
그 $x \in W_1$. 즉, $X_i = W_i$을 설정하면 해당 작업이 수행된다.
\end{proof}

\begin{lemma}
\label{etale-cohomology-lemma-h0-topological}
$Z \subset X$를 위상 공간 $X$의 닫힌 부분집합으로 가정한다.
추정하다
\begin{enumerate}
\item $X$는 스펙트럼 공간이다.
(위상, 정의 \ref{topology-definition-spectral-space}) 및
\item~을 위한$x \in X$교차로$Z \cap \overline{\{x\}}$
연결입니다(특히 비어 있지 않음).
\end{enumerate}
그렇다면 누구에게나층 $\mathcal{F}$~에$X$우리는
$\Gamma(X, \mathcal{F}) = \Gamma(Z, \mathcal{F}|_Z)$.
\end{lemma}

\begin{proof}
$x \leadsto x'$가 특수화 포인트인 경우
표준 맵 $\mathcal{F}_{x'} \to \mathcal{F}_x$ 호환
절은 $\mathcal{F}$의 열기 및 기능을 초과한다. 매 포인트부터
$X$의 $Z$ 지점을 전문적으로 다루므로 다음과 같습니다.
$\Gamma(X, \mathcal{F}) \to \Gamma(Z, \mathcal{F}|_Z)$는 단사형이다.
어려운 부분은 그것이 전사임을 보여주는 것이다.

\medskip\noindent $\mathcal{B}$는 $X$의 모든 준소형 개방의 집합임을 나타냅니다. $\mathcal{F}$를 여과 여극한 $\mathcal{F} = \colim \mathcal{F}_i$로 작성한다. 여기서 각 $\mathcal{F}_i$는 가군, 방정식(\ref{modules-equation-towards-constructible-sets})과 같습니다. 가군, 보조정리 \ref{modules-lemma-filtered-colimit-constructibles}를 참조하라. 그런 다음 $\mathcal{F}|_Z = \colim \mathcal{F}_i|_Z$는 $Z$에 대한 제한으로 왼쪽 수반 (범주, 보조정리 \ref{categories-lemma-adjoint-exact} 및 층, 보조정리이다. \ref{sheaves-lemma-f-map}). 층, 보조정리 \ref{sheaves-lemma-directed-colimits-sections}에 의해 함자 $\Gamma(X, -)$와 $\Gamma(Z, -)$는 여과 여극한으로 통근한다. 따라서 층 $\mathcal{F}$는 가군, 방정식(\ref{modules-equation-towards-constructible-sets})과 같다고 가정할 수 있다.

\medskip\noindent
$\mathcal{F} \subset \mathcal{G}$ 임베딩이 있다고 가정한다.
그럼 우리는
$$
\Gamma(X, \mathcal{F}) =
\Gamma(Z, \mathcal{F}|_Z) \cap \Gamma(X, \mathcal{G})
$$
$\Gamma(Z, \mathcal{G}|_Z)$에서 교차점이 발생하는 곳이다.
이는 증명의 첫 번째 비고에서 다음과 같습니다.
$\mathcal{G}$의 대역 절단이 $\mathcal{F}$에 있는지 확인하여
줄기에서 그리고 $X$의 모든 지점이 $Z$의 지점에 특화되어 있기 때문이다.

\medskip\noindent
가군, 보조정리 \ref{modules-lemma-constructible-in-constant}에 의해
주입이 있습니다 $\mathcal{F} \to \prod (Z_i \to X)_*\underline{S_i}$
여기서 곱은 유한하고 $Z_i \subset X$는 닫혀 있으며 $S_i$는 유한한다.
따라서 층에 대한 전사성을 증명하는 것으로 충분한다.
$(Z_i \to X)_*\underline{S_i}$. 그것을 관찰하십시오
$$
\Gamma(X, (Z_i \to X)_*\underline{S_i}) = \Gamma(Z_i, \underline{S_i})
\quad\text{및}\quad
\Gamma(X, (Z_i \to X)_*\underline{S_i}|_Z) =
\Gamma(Z \cap Z_i, \underline{S_i})
$$
또한 조건 (1) 및 (2)는 $Z_i$에 상속된다. 이건 분명해
(2)에 대해 다음과 같습니다.
위상, 보조정리 \ref{topology-lemma-spectral-sub} (1)에 대해. 그리하여 그것은
(유한) 상수층의 경우 보조정리를 증명하는 것으로 충분한다.
이 사건은 보조정리 \ref{etale-cohomology-lemma-connected-topological}의 재진술이다.
그러면 증명이 완료된다.
\end{proof}

\begin{example}
\label{etale-cohomology-example-quinard}
보조정리 \ref{etale-cohomology-lemma-h0-topological}는 $X$가 스펙트럼이 아닌 경우 거짓이다.
여기에 예가 있다. $Y$를 $T_1$ 위상 공간으로 두고,
$y \in Y$ 미개방 지점이다. $X = Y \amalg \{ x \}$를 부여하자
닫힌 집합이 $\emptyset$, $\{y\}$인 위상, 그리고 모두
$F \amalg \{ x \}$, 여기서 $F$는 $Y$의 닫힌 부분집합이다. 그 다음에
$Z = \{x, y\}$는 $X$의 닫힌 부분집합이며, 이는 가정 (2)를 충족한다.
보조정리 \ref{etale-cohomology-lemma-h0-topological}. 그러나 $X$는 연결이지만 $Z$는 그렇지 않습니다.
따라서 보조정리의 결론은 상수층에 대해 실패한다.
가치있는$\{0, 1\}$~에$X$.
\end{example}

\begin{lemma}
\label{etale-cohomology-lemma-h0-henselian-pair}
$(A, I)$를 헨셀 쌍이라고 가정한다. 집합 $X = \Spec(A)$ 및
$Z = \Spec(A/I)$. 누구에게나층 $\mathcal{F}$~에$X_\etale$
$\Gamma(X, \mathcal{F}) = \Gamma(Z, \mathcal{F}|_Z)$이 있다.
\end{lemma}

\begin{proof}
환의 스펙트럼은 스펙트럼 공간이라는 점을 기억하세요.
대수학, 보조정리 \ref{algebra-lemma-spec-spectral}. 에 의해
대수학에 대한 추가 정보, 보조정리
\ref{more-algebra-lemma-irreducible-henselian-pair-connected}
$\overline{\{x\}} \cap Z$는 모든 $x \in X$에 대해 연결이다.
보조정리 \ref{etale-cohomology-lemma-h0-topological}에 의해 우리는 다음 문장을 볼 수 있다.
$X_{Zar}$의 층에 대해 참이다. 임의의 유한 사상 $X' \to X$
$X' = \Spec(A')$ 및 $Z \times_X X' = \Spec(A'/IA')$이 있다.
$(A', IA')$ 헨셀 쌍으로, 대수학에 대한 추가 정보, 보조정리를 참조하라.
\ref{more-algebra-lemma-integral-over-henselian-pair}
$(X')_{Zar}$의 층에 대해 동일한 설명을 얻습니다.
따라서 보조정리 \ref{etale-cohomology-lemma-gabber-h0}를 적용하여 결론을 내릴 수 있다.
\end{proof}

\noindent 마지막으로 Gabber의 정리를 진술하고 증명할 수 있다.

\begin{theorem}[개버]
\label{etale-cohomology-theorem-gabber}
$(A, I)$를 헨셀 쌍이라고 가정한다. 집합 $X = \Spec(A)$ 및
$Z = \Spec(A/I)$. 어떤 비틀림에도아벨 층 $\mathcal{F}$~에$X_\etale$
$H^q_\etale(X, \mathcal{F}) = H^q_\etale(Z, \mathcal{F}|_Z)$이 있다.
\end{theorem}

\begin{proof}
결과는 다음과 같습니다.$q = 0$~에 의해보조정리 \ref{etale-cohomology-lemma-h0-henselian-pair}.
$q \geq 1$로 놔두세요. 결과가 모든 각도 $< q$로 표시되었다고 가정한다.
$\mathcal{F}$를 비틀림 아벨 층이라고 가정한다. 허락하다
$\mathcal{F} \to \mathcal{F}'$
비틀림의 주입 맵이 된다. 아벨 층 (나중에 선택됨)
여핵 $\mathcal{Q}$를 사용하여 짧은 완전열을 얻습니다.
$$
0 \to \mathcal{F} \to \mathcal{F}' \to \mathcal{Q} \to 0
$$
$X_\etale$에 비틀림 아벨 층이 있다. 이것은 길고 정확한 지도를 제공합니다
코호몰로지 $X$ 및 $Z$에 대한 시퀀스는 다음과 같습니다.
$$
\xymatrix{
H^{q - 1}_\etale(X, \mathcal{F}') \ar[d] \ar[r] &
H^{q - 1}_\etale(X, \mathcal{Q}) \ar[d] \ar[r] &
H^q_\etale(X, \mathcal{F}) \ar[d] \ar[r] &
H^q_\etale(X, \mathcal{F}') \ar[d] \\
H^{q - 1}_\etale(Z, \mathcal{F}'|_Z) \ar[r] &
H^{q - 1}_\etale(Z, \mathcal{Q}|_Z) \ar[r] &
H^q_\etale(Z, \mathcal{F}|_Z) \ar[r] &
H^q_\etale(Z, \mathcal{F}'|_Z)
}
$$
정확한 행과 함께 아벨 군의 가환 도식 사용
증명을 마치겠습니다.

\medskip\noindent
$\mathcal{F}$에 대한 주입성. $\xi$를 0이 아닌 요소로 설정한다.
$H^q_\etale(X, \mathcal{F})$. 에 의해
보조정리 \ref{etale-cohomology-lemma-efface-cohomology-on-closed-by-finite-cover} 적용됨
$Z = X$ (!) 우리는 $\mathcal{F} \subset \mathcal{F}'$를 찾을 수 있습니다
$\xi$는 오른쪽에 0으로 매핑된다. 그러면 $\xi$는 다음과 같은 이미지이다.
$H^{q - 1}_\etale(X, \mathcal{Q})$ 요소와 전단사성
에 대하여 $q - 1$는 $\xi$가 0으로 매핑되지 않음을 의미한다.
$H^q_\etale(Z, \mathcal{F}|_Z)$.

\medskip\noindent
$\mathcal{F}$에 대한 생존성. $\xi$를 다음 요소로 설정한다.
$H^q_\etale(Z, \mathcal{F}|_Z)$. 에 의해
보조정리 \ref{etale-cohomology-lemma-efface-cohomology-on-closed-by-finite-cover} 적용됨
$Z = Z$ 우리는 $\mathcal{F} \subset \mathcal{F}'$를 찾을 수 있다.
$\xi$는 오른쪽에 0으로 매핑된다. 그러면 $\xi$는 다음과 같은 이미지이다.
$H^{q - 1}_\etale(Z, \mathcal{Q}|_Z)$ 요소와 전단사성
에 대하여 $q - 1$는 $\xi$이 수직 지도의 이미지에 있음을 의미한다.
\end{proof}

\begin{lemma}
\label{etale-cohomology-lemma-vanishing-restriction-injective}
$X$는 다음과 같이 덮을 수 있는 아핀 대각선을 갖는 스킴라고 가정한다.
$n + 1$ 아핀이 열립니다. $Z \subset X$를 폐쇄형 하위 구성이라고 가정한다.
$\mathcal{A}$를 $X_\etale$에서 환의 꼬임층으로 설정한다.
$\mathcal{I}$를 $\mathcal{A}$-가군의 단사 층으로 설정한다.
$X_\etale$에 있다.
그 다음에$H^q_\etale(Z, \mathcal{I}|_Z) = 0$~을 위한$q > n$.
\end{lemma}

\begin{proof} 이를 $n$에 대한 귀납법으로 증명하겠습니다. $n = 0$인 경우 $X$는 유사한다. $X = \Spec(A)$ 및 $Z = \Spec(A/I)$라고 말한다. $A^h$를 에탈 $A$-대수 $B$의 여과 여극한으로 하여 $A/I \to B/IB$가 동형사상이 되도록 한다. 그런 다음 $(A^h, IA^h)$는 헨셀 쌍이고 $A/I = A^h/IA^h$이다. 대수학에 대한 추가 정보, 보조정리 \ref{more-algebra-lemma-henselization} 및 해당 증명을 참조하라. 집합 $X^h = \Spec(A^h)$. 정리 \ref{etale-cohomology-theorem-gabber}에 의해 $$
H^q_\etale(Z, \mathcal{I}|_Z) = H^q_\etale(X^h, \mathcal{I}|_{X^h})
$$ 정리 \ref{etale-cohomology-theorem-colimit}에 의해 $$
H^q_\etale(X^h, \mathcal{I}|_{X^h}) =
\colim_{A \to B} H^q_\etale(\Spec(B), \mathcal{I}|_{\Spec(B)})
$$가 있음을 알 수 있다. 여기서 여극한은 $A$-대수학 위에 있다. $B$ 위와 같습니다. 사상 $\Spec(B) \to \Spec(A)$는 에탈이므로 제한 $\mathcal{I}|_{\Spec(B)}$는 $\mathcal{A}|_{\Spec(B)}$-가군 (코호몰로지 위의 사이트의 분사형 층이다. 보조정리 \ref{sites-cohomology-lemma-cohomology-of-open}). 따라서 오른쪽의 코호몰로지 군은 0이고 이 경우 결과를 얻습니다.

\medskip\noindent
유도 단계. 유도 단계를 수행하기 위해 Mayer-Vietoris를 사용할 수 있다.
즉, $X = U \cup V$ 여기서 $U$는 $n$ 아핀의 합집합이라고 가정한다.
열리고 $V$는 유사한다. 그런 다음 $X$의 대각선이 유사하다는 것을 이용하여,
$U \cap V$는 $n$ 아핀 오픈의 합집합이라는 것을 알 수 있다. 메이어-비에토리스
완전열을 제공한다.
$$
H^{q - 1}_\etale(U \cap V \cap Z, \mathcal{I}|_Z) \to
H^q_\etale(Z, \mathcal{I}|_Z) \to
H^q_\etale(U \cap Z, \mathcal{I}|_Z) \oplus
H^q_\etale(V \cap Z, \mathcal{I}|_Z)
$$
그리고 귀납 가설에 의해 원하는 대로 $q > n$에 대해 소멸을 얻습니다.
\end{proof}





\section{코호몰로지 곡선의 꼬임층 } \label{etale-cohomology-section-vanishing-torsion}

\noindent 이 절의 목표는 곡선에서 꼬임층의 코호몰로지에 대한 기본 유한성과 소멸 결과를 증명하는 것이다. 정리 \ref{etale-cohomology-theorem-vanishing-affine-curves}을 참조하라. 절 \ref{etale-cohomology-section-vanishing-torsion-coefficients}에서는 뇌터 환에 대해 꼬임 가군의 구성가능층에 대해 논의한다.

\begin{situation} \label{etale-cohomology-situation-what-to-prove} 여기서 $k$는 대수적으로 닫힌 체이고, $X$는 차원 $\leq 1$의 분리된 유한 유형 스킴이다. $k$, 그리고 $\mathcal{F}$는 $X_\etale$의 비틀림 아벨 층이다. \end{situation}

\noindent 상황 \ref{etale-cohomology-situation-what-to-prove}에서 다음 명제들을 증명하고자 한다. \begin{enumerate} \item \label{etale-cohomology-item-vanishing} $H^q_\etale(X, \mathcal{F}) = 0$이 $q > 2$에 대하여 성립한다. \item \label{etale-cohomology-item-vanishing-affine} $H^q_\etale(X, \mathcal{F}) = 0$이 $q > 1$에 대하여 성립한다. 단, $X$는 아핀이다. \item \label{etale-cohomology-item-vanishing-p-p} $H^q_\etale(X, \mathcal{F}) = 0$이 $q > 1$에 대하여 성립한다. 단, $p = \text{char}(k) > 0$이고 $\mathcal{F}$는 $p$-멱 꼬임층이다. \item \label{etale-cohomology-item-finite-prime-to-p} $H^q_\etale(X, \mathcal{F})$는 유한하다. 단, $\mathcal{F}$는 구성 가능하고 그 꼬임 차수는 $\text{char}(k)$와 서로소이다. \item \label{etale-cohomology-item-finite-proper} $H^q_\etale(X, \mathcal{F})$는 유한하다. 단, $X$는 고유이고 $\mathcal{F}$는 구성 가능하다. \item \label{etale-cohomology-item-base-change-prime-to-p} $H^q_\etale(X, \mathcal{F}) \to
H^q_\etale(X_{k'}, \mathcal{F}|_{X_{k'}})$는 대수적으로 닫힌 체의 임의의 확대 $k'/k$에 대하여 동형사상이다. 단, $\mathcal{F}$는 $\text{char}(k)$와 서로소인 꼬임층이다. \item \label{etale-cohomology-item-base-change-proper} $H^q_\etale(X, \mathcal{F}) \to
H^q_\etale(X_{k'}, \mathcal{F}|_{X_{k'}})$는 대수적으로 닫힌 체의 임의의 확대 $k'/k$에 대하여 동형사상이다. 단, $X$는 고유이다. \item \label{etale-cohomology-item-surjective} $H^2_\etale(X, \mathcal{F}) \to H^2_\etale(U, \mathcal{F})$는 모든 열린부분 $U \subset X$에 대하여 전사이다. \end{enumerate} 상황 \ref{etale-cohomology-situation-what-to-prove}가 주어졌을 때, 그 상황에 적용되는 위 명제들이 모두 참이면 ``명제 (\ref{etale-cohomology-item-vanishing})--(\ref{etale-cohomology-item-surjective})가 성립한다''고 하겠다. 상수 계수 코호몰로지 계산의 다음 귀결로 증명을 시작한다.

\begin{lemma}
\label{etale-cohomology-lemma-constant-smooth-statements}
상황 \ref{etale-cohomology-situation-what-to-prove}
$X$가 매끄럽고 $\mathcal{F} = \underline{\mathbf{Z}/\ell\mathbf{Z}}$라고 가정한다.
일부 소수 $\ell$에 대해. 그런 다음 진술
(\ref{etale-cohomology-item-vanishing}) -- (\ref{etale-cohomology-item-surjective}) 보류
$\mathcal{F}$의 경우.
\end{lemma}

\begin{proof} $X$는 매끄럽기 때문에 $X$는 매끄러운 곡선과 점의 유한하고 분리된 결합임을 알 수 있다. 점($k$의 스펙트럼)에 있는 임의의 층 중 더 높은 에탈 코호몰로지는 예의 경우 보조정리 \ref{etale-cohomology-lemma-compare-cohomology-point} 또는 \ref{etale-cohomology-lemma-affine-only-closed-points}에 의해 사소한다. 따라서 $X$는 매끄러운 곡선이라고 가정할 수 있다.

\medskip\noindent
사례 I: $\ell$는 $k$의 표수와 다릅니다.
이 사건은 다음과 같습니다.
보조정리 \ref{etale-cohomology-lemma-cohomology-smooth-projective-curve}
(투영 케이스) 및
보조정리 \ref{etale-cohomology-lemma-vanishing-cohomology-mu-smooth-curve}
(아핀 케이스). 명령문 (\ref{etale-cohomology-item-base-change-prime-to-p})
코호몰로지 및 대수적으로 닫힌 땅의 확장
체 속은 $g$이고 수는 다음과 같습니다.
``구멍'' $r$은 $k$에서 $k'$로 전달될 때 변경되지 않습니다.
문(\ref{etale-cohomology-item-surjective})은 다음과 같습니다. $H^2_\etale(U, \mathcal{F})$
$U \not = X$가 되자마자 0이 된다. 왜냐하면 $U$는 유사하기 때문이다.
(다양체, 기본형 \ref{varieties-lemma-proper-minus-point} 및
\ref{varieties-lemma-curve-affine-projective}).

\medskip\noindent 사례 II: $\ell$는 $k$의 표수와 같습니다. 소멸 에 의해 보조정리 \ref{etale-cohomology-lemma-vanishing-variety-char-p-p}. 문 (\ref{etale-cohomology-item-finite-proper}) 및 (\ref{etale-cohomology-item-base-change-proper})는 보조정리 \ref{etale-cohomology-lemma-finiteness-proper-variety-char-p-p}에서 따릅니다. \end{proof}

\begin{remark}[대수적으로 닫힌 땅의 확장에 따른 불변성 체]
\label{etale-cohomology-remark-invariance}
$k$를 표수 $p > 0$의 대수적으로 닫힌 체로 설정한다.
절 \ref{etale-cohomology-section-artin-schreier}에서 우리는
완전열
$$
k[x] \to k[x] \to
H^1_\etale(\mathbf{A}^1_k, \mathbf{Z}/p\mathbf{Z}) \to 0
$$
여기서 첫 번째 화살표는 $f(x)$를 $f^p - f$에 매핑한다. 대표자 집합
여핵은 다항식으로 구성된다.
$$
\sum\nolimits_{p \not | n} \lambda_n x^n
$$
$\lambda_n \in k$로. ($k$가 대수적으로 닫혀 있지 않은 경우
여기에도 몇 가지 상수를 추가해야 한다.) 특히
$k'/k$가 대수적으로 닫힌 확장인 경우
지도
$$
H^1_\etale(\mathbf{A}^1_k, \mathbf{Z}/p\mathbf{Z})
\to
H^1_\etale(\mathbf{A}^1_{k'}, \mathbf{Z}/p\mathbf{Z})
$$
일반적으로 동형사상이 아닙니다. 특히, 지도
$\pi_1(\mathbf{A}^1_{k'}) \to \pi_1(\mathbf{A}^1_k)$
에탈 기본 그룹 간(여기에 향후 참조 삽입)
동형사상도 아닙니다. 따라서 에탈 호모토피 유형
아핀선의 는 대수적으로 닫힌 근거 체에 따라 달라집니다.
위의 보조정리 \ref{etale-cohomology-lemma-constant-smooth-statements}에서 다음을 볼 수 있다.
표수 $p$에서만 일어나는 현상이다.
$p$-전력 꼬임 계수를 사용한다.
\end{remark}

\begin{lemma}
\label{etale-cohomology-lemma-ses-statements}
$k$는 대수적으로 닫힌 체라고 가정한다. $X$를 분리된 유한이라고 하자
유형스킴~ 위에$k$~의차원 $\leq 1$. 허락하다
$0 \to \mathcal{F}_1 \to \mathcal{F} \to \mathcal{F}_2 \to 0$
$X$에 비틀림 아벨 층이 있는 짧은 완전열이다.
만약 문 (\ref{etale-cohomology-item-vanishing}) -- (\ref{etale-cohomology-item-surjective}) 유지
$\mathcal{F}_1$ 및 $\mathcal{F}_2$에 대해 유지된다.
$\mathcal{F}$의 경우.
\end{lemma}

\begin{proof}
이는 대부분 정의와 긴 완전열에서 즉각적으로 나타납니다.
코호몰로지. 또한 $\mathcal{F}$ 구성 가능하다는 점을 확인하세요.
(관련\ 의 표수의 $k$) 필요충분조건은
$\mathcal{F}_1$ 및 $\mathcal{F}_2$ 모두 구성 가능한다.
(관련.\ 토션 프라임의표수~의$k$). 보다
명제 \ref{etale-cohomology-proposition-constructible-over-noetherian}.
일부 세부 사항은 생략되었다.
\end{proof}

\begin{lemma}
\label{etale-cohomology-lemma-finite-pushforward-statements}
$k$는 대수적으로 닫힌 체라고 가정한다. $f : X \to Y$를
유한 사상 분리된 유한 유형 스킴 이상의 $k$
차원 $\leq 1$. $\mathcal{F}$를 $X$의 비틀림 아벨 층으로 설정한다.
만약 문 (\ref{etale-cohomology-item-vanishing}) -- (\ref{etale-cohomology-item-surjective}) 유지
$\mathcal{F}$의 경우 $f_*\mathcal{F}$ 동안 유지된다.
\end{lemma}

\begin{proof}
즉, $H^q_\etale(X, \mathcal{F}) = H^q_\etale(Y, f_*\mathcal{F})$
에 의해소멸~의$R^qf_*$~을 위한$q > 0$
(명제 \ref{etale-cohomology-proposition-finite-higher-direct-image-zero}) 및
르레이 스펙트럼 열
(코호몰로지 사이트, 보조정리 \ref{sites-cohomology-lemma-apply-Leray}).
(\ref{etale-cohomology-item-surjective})의 경우 $f_*$ 통근 형식을 사용한다.
임의의 기본 변경
(보조정리 \ref{etale-cohomology-lemma-finite-pushforward-commutes-with-base-change}).
\end{proof}

\begin{lemma} \label{etale-cohomology-lemma-restrict-to-open} 상황 \ref{etale-cohomology-situation-what-to-prove}에서는 $\mathcal{F}$이 구성 가능하다고 가정한다. $j : X' \to X$를 조밀 개방형 하위 구성표의 포함으로 간주한다. 그런 다음 명령문 (\ref{etale-cohomology-item-vanishing}) -- (\ref{etale-cohomology-item-surjective})는 $\mathcal{F}$ 필요충분조건은에 대해 유지되며 $j_!j^{-1}\mathcal{F}$에 대해 유지된다. \end{lemma}

\begin{proof}
$X'$는 조밀이므로 $Z = X \setminus X'$에는 차원 $0$가 있음을 알 수 있다.
그러므로 는 유한하다집합 $Z = \{x_1, \ldots, x_n\}$~의$k$-합리적인 점.
짧은 완전열을 고려해보세요.
$$
0 \to j_!j^{-1}\mathcal{F} \to \mathcal{F} \to i_*i^{-1}\mathcal{F} \to 0
$$
보조정리 \ref{etale-cohomology-lemma-ses-associated-to-open}. 그것을 관찰하십시오
$H^q_\etale(X, i_*i^{-1}\mathcal{F}) =
H^q_\etale(Z, i^*\mathcal{F})$.
즉, $i : Z \to X$는 닫힌 몰입이므로 유한한다.
그러므로 우리는소멸~의$R^qi_*$~을 위한$q > 0$~에 의해
명제 \ref{etale-cohomology-proposition-finite-higher-direct-image-zero} 및
따라서 평등은 Leray 스펙트럼 열에서 따릅니다.
(코호몰로지 사이트, 보조정리 \ref{sites-cohomology-lemma-apply-Leray}).
$Z$은 대수적으로 닫힌 스펙트럼의 분리된 결합이므로
체, 우리는 $H^q_\etale(Z, i^*\mathcal{F}) = 0$라고 결론을 내립니다.
$q > 0$ 및
$$
H^0_\etale(Z, i^{-1}\mathcal{F}) =
\bigoplus\nolimits_{i = 1, \ldots, n} \mathcal{F}_{x_i}
$$
이는 $\mathcal{F}_{x_i}$로 인해 유한한다.
가정 $\mathcal{F}$은 구성 가능한다.
길고 정확한 코호몰로지 시퀀스는 완전열을 제공한다.
$$
0 \to
H^0_\etale(X, j_!j^{-1}\mathcal{F}) \to
H^0_\etale(X, \mathcal{F}) \to
H^0_\etale(Z, i^{-1}\mathcal{F}) \to
H^1_\etale(X, j_!j^{-1}\mathcal{F}) \to
H^1_\etale(X, \mathcal{F}) \to 0
$$
및 동형사상 $H^q_\etale(X, j_!j^{-1}\mathcal{F}) \to
H^q_\etale(X, \mathcal{F})$ 에 대하여 $q > 1$.

\medskip\noindent 이 시점에서 각각의 (\ref{etale-cohomology-item-vanishing}) - (\ref{etale-cohomology-item-surjective})는 $\mathcal{F}$에 대해 유지되고 필요충분조건은은 $j_!j^{-1}\mathcal{F}$에 대해 유지된다는 것을 쉽게 추론할 수 있다. 독자를 돕기 위해 몇 가지 작은 설명을 하겠습니다. () $\mathcal{F}$가 $k$의 표수에 대한 비틀림 소수라면 $j_!j^{-1}\mathcal{F}$도 마찬가지이다. (b) 층 $j_!j^{-1}\mathcal{F}$는 구성 가능한다. (c) $H^0_\etale(Z, i^{-1}\mathcal{F}) =
H^0_\etale(Z_{k'}, i^{-1}\mathcal{F}|_{Z_{k'}})$, 그리고 (d) $U \subset X$이 개방형이면 $U' = U \cap X'$는 $U$에서 조밀이다. \end{proof}

\begin{lemma} \label{etale-cohomology-lemma-even-easier} 상황 \ref{etale-cohomology-situation-what-to-prove}에서는 $X$가 매끄럽다고 가정한다. $j : U \to X$를 열린 몰입으로 둡니다. $\ell$를 소수라고 하자. $\mathcal{F} = j_!\underline{\mathbf{Z}/\ell\mathbf{Z}}$로 놔두세요. 그런 다음 명령문 (\ref{etale-cohomology-item-vanishing}) - (\ref{etale-cohomology-item-surjective})는 $\mathcal{F}$에 대해 유지된다. \end{lemma}

\begin{proof} $X$는 매끄러운 곡선의 분리된 결합이므로 $X$는 곡선(예: 기약)이라고 가정할 수 있다. 그러면 $U = \emptyset$이고 증명할 것이 없거나 $U \subset X$가 조밀이다. 이 경우 보조정리는 기본정리 \ref{etale-cohomology-lemma-constant-smooth-statements} 및 \ref{etale-cohomology-lemma-restrict-to-open}에서 따릅니다. \end{proof}

\begin{lemma} \label{etale-cohomology-lemma-somewhat-easier} 상황 \ref{etale-cohomology-situation-what-to-prove}에서는 $X$가 감소했다고 가정한다. $j : U \to X$를 열린 몰입으로 둡니다. $\ell$를 소수, $\mathcal{F} = j_!\underline{\mathbf{Z}/\ell\mathbf{Z}}$로 설정한다. 그런 다음 명령문 (\ref{etale-cohomology-item-vanishing}) - (\ref{etale-cohomology-item-surjective})는 $\mathcal{F}$에 대해 유지된다. \end{lemma}

\begin{proof}
보조정리 \ref{etale-cohomology-lemma-even-easier}와의 차이점은 여기서는 그렇지 않다는 것이다.
$X$가 매끄럽다고 가정한다. $\nu : X^\nu \to X$를 정규화라고 하자.
사상. 그러면 $\nu$는 유한한다.
(다양체, 보조정리 \ref{varieties-lemma-normalization-locally-algebraic})
그리고 $X^\nu$는 부드럽습니다.
(다양체, 보조정리 \ref{varieties-lemma-regular-point-on-curve}).
$j^\nu : U^\nu \to X^\nu$를 $U$의 역상으로 설정한다.
보조정리 \ref{etale-cohomology-lemma-even-easier}에 의해 결과는 다음과 같이 유지된다.
$j^\nu_!\underline{\mathbf{Z}/\ell\mathbf{Z}}$.
보조정리 \ref{etale-cohomology-lemma-finite-pushforward-statements}에 의해
결과는 $\nu_*j^\nu_!\underline{\mathbf{Z}/\ell\mathbf{Z}}$에 적용된다.
일반적으로 그것은 사실이 아닐 것이다.
$\nu_*j^\nu_!\underline{\mathbf{Z}/\ell\mathbf{Z}}$는 다음과 같습니다.
$j_!\underline{\mathbf{Z}/\ell\mathbf{Z}}$
하지만 우리는 다음과 같이 이 문제를 해결할 수 있다.
$X$가 감소함에 따라 사상 $\nu : X^\nu \to X$
조밀 열린 $j' : X' \to X$에 대한 동형사상이다.
(다양체, 보조정리 \ref{varieties-lemma-normalization-locally-algebraic}).
이 공개에 대해 우리는 합의했다
$$
(j')^{-1}(\nu_*j^\nu_!\underline{\mathbf{Z}/\ell\mathbf{Z}})
=
(j')^{-1}(j_!\underline{\mathbf{Z}/\ell\mathbf{Z}})
$$
보조정리 \ref{etale-cohomology-lemma-restrict-to-open} 사용
$j' : X' \to X$ 및 위의 층에 대해 두 번
우리는 결론을 내린다.
\end{proof}

\begin{lemma}
\label{etale-cohomology-lemma-vanishing-easier}
상황 \ref{etale-cohomology-situation-what-to-prove}에서는 $X$이 감소했다고 가정한다.
$j : U \to X$를 열린 몰입으로 $U$ 연결로 설정한다. 허락하다
$\ell$는 소수여야 한다. $\mathcal{G}$를 로컬에서 유한하게 놔두세요
상수층~의$\mathbf{F}_\ell$-벡터공간~에$U$. 허락하다
$\mathcal{F} = j_!\mathcal{G}$. 그런 다음 진술
(\ref{etale-cohomology-item-vanishing}) -- (\ref{etale-cohomology-item-surjective}) $\mathcal{F}$ 동안 유지된다.
\end{lemma}

\begin{proof}
$f : V \to U$를 $\ell$에 대한 소수의 유한 에탈 사상라고 하자
보조정리 \ref{etale-cohomology-lemma-pullback-filtered}에서와 같습니다. 의 토론
절 \ref{etale-cohomology-section-trace-method}는 지도를 제공한다.
$$
\mathcal{G} \to f_*f^{-1}\mathcal{G} \to \mathcal{G}
$$
그 구성은 동형사상이다. 그러므로 증명하는 것으로 충분하다.
보조정리와 $\mathcal{F} = j_!f_*f^{-1}\mathcal{G}$.
자리스키의 메인으로 정리
(사상, 보조정리에 대한 추가 정보
\ref{more-morphisms-lemma-quasi-finite-separated-pass-through-finite})
우리는 다이어그램을 선택할 수 있습니다
$$
\xymatrix{
V \ar[r]_{j'} \ar[d]_f & Y \ar[d]^{\overline{f}} \\
U \ar[r]^j & X
}
$$
$\overline{f} : Y \to X$ 유한 및 $j'$ 및 열린 몰입 사용
조밀 이미지로. $Y$을 축소하여 대체할 수 있습니다(이것은
$V$을 변경하지 마십시오. $V$이 $U$)에 비해 에탈로 줄어들기 때문이다.
$f$는 유한하고 $V$ 조밀은 $Y$이므로 $V = U \times_X Y$이 있다. 에 의해
보조정리 \ref{etale-cohomology-lemma-compatible-shriek-push-finite} 우리는
$$
j_!f_*f^{-1}\mathcal{G} = \overline{f}_*j'_!f^{-1}\mathcal{G}
$$
보조정리 \ref{etale-cohomology-lemma-finite-pushforward-statements}로 충분한다.
$j'_!f^{-1}\mathcal{G}$을 고려해보세요.
다음과 같이 주어진 여과의 존재
보조정리 \ref{etale-cohomology-lemma-pullback-filtered},
$j'_!$이 정확하다는 사실
보조정리 \ref{etale-cohomology-lemma-ses-statements}
우리를 사건으로 축소
$\mathcal{F} = j'_!\underline{\mathbf{Z}/\ell\mathbf{Z}}$
이는 보조정리 \ref{etale-cohomology-lemma-somewhat-easier}이다.
\end{proof}


%10.20.09

\begin{theorem}
\label{etale-cohomology-theorem-vanishing-affine-curves}
$k$가 대수적으로 닫힌 체인 경우 $X$는 분리된 유한 유형이다.
스킴의 차원 $\leq 1$는 $k$ 위에 있고 $\mathcal{F}$는 비틀림이다.
$X_\etale$에 아벨 층, 그런 다음
\begin{enumerate}
\item
$H^q_\etale(X, \mathcal{F}) = 0$는 $q > 2$에 대해,
\item
$H^q_\etale(X, \mathcal{F}) = 0$ 에 대하여 $q > 1$ 만약 $X$가 유사하다면,
\item
$H^q_\etale(X, \mathcal{F}) = 0$ 에 대하여 $q > 1$ 만약 $p = \text{char}(k) > 0$
그리고 $\mathcal{F}$는 $p$-파워 비틀림이다.
\item
$H^q_\etale(X, \mathcal{F})$는 $\mathcal{F}$인 경우 유한한다.
건설 가능하고 비틀림 프라임 $\text{char}(k)$,
\item
$H^q_\etale(X, \mathcal{F})$는 $X$가 적절하고
$\mathcal{F}$ 건설 가능,
\item
$H^q_\etale(X, \mathcal{F}) \to
H^q_\etale(X_{k'}, \mathcal{F}|_{X_{k'}})$는 동형사상이다.
대수적으로 닫힌 체의 확장 $k'/k$에 대해
$\mathcal{F}$가 $\text{char}(k)$에 대한 비틀림 프라임인 경우,
\item
$H^q_\etale(X, \mathcal{F}) \to
H^q_\etale(X_{k'}, \mathcal{F}|_{X_{k'}})$는 동형사상이다.
대수적으로 닫힌 체의 확장 $k'/k$에 대해
$X$가 적절하다면,
\item
$H^2_\etale(X, \mathcal{F}) \to H^2_\etale(U, \mathcal{F})$
모든 $U \subset X$ 열려 있음을 나타냅니다.
\end{enumerate}
\end{theorem}

\begin{proof} 정리는 상황 \ref{etale-cohomology-situation-what-to-prove} 문 (\ref{etale-cohomology-item-vanishing}) -- (\ref{etale-cohomology-item-surjective})이 유지됨을 나타냅니다. 첫 번째 단계는 $X$를 축소하여 바꾸는 것이다. 이는 명제 \ref{etale-cohomology-proposition-topological-invariance}에서 허용된다. 보조정리 \ref{etale-cohomology-lemma-torsion-colimit-constructible}에 의해 $\mathcal{F}$를 구성 가능한 아벨 층의 여과 여극한으로 쓸 수 있다. 코호몰로지를 타고 여극한으로 통근한다. 보조정리 \ref{etale-cohomology-lemma-colimit}를 참조하라. 또한, $X_{k'} \to X$를 통한 당김은 여극한을 왼쪽 인접으로 통근한다. 따라서 구성가능층에 대한 진술을 증명하는 것으로 충분한다.

\medskip\noindent
이 단락에서는 더 이상 추가하지 않고 보조정리 \ref{etale-cohomology-lemma-ses-statements}를 사용한다.
언급하다. 글쓰기
$\mathcal{F} = \mathcal{F}_1 \oplus \ldots \oplus \mathcal{F}_r$
여기서 $\mathcal{F}_i$는 $\ell_i$일부 소수 $\ell_i$에 대해 기본이다.
가정하다$\ell^n$죽인다$\mathcal{F}$어떤 소수에게는$\ell$. 이제 고려해보세요
완전열
$$
0 \to \mathcal{F}[\ell] \to \mathcal{F} \to \mathcal{F}/\mathcal{F}[\ell] \to 0.
$$
따라서 우리는 $\mathcal{F}$가 $\ell$-비틀림이라고 가정하는 것으로 충분하다는 것을 알 수 있다.
이는 $\mathcal{F}$가 다음의 구성가능층임을 의미한다.
$\mathbf{F}_\ell$-벡터공간 일부 소수 $\ell$에 대해.

\medskip\noindent 정의는 조밀 열린 $U \subset X$가 있음을 의미하며 $\mathcal{F}|_U$는 유한한 국소상수층이다. $\mathbf{F}_\ell$-벡터공간. $\dim(X) \leq 1$ $U$를 축소한 후 $U = U_1 \amalg \ldots \amalg U_n$는 기약 스킴의 분리된 합집합이라고 가정할 수 있으므로 ($\geq 2$ 구성 요소의 교차점에 있는 닫힌 점만 제거하면 된다. $U$). 보조정리 \ref{etale-cohomology-lemma-restrict-to-open}에 의해 $\mathcal{F} = j_!\mathcal{G}$의 경우로 축소한다. 여기서 $\mathcal{G}$는 $\mathbf{F}_\ell$-벡터공간의 $U$의 유한한 국소상수층이다.

\medskip\noindent
우리가 선택한 이후로$U = U_1 \amalg \ldots \amalg U_n$~와 함께$U_i$ 기약
우리는
$$
j_!\mathcal{G} =
j_{1!}(\mathcal{G}|_{U_1}) \oplus \ldots \oplus
j_{n!}(\mathcal{G}|_{U_n})
$$
여기서 $j_i : U_i \to X$는 사상을 포함한다.
$j_{i!}(\mathcal{G}|_{U_i})$의 경우는 다음에서 처리된다.
보조정리 \ref{etale-cohomology-lemma-vanishing-easier}.
\end{proof}

\begin{theorem} \label{etale-cohomology-theorem-vanishing-curves} $X$를 유한 유형으로 두고, 차원 $1$ 스킴은 대수적으로 닫힌 체 $k$에 대해 설명한다. $\mathcal{F}$를 $X_\etale$에서 꼬임층으로 설정한다. 그런 다음 $$
H_\etale^q(X, \mathcal{F}) = 0, \quad \forall q \geq 3.
$$ 만약 $X$ 아핀이면 또한 $H_\etale^2(X, \mathcal{F}) = 0$이다. \end{theorem}

\begin{proof} $X$이 분리된 경우 이는 보다 정확한 정리 \ref{etale-cohomology-theorem-vanishing-affine-curves}에서 바로 이어집니다. $X$가 분리되지 않은 경우 아핀 개방형 피복 $X = X_1 \cup \ldots \cup X_n$를 선택한다. $n$에 대한 유도를 통해 소멸이 $U = X_1 \cup \ldots \cup X_{n - 1}$에 대해 유지된다고 가정할 수 있다. 그런 다음 Mayer-Vietoris (보조정리 \ref{etale-cohomology-lemma-mayer-vietoris})는 $$
H^2_\etale(U, \mathcal{F}) \oplus H^2_\etale(X_n, \mathcal{F}) \to
H^2_\etale(U \cap X_n, \mathcal{F}) \to
H^3_\etale(X, \mathcal{F}) \to 0
$$를 제공한다. 그러나 $U \cap X_n$는 아핀 스킴의 개방형이므로 차원 가정과 관련이 있으므로 그룹은 $H^2_\etale(U \cap X_n, \mathcal{F})$는 정리 \ref{etale-cohomology-theorem-vanishing-affine-curves}에 의해 사라집니다. \end{proof}

\begin{lemma}
\label{etale-cohomology-lemma-base-change-dim-1-separably-closed}
$k'/k$는 분리 가능하게 닫힌 체의 확장이라고 가정한다.
허락하다$X$적절한 사람이 되다스킴~ 위에$k$~의차원 $\leq 1$.
$\mathcal{F}$를 $X$의 비틀림 아벨 층이라고 가정한다.
그런 다음 지도 $H^q_\etale(X, \mathcal{F}) \to
H^q_\etale(X_{k'}, \mathcal{F}|_{X_{k'}})$는 동형사상이다.
$q \geq 0$의 경우.
\end{lemma}

\begin{proof}
우리는 대수적으로 닫힌 체에 대해 이것을 보았습니다.
정리 \ref{etale-cohomology-theorem-vanishing-affine-curves}.
보조정리 명령문에서와 같이 $k \subset k'$이 주어지면 다음을 수행할 수 있다.
다이어그램을 선택하세요
$$
\xymatrix{
k' \ar[r] & \overline{k}' \\
k \ar[u] \ar[r] & \overline{k} \ar[u]
}
$$
여기서 $k \subset \overline{k}$ 및 $k' \subset \overline{k}'$는
대수적 종결. $k$와 $k'$는 분리 가능하게 닫혀 있으므로
체 확장
$\overline{k}/k$ 및 $\overline{k}'/k'$
대수적이고 완전히 분리될 수 없다. 이 경우 사상
$X_{\overline{k}} \to X$ 및 $X_{\overline{k}'} \to X_{k'}$
보편적인 동형이다. 따라서 $\mathcal{F}$의 코호몰로지
$X_{\overline{k}}$ 및 코호몰로지에서 계산될 수 있다.
$\mathcal{F}|_{X_{k'}}$는 $X_{\overline{k}'}$에서 계산될 수 있다.
명제 \ref{etale-cohomology-proposition-topological-invariance}를 참조하라.
따라서 우리는 대수적으로 다음의 경우로부터 일반적인 경우를 추론한다.
체를 닫았습니다.
\end{proof}






\section{코호몰로지 곡선의 꼬임 가군 } \label{etale-cohomology-section-vanishing-torsion-coefficients}

\noindent 이 절에서 우리는 비틀림인 뇌터 환에 대해 가군의 구성가능층에 대한 절 \ref{etale-cohomology-section-vanishing-torsion}의 인수를 반복한다. 우리는 가장 흥미로운 단계부터 시작한다.

\begin{lemma}
\label{etale-cohomology-lemma-constant-statements-coefficients}
$\Lambda$를 뇌터 환으로 두고, $M$를 유한 $\Lambda$-가군으로 가정한다.
는 정수 $n > 0$에 의해 소멸된다. $k$는 대수적으로 닫힌 체라고 가정한다.
그리고 보자$X$분리된 유한 유형이어야 한다.스킴~의차원 $\leq 1$~ 위에$k$.
그 다음에
\begin{enumerate}
\item $H^q_\etale(X, \underline{M})$는 유한 $\Lambda$-가군이다.
$n$가 $\text{char}(k)$의 소수인 경우,
\item $H^q_\etale(X, \underline{M})$는 유한 $\Lambda$-가군이다.
$X$가 적절한 경우.
\end{enumerate}
\end{lemma}

\begin{proof}
만약에$n = \ell n'$일부 소수의 경우$\ell$, 그 다음에
우리는 짧은 완전열 $0 \to M[\ell] \to M \to M' \to 0$를 얻습니다.
유한한 $\Lambda$-가군 및 $M'$는 $n'$에 의해 소멸된다.
이는 상수의 해당 짧은 완전열을 생성한다.
층, 이는 결국 완전열을 발생시킵니다.
코호몰로지 가군
$$
H^q_\etale(X, \underline{M[n]}) \to
H^q_\etale(X, \underline{M}) \to
H^q_\etale(X, \underline{M'})
$$
따라서 $M$이 소멸된 경우의 결과를 보여드릴 수 있다면
소수이면 $n$에 대한 귀납법으로 승리한다.

\medskip\noindent
$\ell$는 $\ell$가 $M$를 소멸시키는 소수라고 가정한다.
그런 다음 $\Lambda$를 $\mathbf{F}_\ell$-대수로 바꿀 수 있다.
$\Lambda/\ell \Lambda$. 즉, $\mathcal{F}$의 코호몰로지
$\Lambda$-가군의 층은 코호몰로지와 동일한다.
$\mathcal{F}$를 $\Lambda/\ell \Lambda$-가군의 층으로,
예에 대해
코호몰로지 사이트, 보조정리에
\ref{sites-cohomology-lemma-cohomology-modules-abelian-agree}.

\medskip\noindent
$\ell$이 $\ell$가 $M$를 소멸하는 소수라고 가정한다.
및 $\Lambda$.
$M$가 유한 자유 $\Lambda$-가군인 경우로 줄여보자.
즉, 짧은 완전열을 선택한다.
$$
0 \to N \to \Lambda^{\oplus m} \to M \to 0
$$
이는 완전열을 결정한다.
$$
H^q_\etale(X, \underline{\Lambda^{\oplus m}}) \to
H^q_\etale(X, \underline{M}) \to
H^{q + 1}_\etale(X, \underline{N})
$$
$q$에 대한 귀납법을 내림차순으로 $M$에 대한 결과를 얻습니다.
$\Lambda^{\oplus m}$에 대한 결과를 알고 있다면.
여기서는 코호몰로지 군이
도 단위로 사라지다 $> 2$
정리 \ref{etale-cohomology-theorem-vanishing-affine-curves}.

\medskip\noindent
$\ell$를 소수로 두고 $\ell$가 소멸한다고 가정한다.
$\Lambda$. 코호몰로지 군이
$H^q_\etale(X, \underline{\Lambda})$는 유한한 $\Lambda$-가군이다.
우리는 이것을 보여주기 위해 트릭을 사용할 것이다. '올바른'' 인수는 다음을 사용한다.
나중에 보여드릴 계수 정리 이다.
기준을 선택하세요 $\Lambda = \bigoplus_{i \in I} \mathbf{F}_\ell e_i$
그렇게$e_0 = 1$어떤 사람들에게는$0 \in I$.
이 기준의 선택에 따라 동형사상이 결정된다.
$$
\underline{\Lambda} = \bigoplus \underline{\mathbf{F}_\ell} e_i
$$
$X_\etale$의 층. 따라서 우리는
$$
H^q_\etale(X, \underline{\Lambda}) =
H^q_\etale(X, \bigoplus \underline{\mathbf{F}_\ell} e_i) =
\bigoplus H^q_\etale(X, \underline{\mathbf{F}_\ell})e_i
$$
코호몰로지를 $X$ 위로 직합으로 통근하므로
정리 \ref{etale-cohomology-theorem-colimit} 기준
(또는 보조정리 \ref{etale-cohomology-lemma-colimit} 또는
보조정리 \ref{etale-cohomology-lemma-direct-sum-bounded-below-cohomology}).
우리는 이미 $H^q_\etale(X, \underline{\mathbf{F}_\ell})$를 알고 있기 때문에
유한 차원 $\mathbf{F}_\ell$-벡터공간
(정리 \ref{etale-cohomology-theorem-vanishing-affine-curves} 기준),
$H^q_\etale(X, \underline{\Lambda})$
같은 순위의 $\Lambda$ 이상은 무료이다. 즉, 주어진
기초$\xi_1, \ldots, \xi_m$~의$H^q_\etale(X, \underline{\mathbf{F}_\ell})$
우리는 $\xi_1 e_0, \ldots, \xi_m e_0$가 $\Lambda$ 기반을 형성한다는 것을 알 수 있다.
$H^q_\etale(X, \underline{\Lambda})$.
\end{proof}

\begin{lemma} \label{etale-cohomology-lemma-finite-pushforward-coefficients} $\Lambda$는 뇌터 환, $k$는 대수적으로 닫힌 체, $f : X \to Y$는 유한이라고 가정 사상의 분리된 유한 유형 스킴의 $k$의 차원 $\leq 1$에 대해 $\mathcal{F}$를 층으로 설정한다. $\Lambda$-가군 $X_\etale$에 있다. $H^q_\etale(X, \mathcal{F})$가 유한한 $\Lambda$-가군인 경우 $H^q_\etale(Y, f_*\mathcal{F})$도 마찬가지이다. \end{lemma}

\begin{proof}
즉, $H^q_\etale(X, \mathcal{F}) = H^q_\etale(Y, f_*\mathcal{F})$
에 의해소멸~의$R^qf_*$~을 위한$q > 0$
(명제 \ref{etale-cohomology-proposition-finite-higher-direct-image-zero}) 및
르레이 스펙트럼 열
(코호몰로지 사이트, 보조정리 \ref{sites-cohomology-lemma-apply-Leray}).
\end{proof}

\begin{lemma} \label{etale-cohomology-lemma-restrict-to-open-coefficients} $\Lambda$를 뇌터 환으로 두고, $k$를 대수적으로 닫힌 체로 두고, $X$를 분리된 유한 유형으로 둡니다. 스킴 $k$의 차원 $\leq 1$에 대해, $\mathcal{F}$를 $\Lambda$-가군의 $X_\etale$의 구성가능층으로 둡니다. $j : X' \to X$는 조밀 개방형 하위 구성표를 포함한다고 가정한다. 그러면 $H^q_\etale(X, \mathcal{F})$는 유한 $\Lambda$-가군 필요충분조건은 $H^q_\etale(X, j_!j^{-1}\mathcal{F})$는 유한 $\Lambda$-가군이다. \end{lemma}

\begin{proof}
$X'$는 조밀이므로 $Z = X \setminus X'$에는 차원 $0$가 있음을 알 수 있다.
그러므로 는 유한하다집합 $Z = \{x_1, \ldots, x_n\}$~의$k$-합리적인 점.
짧은 완전열을 고려해보세요.
$$
0 \to j_!j^{-1}\mathcal{F} \to \mathcal{F} \to i_*i^{-1}\mathcal{F} \to 0
$$
보조정리 \ref{etale-cohomology-lemma-ses-associated-to-open}. 그것을 관찰하십시오
$H^q_\etale(X, i_*i^{-1}\mathcal{F}) =
H^q_\etale(Z, i^*\mathcal{F})$.
즉, $i : Z \to X$는 닫힌 몰입이므로 유한한다.
그러므로 우리는소멸~의$R^qi_*$~을 위한$q > 0$~에 의해
명제 \ref{etale-cohomology-proposition-finite-higher-direct-image-zero} 및
따라서 평등은 Leray 스펙트럼 열에서 따릅니다.
(코호몰로지 사이트, 보조정리 \ref{sites-cohomology-lemma-apply-Leray}).
$Z$은 대수적으로 닫힌 스펙트럼의 분리된 결합이므로
체, 우리는 $H^q_\etale(Z, i^*\mathcal{F}) = 0$라고 결론을 내립니다.
$q > 0$ 및
$$
H^0_\etale(Z, i^{-1}\mathcal{F}) =
\bigoplus\nolimits_{i = 1, \ldots, n} \mathcal{F}_{x_i}
$$
이는 유한한 $\Lambda$-가군 $\mathcal{F}_{x_i}$이다.
가정 $\mathcal{F}$는 $\Lambda$-가군의 구성가능층이다.
길고 정확한 코호몰로지 시퀀스는 완전열을 제공한다.
$$
0 \to
H^0_\etale(X, j_!j^{-1}\mathcal{F}) \to
H^0_\etale(X, \mathcal{F}) \to
H^0_\etale(Z, i^{-1}\mathcal{F}) \to
H^1_\etale(X, j_!j^{-1}\mathcal{F}) \to
H^1_\etale(X, \mathcal{F}) \to 0
$$
및 동형사상 $H^0_\etale(X, j_!j^{-1}\mathcal{F}) \to
H^0_\etale(X, \mathcal{F})$ 에 대하여 $q > 1$. 보조정리는 이것으로부터 쉽게 따릅니다.
\end{proof}

\begin{lemma}
\label{etale-cohomology-lemma-somewhat-easier-coefficients}
$\Lambda$를 뇌터 환으로 두고, $M$를 유한 $\Lambda$-가군으로 가정한다.
는 정수 $n > 0$에 의해 소멸된다. $k$는 대수적으로 닫힌 체라고 가정한다.
허락하다$X$분리된 유한 유형이어야 한다.스킴~의차원 $\leq 1$~ 위에$k$, 그리고
$j : U \to X$를 열린 몰입으로 설정한다. 그 다음에
\begin{enumerate}
\item $H^q_\etale(X, j_!\underline{M})$는 유한 $\Lambda$-가군이다.
$n$가 $\text{char}(k)$의 소수인 경우,
\item $H^q_\etale(X, j_!\underline{M})$는 유한 $\Lambda$-가군이다.
$X$가 적절한 경우.
\end{enumerate}
\end{lemma}

\begin{proof} $\dim(X) \leq 1$ 열린 $V \subset X$가 있으므로 $U$와 분리되어 $X' = U \cup V$는 조밀이 $X$에서 열려 있습니다(자세한 내용은 생략됨). $j' : X' \to X$가 사상 포함을 나타내는 경우 $j_!\underline{M}$는 $j'_!\underline{M}$의 직접 합산임을 알 수 있다. 따라서 $U$가 열려 있고 조밀이 $X$인 경우에는 보조정리를 증명하면 충분한다. 이 사례는 보조정리 \ref{etale-cohomology-lemma-restrict-to-open-coefficients} 및 \ref{etale-cohomology-lemma-constant-statements-coefficients}에 따릅니다. \end{proof}

\begin{lemma}
\label{etale-cohomology-lemma-ses-statements-coefficients}
$\Lambda$를 뇌터 환으로 두고, $k$를 대수적으로 닫힌 체로 두고,
허락하다$X$분리된 유한 유형이어야 함스킴~ 위에$k$~의차원 $\leq 1$, 그리고
$0 \to \mathcal{F}_1 \to \mathcal{F} \to \mathcal{F}_2 \to 0$
가 되다짧은 완전열~의층~의$\Lambda$-가군~에$X_\etale$. 만약에
$H^q_\etale(X, \mathcal{F}_i)$, $i = 1, 2$는 유한한다. $\Lambda$-가군
그러면 $H^q_\etale(X, \mathcal{F})$는 유한한 $\Lambda$-가군이다.
\end{lemma}

\begin{proof} 코호몰로지의 긴 완전열에서 즉시이다. \end{proof}

\begin{lemma}
\label{etale-cohomology-lemma-vanishing-easier-coefficients}
$\Lambda$를 뇌터 환으로 두고, $k$를 대수적으로 닫힌 체로 두고,
허락하다$X$분리된 유한 유형이어야 한다.스킴~의차원 $\leq 1$~ 위에$k$,
$j : U \to X$를 $U$ 연결로 열린 몰입으로 설정하고,
$\ell$는 소수이고, $n > 0$라고 하고, $\mathcal{G}$는 유한 유형이라고 두고,
국소상수층의 $\Lambda$-가군의 $U_\etale$에 의해 소멸되었다.
$\ell^n$. 그 다음에
\begin{enumerate}
\item $H^q_\etale(X, j_!\mathcal{G})$는 유한 $\Lambda$-가군이다.
$\ell$가 $\text{char}(k)$의 소수인 경우,
\item $H^q_\etale(X, j_!\mathcal{G})$는 유한 $\Lambda$-가군이다.
$X$가 적절한 경우.
\end{enumerate}
\end{lemma}

\begin{proof}
$f : V \to U$를 $\ell$에 대한 소수의 유한 에탈 사상라고 하자
보조정리 \ref{etale-cohomology-lemma-pullback-filtered-modules}에서와 같습니다. 의 토론
절 \ref{etale-cohomology-section-trace-method}는 지도를 제공한다.
$$
\mathcal{G} \to f_*f^{-1}\mathcal{G} \to \mathcal{G}
$$
그 구성은 동형사상이다. 그러므로 증명하는 것으로 충분하다.
$H^q_\etale(X, j_!f_*f^{-1}\mathcal{G})$의 유한성.
자리스키의 메인으로 정리
(사상, 보조정리에 대한 추가 정보
\ref{more-morphisms-lemma-quasi-finite-separated-pass-through-finite})
우리는 다이어그램을 선택할 수 있습니다
$$
\xymatrix{
V \ar[r]_{j'} \ar[d]_f & Y \ar[d]^{\overline{f}} \\
U \ar[r]^j & X
}
$$
$\overline{f} : Y \to X$ 유한 및 $j'$ 및 열린 몰입 포함
조밀 이미지. $f$는 유한하고 $V$ 조밀은 $Y$이므로 다음과 같습니다.
$V = U \times_X Y$. 보조정리 \ref{etale-cohomology-lemma-compatible-shriek-push-finite}에 의해 우리는
$$
j_!f_*f^{-1}\mathcal{G} = \overline{f}_*j'_!f^{-1}\mathcal{G}
$$
보조정리 \ref{etale-cohomology-lemma-finite-pushforward-coefficients}로 충분한다.
$j'_!f^{-1}\mathcal{G}$을 고려해보세요. 다음과 같이 주어진 여과의 존재
보조정리 \ref{etale-cohomology-lemma-pullback-filtered-modules},
$j'_!$이 정확하다는 사실
보조정리 \ref{etale-cohomology-lemma-ses-statements-coefficients}
우리를 사건으로 축소
$\mathcal{F} = j'_!\underline{M}$ 유한 $\Lambda$-가군 $M$
이는 보조정리 \ref{etale-cohomology-lemma-somewhat-easier-coefficients}이다.
\end{proof}

\begin{theorem}
\label{etale-cohomology-theorem-vanishing-affine-curves-coefficients}
$\Lambda$를 뇌터 환으로 두고, $k$를 대수적으로 닫힌 체로 두고,
허락하다$X$분리된 유한 유형이어야 한다.스킴~의차원 $\leq 1$~ 위에$k$,
그리고 $\mathcal{F}$를 $\Lambda$-가군의 구성가능층으로 설정한다.
비틀림인 $X_\etale$에서. 그 다음에
\begin{enumerate}
\item
\label{etale-cohomology-item-finite-prime-to-p-coefficients}
$H^q_\etale(X, \mathcal{F})$는 유한한 $\Lambda$-가군이다.
$\mathcal{F}$가 $\text{char}(k)$에 대한 비틀림 프라임인 경우,
\item
\label{etale-cohomology-item-finite-proper-coefficients}
$H^q_\etale(X, \mathcal{F})$유한하다$\Lambda$-가군만약에$X$적절하다.
\end{enumerate}
\end{theorem}

\begin{proof}
더 이상 언급하지 않고.
$\mathcal{F} = \mathcal{F}_1 \oplus \ldots \oplus \mathcal{F}_r$ 쓰기
여기서 $\mathcal{F}_i$는 $\ell_i^{n_i}$에 의해 소멸된다.
일부 소수 $\ell_i$ 및 정수 $n_i > 0$의 경우.
보조정리 \ref{etale-cohomology-lemma-ses-statements-coefficients}이면 충분한다.
$\mathcal{F}_i$에 대해 정리를 증명한다. 따라서 우리는 할 수 있고 할 수 있습니다
가정하다$\ell^n$죽인다$\mathcal{F}$어떤 소수에게는$\ell$
및 정수 $n > 0$.

\medskip\noindent
$\mathcal{F}$은 $\Lambda$-가군의 층으로 구성 가능하므로,
조밀 열려 있는 $U \subset X$가 있으므로 $\mathcal{F}|_U$은 다음과 같습니다.
유한 유형, $\Lambda$-가군의 국소상수층.
$\dim(X) \leq 1$ $U$을 축소한 후 다음과 같이 가정할 수 있다.
$U = U_1 \amalg \ldots \amalg U_n$는 기약의 분리된 합집합이다.
스킴 (교차점에 있는 닫힌 점만 제거하세요.
~의$\geq 2$구성 요소$U$). 에 의해
보조정리 \ref{etale-cohomology-lemma-restrict-to-open-coefficients}
우리는 사건을 줄인다
$\mathcal{F} = j_!\mathcal{G}$ 여기서 $\mathcal{G}$는 유한 유형이다.
국소상수층 $\Lambda$-가군의 $U$ (그리고 전멸됨)
$\ell^n$).

\medskip\noindent
우리가 선택한 이후로$U = U_1 \amalg \ldots \amalg U_n$~와 함께$U_i$ 기약
우리는
$$
j_!\mathcal{G} =
j_{1!}(\mathcal{G}|_{U_1}) \oplus \ldots \oplus
j_{n!}(\mathcal{G}|_{U_n})
$$
여기서 $j_i : U_i \to X$는 사상을 포함한다.
$j_{i!}(\mathcal{G}|_{U_i})$의 경우는 다음에서 처리된다.
보조정리 \ref{etale-cohomology-lemma-vanishing-easier-coefficients}.
\end{proof}








\section{적절한 스킴} \label{etale-cohomology-section-finite-etale-over-proper}의 첫 번째 코호몰로지

\noindent 기본 그룹에서 절 \ref{pione-section-finite-etale-over-proper} 우리는 어떤 의미에서 $R^1f_*\underline{G}$를 복용하면 $f : X \to Y$가 고유 사상이고 $G$가 유한 그룹인 경우 기본 변경으로 통근하는 것을 확인했다. (반드시 교환 가능한 것은 아닙니다). 이 절에서 우리는 이러한 결과의 유용한 결과를 추론한다.

\begin{lemma}
\label{etale-cohomology-lemma-proper-over-henselian-and-h1}
$A$를 헨셀리안 국소환이라고 가정한다. $X$를 $A$에 대해 적절한 스킴으로 설정한다.
닫힌 올 $X_0$. $M$를 유한 아벨 군으로 둡니다.
그런 다음 $H^1_\etale(X, \underline{M}) = H^1_\etale(X_0, \underline{M})$.
\end{lemma}

\begin{proof}
사이트에 코호몰로지로, 보조정리 \ref{sites-cohomology-lemma-torsors-h1}
$H^1_\etale(X, \underline{M})$의 요소는
$\underline{M}$-$\mathcal{F}$를 $X_\etale$에 저장한다.
그러한 토서는 분명히 유한한 국소상수층이다.
따라서 $\mathcal{F}$는 스킴 $V$ 유한으로 표현 가능한다.
$X$, 보조정리 \ref{etale-cohomology-lemma-characterize-finite-locally-constant}에 대한 에탈이다.
반대로,스킴 $V$유한 에탈 위의$X$와$M$-행동
그것은 그것을$M$-토르소 오버$X$을 발생시킨다코호몰로지
수업. $X_0$에 대한 코호몰로지 클래스 간의 동일한 번역과
토서들 유한 에탈은 $X_0$에 걸쳐 유지된다. 따라서 보조정리
범주의 동치의 결과이다.
기본 그룹, 보조정리
\ref{pione-lemma-finite-etale-on-proper-over-henselian}.
\end{proof}

\noindent
다음 기술 보조정리는 증명의 핵심 요소이다.
고유 기저변환 정리. 주장은 말 그대로 작동한다.
올에 차원이 있는 $A$에 대한 적절한 스킴에 대해
$\leq 1$, 그러나 실제로 결론은
고유 기저변환 정리 이 특정 버전만 필요한다.
증명에 있다.

\begin{lemma}
\label{etale-cohomology-lemma-efface-cohomology-on-fibre-by-finite-cover}
$A$를 헨셀리안 국소환이라고 가정한다. $X = \mathbf{P}^1_A$로 놔두세요.
$X_0 \subset X$를 닫힌 올로 둡니다. $\ell$를 소수라고 하자
숫자. $\mathcal{I}$를 다음의 단사 층으로 가정한다.
$\mathbf{Z}/\ell\mathbf{Z}$-가군 $X_\etale$에 있다. 그 다음에
$H^q_\etale(X_0, \mathcal{I}|_{X_0}) = 0$~을 위한$q > 0$.
\end{lemma}

\begin{proof}
$X$은 $2$로 처리될 수 있는 분리된 스킴이다.
아핀이 열립니다. 따라서 $q > 1$에 대해 이는 Gabber의 아핀에서 따릅니다.
고유 기저변환 정리의 변형, 참조
보조정리 \ref{etale-cohomology-lemma-vanishing-restriction-injective}.
따라서 우리는 $q = 1$라고 가정할 수 있다. 허락하다
$\xi \in H^1_\etale(X_0, \mathcal{I}|_{X_0})$.
목표: $\xi$가 $0$임을 보여줍니다.
기본형 \ref{etale-cohomology-lemma-torsion-colimit-constructible} 및
\ref{etale-cohomology-lemma-colimit} 지도를 찾을 수 있어요 $\mathcal{F} \to \mathcal{I}$
$\mathcal{F}$ 및 구성가능층
$\mathbf{Z}/\ell\mathbf{Z}$-가군
및 $\xi$는 다음의 $\zeta$ 요소에서 나옵니다.
$H^1_\etale(X_0, \mathcal{F}|_{X_0})$. 주입 맵이 있다고 가정한다.
$\mathcal{F} \to \mathcal{F}'$ 의 층 의
$\mathbf{Z}/\ell\mathbf{Z}$-가군 $X_\etale$에 있다.
$\mathcal{I}$은 단사형이므로 주어진 지도를 확장할 수 있다.
$\mathcal{F} \to \mathcal{I}$를 지도 $\mathcal{F}' \to \mathcal{I}$에 추가한다.
이 상황에서는 $\mathcal{F}$을 $\mathcal{F}'$로 바꿀 수 있다.
그리고$\zeta$의 이미지로$\zeta$~에$H^1_\etale(X_0, \mathcal{F}'|_{X_0})$.
또한, $\mathcal{F} = \mathcal{F}_1 \oplus \mathcal{F}_2$이 직합인 경우,
그러면 $\mathcal{F}$를 $\mathcal{F}_i$로 바꿀 수 있다.
그리고$\zeta$의 이미지로$\zeta$~에$H^1_\etale(X_0, \mathcal{F}_i|_{X_0})$.

\medskip\noindent
보조정리 \ref{etale-cohomology-lemma-constructible-maps-into-constant-general}에 의해
위의 설명은 $\mathcal{F}$라고 가정할 수 있다.
$f_*\underline{M}$ 형식이다. 여기서 $M$는 유한한다.
$\mathbf{Z}/\ell\mathbf{Z}$-가군
그리고 $f : Y \to X$는 유한한 표현의 유한한 사상이다.
(층은 여전히 다음과 같이 구성 가능한다.
보조정리 \ref{etale-cohomology-lemma-finite-pushforward-constructible}
하지만 우리는 이것이 필요하지 않습니다).
$f_*$의 형성으로 인해 어떠한 베이스 변화에도 통근이 가능한다.
(보조정리 \ref{etale-cohomology-lemma-finite-pushforward-commutes-with-base-change})
$f_*\underline{M}$에서 $X_0$로의 제한은 다음과 같습니다.
유도된 사상을 통한 $\underline{M}$의 직상과 동일
$Y_0 \to X_0$ 특수 올. 르레이의 스펙트럼 열
(명제 \ref{etale-cohomology-proposition-leray})
그리고 고차 직상의 소멸
(명제 \ref{etale-cohomology-proposition-finite-higher-direct-image-zero}),
우리는 발견
$$
H^1_\etale(X_0, f_*\underline{M}|_{X_0}) = H^1_\etale(Y_0, \underline{M}).
$$
$Y \to \Spec(A)$이 적절하므로 다음을 사용할 수 있다.
보조정리 \ref{etale-cohomology-lemma-proper-over-henselian-and-h1} 그걸 보시려면
$H^1_\etale(Y_0, \underline{M})$는 다음과 같습니다.
$H^1_\etale(Y, \underline{M})$. 따라서 우리는 코호몰로지
클래스 $\zeta$는 코호몰로지 클래스로 리프트된다.
$$
\tilde\zeta \in H^1_\etale(Y, \underline{M}) = H^1_\etale(X, f_*\underline{M})
$$
그러나 $\tilde \zeta$는 0으로 매핑된다.
$H^1_\etale(X, \mathcal{I})$는 $\mathcal{I}$가 단사이므로
그리고 교환성에 의해
$$
\xymatrix{
H^1_\etale(X, f_*\underline{M}) \ar[r] \ar[d] &
H^1_\etale(X, \mathcal{I}) \ar[d] \\
H^1_\etale(X_0, (f_*\underline{M})|_{X_0}) \ar[r] &
H^1_\etale(X_0, \mathcal{I}|_{X_0})
}
$$
우리는 이미지가$\xi$~의$\zeta$역시 0이다.
\end{proof}











\section{거점 변경 예비} \label{etale-cohomology-section-base-change-preliminaries}

\noindent
매끄러운 기저변환 정리에 관심이 있으시면
또는 고유 기저변환 정리로 바로 건너뛰어야 한다.
해당 절.
이 절과 다음 몇 가지 절에서는 가환 도식을 고려한다.
$$
\xymatrix{
X \ar[d]_f & Y \ar[l]^h \ar[d]^e \\
S & T \ar[l]_g
}
$$
스킴; 우리는 일반적으로 이 다이어그램이 데카르트 다이어그램이라고 가정한다.
즉, $Y = X \times_S T$이다. 가환 도식 위와 같음
가환 도식이 발생한다.
$$
\xymatrix{
X_\etale \ar[d]_{f_{small}} & Y_\etale \ar[d]^{e_{small}} \ar[l]^{h_{small}} \\
S_\etale & T_\etale \ar[l]_{g_{small}}
}
$$
작은 에탈 사이트. 표기법을 활용해보자
$$
f^{-1} = f_{small}^{-1}, \quad
g_* = g_{small, *}, \quad
e^{-1} = e_{small}^{-1}, \text{이고}\quad
h_* = h_{small, *}.
$$
사이트, 절 \ref{sites-section-pullback}에 의해
기본 변경 또는 당김 맵을 얻습니다.
$$
f^{-1}g_*\mathcal{F}
\longrightarrow
h_*e^{-1}\mathcal{F}
$$
한 동안층 $\mathcal{F}$~에$T_\etale$. 만약에$\mathcal{F}$아벨리안이다
층 위의 $T_\etale$, 파생된 기본 변경 맵을 얻습니다.
$$
f^{-1}Rg_*\mathcal{F}
\longrightarrow
Rh_*e^{-1}\mathcal{F}
$$
사이트, 보조정리의 코호몰로지 참조
\ref{sites-cohomology-lemma-base-change-map-flat-case}.
마지막으로, $K$가 $D(T_\etale)$의 임의의 객체인 경우
기지 변경 맵이 있습니다
$$
f^{-1}Rg_*K
\longrightarrow
Rh_*e^{-1}K
$$
보다
코호몰로지 사이트, 비고 \ref{sites-cohomology-remark-base-change}.

\begin{lemma}
\label{etale-cohomology-lemma-base-change-local}
스킴의 데카르트 다이어그램을 고려해보세요.
$$
\xymatrix{
X \ar[d]_f & Y \ar[l]^h \ar[d]^e \\
S & T \ar[l]_g
}
$$
$\{U_i \to X\}$를 에탈 피복으로 하여 $U_i \to S$
$U_i \to V_i \to S$ 와 $V_i \to S$ 에탈의 요소
데카르트 다이어그램을 고려하라.
$$
\xymatrix{
U_i \ar[d]_{f_i} & U_i \times_X Y \ar[l]^{h_i} \ar[d]^{e_i} \\
V_i & V_i \times_S T \ar[l]_{g_i}
}
$$
허락하다$\mathcal{F}$가 되다층~에$T_\etale$. 허락하다$K$~에$D(T_\etale)$.
집합 $K_i = K|_{V_i \times_S T}$ 및
$\mathcal{F}_i = \mathcal{F}|_{V_i \times_S T}$.
\begin{enumerate}
\item 만약 $f_i^{-1}g_{i, *}\mathcal{F}_i = h_{i, *}e_i^{-1}\mathcal{F}_i$
모든 $i$에 대해, 그 다음에는 $f^{-1}g_*\mathcal{F} = h_*e^{-1}\mathcal{F}$이다.
\item 만약 $f_i^{-1}Rg_{i, *}K_i = Rh_{i, *}e_i^{-1}K_i$
모든 $i$에 대해, 그 다음에는 $f^{-1}Rg_*K = Rh_*e^{-1}K$이다.
\item $\mathcal{F}$가 아벨 층이고
$f_i^{-1}R^qg_{i, *}\mathcal{F}_i = R^qh_{i, *}e_i^{-1}\mathcal{F}_i$
모든 $i$에 대해
$f^{-1}R^qg_*\mathcal{F} = R^qh_*e^{-1}\mathcal{F}$.
\end{enumerate}
\end{lemma}

\begin{proof} 증명 (1). 먼저 $$
(f^{-1}g_*\mathcal{F})|_{U_i} =
f_i^{-1}(g_*\mathcal{F}|_{V_i}) =
f_i^{-1}g_{i, *}\mathcal{F}_i
$$ $U_i \to X \to S$가 $U_i \to V_i \to S$와 동일하기 때문에 첫 번째 동일성과 $g_*\mathcal{F}|_{V_i} = g_{i, *}\mathcal{F}_i$가 사이트, 보조정리 \ref{sites-lemma-localize-morphism-strong}에 의해 동일하기 때문에 두 번째 동일함을 관찰한다. 마찬가지로 $$
(h_*e^{-1}\mathcal{F})|_{U_i} =
h_{i, *}(e^{-1}\mathcal{F}|_{U_i \times_X Y}) =
h_{i, *}e_i^{-1}\mathcal{F}_i
$$가 있다. 따라서 기본 변경 맵 $f_i^{-1}g_{i, *}\mathcal{F}_i \to h_{i, *}e_i^{-1}\mathcal{F}_i$이 모든 $i$에 대해 동형사상인 경우 기본 변경 맵 $f^{-1}g_*\mathcal{F} \to h_*e^{-1}\mathcal{F}$는 모든 $U_i$에 대해 동형사상으로 제한된다. $i$ 그리고 $\{U_i \to X\}$는 에탈 피복이므로 동형사상라고 결론을 내립니다.

\medskip\noindent 다른 두 진술의 경우 사이트, 보조정리 \ref{sites-lemma-localize-morphism-strong}에 대한 호소를 사이트, 보조정리에 대한 코호몰로지에 대한 호소로 대체한다. \ref{sites-cohomology-lemma-restrict-direct-image-open}. \end{proof}

\begin{lemma} \label{etale-cohomology-lemma-base-change-compose} 스킴 $$
\xymatrix{
W \ar[d]_i & Z \ar[l]^j \ar[d]^k \\
X \ar[d]_f & Y \ar[l]^h \ar[d]^e \\
S & T \ar[l]_g
}
$$의 직교 다이어그램 타워를 고려하라. $K$를 $D(T_\etale)$에 두십시오. $$
f^{-1}Rg_*K \to Rh_*e^{-1}K
\quad\text{그리고}\quad
i^{-1}Rh_*e^{-1}K \to Rj_*k^{-1}e^{-1}K
$$가 동형사상인 경우 $(f \circ i)^{-1}Rg_*K \to Rj_*(e \circ k)^{-1}K$는 동형사상이다. 마찬가지로, $\mathcal{F}$가 $T_\etale$에서 아벨 층이고 $$
f^{-1}R^qg_*\mathcal{F} \to R^qh_*e^{-1}\mathcal{F}
\quad\text{그리고}\quad
i^{-1}R^qh_*e^{-1}\mathcal{F} \to R^qj_*k^{-1}e^{-1}\mathcal{F}
$$가 동형사상인 경우 $(f \circ i)^{-1}R^qg_*\mathcal{F} \to R^qj_*(e \circ k)^{-1}\mathcal{F}$는 동형사상이다. \end{lemma}

\begin{proof} 이는 이러한 기본 변경 맵의 구성이 외부 직사각형에 대한 기본 변경 맵인지 확인하는 경우 형식적이다. 사이트, 비고 \ref{sites-cohomology-remark-compose-base-change-horizontal}의 코호몰로지를 참조하라. \end{proof}

\begin{lemma}
\label{etale-cohomology-lemma-base-change-Rf-star-colim}
$I$를 집합으로 지정한다. 역 시스템을 고려하라.
스킴의 데카르트 다이어그램
$$
\xymatrix{
X_i \ar[d]_{f_i} & Y_i \ar[l]^{h_i} \ar[d]^{e_i} \\
S_i & T_i \ar[l]_{g_i}
}
$$
아핀 전이 사상 및 $g_i$ 준컴팩트 및
준분리. 집합 $X = \lim X_i$,
$S = \lim S_i$, $T = \lim T_i$ 및 $Y = \lim Y_i$ ~
데카르트 다이어그램을 얻으십시오
$$
\xymatrix{
X \ar[d]_f & Y \ar[l]^h \ar[d]^e \\
S & T \ar[l]_g
}
$$
$(\mathcal{F}_i, \varphi_{i'i})$를 층의 시스템으로 설정한다.
$(T_i)$ 정의 \ref{etale-cohomology-definition-inverse-system-sheaves}에서와 같습니다. 집합
$\mathcal{F} = \colim p_i^{-1}\mathcal{F}_i$ 의 $T$
여기서 $p_i : T \to T_i$는 투영이다.
그러면 다음이 있다.
\begin{enumerate}
\item 만약 $f_i^{-1}g_{i, *}\mathcal{F}_i = h_{i, *}e_i^{-1}\mathcal{F}_i$
모든 $i$에 대해
$f^{-1}g_*\mathcal{F} = h_*e^{-1}\mathcal{F}$.
\item $\mathcal{F}_i$가 모든 $i$에 대해 아벨 층이고
$f_i^{-1}R^qg_{i, *}\mathcal{F}_i = R^qh_{i, *}e_i^{-1}\mathcal{F}_i$
모든 $i$에 대해
$f^{-1}R^qg_*\mathcal{F} = R^qh_*e^{-1}\mathcal{F}$.
\end{enumerate}
\end{lemma}

\begin{proof} (2)를 증명하고 (1)의 증명을 생략한다. 층의 당김은 왼쪽 수반이므로 여극한으로 통근한다는 점은 더 이상 언급하지 않고 사용하겠습니다. $h_i$은 $g_i$의 기본 변경으로 준컴팩트하고 준분리되어 있음을 확인하세요. $q_i : Y \to Y_i$ 투영을 표시하고 $e^{-1}\mathcal{F} = \colim e^{-1}p_i^{-1}\mathcal{F}_i =
\colim q_i^{-1}e_i^{-1}\mathcal{F}_i$을 관찰한다. 보조정리 \ref{etale-cohomology-lemma-relative-colimit-general}에 의해 $$
R^qh_*e^{-1}\mathcal{F} = \colim r_i^{-1}R^qh_{i, *}e_i^{-1}\mathcal{F}_i
$$가 제공된다. 여기서 $r_i : X \to X_i$는 투영이다. 마찬가지로 $$
f^{-1}Rg_*\mathcal{F} =
f^{-1}\colim s_i^{-1}R^qg_{i, *}\mathcal{F}_i =
\colim r_i^{-1}f_i^{-1}R^qg_{i, *}\mathcal{F}_i
$$가 있는데, 여기서 $s_i : S \to S_i$는 투영이다. 보조정리가 이어집니다. \end{proof}

\begin{lemma}
\label{etale-cohomology-lemma-base-change-Rf-star-colim-complexes}
$I$, $X_i$, $Y_i$, $S_i$, $T_i$, $f_i$, $h_i$, $e_i$, $g_i$,
$X$, $Y$, $S$, $T$, $f$, $h$, $e$, $g$는 다음과 같습니다.
보조정리 \ref{etale-cohomology-lemma-base-change-Rf-star-colim}.
$0 \in I$ 및 $K_0 \in D^+(T_{0, \etale})$를 허용한다.
$i \in I$의 경우 $i \geq 0$는 $K_i$의 하락을 나타냅니다.
$K_0$ ~ $T_i$. $K$는 $K_0$에서 $T$로의 철회를 나타냅니다.
$f_i^{-1}Rg_{i, *}K_i = Rh_{i, *}e_i^{-1}K_i$인 경우
모든 $i \geq 0$에 대해, 그 다음에는 $f^{-1}Rg_*K = Rh_*e^{-1}K$이다.
\end{lemma}

\begin{proof}
베이스 변경 맵을 보여주는 것으로 충분한다.
$f^{-1}Rg_*K \to Rh_*e^{-1}K$는 동형사상을 유도한다.
코호몰로지 층에 있다. 즉, 우리는 보여줘야 한다
$f^{-1}R^pg_*K \to R^ph_*e^{-1}K$는 동형사상이다.
모든 $p \in \mathbf{Z}$에 대해 다음과 같이 주어진다면
$f_i^{-1}R^pg_{i, *}K_i \to R^ph_{i, *}e_i^{-1}K_i$
모든 $i \geq 0$ 및 $p \in \mathbf{Z}$에 대한 동형사상이다.
이 시점에서 우리는 증명에서와 똑같이 논쟁할 수 있다.
보조정리 \ref{etale-cohomology-lemma-base-change-Rf-star-colim}
참조 대체
보조정리 \ref{etale-cohomology-lemma-relative-colimit-general}
에 대한 참조로
보조정리 \ref{etale-cohomology-lemma-relative-colimit-general-complexes}.
\end{proof}

\begin{lemma}
\label{etale-cohomology-lemma-base-change-f-star-general-stalks}
스킴의 데카르트 다이어그램을 고려해보세요.
$$
\xymatrix{
X \ar[d]_f & Y \ar[l]^h \ar[d]^e \\
S & T \ar[l]_g
}
$$
여기서 $g : T \to S$는 준컴팩트하고 준분리되어 있다.
$\mathcal{F}$를
아벨 층 - $T_\etale$. $q \geq 0$로 놔두세요. 다음은 동일한다.
\begin{enumerate}
\item 모든 기하점 $\overline{x}$ 중 $X$의 것과 그 상
$\overline{s} = f(\overline{x})$ 우리는
$$
H^q(\Spec(\mathcal{O}^{sh}_{X, \overline{x}}) \times_S T, \mathcal{F})
=
H^q(\Spec(\mathcal{O}^{sh}_{S, \overline{s}}) \times_S T, \mathcal{F})
$$
\item $f^{-1}R^qg_*\mathcal{F} \to R^qh_*e^{-1}\mathcal{F}$
동형사상이다.
\end{enumerate}
\end{lemma}

\begin{proof}
$Y = X \times_S T$ 이후 우리는
$\Spec(\mathcal{O}^{sh}_{X, \overline{x}}) \times_X Y =
\Spec(\mathcal{O}^{sh}_{X, \overline{x}}) \times_S T$. 따라서
(1)의 지도는 해당 지도의 $\overline{x}$에 있는 줄기의 지도이다.
(2)에서 정리 \ref{etale-cohomology-theorem-higher-direct-images}로 (그리고
보조정리 \ref{etale-cohomology-lemma-stalk-pullback}).
따라서 정리 \ref{etale-cohomology-theorem-exactness-stalks}에 의한 결과이다.
\end{proof}

\begin{lemma}
\label{etale-cohomology-lemma-check-stalks-better}
$f : X \to S$를 스킴의 사상으로 설정한다.
허락하다$\overline{x}$가 되다기하점~의$X$이미지 포함$\overline{s}$~에$S$.
$\Spec(K) \to \Spec(\mathcal{O}^{sh}_{S, \overline{s}})$
사상이고 $K$는 분리 가능하게 닫힌 체이다. $\mathcal{F}$를
아벨 층 - $\Spec(K)_\etale$. $q \geq 0$로 놔두세요. 다음은
동등한
\begin{enumerate}
\item
$H^q(\Spec(\mathcal{O}^{sh}_{X, \overline{x}}) \times_S \Spec(K), \mathcal{F}) =
H^q(\Spec(\mathcal{O}^{sh}_{S, \overline{s}}) \times_S \Spec(K), \mathcal{F})$
\item
$H^q(\Spec(\mathcal{O}^{sh}_{X, \overline{x}})
\times_{\Spec(\mathcal{O}^{sh}_{S, \overline{s}})} \Spec(K), \mathcal{F}) =
H^q(\Spec(K), \mathcal{F})$
\end{enumerate}
\end{lemma}

\begin{proof}
$\Spec(K) \times_S \Spec(\mathcal{O}^{sh}_{S, \overline{s}})$를 확인하세요.
는 $K$에 대한 에탈 대수학의 여과 여극한 스펙트럼이다.
$K$는 분리 가능하게 닫혀 있으므로 각 에탈 $K$-대수
는 $K$ 사본의 유한 곱이다. 따라서 우리는 쓸 수 있습니다
$$
\Spec(K) \times_S \Spec(\mathcal{O}^{sh}_{S, \overline{s}}) =
\lim_{i \in I} \coprod\nolimits_{a \in A_i} \Spec(K)
$$
각 용어가 분리된 공용체인 공동 필터링된 극한
사본의$\Spec(K)$유한한 이상집합 $A_i$.
$A_i$는 주어진 대로 비어 있지 않습니다.
$\Spec(K) \to \Spec(\mathcal{O}^{sh}_{S, \overline{s}})$.
그것은 다음과 같습니다
\begin{align*}
\Spec(\mathcal{O}^{sh}_{X, \overline{x}}) \times_S \Spec(K)
& =
\Spec(\mathcal{O}^{sh}_{X, \overline{x}})
\times_{\Spec(\mathcal{O}^{sh}_{S, \overline{s}})}
\left(
\Spec(\mathcal{O}^{sh}_{S, \overline{s}})
\times_S \Spec(K)\right) \\
& =
\lim_{i \in I} \coprod\nolimits_{a \in A_i}
\Spec(\mathcal{O}^{sh}_{X, \overline{x}})
\times_{\Spec(\mathcal{O}^{sh}_{S, \overline{s}})} \Spec(K)
\end{align*}
설정에서 코호몰로지를 타고 극한으로 통근하므로
스킴 (정리 \ref{etale-cohomology-theorem-colimit}) 결론을 내립니다.
\end{proof}













\section{직상 기본 변경} \label{etale-cohomology-section-base-change-f-star}

\noindent
이 절은 예비적이므로 첫 독해에서는 건너뛰어야 한다.
이 절에서는 사상 $f : X \to S$에 대해 논의한다.
$f^{-1}g_* = h_*e^{-1}$ 모든 층 (집합) 모든 직교 좌표에 대해
도표
$$
\xymatrix{
X \ar[d]_f & Y \ar[l]^h \ar[d]^e \\
S & T \ar[l]_g
}
$$
$g$ 준컴팩트 및 준분리형을 사용한다.

\begin{lemma}
\label{etale-cohomology-lemma-base-change-f-star-general}
스킴의 데카르트 다이어그램을 고려해보세요.
$$
\xymatrix{
X \ar[d]_f & Y \ar[l]^h \ar[d]^e \\
S & T \ar[l]_g
}
$$
가정$f$평평하고 모든 물체는$U$~의$X_\etale$가지다
 피복 $\{U_i \to U\}$ $U_i \to S$
요인$U_i \to V_i \to S$~와 함께$V_i \to S$
에탈 및 $U_i \to V_i$ 준 컴팩트
기하학적으로 연결 올.
그렇다면 누구에게나층 $\mathcal{F}$~의집합~에$T_\etale$우리는
$f^{-1}g_*\mathcal{F} = h_*e^{-1}\mathcal{F}$.
\end{lemma}

\begin{proof}
$U \to X$를 에탈 사상으로 하여 $U \to S$가 다음과 같이 고려되도록 한다.
$U \to V \to S$ 와 $V \to S$ 에탈 및 $U \to V$ 준 컴팩트
기하학적으로 연결 올이다. $U \to V$가 평평한지 확인하세요.
(평탄도에 대한 자세한 내용, 보조정리 \ref{flat-lemma-etale-flat-up-down}).
우리는 주장
\begin{align*}
f^{-1}g_*\mathcal{F}(U)
& = g_*\mathcal{F}(V) \\
& = \mathcal{F}(V \times_S T) \\
& = e^{-1}\mathcal{F}(U \times_X Y) \\
& = h_*e^{-1}\mathcal{F}(U)
\end{align*}
즉, $U$을 $X_\etale$의 객체로 생각하고,
$V$ $S_\etale$의 객체로서 우리는 첫 번째 평등을 봅니다.
엄밀히 말하면 보조정리 \ref{etale-cohomology-lemma-sections-upstairs}\footnote{에서 따릅니다.
말하자면, 우리는 또한 $f^{-1}g_*\mathcal{F}$의
$U_\etale$로의 제한이 $U \to V$를 따라
$g_*\mathcal{F}$의 $V_\etale$로의 제한을 당긴 것임을 사용하고 있다. 다음을 보라.
사이트, 보조정리 \ref{sites-lemma-localize-morphism-strong}.}.
$V \times_S T$을 $T_\etale$의 객체로 생각
두 번째 동일성은 $g_*$의 정의에서 따릅니다.
$U \times_X Y = U \times_S T$ ($Y = X \times_S T$ 때문에)
따라서 $U \times_X Y \to V \times_S T$
$U \to V$의 기본 변경으로 기하학적으로 연결 올이 있다.
$U \times_X Y$을 $Y_\etale$의 객체로 생각하면 다음과 같습니다.
세 번째 등식은 보조정리 \ref{etale-cohomology-lemma-sections-upstairs}에서 따릅니다.
이전처럼. 마지막으로, 네 번째 동등성은 정의에서 따릅니다.
$h_*$.

\medskip\noindent 가정에 의해 $X_\etale$의 모든 객체에는 이전 단락의 인수가 적용되는 에탈 피복이 있으므로 보조정리가 참임을 알 수 있다. \end{proof}

\begin{lemma}
\label{etale-cohomology-lemma-fppf-reduced-fibres-base-change-f-star}
스킴의 데카르트 다이어그램을 고려해보세요.
$$
\xymatrix{
X \ar[d]_f & Y \ar[l]^h \ar[d]^e \\
S & T \ar[l]_g
}
$$
여기서 $f$는 평평하고 국부적으로 유한하게 표현된다.
기하학적으로 축소된 올.
그런 다음 $f^{-1}g_*\mathcal{F} = h_*e^{-1}\mathcal{F}$
누구에게나층 $\mathcal{F}$~에$T_\etale$.
\end{lemma}

\begin{proof} 보조정리 \ref{etale-cohomology-lemma-base-change-f-star-general}와 사상, 보조정리 \ref{more-morphisms-lemma-cover-by-geometrically-connected}에 대한 자세한 내용을 결합한다. \end{proof}

\begin{lemma}
\label{etale-cohomology-lemma-base-change-f-star-field}
스킴의 데카르트 다이어그램을 고려해보세요.
$$
\xymatrix{
X \ar[d]_f & Y \ar[l]^h \ar[d]^e \\
S & T \ar[l]_g
}
$$
$S$는 분리 가능하게 닫힌 체의 스펙트럼이라고 가정한다.
그런 다음 $f^{-1}g_*\mathcal{F} = h_*e^{-1}\mathcal{F}$
누구에게나층 $\mathcal{F}$~에$T_\etale$.
\end{lemma}

\begin{proof}
$X$에서 로컬로 작업할 수 있다. 그러므로 우리는 $X$가 유사하다고 가정할 수 있다.
그런 다음 $X$을 아핀 스킴의 공동 필터링된 극한으로 쓸 수 있다.
$S$에 대한 유한 유형이다. 보조정리 \ref{etale-cohomology-lemma-base-change-Rf-star-colim}에 의해
$X$는 $S$에 대해 유한 유형이라고 가정할 수 있다.
그러면 보조정리 \ref{etale-cohomology-lemma-base-change-f-star-general}
유한의 스킴 때문에 적용된다.
분리 가능하게 닫힌 체에 대한 유형은 유한 분리이다.
연결과 기하학적으로 연결 스킴의 합집합
(다양체, 보조정리 참조)
\ref{varieties-lemma-separably-closed-field-connected-components}).
\end{proof}

\begin{lemma}
\label{etale-cohomology-lemma-base-change-f-star-valuation}
스킴의 데카르트 다이어그램을 고려해보세요.
$$
\xymatrix{
X \ar[d]_f & Y \ar[l]^h \ar[d]^e \\
S & T \ar[l]_g
}
$$
가정
\begin{enumerate}
\item $f$는 평평하고 개방적이다.
\item $S$의 잔기 체는 분리 가능하게 대수적으로 닫혀 있다.
\item 에탈 사상 $U \to X$와 $U$ 아핀이 주어졌습니다.
$U$를 개방형 하위 구성표의 유한 분리 결합으로 작성할 수 있다.
의 $X$ (예의 경우 $X$가 정규 정역 스킴인 경우
분리 가능하게 닫힌 기능 포함 체),
\item 올 $X_s$ 중 $f$의 것의 임의의 공집합이 아닌 열린부분은 연결이다.
(예의 경우 $X_s$가 기약이거나 비어 있는 경우).
\end{enumerate}
그렇다면 누구에게나층 $\mathcal{F}$~의집합~에$T_\etale$우리는
$f^{-1}g_*\mathcal{F} = h_*e^{-1}\mathcal{F}$.
\end{lemma}

\begin{proof} 생략한다. 힌트: 가정은 보조정리 \ref{etale-cohomology-lemma-base-change-f-star-general}의 조건을 거의 의미하지 않습니다. 예 부분(3)은 보조정리 \ref{etale-cohomology-lemma-normal-scheme-with-alg-closed-function-field}에서 따릅니다. \end{proof}

\noindent 다음 보조정리는 여기에 속하지 않지만 어디에도 적합한 곳은 없는 것 같습니다.

\begin{lemma} \label{etale-cohomology-lemma-fppf-reduced-fibres-pullback-products} $f : X \to S$는 기하학적으로 축소된 올을 사용하여 평평하고 국부적으로 유한하게 표현되는 스킴의 사상라고 가정한다. 그런 다음 $f^{-1} : \Sh(S_\etale) \to \Sh(X_\etale)$는 제품을 가지고 출퇴근한다. \end{lemma}

\begin{proof}
$I$를 집합으로 두고 $\mathcal{G}_i$를 $S_\etale$에서 층으로 둡니다.
$i \in I$의 경우.
$U \to X$를 에탈 사상으로 하여 $U \to S$가 다음과 같이 고려되도록 한다.
$U \to V \to S$ 와 $V \to S$ 에탈 및 $U \to V$ 유한 평면
기하학적으로 연결 올로 표현된다. 그럼 우리는
\begin{align*}
f^{-1}(\prod \mathcal{G}_i)(U)
& =
(\prod \mathcal{G}_i)(V) \\
& =
\prod \mathcal{G}_i(V) \\
& =
\prod f^{-1}\mathcal{G}_i(U) \\
& =
(\prod f^{-1}\mathcal{G}_i)(U)
\end{align*}
여기서는 보조정리 \ref{etale-cohomology-lemma-sections-upstairs}를 사용했다.
첫 번째와 세 번째 평등에서
(우리는 또한 $f^{-1}\mathcal{G}$의
$U_\etale$로의 제한이 $U \to V$를 따라
$\mathcal{G}$의 $V_\etale$로의 제한을 당긴 것임을 사용한다. 다음을 참조하라.
사이트, 보조정리 \ref{sites-lemma-localize-morphism-strong}).
사상, 보조정리에 대한 추가 정보
\ref{more-morphisms-lemma-cover-by-geometrically-connected}
모든 물건$U$~의$X_\etale$에탈이 있다피복
$\{U_i \to U\}$ 이전의 논의와 같이
단락은 $U_i$에 적용된다. 보조정리가 이어집니다.
\end{proof}

\begin{lemma}
\label{etale-cohomology-lemma-base-change-f-star}
$f : X \to S$를 스킴의 평탄 사상으로 설정한다.
그것은 모든 사람에게기하점 $\overline{x}$~의$X$지도
$$
\mathcal{O}_{S, f(\overline{x})}^{sh}
\longrightarrow
\mathcal{O}_{X, \overline{x}}^{sh}
$$
기하학적으로 연결 올을 갖습니다. 그러면 매
스킴의 데카르트 다이어그램
$$
\xymatrix{
X \ar[d]_f & Y \ar[l]^h \ar[d]^e \\
S & T \ar[l]_g
}
$$
$g$ 준컴팩트하고 준분리된 방식으로 우리는
$f^{-1}g_*\mathcal{F} = h_*e^{-1}\mathcal{F}$
누구에게나층 $\mathcal{F}$~의집합~에$T_\etale$.
\end{lemma}

\begin{proof}
줄기에서 동등성을 확인하는 것으로 충분한다.
정리 \ref{etale-cohomology-theorem-exactness-stalks}.
정리 \ref{etale-cohomology-theorem-higher-direct-images}에 의해 우리는
$$
(h_*e^{-1}\mathcal{F})_{\overline{x}} =
\Gamma(\Spec(\mathcal{O}_{X, \overline{x}}^{sh}) \times_X Y, e^{-1}\mathcal{F})
$$
그리고 우리도 비슷하게
$$
(f^{-1}g_*^{-1}\mathcal{F})_{\overline{x}} =
(g_*^{-1}\mathcal{F})_{f(\overline{x})} =
\Gamma(\Spec(\mathcal{O}_{S, f(\overline{x})}^{sh}) \times_S T, \mathcal{F})
$$
이 집합은 보조정리 \ref{etale-cohomology-lemma-sections-upstairs}를 적용하여 동일한다.
사상으로
$$
\Spec(\mathcal{O}_{X, \overline{x}}^{sh}) \times_X Y
\longrightarrow
\Spec(\mathcal{O}_{S, f(\overline{x})}^{sh}) \times_S T
$$
이는 기본 변경 사항이다.
$\Spec(\mathcal{O}_{X, \overline{x}}^{sh}) \to
\Spec(\mathcal{O}_{S, f(\overline{x})}^{sh})$
왜냐하면 $Y = X \times_S T$ 때문이다.
\end{proof}







\section{고차 직상} \label{etale-cohomology-section-base-change}에 대한 기본 변경

\noindent 이 절은 고차 직상에 대한 절 \ref{etale-cohomology-section-base-change-f-star}와 유사한다. 이 절은 예비적이므로 첫 독해에서는 건너뛰어야 한다.

\begin{remark}
\label{etale-cohomology-remark-base-change-holds}
$f : X \to S$를 스킴의 사상으로 설정한다. $n$를 정수로 둡니다.
모든 가환 도식에 대해 $BC(f, n, q_0)$가 참이라고 말할 것이다.
$$
\xymatrix{
X \ar[d]_f & X' \ar[l] \ar[d]_{f'} & Y \ar[l]^h \ar[d]^e \\
S & S' \ar[l] & T \ar[l]_g
}
$$
$X' = X \times_S S'$ 및 $Y = X' \times_{S'} T$ 및
$g$ 준컴팩트하고 준분리형이며, 모든 아벨 층
$\mathcal{F}$ $T_\etale$의 $n$ 기지 변경 맵에 의해 전멸되었다.
$$
(f')^{-1}R^qg_*\mathcal{F}
\longrightarrow
R^qh_*e^{-1}\mathcal{F}
$$
$q \leq q_0$에 대한 동형사상이다.
\end{remark}

\begin{lemma}
\label{etale-cohomology-lemma-base-change-q-injective}
와 함께$f : X \to S$그리고$n$에서와 같이비고 \ref{etale-cohomology-remark-base-change-holds}
일부 $q \geq 1$에 대해 $BC(f, n, q - 1)$이 있다고 가정한다. 그 다음에
가환 도식마다
$$
\xymatrix{
X \ar[d]_f & X' \ar[l] \ar[d]_{f'} & Y \ar[l]^h \ar[d]^e \\
S & S' \ar[l] & T \ar[l]_g
}
$$
$X' = X \times_S S'$ 및 $Y = X' \times_{S'} T$ 및
$g$ 준컴팩트하고 준분리형이며, 모든 아벨 층
$\mathcal{F}$ $T_\etale$의 $n$에 의해 소멸됨
\begin{enumerate}
\item 기지 변경 맵
$(f')^{-1}R^qg_*\mathcal{F}\to R^qh_*e^{-1}\mathcal{F}$
주입식이다,
\item 만약 $\mathcal{F} \subset \mathcal{G}$ 여기서 $\mathcal{G}$
$T_\etale$에서 $n$에 의해 소멸되고,
$$
\Coker\left(
(f')^{-1}R^qg_*\mathcal{F}\to R^qh_*e^{-1}\mathcal{F}
\right)
\subset
\Coker\left(
(f')^{-1}R^qg_*\mathcal{G}\to R^qh_*e^{-1}\mathcal{G}
\right)
$$
\item 만약 에서 (2) 층 $\mathcal{G}$는 분사형 층이다.
$\mathbf{Z}/n\mathbf{Z}$-가군의 다음
$$
\Coker\left((f')^{-1}R^qg_*\mathcal{F}\to R^qh_*e^{-1}\mathcal{F} \right)
\subset R^qh_*e^{-1}\mathcal{G}
$$
\end{enumerate}
\end{lemma}

\begin{proof}
짧은 완전열을 선택하세요.
$0 \to \mathcal{F} \to \mathcal{I} \to \mathcal{Q} \to 0$
여기서 $\mathcal{I}$는 $\mathbf{Z}/n\mathbf{Z}$-가군의 단사 층이다.
유도된 다이어그램을 고려해보세요
$$
\xymatrix{
(f')^{-1}R^{q - 1}g_*\mathcal{I} \ar[d]_{\cong} \ar[r] &
(f')^{-1}R^{q - 1}g_*\mathcal{Q} \ar[d]_{\cong} \ar[r] &
(f')^{-1}R^qg_*\mathcal{F} \ar[d] \ar[r] &
0 \ar[d] \\
R^{q - 1}h_*e^{-1}\mathcal{I} \ar[r] &
R^{q - 1}h_*e^{-1}\mathcal{Q} \ar[r] &
R^qh_*e^{-1}\mathcal{F} \ar[r] &
R^qh_*e^{-1}\mathcal{I}
}
$$
정확한 행이 있다. 오른쪽 상단 모서리에 0이 있다.
$\mathcal{I}$는 단사형이다. 왼쪽 두 개의 수직 화살표는
동형사상~에 의해$BC(f, n, q - 1)$. 우리는 그 부분을 결론짓는다(1) 보유하고 있다.
위의 내용도 이를 보여줍니다.
$$
\Coker\left(
(f')^{-1}R^qg_*\mathcal{F}\to R^qh_*e^{-1}\mathcal{F}
\right)
\subset
R^qh_*e^{-1}\mathcal{I}
$$
따라서 부분 (3)이 유지된다. (2)를 증명하려면 다음을 선택하세요.
$\mathcal{F} \subset \mathcal{G} \subset \mathcal{I}$.
\end{proof}

\begin{lemma}
\label{etale-cohomology-lemma-base-change-q-integral-top}
와 함께$f : X \to S$그리고$n$에서와 같이비고 \ref{etale-cohomology-remark-base-change-holds}
일부 $q \geq 1$에 대해 $BC(f, n, q - 1)$이 있다고 가정한다. 고려하다
가환 도식들
$$
\vcenter{
\xymatrix{
X \ar[d]_f &
X' \ar[d]_{f'} \ar[l] &
Y \ar[l]^h \ar[d]^e &
Y' \ar[l]^{\pi'} \ar[d]^{e'} \\
S &
S' \ar[l] &
T \ar[l]_g &
T' \ar[l]_\pi
}
}
\quad\text{및}\quad
\vcenter{
\xymatrix{
X' \ar[d]_{f'} & & Y' \ar[ll]^{h' = h \circ \pi'} \ar[d]^{e'} \\
S' & & T' \ar[ll]_{g' = g \circ \pi}
}
}
$$
여기서 모든 정사각형은 데카르트이고, $g$ 준간밀하고 준분리되어 있으며,
$\pi$는 정역 전사이다. $\mathcal{F}$를 아벨 층으로 설정한다.
$T_\etale$에서 $n$와 집합 $\mathcal{F}' = \pi^{-1}\mathcal{F}$에 의해 소멸되었다.
베이스 변경 맵의 경우
$$
(f')^{-1}R^qg'_*\mathcal{F}' \longrightarrow R^qh'_*(e')^{-1}\mathcal{F}'
$$
동형사상이면 기본 변경 맵이다.
$(f')^{-1}R^qg_*\mathcal{F} \to R^qh_*e^{-1}\mathcal{F}$
동형사상이다.
\end{lemma}

\begin{proof}
$\mathcal{F} \to \pi_*\pi^{-1}\mathcal{F}'$는 분사형이라는 점에 유의하세요.
$\pi$은 전사입니다(줄기를 확인하세요). 따라서
보조정리 \ref{etale-cohomology-lemma-base-change-q-injective}
우리는 기본 변경 맵을 보여주는 것으로 충분하다는 것을 알 수 있다.
$$
(f')^{-1}R^qg_*\pi_*\mathcal{F}'
\longrightarrow
R^qh_*e^{-1}\pi_*\mathcal{F}'
$$
동형사상이다. 이는 가정에서 따온 것이다.
$R^qg_*\pi_*\mathcal{F}' = R^qg'_*\mathcal{F}'$이 있다.
$e^{-1}\pi_*\mathcal{F}' =\pi'_*(e')^{-1}\mathcal{F}'$이 있고,
$R^qh_*\pi'_*(e')^{-1}\mathcal{F}' = R^qh'_*(e')^{-1}\mathcal{F}'$이 있다.
이것은 보조정리에서 따온 것이다.
\ref{etale-cohomology-lemma-integral-pushforward-commutes-with-base-change} 그리고
\ref{etale-cohomology-lemma-what-integral} 및 상대 레이 스펙트럼 열
(코호몰로지 사이트, 보조정리 \ref{sites-cohomology-lemma-relative-Leray}).
\end{proof}

\begin{lemma}
\label{etale-cohomology-lemma-base-change-q-integral-bottom}
와 함께$f : X \to S$그리고$n$에서와 같이비고 \ref{etale-cohomology-remark-base-change-holds}
일부 $q \geq 1$에 대해 $BC(f, n, q - 1)$이 있다고 가정한다. 고려하다
가환 도식들
$$
\vcenter{
\xymatrix{
X \ar[d]_f &
X' \ar[d]_{f'} \ar[l] &
X'' \ar[l]^{\pi'} \ar[d]_{f''} &
Y \ar[l]^{h'} \ar[d]^e \\
S &
S' \ar[l] &
S'' \ar[l]_\pi &
T \ar[l]_{g'}
}
}
\quad\text{및}\quad
\vcenter{
\xymatrix{
X' \ar[d]_{f'} & & Y \ar[ll]^{h = h' \circ \pi'} \ar[d]^e \\
S' & & T \ar[ll]_{g = g' \circ \pi}
}
}
$$
여기서 모든 정사각형은 데카르트이고, $g'$ 준간밀하고 준분리되어 있으며,
$\pi$은 정역이다. $\mathcal{F}$를 아벨 층으로 설정한다.
$T_\etale$에서 $n$에 의해 전멸되었다. 베이스 변경 맵의 경우
$$
(f')^{-1}R^qg_*\mathcal{F} \longrightarrow R^qh_*e^{-1}\mathcal{F}
$$
동형사상이면 기본 변경 맵이다.
$(f'')^{-1}R^qg'_*\mathcal{F} \to R^qh'_*e^{-1}\mathcal{F}$
동형사상이다.
\end{lemma}

\begin{proof} $\pi$ 및 $\pi'$는 정역이므로 $R\pi_* = \pi_*$ 및 $R\pi'_* = \pi'_*$가 있다. 보조정리 \ref{etale-cohomology-lemma-what-integral}를 참조하라. $(f')^{-1}\pi_* = \pi'_*(f'')^{-1}$도 있다. 따라서 $\pi'_*(f'')^{-1}R^qg'_*\mathcal{F} = (f')^{-1}R^qg_*\mathcal{F}$ 및 $\pi'_*R^qh'_*e^{-1}\mathcal{F} = R^qh_*e^{-1}\mathcal{F}$가 표시된다. 따라서 가정은 함자 $\pi'_*$를 적용한 후 지도가 동형사상이 됨을 의미한다. 따라서 보조정리 \ref{etale-cohomology-lemma-what-integral}에 의한 동형사상임을 알 수 있다. \end{proof}

\begin{lemma}
\label{etale-cohomology-lemma-formal-argument}
$T$를 준간소화하고 준분리된 스킴라고 하자.
준컴팩트 및 준분리형의 경우 $P$를 성질로 설정한다.
$T$에 대한 스킴. 추정하다
\begin{enumerate}
\item 만약 $T'' \to T'$이 준 컴팩트하고 두꺼워지면
$T$에 걸쳐 준분리된 스킴, 그 다음 $P(T'')$ 필요충분조건은 $P(T')$.
\item $T' = \lim T_i$가 역 시스템의 극한인 경우
아핀을 사용하여 $T$ 위에 준 컴팩트하고 준 분리된 스킴
전이 사상 및 $P(T_i)$는 모든 $i$에 대해 유지된다.
$P(T')$이 유지된다.
\item $Z \subset T'$가 다음과 같은 폐쇄형 하위 체계인 경우
준컴팩트 보수 $V \subset T'$ 및 $P(T')$가 유지된다.
그러면 $P(V)$ 또는 $P(Z)$ 중 하나가 유지된다.
\end{enumerate}
그 다음에$P(T)$암시한다$P(\Spec(K))$어떤 사람들에게는사상 $\Spec(K) \to T$
여기서 $K$는 체이다.
\end{lemma}

\begin{proof}
다음을 고려하라.집합 $\mathfrak T$폐쇄된 하위 계획의$T' \subset T$
$P(T')$와 같습니다. 가정 (2)에 의해 이 집합에는 최소한의 요소가 있다.
$T'$라고 말하세요. 가정 (1)에 의해 $T'$가 감소하는 것을 볼 수 있다.
$\eta \in T'$를 기약의 일반 지점으로 설정한다.
$T'$의 구성요소이다. 그런 다음 일부 체에 대해 $\eta = \Spec(K)$
$K$ 및 $\eta = \lim V$ 여기서 극한은 아핀 위에 있다.
개방형 하위 제도$V \subset T'$포함하는$\eta$.
가정 (3)와 $T'$의 최소값을 보면
$P(V)$은 모든 $V$에 적용된다. 따라서 $P(\eta)$
(2)에 의해 증명이 완료되었다.
\end{proof}

\begin{lemma}
\label{etale-cohomology-lemma-base-change-does-not-hold-pre}
와 함께$f : X \to S$그리고$n$에서와 같이비고 \ref{etale-cohomology-remark-base-change-holds}
일부 $q \geq 1$에 대해 $BC(f, n, q - 1)$이 참이라고 가정하지만
$BC(f, n, q)$는 아닙니다. 그러면 가환 도식이 존재한다.
$$
\xymatrix{
X \ar[d]_f & X' \ar[d]_{f'} \ar[l] & Y \ar[l]^h \ar[d]^e \\
S & S' \ar[l] & \Spec(K) \ar[l]_g
}
$$
여기서 $X' = X \times_S S'$, $Y = X' \times_{S'} \Spec(K)$,
$K$은 체이고, $\mathcal{F}$는 아벨 층이다.
$\Spec(K)$에서 $n$에 의해 소멸되었다.
$(f')^{-1}R^qg_*\mathcal{F} \to R^qh_*e^{-1}\mathcal{F}$
동형사상이 아닙니다.
\end{lemma}

\begin{proof}
가환 도식을 선택하세요.
$$
\xymatrix{
X \ar[d]_f & X' \ar[l] \ar[d]_{f'} & Y \ar[l]^h \ar[d]^e \\
S & S' \ar[l] & T \ar[l]_g
}
$$
$X' = X \times_S S'$ 및 $Y = X' \times_{S'} T$ 및
$g$ 준컴팩트 및 준분리형, 그리고 아벨 층
$\mathcal{F}$ $T_\etale$가 $n$에 의해 소멸되었다.
베이스 변경 맵
$(f')^{-1}R^qg_*\mathcal{F} \to R^qh_*e^{-1}\mathcal{F}$
동형사상이 아닙니다. 물론 우리는 $S'$를 대체할 수도 있고 실제로 대체할 수도 있다.
$S'$의 아핀 오픈에 의해; 이는 $T$가 준컴팩트함을 의미한다.
그리고 준분리. 보조정리 \ref{etale-cohomology-lemma-base-change-q-injective}에 의해 우리는
$(f')^{-1}R^qg_*\mathcal{F} \to R^qh_*e^{-1}\mathcal{F}$
주입식이다. 고르다기하점 $\overline{x}$~의$X'$
그리고 요소$\xi$~의$(R^qh_*q^{-1}\mathcal{F})_{\overline{x}}$
지도의 이미지에 없는 것
$((f')^{-1}R^qg_*\mathcal{F})_{\overline{x}} \to
(R^qh_*e^{-1}\mathcal{F})_{\overline{x}}$.

\medskip\noindent
다음을 고려해보세요.사상 $\pi : T' \to T$~와 함께$T'$준컴팩트
준분리되어 $\mathcal{F}' = \pi^{-1}\mathcal{F}$을 나타낸다.
$\pi' : Y' = Y \times_T T' \to Y$ $\pi$의 기본 변경을 나타냅니다.
및 $e' : Y' \to T'$ $e$의 기본 변경 사항이다. 그림
$$
\vcenter{
\xymatrix{
X' \ar[d]_{f'} & Y \ar[l]^h \ar[d]^e & Y' \ar[l]^{\pi'} \ar[d]^{e'} \\
S' & T \ar[l]_g & T' \ar[l]_\pi
}
}
\quad\text{및}\quad
\vcenter{
\xymatrix{
X' \ar[d]_{f'} & & Y' \ar[ll]^{h' = h \circ \pi'} \ar[d]^{e'} \\
S' & & T' \ar[ll]_{g' = g \circ \pi}
}
}
$$
당김 맵을 사용하여 표준 가환 도식을 얻습니다.
$$
\xymatrix{
(f')^{-1}R^qg_*\mathcal{F} \ar[r] \ar[d] &
(f')^{-1}R^qg'_*\mathcal{F}' \ar[d] \\
R^qh_*e^{-1}\mathcal{F} \ar[r] &
R^qh'_*(e')^{-1}\mathcal{F}'
}
$$
$X'$의 아벨 층. $P(T')$를 성질로 설정한다.
\begin{itemize}
\item $\xi'$를 $\xi$의 $(Rh'_*(e')^{-1}\mathcal{F}')_{\overline{x}}$에서의 상이라 하자. 이는 다음 사상의 상에 속하지 않는다.
$(f^{-1}R^qg'_*\mathcal{F}')_{\overline{x}} \to
(R^qh'_*(e')^{-1}\mathcal{F}')_{\overline{x}}$.
\end{itemize}
우리는 주장 (1), (2) 및 (3)를 가정한다.
보조정리 \ref{etale-cohomology-lemma-formal-argument} $P$ 동안 보류
이는 우리의 보조정리를 증명한다.

\medskip\noindent 조건 (1) 의 보조정리 \ref{etale-cohomology-lemma-formal-argument}는 $P$에 대해 유지된다. 왜냐하면 의 에탈 위상 스킴과 스킴의 두꺼워짐은 동일한다. 명제 \ref{etale-cohomology-proposition-topological-invariance}를 참조하라.

\medskip\noindent
$I$가 방향성 집합이고 $T_i$라고 가정한다.
준 컴팩트 및 준 분리의 $I$에 대한 역 시스템이다.
스킴은 아핀 전이 사상을 사용하여 $T$에 걸쳐 있다.
집합 $T' = \lim T_i$. $\mathcal{F}'$ 및 $\mathcal{F}_i$를 나타냅니다.
$\mathcal{F}$을 $T'$로 당김, 즉 \ $T_i$. 고려하다
다이어그램
$$
\vcenter{
\xymatrix{
X \ar[d]_{f'} & Y \ar[l]^h \ar[d]^e & Y_i \ar[l]^{\pi_i'} \ar[d]^{e_i} \\
S & T \ar[l]_g & T_i \ar[l]_{\pi_i}
}
}
\quad\text{및}\quad
\vcenter{
\xymatrix{
X \ar[d]_{f'} & & Y_i \ar[ll]^{h_i = h \circ \pi_i'} \ar[d]^{e_i} \\
S & & T_i \ar[ll]_{g_i = g \circ \pi_i}
}
}
$$
이전 단락에서와 같이. $\mathcal{F}'$이 켜져 있는 것이 분명한다.
$T'$는 $\mathcal{F}_i$에서 $T'$까지의 하락세 중 여극한이다.
그리고 $(e')^{-1}\mathcal{F}'$는 하락세의 여극한입니다
$e_i^{-1}\mathcal{F}_i$부터 $Y'$까지.
보조정리 \ref{etale-cohomology-lemma-relative-colimit-general}에 의해
우리는
$$
R^qh'_*(e')^{-1}\mathcal{F}' = \colim R^qh_{i, *}e_i^{-1}\mathcal{F}_i
\quad\text{및}\quad
(f')^{-1}R^qg'_*\mathcal{F}' = \colim (f')^{-1}R^qg_{i, *}\mathcal{F}_i
$$
다음과 같은 경우$P(T_i)$모두에게 해당됩니다$i$, 그 다음에
$P(T')$이 유지된다. 따라서 조건 (2) 보조정리 \ref{etale-cohomology-lemma-formal-argument}
$P$에 대해 유지된다.

\medskip\noindent
가장 흥미로운 것은조건 (3) 의보조정리 \ref{etale-cohomology-lemma-formal-argument}.
$T'$는 $T$에 대해 준간소하고 준분리된 스킴라고 가정한다.
$P(T')$이 참이 된다.
$Z \subset T'$를 보완 $V \subset T'$가 있는 폐쇄형 하위 구성이라고 가정한다.
준컴팩트. 다이어그램을 고려하십시오
$$
\xymatrix{
Y' \times_{T'} Z \ar[d]_{e_Z} \ar[r]_{i'} &
Y' \ar[d]_{e'} &
Y' \times_{T'} V \ar[l]^{j'} \ar[d]^{e_V} \\
Z \ar[r]^i &
T' &
V \ar[l]_j
}
$$
주입형 맵을 선택하세요. $j^{-1}\mathcal{F}' \to \mathcal{J}$
여기서 $\mathcal{J}$는 $\mathbf{Z}/n\mathbf{Z}$-가군의 단사 층이다.
$V$에 있다. 줄기를 보면 지도가
$$
\mathcal{F}' \to \mathcal{G} = j_*\mathcal{J} \oplus i_*i^{-1}\mathcal{F}'
$$
주입식이다. 따라서 $\xi'$는 0이 아닌 요소에 매핑된다.
\begin{align*}
& \Coker\left(
((f')^{-1}R^qg'_*\mathcal{G})_{\overline{x}}
\to
(R^qh'_*(e')^{-1}\mathcal{G})_{\overline{x}}
\right) = \\
&
\Coker\left(
((f')^{-1}R^qg'_*j_*\mathcal{J})_{\overline{x}}
\to
(R^qh'_*(e')^{-1}j_*\mathcal{J})_{\overline{x}}
\right) \oplus \\
& \Coker\left(
((f')^{-1}R^qg'_*i_*i^{-1}\mathcal{F}')_{\overline{x}}
\to
(R^qh'_*(e')^{-1}i_*i^{-1}\mathcal{F}')_{\overline{x}}
\right)
\end{align*}
부분적으로 (2) 의보조정리 \ref{etale-cohomology-lemma-base-change-q-injective}.
$\xi'$가 두 번째 피합수에서 0으로 매핑되지 않으면
우리는 사용
$$
(f')^{-1}R^qg'_*i_*i^{-1}\mathcal{F}' =
(f')^{-1}R^q(g' \circ i)_*i^{-1}\mathcal{F}'
$$
($Ri_* = i_*$ 때문에
명제 \ref{etale-cohomology-proposition-finite-higher-direct-image-zero}) 및
$$
R^qh'_*(e')^{-1}i_*i^{-1}\mathcal{F} =
R^qh'_*i'_*e_Z^{-1}i^{-1}\mathcal{F} =
R^q(h' \circ i')_*e_Z^{-1}i^{-1}\mathcal{F}'
$$
(첫 번째 평등
보조정리 \ref{etale-cohomology-lemma-finite-pushforward-commutes-with-base-change}
그리고 두 번째 이유는
$Ri'_* = i'_*$ 에 의해
명제 \ref{etale-cohomology-proposition-finite-higher-direct-image-zero})
$P(Z)$이 있음을 알 수 있다.
마지막으로, $\xi'$가 첫 번째 피합수에서 0으로 매핑되지 않는다고 가정한다.
우리는
$$
(e')^{-1}j_*\mathcal{J} = j'_*e_V^{-1}\mathcal{J}
\quad\text{및}\quad
R^aj'_*e_V^{-1}\mathcal{J} = 0, \quad a = 1, \ldots, q - 1
$$
다이어그램에 $BC(f, n, q - 1)$ 적용
$$
\xymatrix{
X \ar[d]_f & Y' \ar[l] \ar[d]_{e'} & Y \ar[l]^{j'} \ar[d]^{e_V} \\
S & T' \ar[l] & V \ar[l]_j
}
$$
그리고 $\mathcal{J}$는 단사라는 사실.
$h' \circ j'$에 대한 상대 Leray 스펙트럼 열 기준
(코호몰로지 사이트, 보조정리 \ref{sites-cohomology-lemma-relative-Leray})
우리는 그것을 추론
$$
R^qh'_*(e')^{-1}j_*\mathcal{J} =
R^qh'_*j'_*e_V^{-1}\mathcal{J}
\longrightarrow
R^q(h' \circ j')_* e_V^{-1}\mathcal{J}
$$
주입식이다. 따라서 $\xi$는 0이 아닌 요소에 매핑된다.
$(R^q(h' \circ j')_* e_V^{-1}\mathcal{J})_{\overline{x}}$.
적용부분(3) 의보조정리 \ref{etale-cohomology-lemma-base-change-q-injective}
주입에 $j^{-1}\mathcal{F}' \to \mathcal{J}$
우리는 $P(V)$가 성립한다고 결론을 내립니다.
\end{proof}

\begin{lemma}
\label{etale-cohomology-lemma-base-change-does-not-hold}
와 함께$f : X \to S$그리고$n$에서와 같이비고 \ref{etale-cohomology-remark-base-change-holds}
일부 $q \geq 1$에 대해 가정해 보자.
$BC(f, n, q - 1)$은 참이지만 $BC(f, n, q)$는 그렇지 않습니다.
그러면 가환 도식이 존재한다.
$$
\xymatrix{
X \ar[d]_f & X' \ar[d] \ar[l] & Y \ar[l]^h \ar[d] \\
S & S' \ar[l] & \Spec(K) \ar[l]
}
$$
두 정사각형 데카르트식을 사용하는 경우, 여기서
\begin{enumerate}
\item $S'$는 유사 관계이고, 정역 및 정규는 대수적으로
닫힌 함수 체,
\item $K$는 대수적으로 닫혀 있으며 $\Spec(K) \to S'$
지배적입니다(즉, $K$는
$S'$의 체 기능)
\end{enumerate}
그리고 정수 $d | n$가 존재한다.
$R^qh_*(\mathbf{Z}/d\mathbf{Z})$는 0이 아닙니다.
\end{lemma}

\noindent 반대로, 보조정리에서 $R^qh_*(\mathbf{Z}/d\mathbf{Z})$이 사라지지 않는다는 것은 $BC(f, n, q)$가 사실이 아니라는 것을 의미한다. 보조정리 \ref{etale-cohomology-lemma-Rf-star-zero-normal-with-alg-closed-function-field}는 $R^q(\Spec(K) \to S')_*\mathbf{Z}/d\mathbf{Z} = 0$를 보여주기 때문이다.

\begin{proof}
먼저 다음과 같이 다이어그램과 $\mathcal{F}$를 선택한다.
보조정리 \ref{etale-cohomology-lemma-base-change-does-not-hold-pre}.
우리는 $S'$가 유사하다고 가정할 수도 있고 실제로 가정합니다(이것은 명백하지만
의심스러운 경우 보조정리의 증명을 참조하라.
보조정리 \ref{etale-cohomology-lemma-base-change-q-integral-top}에 의해
$K$이 대수적으로 닫혀 있다고 가정할 수 있다.
그러면 $\mathcal{F}$는 $\mathbf{Z}/n\mathbf{Z}$-가군에 해당한다.
이러한 가군은 $\mathbf{Z}/d\mathbf{Z}$ 사본의 직합이다.
$d | n$을 변경하기 위해 $\mathcal{F}$이 다음과 같다고 가정할 수 있다.
값 $\mathbf{Z}/d\mathbf{Z}$의 상수이다.
보조정리 \ref{etale-cohomology-lemma-base-change-q-integral-bottom}에 의해
$S'$을 정규화로 대체할 수 있다.
~의$S'$~에$\Spec(K)$마무리하는 것증명.
\end{proof}










\section{매끄러운 기저변환}
\label{etale-cohomology-section-smooth-base-change}

\noindent 이 절에서 매끄러운 기저변환 정리를 증명한다.

\begin{lemma}
\label{etale-cohomology-lemma-smooth-base-change-fields}
$K/k$를 체의 확장으로 설정한다. $X$를 부드러운 아핀 곡선으로 설정한다.
$k$를 유리점 $x \in X(k)$로 오버한다. $\mathcal{F}$를 아벨리언이라고 ​​하자
층~에$\Spec(K)$정수로 소멸$n$에서 반전 가능$k$.
$q > 0$ 그리고
$$
\xi \in H^q(X_K, (X_K \to \Spec(K))^{-1}\mathcal{F})
$$
존재한다
\begin{enumerate}
\item 유한 확장 $K'/K$ 및 $k'/k$와 $k' \subset K'$,
\item 유한 에탈 갈루아 피복 $Z \to X_{k'}$와 그 군 $G$
\end{enumerate}
$G$의 차수는 $n$의 거듭제곱을 나누어서 다음과 같습니다.
$Z \to X_{k'}$은 $x_{k'}$로 분할된다.
$\xi$는 $H^q(Z_{K'}, (Z_{K'} \to \Spec(K))^{-1}\mathcal{F})$에서 죽습니다.
\end{lemma}

\begin{proof}
$q > 1$에 대해 우리는 $\xi$이(가) 죽는다는 것을 알고 있다.
$H^q(X_{\overline{K}}, (X_{\overline{K}} \to \Spec(K))^{-1}\mathcal{F})$
(정리 \ref{etale-cohomology-theorem-vanishing-affine-curves}).
보조정리 \ref{etale-cohomology-lemma-directed-colimit-cohomology}에 의해 우리는
이는 다음과 같은 유한한 확장 $K'/K$이 있음을 의미한다.
$\xi$가 $H^q(X_{K'}, (X_{K'} \to \Spec(K))^{-1}\mathcal{F})$에서 사망한다.
따라서 이 경우에는 $k' = k$ 및 $Z = X$를 사용할 수 있다.

\medskip\noindent
$q = 1$를 가정한다. $\mathcal{F}$는
이산 가군 $M$ 연속 $\text{Gal}_K$ 동작 포함, 참조
보조정리 \ref{etale-cohomology-lemma-equivalence-abelian-sheaves-point}.
$M$는 $n$-비틀림이므로 유한합집합이다.
$\text{Gal}_K$-안정적인 하위 그룹. 따라서
$M$가 $n$에 의해 소멸된 유한 아벨 군인 경우로 축소한다.
보조정리 \ref{etale-cohomology-lemma-colimit}를 참조하라. $K$을 유한 확장으로 대체한 후
우리는 다음의 행동을 가정할 수 있다.$\text{Gal}_K$~에$M$사소한 것이다.
따라서 우리는 $\mathcal{F} = \underline{M}$가 상수라고 가정할 수 있다.
층 값이 유한한 아벨 군 $M$가 $n$에 의해 소멸되었다.

\medskip\noindent $M$를 순환 그룹의 직합으로 쓸 수 있다. 갈루아 그룹이 $k$에서 반전 가능한 순서를 갖는 두 개의 유한 에탈 갈루아 피복은 갈루아 그룹이 $k$에서 반전 가능한 순서를 갖는 세 번째 그룹에 의해 지배될 수 있습니다(기본 그룹, 절 \ref{pione-section-finite-etale-under-galois}). 따라서 $M = \mathbf{Z}/d\mathbf{Z}$ 여기서 $d | n$일 때 보조정리를 증명하는 것으로 충분한다.

\medskip\noindent $M = \mathbf{Z}/d\mathbf{Z}$ 여기서 $d | n$라고 가정한다. 이 경우 $\overline{\xi} = \xi|_{X_{\overline{K}}}$는 $$
H^1(X_{\overline{k}}, \mathbf{Z}/d\mathbf{Z}) =
H^1(X_{\overline{K}}, \mathbf{Z}/d\mathbf{Z})
$$의 요소이다. 정리 \ref{etale-cohomology-theorem-vanishing-affine-curves}를 참조하라. 이 그룹은 $\mathbf{Z}/d\mathbf{Z}$-토서를 분류한다. 사이트, 보조정리 \ref{sites-cohomology-lemma-torsors-h1}에 대한 코호몰로지를 참조하라. $\overline{\xi}$에 해당하는 토르서($X_{\overline{k}, \etale}$에서 층으로 표시)는 유한 에탈 사상 $T \to X_{\overline{k}}$에 전이적인 $\mathbf{Z}/d\mathbf{Z}$ 동작을 부여한다. 올의 $T$에 대한 $x_{\overline{k}}$, 보조정리 \ref{etale-cohomology-lemma-characterize-finite-locally-constant}를 참조하라. 연결 구성요소 $T' \subset T$를 선택합니다($\overline{\xi}$에 $d$ 순서가 있는 경우 $T$는 이미 연결입니다). 그런 다음 $T' \to X_{\overline{k}}$는 Galois 그룹이 하위 그룹 $G \subset \mathbf{Z}/d\mathbf{Z}$인 유한 에탈 Galois 표지입니다(작은 세부 사항 생략). 또한 $\overline{\xi}$ 요소는 $H^1(X_{\overline{k}}, \mathbf{Z}/d\mathbf{Z}) \to
H^1(T', \mathbf{Z}/d\mathbf{Z})$ 맵에서 0으로 매핑된다. 이는 $T$의 정의 성질 중 하나이기 때문이다.

\medskip\noindent
다음으로, 극한 인수를 사용하여 유한 확장 $k'/k$을 선택한다.
$\overline{k}$에 포함되어 $T' \to X_{\overline{k}}$
유한한 에탈 갈루아 커버로 하강함$Z \to X_{k'}$그룹과 함께$G$.
극한, 기본형 \ref{limits-lemma-descend-finite-presentation},
\ref{limits-lemma-descend-finite-finite-presentation} 그리고
\ref{limits-lemma-descend-etale}.
$k'$를 증가시킨 후 $Z$가 $x_{k'}$로 분할된다고 가정할 수 있다.
$\xi$의 이미지는
$H^1(Z_{\overline{K}}, \mathbf{Z}/d\mathbf{Z})$은 구조상 0이다.
따라서 보조정리 \ref{etale-cohomology-lemma-directed-colimit-cohomology}
유한한 하위 확장 $\overline{K}/K'/K$을 찾을 수 있다.
포함하는$k'$그렇게$\xi$사망하다$H^1(Z_{K'}, \mathbf{Z}/d\mathbf{Z})$
이것으로 증명이 끝납니다.
\end{proof}

\begin{theorem}[매끄러운 기저변환]
\label{etale-cohomology-theorem-smooth-base-change}
스킴의 데카르트 다이어그램을 고려해보세요.
$$
\xymatrix{
X \ar[d]_f & Y \ar[l]^h \ar[d]^e \\
S & T \ar[l]_g
}
$$
여기서 $f$는 매끄럽고 $g$는 준컴팩트하고 준분리되어 있다. 그 다음에
$$
f^{-1}R^qg_*\mathcal{F} = R^qh_*e^{-1}\mathcal{F}
$$
모든 $q$ 및 아벨 층 $\mathcal{F}$에 대해
$T_\etale$에서 기하점의 줄기가 모두 비틀림이다.
$S$에 취소할 수 있는 주문이다.
\end{theorem}

\begin{proof}[매끄러운 기저변환]의 첫 번째 증명 이 증명은 매우 길지만 아래에 제공된 두 번째 증명보다 더 직접적입니다(덜 일반적인 이론을 사용함).

\medskip\noindent
정리는 $X_\etale$에 로컬이다. 좀 더 정확하게 말하자면,
$U \to X$ 에탈은 $U \to S$를 $U \to V \to S$로 고려한다.
$V \to S$ 에탈. 그러면 우리는 데카르트 제곱을 생각해 볼 수 있다.
$$
\xymatrix{
U \ar[d]_{f'} & U \times_X Y \ar[l]^{h'} \ar[d]^{e'} \\
V & V \times_S T \ar[l]_{g'}
}
$$
그리고 $\mathcal{F}' = \mathcal{F}|_{V \times_S T}$ 설정
우리는 $f^{-1}R^qg_*\mathcal{F}|_U = (f')^{-1}R^qg'_*\mathcal{F}'$
그리고 $R^qh_*e^{-1}\mathcal{F}|_U = R^qh'_*(e')^{-1}\mathcal{F}'$
(국소화와 사상의 사이트 호환성에 대해 다음과 같이 설명한다.
사이트, 보조정리 \ref{sites-lemma-localize-morphism-strong} 및
그리고
코호몰로지 사이트, 보조정리에
\ref{sites-cohomology-lemma-restrict-direct-image-open}).
따라서 다음과 같이 $X$의 에탈 피복을 생성하는 것으로 충분한다.
$U \to X$ 및 인수분해 $U \to V \to S$
위와 같이 정리가 다음 다이어그램에 대해 유지된다.
$f'$, $h'$, $g'$, $e'$.

\medskip\noindent 매끄러운 사상의 로컬 구조에 따라 사상, 보조정리 \ref{morphisms-lemma-smooth-etale-over-affine-space}를 참조하여 $X$ 및 $S$가 유사하고 $X \to S$ 에탈 사상 $X \to \mathbf{A}^d_S$을 통해 인수분해한다. 직교 다이어그램 $$
\xymatrix{
W \ar[d]_i & Z \ar[l]^j \ar[d]^k \\
X \ar[d]_f & Y \ar[l]^h \ar[d]^e \\
S & T \ar[l]_g
}
$$의 타워가 있고 정리가 아래쪽 및 위쪽 사각형에 대해 유지되는 경우 정리는 외부 직사각형에 대해 유지된다. 이것은 공식적인 것이다. $X \to S$를 구성 $$
X \to \mathbf{A}^{d - 1}_S \to \mathbf{A}^{d - 2}_S \to \ldots \to
\mathbf{A}^1_S \to S
$$로 작성하면 $X$와 $S$가 유사하고 $X \to S$가 상대 차원을 가질 때 정리를 증명하는 것으로 충분하다고 결론을 내립니다. $1$.

\medskip\noindent
모든$n \geq 1$반전 가능$S$, 허락하다$\mathcal{F}[n]$
의 하위 단이 되다절~의$\mathcal{F}$에 의해 전멸$n$. 그 다음에
$\mathcal{F} = \colim \mathcal{F}[n]$ 우리의 가정에 의해
$\mathcal{F}$의 줄기. 함자 $e^{-1}$ 및 $f^{-1}$
여극한이 인접해 있으므로 통근한다. 함자
$R^qh_*$ 및 $R^qg_*$는 여과 여극한으로 출퇴근한다.
보조정리 \ref{etale-cohomology-lemma-relative-colimit}.
따라서 $\mathcal{F}[n]$에 대해 정리를 증명하면 충분한다.
이제부터 정수 $n$를 수정한다.
층 중 $\mathbf{Z}/n\mathbf{Z}$-가군 및
$S$는 $\Spec(\mathbf{Z}[1/n])$에 대한 스킴라고 가정한다.

\medskip\noindent
다음으로 $T$이 아핀인 경우로 축소한다. $g$는 준컴팩트하므로
그리고 준 분리되어 있고 $S$는 유사하며, 스킴 $T$는 준 컴팩트하고
준분리. 따라서 우리는 유도 원리를 사용할 수 있습니다
코호몰로지 중 스킴, 보조정리 \ref{coherent-lemma-induction-principle}.
따라서 $T = W \cup W'$인 경우를 보여주는 것으로 충분한다.
열린 피복이고 정리는 사각형에 대해 유지된다.
$$
\xymatrix{
X \ar[d] & e^{-1}(W) \ar[l]^i \ar[d] \\
S & W \ar[l]_a
}
\quad
\xymatrix{
X \ar[d] & e^{-1}(W') \ar[l]^j \ar[d] \\
S & W' \ar[l]_b
}
\quad
\xymatrix{
X \ar[d] & e^{-1}(W \cap W') \ar[l]^-k \ar[d] \\
S & W \cap W' \ar[l]_c
}
$$
그러면 정리는 원래 다이어그램에 적용된다. 이것을 보기 위해 우리는 다음을 고려한다.
도표
$$
\xymatrix{
f^{-1}R^{q - 1}c_*\mathcal{F}|_{W \cap W'} \ar[d]^{\cong} \ar[r] &
f^{-1}R^qg_*\mathcal{F} \ar[d] \ar[r] &
f^{-1}R^qa_*\mathcal{F}|_W \oplus
f^{-1}R^qb_*\mathcal{F}|_{W'} \ar[d]_{\cong} \\
R^qk_*e^{-1}\mathcal{F}|_{e^{-1}(W \cap W')} \ar[r] &
R^qh_*e^{-1}\mathcal{F} \ar[r] &
R^qi_*e^{-1}\mathcal{F}|_{e^{-1}(W)} \oplus
R^qj_*e^{-1}\mathcal{F}|_{e^{-1}(W')}
}
$$
그 행은 다음의 긴 완전열이다.
보조정리 \ref{etale-cohomology-lemma-relative-mayer-vietoris}.
따라서 $5$-보조정리는 원하는 결론을 제공한다.

\medskip\noindent 요약하면, $S$, $X$, $T$ 및 $Y$ 아핀을 가정할 수 있으며, $\mathcal{F}$는 $n$ 비틀림이다. $X \to S$는 상대 차원 $1$의 부드러움이고, $S$는 $\mathbf{Z}[1/n]$에 대한 스킴이다. $q$에 대한 귀납법으로 정리를 증명하겠습니다. 기본 사례 $q = 0$는 보조정리 \ref{etale-cohomology-lemma-fppf-reduced-fibres-base-change-f-star}에 의해 처리된다. $q > 0$ 및 정리가 모든 더 작은 각도에 대해 유지된다고 가정한다. 짧은 완전열 $0 \to \mathcal{F} \to \mathcal{I} \to \mathcal{Q} \to 0$를 선택한다. 여기서 $\mathcal{I}$는 $\mathbf{Z}/n\mathbf{Z}$-가군의 분사형 층이다. 정확한 행을 포함하는 유도된 다이어그램 $$
\xymatrix{
f^{-1}R^{q - 1}g_*\mathcal{I} \ar[d]_{\cong} \ar[r] &
f^{-1}R^{q - 1}g_*\mathcal{Q} \ar[d]_{\cong} \ar[r] &
f^{-1}R^qg_*\mathcal{F} \ar[d] \ar[r] &
0 \ar[d] \\
R^{q - 1}h_*e^{-1}\mathcal{I} \ar[r] &
R^{q - 1}h_*e^{-1}\mathcal{Q} \ar[r] &
R^qh_*e^{-1}\mathcal{F} \ar[r] &
R^qh_*e^{-1}\mathcal{I}
}
$$을 고려해보세요. $\mathcal{I}$는 분사형이므로 오른쪽 상단에 0이 있다. 왼쪽 두 개의 수직 화살표는 귀납 가설에 의한 동형사상이다. 따라서 $R^qh_*e^{-1}\mathcal{I} = 0$을 증명하면 충분한다.

\medskip\noindent
$S = \Spec(A)$ 및 $T = \Spec(B)$라고 쓰고 사상라고 말한다.
$T \to S$는 환 맵 $A \to B$에 의해 제공된다. 우리는 쓸 수 있다
$A \to B = \colim_{i \in I} (A_i \to B_i)$ 필터링됨
$\mathbf{Z}[1/n]$에 대한 유한 유형의 환 맵의 여극한
(대수학, 보조정리 \ref{algebra-lemma-limit-no-condition} 참조).
$i \in I$의 경우 집합 $S_i = \Spec(A_i)$ 및 $T_i = \Spec(B_i)$이다.
$i$가 충분히 큰 경우 매끄러운 사상 $X_i \to S_i$를 찾을 수 있다.
친척의차원 $1$그렇게$X = X_i \times_{S_i} S$, 보다
극한, 기본형 \ref{limits-lemma-descend-finite-presentation},
\ref{limits-lemma-descend-smooth} 그리고
\ref{limits-lemma-descend-dimension-d}. 두자 $Y_i = X_i \times_{S_i} T_i$
사각형을 얻으려면
$$
\xymatrix{
X_i \ar[d]_{f_i} & Y_i \ar[l]^{h_i} \ar[d]^{e_i} \\
S_i & T_i \ar[l]_{g_i}
}
$$
$\mathcal{I}_i = (T \to T_i)_*\mathcal{I}$를 확인하세요.
단사이다 층 의 $\mathbf{Z}/n\mathbf{Z}$-가군 위의
$T_i$, 사이트, 보조정리의 코호몰로지를 참조하라.
\ref{sites-cohomology-lemma-pushforward-injective-flat}.
우리는 $\mathcal{I} = \colim (T \to T_i)^{-1}\mathcal{I}_i$
보조정리 \ref{etale-cohomology-lemma-linus-hamann} 기준. $e$로 뒤로 당기면 다음과 같은 결과가 나옵니다.
$e^{-1}\mathcal{I} = \colim (Y \to Y_i)^{-1}e_i^{-1}\mathcal{I}_i$.
보조정리 \ref{etale-cohomology-lemma-relative-colimit-general}에 의해 시스템에 적용됨
사상 $Y_i \to X_i$와 극한 $Y \to X$가 있다.
$$
R^qh_*e^{-1}\mathcal{I} =
\colim (X \to X_i)^{-1} R^qh_{i, *} e_i^{-1}\mathcal{I}_i
$$
이는 $T$ 및 $S$가 유한한 경우로 축소된다.
$\mathbf{Z}[1/n]$ 위에 입력하세요.

\medskip\noindent
요약하면, 정리가 유지되는 정수 $q \geq 1$가 있다.
도 단위 $< q$, 유한 유형의 스킴 $S$ 및 $T$ 아핀
$\mathbf{Z}[1/n]$ 위에 입력하면
$X \to S$ 상대 차원 $1$를 $X$ 아핀으로 매끄럽게 하고,
$\mathcal{I}$는 $\mathbf{Z}/n\mathbf{Z}$-가군의 분사형 층이다.
그리고 우리는 $R^qh_*e^{-1}\mathcal{I} = 0$를 보여주어야 한다.
$\dim(T)$에 대한 유도를 통해 이 작업을 수행하겠습니다.

\medskip\noindent 기본 사례는 $T = \emptyset$, 즉 $\dim(T) < 0$이다. 이것이 마음에 들지 않으면 기본 케이스로 $\dim(T) = 0$ 케이스를 사용할 수 있다. 이 경우 $T \to S$는 유한합니다(실제로 $T \to \Spec(\mathbf{Z}[1/n])$도 대상이 Jacobson이므로 유한한다. 자세한 내용은 생략됨). 따라서 $h$도 유한하므로 소멸 고차 직상을 갖습니다(아래 참조 참조).

\medskip\noindent
$\dim(T) = d \geq 0$를 가정하면 $T$가 발생하는 모든 상황에 대한 결과를 알 수 있다.
더 낮았다차원. 선택하다$U$아핀 및 에탈 오버$X$그리고절 $\xi$
~의$R^qh_*q^{-1}\mathcal{I}$~ 위에$U$. 우리는 보여줘야 해
$\xi$는 0이다. 물론 $X$를 $U$로 바꿀 수도 있다.
(그리고 그에 따라 $Y$ 에 의해 $U \times_X Y$)
그리고 $\xi \in H^0(X, R^qh_*e^{-1}\mathcal{I})$라고 가정한다.
게다가 $R^qh_*e^{-1}\mathcal{I}$는 층이므로,
$\xi$이 $X$에서 로컬로 0이라는 것을 증명하는 것으로 충분한다.
따라서 $X$를 에탈 피복의 멤버로 대체할 수 있다.
특히 보조정리 \ref{etale-cohomology-lemma-higher-direct-images}를 사용하면
$\xi$는 요소의 이미지라고 가정할 수 있다.
$\tilde \xi \in H^q(Y, e^{-1}\mathcal{I})$.
$\tilde \xi$ 측면에서 우리의 임무는 다음을 보여주는 것이다.
$\tilde \xi$가 $H^q(U_i \times_X Y, e^{-1}\mathcal{I})$에서 사망한다.
일부 에탈 피복 $\{U_i \to X\}$.

\medskip\noindent
사상, 보조정리 \ref{more-morphisms-lemma-cover-smooth-by-special}에 대한 자세한 내용
$X \to S$가 $X \to V \to S$로 간주된다고 가정할 수 있다.
여기서 $V \to S$는 에탈이고 $X \to V$는 매끄러운 사상이다.
아핀 스킴의 상대 차원 $1$ 중 절이 있고
기하학적으로 연결 올을 갖습니다. 그것을 관찰하십시오
$\dim(V \times_S T) \leq \dim(T) = d$
대수학에 대한 추가 정보, 보조정리의 예
\ref{more-algebra-lemma-dimension-etale-extension}.
따라서 $S$을 $V$로, $T$를 $V \times_S T$로 바꿀 수 있다.
(정확히 증명의 첫 번째 문단의 논의와 같습니다).
따라서 우리는 $X \to S$가 상대적인 차원 $1$에 대해 부드럽다고 가정할 수 있다.
기하학적으로 연결 올이고 절 $\sigma : S \to X$을 갖습니다.

\medskip\noindent
$\pi : T' \to T$를 유한 전사 사상으로 둡니다.
아래에서 $\dim(T') \leq \dim(T) = d$를 사용하겠습니다.
대수학, 보조정리 \ref{algebra-lemma-integral-dim-up}.
주입형 맵을 선택하세요. $\pi^{-1}\mathcal{I} \to \mathcal{I}'$
$\mathbf{Z}/n\mathbf{Z}$-가군의 분사형 층으로.
그러면 $\mathcal{I} \to \pi_*\mathcal{I}'$ 는 단사입니다
따라서 분할이 있습니다($\mathcal{I}$는 분사형이므로).
층 중 $\mathbf{Z}/n\mathbf{Z}$-가군).
$\pi' : Y' = Y \times_T T' \to Y$ $\pi$의 기본 변경을 나타냅니다.
및 $e' : Y' \to T'$ $e$의 기본 변경 사항이다. 그림
$$
\xymatrix{
X \ar[d]_f & Y \ar[l]^h \ar[d]^e & Y' \ar[l]^{\pi'} \ar[d]^{e'} \\
S & T \ar[l]_g & T' \ar[l]_\pi
}
$$
명제 \ref{etale-cohomology-proposition-finite-higher-direct-image-zero} 및
보조정리 \ref{etale-cohomology-lemma-finite-pushforward-commutes-with-base-change} 우리는
$R\pi'_*(e')^{-1}\mathcal{I}' = e^{-1}\pi_*\mathcal{I}'$.
따라서 Leray에 의해 스펙트럼 열
(코호몰로지 사이트, 보조정리 \ref{sites-cohomology-lemma-Leray})
우리는
$$
H^q(Y', (e')^{-1}\mathcal{I}') =
H^q(Y, e^{-1}\pi_*\mathcal{I}') \supset H^q(Y, e^{-1}\mathcal{I})
$$
그리고 이는 $U \to X$ 에탈에 의한 염기 변경 후에도 그대로 유지된다.
따라서 $T$을 $T'$로, $\mathcal{I}$를 $\mathcal{I}'$로 바꿀 수 있다.
및 $\tilde \xi$는 $H^q(Y', (e')^{-1}\mathcal{I}')$의 이미지로 표시된다.

\medskip\noindent
$T \to S' \to S$ 인수분해가 있다고 가정한다.
$\pi : S' \to S$은 유한한다.
$X' = S' \times_S X$을 설정하면 유도된 다이어그램을 고려할 수 있다.
$$
\xymatrix{
X \ar[d]_f & X' \ar[l]^{\pi'} \ar[d]^{f'} & Y \ar[l]^{h'} \ar[d]^e \\
S & S' \ar[l]_\pi & T \ar[l]_g
}
$$
$\pi'$에는 소멸 고차 직상이 있으므로 다음을 알 수 있다.
$R^qh_*e^{-1}\mathcal{I} = \pi'_*R^qh'_*e^{-1}\mathcal{I}$
르레이의 스펙트럼 열. 따라서
$H^0(X, R^qh_*e^{-1}\mathcal{I}) = H^0(X', R^qh'_*e^{-1}\mathcal{I})$.
따라서 $\xi$는 0이다. 필요충분조건은 해당 절
$R^qh'_*e^{-1}\mathcal{I}$는 0이다.\footnote{이것은
단계는 다른 방식으로도 볼 수 있다. 즉, 우리는 그것을 보여주어야 한다.
에탈 피복 $\{U_i \to X\}$이 있다.
$\tilde \xi$가 $H^q(U_i \times_X Y, e^{-1}\mathcal{I})$에서 사망한다.
그러나 에탈이 있음을 증명하면 피복
$\{U'_j \to X'\}$ 그래서 $\tilde \xi$가 죽습니다.
$H^q(U'_i \times_{X'} Y, e^{-1}\mathcal{I})$, 그런 다음 성질 (B)
에 대하여 $X' \to X$ (보조정리 \ref{etale-cohomology-lemma-finite-B})에 에탈 피복이 존재한다.
$\{U_i \to X\}$ $U_i \times_X X'$는 분리된 합집합이다.
~의스킴~ 위에$X'$각각의 요인을 통해$U'_j$어떤 사람들에게는$j$.
따라서 우리는 $\tilde \xi$가 $H^q(U_i \times_X Y, e^{-1}\mathcal{I})$에서 죽는 것을 볼 수 있다.
원하는 대로.}.
따라서 $S$을 $S'$로, $X$를 $X'$로 바꿀 수 있다.
$\sigma : S \to X$ 베이스가 $\sigma' : S' \to X'$로 변경되는 것을 확인하세요.
따라서 이 교체 후에도 $X \to S$는 여전히 사실이다.
절 $\sigma$ 및 기하학적으로 연결 올을 갖습니다.

\medskip\noindent
$S$ 및 $T$는 Nagata 스킴임을 사용한다.
대수학, 명제 \ref{algebra-proposition-ubiquity-nagata}
보장하는 것
다양한 정규화는 유한한다.
사상, 기본형 \ref{morphisms-lemma-nagata-normalization-finite} 및
\ref{morphisms-lemma-nagata-normalization}.
특히, 먼저 정규화를 통해 $T$을 대체할 수 있다.
그런 다음 교체$S$정규화에 의해$S$~에$T$.
그러면 $T \to S$는 지배적인 결합의 분리된 결합이다.
사상 정역 정규 스킴, 참조
사상, 보조정리 \ref{morphisms-lemma-normal-normalization}.
분명히 우리는 하나의 연결 구성요소를 주장할 수 있다.
그러므로 우리는 $T \to S$이
정역 정규 스킴의 지배적인 사상.

\medskip\noindent
$s \in S$ 및 $t \in T$를 일반 포인트로 설정한다. 에 의해
보조정리 \ref{etale-cohomology-lemma-smooth-base-change-fields} 유한한 존재 체
확장$K/\kappa(t)$그리고$k/\kappa(s)$그렇게$k$에 포함되어 있습니다$K$
그리고 유한 에탈 갈루아 피복 $Z \to X_k$와 갈루아 그룹 $G$
$n$의 거듭제곱을 $\sigma(\Spec(k))$로 나누는 순서
$\tilde \xi$는 $H^q(Z_K, e^{-1}\mathcal{I}|_{Z_K})$에서 0으로 매핑된다.
허락하다$T' \to T$의 정규화가 되다$T$~에$\Spec(K)$
그리고 보자$S' \to S$의 정규화가 되다$S$~에$\Spec(k)$.
그런 다음 가환 도식을 얻습니다.
$$
\xymatrix{
S' \ar[d] & T' \ar[l] \ar[d] \\
S & T \ar[l]
}
$$
수직 화살표는 유한한다. 위에 주어진 주장에 의해 우리는
$S$ 및 $T$를 $S'$ 및 $T'$로 대체할 수 있다.
$X$ 에 의해 $X \times_S S'$ 및 $Y$ 에 의해 $Y \times_T T'$). 이번 교체 이후
우리는 유한한 에탈 갈루아 피복
$Z \to X_s$ 일반 올의 $X \to S$
갈루아 그룹과 함께$G$권력을 나누는 질서의$n$
갈라지다$\sigma(s)$그렇게$\tilde \xi$0으로 매핑
$H^q(Z_t, (Z_t \to Y)^{-1}e^{-1}\mathcal{I})$.
여기 $Z_t = Z \times_S t = Z \times_s t = Z \times_{X_s} Y_t$.
$n$는 $S$에서 반전 가능하므로,
기초 그룹별, 보조정리 \ref{pione-lemma-extend-covering}
우리는 유한한 에탈을 찾을 수 있습니다사상 $U \to X$누구의 제한$X_s$
$Z$이다.

\medskip\noindent 이 시점에서 $X$를 $U$로 바꾸고 $Y$를 $U \times_X Y$로 바꿉니다. 이 교체 후에는 $X \to S$의 올이 기하학적으로 연결인 경우가 더 이상 없을 수 있습니다(여전히 절이 있지만 우리는 이것을 사용하지 않을 것입니다). 그러나 우리가 얻는 것은 이 교체 후 $\tilde \xi$가 $H^q(Y_t, e^{-1}\mathcal{I})$에서 0으로 매핑된다는 것이다. 즉, $\tilde \xi$는 $Y \to T$의 일반 올에서 0으로 제한된다.

\medskip\noindent
$t$는 $T$의 체 함수의 스펙트럼이다. 즉,
스킴 $t$는 $T$의 비어 있지 않은 아핀 개방형 하위 체계의 극한이다.
보조정리 \ref{etale-cohomology-lemma-directed-colimit-cohomology}에 의해 우리는 존재한다고 결론을 내립니다.
비어 있지 않은 개방형 하위 구성표$V \subset T$그렇게$\tilde \xi$0으로 매핑
$H^q(Y \times_T V, e^{-1}\mathcal{I}|_{Y \times_T V})$.

\medskip\noindent
$Z = T \setminus V$을 표시한다. 다이어그램을 고려하십시오
$$
\xymatrix{
Y \times_T Z \ar[d]_{e_Z} \ar[r]_{i'} &
Y \ar[d]_e &
Y \times_T V \ar[l]^{j'} \ar[d]^{e_V} \\
Z \ar[r]^i &
T &
V \ar[l]_j
}
$$
주사제 선택 $i^{-1}\mathcal{I} \to \mathcal{I}'$
주사제로층~의$\mathbf{Z}/n\mathbf{Z}$-가군~에$Z$.
줄기를 보면 지도가
$$
\mathcal{I} \to j_*\mathcal{I}|_V \oplus i_*\mathcal{I}'
$$
는 단사이므로 $\mathcal{I}$는 단사 층으로 분할된다.
$\mathbf{Z}/n\mathbf{Z}$-가군. 따라서 다음을 보여주는 것으로 충분한다.
$\tilde \xi$는 0으로 매핑된다.
$$
H^q(Y, e^{-1}j_*\mathcal{I}|_V) \oplus
H^q(Y, e^{-1}i_*\mathcal{I}')
$$
적어도 $X$을 에탈 피복의 멤버로 바꾼 후에는 말이죠.
그것을 관찰하십시오
$$
e^{-1}j_*\mathcal{I}|_V = j'_*e_V^{-1}\mathcal{I}|_V,\quad
e^{-1}i_*\mathcal{I}' = i'_*e_Z^{-1}\mathcal{I}'
$$
$q$에 대한 귀납 가설을 통해 우리는 다음을 알 수 있다.
$$
R^aj'_*e_V^{-1}\mathcal{I}|_V = 0, \quad a = 1, \ldots, q - 1
$$
$j'$ 및 소멸에 대한 Leray 스펙트럼 열에 의해
그 위에는 다음과 같다
$$
H^q(Y, j'_*(e_V^{-1}\mathcal{I}|_V))
\longrightarrow
H^q(Y \times_T V, e_V^{-1}\mathcal{I}_V) =
H^q(Y \times_T V, e^{-1}\mathcal{I}|_{Y \times_T V})
$$
주입식이다. 따라서 $\tilde \xi$ 이미지의 소멸
위의 첫 번째 요약에서 $\tilde \xi$가 사라진다는 것을 알고 있기 때문이다.
$H^q(Y \times_T V, e^{-1}\mathcal{I}|_{Y \times_T V})$에서.
$\dim(Z) < \dim(T) = d$ 유도에 의해 $\tilde \xi$의 이미지가 생성되었기 때문에
두 번째 소환에서
$$
H^q(Y, e^{-1}i_*\mathcal{I}') =
H^q(Y, i'_*e_Z^{-1}\mathcal{I}') =
H^q(Y \times_T Z, e_Z^{-1}\mathcal{I}')
$$
$X$을 적합한 에탈 피복의 멤버로 교체한 후 사망한다.
이로써 매끄러운 기저변환 정리의 증명이 종료된다.
\end{proof}

\begin{proof}[매끄러운 기저변환]의 두 번째 증명 이 증명은 긴 첫 번째 증명과 동일한다. 여러 보조정리에 사용된 인수를 분리했다는 점에서만 더 짧습니다.

\medskip\noindent $q = 0$의 경우는 보조정리 \ref{etale-cohomology-lemma-fppf-reduced-fibres-base-change-f-star}이다. 따라서 우리는 $q > 0$라고 가정할 수 있으며 그 결과는 모든 작은 각도에 대해 참이다.

\medskip\noindent
모든$n \geq 1$반전 가능$S$, 허락하다$\mathcal{F}[n]$
의 하위 단이 되다절~의$\mathcal{F}$에 의해 전멸$n$. 그 다음에
$\mathcal{F} = \colim \mathcal{F}[n]$ 우리의 가정에 의해
$\mathcal{F}$의 줄기. 함자 $e^{-1}$ 및 $f^{-1}$
여극한이 인접해 있으므로 통근한다. 함자
$R^qh_*$ 및 $R^qg_*$는 여과 여극한으로 출퇴근한다.
보조정리 \ref{etale-cohomology-lemma-relative-colimit}.
따라서 $\mathcal{F}[n]$에 대해 정리를 증명하면 충분한다.
이제부터 정수를 수정한다.$n$반전 가능$S$그리고 우리는 함께 일해요
층 중 $\mathbf{Z}/n\mathbf{Z}$-가군.

\medskip\noindent
에 의해 보조정리 \ref{etale-cohomology-lemma-base-change-local} 질문은 에탈 로컬이다.
$X$ 및 $S$에 있다. 매끄러운 사상의 로컬 구조를 참조하라.
사상, 보조정리 \ref{morphisms-lemma-smooth-etale-over-affine-space},
$X$ 및 $S$가 유사하고 $X \to S$라고 가정할 수 있다.
에탈 사상 $X \to \mathbf{A}^d_S$을 통해 요인을 계산한다.
$X \to S$를 컴포지션으로 작성
$$
X \to \mathbf{A}^{d - 1}_S \to \mathbf{A}^{d - 2}_S \to \ldots \to
\mathbf{A}^1_S \to S
$$
보조정리 \ref{etale-cohomology-lemma-base-change-compose}로 결론을 내립니다.
$X$ 및 $S$일 때 정리를 증명하는 것으로 충분한다.
는 유사하며 $X \to S$에는 상대적인 차원 $1$이 있다.

\medskip\noindent
보조정리 \ref{etale-cohomology-lemma-base-change-does-not-hold}에 의해 다음을 보여주는 것으로 충분한다.
$R^qh_*\mathbf{Z}/d\mathbf{Z} = 0$ 에 대하여 $d | n$
데카르트 다이어그램
$$
\xymatrix{
X \ar[d] & Y \ar[d] \ar[l]^h \\
S & \Spec(K) \ar[l]
}
$$
여기서 $X \to S$는 상대 차원 $1$와 유사하고 매끄럽습니다.
$S$는 정규의 스펙트럼이다.
도메인$A$대수적으로 닫힌 분수체 $L$,
그리고 $K/L$는 대수적으로 닫힌 체의 확장이다.

\medskip\noindent
$R^qh_*\mathbf{Z}/d\mathbf{Z}$는 층과 연관되어 있음을 기억하세요.
전층으로
$$
U \longmapsto
H^q(U \times_X Y, \mathbf{Z}/d\mathbf{Z}) =
H^q(U \times_S \Spec(K), \mathbf{Z}/d\mathbf{Z})
$$
$X_\etale$ (보조정리 \ref{etale-cohomology-lemma-higher-direct-images})에 있다.
따라서 다음을 표시하는 것으로 충분합니다: 주어진 $U$ 및
$\xi \in H^q(U \times_S \Spec(K), \mathbf{Z}/d\mathbf{Z})$
에탈이 존재한다 피복 $\{U_i \to U\}$
$\xi$는 $H^q(U_i \times_S \Spec(K), \mathbf{Z}/d\mathbf{Z})$에서 죽습니다.

\medskip\noindent
물론 우리는 $U$ 아핀을 취할 수도 있다. 그런 다음 $U \times_S \Spec(K)$
는 $K$에 대한 (부드러운) 아핀 곡선이므로 소멸이 된다.
$q > 1$에 대해 정리 \ref{etale-cohomology-theorem-vanishing-affine-curves}에 의해.

\medskip\noindent
최종 사례: $q = 1$. $U$를 다음 멤버로 대체할 수 있다.
에탈 피복 사상, 보조정리에 대한 자세한 내용과 같습니다.
\ref{more-morphisms-lemma-cover-smooth-by-special}.
그러면 $U \to S$는 $U \to V \to S$로 인수분해된다.
$U \to V$는 기하학적으로 연결 올을 가지며,
$U$, $V$는 아핀이고, $V \to S$는 에탈이고,
절 $\sigma : V \to U$. 에 의해
보조정리 \ref{etale-cohomology-lemma-normal-scheme-with-alg-closed-function-field}
우리는 $V$가 다음의 (유한) 분리 결합과 동형임을 알 수 있다.
(아핀) $S$의 하위 구성을 엽니다.
분명히 $S$을 다음 중 하나로 바꿀 수 있다.
및 $X$는 $U$의 해당 구성요소에 의해 결정된다.
따라서 우리는 $X \to S$가 기하학적으로 연결을 갖는다고 가정할 수 있다.
올에는 절 $\sigma$이 있으며,
$\xi \in H^1(Y, \mathbf{Z}/d\mathbf{Z})$.
$K$ 및 $L$는 대수적으로 닫혀 있으므로 다음과 같습니다.
$$
H^1(X_L, \mathbf{Z}/d\mathbf{Z}) =
H^1(Y, \mathbf{Z}/d\mathbf{Z})
$$
보조정리 \ref{etale-cohomology-lemma-base-change-dim-1-separably-closed}를 참조하라.
따라서 유한한 에탈 갈루아 피복이 있다.
$Z \to X_L$ 갈루아 그룹 $G \subset \mathbf{Z}/d\mathbf{Z}$
$\xi$를 전멸시킵니다.
이는 성명서를 보면 알 수 있다.
증명 중 보조정리 \ref{etale-cohomology-lemma-smooth-base-change-fields}
또는 $\xi$에 해당하는 것을 직접 사용하여
$\mathbf{Z}/d\mathbf{Z}$-$X_L$ 위에 토서.
마지막으로,
기본 그룹, 보조정리 \ref{pione-lemma-extend-covering-general}
우리는 (필수적으로 전사)를 찾습니다
유한 에탈 사상 $X' \to X$
$X_L$에 대한 제한은 $Z \to X_L$이다.
$\xi$가 $X'_K$에서 죽으므로 증명이 종료된다.
\end{proof}

\noindent 매끄러운 기저변환 정리의 다음과 같은 즉각적인 결과는 실제로 자주 사용되는 것이다.

\begin{lemma}
\label{etale-cohomology-lemma-smooth-base-change-general}
$S$를 스킴으로 설정한다. $S' = \lim S_i$를 유향 역원이라고 하자
극한 의 스킴 $S_i$ 아핀 전환을 사용하여 $S$를 부드럽게 처리한다.
사상. $f : X \to S$를 준컴팩트하고 준분리되도록 하세요.
그리고 올 정사각형을 형성하세요
$$
\xymatrix{
X' \ar[d]_{f'} \ar[r]_{g'} & X \ar[d]^f \\
S' \ar[r]^g & S
}
$$
그 다음에
$$
g^{-1}Rf_*E = R(f')_*(g')^{-1}E
$$
누구에게나$E \in D^+(X_\etale)$누구의코호몰로지 층 $H^q(E)$
$S$에서 반전 가능한 순서의 비틀림인 줄기를 갖습니다.
\end{lemma}

\begin{proof} 스펙트럼 열 $$
E_2^{p, q} = R^pf_*H^q(E)
\quad\text{그리고}\quad
{E'}_2^{p, q} = R^pf'_*H^q((g')^{-1}E) = R^pf'_*(g')^{-1}H^q(E)
$$가 $R^nf_*E$와 $R^nf'_*(g')^{-1}E$로 수렴되는 것을 생각해 보세요. 이 스펙트럼 열은 유도 범주, 보조정리 \ref{derived-lemma-two-ss-complex-functor}로 구성된다. 매끄러운 기저변환 정리 (정리 \ref{etale-cohomology-theorem-smooth-base-change})를 보조정리 \ref{etale-cohomology-lemma-base-change-Rf-star-colim}와 결합하면 $$
g^{-1}R^pf_*H^q(E) = R^p(f')_*(g')^{-1}H^q(E)
$$를 볼 수 있다. 위의 모든 항목을 결합하면 보조정리를 얻습니다. \end{proof}












\section{매끄러운 기저변환} \label{etale-cohomology-section-applications-smooth-base-change}의 응용

\noindent 이 절에서는 매끄러운 기저변환 정리의 다소 즉각적인 결과에 대해 논의한다.

\begin{lemma} \label{etale-cohomology-lemma-base-change-field-extension} $L/K$를 체의 확장으로 둡니다. $g : T \to S$를 $K$에 걸쳐 스킴의 준컴팩트 및 준분리 사상으로 설정한다. $g_L : T_L \to S_L$는 $g$에서 $\Spec(L)$로의 기본 변경을 나타냅니다. $E \in D^+(T_\etale)$에는 코호몰로지 층이 있으며, 그 줄기는 $K$에서 반전 가능한 순서의 비틀림이다. $E_L$는 $E$에서 $(T_L)_\etale$로의 당김이다. 그러면 $Rg_{L, *}E_L$는 $Rg_*E$에서 $S_L$로의 철회이다. \end{lemma}

\begin{proof}
$L/K$가 분리 가능한 경우 $L$는 부드러운 여과 여극한이다.
$K$-대수학, 참조
대수학, 보조정리 \ref{algebra-lemma-colimit-syntomic}.
따라서 이 경우 보조정리는
보조정리 \ref{etale-cohomology-lemma-smooth-base-change-general}.
일반적인 경우에는 $K'$ 및 $L'$를 완벽한 클로저로 둡니다.
(대수학, 정의 \ref{algebra-definition-perfection})
$K$ 및 $L$. 그런 다음 $\Spec(K') \to \Spec(K)$ 그리고
$\Spec(L') \to \Spec(L)$는 보편적인 동형이다.
$K'/K$와 $L'/L$는 완전히 분리될 수 없기 때문이다.
(대수학, 보조정리 \ref{algebra-lemma-p-ring-map} 참조).
따라서 우리는 $(T_{K'})_\etale = T_\etale$,
$(S_{K'})_\etale = S_\etale$,
$(T_{L'})_\etale = (T_L)\etale$ 그리고
$(S_{L'})_\etale = (S_L)_\etale$ 에 의해
에탈 코호몰로지의 위상학적 불변성 참조
명제 \ref{etale-cohomology-proposition-topological-invariance}.
이는 보조정리를 체의 경우로 축소한다.
확장자 $L'/K'$ (정의로 분리 가능)
완벽함 체, 대수학, 정의 \ref{algebra-definition-perfect})을 참조하라.
\end{proof}

\begin{lemma} \label{etale-cohomology-lemma-smooth-base-change-separably-closed} $K/k$는 분리 가능하게 닫힌 체의 확장이라고 가정한다. $X$를 $k$에 걸쳐 준간소하고 준분리된 스킴으로 설정한다. $E \in D^+(X_\etale)$에는 코호몰로지 층이 있고 그 줄기는 $k$에서 반전 가능한 순서의 비틀림이다. 그러면 \begin{enumerate} \item 지도 $H^q_\etale(X, E) \to H^q_\etale(X_K, E|_{X_K})$는 동형사상이고, \item $E \to R(X_K \to X)_*E|_{X_K}$는 동형사상이다. \end{enumerate} \end{lemma}

\begin{proof}
증명 (1). 먼저 $\overline{k}$ 및 $\overline{K}$
대수적 종결이 되다
$k$ 및 $K$. 사상 $\Spec(\overline{k}) \to \Spec(k)$ 및
$\Spec(\overline{K}) \to \Spec(K)$는 보편적인 동형이다.
$\overline{k}/k$와 $\overline{K}/K$는 완전히 분리될 수 없기 때문이다.
(대수학, 보조정리 \ref{algebra-lemma-p-ring-map} 참조).
따라서 $H^q_\etale(X, \mathcal{F}) =
H^q_\etale(X_{\overline{k}}, \mathcal{F}_{X_{\overline{k}}})$ 기준
에탈 코호몰로지의 위상학적 불변성 참조
명제 \ref{etale-cohomology-proposition-topological-invariance}.
$X_K$ 및 $X_{\overline{K}}$도 마찬가지이다.
따라서 $k$ 및 $K$이 대수적으로 닫혀 있다고 가정할 수 있다.
이 경우 $K$는 부드러운 $k$-대수의 극한이다.
대수학, 보조정리 \ref{algebra-lemma-colimit-syntomic}.
우리는 보조정리가 다음의 특별한 경우라고 결론을 내립니다.
정리 \ref{etale-cohomology-theorem-smooth-base-change}
보조정리 \ref{etale-cohomology-lemma-smooth-base-change-general}.

\medskip\noindent
증명의 (2). 모든 준컴팩트 및 준분리형에 적합$U$~에$X_\etale$
위의 내용은 지도의 제한 사항을 보여줍니다.
$E \to R(X_K \to X)_*E|_{X_K}$는 코호몰로지에서 동형사상을 결정한다.
$X_\etale$의 모든 객체는 $U$에 의해 에탈 피복을 가지므로
이것은 원하는 진술을 증명한다.
\end{proof}

\begin{lemma} \label{etale-cohomology-lemma-base-change-does-not-hold-post} $f : X \to S$ 및 $n$를 비고 \ref{etale-cohomology-remark-base-change-holds}에서와 같이 사용하면 $n$가 $S$에서 반전 가능하다고 가정한다. $q \geq 1$ $BC(f, n, q - 1)$은 사실이지만 $BC(f, n, q)$는 사실이 아닙니다. 그런 다음 가환 도식 $$
\xymatrix{
X \ar[d]_f & X' \ar[d] \ar[l] & Y \ar[l]^h \ar[d] \\
S & S' \ar[l] & \Spec(K) \ar[l]
}
$$가 두 사각형 데카르트를 포함하고, 여기서 $S'$는 아핀이고, 정역 및 정규는 대수적으로 닫힌 함수 체 $K$가 있으며 정수가 존재한다. $d | n$ $R^qh_*(\mathbf{Z}/d\mathbf{Z})$는 0이 아닙니다. \end{lemma}

\begin{proof}
먼저 다음과 같이 다이어그램과 $\mathcal{F}$를 선택한다.
보조정리 \ref{etale-cohomology-lemma-base-change-does-not-hold}.
우리는 $S'$가 유사하다고 가정할 수도 있고 실제로 가정합니다(이것은 명백하지만
의심스러운 경우 보조정리의 증명을 참조하라.
$K'$를 $S'$의 체 함수로 설정한다.
그리고 $Y' = X' \times_{S'} \Spec(K')$를 사용하여 다이어그램을 얻습니다.
$$
\xymatrix{
X \ar[d]_f &
X' \ar[d] \ar[l] &
Y' \ar[l]^{h'} \ar[d] &
Y \ar[l] \ar[d] \\
S &
S' \ar[l] &
\Spec(K') \ar[l] &
\Spec(K) \ar[l]
}
$$
보조정리 \ref{etale-cohomology-lemma-smooth-base-change-separably-closed}에 의해
총계 직상 $R(Y \to Y')_*\mathbf{Z}/d\mathbf{Z}$
와 동형이다$\mathbf{Z}/d\mathbf{Z}$~에$D(Y'_\etale)$; 여기
$n$는 $S$에서 반전 가능하다는 것을 사용한다.
따라서 $Rh'_*\mathbf{Z}/d\mathbf{Z} = Rh_*\mathbf{Z}/d\mathbf{Z}$
상대 Leray 스펙트럼 열에 의해. 이로써 증명이 완료된다.
\end{proof}










\section{고유 기저변환 정리} \label{etale-cohomology-section-proper-base-change}

\noindent 고유 기저변환 정리는 이 절에서 진술되고 증명된다. 우리의 접근 방식은 인수를 약간 단순화하기 위해 Gabber의 아이디어(아핀 사례에서 나온)를 사용하여 \cite[XII, Theorem 5.1]{SGA4}의 증명을 대략 따릅니다.

\begin{lemma}
\label{etale-cohomology-lemma-zariski-h0-proper-over-henselian-pair}
$(A, I)$를 헨셀 쌍이라고 가정한다. $f : X \to \Spec(A)$를 고유 사상으로 설정한다.
스킴. $Z = X \times_{\Spec(A)} \Spec(A/I)$로 놔두세요. 누구에게나
층 $\mathcal{F}$ $X$와 연관된 위상 공간에서 우리는
$\Gamma(X, \mathcal{F}) = \Gamma(Z, \mathcal{F}|_Z)$이 있다.
\end{lemma}

\begin{proof}
이를 증명하기 위해 보조정리 \ref{etale-cohomology-lemma-h0-topological}를 사용하겠습니다. 먼저 관찰하다
$X$의 기본 위상 공간은 성질, 보조정리에 의해 스펙트럼이다.
\ref{properties-lemma-quasi-compact-quasi-separated-spectral}.
$Y \subset X$를 기약 폐쇄형 하위 체계라고 가정한다. 증명을 끝내려면
$Y \cap Z = Y \times_{\Spec(A)} \Spec(A/I)$는 연결임을 보여줍니다.
$X$을 $Y$로 교체
$X$이 기약라고 가정할 수 있으며 $Z$를 보여야 한다.
연결이다. $X \to \Spec(B) \to \Spec(A)$를 Stein 인수분해라고 하자.
의 $f$ (사상, 정리에 대한 추가 정보
\ref{more-morphisms-theorem-stein-factorization-general}).
그러면 $A \to B$는 정역이고 $(B, IB)$는 헨셀 쌍이다.
(대수학에 대한 자세한 내용, 보조정리 \ref{more-algebra-lemma-integral-over-henselian-pair}).
따라서 $X \to \Spec(A)$의 올이 기하학적으로 가정할 수 있다.
연결. 반면에 이미지는$T \subset \Spec(A)$~의$f$
기약이고 $X$이 $A$보다 적절하므로 닫혀 있다. 따라서 $T \cap V(I)$
대수학에 대한 추가 정보, 보조정리의 연결이다.
\ref{more-algebra-lemma-irreducible-henselian-pair-connected}.
이제 $Y \times_{\Spec(A)} \Spec(A/I) \to T \cap V(I)$
은 연결 올을 사용하는 전사적 폐쇄 지도이다.
이제 결과는 위상, 보조정리의 결과를 따릅니다.
\ref{topology-lemma-connected-fibres-quotient-topology-connected-components}.
\end{proof}

\begin{lemma}
\label{etale-cohomology-lemma-h0-proper-over-henselian-pair}
$(A, I)$를 헨셀 쌍이라고 가정한다. $f : X \to \Spec(A)$를 고유 사상으로 설정한다.
스킴. $i : Z \to X$를 닫힌 몰입으로 설정한다.
$X \times_{\Spec(A)} \Spec(A/I)$을 $X$로 변환한다. 누구에게나
층 $\mathcal{F}$ $X_\etale$ 우리는
$\Gamma(X, \mathcal{F}) = \Gamma(Z, i_{small}^{-1}\mathcal{F})$이 있다.
\end{lemma}

\begin{proof} 이는 보조정리 \ref{etale-cohomology-lemma-gabber-h0} 및 \ref{etale-cohomology-lemma-zariski-h0-proper-over-henselian-pair}에서 파생되며 $X$에 대한 모든 스킴 유한은 $\Spec(A)$에 대해 적절하다는 사실이다. \end{proof}

\begin{lemma}
\label{etale-cohomology-lemma-h0-proper-over-henselian-local}
$A$를 헨셀리안 국소환이라고 가정한다. $f : X \to \Spec(A)$
스킴의 고유 사상이어야 한다. $X_0 \subset X$를 올로 설정한다.
$f$닫힌 지점 위에. 누구에게나층 $\mathcal{F}$~에$X_\etale$우리
$\Gamma(X, \mathcal{F}) =
\Gamma(X_0, \mathcal{F}|_{X_0})$이 있다.
\end{lemma}

\begin{proof} 이는 보조정리 \ref{etale-cohomology-lemma-h0-proper-over-henselian-pair}의 특별한 경우이다. \end{proof}

\noindent
$f : X \to S$를 스킴의 사상으로 설정한다. 허락하다
$\overline{s} : \Spec(k) \to S$는 기하점이다. 올
~의$f$~에$\overline{s}$은스킴
$X_{\overline{s}} =
\Spec(k) \times_{\overline{s}, S} X$ 조회됨
$\Spec(k)$에 대한 스킴으로 표시된다. $\mathcal{F}$가 층인 경우
$X_\etale$, 다음을 나타냅니다.
$\mathcal{F}_{\overline{s}} = p_{small}^{-1}\mathcal{F}$
$\mathcal{F}$을 $(X_{\overline{s}})_\etale$로 철회한다.
다음에서는 집합을 고려해 보자.
$$
\Gamma(X_{\overline{s}}, \mathcal{F}_{\overline{s}})
$$
$s \in S$를 $\overline{s}$의 이미지 포인트로 둡니다. $\kappa(s)^{sep}$
의 분리 가능한 대수적 종결이 되십시오.$\kappa(s)$~에$k$에서와 같이
정의 \ref{etale-cohomology-definition-algebraic-geometric-point}.
보조정리 \ref{etale-cohomology-lemma-sections-base-field-extension}에 의해
당김은 전단을 정의합니다
$$
\Gamma(X_{\kappa(s)^{sep}},
p_{sep}^{-1} \mathcal{F})
\longrightarrow
\Gamma(X_{\overline{s}}, \mathcal{F}_{\overline{s}})
$$
여기서 $p_{sep} : X_{\kappa(s)^{sep}} =
\Spec(\kappa(s)^{sep}) \times_S X \to X$
투영이다.

\begin{lemma}
\label{etale-cohomology-lemma-proper-pushforward-stalk}
$f : X \to S$를 스킴의 고유 사상으로 설정한다. 허락하다
$\overline{s} \to S$는 기하점이다.
누구에게나층 $\mathcal{F}$~에$X_\etale$
정식 지도
$$
(f_*\mathcal{F})_{\overline{s}} \longrightarrow
\Gamma(X_{\overline{s}}, \mathcal{F}_{\overline{s}})
$$
전단적이다.
\end{lemma}

\begin{proof} 정리 \ref{etale-cohomology-theorem-higher-direct-images} (집합의 층)에 의해 $$
(f_*\mathcal{F})_{\overline{s}} =
\Gamma(X \times_S \Spec(\mathcal{O}_{S, \overline{s}}^{sh}),
p_{small}^{-1}\mathcal{F})
$$가 있다. 여기서 $p : X \times_S \Spec(\mathcal{O}_{S, \overline{s}}^{sh}) \to X$는 투영이다. 엄격한 헨셀식 국소환 $\mathcal{O}_{S, \overline{s}}^{sh}$의 나머지 체는 $\kappa(s)^{sep}$이므로 보조정리 및 보조정리 \ref{etale-cohomology-lemma-h0-proper-over-henselian-local} 위의 논의에서 결론을 내립니다. \end{proof}

\begin{lemma} \label{etale-cohomology-lemma-proper-base-change-f-star} $f : X \to Y$를 스킴의 고유 사상으로 둡니다. $g : Y' \to Y$를 스킴의 사상으로 설정한다. 집합 $X' = Y' \times_Y X$ 투영 $f' : X' \to Y'$ 및 $g' : X' \to X$이 포함된다. $\mathcal{F}$를 $X_\etale$의 층으로 설정한다. 그런 다음 $g^{-1}f_*\mathcal{F} = f'_*(g')^{-1}\mathcal{F}$. \end{lemma}

\begin{proof}
표준 맵 $g^{-1}f_*\mathcal{F} \to f'_*(g')^{-1}\mathcal{F}$이 있다.
즉, 지도와 인접해 있다.
$$
f_*\mathcal{F} \longrightarrow
g_*f'_*(g')^{-1}\mathcal{F} = f_*g'_*(g')^{-1}\mathcal{F}
$$
이는 표준 맵에 적용되는 $f_*$이다.
$\mathcal{F} \to g'_*(g')^{-1}\mathcal{F}$. 이 지도를 확인하려면
동형사상 줄기에서 무슨 일이 일어나는지 계산할 수 있다.
허락하다$y' : \Spec(k) \to Y'$가 되다기하점이미지 포함$y$~에$Y$.
보조정리 \ref{etale-cohomology-lemma-proper-pushforward-stalk}에 의해 줄기는
$\Gamma(X'_{y'}, \mathcal{F}_{y'})$ 및 $\Gamma(X_y, \mathcal{F}_y)$
각기. 여기에서는 층 $\mathcal{F}_y$ 및 $\mathcal{F}_{y'}$
의 후퇴이다$\mathcal{F}$예측에 의해$X_y \to X$
및 $X'_{y'} \to X$. $X_y = X'_{y'}$ 이후 스킴
줄기가 동의한다는 것을 알 수 있다. 우리는
이 동형사상이 우리 지도와 호환되는지 확인하세요.
\end{proof}

\noindent
이 시점에서 고유 기저변환 정리에 대해 논의하기 시작한다.
이를 위해 몇 가지 표기법을 소개한다. 가환 도식을 고려해보세요
\begin{equation}
\label{etale-cohomology-equation-base-change-diagram}
\vcenter{
\xymatrix{
X' \ar[r]_{g'} \ar[d]_{f'} & X \ar[d]^f \\
Y' \ar[r]^g & Y
}
}
\end{equation}
사상의 스킴. 그런 다음 사이트의 가환 도식을 얻습니다.
$$
\xymatrix{
X'_\etale \ar[r]_{g'_{small}} \ar[d]_{f'_{small}} &
X_\etale \ar[d]^{f_{small}} \\
Y'_\etale \ar[r]^{g_{small}} &
Y_\etale
}
$$
어떤 물체에 대해서도$E$~의$D(X_\etale)$표준 베이스 변경 맵을 얻습니다.
\begin{equation}
\label{etale-cohomology-equation-base-change}
g_{small}^{-1}Rf_{small, *}E \longrightarrow Rf'_{small, *}(g'_{small})^{-1}E
\end{equation}
$D(Y'_\etale)$에서. 사이트, 비고의 코호몰로지를 참조하라.
\ref{sites-cohomology-remark-base-change} 상수를 사용하는 곳
층 $\mathbf{Z}$를 환의 층으로 사용한다.
일반적으로 이 수식에서는 아래 첨자 ${}_{small}$를 생략한다.
예의 경우, $E = \mathcal{F}[0]$ 여기서 $\mathcal{F}$는 아벨리안이다.
층 위의 $X_\etale$, 베이스 변경 맵은 맵이다.
\begin{equation}
\label{etale-cohomology-equation-base-change-sheaf}
g^{-1}Rf_*\mathcal{F} \longrightarrow Rf'_*(g')^{-1}\mathcal{F}
\end{equation}
$D(Y'_\etale)$에서.

\medskip\noindent
지도(\ref{etale-cohomology-equation-base-change})는 동형사상이 될 가능성이 없다.
위에 주어진 일반론에서. 보여드리는 것이 목표입니다
다이어그램 (\ref{etale-cohomology-equation-base-change-diagram})인 경우 동형사상이다.
데카르트, $f : X \to Y$ 고유, 코호몰로지 층 $E$
비틀림이고 $E$는 아래로 제한된다.
이 질문을 연구하기 위해 다음 용어를 소개한다.
{\it 코호몰로지가 베이스 변경으로 통근한다고 가정해 보자.
에 대하여 $f : X \to Y$} 만약 (\ref{etale-cohomology-equation-base-change-sheaf})
모든 다이어그램에 대한 동형사상입니다(\ref{etale-cohomology-equation-base-change-diagram}).
여기서 $X' = Y' \times_Y X$ 및 모든 비틀림 아벨 층 $\mathcal{F}$.

\begin{lemma} \label{etale-cohomology-lemma-proper-base-change-in-terms-of-injectives} $f : X \to Y$를 스킴의 고유 사상이라 하자. 다음은 동치이다. \begin{enumerate} \item $f$에 대하여 코호몰로지는 밑변환과 가환한다(위 참조). \item 모든 소수 $\ell$과 모든 단사 $\mathbf{Z}/\ell\mathbf{Z}$-가군층 $\mathcal{I}$ 중 $X_\etale$ 위의 것, 그리고 $X' = Y' \times_Y X$인 모든 도식 (\ref{etale-cohomology-equation-base-change-diagram})에 대하여, 층 $R^qf'_*(g')^{-1}\mathcal{I}$는 $q > 0$이면 0이다. \end{enumerate} \end{lemma}

\begin{proof}
(1)는 (2)를 의미하는 것이 분명한다. 반대로, (2)를 가정하고
$\mathcal{F}$는 $X_\etale$에 비틀림 아벨 층이 있다. $Y' \to Y$
스킴의 사상이고 $X' = Y' \times_Y X$
다음과 같이 $g' : X' \to X$ 및 $f' : X' \to Y'$ 투영을 사용한다.
다이어그램(\ref{etale-cohomology-equation-base-change-diagram}).
층의 지도를 보여주고 싶다.
$$
g^{-1}R^qf_*\mathcal{F} \longrightarrow R^qf'_*(g')^{-1}\mathcal{F}
$$
모든 $q \geq 0$에 대해 동형사상이다.

\medskip\noindent
모든 $n \geq 1$에 대해 $\mathcal{F}[n]$를 절의 하위 묶음으로 둡니다.
$\mathcal{F}$이 $n$에 의해 전멸되었다. 그 다음에
$\mathcal{F} = \colim \mathcal{F}[n]$.
함자 $g^{-1}$ 및 $(g')^{-1}$는 임의의 여극한으로 통근한다.
(왼쪽 인접). $f$ 또는 $f'$를 따라 고차 직상을 사용한다.
보조정리 \ref{etale-cohomology-lemma-relative-colimit}로 여과 여극한으로 출퇴근한다.
그러므로 우리는
$$
g^{-1}R^qf_*\mathcal{F} = \colim g^{-1}R^qf_*\mathcal{F}[n]
\quad\text{및}\quad
R^qf'_*(g')^{-1}\mathcal{F} =
\colim R^qf'_*(g')^{-1}\mathcal{F}[n]
$$
따라서 $\mathcal{F}$의 경우 결과를 증명하면 충분한다.
양의 정수 $n$에 의해 소멸된다.

\medskip\noindent
만약에$n = \ell n'$일부 소수의 경우$\ell$, 그러면 우리는 짧은 것을 얻습니다
완전열
$$
0 \to \mathcal{F}[\ell] \to \mathcal{F} \to
\mathcal{F}/\mathcal{F}[\ell] \to 0
$$
$\mathcal{F}/\mathcal{F}[\ell]$이 $n'$에 의해 소멸되는 것을 관찰하세요.
또한 결과가 $\mathcal{F}[\ell]$ 및
$\mathcal{F}/\mathcal{F}[\ell]$, 결과는 다음과 같이 유지된다.
고차 직상의 긴 완전열 (및 $5$ 보조정리).
이런 식으로 $\mathcal{F}$가 소멸되는 경우로 축소한다.
소수 $\ell$로 표시된다.

\medskip\noindent
$\mathcal{F}$이 소수 $\ell$에 의해 소멸된다고 가정한다.
단사 분해 $\mathcal{F} \to \mathcal{I}^\bullet$를 선택하세요.
$D(X_\etale, \mathbf{Z}/\ell\mathbf{Z})$에서. 가정 적용 중
(2) 및 Leray의 비순환성 보조정리
(유도 범주, 보조정리 \ref{derived-lemma-leray-acyclicity})
우리는 그것을 본다
$$
f'_*(g')^{-1}\mathcal{I}^\bullet
$$
$Rf'_*(g')^{-1}\mathcal{F}$를 계산한다. 적용하여 마무리합니다
보조정리 \ref{etale-cohomology-lemma-proper-base-change-f-star}.
\end{proof}

\begin{lemma} \label{etale-cohomology-lemma-sandwich} $f : X \to Y$와 $g : Y \to Z$를 스킴의 고유 사상이라 하자. 다음을 가정한다. \begin{enumerate} \item 코호몰로지는 $f$에 대한 밑변환과 가환한다. \item 코호몰로지는 $g \circ f$에 대한 밑변환과 가환한다. \item $f$는 전사이다. \end{enumerate} 그러면 코호몰로지는 $g$에 대한 밑변환과 가환한다. \end{lemma}

\begin{proof}
우리는 동등성을 사용합니다
보조정리 \ref{etale-cohomology-lemma-proper-base-change-in-terms-of-injectives}
더 이상 언급하지 않고. $\ell$를 소수로 둡니다.
$\mathcal{I}$를 다음의 단사 층으로 가정한다.
$\mathbf{Z}/\ell\mathbf{Z}$-가군 $Y_\etale$에 있다.
층 $f^{-1}\mathcal{I} \to \mathcal{J}$의 주입형 맵을 선택하세요.
여기서 $\mathcal{J}$는 다음의 단사 층이다.
$\mathbf{Z}/\ell\mathbf{Z}$-가군 $X_\etale$에 있다.
$f$는 전사이므로 지도 $\mathcal{I} \to f_*\mathcal{J}$
분사형입니다(기하점의 줄기 참조).
$\mathcal{I}$는 단사형이므로 $\mathcal{I}$
는 $f_*\mathcal{J}$의 직접적인 합이다. 그러므로 충분하다
$f_*\mathcal{J}$에 대해 원하는 소멸을 증명한다.

\medskip\noindent
$Z' \to Z$를 스킴 및 집합의 사상으로 설정한다.
$Y' = Z' \times_Z Y$ 및 $X' = Z' \times_Z X = Y' \times_ Y X$.
$a : X' \to X$, $b : Y' \to Y$ 및 $c : Z' \to Z$를 나타냅니다.
예측. $f' : X' \to Y'$ 및 $g' : Y' \to Z'$도 마찬가지이다.
보조정리 \ref{etale-cohomology-lemma-proper-base-change-f-star}에 의해 우리는
$b^{-1}f_*\mathcal{J} = f'_*a^{-1}\mathcal{J}$.
반면에 우리는 $R^qf'_*a^{-1}\mathcal{J}$와
$R^q(g' \circ f')_*a^{-1}\mathcal{J}$는 $q > 0$에 대해 0이다.
스펙트럼 열 사용
(코호몰로지 사이트, 보조정리 \ref{sites-cohomology-lemma-relative-Leray})
$$
R^pg'_* R^qf'_* a^{-1}\mathcal{J} \Rightarrow
R^{p + q}(g' \circ f')_* a^{-1}\mathcal{J}
$$
우리는 다음과 같이 결론을 내립니다.
$ R^pg'_*(b^{-1}f_*\mathcal{J}) = R^pg'_*(f'_*a^{-1}\mathcal{J}) = 0$
원하는 대로 $p > 0$에 대해.
\end{proof}

\begin{lemma} \label{etale-cohomology-lemma-composition} $f : X \to Y$ 및 $g : Y \to Z$를 스킴의 고유 사상으로 설정한다. \begin{enumerate} \item 코호몰로지는 $f$에 대한 베이스 변경으로 통근하고, \item 코호몰로지는 $g$에 대한 베이스 변경으로 통근한다고 가정한다. \end{enumerate} 그러면 코호몰로지는 $g \circ f$에 대한 베이스 변경으로 통근한다. \end{lemma}

\begin{proof}
우리는 동등성을 사용합니다
보조정리 \ref{etale-cohomology-lemma-proper-base-change-in-terms-of-injectives}
더 이상 언급하지 않고. $\ell$를 소수로 둡니다.
$\mathcal{I}$를 다음의 단사 층으로 가정한다.
$\mathbf{Z}/\ell\mathbf{Z}$-가군 $X_\etale$에 있다.
그러면 $f_*\mathcal{I}$는 다음의 단사 층이다.
$\mathbf{Z}/\ell\mathbf{Z}$-가군 $Y_\etale$의
(코호몰로지 위의 사이트, 보조정리
\ref{sites-cohomology-lemma-pushforward-injective-flat}).
결과는 공식적으로 이것으로부터 이어지지만, 우리는 또한
그것을 철자하라.

\medskip\noindent
$Z' \to Z$를 스킴 및 집합의 사상으로 설정한다.
$Y' = Z' \times_Z Y$ 및 $X' = Z' \times_Z X = Y' \times_ Y X$.
$a : X' \to X$, $b : Y' \to Y$ 및 $c : Z' \to Z$를 나타냅니다.
예측. $f' : X' \to Y'$ 및 $g' : Y' \to Z'$도 마찬가지이다.
보조정리 \ref{etale-cohomology-lemma-proper-base-change-f-star}에 의해 우리는
$b^{-1}f_*\mathcal{I} = f'_*a^{-1}\mathcal{I}$.
반면에 우리는 $R^qf'_*a^{-1}\mathcal{I}$와
$R^q(g')_*b^{-1}f_*\mathcal{I}$는 $q > 0$에 대해 0이다.
스펙트럼 열 사용
(코호몰로지 사이트, 보조정리 \ref{sites-cohomology-lemma-relative-Leray})
$$
R^pg'_* R^qf'_* a^{-1}\mathcal{I} \Rightarrow
R^{p + q}(g' \circ f')_* a^{-1}\mathcal{I}
$$
우리는 $R^p(g' \circ f')_*a^{-1}\mathcal{I} = 0$에 대해 결론을 내렸습니다.
$p > 0$ 원하는 대로.
\end{proof}

\begin{lemma}
\label{etale-cohomology-lemma-finite}
\begin{slogan}
고유 기저변환 정리는 에탈 코호몰로지에서 유한 사상에 대하여 성립한다.
\end{slogan}
$f : X \to Y$를 스킴의 유한 사상이라 하자. 그러면 코호몰로지는 $f$에 대한 밑변환과 가환한다.
\end{lemma}

\begin{proof} 유한 사상이 적절한지 확인하세요. 사상, 보조정리 \ref{morphisms-lemma-finite-proper}를 참조하라. 또한 유한 사상의 기본 변경은 유한한다. 사상, 보조정리 \ref{morphisms-lemma-base-change-finite}를 참조하라. 따라서 보조정리 \ref{etale-cohomology-lemma-proper-base-change-in-terms-of-injectives}와 명제 \ref{etale-cohomology-proposition-finite-higher-direct-image-zero}를 결합한 결과는 다음과 같습니다. \end{proof}

\begin{lemma} \label{etale-cohomology-lemma-reduce-to-P1} 코호몰로지가 스킴의 모든 고유 사상에 대해 기본 변경으로 통근한다는 것을 증명하려면 모든 사상 $\mathbf{P}^1_S \to S$에 대해 유지된다는 것을 증명하면 충분한다. 스킴 $S$. \end{lemma}

\begin{proof} $f : X \to Y$를 스킴의 고유 사상으로 설정한다. $Y = \bigcup Y_i$를 아핀 개방형 피복과 집합 $X_i = f^{-1}(Y_i)$로 설정한다. 코호몰로지가 $X_i \to Y_i$에 대한 기본 변경으로 통근한다는 것을 증명할 수 있다면 코호몰로지는 $f$에 대한 기본 변경으로 통근한다. 즉, 고차 직상의 형성은 자리스키(심지어 에탈)와 통근하는 국소화를 베이스로 보조정리 \ref{etale-cohomology-lemma-higher-direct-images}를 참조하라. 따라서 우리는 $Y$가 유사하다고 가정할 수 있다.

\medskip\noindent $Y$를 아핀 스킴으로 두고 $X \to Y$를 고유 사상으로 둡니다. Chow의 보조정리에 따르면 가환 도식 $$
\xymatrix{
X \ar[rd] & X' \ar[d] \ar[l]^\pi \ar[r] & \mathbf{P}^n_Y \ar[dl] \\
& Y &
}
$$가 존재한다. 여기서 $X' \to \mathbf{P}^n_Y$는 몰입이고 $\pi : X' \to X$는 고유 및 전사이다. 극한, 보조정리를 참조하라. \ref{limits-lemma-chow-finite-type}. $X \to Y$가 적절하므로 $X' \to Y$가 적절하다고 판단됩니다(사상, 보조정리 \ref{morphisms-lemma-composition-proper}). 따라서 $X' \to \mathbf{P}^n_Y$는 닫힌 몰입 (사상, 보조정리 \ref{morphisms-lemma-image-proper-scheme-closed})이다. $X' \to X \times_Y \mathbf{P}^n_Y = \mathbf{P}^n_X$는 닫힌 몰입입니다(닫힌 이미지의 몰입).

\medskip\noindent 보조정리 \ref{etale-cohomology-lemma-sandwich}에 의해 코호몰로지가 $\pi$ 및 $X' \to Y$에 대한 기본 변경으로 통근한다는 것을 증명하는 것으로 충분한다. 이러한 사상은 둘 다 닫힌 몰입과 투사 $\mathbf{P}^n_S \to S$(일부 $S$의 경우)로 간주된다. 보조정리 \ref{etale-cohomology-lemma-finite}에 의해 결과는 닫힌 침수에 대해 유지됩니다(닫힌 몰입은 유한하므로). 보조정리 \ref{etale-cohomology-lemma-composition}에 의해 투영 $\mathbf{P}^n_S \to S$에 대한 결과를 증명하는 것으로 충분한다.

\medskip\noindent
모든 $n \geq 1$에는 유한 전사 사상이 있다.
$$
\mathbf{P}^1_S \times_S \ldots \times_S \mathbf{P}^1_S
\longrightarrow
\mathbf{P}^n_S
$$
좌표에 의해 주어진
$$
((x_1 : y_1), (x_2 : y_2), \ldots, (x_n : y_n))
\longmapsto
(F_0 : \ldots : F_n)
$$
어디$F_0, \ldots, F_n$~에$x_1, \ldots, y_n$
다음과 같은 정수 계수를 갖는 다항식은 다음과 같습니다.
$$
\prod (x_i t + y_i) = F_0 t^n + F_1 t^{n - 1} + \ldots + F_n
$$
신청
기본정리 \ref{etale-cohomology-lemma-sandwich}, \ref{etale-cohomology-lemma-finite} 및 \ref{etale-cohomology-lemma-composition}
다시 한 번 보조정리가 참이라는 결론을 내렸습니다.
\end{proof}

\begin{theorem}[고유 기저변환] \label{etale-cohomology-theorem-proper-base-change} $f : X \to Y$를 스킴의 고유 사상으로 둡니다. $g : Y' \to Y$를 스킴의 사상으로 설정한다. 집합 $X' = Y' \times_Y X$ 그리고 데카르트 다이어그램 $$
\xymatrix{
X' \ar[r]_{g'} \ar[d]_{f'} & X \ar[d]^f \\
Y' \ar[r]^g & Y
}
$$을 고려하라. $\mathcal{F}$를 $X_\etale$의 아벨 꼬임층으로 둡니다. 그러면 기본 변경 맵 $$
g^{-1}Rf_*\mathcal{F} \longrightarrow Rf'_*(g')^{-1}\mathcal{F}
$$은 동형사상이다. \end{theorem}

\begin{proof}
위에서 소개한 용어에서 이는 코호몰로지 통근을 의미한다.
스킴의 고유 사상마다 베이스가 변경된다. 에 의해
보조정리 \ref{etale-cohomology-lemma-reduce-to-P1}
코호몰로지가 베이스 변경으로 통근한다는 것을 증명하는 것으로 충분한다.
사상 $\mathbf{P}^1_S \to S$마다 스킴 $S$마다.

\medskip\noindent
$S$를 닫힌 헨셀리언 국소환의 스펙트럼으로 정의한다.
포인트 $s$. 집합 $X = \mathbf{P}^1_S$ 및 $X_0 = X_s = \mathbf{P}^1_s$.
$\mathcal{F}$를 $\mathbf{Z}/\ell\mathbf{Z}$-가군의 층으로 설정한다.
$X_\etale$에 있다. 증명의 핵심은 다음과 같습니다.
$$
H^q_\etale(X, \mathcal{F}) = H^q_\etale(X_0, \mathcal{F}|_{X_0}).
$$
즉, 해상도 $\mathcal{F} \to \mathcal{I}^\bullet$를 선택한다.
$\mathbf{Z}/\ell\mathbf{Z}$-가군의 층 분사로.
그러면 $\mathcal{I}^\bullet|_{X_0}$는 $\mathcal{F}|_{X_0}$의 해상도이다.
오른쪽 $H^0_\etale(X_0, -)$-비순환 객체에 대해서는 참조하라.
보조정리 \ref{etale-cohomology-lemma-efface-cohomology-on-fibre-by-finite-cover}.
Leray의 비순환성 보조정리는 우변이 다음과 같이 계산된다는 것을 알려줍니다.
복합체 $H^0_\etale(X_0, \mathcal{I}^\bullet|_{X_0})$
이는 $H^0_\etale(X, \mathcal{I}^\bullet)$와 같습니다.
보조정리 \ref{etale-cohomology-lemma-h0-proper-over-henselian-local}. 이 복합체
좌변을 계산한다.

\medskip\noindent
$S$이 일반적이고 $\mathcal{F}$가 층라고 가정한다.
$\mathbf{Z}/\ell\mathbf{Z}$-가군 $X_\etale$에 있다.
$\overline{s} : \Spec(k) \to S$를 기하점으로 설정한다.
$S$이 $s \in S$ 위에 놓여 있다. 우리는
$$
(R^qf_*\mathcal{F})_{\overline{s}} =
H^q_\etale(\mathbf{P}^1_{\mathcal{O}_{S, \overline{s}}^{sh}},
\mathcal{F}|_{\mathbf{P}^1_{\mathcal{O}_{S, \overline{s}}^{sh}}}) =
H^q_\etale(\mathbf{P}^1_{\kappa(s)^{sep}},
\mathcal{F}|_{\mathbf{P}^1_{\kappa(s)^{sep}}})
$$
여기서 $\kappa(s)^{sep}$는 다음의 잔기 체이다.
$\mathcal{O}_{S, \overline{s}}^{sh}$, 즉 분리 가능한 대수
폐쇄$\kappa(s)$~에$k$.
정리 \ref{etale-cohomology-theorem-higher-direct-images}에 의한 첫 번째 동등성
그리고 표시된 공식에 의한 두 번째 동등성은
이전 단락.

\medskip\noindent
마지막으로 스킴 $g : T \to S$ 중 사상을 고려하세요.
$S$ 과 $\mathcal{F}$ 는 위와 같습니다.
집합 $f' : \mathbf{P}^1_T \to T$ 투영하고
$g' : \mathbf{P}^1_T \to \mathbf{P}^1_S$ 사상 유도
$g$로. 기본 변경 맵을 고려하세요.
$$
g^{-1}R^qf_*\mathcal{F}
\longrightarrow
R^qf'_*(g')^{-1}\mathcal{F}
$$
$\overline{t}$를 이미지가 있는 $T$의 기하점으로 설정한다.
$\overline{s} = g(\overline{t})$. 우리의 토론으로
$\overline{t}$에 있는 줄기 지도 위는 지도이다.
$$
H^q_\etale(\mathbf{P}^1_{\kappa(s)^{sep}},
\mathcal{F}|_{\mathbf{P}^1_{\kappa(s)^{sep}}})
\longrightarrow
H^q_\etale(\mathbf{P}^1_{\kappa(t)^{sep}},
\mathcal{F}|_{\mathbf{P}^1_{\kappa(t)^{sep}}})
$$
$\kappa(s)^{sep} \subset \kappa(t)^{sep}$ 이 지도는
동형사상 에 의해 보조정리 \ref{etale-cohomology-lemma-base-change-dim-1-separably-closed}.

\medskip\noindent
이는 코호몰로지가 다음의 기본 변경으로 통근한다는 것을 증명한다.
$\mathbf{P}^1_S \to S$ 및 층의 $\mathbf{Z}/\ell\mathbf{Z}$-가군.
특히, $\mathbf{Z}/\ell\mathbf{Z}$-가군의 단사 층의 경우
모든 염기 변화의 고차 직상은 0이다.
즉, 조건 (2)
보조정리 \ref{etale-cohomology-lemma-proper-base-change-in-terms-of-injectives}
유지되고 증명이 완료된다.
\end{proof}

\begin{lemma} \label{etale-cohomology-lemma-proper-base-change} $f : X \to Y$를 스킴의 고유 사상으로 둡니다. $g : Y' \to Y$를 스킴의 사상으로 설정한다. 집합 $X' = Y' \times_Y X$ 및 $f' : X' \to Y'$ 및 $g' : X' \to X$ 투영을 나타냅니다. $E \in D^+(X_\etale)$에는 비틀림 코호몰로지 층이 있다고 가정한다. 그러면 기본 변경 맵(\ref{etale-cohomology-equation-base-change}) $g^{-1}Rf_*E \to Rf'_*(g')^{-1}E$는 동형사상이다. \end{lemma}

\begin{proof}
이는 고유 기저변환 정리의 간단한 결과이다.
(정리 \ref{etale-cohomology-theorem-proper-base-change}) 스펙트럼 사용
시퀀스
$$
E_2^{p, q} = R^pf_*H^q(E)
\quad\text{및}\quad
{E'}_2^{p, q} = R^pf'_*(g')^{-1}H^q(E)
$$
$R^nf_*E$ 및 $R^nf'_*(g')^{-1}E$로 수렴된다.
스펙트럼 열은 다음과 같이 구성된다.
유도 범주, 보조정리 \ref{derived-lemma-two-ss-complex-functor}.
일부 세부 사항은 생략되었다.
\end{proof}

\begin{lemma} \label{etale-cohomology-lemma-proper-base-change-stalk} $f : X \to Y$를 스킴의 고유 사상으로 둡니다. $\overline{y} \to Y$를 기하점으로 설정한다. \begin{enumerate} \item 비틀림 아벨 층 $\mathcal{F}$ $X_\etale$에 대해 $(R^nf_*\mathcal{F})_{\overline{y}} =
H^n_\etale(X_{\overline{y}}, \mathcal{F}_{\overline{y}})$이 있다. \item 비틀림 코호몰로지 층이 있는 $E \in D^+(X_\etale)$의 경우 $(R^nf_*E)_{\overline{y}} =
H^n_\etale(X_{\overline{y}}, E|_{X_{\overline{y}}})$이 있다. \end{enumerate} \end{lemma}

\begin{proof} 진술에서 $\mathcal{F}_{\overline{y}}$는 $\mathcal{F}$의 스킴 이론 올 $X_{\overline{y}} = \overline{y} \times_Y X$에 대한 후퇴를 나타냅니다. $\overline{y} \to Y$로 뒤로 당기면 $\mathcal{F}$의 줄기가 생성되므로 보조정리의 첫 번째 명령문은 정리 \ref{etale-cohomology-theorem-proper-base-change}의 특별한 경우이다. 두 번째는 보조정리 \ref{etale-cohomology-lemma-proper-base-change}의 특별한 경우이다. \end{proof}






\section{고유 기저변환} \label{etale-cohomology-section-applications-proper-base-change}의 응용

\noindent 이 절에서는 고유 기저변환 정리의 다소 즉각적인 결과에 대해 논의한다.

\begin{lemma} \label{etale-cohomology-lemma-base-change-separably-closed} $K/k$는 분리 가능하게 닫힌 체의 확장이라고 가정한다. $X$를 $k$에 대해 적절한 스킴으로 설정한다. $\mathcal{F}$를 $X_\etale$의 비틀림 아벨 층으로 설정한다. 그러면 지도 $H^q_\etale(X, \mathcal{F}) \to
H^q_\etale(X_K, \mathcal{F}|_{X_K})$는 $q \geq 0$에 대한 동형사상이다. \end{lemma}

\begin{proof} 줄기를 보면 이것이 정리 \ref{etale-cohomology-theorem-proper-base-change}의 특별한 경우라는 것을 알 수 있다. \end{proof}

\begin{lemma}
\label{etale-cohomology-lemma-cohomological-dimension-proper}
$f : X \to Y$를 스킴의 고유 사상으로 설정한다.
올에는 모두 차원 $\leq n$이 있다.
그러면 어떤 아벨리안에 대해서도꼬임층 $\mathcal{F}$~에$X_\etale$
우리는$R^qf_*\mathcal{F} = 0$~을 위한$q > 2n$.
\end{lemma}

\begin{proof} 모든 고유 사상에 대해 $n$에 대한 귀납법을 통해 이를 증명하겠습니다.

\medskip\noindent 만약 $n = 0$이면 $f$는 유한한 사상(사상, 보조정리 \ref{more-morphisms-lemma-characterize-finite}에 대한 자세한 내용)이며 결과는 다음과 같이 참이다. 명제 \ref{etale-cohomology-proposition-finite-higher-direct-image-zero}.

\medskip\noindent
$n > 0$인 경우 보조정리 \ref{etale-cohomology-lemma-proper-base-change-stalk}를 사용한다.
$H^i_\etale(X, \mathcal{F}) = 0$을 증명하는 것으로 충분하다는 것을 알 수 있다.
$i > 2n$ 및 $X$에 대해 적절한 스킴, $\dim(X) \leq n$
대수적으로 닫힌 체 $k$
그리고 $\mathcal{F}$는 $X$의 비틀림 아벨 층이다.

\medskip\noindent
$n = 1$인 경우 이는 정리 \ref{etale-cohomology-theorem-vanishing-curves}에서 따릅니다.
$n > 1$를 가정한다. 명제 \ref{etale-cohomology-proposition-topological-invariance}에 의해
$X$을 축소하여 대체할 수 있다.
$\nu : X^\nu \to X$를 정규화로 둡니다.
이것은 전사적 쌍합리 유한 사상이다.
(다양체, 보조정리 \ref{varieties-lemma-normalization-locally-algebraic} 참조)
따라서 조밀에 대한 동형사상은 $U \subset X$를 엽니다.
(사상, 보조정리 \ref{morphisms-lemma-birational-birational}).
그러면 우리는 그것을 본다
$c : \mathcal{F} \to \nu_*\nu^{-1}\mathcal{F}$는 단사형입니다
($\nu$은 전사이기 때문에) 및 $U$ 위에 동형사상이 있다.
$i : Z \to X$은 $U$의 보수 포함을 나타냅니다.
$U$는 $X$의 조밀이므로 $\dim(Z) < \dim(X) = n$가 있다.
명제 \ref{etale-cohomology-proposition-closed-immersion-pushforward}에 의해
$\Coker(c) = i_*\mathcal{G}$ 일부에게는
아벨리안 꼬임층 $\mathcal{G}$ 위의 $Z_\etale$.
그런 다음 $H^q_\etale(X, \Coker(c)) = H^q_\etale(Z, \mathcal{G})$
(명제 \ref{etale-cohomology-proposition-finite-higher-direct-image-zero} 기준
그리고 Leray 스펙트럼 열) 귀납 가설에 의해 우리는 다음과 같이 결론을 내립니다.
$c$의 여핵에는 $\leq 2(n - 1)$ 단위의 코호몰로지가 있다.
따라서 $\nu_*\nu^{-1}\mathcal{F}$에 대한 결과를 증명하면 충분한다.
$\nu$는 유한하므로 이는 다음과 같이 표시된다.
$H^i_\etale(X^\nu, \nu^{-1}\mathcal{F})$는
$i > 2n$의 경우 0이다. 이 경우는 다음 단락에서 다룬다.

\medskip\noindent
추정하다$X$~이다정역 정규적절한스킴~ 위에$k$~의차원 $n$.
일정하지 않은 유리함수를 선택하세요.$f$~에$X$. 그래프
$X' \subset X \times \mathbf{P}^1_k$~의$f$다이어그램에 앉아
$$
X \xleftarrow{b} X' \xrightarrow{f} \mathbf{P}^1_k
$$
$b$가 개방형 하위 구성표에 대한 동형사상인지 확인하세요.
$U \subset X$ 그 보완물은 닫힌 하위 체계이다.
$Z \subset X$ 공동차원 $\geq 2$. 즉, $U$는
도메인정의~의$f$모든 코드 치수가 포함되어 있다.$1$
$X$ 포인트, 참조
사상, 기본형 \ref{morphisms-lemma-rational-map-from-reduced-to-separated}
그리고 \ref{morphisms-lemma-extend-across}
(Serre의 정규성 기준과 결합하여 다음을 참조하라.
성질, 보조정리 \ref{properties-lemma-criterion-normal}).
또한 $b$의 올에는 차원 $\leq 1$가 있다.
$\mathbf{P}^1$). 따라서 $R^ib_*b^{-1}\mathcal{F}$는 0이 아닙니다.
유도에 의해 $i \in \{0, 1, 2\}$인 경우. 구별되는 삼각형을 선택하세요
$$
\mathcal{F} \to Rb_*b^{-1}\mathcal{F} \to Q \to \mathcal{F}[1]
$$
$\mathcal{F} \to b_*b^{-1}\mathcal{F}$를 사용하는 것은 단사적입니다
이전과 마찬가지로 방금 말한 내용을 사용하면 $Q$에 0이 아닌 것을 알 수 있다.
코호몰로지 층 $0, 1, 2$ $Z$에 앉아 있는 각도로만 표시된다.
게다가, 이들 코호몰로지 층은 비틀림에 의해
보조정리 \ref{etale-cohomology-lemma-torsion-direct-image}.
유도와 스펙트럼 열에 의해
유도 범주, 보조정리 \ref{derived-lemma-two-ss-complex-functor}
$H^i(X, Q)$는 0이다.
$i > 2 + 2\dim(Z) \leq 2 + 2(n - 2) = 2n - 2$. 그러므로 충분하다
$H^i(X', b^{-1}\mathcal{F}) = 0$에 대해 증명하기 위해
$i > 2n$. 이 시점에서 우리는 사상을 사용한다.
$$
f : X' \to \mathbf{P}^1_k
$$
올에는 차원 $< n$이 있다. 따라서 귀납법을 통해 우리는 다음을 알 수 있다.
$R^if_*b^{-1}\mathcal{F} = 0$~을 위한$i > 2(n - 1)$.
Leray 스펙트럼 열로 결론을 내립니다.
$$
H^i(\mathbf{P}^1_k, R^jf_*b^{-1}\mathcal{F})
\Rightarrow
H^{i + j}(X', b^{-1}\mathcal{F})
$$
그리고 $\dim(\mathbf{P}^1_k) = 1$라는 사실.
\end{proof}

\noindent 모드 $n$ 계수로 작업할 때 제한되지 않은 복합체에 대한 적절한 기본 변경을 수행할 수 있다.

\begin{lemma} \label{etale-cohomology-lemma-proper-base-change-mod-n} $f : X \to Y$를 스킴의 고유 사상으로 둡니다. $g : Y' \to Y$를 스킴의 사상으로 설정한다. 집합 $X' = Y' \times_Y X$ 및 $f' : X' \to Y'$ 및 $g' : X' \to X$ 투영을 나타냅니다. $n \geq 1$를 정수로 둡니다. $E \in D(X_\etale, \mathbf{Z}/n\mathbf{Z})$로 놔두세요. 그러면 기본 변경 맵(\ref{etale-cohomology-equation-base-change}) $g^{-1}Rf_*E \to Rf'_*(g')^{-1}E$은 동형사상이다. \end{lemma}

\begin{proof}
$Y$과 $Y'$가 준컴팩트한 경우 이를 증명하는 것으로 충분한다.
사상, 보조정리에 의해
\ref{morphisms-lemma-morphism-finite-type-bounded-dimension}
올의 차원이
$f : X \to Y$ 및 $f' : X' \to Y'$은 경계가 있다. 따라서
보조정리 \ref{etale-cohomology-lemma-cohomological-dimension-proper}는 다음을 의미한다.
$$
f_* :
\textit{Mod}(X_\etale, \mathbf{Z}/n\mathbf{Z})
\longrightarrow
\textit{Mod}(Y_\etale, \mathbf{Z}/n\mathbf{Z})
$$
그리고
$$
f'_* :
\textit{Mod}(X'_\etale, \mathbf{Z}/n\mathbf{Z})
\longrightarrow
\textit{Mod}(Y'_\etale, \mathbf{Z}/n\mathbf{Z})
$$
유한한 코호몰로지 차원을 갖는다는 의미에서
유도 범주, 보조정리 \ref{derived-lemma-unbounded-right-derived}.
K 단사를 선택하세요. 복합체 $\mathcal{I}^\bullet$
$\mathbf{Z}/n\mathbf{Z}$-가군 각각의
용어 $\mathcal{I}^n$는 다음의 단사 층이다.
$\mathbf{Z}/n\mathbf{Z}$-가군은 $E$을 나타냅니다.
분사 참조, 정리
\ref{injectives-theorem-K-injective-embedding-grothendieck}.
일반적인 고유 기저변환 정리로 우리는 다음을 찾습니다.
저것$R^qf'_*(g')^{-1}\mathcal{I}^n = 0$~을 위한$q > 0$, 보다
정리 \ref{etale-cohomology-theorem-proper-base-change}.
따라서 우리는 다음과 같이 결론을 내립니다.
유도 범주, 보조정리 \ref{derived-lemma-unbounded-right-derived}
복합체로 $Rf'_*(g')^{-1}E$를 계산할 수 있다.
$f'_*(g')^{-1}\mathcal{I}^\bullet$. 또 다른 응용
일반적인 고유 기저변환 정리는 다음을 보여줍니다.
이는 원하는 대로 $g^{-1}f_*\mathcal{I}^\bullet$와 같습니다.
\end{proof}

\begin{lemma} \label{etale-cohomology-lemma-pull-out-constant} $X$를 준컴팩트하고 준분리된 스킴라고 하자. $E \in D^+(X_\etale)$ 및 $K \in D^+(\mathbf{Z})$로 설정한다. 그러면 $$
R\Gamma(X, E \otimes_\mathbf{Z}^\mathbf{L} \underline{K}) =
R\Gamma(X, E) \otimes_\mathbf{Z}^\mathbf{L} K
$$ \end{lemma}

\begin{proof}
말하다$H^i(E) = 0$~을 위한$i \geq a$그리고$H^j(K) = 0$~을 위한$j \geq b$.
복합체 아래의 경계로 $K$을 나타낼 수 있다.
$K^\bullet$ 비틀림 없음 $\mathbf{Z}$-가군.
(K-플랫을 선택하세요.복합체 $L^\bullet$대표하는$K$그런 다음
$K^\bullet = \tau_{\geq b - 1}L^\bullet$을 사용한다. 이것은 작동한다.
$\mathbf{Z}$에는 전역 차원 $1$이 있다. 보다
대수학에 대한 추가 정보, 보조정리 \ref{more-algebra-lemma-last-one-flat}.)
복합체 아래의 경계로 $E$을 나타낼 수 있다.
$\mathcal{E}^\bullet$.
그러면 $E \otimes_\mathbf{Z}^\mathbf{L} \underline{K}$는
로 표현
$$
\text{Tot}(\mathcal{E}^\bullet \otimes_\mathbf{Z} \underline{K}^\bullet)
$$
구분된 삼각형 사용
$$
\sigma_{\geq -b + n + 1}K^\bullet
\to K^\bullet \to
\sigma_{\leq -b + n}K^\bullet
$$
그리고 사소한 소멸
$$
H^n(X,
\text{Tot}(\mathcal{E}^\bullet \otimes_\mathbf{Z}
\sigma_{\geq -a + n + 1}\underline{K}^\bullet) = 0
$$
그리고
$$
H^n(R\Gamma(X, E)
\otimes_\mathbf{Z}^\mathbf{L} \sigma_{\geq -a + n + 1}K^\bullet) = 0
$$
$K^\bullet$가 제한된 복합체인 경우로 축소한다.
플랫 $\mathbf{Z}$-가군. 주장을 반복하여 우리는
$K^\bullet$가 단일 플랫 $\mathbf{Z}$-가군과 동일한 경우
어느 정도 앉아 있다. 다음은 바보를 이용해서
$\mathcal{E}^\bullet$에 대한 잘림
우리는 다음과 같은 경우와 똑같은 방식으로 축소한다.
$\mathcal{E}^\bullet$는 아벨 층 한 마리가 앉아 있는 것이다.
어느 정도. 따라서 다음을 보여주는 것으로 충분한다.
$$
H^n(X, \mathcal{E} \otimes_\mathbf{Z} \underline{M}) =
H^n(X, \mathcal{E}) \otimes_\mathbf{Z} M
$$
$M$가 평평한 $\mathbf{Z}$-가군 및 $\mathcal{E}$인 경우
$X$의 아벨 층이다. 이 경우 우리는
쓰기 $M$는 유한 자유 $\mathbf{Z}$-가군의 여과 여극한이다.
(Lazard의 정리, 대수학, 정리 \ref{algebra-theorem-lazard} 참조).
정리 \ref{etale-cohomology-theorem-colimit}에 의해 이는 우리를 유한 자유의 경우로 축소시킵니다.
$\mathbf{Z}$-가군 $M$ 이 경우 결과는 간단하게 참이다.
\end{proof}

\begin{lemma}
\label{etale-cohomology-lemma-projection-formula-proper}
$f : X \to Y$를 스킴의 고유 사상으로 설정한다.
$E \in D^+(X_\etale)$에 비틀림 코호몰로지 층이 있다고 가정한다.
$K \in D^+(Y_\etale)$로 놔두세요. 그 다음에
$$
Rf_*E \otimes_\mathbf{Z}^\mathbf{L} K =
Rf_*(E \otimes_\mathbf{Z}^\mathbf{L} f^{-1}K)
$$
$D^+(Y_\etale)$에서.
\end{lemma}

\begin{proof}
왼쪽에서 오른쪽으로 표준 지도가 있다.
코호몰로지 사이트, 절 \ref{sites-cohomology-section-projection-formula}.
줄기에서 동등성을 확인하겠습니다.
파생된 텐서곱 계산은 당김과 함께 통근한다는 점을 기억하세요.
사이트, 보조정리의 코호몰로지를 참조하라.
\ref{sites-cohomology-lemma-pullback-tensor-product}.
따라서 우리는
$$
(E \otimes_\mathbf{Z}^\mathbf{L} f^{-1}K)_{\overline{x}}
=
E_{\overline{x}} \otimes_\mathbf{Z}^\mathbf{L} K_{\overline{y}}
$$
어디$\overline{y}$의 이미지이다$\overline{x}$~에$Y$.
$\mathbf{Z}$에는 전역 차원 $1$이 있으므로 다음을 알 수 있다.
이 복합체에는 소멸 코호몰로지 정도가 $< - 1 + a + b$ 있다.
만약에$H^i(E) = 0$~을 위한$i \geq a$그리고$H^j(K) = 0$~을 위한$j \geq b$.
게다가 $H^i(E)$는 비틀림 아벨식이므로
층 각각의 $i$에 대해, 코호몰로지에 대해서도 마찬가지이다.
복합체 $E \otimes_\mathbf{Z}^\mathbf{L} K$의 층.
즉, 우리는
$$
(E \otimes_\mathbf{Z}^\mathbf{L} f^{-1}K)
\otimes_{\mathbf{Z}}^\mathbf{L} \mathbf{Q} =
(E \otimes_\mathbf{Z}^\mathbf{L} \mathbf{Q})
\otimes_{\mathbf{Q}}^\mathbf{L}
(f^{-1}K \otimes_{\mathbf{Z}}^\mathbf{L} \mathbf{Q})
$$
이는 유도 범주에서 0이다.
이런 식으로 우리는 보조정리 \ref{etale-cohomology-lemma-proper-base-change-stalk}를 볼 수 있다.
보여주기에 충분하다는 것을 확인하기 위해 양쪽에 적용된다.
$$
R\Gamma(X_{\overline{y}},
E|_{X_{\overline{y}}} \otimes_\mathbf{Z}^\mathbf{L}
(X_{\overline{y}} \to \overline{y})^{-1}K_{\overline{y}}) =
R\Gamma(X_{\overline{y}},
E|_{X_{\overline{y}}}) \otimes_\mathbf{Z}^\mathbf{L} K_{\overline{y}}
$$
이는 보조정리 \ref{etale-cohomology-lemma-pull-out-constant}에 표시된다.
\end{proof}








\section{국소 비순환성}
\label{etale-cohomology-section-local-acyclicity}

\noindent 이 절에서는 매끄러운 기저변환 정리에서 매끄러운 사상의 국소 비순환성을 추론한다. SGA 4 또는 SGA 4.5에서 저자는 먼저 매끄러운 사상에 대한 국소 비순환성 버전을 증명한 다음 매끄러운 기저변환 정리를 추론한다.

\medskip\noindent
우리는 다음과 같이 주어진 국소 비순환성의 공식을 사용할 것이다.
\cite[Definition 2.12, page 242]{SGA4.5}를 정렬한다.
$f : X \to S$를 스킴의 사상으로 설정한다.
$\overline{x}$를 이미지가 있는 $X$의 기하점으로 설정한다.
$\overline{s} = f(\overline{x})$~에$S$.
$\overline{t}$를 기하점으로 설정한다.
$\Spec(\mathcal{O}^{sh}_{S, \overline{s}})$.
우리는 가환 도식을 얻습니다.
$$
\xymatrix{
F_{\overline{x}, \overline{t}} =
\overline{t} \times_{\Spec(\mathcal{O}^{sh}_{S, \overline{s}})}
\Spec(\mathcal{O}^{sh}_{X, \overline{x}}) \ar[r] \ar[d] &
\Spec(\mathcal{O}^{sh}_{X, \overline{x}}) \ar[r] \ar[d] &
X \ar[d] \\
\overline{t} \ar[r] &
\Spec(\mathcal{O}^{sh}_{S, \overline{s}}) \ar[r] &
S
}
$$
스킴 $F_{\overline{x}, \overline{t}}$가 호출된다.
{\it 다양체의 소멸 주기의 $f$ $\overline{x}$}.
$K$를 $D(X_\etale)$의 객체로 둡니다. 스킴 중 사상에 대해
$g : Y\to X$ 대신에 $R\Gamma(Y, K)$라고 씁니다.
$R\Gamma(Y_\etale, g_{small}^{-1}K)$. 부터
$\mathcal{O}^{sh}_{X, \overline{x}}$는 엄격히 헨셀적이다.
우리는
$K_{\overline{x}} = R\Gamma(\Spec(\mathcal{O}^{sh}_{X, \overline{x}}), K)$.
따라서 우리는 표준 맵을 얻습니다.
\begin{equation}
\label{etale-cohomology-equation-alpha-K}
\alpha_{K, \overline{x}, \overline{t}} :
K_{\overline{x}}
\longrightarrow
R\Gamma(F_{\overline{x}, \overline{t}}, K)
\end{equation}
코호몰로지를 뒤로 당겨서
$F_{\overline{x}, \overline{t}} \to
\Spec(\mathcal{O}^{sh}_{X, \overline{x}})$.

\begin{definition}
\label{etale-cohomology-definition-locally-acyclic}
\begin{reference}
\cite[Definition 2.12, page 242]{SGA4.5} 그리고
\cite[Definition (1.3), page 54]{SGA4.5}
\end{reference}
$f : X \to S$를 스킴의 사상으로 설정한다.
$K$를 $D(X_\etale)$의 객체로 둡니다.
\begin{enumerate}
\item
$\overline{x}$를 이미지가 있는 $X$의 기하점으로 설정한다.
$\overline{s} = f(\overline{x})$.
$f$가 $\overline{x}$에서 $K$에 관하여 {\it 국소 비순환적}이라고 함은
만약 기하점 $\overline{t}$마다
$\Spec(\mathcal{O}^{sh}_{S, \overline{s}})$
지도(\ref{etale-cohomology-equation-alpha-K})는 동형사상\footnote{이다.
$\overline{t}$는 다음의 대수적 기하점이라고 가정한다.
$\Spec(\mathcal{O}^{sh}_{S, \overline{s}})$. 자주 사용하는
보조정리 \ref{etale-cohomology-lemma-smooth-base-change-separably-closed}
이 경우로 줄일 수 있다.}.
\item $f$는 {\it 국소 비순환을 $K$}에 상대적이라고 말한다.
만약에$f$~이다국소 비순환~에$\overline{x}$상대적인$K$
모든기하점 $\overline{x}$~의$X$.
\item $f$가 $K$에 관하여 {\it 보편적으로 국소 비순환적}이라고 함은
스킴의 사상 $S' \to S$에 대해 기본 변경 $f' : X' \to S'$이 있는 경우
$K$에서 $X'$로의 철회에 상대적인 국소 비순환이다.
\item 모든 기하학에 대해 $f$는 {\it 국소 비순환}라고 말한다.
전철기$\overline{x}$~의$X$그리고 임의의 정수$n$에 프라임표수
$\kappa(\overline{x})$의 경우, 사상 $f$는 국소 비순환이다.
값이 있는 상수층을 기준으로 $\overline{x}$에서
$\mathbf{Z}/n\mathbf{Z}$.
\item $f$가 {\it 보편적으로 국소 비순환적}이라고 함은
스킴의 사상 $S' \to S$에 대해 기본 변경 $f' : X' \to S'$
국소 비순환이다.
\end{enumerate}
\end{definition}

\noindent
$M$를 아벨 군으로 설정한다. 그러면 $f : X \to S$의 국소 비순환성
상수층 $\underline{M}$에 대하여
그 요구 사항에
$$
H^q(F_{\overline{x}, \overline{t}}, \underline{M}) =
\left\{
\begin{matrix}
M & \text{만약에} & q = 0 \\
0 & \text{만약에} & q \not = 0
\end{matrix}
\right.
$$
누구에게나기하점 $\overline{x}$~의$X$그리고 어떤 기하학적
가리키다$\overline{t}$~의$\Spec(\mathcal{O}^{sh}_{S, f(\overline{x})})$.
이런 식으로 우리는 국소 비순환이 다음에 해당한다는 것을 알 수 있다.
소멸 기하학 올의 상위 코호몰로지 군 중
$F_{\overline{x}, \overline{t}}$ 사이의 지도
$\overline{x}$ 및 $\overline{s}$에서 엄격한 헨셀화를 수행한다.

\begin{proposition} \label{etale-cohomology-proposition-smooth-locally-acyclic} $f : X \to S$를 스킴의 매끄러운 사상으로 둡니다. 그러면 $f$는 보편적으로 국소 비순환이다. \end{proposition}

\begin{proof}
매끄러운 사상의 베이스 변경이 원활하므로 충분한다.
매끄러운 사상이 국소 비순환임을 보여줍니다.
$\overline{x}$를 이미지가 있는 $X$의 기하점으로 설정한다.
$\overline{s} = f(\overline{x})$. 허락하다
$\overline{t}$는 다음의 기하점이다.
$\Spec(\mathcal{O}^{sh}_{S, f(\overline{x})})$.
우리는 환 맵의 성질을 증명하려고 하기 때문에
$\mathcal{O}^{sh}_{S, \overline{s}} \to \mathcal{O}^{sh}_{X, \overline{x}}$
(정의 \ref{etale-cohomology-definition-locally-acyclic} 다음 토론 참조)
기본 변경으로 $f : X \to S$을 대체할 수 있다.
$X \times_S \Spec(\mathcal{O}^{sh}_{S, \overline{s}}) \to
\Spec(\mathcal{O}^{sh}_{S, \overline{s}})$.
따라서 우리는 $S$가 엄밀히 말하면
헨셀리안 국소환 그리고 그 $\overline{s}$는 위에 놓여 있습니다
$S$의 폐점.

\medskip\noindent 보조정리 \ref{etale-cohomology-lemma-base-change-f-star-general-stalks}를 다이어그램 $$
\xymatrix{
X \ar[d]_f & X_{\overline{t}} \ar[l]^h \ar[d]^e \\
S & \overline{t} \ar[l]_g
}
$$에 적용하고 층 $\mathcal{F} = \underline{M}$를 적용한다. 여기서 $M = \mathbf{Z}/n\mathbf{Z}$는 일부 정수에 대해 적용된다. $n$ $\overline{x}$의 잔기 체의 표수를 프라임한다. 우리는 $f^{-1}R^qg_*\mathcal{F} \to R^qh_*e^{-1}\mathcal{F}$ 맵이 매끄러운 기저변환에 의한 동형사상라는 것을 알고 있다. 정리 \ref{etale-cohomology-theorem-smooth-base-change} ($n$를 선택하여 토션 홀드의 가정을 참조하라.) 따라서 보조정리 \ref{etale-cohomology-lemma-base-change-f-star-general-stalks}는 $$
H^q(F_{\overline{x}, \overline{t}}, \underline{M}) =
H^q(\Spec(\mathcal{O}^{sh}_{X, \overline{x}}) \times_S \overline{t},
\underline{M})
=
H^q(\Spec(\mathcal{O}^{sh}_{S, \overline{s}}) \times_S \overline{t},
\underline{M}) =
H^q(\overline{t}, \underline{M})
$$에서 중간 등식을 제공한다. 외부 두 등식에 대해서는 $S = \Spec(\mathcal{O}^{sh}_{S, \overline{s}})$을 사용한다. $\overline{t}$는 분리 가능하게 닫힌 체의 스펙트럼이므로 우리는 이것이 우리가 보여야 했던 $$
H^q(F_{\overline{x}, \overline{t}}, \underline{M}) =
\left\{
\begin{matrix}
M & \text{만약에} & q = 0 \\
0 & \text{만약에} & q \not = 0
\end{matrix}
\right.
$$라고 결론을 내립니다(정의 \ref{etale-cohomology-definition-locally-acyclic} 다음 논의 참조). \end{proof}

\begin{lemma}
\label{etale-cohomology-lemma-locally-acyclic-locally-constant}
$f : X \to S$를 스킴의 사상으로 설정한다. $\mathcal{F}$를
$X_\etale$에 대한 지역적 상수 아벨 층
가리키다$\overline{x}$~의$X$그만큼아벨 군
$\mathcal{F}_{\overline{x}}$은 모든 요소가 다음을 갖는 비틀림 그룹이다.
의 표수에 소수를 지정한다.
$\overline{x}$의 잔기 체. $f$가 국소 비순환인 경우 $f$는
국소 비순환은 $\mathcal{F}$에 상대적이다.
\end{lemma}

\begin{proof}
즉, $\overline{x}$를 $X$의 기하점라고 하자.
$\mathcal{F}$는 지역적으로 일정하므로 다음을 알 수 있다.
$\mathcal{F}$를 $\Spec(\mathcal{O}^{sh}_{X, \overline{x}})$로 제한
는 상수층 $\underline{M}$와 동형이다.
$M = \mathcal{F}_{\overline{x}}$. 가정으로 다음과 같이 쓸 수 있다.
$M = \colim M_i$ 유한 아벨 군의 여과 여극한으로
$M_i$의 잔기 체의 표수에 소수의
$\overline{x}$. 기하점 $\overline{t}$을 고려해보세요.
$\Spec(\mathcal{O}^{sh}_{S, f(\overline{x})})$.
$F_{\overline{x}, \overline{t}}$는 유사하므로,
$$
H^q(F_{\overline{x}, \overline{t}}, \underline{M}) =
\colim H^q(F_{\overline{x}, \overline{t}}, \underline{M_i})
$$
보조정리 \ref{etale-cohomology-lemma-colimit} 기준.
각 $i$에 대해 $M_i = \bigoplus \mathbf{Z}/n_{i, j}\mathbf{Z}$라고 쓸 수 있다.
일부 정수에 대해 유한 직합으로 소수 $n_{i, j}$
$\overline{x}$의 잔기 체의 표수.
$f$는 국소 비순환이므로 다음을 알 수 있다.
$$
H^q(F_{\overline{x}, \overline{t}},
\underline{\mathbf{Z}/n_{i, j}\mathbf{Z}}) =
\left\{
\begin{matrix}
\mathbf{Z}/n_{i, j}\mathbf{Z} & \text{만약에} & q = 0 \\
0 & \text{만약에} & q \not = 0
\end{matrix}
\right.
$$
정의 \ref{etale-cohomology-definition-locally-acyclic} 다음 토론을 참조하라.
직합 및 여극한을 사용하여 다음과 같은 결론을 내립니다.
$$
H^q(F_{\overline{x}, \overline{t}}, \underline{M}) =
\left\{
\begin{matrix}
M & \text{만약에} & q = 0 \\
0 & \text{만약에} & q \not = 0
\end{matrix}
\right.
$$
그리고 우리는 승리한다.
\end{proof}

\begin{lemma}
\label{etale-cohomology-lemma-locally-acyclic-quasi-finite-base-change}
허락하다
$$
\xymatrix{
X' \ar[r]_{g'} \ar[d]_{f'} & X \ar[d]^f \\
S' \ar[r]^g & S
}
$$
스킴의 데카르트 다이어그램이어야 한다. $K$를 $D(X_\etale)$의 객체로 둡니다.
$\overline{x}'$를 $X'$ 이미지 $\overline{x}$의 기하점으로 설정한다.
$X$에서. 만약에
\begin{enumerate}
\item $f$~이다국소 비순환~에$\overline{x}$상대적인$K$그리고
\item $g$는 국지적으로 준유한이거나 $S' = \lim S_i$이다.
$S$에 대한 스킴의 방향성 역극한이다.
아핀 전이 사상 또는 $g : S' \to S$가 정역이다.
\end{enumerate}
그 다음에$f'$ 국소 비순환~에$\overline{x}'$상대적인$(g')^{-1}K$.
\end{lemma}

\begin{proof}
$\overline{s}'$ 및 $\overline{s}$ 이미지를 나타냅니다.
$\overline{x}'$ 및 $\overline{x}$의 $S'$ 및 $S$.
$\overline{t}'$를 스펙트럼의 기하점으로 설정한다.
$\Spec(\mathcal{O}^{sh}_{S', \overline{s}'})$ 및 표시
$\overline{t}$ $\Spec(\mathcal{O}^{sh}_{S, \overline{s}})$의 이미지이다.
대수학으로 보조정리
\ref{algebra-lemma-base-change-strict-henselization-quasi-finite}
그리고 $g$에 대한 가정은 다음과 같습니다.
$$
\mathcal{O}^{sh}_{X, \overline{x}}
\otimes_{\mathcal{O}^{sh}_{S, \overline{s}}}
\mathcal{O}^{sh}_{S', \overline{s}'}
\longrightarrow
\mathcal{O}^{sh}_{X', \overline{x}'}
$$
동형사상이다. 우리의 관례에 따르면
$\kappa(\overline{t}) = \kappa(\overline{t}')$
우리는 다음과 같이 결론을 내립니다.
$$
F_{\overline{x}', \overline{t}'} =
\Spec\left(
\mathcal{O}^{sh}_{X', \overline{x}'}
\otimes_{\mathcal{O}^{sh}_{S', \overline{s}'}}
\kappa(\overline{t}')\right) =
\Spec\left(
\mathcal{O}^{sh}_{X, \overline{x}}
\otimes_{\mathcal{O}^{sh}_{S, \overline{s}}}
\kappa(\overline{t})\right) =
F_{\overline{x}, \overline{t}}
$$
즉, $f'$의 $\overline{x}'$에서의 소멸 사이클은
$f$의 $\overline{x}$에서의 소멸 사이클의 한 예이다. 보조정리는
이것과 정의에서 즉시 따른다.
\end{proof}










\section{쌍대특수화 사상} \label{etale-cohomology-section-cospecialization}

\noindent
$f : X \to S$를 스킴의 사상으로 설정한다. $\overline{x}$를
기하점~의$X$이미지 포함$\overline{s} = f(\overline{x})$~에$S$.
$\overline{t}$를 기하점으로 설정한다.
$\Spec(\mathcal{O}^{sh}_{S, \overline{s}})$.
$K \in D(X_\etale)$로 놔두세요. 스킴의 사상 $g : Y \to X$에 대해
우리는 쓴다$K|_Y$대신에$g_{small}^{-1}K$그리고$R\Gamma(Y, K)$
$R\Gamma(Y_\etale, g_{small}^{-1}K)$ 대신.
우리는 주장 만약
\begin{enumerate}
\item $K$는 아래로 제한된다. 즉, $K \in D^+(X_\etale)$이다.
\item $f$는 $K$를 기준으로 국소 비순환이다.
\end{enumerate}
그러면 {\it 쌍대특수화 사상}가 있다.
$$
cosp :
R\Gamma(X_{\overline{t}}, K)
\longrightarrow
R\Gamma(X_{\overline{s}}, K)
$$
이는 특수화 지도와 밀접한 관련이 있다.
절 \ref{etale-cohomology-section-specialization}에서 고려되며 특히
비고 \ref{etale-cohomology-remark-specialization-map-and-fibres}.

\medskip\noindent
지도를 구성하려면 사상을 고려한다.
$$
X_{\overline{t}}
\xrightarrow{h}
X \times_S \Spec(\mathcal{O}^{sh}_{S, \overline{s}})
\xleftarrow{i}
X_{\overline{s}}
$$
$h^{-1}$와 $Rh_*$ 사이의 부속 단위는 지도를 제공한다.
$$
\beta_{K, \overline{s}, \overline{t}} :
K|_{X \times_S \Spec(\mathcal{O}^{sh}_{S, \overline{s}})}
\longrightarrow
Rh_*(K|_{X_{\overline{t}}})
$$
$D((X \times_S \Spec(\mathcal{O}^{sh}_{S, \overline{s}}))_\etale)$에서.
보조정리 \ref{etale-cohomology-lemma-beta-is-isomorphism} 아래는 하락세를 보여줍니다.
$i^{-1}\beta_{K, \overline{s}, \overline{t}}$
위의 가정에 따르면 동형사상이다.
따라서 쌍대특수화 사상을 구성으로 정의할 수 있다.
\begin{align*}
R\Gamma(X_{\overline{t}}, K)
& =
R\Gamma(X \times_S \Spec(\mathcal{O}^{sh}_{S, \overline{s}}),
Rh_*(K|_{X_{\overline{t}}})) \\
&
\xrightarrow{i^{-1}}
R\Gamma(X_{\overline{s}}, i^{-1}Rh_*(K|_{X_{\overline{t}}})) \\
&
\xrightarrow{(i^{-1}\beta_{K, \overline{s}, \overline{t}})^{-1}}
R\Gamma(X_{\overline{s}},
i^{-1}(K|_{X \times_S \Spec(\mathcal{O}^{sh}_{S, \overline{s}})})) \\
& =
R\Gamma(X_{\overline{s}}, K)
\end{align*}

\begin{lemma} \label{etale-cohomology-lemma-beta-is-isomorphism} 지도 $i^{-1}\beta_{K, \overline{s}, \overline{t}}$는 동형사상이다. \end{lemma}

\begin{proof}
지도 구성 $h$, $i$, $\beta_{K, \overline{s}, \overline{t}}$
에만 의존한다
$X$ 및 $K$을 $\Spec(\mathcal{O}^{sh}_{S, \overline{s}})$로 기본 변경한다.
따라서 우리는 $S$가 엄격하게 헨셀적인 스킴라고 가정할 수 있고 실제로 가정한다.
닫힌 지점 $\overline{s}$을 사용한다. 국소 비순환성
~의$f$상대적인$K$이 기본 변경으로 보존됩니다(예:예
보조정리 \ref{etale-cohomology-lemma-locally-acyclic-quasi-finite-base-change} 또는 그냥
이 매우 특별한 경우에는 Hensel 환을 엄격하게 비교하여 직접적으로 비교한다.

\medskip\noindent
$\overline{x}$를 $X_{\overline{s}}$의 기하점으로 설정한다.
또는 동등하게 $\overline{x}$를 기하점으로 설정한다.
$f$에 의한 것은 $\overline{s}$이다. 줄기를 계산해 보자.
$i^{-1}\beta_{K, \overline{s}, \overline{t}}$
$\overline{x}$에 있다. 첫째, 우리는
$$
(i^{-1}\beta_{K, \overline{s}, \overline{t}})_{\overline{x}} =
(\beta_{K, \overline{s}, \overline{t}})_{\overline{x}}
$$
당김은 줄기를 유지하므로 보조정리 \ref{etale-cohomology-lemma-stalk-pullback}를 참조하라.
$S = \Spec(\mathcal{O}^{sh}_{S, \overline{s}})$ 상황에 처해있으니
$h : X_{\overline{t}} \to X$에는 성질이 있다.
$X_{\overline{t}} \times_X \Spec(\mathcal{O}^{sh}_{X, \overline{x}})
= F_{\overline{x}, \overline{t}}$. 따라서 우리는
$$
(\beta_{K, \overline{s}, \overline{t}})_{\overline{x}} :
K_{\overline{x}}
\longrightarrow
Rh_*(K|_{X_{\overline{t}}})_{\overline{x}} =
R\Gamma(F_{\overline{x}, \overline{t}}, K)
$$
여기서 등호는 정리 \ref{etale-cohomology-theorem-higher-direct-images}이다.
지도 $(\beta_{K, \overline{s}, \overline{t}})_{\overline{x}}$
다름 아닌 지도 $\alpha_{K, \overline{x}, \overline{t}}$가 사용된 것입니다
정의 \ref{etale-cohomology-definition-locally-acyclic}에서. 결과는 다음과 같습니다.
지도가 줄기의 동형사상인지 확인하려면 다음을 수행하세요.
정리 \ref{etale-cohomology-theorem-exactness-stalks}.
\end{proof}

\noindent 쌍대특수화 사상이 존재할 때 특수화 맵의 반대가 되려고 한다.

\begin{lemma} \label{etale-cohomology-lemma-specialization-cospecialization} 위의 상황에서 추가로 $f$가 준 컴팩트하고 준 분리된 경우 다이어그램 $$
\xymatrix{
(Rf_*K)_{\overline{s}} \ar[r] \ar[d]_{sp} &
R\Gamma(X_{\overline{s}}, K) \\
(Rf_*K)_{\overline{t}} \ar[r] &
R\Gamma(X_{\overline{t}}, K) \ar[u]_{cosp}
}
$$은 교환 가능한다. \end{lemma}

\begin{proof}
보조정리 \ref{etale-cohomology-lemma-beta-is-isomorphism}의 증명과 같습니다.
$S$을 $\Spec(\mathcal{O}^{sh}_{S, \overline{s}})$로 바꿀 수 있다.
그런 다음 지도는 $h : X_{\overline{t}} \to X$로 단순화된다.
$i : X_{\overline{s}} \to X$ 그리고
$\beta_{K, \overline{s}, \overline{t}} : K \to Rh_*(K|_{X_{\overline{t}}})$.
$(Rf_*K)_{\overline{s}} = R\Gamma(X, K)$을 사용하여
정리 \ref{etale-cohomology-theorem-higher-direct-images}
$sp$와 베이스 변경 맵의 구성
$(Rf_*K)_{\overline{t}} \to R\Gamma(X_{\overline{t}}, K)$
$h$에 따른 코호몰로지의 당김이다.
지도랑 똑같네요
$$
R\Gamma(X, K) \xrightarrow{\beta_{K, \overline{s}, \overline{t}}}
R\Gamma(X, Rh_*(K|_{X_{\overline{t}}})) =
R\Gamma(X_{\overline{t}}, K)
$$
이제 지도 $cosp$는 먼저 $=$ 기호를 반전시킵니다.
수식을 사용한 다음 $i$를 따라 뒤로 당겨서 마지막으로 적용한다.
$i^{-1}\beta_{K, \overline{s}, \overline{t}}$의 반대이다.
따라서 우리는 원하는 교환성을 얻습니다.
\end{proof}

\begin{lemma}
\label{etale-cohomology-lemma-sp-isom-proper-torsion-loc-ac}
$f : X \to S$를 스킴의 사상으로 설정한다. $K \in D(X_\etale)$로 놔두세요.
추정하다
\begin{enumerate}
\item $K$는 아래로 제한된다. 즉, $K \in D^+(X_\etale)$이다.
\item $f$는 $K$를 기준으로 국소 비순환이다.
\item $f$가 적절하고
\item $K$에는 비틀림 코호몰로지 층이 있다.
\end{enumerate}
그러면 매기하점 $\overline{s}$~의$S$그리고 모든 기하학적
가리키다$\overline{t}$~의$\Spec(\mathcal{O}^{sh}_{S, \overline{s}})$
특수화 지도 모두
$sp : (Rf_*K)_{\overline{s}} \to (Rf_*K)_{\overline{t}}$
그리고 쌍대특수화 사상
$cosp : R\Gamma(X_{\overline{t}}, K) \to R\Gamma(X_{\overline{s}}, K)$
동형사상이다.
\end{lemma}

\begin{proof}
고유 기저변환 정리에 의해 (다음과 같은 형식으로)
보조정리 \ref{etale-cohomology-lemma-proper-base-change-stalk}) 우리는
$(Rf_*K)_{\overline{s}} = R\Gamma(X_{\overline{s}}, K)$
$\overline{t}$의 경우에도 마찬가지이다. ``올바른'' 증명 은 다음과 같습니다.
의 주장을 보여주기 위해
보조정리 \ref{etale-cohomology-lemma-specialization-cospecialization}
이 경우 $sp$와 $cosp$는 역 동형사상임을 보여줍니다.
대신 $cosp$가 동형사상임을 직접 보여드리겠습니다.
위의 논의에서 $cosp$는 동형사상임을 알 수 있다.
필요충분조건은 $i$에 의한 철수
$$
R\Gamma(X \times_S \Spec(\mathcal{O}^{sh}_{S, \overline{s}}),
Rh_*(K|_{X_{\overline{t}}}))
\longrightarrow
R\Gamma(X_{\overline{s}}, i^{-1}Rh_*(K|_{X_{\overline{t}}}))
$$
$D^+(\textit{Ab})$의 동형사상이다. 이것은 적절한 사실입니다
고유 사상에 대한 기본 변경 정리
$f' : X \times_S \Spec(\mathcal{O}^{sh}_{S, \overline{s}}) \to
\Spec(\mathcal{O}^{sh}_{S, \overline{s}})$
사상 $\overline{s} \to \Spec(\mathcal{O}^{sh}_{S, \overline{s}})$에 의해
그리고 복합체 $K' = Rh_*(K|_{X_{\overline{t}}})$. 복합체
$K'$는 아래 경계를 가지며 비틀림이 있다. 코호몰로지 층
보조정리 \ref{etale-cohomology-lemma-torsion-direct-image} 기준.
$\Spec(\mathcal{O}^{sh}_{S, \overline{s}})$은 엄격하게 헨셀식이므로
$\overline{s}$이 닫힌 점 위에 놓여 있고,
표시된 화살표의 소스는 다음과 같습니다.
$(Rf'_*K')_{\overline{s}}$ 그리고 대상은 다음과 같습니다.
$R\Gamma(X_{\overline{s}}, K')$ 그리고 표시된 지도는 동형사상이다.
이미 사용된
보조정리 \ref{etale-cohomology-lemma-proper-base-change-stalk}.
따라서 우리는 다이어그램에 있는 네 개의 화살표 중 세 개를 볼 수 있다.
보조정리 \ref{etale-cohomology-lemma-specialization-cospecialization} 중 동형사상
그리고 우리는 결론을 내립니다.
\end{proof}

\begin{lemma}
\label{etale-cohomology-lemma-sp-isom-proper-loc-cst-torsion}
$f : X \to S$를 스킴의 사상으로 설정한다. $\mathcal{F}$
$X_\etale$에서 아벨 층이 되어야 한다. 추정하다
\begin{enumerate}
\item $f$ 부드럽고 적절함
\item $\mathcal{F}$는 지역적으로 일정하며,
\item $\mathcal{F}_{\overline{x}}$는 비틀림 그룹이다.
그 요소는 다음의 잔기 표수에 대해 소수 차수를 갖습니다.
$\overline{x}$모든기하점 $\overline{x}$~의$X$.
\end{enumerate}
그러면 매기하점 $\overline{s}$~의$S$그리고 모든 기하학적
가리키다$\overline{t}$~의$\Spec(\mathcal{O}^{sh}_{S, \overline{s}})$
특수화 지도
$sp : (Rf_*\mathcal{F})_{\overline{s}} \to (Rf_*\mathcal{F})_{\overline{t}}$
동형사상이다.
\end{lemma}

\begin{proof} 이는 기본정리 \ref{etale-cohomology-lemma-sp-isom-proper-torsion-loc-ac} 및 \ref{etale-cohomology-lemma-locally-acyclic-locally-constant} 및 명제 \ref{etale-cohomology-proposition-smooth-locally-acyclic}에서 따릅니다. \end{proof}

















\section{코호몰로지 차원}
\label{etale-cohomology-section-cd}

\noindent 절 \ref{etale-cohomology-section-vanishing-torsion}의 결과를 사용하여 스킴의 코호몰로지 차원과 체의 코호몰로지 차원에 대한 일부 경계를 추론할 수 있다. 정리 (명제 \ref{etale-cohomology-proposition-cd-affine}의 증명에서)를 변경한다.

\begin{definition}
\label{etale-cohomology-definition-cd}
$X$를 준간소화하고 준분리된 스킴라고 하자.
{\it 코호몰로지 차원의 $X$}가 가장 작습니다.
요소
$$
\text{cd}(X) \in \{0, 1, 2, \ldots\} \cup \{\infty\}
$$
모든 아벨에 대해 꼬임층 $\mathcal{F}$
$X_\etale$에 $H^i_\etale(X, \mathcal{F}) = 0$이 있다.
$i > \text{cd}(X)$의 경우. $X = \Spec(A)$ 우리가 때때로
이것을 $A$의 코호몰로지 차원라고 부릅니다.
\end{definition}

\noindent 스킴이 표수 $p$에 있는 경우 코호몰로지의 소멸에 대한 더 선명한 경계를 $p$-전력 꼬임층으로 얻을 수 있는 경우가 많습니다. 이 문제는 다른 곳에서 다룰 것입니다(여기에 향후 참조를 삽입하십시오).

\begin{lemma} \label{etale-cohomology-lemma-cd-limit} $X = \lim X_i$는 아핀 전이 사상을 갖는 준 컴팩트 및 준 분리 스킴 시스템의 방향성 극한라고 가정한다. 그런 다음 $\text{cd}(X) \leq \max \text{cd}(X_i)$. \end{lemma}

\begin{proof} 투영을 $f_i : X \to X_i$ 나타냅니다. $\mathcal{F}$를 $X_\etale$의 아벨 꼬임층으로 설정한다. 그러면 $\mathcal{F} = \lim f_i^{-1}f_{i, *}\mathcal{F}$ 에 의해 보조정리 \ref{etale-cohomology-lemma-linus-hamann}가 된다. 따라서 $H^q_\etale(X, \mathcal{F}) =
\colim H^q_\etale(X_i, f_{i, *}\mathcal{F})$는 정리 \ref{etale-cohomology-theorem-colimit}에 의해 이루어집니다. 보조정리가 이어집니다. \end{proof}

\begin{lemma}
\label{etale-cohomology-lemma-cd-curve-over-field}
$K$를 체로 설정한다. $X$를 $1$차원이라고 합시다.
$K$에 대한 유한 유형의 아핀 스킴이다. 그 다음에
$\text{cd}(X) \leq 1 + \text{cd}(K)$.
\end{lemma}

\begin{proof}
$\mathcal{F}$를 $X_\etale$의 아벨 꼬임층으로 설정한다.
사상에 대해 Leray 스펙트럼 열을 고려하라.
$f : X \to \Spec(K)$. 우리는 얻는다
$$
E_2^{p, q} = H^p(\Spec(K), R^qf_*\mathcal{F})
$$
$H^{p + q}_\etale(X, \mathcal{F})$로 수렴된다.
기하점에서 $R^qf_*\mathcal{F}$의 줄기
$\Spec(\overline{K}) \to \Spec(K)$는 코호몰로지이다.
$\mathcal{F}$을 $X_{\overline{K}}$로 철회한다.
따라서 $\geq 2$ 단위로 사라집니다.
정리 \ref{etale-cohomology-theorem-vanishing-affine-curves}.
\end{proof}

\begin{lemma} \label{etale-cohomology-lemma-cd-field-extension} $L/K$를 체 확장자로 둡니다. 그런 다음 $\text{cd}(L) \leq \text{cd}(K) + \text{trdeg}_K(L)$이 있다. \end{lemma}

\begin{proof}
$\text{trdeg}_K(L) = \infty$이면 이는 분명한다.
그렇지 않다면 일련의 확장을 찾을 수 있다.
$L= L_r/L_{r - 1}/ \ldots /L_1/L_0 = K$ 이렇게
$\text{trdeg}_{L_i}(L_{i + 1}) = 1$ 및 $r = \text{trdeg}_K(L)$.
따라서 $r = 1$인 경우에는 보조정리를 증명하면 충분한다.
이 경우 $L = \colim A_i$라고 쓸 수 있다.
유한 유형 $K$-부분대수학의 여과 여극한으로.
보조정리 \ref{etale-cohomology-lemma-cd-limit}에 의해 다음을 증명하는 것으로 충분한다.
$\text{cd}(A_i) \leq 1 + \text{cd}(K)$. 이는 다음과 같습니다
보조정리 \ref{etale-cohomology-lemma-cd-curve-over-field}에서.
\end{proof}

\begin{lemma}
\label{etale-cohomology-lemma-strictly-henselian}
$K$를 체로 설정한다. $X$를 $K$에 대해 유한 유형의 스킴으로 둡니다.
$x \in X$로 놔두세요. 집합 $a = \text{trdeg}_K(\kappa(x))$
및 $d = \dim_x(X)$. 그러면 지도가 있어요
$$
K(t_1, \ldots, t_a)^{sep} \longrightarrow \mathcal{O}_{X, x}^{sh}
$$
그렇게
\begin{enumerate}
\item $\mathcal{O}_{X, x}^{sh}$의 잔기 체는 완전히 분리할 수 없다.
$K(t_1, \ldots, t_a)^{sep}$ 확장자,
\item $\mathcal{O}_{X, x}^{sh}$는 유한의 여과 여극한이다.
유형$K(t_1, \ldots, t_a)^{sep}$-대수학차원 $\leq d - a$.
\end{enumerate}
\end{lemma}

\begin{proof}
우리는 $X$가 유사하다고 가정할 수 있다. Noether 정규화를 통해 가능하면
$X$를 다시 축소하면 유한한 사상을 선택할 수 있다.
$\pi : X \to \mathbf{A}^d_K$, 참조
대수학, 보조정리 \ref{algebra-lemma-Noether-normalization-at-point}.
$\kappa(x)$는 $\pi(x)$의 잔여물 체의 유한 확장이므로,
이 잔류물체초월등급이 있다$a$~ 위에$K$또한.
따라서 우리는 유한한 사상 $\pi' : \mathbf{A}^d_K \to \mathbf{A}^d_K$를 찾을 수 있다.
$\pi'(\pi(x))$는 선형의 일반 점에 해당한다.
설정에 의해 주어진 부분 공간 $\mathbf{A}^a_K \subset \mathbf{A}^d_K$
마지막 $d - a$ 좌표는 0과 같습니다. 따라서 구성
$$
X \xrightarrow{\pi' \circ \pi} \mathbf{A}^d_K \xrightarrow{p} \mathbf{A}^a_K
$$
$\pi' \circ \pi$ 및 $p$을 첫 번째 $a$ 좌표에 투영
$x$을 일반 포인트 $\eta \in \mathbf{A}^a_K$에 매핑한다. 유도된 지도
$$
K(t_1, \ldots, t_a)^{sep} =
\mathcal{O}_{\mathbf{A}^a_k, \eta}^{sh}
\longrightarrow \mathcal{O}_{X, x}^{sh}
$$
위의 에탈 국소환은 잔차 체가 분명하므로 (1)를 만족한다.
의 $\mathcal{O}_{X, x}^{sh}$ 은 분리 가능하게 닫혀 있는 식의 대수적 확장이다.
체 $K(t_1, \ldots, t_a)^{sep}$. 반면, $X = \Spec(B)$인 경우,
그러면 $\mathcal{O}_{X, x}^{sh} = \colim B_j$는 여과 여극한이다.
에탈 $B$-대수학 $B_j$. $B_j$가 준유한인 것을 관찰하세요.
$K[t_1, \ldots, t_d]$는 $B$가 $K[t_1, \ldots, t_d]$에 대해 유한하므로
우리도 비슷하게 쓸 수 있습니다
$K(t_1, \ldots, t_a)^{sep} = \colim A_i$를 여과 여극한으로
에탈 $K[t_1, \ldots, t_a]$-대수학. 모든 $i$에 대해 우리는
다음과 같은 $j$를 찾으세요.
$A_i \to K(t_1, \ldots, t_a)^{sep} \to \mathcal{O}_{X, x}^{sh}$
지도 $\psi_{i, j} : A_i \to B_j$를 통해 요인을 분석한다. 그러면 $B_j$는 준 유한이다.
$A_i[t_{a + 1}, \ldots, t_d]$ 이상. 따라서
$$
B_{i, j} = B_j \otimes_{\psi_{i, j}, A_i} K(t_1, \ldots, t_a)^{sep}
$$
준 유한이므로 차원 $\leq d - a$이 있다.
$K(t_1, \ldots, t_a)^{sep}[t_{a + 1}, \ldots, t_d]$.
(2)의 증명은 이제 $\mathcal{O}_{X, x}^{sh}$로 완료되었다.
여과 여극한\footnote{$R$를 환으로 둡니다.
$A = \colim_{i \in I} A_i$
유한하게 제시된 $R$-대수의 여과 여극한이다.
$B = \colim_{j \in J} B_j$를 $R$-대수학의 여과 여극한으로 설정한다.
$A \to B$를 $R$ 대수 맵으로 설정한다.
모든 $i \in I$에 대해 $j \in J$ 및 $R$ 대수 맵이 있다고 가정한다.
$\psi_{i, j} : A_i \to B_j$.
다음과 같은 경우 $(i', j', \psi_{i', j'}) \geq (i, j, \psi_{i, j})$라고 말하세요.
$i' \geq i$, $j' \geq j$ 및 $\psi_{i, j}$ 및 $\psi_{i', j'}$
호환된다. 그런 다음 트리플의 컬렉션은 다음을 형성한다.
집합 및 $B = \colim B_j \otimes_{\psi_{i, j} A_i} A$를 지시한다.}
대수학 $B_{i, j}$. 일부 세부 사항은 생략되었다.
\end{proof}

\begin{proposition} \label{etale-cohomology-proposition-cd-affine} $K$를 체로 둡니다. $X$를 $K$ 위에 유한 유형의 아핀 스킴으로 둡니다. 그런 다음 $\text{cd}(X) \leq \dim(X) + \text{cd}(K)$가 있다. \end{proposition}

\begin{proof} 이를 $\dim(X)$에 대한 귀납법으로 증명하겠습니다. $\mathcal{F}$를 $X_\etale$의 아벨 꼬임층으로 설정한다.

\medskip\noindent
$\dim(X) = 0$ 사례이다. 이 경우 구조는 사상
$f : X \to \Spec(K)$은 유한한다. 따라서 우리는 $R^if_*\mathcal{F} = 0$
$i > 0$에 대해서는 명제 \ref{etale-cohomology-proposition-finite-higher-direct-image-zero}를 참조하라.
따라서 $H^i_\etale(X, \mathcal{F}) = H^i_\etale(\Spec(K), f_*\mathcal{F})$
$f$에 대한 Leray의 스펙트럼 열
(코호몰로지 사이트, 보조정리 \ref{sites-cohomology-lemma-Leray})
결과는 분명한다.

\medskip\noindent 사례 $\dim(X) = 1$이다. 이것은 보조정리 \ref{etale-cohomology-lemma-cd-curve-over-field}이다.

\medskip\noindent
$d = \dim(X) > 1$ 및 명제가 유한 유형에 대해 유지된다고 가정한다.
아핀 스킴의 차원 $< d$의 체.
Noether 정규화에 대해서는 예를 참조하라.
다양체, 보조정리 \ref{varieties-lemma-noether-normalization-affine},
유한한 사상 $f : X \to \mathbf{A}^d_K$가 존재한다.
그것을 기억해$R^if_*\mathcal{F} = 0$~을 위한$i > 0$~에 의해
명제 \ref{etale-cohomology-proposition-finite-higher-direct-image-zero}.
$f$에 대한 Leray의 스펙트럼 열 작성
(코호몰로지 사이트, 보조정리 \ref{sites-cohomology-lemma-Leray})
$\pi_*\mathcal{F}$에 대한 결과를 증명하는 것으로 충분하다고 결론을 내렸습니다.
$\mathbf{A}^d_K$에 있다.

\medskip\noindent
막간 I. $j : X \to Y$를 부드러운 열린 몰입으로 설정한다.
$d$차원 다양체 위의 $K$ (반드시 유사하지는 않음)
그 보수는 유효 까르띠에 제수 $D$의 지지집합이다.
그만큼층 $R^qj_*\mathcal{F}$~을 위한$q > 0$지원됩니다$D$.
우리는 주장 $(R^qj_*\mathcal{F})_{\overline{y}} = 0$
$a = \text{trdeg}_K(\kappa(y)) > d - q$의 경우. 즉,
정리 \ref{etale-cohomology-theorem-higher-direct-images} 우리는
$$
(R^qj_*\mathcal{F})_{\overline{y}} =
H^q(\Spec(\mathcal{O}_{Y, y}^{sh}) \times_Y X, \mathcal{F})
$$
지역 방정식 선택 $f \in \mathfrak m_y \subset \mathcal{O}_{Y, y}$
$D$의 경우. 그럼 우리는
$$
\Spec(\mathcal{O}_{Y, y}^{sh}) \times_Y X =
\Spec(\mathcal{O}_{Y, y}^{sh}[1/f])
$$
보조정리 \ref{etale-cohomology-lemma-strictly-henselian}를 사용하여 임베딩을 얻습니다.
$$
K(t_1, \ldots, t_a)^{sep}(x) =
K(t_1, \ldots, t_a)^{sep}[x]_{(x)}[1/x]
\longrightarrow
\mathcal{O}_{Y, y}^{sh}[1/f]
$$
분수 체의 $K$에 대한 초월도가
$\mathcal{O}_{Y, y}^{sh}$는 $d$이며, $\mathcal{O}_{Y, y}^{sh}[1/f]$임을 알 수 있다.
는 $(d - a - 1)$차원 유한형 대수 중 여과 여극한이다.
체 $K(t_1, \ldots, t_a)^{sep}(x)$ 그 자체는 동질성을 갖는다.
차원 $1$ 에 의해 보조정리 \ref{etale-cohomology-lemma-cd-field-extension}. 그래서 인덕션으로
가설과 보조정리 \ref{etale-cohomology-lemma-cd-limit}
원하는 소멸을 얻습니다.

\medskip\noindent
막간 II. 허락하다$Z$순조롭게 지내다다양체~ 위에$K$~의차원 $d - 1$.
$E_a \subset Z$를 $z \in Z$ 점의 집합으로 설정한다.
$\text{trdeg}_K(\kappa(z)) \leq a$. $E_a$이 닫혀 있는지 확인하세요.
특수화 아래를 참조하라.
다양체, 보조정리 \ref{varieties-lemma-dimension-locally-algebraic}.
$\mathcal{G}$가 $Z$의 비틀림 아벨 층이라고 가정한다.
지지집합이 $E_a$에 포함되어 있다. 그러면 우리는 주장
$H^b_\etale(Z, \mathcal{G}) = 0$~을 위한$b > a + \text{cd}(K)$.
즉, $\mathcal{G} = \colim \mathcal{G}_i$라고 쓸 수 있다.
$\mathcal{G}_i$ 비틀림 아벨 층
폐쇄형 하위 체계에서 지원됨$Z_i$에 포함된$E_a$, 보다
보조정리 \ref{etale-cohomology-lemma-support-in-subset}.
그러면 귀납 가설이 시작되어 다음을 암시한다.
$\mathcal{G}_i$\footnote{에 대해 원하는 소멸은 여기
먼저 명제 \ref{etale-cohomology-proposition-closed-immersion-pushforward}를 사용한다.
$\mathcal{G}_i$을 $Z_i$에 층의 직상으로 쓰려면,
귀납 가설은 $Z_i$에서 이 층에 대해 소멸을 제공하며,
$Z_i \to Z$에 대한 Leray 스펙트럼 열은 소멸을 제공한다.
$\mathcal{G}_i$.}의 경우. 마지막으로 우리는
정리 \ref{etale-cohomology-theorem-colimit}로 결론을 내립니다.

\medskip\noindent
가환 도식을 고려해보세요.
$$
\xymatrix{
\mathbf{A}^d_K \ar[rd]_f \ar[rr]_-j & &
\mathbf{P}^1_K \times_K \mathbf{A}^{d - 1}_K \ar[ld]^g \\
& \mathbf{A}^{d - 1}_K
}
$$
$j$는 부드러운 $d$차원의 열린 몰입임을 확인하세요.
다양체 그 보수는 효과적인 까르띠에 제수 $D$이다. 따라서
우리는 막간 I에서 얻은 결과를 사용할 수 있다.
우리는 친척 Leray 스펙트럼 열을 연구하려고 한다.
$$
E_2^{p, q} = R^pg_*R^qj_*\mathcal{F} \Rightarrow R^{p + q}f_*\mathcal{F}
$$
부터$R^qj_*\mathcal{F}$~을 위한$q > 0$에서 지원됩니다$D$그리고 그 이후로
$g|_D : D \to \mathbf{A}^{d - 1}_K$는 동형사상이다.
$R^pg_*R^qj_*\mathcal{F} = 0$~을 위한$p > 0$그리고$q > 0$. 더욱이, 우리는
$R^qj_*\mathcal{F} = 0$~을 위한$q > d$. 반면에,$g$는
고유 사상 상대 차원 $1$. 따라서
보조정리 \ref{etale-cohomology-lemma-cohomological-dimension-proper}
우리는 그것을 본다$R^pg_*j_*\mathcal{F} = 0$~을 위한$p > 2$. 따라서$E_2$-페이지
스펙트럼 열의 모양은 다음과 같습니다
$$
\begin{matrix}
g_*R^dj_*\mathcal{F} & 0 & 0 \\
\ldots & \ldots & \ldots \\
g_*R^2j_*\mathcal{F} & 0 & 0 \\
g_*R^1j_*\mathcal{F} & 0 & 0 \\
g_*j_*\mathcal{F} & R^1g_*j_*\mathcal{F} & R^2g_*j_*\mathcal{F}
\end{matrix}
$$
우리는 다음과 같이 결론을 내립니다.$R^qf_*\mathcal{F} = g_*R^qj_*\mathcal{F}$~을 위한$q > 2$.
막간에서 우리는지지집합~의$R^qf_*\mathcal{F}$~을 위한$q > 2$
$\mathbf{A}^{d - 1}_K$ 포인트의 집합에 포함되어 있다.
그 잔기 체는 초월 등급 $\leq d - q$을 갖습니다.
막간 II로
$$
H^p(\mathbf{A}^{d - 1}_K, R^qf_*\mathcal{F}) = 0
\text{~을 위한}p > d - q + \text{cd}(K)\text{ 및 }q > 2
$$
한편, 정리 \ref{etale-cohomology-theorem-higher-direct-images}에 의해
우리에겐 $R^2f_*\mathcal{F}_{\overline{\eta}} =
H^2(\mathbf{A}^1_{\overline{\eta}}, \mathcal{F}) = 0$
(소멸 차원 $1$의 경우 $\eta$는
$\mathbf{A}^{d - 1}_K$의 일반 포인트이다.
그러므로 막간 II에서 다시 우리는 본다.
$$
H^p(\mathbf{A}^{d - 1}_K, R^2f_*\mathcal{F}) = 0
\text{~을 위한}p > d - 2 + \text{cd}(K)
$$
마지막으로, 우리는
$$
H^p(\mathbf{A}^{d - 1}_K, R^qf_*\mathcal{F}) = 0
\text{~을 위한}p > d - 1 + \text{cd}(K)\text{ 및 }q = 0, 1
$$
유도 가설에 의해. 방금 말한 모든 내용을 종합해 보면
르레이와 함께 스펙트럼 열
$H^p(\mathbf{A}^{d - 1}_K, R^qf_*\mathcal{F}) \Rightarrow
H^{p + q}(\mathbf{A}^d_K, \mathcal{F})$ 결론을 내립니다.
\end{proof}

\begin{lemma}
\label{etale-cohomology-lemma-interlude-II}
$K$를 체로 설정한다. $X$를 $K$ 위에 유한 유형의 아핀 스킴으로 둡니다.
$E_a \subset X$를 포인트의 집합으로 설정한다.
$x \in X$와 $\text{trdeg}_K(\kappa(x)) \leq a$.
$\mathcal{F}$를 $X_\etale$에서 아벨 꼬임층으로 설정한다.
지지집합이 $E_a$에 포함되어 있다. 그 다음에
$H^b_\etale(X, \mathcal{F}) = 0$~을 위한$b > a + \text{cd}(K)$.
\end{lemma}

\begin{proof}
$\mathcal{F} = \colim \mathcal{F}_i$라고 쓸 수 있다.
$\mathcal{F}_i$ 비틀림 아벨 층
폐쇄형 하위 체계에서 지원됨$Z_i$에 포함된$E_a$, 보다
보조정리 \ref{etale-cohomology-lemma-support-in-subset}.
그런 다음 명제 \ref{etale-cohomology-proposition-cd-affine}는 다음을 제공한다.
$\mathcal{F}_i$에 대해 원하는 소멸. 자세한 내용은 생략했다.
힌트: 먼저 명제 \ref{etale-cohomology-proposition-closed-immersion-pushforward}를 사용하세요.
$\mathcal{F}_i$을 $Z_i$에 층의 직상으로 쓰려면,
$Z_i$에서 이 층에 대해 소멸을 사용하고
$Z_i \to Z$에 대한 Leray 스펙트럼 열을 사용하여 소멸을 얻습니다.
$\mathcal{F}_i$의 경우. 마지막으로 우리는
정리 \ref{etale-cohomology-theorem-colimit}로 결론을 내립니다.
\end{proof}

\begin{lemma}
\label{etale-cohomology-lemma-interlude-I}
$f : X \to Y$를 유한의 스킴의 아핀 사상으로 설정한다.
체 $K$ 위에 입력하세요. $E_a(X)$를 점 $x \in X$의 집합으로 설정한다.
$\text{trdeg}_K(\kappa(x)) \leq a$로.
$\mathcal{F}$를 $X_\etale$에서 아벨 꼬임층으로 설정한다.
지지집합이 $E_a$에 포함되어 있다. 그 다음에
$R^qf_*\mathcal{F}$에는 지지집합이 포함되어 있다.
$E_{a - q}(Y)$.
\end{lemma}

\begin{proof}
질문은 $Y$에 대한 로컬이므로 $Y$이 유사하다고 가정할 수 있다.
그러면 $X$도 유사하므로 다이어그램을 선택할 수 있다.
$$
\xymatrix{
X \ar[d]_f \ar[r]_i & \mathbf{A}^{n + m}_K \ar[d]^{\text{pr}} \\
Y \ar[r]^j & \mathbf{A}^n_K
}
$$
여기서 수평 화살표는 닫힌 몰입이고 수직 화살표는
오른쪽 화살표는 투영입니다(자세한 내용은 생략).
그런 다음 $j_*R^qf_*\mathcal{F} = R^q\text{pr}_*i_*\mathcal{F}$
$i$ 및 $j$의 고차 직상의 소멸에 의해, 참조
명제 \ref{etale-cohomology-proposition-finite-higher-direct-image-zero}.
또한, 명제의 $j_*$의 줄기에 대한 설명은
$j_*R^qf_*\mathcal{F}$에 대해 소멸을 증명하면 충분함을 보여줍니다.
따라서 우리는 $f$이 투영이라고 가정할 수 있다.
사상 $\text{pr} : \mathbf{A}^{n + m}_K \to \mathbf{A}^n_K$
그리고 아벨리안꼬임층 $\mathcal{F}$~에$\mathbf{A}^{n + m}_K$
보조정리의 문에서 가정을 만족한다.

\medskip\noindent
$y$를 $\mathbf{A}^n_K$의 한 점으로 둡니다.
정리 \ref{etale-cohomology-theorem-higher-direct-images}에 의해 우리는
$$
(R^q\text{pr}_*\mathcal{F})_{\overline{y}} =
H^q(\mathbf{A}^{n + m}_K \times_{A^n_K} \Spec(\mathcal{O}_{Y, y}^{sh}),
\mathcal{F}) =
H^q(\mathbf{A}^m_{\mathcal{O}_{Y, y}^{sh}}, \mathcal{F})
$$
$b = \text{trdeg}_K(\kappa(y))$라고 말하세요. 보조정리 \ref{etale-cohomology-lemma-strictly-henselian}에서
우리는 임베딩을 얻습니다
$$
L = K(t_1, \ldots, t_b)^{sep} \longrightarrow \mathcal{O}_{Y, y}^{sh}
$$
$\mathcal{O}_{Y, y}^{sh} = \colim B_i$을 필터링된 것으로 씁니다.
여극한 유한 유형 $L$-부분대수 $B_i \subset \mathcal{O}_{Y, y}^{sh}$
일반 함수의 환 $K[T_1, \ldots, T_n]$를 포함한다.
$\mathbf{A}^n_K$에 있다. 그러면 우리는 얻는다
$$
\mathbf{A}^m_{\mathcal{O}_{Y, y}^{sh}} =
\lim \mathbf{A}^m_{B_i}
$$
$z \in \mathbf{A}^m_{B_i}$가 지지집합에 있는 지점인 경우
$\mathcal{F}$, 이미지$x$~의$z$~에$\mathbf{A}^{m + n}_K$
가정에 의해 $\text{trdeg}_K(\kappa(x)) \leq a$을 충족한다.
보조정리의 $\mathcal{F}$.
$\mathcal{O}_{Y, y}^{sh}$는 다음의 여과 여극한이므로
에탈 대수는 $K[T_1, \ldots, T_n]$ 이후
$B_i \subset \mathcal{O}_{Y, y}^{sh}$
$\kappa(z)/\kappa(x)$는 대수적이라는 것을 알 수 있습니다
(일부 세부사항 생략). 그런 다음 $\text{trdeg}_K(\kappa(z)) \leq a$
그러므로 $\text{trdeg}_L(\kappa(z)) \leq a - b$.
보조정리 \ref{etale-cohomology-lemma-interlude-II}에 의해 우리는
$$
H^q(\mathbf{A}^m_{B_i}, \mathcal{F}) = 0\text{~을 위한}q > a - b
$$
따라서 정리 \ref{etale-cohomology-theorem-colimit}에 의해 우리는
$(Rf_*\mathcal{F})_{\overline{y}} = 0$ 에 대하여 $q > a - b$
원하는대로.
\end{proof}








\section{유한 코호몰로지 차원} \label{etale-cohomology-section-finite-cd}

\noindent 절 \ref{etale-cohomology-section-cd}에서 시작된 논의를 계속한다.

\begin{definition}
\label{etale-cohomology-definition-cd-f}
$f : X \to Y$를 준컴팩트하고 준분리된 것으로 하자.
스킴의 사상.
{\it 코호몰로지 차원의 $f$}가 가장 작습니다.
요소
$$
\text{cd}(f) \in \{0, 1, 2, \ldots\} \cup \{\infty\}
$$
모든 아벨에 대해 꼬임층 $\mathcal{F}$
$X_\etale$에 $R^if_*\mathcal{F} = 0$이 있다.
$i > \text{cd}(f)$의 경우.
\end{definition}

\begin{lemma}
\label{etale-cohomology-lemma-finite-cd}
$K$를 체로 설정한다.
\begin{enumerate}
\item $f : X \to Y$가 $K$에 대한 유한 유형 스킴의 사상인 경우,
그런 다음 $\text{cd}(f) < \infty$.
\item 만약 $\text{cd}(K) < \infty$이면 $\text{cd}(X) < \infty$
모든 유한 유형에 대해스킴 $X$~ 위에$K$.
\end{enumerate}
\end{lemma}

\begin{proof} 증명 (1). 우리는 $Y$가 유사하다고 가정할 수 있다. 이를 증명하기 위해 스킴, 보조정리 \ref{coherent-lemma-induction-principle}의 코호몰로지의 유도 원리를 사용하겠습니다. $X$도 유사하면 결과는 보조정리 \ref{etale-cohomology-lemma-interlude-I}에 의해 유지된다. 따라서 $X = U \cup V$이고 결과가 $U \to Y$, $V \to Y$ 및 $U \cap V \to Y$에 대해 참이면 $f$에 대해서도 참이라는 것을 보여주는 것으로 충분한다. 이는 상대적인 Mayer-Vietoris 순서를 따릅니다. 보조정리 \ref{etale-cohomology-lemma-relative-mayer-vietoris}를 참조하라.

\medskip\noindent 증명 (2). 이를 증명하기 위해 스킴, 보조정리 \ref{coherent-lemma-induction-principle}의 코호몰로지의 유도 원리를 사용하겠습니다. $X$가 유사하면 결과는 명제 \ref{etale-cohomology-proposition-cd-affine}만큼 유지된다. 따라서 $X = U \cup V$이고 결과가 $U$, $V$ 및 $U \cap V $에 대해 참이면 $X$에 대해서도 참이라는 것을 보여주는 것으로 충분한다. 이는 Mayer-Vietoris 수열에 따른 것이다. 보조정리 \ref{etale-cohomology-lemma-mayer-vietoris}를 참조하라. \end{proof}

\begin{lemma} \label{etale-cohomology-lemma-finite-cd-mod-n-direct-sums} 코호몰로지 및 직합. $n \geq 1$를 정수로 둡니다. \begin{enumerate} \item $f : X \to Y$를 $\text{cd}(f) < \infty$을 사용하여 스킴의 준컴팩트 및 준분리 사상으로 설정한다. 그러면 함자 $$
Rf_* :
D(X_\etale, \mathbf{Z}/n\mathbf{Z})
\longrightarrow
D(Y_\etale, \mathbf{Z}/n\mathbf{Z})
$$는 직합으로 통근한다. \item $X$는 $\text{cd}(X) < \infty$를 사용하여 준간소하고 준분리된 스킴라고 한다. 그러면 함자 $$
R\Gamma(X, -) :
D(X_\etale, \mathbf{Z}/n\mathbf{Z})
\longrightarrow
D(\mathbf{Z}/n\mathbf{Z})
$$는 직합으로 통근한다. \end{enumerate} \end{lemma}

\begin{proof}
증명 (1). $\text{cd}(f) < \infty$ 이후 우리는
$$
f_* :
\textit{Mod}(X_\etale, \mathbf{Z}/n\mathbf{Z})
\longrightarrow
\textit{Mod}(Y_\etale, \mathbf{Z}/n\mathbf{Z})
$$
의미에서 유한 코호몰로지 차원 가 있다.
유도 범주, 보조정리 \ref{derived-lemma-unbounded-right-derived}.
$I$를 집합으로 하고 $i \in I$의 경우 $E_i$를
$D(X_\etale, \mathbf{Z}/n\mathbf{Z})$의 객체이다.
K 단사를 선택하세요. 복합체 $\mathcal{I}_i^\bullet$
$\mathbf{Z}/n\mathbf{Z}$-가군 각각의
용어 $\mathcal{I}_i^n$는 다음의 단사 층이다.
$\mathbf{Z}/n\mathbf{Z}$-가군은 $E_i$을 나타냅니다.
분사 참조, 정리
\ref{injectives-theorem-K-injective-embedding-grothendieck}.
그런 다음 $\bigoplus E_i$는 복합체로 표시된다.
$\bigoplus \mathcal{I}_i^\bullet$ (항별 직합), 참조
분사, 보조정리 \ref{injectives-lemma-derived-products}.
보조정리 \ref{etale-cohomology-lemma-relative-colimit}에 의해 우리는
$$
R^qf_*(\bigoplus \mathcal{I}_i^n) =
\bigoplus R^qf_*(\mathcal{I}_i^n) = 0
$$
$q > 0$ 및 모든 $n$에 대해. 따라서 우리는 다음과 같이 결론을 내립니다.
유도 범주, 보조정리 \ref{derived-lemma-unbounded-right-derived}
복합체로 $Rf_*(\bigoplus E_i)$를 계산할 수 있다.
$$
f_*(\bigoplus \mathcal{I}_i^\bullet) =
\bigoplus f_*(\mathcal{I}_i^\bullet)
$$
(보조정리 \ref{etale-cohomology-lemma-relative-colimit}에 의해 다시 동등)
이미 사용된 $\bigoplus Rf_*E_i$을 나타냅니다.
분사, 보조정리 \ref{injectives-lemma-derived-products}.

\medskip\noindent 증명 (2). 이는 (1)의 증명과 동일하므로 생략한다. \end{proof}

\begin{lemma} \label{etale-cohomology-lemma-proper-mod-n-direct-sums} $f : X \to Y$를 스킴의 고유 사상으로 둡니다. $n \geq 1$를 정수로 둡니다. 그런 다음 함자 $$
Rf_* :
D(X_\etale, \mathbf{Z}/n\mathbf{Z})
\longrightarrow
D(Y_\etale, \mathbf{Z}/n\mathbf{Z})
$$는 직합으로 통근한다. \end{lemma}

\begin{proof} $Y$이 준컴팩트한 경우 이를 증명하는 것으로 충분한다. 사상, 보조정리 \ref{morphisms-lemma-morphism-finite-type-bounded-dimension}에 의해 $f : X \to Y$의 올의 차원이 제한되어 있음을 알 수 있다. 따라서 보조정리 \ref{etale-cohomology-lemma-cohomological-dimension-proper}는 $\text{cd}(f) < \infty$를 의미한다. 따라서 보조정리 \ref{etale-cohomology-lemma-finite-cd-mod-n-direct-sums}의 결과이다. \end{proof}

\begin{lemma} \label{etale-cohomology-lemma-pull-out-constant-mod-n} $X$를 $\text{cd}(X) < \infty$이 되도록 준간소하고 준분리된 스킴으로 둡니다. $\Lambda$를 비틀림 환이라고 가정한다. $E \in D(X_\etale, \Lambda)$ 및 $K \in D(\Lambda)$로 설정한다. 그러면 $$
R\Gamma(X, E \otimes_\Lambda^\mathbf{L} \underline{K}) =
R\Gamma(X, E) \otimes_\Lambda^\mathbf{L} K
$$ \end{lemma}

\begin{proof}
오른쪽에서 왼쪽으로 표준 지도가 있다.
코호몰로지 사이트, 절 \ref{sites-cohomology-section-projection-formula}.
$T(K)$를 성질로 설정한다.
보조정리는 $K \in D(\Lambda)$에 대해 유지된다.
(1), (2), (3)의 조건을 확인하겠습니다.
대수학에 대한 추가 정보, 비고 \ref{more-algebra-remark-P-resolution}
$T$이 끝날 때까지 기다리세요.
성질 (1)는 양쪽이 동등하기 때문에 유지된다.
직합으로 출퇴근한다.
보조정리 \ref{etale-cohomology-lemma-finite-cd-mod-n-direct-sums}.
성질 (2)는 완전 함자를 비교하고 있으므로 유지된다.
삼각측량된 범주 사이에서 다음을 사용할 수 있다.
유도 범주, 보조정리 \ref{derived-lemma-third-isomorphism-triangle}.
성질 (3)는 $K = \Lambda[k]$일 때 보조정리가 유지된다고 말한다.
모든 교대조 $k \in \mathbf{Z}$에 대해 이는 명백한다.
\end{proof}

\begin{lemma}
\label{etale-cohomology-lemma-projection-formula-proper-mod-n}
$f : X \to Y$를 스킴의 고유 사상으로 설정한다. $\Lambda$
비틀림 환이어야 한다. $E \in D(X_\etale, \Lambda)$ 그리고
$K \in D(Y_\etale, \Lambda)$. 그 다음에
$$
Rf_*E \otimes_\Lambda^\mathbf{L} K =
Rf_*(E \otimes_\Lambda^\mathbf{L} f^{-1}K)
$$
$D(Y_\etale, \Lambda)$에서.
\end{lemma}

\begin{proof} 사이트, 절 \ref{sites-cohomology-section-projection-formula}에는 코호몰로지에 의해 왼쪽에서 오른쪽으로 표준 지도가 있다. $\overline{y}$에서 줄기의 동등성을 확인하겠습니다. 고유 기저변환 (보조정리 \ref{etale-cohomology-lemma-proper-base-change-mod-n} 여기서 $Y' = \overline{y}$ 형식)에 의해 이는 $Y$가 대수적으로 닫힌 체의 스펙트럼인 경우로 축소된다. 이는 보조정리 \ref{etale-cohomology-lemma-pull-out-constant-mod-n}에 표시되어 있으며 여기서 $\text{cd}(X) < \infty$를 보조정리 \ref{etale-cohomology-lemma-cohomological-dimension-proper}로 사용한다. \end{proof}






\section{에탈의 퀴네스 코호몰로지} \label{etale-cohomology-section-kunneth}

\noindent 먼저 요소 중 하나가 적절한 경우 Künneth 공식을 증명한다. 그런 다음 이 공식을 사용하여 열린 몰입에 대한 염기 변화 성질을 증명한다. 그러면 ``체'의 스펙트럼을 향한 사상에 의한 기본 변화가 발생합니다(매끄러운 기저변환과 유사). 마지막으로 우리는 이것을 사용하여 보다 일반적인 Künneth 공식을 얻습니다.

\begin{remark}
\label{etale-cohomology-remark-define-kunneth-map}
스킴의 범주에 있는 데카르트 다이어그램을 고려해보세요.
$$
\xymatrix{
X \times_S Y \ar[d]_p \ar[r]_q \ar[rd]_c & Y \ar[d]^g \\
X \ar[r]^f & S
}
$$
$\Lambda$를 환으로 두고 $E \in D(X_\etale, \Lambda)$로 설정한다.
및 $K \in D(Y_\etale, \Lambda)$. 그런 다음 표준 맵이 있다.
$$
Rf_*E \otimes_\Lambda^\mathbf{L} Rg_*K
\longrightarrow
Rc_*(p^{-1}E \otimes_\Lambda^\mathbf{L} q^{-1}K)
$$
예의 경우 표준 맵을 사용하여 이를 정의할 수 있다.
$Rf_*E \to Rc_*p^{-1}E$ 및 $Rg_*K \to Rc_*q^{-1}K$ 및
사이트의 코호몰로지에 정의된 관련 컵 제품,
비고 \ref{sites-cohomology-remark-cup-product}.
아니면 지도에 인접한 것을 사용할 수도 있습니다
$$
c^{-1}(Rf_*E \otimes_\Lambda^\mathbf{L} Rg_*K)
=
p^{-1}f^{-1}Rf_*E \otimes_\Lambda^\mathbf{L} q^{-1} g^{-1}Rg_*K
\to
p^{-1}E \otimes_\Lambda^\mathbf{L} q^{-1}K
$$
이는 부속 맵 $f^{-1}Rf_*E \to E$을 사용하고
$g^{-1}Rg_*K \to K$.
\end{remark}


\begin{lemma} \label{etale-cohomology-lemma-kunneth-one-proper} $k$를 분리 가능하게 닫힌 체라고 가정한다. $X$를 $k$에 대해 적절한 스킴으로 설정한다. $Y$는 $k$에 걸쳐 준간소하고 준분리된 스킴라고 한다. \begin{enumerate} \item 만약 $E \in D^+(X_\etale)$에 비틀림 코호몰로지 층 및 $K \in D^+(Y_\etale)$가 있으면 $$
R\Gamma(X \times_{\Spec(k)} Y,
\text{pr}_1^{-1}E
\otimes_\mathbf{Z}^\mathbf{L}
\text{pr}_2^{-1}K
)
=
R\Gamma(X, E)
\otimes_\mathbf{Z}^\mathbf{L}
R\Gamma(Y, K)
$$ \item 만약 $n \geq 1$ 는 정수이고, $Y$는 $k$, $E \in D(X_\etale, \mathbf{Z}/n\mathbf{Z})$, $K \in D(Y_\etale, \mathbf{Z}/n\mathbf{Z})$에 대해 유한 유형이고, $$
R\Gamma(X \times_{\Spec(k)} Y,
\text{pr}_1^{-1}E
\otimes_{\mathbf{Z}/n\mathbf{Z}}^\mathbf{L}
\text{pr}_2^{-1}K
)
=
R\Gamma(X, E)
\otimes_{\mathbf{Z}/n\mathbf{Z}}^\mathbf{L}
R\Gamma(Y, K)
$$ \end{enumerate} \end{lemma}이다.

\begin{proof}
증명 (1). 보조정리 \ref{etale-cohomology-lemma-projection-formula-proper}에 의해 우리는
$$
R\text{pr}_{2, *}(
\text{pr}_1^{-1}E
\otimes_\mathbf{Z}^\mathbf{L}
\text{pr}_2^{-1}K) =
R\text{pr}_{2, *}(\text{pr}_1^{-1}E)
\otimes_\mathbf{Z}^\mathbf{L}
K
$$
고유 기저변환 기준 (보조정리 \ref{etale-cohomology-lemma-proper-base-change} 형식)
이것은 객체와 동일합니다
$$
\underline{R\Gamma(X, E)}
\otimes_\mathbf{Z}^\mathbf{L}
K
$$
$D(Y_\etale)$. 이 물체에 $R\Gamma(Y, -)$을 사용하면
Leray 스펙트럼 열에 의한 (1)의 평등의 왼쪽
$\text{pr}_2$의 경우. 따라서 우리는 다음과 같이 결론을 내립니다.
보조정리 \ref{etale-cohomology-lemma-pull-out-constant}.

\medskip\noindent 증명 (2). 이는 1)의 증명과 정확히 동일한다. 단, 기본정리 \ref{etale-cohomology-lemma-projection-formula-proper-mod-n}, \ref{etale-cohomology-lemma-proper-base-change-mod-n} 및 \ref{etale-cohomology-lemma-pull-out-constant-mod-n}와 $\text{cd}(Y) < \infty$를 보조정리로 사용한다. \ref{etale-cohomology-lemma-finite-cd}. \end{proof}

\begin{lemma}
\label{etale-cohomology-lemma-supported-in-closed-points}
$K$는 분리 가능하게 닫힌 체이라고 가정한다. $X$를 유한의 스킴으로 설정한다.
$K$ 위에 입력하세요. $\mathcal{F}$를 $X_\etale$에서 아벨 층으로 설정한다.
지지집합은 $X$의 폐점 집합에 포함되어 있다.
그 다음에$H^q(X, \mathcal{F}) = 0$~을 위한$q > 0$그리고$\mathcal{F}$
전역적으로 생성된다.
\end{lemma}

\begin{proof}
($\mathcal{F}$가 비틀림이면 소멸이 즉시 뒤따릅니다.
보조정리 \ref{etale-cohomology-lemma-interlude-II}에서.)
보조정리 \ref{etale-cohomology-lemma-support-in-subset}로 $\mathcal{F}$라고 쓸 수 있다.
구성가능층 $\mathcal{F}_i$의 여과 여극한으로
~의$\mathbf{Z}$-가군누구의 지원$Z_i \subset X$유한하다
닫힌 점의 집합. 에 의해
명제 \ref{etale-cohomology-proposition-closed-immersion-pushforward}
이러한 층은 $(Z_i \to X)_*\mathcal{G}_i$ 형식이다.
여기서 $\mathcal{G}_i$는 $Z_i$의 층이다. $K$는 분리 가능하게 닫혀 있으므로,
스킴 $Z_i$는 분리 가능하게 닫힌 스펙트럼의 유한 분리 결합이다.
체. $H^q(Z_i, \mathcal{G}_i) = H^q(X, \mathcal{F}_i)$를 기억하세요.
$Z_i \to X$에 대한 Leray 스펙트럼 열에 의해 더 높은 것의 소멸
직상 이 사상에 대해
(명제 \ref{etale-cohomology-proposition-finite-higher-direct-image-zero}).
기본형 \ref{etale-cohomology-lemma-equivalence-abelian-sheaves-point} 및
\ref{etale-cohomology-lemma-compare-cohomology-point}
$H^q(Z_i, \mathcal{G}_i)$는 $q > 0$에 대해 0이다.
그리고 $H^0(Z_i, \mathcal{G}_i)$는 $\mathcal{G}_i$을 생성한다.
우리는 결론을 내린다소멸~의$H^q(X, \mathcal{F}_i)$~을 위한$q > 0$
$\mathcal{F}_i$은 대역 절단에 의해 생성된다.
정리 \ref{etale-cohomology-theorem-colimit}에 의해 우리는
$H^q(X, \mathcal{F}) = 0$~을 위한$q > 0$.
이제 증명이 완료되었다.
여과 여극한 전역적으로 생성된 층의 아벨 군
전역적으로 생성됩니다(자세한 내용은 생략됨).
\end{proof}

\begin{lemma}
\label{etale-cohomology-lemma-vanishing-closed-points}
$K$는 분리 가능하게 닫힌 체이라고 가정한다. $X$를 유한의 스킴으로 설정한다.
$K$ 위에 입력하세요. $Q \in D(X_\etale)$로 놔두세요. $Q_{\overline{x}}$
$x$이 $X$의 폐점인 경우에만 0이 아닙니다. 그 다음에
$$
Q = 0 \Leftrightarrow H^i(X, Q) = 0 \text{(임의의)}i
$$
\end{lemma}

\begin{proof} 왼쪽에서 오른쪽으로의 의미는 사소한다. 따라서 우리는 역함수를 증명해야 한다.

\medskip\noindent
$Q$이 아래로 제한된다고 가정한다. 이 경우는 거의 모든 경우에 충분한다.
응용 프로그램. $Q$가 0이 아니면 가장 작은 $i$를 볼 수 있다.
코호몰로지 층 $H^i(Q)$는 0이 아닙니다.
보조정리 \ref{etale-cohomology-lemma-supported-in-closed-points}에 의해 우리는
$H^i(X, Q) = H^0(X, H^i(Q)) \not = 0$ 그리고 결론을 내립니다.

\medskip\noindent
일반적인 경우. $\mathcal{B} \subset \Ob(X_\etale)$를
준컴팩트한 물체. 보조정리 \ref{etale-cohomology-lemma-supported-in-closed-points}에 의해
사이트, 보조정리에 대한 코호몰로지의 가정
\ref{sites-cohomology-lemma-cohomology-over-U-trivial}
만족한다. 우리는 $H^q(U, Q) = H^0(U, H^q(Q))$라고 결론을 내렸습니다.
모든 $U \in \mathcal{B}$에 대해. 특히 이는 $U = X$에 적용된다.
따라서 보조정리 \ref{etale-cohomology-lemma-supported-in-closed-points}에 의한 결론은 다음과 같습니다.
$Q$는 $D(X_\etale)$에서 0이므로 필요충분조건은 $H^q(Q)$는 0이다.
모든 $q$에 대해.
\end{proof}

\begin{lemma} \label{etale-cohomology-lemma-kunneth-localize-on-X} $K$를 체로 둡니다. $j : U \to X$를 $K$ 위에 유한형의 스킴의 열린 몰입라고 하자. $Y$를 $K$에 대해 유한 유형의 스킴라고 가정한다. $$
\xymatrix{
Y \times_{\Spec(K)} X \ar[d]_q  &
Y \times_{\Spec(K)} U \ar[l]^h \ar[d]^p \\
X & U \ar[l]_j
}
$$ 다이어그램을 고려하면 기본 변경 맵 $q^{-1}Rj_*\mathcal{F} \to Rh_*p^{-1}\mathcal{F}$은 $\mathcal{F}$에 대한 동형사상이고 $U_\etale$의 아벨 층은 줄기에서 반전 가능한 순서의 비틀림이다. $K$. \end{lemma}

\begin{proof}
여극한 어디에 $\mathcal{F} = \colim \mathcal{F}[n]$를 쓰세요.
는 $K$에서 가역이 가능한 정수의 곱셈 체계에 관한 것이다.
우리 상황에서는 코호몰로지가 여과 여극한으로 통근하므로
(정확한 참조는 보조정리 \ref{etale-cohomology-lemma-base-change-Rf-star-colim} 참조),
$\mathcal{F}[n]$에 대해 보조정리를 증명하는 것으로 충분한다. 따라서 우리는
$\mathcal{F}$는 $\mathbf{Z}/n\mathbf{Z}$-가군의 층라고 가정한다.
어떤 사람들에게는$n$에서 반전 가능$K$(우리는 이것을 마지막에 사용할 것이다.
증명). 증명에서는 올에 약칭 $X \times_K Y$를 사용한다.
$\Spec(K)$ 이상의 제품이다.
$\dim(X) + \dim(Y)$에 대한 귀납법으로 보조정리를 증명하겠습니다.
보조정리는 $\dim(X) \leq 0$인 경우 사소한다. 이 경우
$U$은 $X$의 개방형 및 폐쇄형 하위 구성표이다.
점 $z \in X \times_K Y$을 선택한다. 우리는 보여줄 것이다
$\overline{z}$의 줄기는 동형사상이다.

\medskip\noindent
$z \mapsto x \in X$라고 가정하고 $\text{trdeg}_K(\kappa(x)) > 0$라고 가정한다.
두자 $X' = \Spec(\mathcal{O}_{X, x}^{sh})$
$U' \subset X'$는 $U$의 역상을 나타냅니다.
기본 변경을 고려하세요.
$$
\xymatrix{
Y \times_K X' \ar[d]_{q'}  &
Y \times_K U' \ar[l]^{h'} \ar[d]^{p'} \\
X' & U' \ar[l]_{j'}
}
$$
$X' \to X$로 다이어그램을 완성했다. $X' \to X$이 필터링된 것을 확인하세요.
에탈 사상의 여극한. 매끄러운 기저변환 형식으로
보조정리 \ref{etale-cohomology-lemma-smooth-base-change-general}
$q^{-1}Rj_*\mathcal{F} \to Rh_*p^{-1}\mathcal{F}$ ~ $X'$ ~
$Y \times_K X'$는 지도입니다
$(q')^{-1}Rj'_*\mathcal{F}' \to Rj'_*(p')^{-1}\mathcal{F}'$ 여기서
$\mathcal{F}'$은 $\mathcal{F}$에서 $U'$로의 철회이다.
(이 단계에서는 에탈 기본 변경을 사용하는 것으로 충분한다.
본질적으로 사소한 결과이다.) 따라서 보여주는 것으로 충분한다.
그 $(q')^{-1}Rj'_*\mathcal{F}' \to Rj'_*(p')^{-1}\mathcal{F}'$
은동형사상우리의 원본을 증명하기 위해
지도는 $\overline{z}$의 줄기에 있는 동형사상이다.
보조정리 \ref{etale-cohomology-lemma-strictly-henselian}에 의해
분리 가능하게 닫힌 체 $L/K$가 있다.
$X' = \lim X_i$ $X_i$ $L$ 위에 유한 유형의 아핀 포함
및 $\dim(X_i) < \dim(X)$. $i$의 경우 충분히 큽니다.
열려있는 존재$U_i \subset X_i$다음으로 제한$U'$~에$X'$.
귀납 가설을 다이어그램에 적용할 수 있다.
$$
\vcenter{
\xymatrix{
Y \times_K X_i \ar[d]_{q_i}  &
Y \times_K U_i \ar[l]^{h_i} \ar[d]^{p_i} \\
X_i & U_i \ar[l]_{j_i}
}
}
\quad\text{즉}\quad
\vcenter{
\xymatrix{
Y_L \times_L X_i \ar[d]_{q_i}  &
Y_L \times_L U_i \ar[l]^{h_i} \ar[d]^{p_i} \\
X_i & U_i \ar[l]_{j_i}
}
}
$$
체 $L$ 및 $\mathcal{F}$의 이 다이어그램에 대한 철수.
보조정리 \ref{etale-cohomology-lemma-base-change-Rf-star-colim}에 의해 우리는 지도가
$(q')^{-1}Rj'_*\mathcal{F}' \to Rj'_*(p')^{-1}\mathcal{F}$은 동형사상이다.

\medskip\noindent
$z \mapsto y \in Y$라고 가정하고 $\text{trdeg}_K(\kappa(y)) > 0$라고 가정한다.
$Y' = \Spec(\mathcal{O}_{X, x}^{sh})$로 놔두세요.
보조정리 \ref{etale-cohomology-lemma-strictly-henselian}에 의해 분리 가능하게 닫힌 체가 있다.
$L/K$ $Y' = \lim Y_i$와 $Y_i$가 $L$ 위에 유한 유형의 유사성을 갖도록
및 $\dim(Y_i) < \dim(Y)$. 특히 $Y'$는 $L$에 대한 스킴이다.
아래첨자로 표시$L$기본 변경 사항스킴~ 위에$K$
$L$를 통해 스킴으로. 가환 도식을 고려해보세요.
$$
\vcenter{
\xymatrix{
Y' \times_K X \ar[d]_f  &
Y' \times_K U \ar[l]^{h'} \ar[d]^{f'} \\
Y \times_K X \ar[d]_q  &
Y \times_K U \ar[l]^h \ar[d]^p \\
X & U \ar[l]_j
}
}
\quad\text{및}\quad
\vcenter{
\xymatrix{
Y' \times_L X_L \ar[d]_{q'}  &
Y' \times_L U_L \ar[l]^{h'} \ar[d]^{p'} \\
X_L \ar[d] &
U_L \ar[l]^{j_L} \ar[d] \\
X & U \ar[l]_j
}
}
$$
위쪽과 아래쪽 행이 왼쪽과 오른쪽에서 동일하다는 것을 관찰하세요.
매끄러운 기저변환에 의해 우리는 그것을 알 수 있습니다
$f^{-1}Rh_*p^{-1}\mathcal{F} = Rh'_*(f')^{-1}p^{-1}\mathcal{F}$
(이전 단락과 유사).
$\Spec(L) \to \Spec(K)$에 대해 매끄러운 기저변환에 의해
(보조정리 \ref{etale-cohomology-lemma-base-change-field-extension})
우리는 $Rj_{L, *}\mathcal{F}_L$가
$Rj_*\mathcal{F}$ ~ $X_L$. 이 두 가지 관찰을 결합하면,
우리는 기본 변화 맵을 증명하는 것으로 충분하다고 결론을 내립니다.
오른쪽 다이어그램의 위쪽 사각형에 대해
은동형사상우리의 원본을 증명하기 위해
지도는 $\overline{z}$\footnote{에 있는 줄기의 동형사상이다.
기본 변경 맵의 '수직 구성'이 기본 변경이라는 점을 사용하세요.
사이트, 비고의 코호몰로지에 설명된 지도
\ref{sites-cohomology-remark-compose-base-change}.}.
그런 다음 $Y' = \lim Y_i$를 사용하여 정확하게 논쟁한다.
이전 단락에서와 같이 유도는 다음과 같습니다.
가설은 $Y' \times_K X$에 대한 지도를 강제로
동형사상.

\medskip\noindent
따라서 $\dim(X) + \dim(Y)$ 최소값을 갖는 모든 카운터 예
비동형성만 가질 것이다
$q^{-1}Rj_*\mathcal{F} \to Rh_*p^{-1}\mathcal{F}$
$X \times_K Y$의 닫힌 점에서 줄기에 (왜냐하면 점
$z$~의$X \times_K Y$폐쇄 지점이다필요충분조건은
둘 다의 이미지$z$~에$X$그리고$Y$닫혀 있습니다).
동형사상을 로컬에서 증명하면 충분하므로,
$X$와 $Y$가 유사하다고 가정할 수 있다. 그러나 그렇다면
우리는 개방형을 선택할 수 있어요조밀담금$Y \to Y'$~와 함께$Y'$
투영. (닫힌 몰입 $Y \to \mathbf{A}^n_K$를 선택하세요.
그리고 $Y'$는 $Y$의 스킴 이론적 종결이라고 가정한다.
$\mathbf{P}^n_K$.) 그런 다음 $\dim(Y') = \dim(Y)$ 그리고
우리는 '최소' 카운터를 얻습니다예~와 함께$Y$투영 이상$K$.
다음 단락에서 우리는 이것이 일어날 수 없음을 보여줍니다.

\medskip\noindent
보조정리의 설명과 같은 다이어그램을 생각해 보세요.
$q^{-1}Rj_*\mathcal{F} \to Rh_*p^{-1}\mathcal{F}$
$X \times_K Y$의 닫히지 않은 모든 지점에서 동형사상이다.
그리고 $Y$는 투영형이다.
지도 제한은 $(X \times_K Y)_{K^{sep}}$
는 보조정리 다이어그램에 해당하는 맵이다.
베이스가 $K^{sep}$로 변경되었다. 따라서 우리는 $K$를 가정할 수 있고 실제로 가정한다.
분리가능하게 대수적으로 닫혀있다. 구별되는 삼각형을 선택하세요
$$
q^{-1}Rj_*\mathcal{F} \to Rh_*p^{-1}\mathcal{F} \to Q \to
(q^{-1}Rj_*\mathcal{F})[1]
$$
$D((X \times_K Y)_\etale)$에서. $Q$는 다음에서 지원된다.
닫힌 점은 증명하기에 충분하다는 것을 알 수 있다.
$H^i(X \times_K Y, Q) = 0$ 모든 $i$에 대해 참조
보조정리 \ref{etale-cohomology-lemma-vanishing-closed-points}.
따라서 다음을 증명하는 것으로 충분한다.
$q^{-1}Rj_*\mathcal{F} \to Rh_*p^{-1}\mathcal{F}$
코호몰로지에 동형사상을 유도한다. $\mathcal{F}$를 기억하세요.
에 의해 소멸된다$n$에서 반전 가능$K$.
보조정리 \ref{etale-cohomology-lemma-kunneth-one-proper}의 퀴네스 공식에 따라
우리는
\begin{align*}
R\Gamma(X \times_K Y, q^{-1}Rj_*\mathcal{F})
&=
R\Gamma(X, Rj_*\mathcal{F})
\otimes_{\mathbf{Z}/n\mathbf{Z}}^\mathbf{L}
R\Gamma(Y, \mathbf{Z}/n\mathbf{Z}) \\
& =
R\Gamma(U, \mathcal{F})
\otimes_{\mathbf{Z}/n\mathbf{Z}}^\mathbf{L}
R\Gamma(Y, \mathbf{Z}/n\mathbf{Z})
\end{align*}
그리고
$$
R\Gamma(X \times_K Y, Rh_*p^{-1}\mathcal{F}) =
R\Gamma(U \times_K Y, p^{-1}\mathcal{F}) =
R\Gamma(U, \mathcal{F})
\otimes_{\mathbf{Z}/n\mathbf{Z}}^\mathbf{L}
R\Gamma(Y, \mathbf{Z}/n\mathbf{Z})
$$
이로써 증명이 완료된다.
\end{proof}

\begin{lemma}
\label{etale-cohomology-lemma-punctual-base-change}
$K$를 체로 설정한다. 가환 도식
$$
\xymatrix{
X \ar[d] & X' \ar[l] \ar[d]_{f'} & Y \ar[l]^h \ar[d]^e \\
\Spec(K) & S' \ar[l] & T \ar[l]_g
}
$$
$K$에 대한 스킴의
$X' = X \times_{\Spec(K)} S'$ 및 $Y = X' \times_{S'} T$ 및
$g$ 준컴팩트하고 준분리형이며, 모든 아벨 층
$\mathcal{F}$~에$T_\etale$누구의줄기주문이 비틀어지고 있다
$K$ 기본 변경 맵에서 반전 가능
$$
(f')^{-1}Rg_*\mathcal{F}
\longrightarrow
Rh_*e^{-1}\mathcal{F}
$$
동형사상이다.
\end{lemma}

\begin{proof}
질문은 $X$에 대한 로컬이므로 $X$이 유사하다고 가정할 수 있다.
극한, 보조정리 \ref{limits-lemma-relative-approximation}에 의해
$X = \lim X_i$를 아핀과 함께 공동 필터링된 극한으로 쓸 수 있다.
이행사상~의스킴 $X_i$유한 유형의$K$.
$X'_i =  X_i \times_{\Spec(K)} S'$ 및 $Y_i = X'_i \times_{S'} T$을 나타냅니다.
보조정리 \ref{etale-cohomology-lemma-base-change-Rf-star-colim}에 의해
제곱에 대한 진술을 증명하는 것으로 충분한다.
모서리 $X_i, Y_i, S_i, T_i$가 있다.
따라서 $X$은 $K$에 대해 유한 유형이라고 가정할 수 있다.
마찬가지로, 우리는 다음과 같이 쓸 수 있다.
$\mathcal{F} = \colim \mathcal{F}[n]$ 여기서 여극한
는 $K$에서 가역이 가능한 정수의 곱셈 체계에 관한 것이다.
위에서 사용된 동일한 보조정리는 다음과 같은 경우로 우리를 줄여줍니다.
$\mathcal{F}$는 $\mathbf{Z}/n\mathbf{Z}$-가군의 층이다.
일부 $n$의 경우 $K$에서 되돌릴 수 있다.

\medskip\noindent
$K$을 대수적 종결자 $\overline{K}$로 대체할 수 있다.
즉, 직상의 형성은 염기변화에 따른 통근이다.
$\overline{K}$에 따라
보조정리 \ref{etale-cohomology-lemma-base-change-field-extension} (둘 다 작동
$g$ 및 $h$). 그리고 그 후에 합의를 증명하는 것으로 충분한다.
$X'_{\overline{K}}$로 제한된다. 다음으로 교체할 수 있다.
$X$ 에탈의 위상적 불변성을 가지므로 이를 감소시킴으로써
코호몰로지, 참조
명제 \ref{etale-cohomology-proposition-topological-invariance}.
이 교체 후 사상 $X \to \Spec(K)$
기하학적으로 축소된 올 평면적이고 유한한 표현이다.
모든 염기 변경, 특히 $X' \to S'$의 경우에도 마찬가지이다.
그러므로 $(f')^{-1}g_*\mathcal{F} \to Rh_*e^{-1}\mathcal{F}$
보조정리 \ref{etale-cohomology-lemma-fppf-reduced-fibres-base-change-f-star}에 의한 동형사상이다.

\medskip\noindent 이 시점에서 보조정리 \ref{etale-cohomology-lemma-base-change-does-not-hold-post}를 적용하여 증명하기에 충분하다는 것을 확인할 수 있다. 가환 도식 $$
\xymatrix{
X \ar[d]_f & X' \ar[d] \ar[l] & Y \ar[l]^h \ar[d] \\
\Spec(K) & S' \ar[l] & \Spec(L) \ar[l]
}
$$가 두 사각형 데카르트식을 갖는 경우, 여기서 $S'$는 다음과 같습니다. 대수적으로 닫힌 함수 체 $K$를 사용하여 아핀, 정역 및 정규를 사용하는 경우 $R^qh_*(\mathbf{Z}/d\mathbf{Z})$는 $q > 0$ 및 $d | n$에 대해 0이다. 이 소멸은 $$
(f')^{-1}R^q(\Spec(L) \to S')_*\mathbf{Z}/d\mathbf{Z}
\longrightarrow
R^qh_*\mathbf{Z}/d\mathbf{Z}
$$가 동형사상라는 진술과 동일한다. 왜냐하면 보조정리 \ref{etale-cohomology-lemma-Rf-star-zero-normal-with-alg-closed-function-field}에 의해 예에 대해 왼쪽이 0이기 때문이다.

\medskip\noindent
쓰다$S' = \Spec(B)$~하도록 하다$L$분수입니다체~의$B$.
$B = \bigcup_{i \in I} B_i$을 유한 유형의 합집합으로 작성하세요.
$K$-하위대수 $B_i$. $J$를 $(i, g)$ 쌍의 집합라고 가정한다.
$i \in I$ 및 $g \in B_i$ 0이 아닌 순서 지정
$(i', g') \geq (i, g)$ 필요충분조건은 $i' \geq i$ 및
$g$는 $(B_{i'})_{g'}$의 반전 가능한 요소에 매핑된다.
그런 다음 $L = \colim_{(i, g) \in J} (B_i)_g$.
$j = (i, g) \in J$ 집합 $S_j = \Spec(B_i)$의 경우
및 $U_j = \Spec((B_i)_g)$.
그 다음에
$$
\vcenter{
\xymatrix{
X' \ar[d] & Y \ar[l]^h \ar[d] \\
S' & \Spec(L) \ar[l]
}
}
\quad\text{는 다음 도식의 여극한이다}\quad
\vcenter{
\xymatrix{
X \times_K S_j \ar[d] & X \times_K U_j \ar[l]^{h_j} \ar[d] \\
S_j & U_j \ar[l]
}
}
$$
따라서 보조정리 \ref{etale-cohomology-lemma-base-change-Rf-star-colim}를 적용할 수 있다.
증명하는 것으로 충분하다는 것을 확인하기 위해
기본 변경 사항은 오른쪽 다이어그램에 있다.
보조정리 \ref{etale-cohomology-lemma-kunneth-localize-on-X}에서 입증되었다.
\end{proof}

\begin{lemma}
\label{etale-cohomology-lemma-punctual-base-change-upgrade}
$K$를 체로 설정한다. $n \geq 1$를 $K$에서 반전 가능하게 만듭니다.
가환 도식을 고려해보세요.
$$
\xymatrix{
X \ar[d] & X' \ar[l]^p \ar[d]_{f'} & Y \ar[l]^h \ar[d]^e \\
\Spec(K) & S' \ar[l] & T \ar[l]_g
}
$$
스킴의
$X' = X \times_{\Spec(K)} S'$ 및 $Y = X' \times_{S'} T$ 및
$g$ 준컴팩트하고 준분리형이다. 표준 지도
$$
p^{-1}E \otimes_{\mathbf{Z}/n\mathbf{Z}}^\mathbf{L} (f')^{-1}Rg_*F
\longrightarrow
Rh_*(h^{-1}p^{-1}E \otimes_{\mathbf{Z}/n\mathbf{Z}}^\mathbf{L} e^{-1}F)
$$
은동형사상만약에$E$~에$D^+(X_\etale, \mathbf{Z}/n\mathbf{Z})$
토르 진폭이 있음$[a, \infty]$어떤 사람들에게는$a \in \mathbf{Z}$그리고
$F$~에$D^+(T_\etale, \mathbf{Z}/n\mathbf{Z})$.
\end{lemma}

\begin{proof} 이 보조정리는 보조정리 \ref{etale-cohomology-lemma-punctual-base-change}를 유도 범주의 객체에 대한 일반화이다. 보조정리의 주장은 보조정리 \ref{etale-cohomology-lemma-punctual-base-change}에서 스킴 $X$가 $K$에 대해 임의적이기 때문에 참이다. 우리는 독자들이 힘든 증명을 건너뛰기를 강력히 촉구합니다(대안: 마지막 단락만 읽으십시오).

\medskip\noindent
우리는 K-플랫 아래의 경계 복합체로 $E$를 표현할 수 있다.
$\mathcal{E}^\bullet$ 플랫 $\mathbf{Z}/n\mathbf{Z}$-가군으로 구성된다.
사이트, 보조정리의 코호몰로지를 참조하라.
\ref{sites-cohomology-lemma-bounded-below-tor-amplitude}.
정수를 선택하세요$b$그렇게$H^i(F) = 0$~을 위한$i < b$.
큰 정수 $N$를 선택하고 짧은 완전열을 고려한다.
$$
0 \to \sigma_{\geq N + 1}\mathcal{E}^\bullet \to
\mathcal{E}^\bullet \to
\sigma_{\leq N}\mathcal{E}^\bullet \to 0
$$
어리석은 잘림. 이렇게 하면 구별되는 삼각형이 생성된다.
$E'' \to E \to E' \to E''[1]$~에$D(X_\etale, \mathbf{Z}/n\mathbf{Z})$.
고정된 $F$ 화살표 양쪽
보조정리의 문에서 $E$의 완전 함자이다. 그것을 관찰하십시오
$$
p^{-1}E'' \otimes_{\mathbf{Z}/n\mathbf{Z}}^\mathbf{L} (f')^{-1}Rg_*F
\quad\text{및}\quad
Rh_*(h^{-1}p^{-1}E'' \otimes_{\mathbf{Z}/n\mathbf{Z}}^\mathbf{L} e^{-1}F)
$$
각도 $\geq N + b$에 앉아 있다. 따라서 보조정리를 증명할 수 있다면
객체 $E'$에 대해 보조정리가 각도 단위로 유지되는 것을 볼 수 있다.
$\leq N + b$ 결론을 내리겠습니다. 일부 세부 사항은 생략되었다.
따라서 우리는 $E$가 표현된다고 가정할 수 있다.
평평한 $\mathbf{Z}/n\mathbf{Z}$-가군의 경계 복합체에 의해.
같은 성격의 또 다른 논증을 해보면, 우리는 다음과 같이 가정할 수 있다.
$E$는 단일 플랫 $\mathbf{Z}/n\mathbf{Z}$-가군으로 제공된다.
$\mathcal{E}$.

\medskip\noindent
다음으로 $F$ 변수에 동일한 인수를 사용한다.
$F$이 단일로 제공되는 경우로 축소
층 중 $\mathbf{Z}/n\mathbf{Z}$-가군 $\mathcal{F}$.
$\mathcal{F}$이 정수 $m | n$에 의해 소멸된다고 가정해 보자.
$\ell$가 $m$과 $m > \ell$를 나누는 소수인 경우,
그러면 짧은 완전열을 볼 수 있다.
$0 \to \mathcal{F}[\ell] \to \mathcal{F} \to
\mathcal{F}/\mathcal{F}[\ell] \to 0$
더 작은 $m$로 줄이다. 이것은 마침내 우리를 다음과 같이 줄이다.
$\mathcal{F}$가 소수에 의해 소멸되는 경우
숫자 $\ell$를 $n$로 나누는 것이다.
이 경우 다음을 관찰하라.
$$
p^{-1}\mathcal{E}
\otimes_{\mathbf{Z}/n\mathbf{Z}}^\mathbf{L}
(f')^{-1}Rg_*\mathcal{F}
=
p^{-1}(\mathcal{E}/\ell \mathcal{E})
\otimes_{\mathbf{F}_\ell}^\mathbf{L}
(f')^{-1}Rg_*\mathcal{F}
$$
$\mathcal{E}$의 평탄도에 따라. 다른 용어도 마찬가지이다.
이는 우리가 층으로 작업하는 경우로 축소된다.
논의되는 $\mathbf{F}_\ell$-벡터공간

\medskip\noindent
$\ell$는 $K$에서 반전 가능한 소수라고 가정한다.
$\mathcal{E}$, $\mathcal{F}$는 층라고 가정한다.
$\mathbf{F}_\ell$-벡터공간~에$X_\etale$그리고$T_\etale$.
우리는 그것을 보여주고 싶다
$$
p^{-1}\mathcal{E} \otimes_{\mathbf{F}_\ell} (f')^{-1}R^qg_*\mathcal{F}
\longrightarrow
R^qh_*(h^{-1}p^{-1}\mathcal{E} \otimes_{\mathbf{F}_\ell} e^{-1}\mathcal{F})
$$
모든 $q \geq 0$에 대해 동형사상이다. 이 질문은 $X$의 로컬 질문이다.
그러므로 우리는 $X$가 유사하다고 가정할 수 있다. $\mathcal{E}$라고 쓸 수 있다.
여과 여극한의 구성가능층으로
$\mathbf{F}_\ell$-벡터공간 $X_\etale$ 참조
보조정리 \ref{etale-cohomology-lemma-torsion-colimit-constructible}.
텐서곱 통근 이후
여과 여극한을 사용하고 이후
고차 직상 너무 합니다 (보조정리 \ref{etale-cohomology-lemma-relative-colimit})
$\mathcal{E}$는 다음의 구성가능층라고 가정할 수 있다.
$\mathbf{F}_\ell$-벡터공간 $X_\etale$에 있다.
그런 다음 정수 $m$와 유한 및 유한하게 표시되는 사상을 선택할 수 있다.
$\pi_i : X_i \to X$, $i = 1, \ldots, m$
주입식 지도가 있도록
$$
\mathcal{E} \to
\bigoplus\nolimits_{i = 1, \ldots, m}
\pi_{i, *}\mathbf{F}_\ell
$$
보조정리 \ref{etale-cohomology-lemma-constructible-maps-into-constant-general}를 참조하라.
직합도 구성가능층임을 확인하세요.
(보조정리 \ref{etale-cohomology-lemma-finite-pushforward-constructible}).
따라서 여핵도 구성 가능한다.
(보조정리 \ref{etale-cohomology-lemma-constructible-abelian}).
차원 변속, 즉 $q$에 대한 유도에 의해,
구성가능층의 범주에
$\mathbf{F}_\ell$-벡터공간 $X_\etale$에 대해 다음을 증명하는 것으로 충분한다.
층 $\pi_{i, *}\mathbf{F}_\ell$에 대한 결과
(세부 사항 생략; 힌트: $q = 0$에 대한 주입성 증명부터 시작하세요.
모든 건설 가능한 $\mathcal{E}$).
이 사례를 증명하기 위해 보조정리의 다이어그램을 다음과 같이 확장한다.
$$
\xymatrix{
X_i \ar[d]^{\pi_i} &
X'_i \ar[l]^{p_i} \ar[d]^{\pi'_i} &
Y_i \ar[l]^{h_i} \ar[d]^{\rho_i} \\
X \ar[d] & X' \ar[l]^p \ar[d]_{f'} & Y \ar[l]^h \ar[d]^e \\
\Spec(K) & S' \ar[l] & T \ar[l]_g
}
$$
모든 사각형 데카르트를 사용한다. 아래 방정식에서 우리는
$R\pi_{i, *} = \pi_{i, *}$를 사용하고 유사하게
$\pi'_i$, $\rho_i$에 대해서는 Leray 스펙트럼 열을 사용하겠습니다.
우리는 사용할 것입니다
보조정리 \ref{etale-cohomology-lemma-finite-pushforward-commutes-with-base-change} 및
우리는 사용할 것입니다
보조정리 \ref{etale-cohomology-lemma-projection-formula-proper-mod-n}
(이 보조정리는 유한 사상의 경우 거의 사소하지만)
$\pi_i$, $\pi'_i$, $\rho_i$.
그렇게 함으로써 우리는 본다.
\begin{align*}
p^{-1}\pi_{i, *}\mathbf{F}_\ell
\otimes_{\mathbf{F}_\ell} (f')^{-1}R^qg_*\mathcal{F}
& =
\pi'_{i, *}\mathbf{F}_\ell
\otimes_{\mathbf{F}_\ell} (f')^{-1}R^qg_*\mathcal{F} \\
& =
\pi'_{i, *}((\pi'_i)^{-1} (f')^{-1}R^qg_*\mathcal{F})
\end{align*}
마찬가지로, 우리는
\begin{align*}
R^qh_*(h^{-1}p^{-1} \pi_{i, *}\mathbf{F}_\ell
\otimes_{\mathbf{F}_\ell} e^{-1}\mathcal{F})
& =
R^qh_*(\rho_{i, *}\mathbf{F}_\ell
\otimes_{\mathbf{F}_\ell} e^{-1}\mathcal{F}) \\
& =
R^qh_*(\rho_i^{-1}e^{-1}\mathcal{F}) \\
& =
\pi'_{i, *}R^qh_{i, *} \rho_i^{-1}e^{-1}\mathcal{F})
\end{align*}
심체
$R^qh_{i, *} \rho_i^{-1}e^{-1}\mathcal{F} = 
(\pi'_i)^{-1} (f')^{-1}R^qg_*\mathcal{F}$ 에 의해
보조정리 \ref{etale-cohomology-lemma-punctual-base-change}
우리는 결론을 내린다.
\end{proof}

\begin{lemma}
\label{etale-cohomology-lemma-punctual-base-change-upgrade-unbounded}
$K$를 체로 설정한다. $n \geq 1$를 $K$에서 반전 가능하게 만듭니다.
가환 도식을 고려해보세요.
$$
\xymatrix{
X \ar[d] & X' \ar[l]^p \ar[d]_{f'} & Y \ar[l]^h \ar[d]^e \\
\Spec(K) & S' \ar[l] & T \ar[l]_g
}
$$
$K$에 대한 유한 유형의 스킴
$X' = X \times_{\Spec(K)} S'$ 및 $Y = X' \times_{S'} T$.
표준 지도
$$
p^{-1}E \otimes_{\mathbf{Z}/n\mathbf{Z}}^\mathbf{L} (f')^{-1}Rg_*F
\longrightarrow
Rh_*(h^{-1}p^{-1}E \otimes_{\mathbf{Z}/n\mathbf{Z}}^\mathbf{L} e^{-1}F)
$$
은동형사상~을 위한$E$~에$D(X_\etale, \mathbf{Z}/n\mathbf{Z})$
그리고$F$~에$D(T_\etale, \mathbf{Z}/n\mathbf{Z})$.
\end{lemma}

\begin{proof} 함자를 직합으로 통근하는 것을 사용하여 이를 보조정리 \ref{etale-cohomology-lemma-punctual-base-change-upgrade}로 줄이다. 독자는 증명을 건너뛰는 것이 좋다. 유도 텐서곱은 직합으로 통근한다는 점을 기억하세요. (파생된) 당김은 직합을 사용하여 통근한다는 점을 기억하세요. $Rh_*$ 및 $Rg_*$는 직합으로 통근한다는 점을 기억하세요. 정리 \ref{etale-cohomology-lemma-finite-cd} 및 \ref{etale-cohomology-lemma-finite-cd-mod-n-direct-sums}를 참조하라(여기서 스킴은 $K$에 대해 유한 유형입니다).

\medskip\noindent
증명을 마무리하기 위해 다음과 같이 주장할 수 있다.
먼저 $E = \text{hocolim} \tau_{\leq N} E$를 씁니다.
우리 함자는 직합으로 통근하므로 그들은 통근한다.
호모토피 여극한을 사용한다. 그러므로 증명하면 충분하다
위에 제한된 $E$에 대한 보조정리이다. 마찬가지로 $F$에 대해서도 우리는
$F$이 위에 제한되어 있다고 가정할 수 있다.
그런 다음 $E$를 복합체 위의 경계로 표현할 수 있다.
$\mathcal{E}^\bullet$의 층의 $\mathbf{Z}/n\mathbf{Z}$-가군.
그 다음에
$$
\mathcal{E}^\bullet = \colim \sigma_{\geq -N}\mathcal{E}^\bullet
$$
(멍청한 잘림).
따라서 우리는 $\mathcal{E}^\bullet$이 제한된 것이라고 가정할 수 있다.
복합체의 층의 $\mathbf{Z}/n\mathbf{Z}$-가군.
$F$에 대해 우리는 복합체 위의 경계를 선택한다.
$\mathbf{Z}/n\mathbf{Z}$-가군의 플랫(!) 층.
그런 다음 $F$가 표현되는 경우로 축소한다.
$\mathbf{Z}/n\mathbf{Z}$-가군의 평면 층의 경계 복합체에 의해.
이 시점에서 보조정리 \ref{etale-cohomology-lemma-punctual-base-change-upgrade}
시작하고 결론을 내립니다.
\end{proof}

\begin{lemma} \label{etale-cohomology-lemma-kunneth} $k$를 분리 가능하게 닫힌 체라고 가정한다. $X$ 및 $Y$를 $k$에 대해 유한 유형 스킴으로 설정한다. $n \geq 1$를 $k$에서 반전 가능한 정수로 둡니다. 그런 다음 $E \in D(X_\etale, \mathbf{Z}/n\mathbf{Z})$ 및 $K \in D(Y_\etale, \mathbf{Z}/n\mathbf{Z})$에 대해 $$
R\Gamma(X \times_{\Spec(k)} Y,
\text{pr}_1^{-1}E
\otimes_{\mathbf{Z}/n\mathbf{Z}}^\mathbf{L}
\text{pr}_2^{-1}K
)
=
R\Gamma(X, E)
\otimes_{\mathbf{Z}/n\mathbf{Z}}^\mathbf{L}
R\Gamma(Y, K)
$$ \end{lemma}가 있다.

\begin{proof}
보조정리 \ref{etale-cohomology-lemma-punctual-base-change-upgrade-unbounded}에 의해 우리는
$$
R\text{pr}_{1, *}(
\text{pr}_1^{-1}E
\otimes_{\mathbf{Z}/n\mathbf{Z}}^\mathbf{L}
\text{pr}_2^{-1}K) =
E \otimes_{\mathbf{Z}/n\mathbf{Z}}^\mathbf{L}
\underline{R\Gamma(Y, K)}
$$
우리는 다음과 같이 결론을 내립니다.
보조정리 \ref{etale-cohomology-lemma-pull-out-constant-mod-n}
우리가 사용할 수 있는 이유는
$\text{cd}(X) < \infty$ 에 의해 보조정리 \ref{etale-cohomology-lemma-finite-cd}.
\end{proof}














\section{카오스 위상과 자리스키 위상 비교} \label{etale-cohomology-section-compare-chaotic-Zariski}

\noindent
아핀 스킴의 구조층을 구성할 때 먼저
아핀 오픈에 대한 값을 구성한 다음 모든 오픈으로 확장한다.
비슷한 구성이 복합체 구성에 유용할 때가 많습니다.
스킴 $X$의 아벨 군. $X_{affine, Zar}$를 기억하세요.
위상이 주어진 $X$의 아핀 개방의 범주를 나타냅니다.
표준 Zariski 피복에 따라 참조
위상, 정의 \ref{topologies-definition-big-small-Zariski}.
우리는 독자들에게 $X_{affine, Zar}$의 토포스를 상기시킵니다.
$X$의 작은 Zariski 토포스이다. 위상, 보조정리를 참조하라.
\ref{topologies-lemma-alternative-zariski}.
이 절에서 우리는 $X_{affine}$ 동일한 기본을 나타냅니다.
범주 혼란스러운 위상, 즉 층
전층에 동의한다. 우리는 사이트의 사상을 얻습니다.
$$
\epsilon : X_{affine, Zar} \longrightarrow X_{affine}
$$
사이트, 절 \ref{sites-cohomology-section-compare}의 코호몰로지에서와 같습니다.

\begin{lemma}
\label{etale-cohomology-lemma-check-zar}
위의 상황에서 $K$를 $D^+(X_{affine})$의 객체로 둡니다.
그러면 $K$는 (완전히 충실한) 함자의 본질적인 이미지에 있다.
$R\epsilon_* ; D(X_{affine, Zar}) \to D(X_{affine})$ 경우에만
다음 두 가지 조건이 성립한다면
\begin{enumerate}
\item $R\Gamma(\emptyset, K)$는 $D(\textit{Ab})$에서 0이고
\item 만약 $U = V \cup W$ 와 함께 $U, V, W \subset X$ 아핀 개방 및
$V, W \subset U$ 표준 열림
(대수학, 정의 \ref{algebra-definition-Zariski-topology}), 그런 다음
지도 $c^K_{U, V, W, V \cap W}$
코호몰로지 사이트, 보조정리 \ref{sites-cohomology-lemma-c-square}
준동형사상이다.
\end{enumerate}
\end{lemma}

\begin{proof} (함자 $R\epsilon_*$는 사이트, 절 \ref{sites-cohomology-section-compare}에 대한 코호몰로지의 논의에 완전히 충실한다.) 집합, 이는 사이트, 보조정리 \ref{sites-cohomology-lemma-descent-squares}에 대한 매우 일반적인 코호몰로지에서 따릅니다. 그의 가설은 스킴, 보조정리 \ref{schemes-lemma-sheaf-on-affines} 및 코호몰로지 사이트, 보조정리 \ref{sites-cohomology-lemma-descent-squares-helper}에 있다.

\medskip\noindent
혼란를 피하려면 $X_{affine, almost-chaotic}$을 표시하라.
빈 객체 $\emptyset$에 피복이 두 개 있는 사이트,
즉, $\{\emptyset \to \emptyset\}$ 및 빈 피복
(사이트, 예 \ref{sites-example-site-topological} 참조)
논의). 그러면 사이트 중 사상이 있다.
$$
X_{affine, Zar} \to X_{affine, almost-chaotic} \to X_{affine}
$$
위의 인수는 첫 번째 화살표에 적용된다. 그럼 놔두자
독자가 그 물체를 볼 수 있도록$K$~의$D^+(X_{affine})$
(완전히 충실한) 함자의 본질적인 이미지에 있다.
$D(X_{affine, almost-chaotic}) \to D(X_{affine})$ 경우에만
$R\Gamma(\emptyset, K)$가 $D(\textit{Ab})$에서 0인 경우.
\end{proof}









\section{큰 것과 작은 것을 비교 토포스} \label{etale-cohomology-section-compare}

\noindent
$S$를 스킴으로 설정한다. ~ 안에
위상, 보조정리 \ref{topologies-lemma-at-the-bottom-etale}
비교 사상을 도입했다.
$\pi_S : (\Sch/S)_\etale \to S_\etale$ 그리고
$i_S : \Sh(S_\etale) \to \Sh((\Sch/S)_\etale)$
$\pi_S \circ i_S = \text{id}$ 및 $\pi_{S, *} = i_S^{-1}$를 사용한다.
보다 일반적으로 $f : T \to S$가 $(\Sch/S)_\etale$의 객체인 경우,
그러면 사상 $i_f : \Sh(T_\etale) \to \Sh((\Sch/S)_\etale)$가 있다.
$f_{small} = \pi_S \circ i_f$, 참조
위상, 기본정리 \ref{topologies-lemma-put-in-T-etale} 및
\ref{topologies-lemma-morphism-big-small-etale}. ~ 안에
하강, 비고 \ref{descent-remark-change-topologies-ringed}
우리는 이것을 링 사이트의 사상으로 확장했다.
$$
\pi_S : ((\Sch/S)_\etale, \mathcal{O}) \to (S_\etale, \mathcal{O}_S)
$$
그리고 사상 울린 토포스
$$
i_S : (\Sh(S_\etale), \mathcal{O}_S) \to (\Sh((\Sch/S)_\etale), \mathcal{O})
$$
그리고
$$
i_f : (\Sh(T_\etale), \mathcal{O}_T) \to (\Sh((\Sch/S)_\etale, \mathcal{O}))
$$
제한 사항 $i_S^{-1} = \pi_{S, *}$에 유의하세요.
위상, 정의 \ref{topologies-definition-restriction-small-etale})
$\mathcal{O}$을 $\mathcal{O}_S$로 변환한다.
마찬가지로 $i_f^{-1}$는 $\mathcal{O}$를 $\mathcal{O}_T$로 변환한다.
하강, 비고 \ref{descent-remark-change-topologies-ringed}를 참조하라.
따라서 $i_S^*\mathcal{F} = i_S^{-1}\mathcal{F}$ 및
$i_f^*\mathcal{F} = i_f^{-1}\mathcal{F}$ 모든 $\mathcal{O}$-가군에 대해
$\mathcal{F}$ - $(\Sch/S)_\etale$. 특히 $i_S^*$ 및 $i_f^*$
완전 함자이다. 함자 $i_S^*$는 종종 다음과 같이 표시된다.
$\mathcal{F} \mapsto \mathcal{F}|_{S_\etale}$ (그리고 이것은 그렇지 않습니다
표기법과 충돌
위상, 정의 \ref{topologies-definition-restriction-small-etale}).

\begin{lemma}
\label{etale-cohomology-lemma-compare-injectives}
$S$를 스킴으로 설정한다. $T$를 $(\Sch/S)_\etale$의 객체로 둡니다.
\begin{enumerate}
\item $\mathcal{I}$가 $\textit{Ab}((\Sch/S)_\etale)$에서 단사라면,
\begin{enumerate}
\item $i_f^{-1}\mathcal{I}$는 $\textit{Ab}(T_\etale)$에서 단사형이다.
\item $\mathcal{I}|_{S_\etale}$는 $\textit{Ab}(S_\etale)$에서 단사형이다.
\end{enumerate}
\item $\mathcal{I}^\bullet$가 K 단사 복합체인 경우
$\textit{Ab}((\Sch/S)_\etale)$에서
\begin{enumerate}
\item $i_f^{-1}\mathcal{I}^\bullet$는 K 단사 복합체이다.
$\textit{Ab}(T_\etale)$,
\item $\mathcal{I}^\bullet|_{S_\etale}$는 K 단사 복합체이다.
$\textit{Ab}(S_\etale)$,
\end{enumerate}
\end{enumerate}
가군에 해당하는 진술은 유지되지 않습니다.
\end{lemma}

\begin{proof} 부분 (1)(b) 및 (2)(b)는 함자 $\pi_{S, *} = i_S^{-1}$ 제한이 완전 함자의 오른쪽 수반라는 사실로부터 공식적으로 따릅니다. $\pi_S^{-1}$, 호몰로지, 보조정리 \ref{homology-lemma-adjoint-preserve-injectives} 및 유도 범주, 보조정리 \ref{derived-lemma-adjoint-preserve-K-injectives}를 참조하라.

\medskip\noindent
부품 (1)() 및 (2)()는 두 가지 방식으로 볼 수 있다. 첫 번째 증명: 다음을 사용할 수 있다.
$i_f^{-1}$은 완전 함자 $i_{f, !}$의 오른쪽 수반이다.
이 함자는 다음과 같이 구성된다.
위상, 보조정리 \ref{topologies-lemma-put-in-T-etale}
집합의 층 및 아벨 층의 경우
가군 사이트, 보조정리 \ref{sites-modules-lemma-g-shriek-adjoint}.
사이트, 보조정리의 가군에 표시된다.
\ref{sites-modules-lemma-exactness-lower-shriek} 정확한다.
두 번째 증명. 표시된 대로 $i_f = i_T \circ f_{big}$를 사용할 수 있다.
위상에서는 보조정리 \ref{topologies-lemma-morphism-big-small-etale}이다.
$f_{big}$는 국소화이므로 이에 따른 하락세를 볼 수 있다.
단사 및 K-단사를 보존한다. 참조
코호몰로지 사이트, 기본정리 \ref{sites-cohomology-lemma-cohomology-of-open} 및
\ref{sites-cohomology-lemma-restrict-K-injective-to-open}.
그런 다음 이미 입증된 부분 (1)(b) 및 (2)(b)를 적용한다.
함자 $i_T^{-1}$로 결론을 내립니다.

\medskip\noindent
$S = \Spec(\mathbf{Z})$를 두고 지도를 고려해보세요.
$2 : \mathcal{O}_S \to \mathcal{O}_S$. 이것은 주입식 지도이다.
$\mathcal{O}_S$-가군의 $S_\etale$. 그러나 철수는
$\pi_S^*(2) : \mathcal{O} \to \mathcal{O}$는 단사적이지 않습니다
$\Spec(\mathbf{F}_2)$를 평가하면 알 수 있다. 이제 선택하세요
단사 사상 $\alpha : \mathcal{O} \to \mathcal{I}$를 선택하되,
그 공역은 단사 $\mathcal{O}$-가군 $\mathcal{I}$이며 $(\Sch/S)_\etale$ 위에 있다고 하자.
그런 다음 다이어그램을 고려하라.
$$
\xymatrix{
\mathcal{O}_S \ar[d]_2 \ar[rr]_{\alpha|_{S_\etale}} & &
\mathcal{I}|_{S_\etale} \\
\mathcal{O}_S \ar@{..>}[rru]
}
$$
그러면 점선 화살표는 $\mathcal{O}_S$-가군의 범주에 존재할 수 없다.
왜냐하면 그것은 의미할 것이기 때문이다
(부속으로) 분사 맵 $\alpha$은 다음을 통해 고려된다.
비단사적 지도 $\pi_S^*(2)$는 그럴 수 없다.
따라서 $\mathcal{I}|_{S_\etale}$는 단사 $\mathcal{O}_S$-가군이 아닙니다.
\end{proof}

\noindent $f : T \to S$를 스킴의 사상으로 설정한다. 위상의 가환 도식, 보조정리 \ref{topologies-lemma-morphism-big-small-etale} (3)는 링형 사이트 $$
\xymatrix{
(T_\etale, \mathcal{O}_T) \ar[d]_{f_{small}} &
((\Sch/T)_\etale, \mathcal{O}) \ar[d]^{f_{big}} \ar[l]^{\pi_T} \\
(S_\etale, \mathcal{O}_S) &
((\Sch/S)_\etale, \mathcal{O}) \ar[l]_{\pi_S}
}
$$의 가환 도식으로 이어집니다. $f_{small}^\sharp$, $f_{big}^\sharp$, $\pi_S^\sharp$, 그리고 $\pi_T^\sharp$. 특히 이는 \begin{equation}
\label{etale-cohomology-equation-compare-big-small}
(f_{big, *}\mathcal{F})|_{S_\etale} =
f_{small, *}(\mathcal{F}|_{T_\etale})
\end{equation} 모든 층에 대해 $\mathcal{F}$ 위의 $(\Sch/T)_\etale$이고 $\mathcal{F}$가 $\mathcal{O}$-가군의 층인 경우 (\ref{etale-cohomology-equation-compare-big-small})는 동형사상의 $\mathcal{O}_S$-가군 $S_\etale$에 있다.

\begin{lemma}
\label{etale-cohomology-lemma-compare-higher-direct-image}
$f : T \to S$를 스킴의 사상으로 설정한다.
\begin{enumerate}
\item $K$의 $D((\Sch/T)_\etale)$에 대해
$
(Rf_{big, *}K)|_{S_\etale} = Rf_{small, *}(K|_{T_\etale})
$
$D(S_\etale)$에서.
\item $K$의 $D((\Sch/T)_\etale, \mathcal{O})$에 대해
$
(Rf_{big, *}K)|_{S_\etale} = Rf_{small, *}(K|_{T_\etale})
$
$D(\textit{Mod}(S_\etale, \mathcal{O}_S))$에서.
\end{enumerate}
보다 일반적으로 $g : S' \to S$를 $(\Sch/S)_\etale$의 객체로 둡니다.
올곱을 고려해보세요.
$$
\xymatrix{
T' \ar[r]_{g'} \ar[d]_{f'} & T \ar[d]^f \\
S' \ar[r]^g & S
}
$$
그 다음에
\begin{enumerate}
\item[(3)] 을 위한$K$~에$D((\Sch/T)_\etale)$우리는
$i_g^{-1}(Rf_{big, *}K) = Rf'_{small, *}(i_{g'}^{-1}K)$
$D(S'_\etale)$에서.
\item[(4)] 을 위한$K$~에$D((\Sch/T)_\etale, \mathcal{O})$우리는
$i_g^*(Rf_{big, *}K) = Rf'_{small, *}(i_{g'}^*K)$
$D(\textit{Mod}(S'_\etale, \mathcal{O}_{S'}))$에서.
\item[(5)] 을 위한$K$~에$D((\Sch/T)_\etale)$우리는
$g_{big}^{-1}(Rf_{big, *}K) = Rf'_{big, *}((g'_{big})^{-1}K)$
$D((\Sch/S')_\etale)$에서.
\item[(6)] 을 위한$K$~에$D((\Sch/T)_\etale, \mathcal{O})$우리는
$g_{big}^*(Rf_{big, *}K) = Rf'_{big, *}((g'_{big})^*K)$
$D(\textit{Mod}(S'_\etale, \mathcal{O}_{S'}))$에서.
\end{enumerate}
\end{lemma}

\begin{proof} 부분(1)은 보조정리 \ref{etale-cohomology-lemma-compare-injectives} 및 (\ref{etale-cohomology-equation-compare-big-small})에서 아벨 층의 K-주사 복합체를 선택하여 $K$를 나타냅니다.

\medskip\noindent 부분(3)은 보조정리 \ref{etale-cohomology-lemma-compare-injectives} 및 위상, 보조정리 \ref{topologies-lemma-morphism-big-small-cartesian-diagram-etale}에서 K-주사 복합체를 선택한다. 아벨 층은 $K$를 나타냅니다.

\medskip\noindent 부품(5)은 사이트, 보조정리 \ref{sites-cohomology-lemma-localize-cartesian-square}에서 코호몰로지이다.

\medskip\noindent 부품(6)은 사이트, 보조정리 \ref{sites-cohomology-lemma-localize-cartesian-square-modules}에서 코호몰로지이다.

\medskip\noindent
부분(2)은 다음과 같이 증명할 수 있다. 위에서 우리는 보았습니다.
그 $\pi_S \circ f_{big} = f_{small} \circ \pi_T$를 사상으로
울린 사이트. 그러므로 우리는 얻는다
$R\pi_{S, *} \circ Rf_{big, *} = Rf_{small, *} \circ R\pi_{T, *}$
사이트, 보조정리에서 코호몰로지에 의해
\ref{sites-cohomology-lemma-derived-pushforward-composition}.
제한 사항 함자 $\pi_{S, *}$ 및 $\pi_{T, *}$
정확하다고 결론을 내립니다.

\medskip\noindent 부분(4)은 부분(6)과 2)가 $f' : T' \to S'$에 적용된 부분에서 이어집니다. \end{proof}

\noindent
$S$를 스킴으로 두고 $\mathcal{H}$를 아벨 층으로 둡니다.
$(\Sch/S)_\etale$. $H^n_\etale(U, \mathcal{H})$는 다음을 나타냅니다.
그만큼코호몰로지~의$\mathcal{H}$물체 위에$U$~의$(\Sch/S)_\etale$.

\begin{lemma}
\label{etale-cohomology-lemma-compare-cohomology}
$f : T \to S$를 스킴의 사상으로 설정한다. 그 다음에
\begin{enumerate}
\item $K$의 $D(S_\etale)$에 대해
$H^n_\etale(S, \pi_S^{-1}K) = H^n(S_\etale, K)$.
\item $K$의 $D(S_\etale, \mathcal{O}_S)$에 대해
$H^n_\etale(S, L\pi_S^*K) = H^n(S_\etale, K)$.
\item $K$의 $D(S_\etale)$에 대해
$H^n_\etale(T, \pi_S^{-1}K) = H^n(T_\etale, f_{small}^{-1}K)$.
\item $K$의 $D(S_\etale, \mathcal{O}_S)$에 대해
$H^n_\etale(T, L\pi_S^*K) = H^n(T_\etale, Lf_{small}^*K)$.
\item $M$의 $D((\Sch/S)_\etale)$에 대해
$H^n_\etale(T, M) = H^n(T_\etale, i_f^{-1}M)$.
\item $M$의 $D((\Sch/S)_\etale, \mathcal{O})$에 대해
$H^n_\etale(T, M) = H^n(T_\etale, i_f^*M)$.
\end{enumerate}
\end{lemma}

\begin{proof} 증명하려면 (5) $M$를 아벨 층의 K-주사 복합체로 표현하고 보조정리 \ref{etale-cohomology-lemma-compare-injectives}를 적용하고 정의를 풀어보세요. 부분(3)은 $i_f^{-1}\pi_S^{-1} = f_{small}^{-1}$로 이어집니다. 부품(1)은 (3)의 특별한 경우이다.

\medskip\noindent
부분(6)은 사이트, 보조정리의 매우 일반적인 코호몰로지에서 이어집니다.
\ref{sites-cohomology-lemma-pullback-same-cohomology}. 그런 다음 부분
(4)는 $Lf_{small}^* = i_f^* \circ L\pi_S^*$ 때문에 따릅니다.
부분(2)은 (4)의 특별한 경우이다.
\end{proof}

\begin{lemma} \label{etale-cohomology-lemma-cohomological-descent-etale} $S$를 스킴으로 둡니다. $K \in D(S_\etale)$의 경우 지도 $$
K \longrightarrow R\pi_{S, *}\pi_S^{-1}K
$$는 동형사상이다. \end{lemma}

\begin{proof} 이는 $\pi_S^{-1}$ 및 $\pi_{S, *} = i_S^{-1}$가 모두 완전 함자이고 구성 $\pi_{S, *} \circ \pi_S^{-1}$가 항등 함자이기 때문에 참이다. \end{proof}

\begin{lemma}
\label{etale-cohomology-lemma-compare-higher-direct-image-proper}
$f : T \to S$를 스킴의 고유 사상으로 설정한다. 그럼 우리는
\begin{enumerate}
\item $\pi_S^{-1} \circ f_{small, *} = f_{big, *} \circ \pi_T^{-1}$
함자 $\Sh(T_\etale) \to \Sh((\Sch/S)_\etale)$로,
\item $\pi_S^{-1}Rf_{small, *}K = Rf_{big, *}\pi_T^{-1}K$
~을 위한$K$~에$D^+(T_\etale)$누구의코호몰로지 층비틀림이 있고,
\item $\pi_S^{-1}Rf_{small, *}K = Rf_{big, *}\pi_T^{-1}K$
~을 위한$K$~에$D(T_\etale, \mathbf{Z}/n\mathbf{Z})$, 그리고
\item $\pi_S^{-1}Rf_{small, *}K = Rf_{big, *}\pi_T^{-1}K$
모두를 위해$K$~에$D(T_\etale)$만약에$f$유한하다.
\end{enumerate}
\end{lemma}

\begin{proof} 증명 (1). $\mathcal{F}$를 $T_\etale$에서 층으로 설정한다. $g : S' \to S$를 $(\Sch/S)_\etale$의 객체로 둡니다. 올곱 $$
\xymatrix{
T' \ar[r]_{f'} \ar[d]_{g'} & S' \ar[d]^g \\
T \ar[r]^f & S
}
$$를 고려하면 $$
(f_{big, *}\pi_T^{-1}\mathcal{F})(S') =
(\pi_T^{-1}\mathcal{F})(T') =
((g'_{small})^{-1}\mathcal{F})(T')  =
(f'_{small, *}(g'_{small})^{-1}\mathcal{F})(S')
$$ 보조정리 \ref{etale-cohomology-lemma-describe-pullback}에 의한 두 번째 동일성을 갖습니다. 반면에 $$
(\pi_S^{-1}f_{small, *}\mathcal{F})(S') =
(g_{small}^{-1}f_{small, *}\mathcal{F})(S')
$$는 보조정리 \ref{etale-cohomology-lemma-describe-pullback}에 의해 다시 수행된다. 따라서 집합 (보조정리 \ref{etale-cohomology-lemma-proper-base-change-f-star})의 층에 대해 고유 기저변환에 의해 우리는 두 개의 집합이 표준적으로 동형이라고 결론을 내립니다. 동형사상은 제한 매핑과 호환되며 동형사상 $\pi_S^{-1}f_{small, *}\mathcal{F} = f_{big, *}\pi_T^{-1}\mathcal{F}$를 정의한다. 따라서 함자 $\pi_S^{-1} \circ f_{small, *} = f_{big, *} \circ \pi_T^{-1}$의 동형사상이다.

\medskip\noindent
증명 (2). 정식 베이스 변경 맵이 있다.
$\pi_S^{-1}Rf_{small, *}K \to Rf_{big, *}\pi_T^{-1}K$
누구에게나$K$~에$D(T_\etale)$, 보다
코호몰로지 사이트, 비고 \ref{sites-cohomology-remark-base-change}.
그것이 동형사상임을 증명하려면,
$i_g : \Sh(S'_\etale) \to \Sh((\Sch/S)_\etale)$에 의한 기본 변경 맵
은동형사상어떤 대상에 대해서$g : S' \to S$~의$(\Sch/S)_\etale$.
$T', g', f'$를 이전 단락과 같다고 가정한다.
거점 변경 맵의 당김은
\begin{align*}
g_{small}^{-1}Rf_{small, *}K
& =
i_g^{-1}\pi_S^{-1}Rf_{small, *}K \\
& \to
i_g^{-1}Rf_{big, *}\pi_T^{-1}K \\
& =
Rf'_{small, *}(i_{g'}^{-1}\pi_T^{-1}K) \\
& =
Rf'_{small, *}((g'_{small})^{-1}K)
\end{align*}
여기서는 $\pi_S \circ i_g = g_{small}$를 사용했다.
$\pi_T \circ i_{g'} = g'_{small}$ 그리고
보조정리 \ref{etale-cohomology-lemma-compare-higher-direct-image}.
이 지도는 고유 기저변환 정리에 의한 동형사상이다.
(보조정리 \ref{etale-cohomology-lemma-proper-base-change}) $K$가 제한되어 있는 경우
아래와 $K$의 코호몰로지 층은 비틀림이다.

\medskip\noindent 부품 (3)의 증명은 부품 (2)의 증명과 동일한다. 단, 대신 보조정리 \ref{etale-cohomology-lemma-proper-base-change-mod-n}를 사용한다. 보조정리 \ref{etale-cohomology-lemma-proper-base-change}.

\medskip\noindent
증명 (4). $f$가 유한하면 함자
$f_{small, *}$ 및 $f_{big, *}$는 정확한다. 이는 다음과 같습니다
명제 \ref{etale-cohomology-proposition-finite-higher-direct-image-zero}에서
~을 위한$f_{small}$. 기본 변경 이후$f'$~의$f$그것도 유한하고,
보조정리 \ref{etale-cohomology-lemma-compare-higher-direct-image} 부분 (3)에서 결론을 내립니다.
$f_{big, *}$도 정확합니다(더 높은 유도 함자가 0이므로).
따라서 이 경우는 (1) 부분에서 따릅니다.
\end{proof}




\section{fppf 및 에탈 위상 비교} \label{etale-cohomology-section-fppf-etale}

\noindent 이 절에 대한 모델은 사이트의 코호몰로지에서 논의된 바와 같이 로컬 컴팩트 위상 공간의 일반적인 위상과 qc 위상의 비교에 대한 절이다. 절 \ref{sites-cohomology-section-cohomology-LC}. 먼저 위상 절 \ref{topologies-section-change-topologies} 및 \ref{topologies-section-etale}의 일부 자료를 검토한다.

\medskip\noindent
$S$를 스킴으로 두고 $(\Sch/S)_{fppf}$를 fppf 사이트로 둡니다.
동일한 기본 범주에는 두 번째 위상이 있다.
즉 에탈 위상, 따라서 두 번째 사이트
$(\Sch/S)_\etale$. 신원 함자
$(\Sch/S)_\etale \to (\Sch/S)_{fppf}$는 연속적이며 다음을 정의한다.
사이트의 사상
$$
\epsilon_S : (\Sch/S)_{fppf} \longrightarrow (\Sch/S)_\etale
$$
사이트, 절 \ref{sites-cohomology-section-compare}의 코호몰로지를 참조하라.
$\epsilon_{S, *}$는 기본 ID 함자이다.
전층 그리고 $\epsilon_S^{-1}$는 에탈 층에 연결된다.
fppf 층화. $S_\etale$를 작은 에탈 사이트라 하자.
사이트 중 사상이 있다.
$$
\pi_S : (\Sch/S)_\etale \longrightarrow S_\etale
$$
연속 함자에 의해 주어진다.
$S_\etale \to (\Sch/S)_\etale$, $U \mapsto U$.
즉, $S_\etale$에는 올곱과 최종 객체가 있고
함자 위의 차량은 다음 차량으로 출퇴근한다.
사이트, 명제 \ref{sites-proposition-get-morphism}가 적용된다.

\begin{lemma}
\label{etale-cohomology-lemma-describe-pullback-pi-fppf}
위와 같이 표기한다.
$\mathcal{F}$를 $S_\etale$에서 층으로 설정한다. 규칙
$$
(\Sch/S)_{fppf} \longrightarrow \textit{Sets},\quad
(f : X \to S) \longmapsto \Gamma(X, f_{small}^{-1}\mathcal{F})
$$
는 층이고 $(\Sch/S)_\etale$의 경우는 층이다.
실제로 이 층은 다음과 같습니다.
$\pi_S^{-1}\mathcal{F}$ $(\Sch/S)_\etale$ 및
$\epsilon_S^{-1}\pi_S^{-1}\mathcal{F}$ - $(\Sch/S)_{fppf}$.
\end{lemma}

\begin{proof} 에탈 위상에 관한 진술은 보조정리 \ref{etale-cohomology-lemma-describe-pullback}의 내용이다. 증명을 마치려면 $\pi_S^{-1}\mathcal{F}$가 fppf 위상에 대한 층임을 보여주는 것으로 충분한다. 이는 보조정리 \ref{etale-cohomology-lemma-describe-pullback}에도 표시된다. \end{proof}

\noindent 보조정리 \ref{etale-cohomology-lemma-describe-pullback-pi-fppf}의 상황에서 $\epsilon_S$와 $\pi_S$의 구성과 동일성은 사이트 $$
a_S : (\Sch/S)_{fppf} \longrightarrow S_\etale
$$의 사상을 결정한다.

\begin{lemma} \label{etale-cohomology-lemma-push-pull-fppf-etale} 위와 같이 표기한다. $f : X \to Y$를 $(\Sch/S)_{fppf}$의 사상으로 설정한다. 그러면 가환 도식 중 토포스 $$
\xymatrix{
\Sh((\Sch/X)_{fppf}) \ar[rr]_{f_{big, fppf}} \ar[d]_{\epsilon_X} & &
\Sh((\Sch/Y)_{fppf}) \ar[d]^{\epsilon_Y} \\
\Sh((\Sch/X)_\etale) \ar[rr]^{f_{big, \etale}} & &
\Sh((\Sch/Y)_\etale)
}
$$와 $$
\xymatrix{
\Sh((\Sch/X)_{fppf}) \ar[rr]_{f_{big, fppf}} \ar[d]_{a_X} & &
\Sh((\Sch/Y)_{fppf}) \ar[d]^{a_Y} \\
\Sh(X_\etale) \ar[rr]^{f_{small}} & &
\Sh(Y_\etale)
}
$$가 있고 $a_X = \pi_X \circ \epsilon_X$와 $a_Y = \pi_X \circ \epsilon_X$가 있다. \end{lemma}

\begin{proof} 다이어그램의 교환성은 위상 절 \ref{topologies-section-change-topologies}의 논의를 따릅니다. \end{proof}

\begin{lemma} \label{etale-cohomology-lemma-proper-push-pull-fppf-etale} 보조정리 \ref{etale-cohomology-lemma-push-pull-fppf-etale}에서 $f$가 적절하다면 $a_Y^{-1} \circ f_{small, *} = f_{big, fppf, *} \circ a_X^{-1}$가 있다. \end{lemma}

\begin{proof}
다음의 증명을 반복하여 이를 증명할 수 있다.
보조정리 \ref{etale-cohomology-lemma-compare-higher-direct-image-proper} 부분 (1);
대신 우리는 이것으로부터 결과를 추론할 것이다.
$\epsilon_{Y, *}$는 기본 전층의 ID 함자이므로,
동형사상을 반영한다. 설명
보조정리 \ref{etale-cohomology-lemma-describe-pullback-pi-fppf}
$\epsilon_{Y, *} \circ a_Y^{-1} = \pi_Y^{-1}$을 보여줍니다.
$X$의 경우에도 마찬가지이다. 표준지도를 보여주기 위해
$a_Y^{-1}f_{small, *}\mathcal{F} \to f_{big, fppf, *}a_X^{-1}\mathcal{F}$
동형사상이면 다음을 보여주는 것으로 충분한다.
\begin{align*}
\pi_Y^{-1}f_{small, *}\mathcal{F}
& =
\epsilon_{Y, *}a_Y^{-1}f_{small, *}\mathcal{F} \\
& \to 
\epsilon_{Y, *}f_{big, fppf, *}a_X^{-1}\mathcal{F} \\
& =
f_{big, \etale, *} \epsilon_{X, *}a_X^{-1}\mathcal{F} \\
& =
f_{big, \etale, *}\pi_X^{-1}\mathcal{F}
\end{align*}
동형사상이다. 이것은 일부입니다
(1) 중 보조정리 \ref{etale-cohomology-lemma-compare-higher-direct-image-proper}이다.
\end{proof}

\begin{lemma} \label{etale-cohomology-lemma-descent-sheaf-fppf-etale} 보조정리 \ref{etale-cohomology-lemma-push-pull-fppf-etale}에서는 $f$가 단조롭고 국소적으로 유한하게 표현되고 전사라고 가정한다. 그러면 함자 $$
\Sh(Y_\etale) \longrightarrow
\left\{
(\mathcal{G}, \mathcal{H}, \alpha)
\middle|
\begin{matrix}
\mathcal{G} \in \Sh(X_\etale),\ \mathcal{H} \in \Sh((\Sch/Y)_{fppf}), \\
\alpha : a_X^{-1}\mathcal{G} \to f_{big, fppf}^{-1}\mathcal{H}
\text{동형사상}
\end{matrix}
\right\}
$$가 $\mathcal{F}$를 $(f_{small}^{-1}\mathcal{F}, a_Y^{-1}\mathcal{F}, can)$로 보내는 것은 동등한다. \end{lemma}

\begin{proof} 함자 $a_X^{-1}$는 완전히 충실합니다($a_{X, *}a_X^{-1} = \text{id}$ 에 의해 보조정리 \ref{etale-cohomology-lemma-describe-pullback-pi-fppf}). 따라서 망각하는 함자 $(\mathcal{G}, \mathcal{H}, \alpha) \mapsto \mathcal{H}$는 $\Sh((\Sch/Y)_{fppf})$의 전체 하위 범주로 트리플의 범주를 식별한다. 게다가 함자 $a_Y^{-1}$는 완전히 충실하므로 보조정리의 함자도 완전히 충실한다.

\medskip\noindent 에탈 피복 $\{Y_i \to Y\}$가 있다고 가정한다. $f_i : X_i \to Y_i$를 $f$의 기본 변경으로 둡니다. $f_{ij} = f_i \times f_j : X_i \times_X X_j  \to Y_i \times_Y Y_j$를 표시한다. 주장: 보조정리가 $f_i$에 대해 참이고 $f_{ij}$가 모든 $i, j$에 대해 참인 경우, 보조정리는 $f$에 대해 참이다. 이를 보려면 주어진 에탈 피복이 4개의 사이트 $Y_\etale, X_\etale, (\Sch/Y)_{fppf}, (\Sch/X)_{fppf}$ 각각에서 최종 객체의 에탈 피복을 결정한다는 점에 유의하세요. 따라서 층의 범주는 네 가지 경우 모두 이 피복(사이트, 보조정리 \ref{sites-lemma-mapping-property-glue})에 대한 접착 데이터의 범주와 동일한다. 범주의 거대한 가환 도식은 주장의 증명을 마무리한다. 자세한 내용은 생략한다. 주장은 $Y$에서 로컬로 에탈 작업을 수행할 수 있음을 나타냅니다.

\medskip\noindent
$\{X \to Y\}$는 fppf 피복이다. $Y$에서 로컬로 에탈 작업 중이다.
우리는 사상 $s : X' \to X$가 존재한다고 가정할 수 있다.
$f' = f \circ s : X' \to Y$는 지역적으로 자유인 전사 유한이다.
사상, 보조정리에 대한 추가 정보
\ref{more-morphisms-lemma-dominate-fppf-etale-locally}.
주장: $f'$에 대해 보조정리가 참이면 $f$에도 참이다.
즉, 트리플 $(\mathcal{G}, \mathcal{H}, \alpha)$이 주어지면
$f$의 경우 $s$로 당김하여 트리플을 얻을 수 있다.
$(s_{small}^{-1}\mathcal{G}, \mathcal{H}, s_{big, fppf}^{-1}\alpha)$
$f'$의 경우. 이 트리플에 대한 해는 층을 제공한다.
$\mathcal{F}$를 $Y_\etale$에서 $a_Y^{-1}\mathcal{F} = \mathcal{H}$로.
증명의 첫 번째 단락에서 이는 트리플이
본질적인 이미지에 이것은 우리를 다음과 같이 줄이다.
다음 단락에서 설명하는 경우이다.

\medskip\noindent
$f$는 국소적으로 자유인 전사 유한이라고 가정한다.
$(\mathcal{G}, \mathcal{H}, \alpha)$를 삼중으로 둡니다.
이 경우 트리플을 고려하라.
$$
(\mathcal{G}_1, \mathcal{H}_1, \alpha_1) =
(f_{small}^{-1}f_{small, *}\mathcal{G},
f_{big, fppf, *}f_{big, fppf}^{-1}\mathcal{H}, \alpha_1)
$$
여기서 $\alpha_1$는 ID에서 나옵니다.
\begin{align*}
a_X^{-1}f_{small}^{-1}f_{small, *}\mathcal{G}
& =
f_{big, fppf}^{-1}a_Y^{-1}f_{small, *}\mathcal{G} \\
& =
f_{big, fppf}^{-1}f_{big, fppf, *}a_X^{-1}\mathcal{G} \\
& \to
f_{big, fppf}^{-1}f_{big, fppf, *}f_{big, fppf}^{-1}\mathcal{H}
\end{align*}
여기서 세 번째 등식은 보조정리 \ref{etale-cohomology-lemma-proper-push-pull-fppf-etale}이다.
화살표는 $\alpha$로 표시된다.
이 트리플은 함자의 이미지에 있다.
$\mathcal{F}_1 = f_{small, *}\mathcal{F}$는 해결책입니다
(이를 보려면 보조정리 \ref{etale-cohomology-lemma-proper-push-pull-fppf-etale}를 다시 사용하세요.
자세한 내용은 생략) 트리플의 표준 맵이 있습니다
$$
(\mathcal{G}, \mathcal{H}, \alpha)
\to
(\mathcal{G}_1, \mathcal{H}_1, \alpha_1)
$$
$\text{id} \to f_{big, fppf, *}f_{big, fppf}^{-1}$ 단위를 사용한다.
두 번째 항목에 (사상을 처방하는 것으로 충분한다.
증명)의 첫 번째 단락에 의한 두 번째 항목이다. 부터
$\{f : X \to Y\}$는 fppf 피복 지도이다.
$\mathcal{H} \to \mathcal{H}_1$는 단사형입니다(자세한 내용은 생략됨).
두자
$$
\mathcal{G}_2 = \mathcal{G}_1 \amalg_\mathcal{G} \mathcal{G}_1\quad
\mathcal{H}_2 = \mathcal{H}_1 \amalg_\mathcal{H} \mathcal{H}_1
$$
그리고 $\alpha_2$를 유도된 동형사상 (당김 함자로 설정한다.
정확하므로 의미가 있습니다). 그러면 $\mathcal{H}$는
두 맵 $\mathcal{H}_1 \to \mathcal{H}_2$의 이퀄라이저이다.
트리플에 대해 위의 주장을 반복한다.
$(\mathcal{G}_2, \mathcal{H}_2, \alpha_2)$
우리는 삼중 동사 사상을 찾습니다.
$$
(\mathcal{G}_2, \mathcal{H}_2, \alpha_2)
\to
(\mathcal{G}_3, \mathcal{H}_3, \alpha_3)
$$
이 마지막 트리플은 함자의 이미지에 있다.
에 해당한다고 하세요.$\mathcal{F}_3$~에$\Sh(Y_\etale)$.
완전한 충실함으로 우리는 두 개의 지도를 얻습니다.
$\mathcal{F}_1 \to \mathcal{F}_3$ 그리고 우리는
$\mathcal{F}$ 이 두 맵의 이퀄라이저가 된다.
완전성의 철회 함자에 의해 우리는
원하는 대로 $a_Y^{-1}\mathcal{F} = \mathcal{H}$를 찾으세요.
\end{proof}

\begin{lemma}
\label{etale-cohomology-lemma-compare-fppf-etale}
비교를 고려하십시오 사상
$\epsilon : (\Sch/S)_{fppf} \to (\Sch/S)_\etale$.
$\mathcal{P}$는 스킴의 유한 사상 클래스를 나타냅니다.
을 위한$X$~에$(\Sch/S)_\etale$나타내다
$\mathcal{A}'_X \subset \textit{Ab}((\Sch/X)_\etale)$
형식의 층으로 구성된 전체 하위 카테고리
$\pi_X^{-1}\mathcal{F}$~와 함께$\mathcal{F}$~에$\textit{Ab}(X_\etale)$.
그런 다음 사이트, 성질의 코호몰로지
(\ref{sites-cohomology-item-base-change-P}),
(\ref{sites-cohomology-item-restriction-A}),
(\ref{sites-cohomology-item-A-sheaf}),
(\ref{sites-cohomology-item-A-and-P}) 및
(\ref{sites-cohomology-item-refine-tau-by-P})
사이트의 코호몰로지, 상황
\ref{sites-cohomology-situation-compare} 보류하세요.
\end{lemma}

\begin{proof}
먼저 $\mathcal{A}'_X \subset \textit{Ab}((\Sch/X)_\etale)$를 보여줍니다.
(1), (2), (3) 및 (4) 조건을 확인하여 약한 Serre 하위 범주이다.
호몰로지, 보조정리 \ref{homology-lemma-characterize-weak-serre-subcategory}.
부분 (1), (2), (3)는 $\pi_X^{-1}$가 정확하고
보조정리 \ref{etale-cohomology-lemma-cohomological-descent-etale}에 의해 예에 대해 완전히 충실한다. 만약에
$0 \to \pi_X^{-1}\mathcal{F} \to \mathcal{G} \to \pi_X^{-1}\mathcal{F}' \to 0$
$\textit{Ab}((\Sch/X)_\etale)$의 짧은 완전열이다.
그런 다음 $0 \to \mathcal{F} \to \pi_{X, *}\mathcal{G} \to \mathcal{F}' \to 0$
보조정리 \ref{etale-cohomology-lemma-cohomological-descent-etale}로 정확한다.
따라서 $\mathcal{G} = \pi_X^{-1}\pi_{X, *}\mathcal{G}$은
$\mathcal{A}'_X$ 최종 조건을 확인한다.

\medskip\noindent 코호몰로지 위의 사이트, 성질 (\ref{sites-cohomology-item-base-change-P})는 스킴의 올곱의 존재와 유한의 기본 변화가 사상의 스킴은 스킴의 유한 사상이다. 사상, 보조정리 \ref{morphisms-lemma-base-change-finite}를 참조하라.

\medskip\noindent 코호몰로지 위의 사이트, 성질 (\ref{sites-cohomology-item-restriction-A})는 위상의 가환 도식 (3)에서 따릅니다. 보조정리 \ref{topologies-lemma-morphism-big-small-etale}.

\medskip\noindent 코호몰로지는 사이트, 성질 (\ref{sites-cohomology-item-A-sheaf})는 보조정리 \ref{etale-cohomology-lemma-describe-pullback-pi-fppf}이다.

\medskip\noindent 코호몰로지 위의 사이트, 성질 (\ref{sites-cohomology-item-A-and-P})는 보조정리 \ref{etale-cohomology-lemma-compare-higher-direct-image-proper} 부분에 의해 유지된다. (4).

\medskip\noindent 코호몰로지 위의 사이트, 성질 (\ref{sites-cohomology-item-refine-tau-by-P})는 사상, 보조정리에 대한 추가 정보에 의해 암시된다. \ref{more-morphisms-lemma-dominate-fppf-etale-locally}. \end{proof}

\begin{lemma}
\label{etale-cohomology-lemma-V-C-all-n-etale-fppf}
위와 같이 표기한다.
\begin{enumerate}
\item $X \in \Ob((\Sch/S)_{fppf})$ 및 아벨 층 $\mathcal{F}$의 경우
$X_\etale$에 우리는
$\epsilon_{X, *}a_X^{-1}\mathcal{F} = \pi_X^{-1}\mathcal{F}$
그리고$R^i\epsilon_{X, *}(a_X^{-1}\mathcal{F}) = 0$~을 위한$i > 0$.
\item 유한 사상 $f : X \to Y$ $(\Sch/S)_{fppf}$의 경우
그리고 아벨 층 $\mathcal{F}$ $X$에
$a_Y^{-1}(R^if_{small, *}\mathcal{F}) =
R^if_{big, fppf, *}(a_X^{-1}\mathcal{F})$
모든 $i$에 대해.
\item 스킴 $X$ 및 $K$의 경우 $D^+(X_\etale)$ 지도
$\pi_X^{-1}K \to R\epsilon_{X, *}(a_X^{-1}K)$은 동형사상이다.
\item 사상 $f : X \to Y$ 스킴의 유한한 경우
그리고 $K$ 에는 $D^+(X_\etale)$ 가 있다.
$a_Y^{-1}(Rf_{small, *}K) = Rf_{big, fppf, *}(a_X^{-1}K)$.
\item 고유 사상 $f : X \to Y$의 스킴
그리고 $K$ 에서 $D^+(X_\etale)$ 비틀림 코호몰로지 층
$a_Y^{-1}(Rf_{small, *}K) = Rf_{big, fppf, *}(a_X^{-1}K)$.
\end{enumerate}
\end{lemma}

\begin{proof}
보조정리 \ref{etale-cohomology-lemma-compare-fppf-etale}에 의해 기본정리
코호몰로지 사이트, 절 \ref{sites-cohomology-section-compare-general}
모두 현재 설정에 적용된다. 결과를 번역하려면
범주 $\mathcal{A}_X$
코호몰로지 사이트, 보조정리 \ref{sites-cohomology-lemma-A}
필수적인 이미지이다
$a_X^{-1} : \textit{Ab}(X_\etale) \to \textit{Ab}((\Sch/X)_{fppf})$.

\medskip\noindent
부분 (1)는 다음과 같습니다.$(V_n)$모두를 위해$n$이는 다음과 같이 유지된다.
코호몰로지 사이트, 보조정리 \ref{sites-cohomology-lemma-V-C-all-n-general}.

\medskip\noindent
부분(2)은 다음의 결론에 $\epsilon_Y^{-1}$을 적용하여 이어집니다.
코호몰로지 사이트, 보조정리 \ref{sites-cohomology-lemma-V-implies-C-general}.

\medskip\noindent 부분(3)은 사이트, 보조정리 \ref{sites-cohomology-lemma-V-C-all-n-general} 부분(1)의 코호몰로지에서 이어집니다. $\pi_X^{-1}K$는 $D^+_{\mathcal{A}'_X}((\Sch/X)_\etale)$ 및 $a_X^{-1} = \epsilon_X^{-1} \circ a_X^{-1}$에 있다.

\medskip\noindent 부분(4)은 같은 이유로 사이트, 보조정리 \ref{sites-cohomology-lemma-V-C-all-n-general} 부분(2)의 코호몰로지에서 이어집니다.

\medskip\noindent
부품(5). 우리는 그것을 사용합니다
\begin{align*}
R\epsilon_{Y, *}Rf_{big, fppf, *}a_X^{-1}K
& =
Rf_{big, \etale, *}R\epsilon_{X, *}a_X^{-1}K \\
& =
Rf_{big, \etale, *}\pi_X^{-1}K \\
& =
\pi_Y^{-1}Rf_{small, *}K \\
& =
R\epsilon_{Y, *} a_Y^{-1}Rf_{small, *}K
\end{align*}
가환 도식에 의한 첫 번째 동등성
보조정리 \ref{etale-cohomology-lemma-push-pull-fppf-etale}
및 코호몰로지 위의 사이트, 보조정리
\ref{sites-cohomology-lemma-derived-pushforward-composition}.
두 번째 동등성은 (3)이다. 세 번째는
보조정리 \ref{etale-cohomology-lemma-compare-higher-direct-image-proper} 부분(2).
네 번째는 다시 (3)이다. 따라서 기본 변경 맵은
$a_Y^{-1}(Rf_{small, *}K) \to Rf_{big, fppf, *}(a_X^{-1}K)$
동형사상을 유도한다.
$$
R\epsilon_{Y, *}a_Y^{-1}Rf_{small, *}K \to
R\epsilon_{Y, *}Rf_{big, fppf, *}a_X^{-1}K
$$
증명은 다음 비고에 의해 완료된다.
$\alpha : a_Y^{-1}L \to M$ 와 $L$ 에서 $D^+(Y_\etale)$
그리고$M$~에$D^+((\Sch/Y)_{fppf})$그렇게$R\epsilon_{Y, *}\alpha$
은 동형사상이고, 은 동형사상이다. 즉,
$i$에 대한 귀납법을 통해 $H^i(\alpha)$가 동형사상임을 보여줍니다.
이는 충분히 작은 모든 $i$에 해당된다.
$i \leq i_0$에 대해 유지되면 다음을 알 수 있다.
$R^j\epsilon_{Y, *}H^i(M) = 0$ - $j > 0$ 및 $i \leq i_0$
에 의해 (1) 왜냐하면$H^i(M) = a_Y^{-1}H^i(L)$이 범위에서.
그러므로 $\epsilon_{Y, *}H^{i_0 + 1}(M) = H^{i_0 + 1}(R\epsilon_{Y, *}M)$
스펙트럼 열 인수로.
따라서 $\epsilon_{Y, *}H^{i_0 + 1}(M) = \pi_Y^{-1}H^{i_0 + 1}(L) =
\epsilon_{Y, *}a_Y^{-1}H^{i_0 + 1}(L)$.
이는 $H^{i_0 + 1}(\alpha)$가 동형사상임을 의미한다.
($\epsilon_{Y, *}$는 동형사상을 그대로 반영하기 때문이다.
원하는 대로 기본 전층)의 신원을 확인한다.
\end{proof}

\begin{lemma} \label{etale-cohomology-lemma-cohomological-descent-etale-fppf} $X$를 스킴으로 둡니다. $K \in D^+(X_\etale)$의 경우 맵 $$
K \longrightarrow Ra_{X, *}a_X^{-1}K
$$는 위와 같이 $a_X : \Sh((\Sch/X)_{fppf}) \to \Sh(X_\etale)$가 있는 동형사상이다. \end{lemma}

\begin{proof} 먼저 $K$가 단일 아벨 층으로 제공되는 경우로 설명을 축소한다. 즉, $K$를 복합체 $\mathcal{F}^\bullet$ 아래의 경계로 표현한다. 층의 경우 $\mathcal{F}^n = a_{X, *} a_X^{-1} \mathcal{F}^n$와 층 $R^qa_{X, *}a_X^{-1}\mathcal{F}^n$가 $q > 0$에 대해 0임을 알 수 있다. Leray의 비순환성 보조정리 (유도 범주, 보조정리 \ref{derived-lemma-leray-acyclicity})를 $a_X^{-1}\mathcal{F}^\bullet$ 및 함자 $a_{X, *}$에 적용하여 결론을 내립니다. 이제부터 $K = \mathcal{F}$라고 가정한다.

\medskip\noindent
보조정리 \ref{etale-cohomology-lemma-describe-pullback-pi-fppf}에 의해 우리는
$a_{X, *}a_X^{-1}\mathcal{F} = \mathcal{F}$. 따라서 다음을 보여주는 것으로 충분한다.
$R^qa_{X, *}a_X^{-1}\mathcal{F} = 0$~을 위한$q > 0$.
이를 위해 $a_X = \epsilon_X \circ \pi_X$를 사용할 수 있으며
르레이 스펙트럼 열
(코호몰로지 사이트, 보조정리 \ref{sites-cohomology-lemma-relative-Leray}).
보조정리 \ref{etale-cohomology-lemma-V-C-all-n-etale-fppf}에 의해
우리는$R^i\epsilon_{X, *}(a_X^{-1}\mathcal{F}) = 0$~을 위한$i > 0$
그리고
$\epsilon_{X, *}a_X^{-1}\mathcal{F} = \pi_X^{-1}\mathcal{F}$.
보조정리 \ref{etale-cohomology-lemma-cohomological-descent-etale}에 의해 우리는
$R^j\pi_{X, *}(\pi_X^{-1}\mathcal{F}) = 0$~을 위한$j > 0$.
이로써 증명이 종료된다.
\end{proof}

\begin{lemma} \label{etale-cohomology-lemma-compare-cohomology-etale-fppf} 스킴 $X$ 및 $a_X : \Sh((\Sch/X)_{fppf}) \to \Sh(X_\etale)$의 경우 위와 같습니다. \begin{enumerate} \item $H^q(X_\etale, \mathcal{F}) = H^q_{fppf}(X, a_X^{-1}\mathcal{F})$의 경우 아벨 층 $\mathcal{F}$는 $X_\etale$, \item $H^q(X_\etale, K) = H^q_{fppf}(X, a_X^{-1}K)$는 $K \in D^+(X_\etale)$에 대해. \end{enumerate} 예: $A$가 아벨 군이면 $H^q_\etale(X, \underline{A}) = H^q_{fppf}(X, \underline{A})$이다. \end{lemma}

\begin{proof}
이는 보조정리 \ref{etale-cohomology-lemma-cohomological-descent-etale-fppf}에서 따온 것이다.
사이트, 비고 \ref{sites-cohomology-remark-before-Leray}에서 코호몰로지에 의해.
\end{proof}







\section{fppf 및 에탈 위상 비교: 가군} \label{etale-cohomology-section-fppf-etale-modules}

\noindent 절 \ref{etale-cohomology-section-fppf-etale}에서 논의를 계속하지만 이 절에서는 가군층에 대해 간략하게 논의한다.

\medskip\noindent
$S$를 스킴으로 설정한다. 사이트 $\epsilon_S$, $\pi_S$의 사상 및
그들의 구성$a_S$에 소개절 \ref{etale-cohomology-section-fppf-etale}
링형 사이트의 사상이 자연스럽게 향상되었다. 첫 번째
다음과 같이 쓰여진다
$$
\epsilon_S :
((\Sch/S)_{fppf}, \mathcal{O})
\longrightarrow
((\Sch/S)_\etale, \mathcal{O})
$$
실제로 구조층에 동일한 기호를 사용할 수 있다.
층에는 동일한 기본 전층이 있다. 두 번째는
$$
\pi_S :
((\Sch/S)_\etale, \mathcal{O})
\longrightarrow
(S_\etale, \mathcal{O}_S)
$$
세 번째는 사상이다.
$$
a_S :
((\Sch/S)_{fppf}, \mathcal{O})
\longrightarrow
(S_\etale, \mathcal{O}_S)
$$
우리는 이미 준연접 가군의 범주가
스킴 $S$는 준연접 가군의 범주와 동일한다.
$(S_\etale, \mathcal{O}_S)$에서, 참조
하강, 명제 \ref{descent-proposition-equivalence-quasi-coherent}.
우리는 다음과 같은 비교에 관심이 있기 때문에
에탈 및 fppf 코호몰로지, 나머지 부분에서는
절 준연접층을 다음과 같이 생각해보세요.
작은 에탈 사이트.
준일관성에 대해 이미 알고 있는 내용을 검토해 보자.
가군 이 사이트에 대해.

\begin{lemma}
\label{etale-cohomology-lemma-review-quasi-coherent}
$S$를 스킴으로 설정한다. $\mathcal{F}$를 준일관되게 하자
$\mathcal{O}_S$-가군 $S_\etale$에 있다.
\begin{enumerate}
\item 규칙
$$
\mathcal{F}^a : (\Sch/S)_\etale \longrightarrow \textit{Ab},\quad
(f : T \to S) \longmapsto \Gamma(T, f_{small}^*\mathcal{F})
$$
fppf에 대한 층 조건과 fortiori 에탈 피복을 만족한다.
\item $\mathcal{F}^a = \pi_S^*\mathcal{F}$ $(\Sch/S)_\etale$에,
\item $\mathcal{F}^a = a_S^*\mathcal{F}$ $(\Sch/S)_{fppf}$에,
\item $\mathcal{F} \mapsto \mathcal{F}^a$ 규칙은 정의한다.
준일관성 $\mathcal{O}_S$-가군 간의 등가성
그리고 준연접 가군 켜짐
$((\Sch/S)_\etale, \mathcal{O})$,
\item $\mathcal{F} \mapsto \mathcal{F}^a$ 규칙은 정의한다.
준일관성 $\mathcal{O}_S$-가군 간의 등가성
그리고 준연접 가군 켜짐
$((\Sch/S)_{fppf}, \mathcal{O})$,
\item $\epsilon_{S, *}a_S^*\mathcal{F} = \pi_S^*\mathcal{F}$이 있다.
그리고 $a_{S, *}a_S^*\mathcal{F} = \mathcal{F}$,
\item $R^i\epsilon_{S, *}(a_S^*\mathcal{F}) = 0$이 있다.
그리고$R^ia_{S, *}(a_S^*\mathcal{F}) = 0$~을 위한$i > 0$.
\end{enumerate}
\end{lemma}

\begin{proof} 우리는 독자들이 강하, 절 \ref{descent-section-quasi-coherent-sheaves}, \ref{descent-section-quasi-coherent-cohomology} 및 \ref{descent-section-quasi-coherent-sheaves-bis}의 자료를 기반으로 이러한 결과 중 증명을 스스로 찾아볼 것을 촉구한다.

\medskip\noindent
먼저 이 보조정리의 표기법이 우리의 표기법과 일치하는 이유를 설명한다.
이전에 $\mathcal{F}^a$ 표기법을 사용했다.
절 \ref{etale-cohomology-section-quasi-coherent} 및
\ref{etale-cohomology-section-cohomology-quasi-coherent}
그리고
하강, 절 \ref{descent-section-quasi-coherent-sheaves}.
즉, 우리는 다음과 같이 알고 있다.
하강, 명제 \ref{descent-proposition-equivalence-quasi-coherent}
준연접 가군이 존재한다는 것
$\mathcal{F}_0$ 스킴 $S$ (즉, 작은
Zariski 사이트) $\mathcal{F}$는
규칙
$$
\mathcal{F}_0^a : (\Sch/S)_\etale \longrightarrow \textit{Ab},\quad
(f : T \to S) \longmapsto \Gamma(T, f^*\mathcal{F})
$$
하위 카테고리 $S_\etale \subset (\Sch/S)_\etale$
여기서 $f^*$는 스킴에서 가군층의 일반적인 하락을 나타냅니다.
$\mathcal{F}_0^a$은 링드의 사상에 의해 당김되므로
사이트
$$
((\Sch/S)_\etale, \mathcal{O}) \longrightarrow (S_{Zar}, \mathcal{O}_{S_{Zar}})
$$
하강 기준, 비고 \ref{descent-remark-change-topologies-ringed-sites}
(당김 구성에서) 즉시 다음과 같습니다.
$\mathcal{F}^a = \mathcal{F}_0^a$. 이는 층 성질
fpqc 피복의 경우에도
하강, 보조정리 \ref{descent-lemma-sheaf-condition-holds} (또한 참조
명제 \ref{etale-cohomology-proposition-quasi-coherent-sheaf-fpqc}).
그런 다음 (2) 및 (3)가 따릅니다.
다시 하강하여 비고 \ref{descent-remark-change-topologies-ringed-sites}
및 (4) 및 (5)는 다음과 같습니다.
하강, 명제 \ref{descent-proposition-equivalence-quasi-coherent}
(메타 결과도 참조하라.
정리 \ref{etale-cohomology-theorem-quasi-coherent}).

\medskip\noindent 부분(6)은 층 $\mathcal{F}^a = \pi_S^*\mathcal{F} = a_S^*\mathcal{F}$의 설명과 바로 연결된다.

\medskip\noindent
어떤 아벨리안에게도$\mathcal{H}$~에$(\Sch/S)_{fppf}$그만큼
고차 직상 $R^p\epsilon_{S, *}\mathcal{H}$는 층이다.
전층 $U \mapsto H^p_{fppf}(U, \mathcal{H})$와 연관됨
$(\Sch/S)_\etale$에 있다. 보다
코호몰로지 사이트, 보조정리 \ref{sites-cohomology-lemma-higher-direct-images}.
그러므로 증명하기 위해
$R^p\epsilon_{S, *}a_S^*\mathcal{F} = R^p\epsilon_{S, *}\mathcal{F}^a = 0$
~을 위한$p > 0$무엇이든 보여주면 충분하다스킴 $U$~ 위에$S$
에탈 피복 $\{U_i \to U\}_{i \in I}$이 있다.
$H^p_{fppf}(U_i, \mathcal{F}^a) = 0$~을 위한$p > 0$.
아핀으로 열린 피복을 취하면 필수
소멸은 평소 코호몰로지와 비교하여 따름
(하강, 명제 \ref{descent-proposition-same-cohomology-quasi-coherent}
또는
정리 \ref{etale-cohomology-theorem-zariski-fpqc-quasi-coherent})
그리고 소멸의 코호몰로지의 준연접층
스킴의 코호몰로지가 제공하는 아핀 스킴, 보조정리
\ref{coherent-lemma-quasi-coherent-affine-cohomology-zero}.

\medskip\noindent
$R^pa_{S, *}a_S^{-1}\mathcal{F} = R^pa_{S, *}\mathcal{F}^a = 0$임을 보여주기 위해
$p > 0$에 대해 우리는 정확히 같은 방식으로 주장한다. 이로써 증명이 완료된다.
\end{proof}

\begin{lemma} \label{etale-cohomology-lemma-cohomological-descent-etale-fppf-modules} $S$를 스킴으로 둡니다. $\mathcal{F}$의 경우 $\mathcal{O}_S$-가군의 $S_\etale$ 맵은 $$
\pi_S^*\mathcal{F} \longrightarrow R\epsilon_{S, *}(a_S^*\mathcal{F})
\quad\text{그리고}\quad
\mathcal{F} \longrightarrow Ra_{S, *}(a_S^*\mathcal{F})
$$와 동형사상이며 위와 같이 $a_S : \Sh((\Sch/S)_{fppf}) \to \Sh(S_\etale)$이다. \end{lemma}

\begin{proof}
이는 다음과 같은 즉각적인 결과이다.
부품 (6) 및 (7)
보조정리 \ref{etale-cohomology-lemma-review-quasi-coherent}.
\end{proof}

\begin{lemma}
\label{etale-cohomology-lemma-cohomological-descent-complex-modules}
$S = \Spec(A)$를 아핀 스킴으로 설정한다. $M^\bullet$를 복합체로 설정한다.
$A$-가군. 복합체 $\mathcal{F}^\bullet$를 고려하라.
전층 중 $\mathcal{O}$-가군 의
$(\textit{Aff}/S)_{fppf}$ 규칙에 따라 제공됨
$$
(U/S) = (\Spec(B)/\Spec(A)) \longmapsto M^\bullet \otimes_A B
$$
그러면 이것은 가군의 복합체이고 표준 맵이다.
$$
M^\bullet \longrightarrow
R\Gamma((\textit{Aff}/S)_{fppf}, \mathcal{F}^\bullet)
$$
준동형사상이다.
\end{lemma}

\begin{proof}
각 $\mathcal{F}^n$는 다음과 일치하는 가군층이다.
가군 $\mathcal{G}^n = (\widetilde{M}^n)^a$ 제한
보조정리 \ref{etale-cohomology-lemma-review-quasi-coherent} ~
$(\textit{Aff}/S)_{fppf} \subset (\Sch/S)_{fppf}$.
이 포함은 링 토포스의 동등성을 정의하므로
(위상, 보조정리 \ref{topologies-lemma-affine-big-site-fppf}),
우리는
$$
R\Gamma((\textit{Aff}/S)_{fppf}, \mathcal{F}^\bullet) =
R\Gamma((\Sch/S)_{fppf}, \mathcal{G}^\bullet)
$$
우리는 $M^\bullet = R\Gamma(S, \widetilde{M}^\bullet)$
예의 유도 범주에 의한 스킴, 보조정리
\ref{perfect-lemma-affine-compare-bounded}. 그러므로 우리는 노력하고 있습니다
비교 지도를 보여주기 위해
$$
R\Gamma(S, \widetilde{M}^\bullet)
\longrightarrow
R\Gamma((\Sch/S)_{fppf}, (\widetilde{M}^\bullet)^a)
$$
동형사상이다. $M^\bullet$이 아래로 제한되는 경우
하강으로 보유, 명제
\ref{descent-proposition-same-cohomology-quasi-coherent}
그리고 유도 범주, 보조정리의 첫 번째 스펙트럼 열
\ref{derived-lemma-two-ss-complex-functor}.
일반적인 경우에는 $M^\bullet = \lim M_n^\bullet$라고 씁니다.
$M_n^\bullet = \tau_{\geq -n}M^\bullet$로. 시스템은 어디서부터
$M_n^p$는 결국 $M^p$ 값으로 일정해집니다. 우리는 주장
$$
(\widetilde{M}^\bullet)^a = R\lim (\widetilde{M}_n^\bullet)^a
$$
즉, 자연지도가 왼쪽에서 오른쪽으로 표시되는 것으로 충분한다.
모든 아핀 객체에 대해 코호몰로지에 동형사상을 유도한다.
$U = \Spec(B)$ 중 $(\Sch/S)_{fppf}$이다. $i \in \mathbf{Z}$ 및
$n > |i|$ 우리는
$$
H^i(U, (\widetilde{M}_n^\bullet)^a) =
H^i(\tau_{\geq -n}M^\bullet \otimes_A B) =
H^i(M^\bullet \otimes_A B)
$$
위에서 처리된 제한된 아래 사례에 의해 첫 번째 동등성이 유지된다.
따라서 사이트, 보조정리의 코호몰로지에서 볼 수 있다.
\ref{sites-cohomology-lemma-RGamma-commutes-with-Rlim}
주장이 유지된다. 그러면 우리는 마침내 얻습니다.
\begin{align*}
R\Gamma((\Sch/S)_{fppf}, (\widetilde{M}^\bullet)^a)
& =
R\Gamma((\Sch/S)_{fppf}, R\lim (\widetilde{M}_n^\bullet)^a) \\
& =
R\lim R\Gamma((\Sch/S)_{fppf}, (\widetilde{M}_n^\bullet)^a) \\
& =
R\lim M_n^\bullet \\
& =
M^\bullet
\end{align*}
원하는대로. 두 번째 등식은 $R\lim$가 다음으로 통근하므로 유지된다.
$R\Gamma$, 사이트, 보조정리의 코호몰로지를 참조하라.
\ref{sites-cohomology-lemma-RGamma-commutes-with-Rlim}.
\end{proof}










\section{ph 및 에탈 위상 비교} \label{etale-cohomology-section-ph-etale}

\noindent 이 절에 대한 모델은 사이트의 코호몰로지에서 논의된 바와 같이 로컬 컴팩트 위상 공간의 일반적인 위상과 qc 위상의 비교에 대한 절이다. 절 \ref{sites-cohomology-section-cohomology-LC}. 먼저 위상 절 \ref{topologies-section-change-topologies} 및 \ref{topologies-section-etale}의 일부 자료를 검토한다.

\medskip\noindent
$S$는 스킴이고 $(\Sch/S)_{ph}$는 ph 사이트이다.
동일한 기본 범주에는 두 번째 위상이 있다.
즉 에탈 위상, 따라서 두 번째 사이트
$(\Sch/S)_\etale$. 신원 함자
$(\Sch/S)_\etale \to (\Sch/S)_{ph}$ 연속
(사상, 보조정리 \ref{more-morphisms-lemma-fppf-ph}에 대한 추가 정보 제공)
및 위상, 보조정리
\ref{topologies-lemma-zariski-etale-smooth-syntomic-fppf})
사이트의 사상을 정의한다.
$$
\epsilon_S : (\Sch/S)_{ph} \longrightarrow (\Sch/S)_\etale
$$
사이트, 절 \ref{sites-cohomology-section-compare}의 코호몰로지를 참조하라.
$\epsilon_{S, *}$는 기본 ID 함자이다.
전층 그리고 $\epsilon_S^{-1}$는 에탈 층에 연결된다.
ph 층화.
$S_\etale$를 작은 에탈 사이트라 하자.
사이트 중 사상이 있다.
$$
\pi_S : (\Sch/S)_\etale \longrightarrow S_\etale
$$
연속 함자에 의해 주어진다.
$S_\etale \to (\Sch/S)_\etale$, $U \mapsto U$.
즉, $S_\etale$에는 올곱과 최종 객체가 있고
함자 위의 차량은 다음 차량으로 출퇴근한다.
사이트, 명제 \ref{sites-proposition-get-morphism}가 적용된다.

\begin{lemma}
\label{etale-cohomology-lemma-describe-pullback-pi-ph}
위와 같이 표기한다.
$\mathcal{F}$를 $S_\etale$에서 층으로 설정한다. 규칙
$$
(\Sch/S)_{ph} \longrightarrow \textit{Sets},\quad
(f : X \to S) \longmapsto \Gamma(X, f_{small}^{-1}\mathcal{F})
$$
는 층이고 $(\Sch/S)_\etale$의 경우는 층이다.
실제로 이 층은 다음과 같습니다.
$\pi_S^{-1}\mathcal{F}$ $(\Sch/S)_\etale$ 및
$\epsilon_S^{-1}\pi_S^{-1}\mathcal{F}$ - $(\Sch/S)_{ph}$.
\end{lemma}

\begin{proof}
에탈 위상에 관한 진술은 다음과 같습니다.
보조정리 \ref{etale-cohomology-lemma-describe-pullback}. 증명을 끝내려면
$\pi_S^{-1}\mathcal{F}$가 ph에 대해 층임을 보여주기에 충분한다.
위상. 위상별, 보조정리 \ref{topologies-lemma-characterize-sheaf}
적절한 전사 사상이 주어지면 이를 보여주는 것으로 충분한다.
$V \to U$ 스킴 위에 $S$ 등화자 도식이 있다.
$$
\xymatrix{
(\pi_S^{-1}\mathcal{F})(U) \ar[r] &
(\pi_S^{-1}\mathcal{F})(V) \ar@<1ex>[r] \ar@<-1ex>[r] &
(\pi_S^{-1}\mathcal{F})(V \times_U V)
}
$$
두자 $\mathcal{G} = \pi_S^{-1}\mathcal{F}|_{U_\etale}$.
가환 도식을 고려해보세요.
$$
\xymatrix{
V \times_U V \ar[r] \ar[rd]_g \ar[d] & V \ar[d]^f \\
V \ar[r]^f & U
}
$$
우리는
$$
(\pi_S^{-1}\mathcal{F})(V) = \Gamma(V, f^{-1}\mathcal{G}) =
\Gamma(U, f_*f^{-1}\mathcal{G})
$$
여기서 $f_*$ 및 $f^{-1}$를 사용하여 사이의 기능을 나타냅니다.
작은 에탈 사이트. 둘째, 우리는
$$
(\pi_S^{-1}\mathcal{F})(V \times_U V) =
\Gamma(V \times_U V, g^{-1}\mathcal{G}) =
\Gamma(U, g_*g^{-1}\mathcal{G})
$$
등화자 도식에 있는 두 개의 지도는 두 개의 지도에서 나옵니다.
$$
f_*f^{-1}\mathcal{G} \longrightarrow g_*g^{-1}\mathcal{G}
$$
따라서 $\mathcal{G}$이 다음과 같다는 것을 증명하는 것으로 충분한다.
층의 두 맵의 이퀄라이저이다.
$\overline{u}$를 $U$의 기하점으로 설정한다. 집합
$\Omega = \mathcal{G}_{\overline{u}}$.
$\overline{u}$에서 줄기를 사용하여
보조정리 \ref{etale-cohomology-lemma-proper-pushforward-stalk}
우리는 두 개의 지도를 얻습니다
$$
H^0(V_{\overline{u}}, \underline{\Omega}) \longrightarrow
H^0((V \times_U V)_{\overline{u}}, \underline{\Omega}) =
H^0(V_{\overline{u}} \times_{\overline{u}} V_{\overline{u}},
\underline{\Omega})
$$
여기서 $\underline{\Omega}$는 상수층 값을 나타냅니다.
$\Omega$. 물론 이 지도들은 프로젝션 지도의 후퇴이다.
그러면 당김에서 절이 나오는 것이 분명한다.
첫 번째 요소에 투영하면 올에서 일정한다.
첫 번째 투영 및 당김에서 나오는 절
첫 번째 요소에 투영하면 올에서 일정한다.
첫 번째 프로젝션. 이미지 교차점의 절
이러한 당김 맵은 모든 항목에서 일정한다.
$V_{\overline{u}} \times_{\overline{u}} V_{\overline{u}}$ 즉,
이는 원하는 대로 $\Omega$의 요소에서 나옵니다.
\end{proof}

\noindent 보조정리 \ref{etale-cohomology-lemma-describe-pullback-pi-ph}의 상황에서 $\epsilon_S$와 $\pi_S$의 구성과 동일성은 사이트 $$
a_S : (\Sch/S)_{ph} \longrightarrow S_\etale
$$의 사상을 결정한다.

\begin{lemma} \label{etale-cohomology-lemma-push-pull-ph-etale} 위와 같이 표기한다. $f : X \to Y$를 $(\Sch/S)_{ph}$의 사상으로 설정한다. 그러면 가환 도식 중 토포스 $$
\xymatrix{
\Sh((\Sch/X)_{ph}) \ar[rr]_{f_{big, ph}} \ar[d]_{\epsilon_X} & &
\Sh((\Sch/Y)_{ph}) \ar[d]^{\epsilon_Y} \\
\Sh((\Sch/X)_\etale) \ar[rr]^{f_{big, \etale}} & &
\Sh((\Sch/Y)_\etale)
}
$$와 $$
\xymatrix{
\Sh((\Sch/X)_{ph}) \ar[rr]_{f_{big, ph}} \ar[d]_{a_X} & &
\Sh((\Sch/Y)_{ph}) \ar[d]^{a_Y} \\
\Sh(X_\etale) \ar[rr]^{f_{small}} & &
\Sh(Y_\etale)
}
$$가 있고 $a_X = \pi_X \circ \epsilon_X$와 $a_Y = \pi_X \circ \epsilon_X$가 있다. \end{lemma}

\begin{proof} 다이어그램의 교환성은 위상 절 \ref{topologies-section-change-topologies}의 논의를 따릅니다. \end{proof}

\begin{lemma} \label{etale-cohomology-lemma-proper-push-pull-ph-etale} 보조정리 \ref{etale-cohomology-lemma-push-pull-ph-etale}에서 $f$가 적절하다면 $a_Y^{-1} \circ f_{small, *} = f_{big, ph, *} \circ a_X^{-1}$가 있다. \end{lemma}

\begin{proof}
다음의 증명을 반복하여 이를 증명할 수 있다.
보조정리 \ref{etale-cohomology-lemma-compare-higher-direct-image-proper} 부분 (1);
대신 우리는 이것으로부터 결과를 추론할 것이다.
$\epsilon_{Y, *}$는 기본 전층의 ID 함자이므로,
동형사상을 반영한다. 설명
보조정리 \ref{etale-cohomology-lemma-describe-pullback-pi-ph}
$\epsilon_{Y, *} \circ a_Y^{-1} = \pi_Y^{-1}$을 보여줍니다.
$X$의 경우에도 마찬가지이다. 표준지도를 보여주기 위해
$a_Y^{-1}f_{small, *}\mathcal{F} \to f_{big, ph, *}a_X^{-1}\mathcal{F}$
동형사상이면 다음을 보여주는 것으로 충분한다.
\begin{align*}
\pi_Y^{-1}f_{small, *}\mathcal{F}
& =
\epsilon_{Y, *}a_Y^{-1}f_{small, *}\mathcal{F} \\
& \to 
\epsilon_{Y, *}f_{big, ph, *}a_X^{-1}\mathcal{F} \\
& =
f_{big, \etale, *} \epsilon_{X, *}a_X^{-1}\mathcal{F} \\
& =
f_{big, \etale, *}\pi_X^{-1}\mathcal{F}
\end{align*}
동형사상이다. 이것은 일부입니다
(1) 중 보조정리 \ref{etale-cohomology-lemma-compare-higher-direct-image-proper}이다.
\end{proof}

\begin{lemma}
\label{etale-cohomology-lemma-compare-ph-etale}
비교를 고려하십시오 사상
$\epsilon : (\Sch/S)_{ph} \to (\Sch/S)_\etale$.
$\mathcal{P}$는 스킴의 고유 사상 클래스를 나타냅니다.
을 위한$X$~에$(\Sch/S)_\etale$나타내다
$\mathcal{A}'_X \subset \textit{Ab}((\Sch/X)_\etale)$
형식의 층으로 구성된 전체 하위 카테고리
$\pi_X^{-1}\mathcal{F}$ 여기서 $\mathcal{F}$는
비틀림 아벨 층 $X_\etale$
그런 다음 사이트, 성질의 코호몰로지
(\ref{sites-cohomology-item-base-change-P}),
(\ref{sites-cohomology-item-restriction-A}),
(\ref{sites-cohomology-item-A-sheaf}),
(\ref{sites-cohomology-item-A-and-P}) 및
(\ref{sites-cohomology-item-refine-tau-by-P})
사이트의 코호몰로지, 상황
\ref{sites-cohomology-situation-compare} 보류하세요.
\end{lemma}

\begin{proof}
먼저 $\mathcal{A}'_X \subset \textit{Ab}((\Sch/X)_\etale)$를 보여줍니다.
(1), (2), (3) 및 (4) 조건을 확인하여 약한 Serre 하위 범주이다.
호몰로지, 보조정리 \ref{homology-lemma-characterize-weak-serre-subcategory}.
부분 (1), (2), (3)는 $\pi_X^{-1}$가 정확하고
보조정리 \ref{etale-cohomology-lemma-cohomological-descent-etale}에 의해 예에 대해 완전히 충실한다. 만약에
$0 \to \pi_X^{-1}\mathcal{F} \to \mathcal{G} \to \pi_X^{-1}\mathcal{F}' \to 0$
$\textit{Ab}((\Sch/X)_\etale)$의 짧은 완전열이다.
그런 다음 $0 \to \mathcal{F} \to \pi_{X, *}\mathcal{G} \to \mathcal{F}' \to 0$
보조정리 \ref{etale-cohomology-lemma-cohomological-descent-etale}로 정확한다.
특히 우리는 $\pi_{X, *}\mathcal{G}$가 아벨리안임을 알 수 있다.
꼬임층 - $X_\etale$.
따라서 $\mathcal{G} = \pi_X^{-1}\pi_{X, *}\mathcal{G}$은
$\mathcal{A}'_X$ 최종 조건을 확인한다.

\medskip\noindent 코호몰로지 위의 사이트, 성질 (\ref{sites-cohomology-item-base-change-P})는 스킴의 올곱이 존재하고 의 기본 변경이 고유 사상의 스킴은 고유 사상의 스킴이다. 사상, 보조정리 \ref{morphisms-lemma-base-change-proper}를 참조하라.

\medskip\noindent 코호몰로지 위의 사이트, 성질 (\ref{sites-cohomology-item-restriction-A})는 위상의 가환 도식 (3)에서 따릅니다. 보조정리 \ref{topologies-lemma-morphism-big-small-etale}.

\medskip\noindent 코호몰로지는 사이트, 성질 (\ref{sites-cohomology-item-A-sheaf})는 보조정리 \ref{etale-cohomology-lemma-describe-pullback-pi-ph}이다.

\medskip\noindent
코호몰로지 사이트, 성질 (\ref{sites-cohomology-item-A-and-P})에 의해 유지된다.
보조정리 \ref{etale-cohomology-lemma-compare-higher-direct-image-proper} 부분 (2)
그리고 $R^if_{small}\mathcal{F}$
$\mathcal{F}$가 아벨 비틀림인 경우 비틀림이다.
층 $X_\etale$에 대한 내용은 보조정리 \ref{etale-cohomology-lemma-torsion-direct-image}를 참조하라.

\medskip\noindent 코호몰로지 위의 사이트, 성질 (\ref{sites-cohomology-item-refine-tau-by-P})는 사상, 보조정리 \ref{more-morphisms-lemma-dominate-fppf-etale-locally}에 대한 추가 정보를 따릅니다. 유한 사상이 적절하고 전사 고유 사상이 ph 피복을 정의한다는 사실과 결합된다. 위상, 보조정리 \ref{topologies-lemma-surjective-proper-ph}를 참조하라. \end{proof}

\begin{lemma}
\label{etale-cohomology-lemma-V-C-all-n-etale-ph}
위와 같이 표기한다.
\begin{enumerate}
\item $X \in \Ob((\Sch/S)_{ph})$ 및 아벨 꼬임층 $\mathcal{F}$
$X_\etale$에 우리는
$\epsilon_{X, *}a_X^{-1}\mathcal{F} = \pi_X^{-1}\mathcal{F}$
그리고$R^i\epsilon_{X, *}(a_X^{-1}\mathcal{F}) = 0$~을 위한$i > 0$.
\item 고유 사상 $f : X \to Y$의 $(\Sch/S)_{ph}$
그리고 abelian 꼬임층 $\mathcal{F}$ 위의 $X$ 우리는
$a_Y^{-1}(R^if_{small, *}\mathcal{F}) =
R^if_{big, ph, *}(a_X^{-1}\mathcal{F})$
모든 $i$에 대해.
\item 스킴 $X$와 대상 $K$ 중 꼬임 코호몰로지층을 갖는 $D^+(X_\etale)$의 것에 대하여,
코호몰로지 층 지도
$\pi_X^{-1}K \to R\epsilon_{X, *}(a_X^{-1}K)$은 동형사상이다.
\item 고유 사상 $f : X \to Y$의 스킴
그리고 $K$ 에서 $D^+(X_\etale)$ 비틀림 코호몰로지 층
$a_Y^{-1}(Rf_{small, *}K) = Rf_{big, ph, *}(a_X^{-1}K)$.
\end{enumerate}
\end{lemma}

\begin{proof}
보조정리 \ref{etale-cohomology-lemma-compare-ph-etale}에 의해 기본정리
코호몰로지 사이트, 절 \ref{sites-cohomology-section-compare-general}
모두 현재 설정에 적용된다. 결과를 번역하려면
범주 $\mathcal{A}_X$
코호몰로지 사이트, 보조정리 \ref{sites-cohomology-lemma-A}
$\textit{Ab}((\Sch/X)_{ph})$의 전체 하위 범주이다.
$a_X^{-1}\mathcal{F}$ 형식의 층으로 구성됨
여기서 $\mathcal{F}$는 $X_\etale$의 아벨 꼬임층이다.

\medskip\noindent
부분 (1)는 다음과 같습니다.$(V_n)$모두를 위해$n$이는 다음과 같이 유지된다.
코호몰로지 사이트, 보조정리 \ref{sites-cohomology-lemma-V-C-all-n-general}.

\medskip\noindent
부분(2)은 다음의 결론에 $\epsilon_Y^{-1}$을 적용하여 이어집니다.
코호몰로지 사이트, 보조정리 \ref{sites-cohomology-lemma-V-implies-C-general}.

\medskip\noindent 부분(3)은 사이트, 보조정리 \ref{sites-cohomology-lemma-V-C-all-n-general} 부분(1)의 코호몰로지에서 이어집니다. $\pi_X^{-1}K$는 $D^+_{\mathcal{A}'_X}((\Sch/X)_\etale)$ 및 $a_X^{-1} = \epsilon_X^{-1} \circ a_X^{-1}$에 있다.

\medskip\noindent 부분(4)은 같은 이유로 사이트, 보조정리 \ref{sites-cohomology-lemma-V-C-all-n-general} 부분(2)의 코호몰로지에서 이어집니다. \end{proof}

\begin{lemma} \label{etale-cohomology-lemma-cohomological-descent-etale-ph} $X$를 스킴으로 둡니다. 비틀림 코호몰로지 층이 있는 $K \in D^+(X_\etale)$의 경우 맵 $$
K \longrightarrow Ra_{X, *}a_X^{-1}K
$$는 위와 같이 $a_X : \Sh((\Sch/X)_{ph}) \to \Sh(X_\etale)$가 있는 동형사상이다. \end{lemma}

\begin{proof}
먼저 진술을 다음과 같은 경우로 축소한다.
$K$는 단일 아벨 층으로 제공된다. 즉, $K$를 나타낸다.
비틀림의 복합체 $\mathcal{F}^\bullet$ 아래 경계에 의해
아벨 층. 이는 사이트, 보조정리에서 코호몰로지로 가능한다.
\ref{sites-cohomology-lemma-torsion}. 의 경우
층 알겠습니다
$\mathcal{F}^n = a_{X, *} a_X^{-1} \mathcal{F}^n$
그리고 층 $R^qa_{X, *}a_X^{-1}\mathcal{F}^n$
$q > 0$에 대해서는 0이다. Leray의 비주기성 보조정리
(유도 범주, 보조정리 \ref{derived-lemma-leray-acyclicity})
$a_X^{-1}\mathcal{F}^\bullet$에 적용됨
그리고 함자 $a_{X, *}$ 결론을 내립니다. 이제부터 $K = \mathcal{F}$로 가정한다.
여기서 $\mathcal{F}$는 비틀림 아벨 층이다.

\medskip\noindent
보조정리 \ref{etale-cohomology-lemma-describe-pullback-pi-ph}에 의해 우리는
$a_{X, *}a_X^{-1}\mathcal{F} = \mathcal{F}$. 따라서 다음을 보여주는 것으로 충분한다.
$R^qa_{X, *}a_X^{-1}\mathcal{F} = 0$~을 위한$q > 0$.
이를 위해 $a_X = \epsilon_X \circ \pi_X$를 사용할 수 있으며
르레이 스펙트럼 열
(코호몰로지 사이트, 보조정리 \ref{sites-cohomology-lemma-relative-Leray}).
보조정리 \ref{etale-cohomology-lemma-V-C-all-n-etale-ph}에 의해
우리는$R^i\epsilon_{X, *}(a_X^{-1}\mathcal{F}) = 0$~을 위한$i > 0$
및 $\epsilon_{X, *}a_X^{-1}\mathcal{F} = \pi_X^{-1}\mathcal{F}$.
보조정리 \ref{etale-cohomology-lemma-cohomological-descent-etale}에 의해 우리는
$R^j\pi_{X, *}(\pi_X^{-1}\mathcal{F}) = 0$~을 위한$j > 0$.
이로써 증명이 종료된다.
\end{proof}

\begin{lemma} \label{etale-cohomology-lemma-compare-cohomology-etale-ph} 위와 같이 스킴 $X$ 및 $a_X : \Sh((\Sch/X)_{ph}) \to \Sh(X_\etale)$의 경우: \begin{enumerate} \item $H^q(X_\etale, \mathcal{F}) = H^q_{ph}(X, a_X^{-1}\mathcal{F})$ 비틀림에 대한 아벨 층 $\mathcal{F}$ $X_\etale$, \item $H^q(X_\etale, K) = H^q_{ph}(X, a_X^{-1}K)$에 대한 $K \in D^+(X_\etale)$ 비틀림 코호몰로지 층. \end{enumerate} 예: $A$가 비틀림 아벨 군이면 $H^q_\etale(X, \underline{A}) = H^q_{ph}(X, \underline{A})$이다. \end{lemma}

\begin{proof}
이는 보조정리 \ref{etale-cohomology-lemma-cohomological-descent-etale-ph}에서 따온 것이다.
사이트, 비고 \ref{sites-cohomology-remark-before-Leray}에서 코호몰로지에 의해.
\end{proof}












\section{h 및 에탈 위상 비교} \label{etale-cohomology-section-h-etale}

\noindent 이 절에 대한 모델은 사이트의 코호몰로지에서 논의된 바와 같이 로컬 컴팩트 위상 공간의 일반적인 위상과 qc 위상의 비교에 대한 절이다. 절 \ref{sites-cohomology-section-cohomology-LC}. 게다가 이 절은 ph와 에탈 위상을 비교한 절과 거의 단어 그대로이다. 먼저 위상 절 \ref{topologies-section-change-topologies} 및 \ref{topologies-section-etale}와 평탄성에 대한 추가 정보 절 \ref{flat-section-h}의 일부 자료를 검토한다.

\medskip\noindent
$S$를 스킴으로 두고 $(\Sch/S)_h$를 h 사이트로 둡니다.
동일한 기본 범주에는 두 번째 위상이 있다.
즉 에탈 위상, 따라서 두 번째 사이트
$(\Sch/S)_\etale$. 신원 함자
$(\Sch/S)_\etale \to (\Sch/S)_h$ 연속
(평탄도에 대한 추가 정보, 보조정리 \ref{flat-lemma-zariski-h}
및 위상, 보조정리
\ref{topologies-lemma-zariski-etale-smooth-syntomic-fppf})
사이트의 사상을 정의한다.
$$
\epsilon_S : (\Sch/S)_h \longrightarrow (\Sch/S)_\etale
$$
사이트, 절 \ref{sites-cohomology-section-compare}의 코호몰로지를 참조하라.
$\epsilon_{S, *}$는 기본 ID 함자이다.
전층 그리고 $\epsilon_S^{-1}$는 에탈 층에 연결된다.
h 층화. $S_\etale$를 작은 에탈 사이트라 하자.
사이트 중 사상이 있다.
$$
\pi_S : (\Sch/S)_\etale \longrightarrow S_\etale
$$
연속 함자에 의해 주어진다.
$S_\etale \to (\Sch/S)_\etale$, $U \mapsto U$.
즉, $S_\etale$에는 올곱과 최종 객체가 있고
함자 위의 차량은 다음 차량으로 출퇴근한다.
사이트, 명제 \ref{sites-proposition-get-morphism}가 적용된다.

\begin{lemma}
\label{etale-cohomology-lemma-describe-pullback-pi-h}
위와 같이 표기한다.
$\mathcal{F}$를 $S_\etale$에서 층으로 설정한다. 규칙
$$
(\Sch/S)_h \longrightarrow \textit{Sets},\quad
(f : X \to S) \longmapsto \Gamma(X, f_{small}^{-1}\mathcal{F})
$$
는 층이고 $(\Sch/S)_\etale$의 경우는 층이다.
실제로 이 층은 다음과 같습니다.
$\pi_S^{-1}\mathcal{F}$ $(\Sch/S)_\etale$ 및
$\epsilon_S^{-1}\pi_S^{-1}\mathcal{F}$ - $(\Sch/S)_h$.
\end{lemma}

\begin{proof} 에탈 위상에 관한 진술은 보조정리 \ref{etale-cohomology-lemma-describe-pullback}의 내용이다. 증명을 마치려면 $\pi_S^{-1}\mathcal{F}$가 h 위상에 대한 층임을 보여주는 것으로 충분한다. 그러나 보조정리 \ref{etale-cohomology-lemma-describe-pullback-pi-ph}에서는 더 강한 pH 위상에서도 $\pi_S^{-1}\mathcal{F}$가 층임을 보여주었습니다. \end{proof}

\noindent 보조정리 \ref{etale-cohomology-lemma-describe-pullback-pi-h}의 상황에서 $\epsilon_S$와 $\pi_S$의 구성과 동일성은 사이트 $$
a_S : (\Sch/S)_h \longrightarrow S_\etale
$$의 사상을 결정한다.

\begin{lemma} \label{etale-cohomology-lemma-push-pull-h-etale} 위와 같이 표기한다. $f : X \to Y$를 $(\Sch/S)_h$의 사상으로 설정한다. 그러면 가환 도식 중 토포스 $$
\xymatrix{
\Sh((\Sch/X)_h) \ar[rr]_{f_{big, h}} \ar[d]_{\epsilon_X} & &
\Sh((\Sch/Y)_h) \ar[d]^{\epsilon_Y} \\
\Sh((\Sch/X)_\etale) \ar[rr]^{f_{big, \etale}} & &
\Sh((\Sch/Y)_\etale)
}
$$와 $$
\xymatrix{
\Sh((\Sch/X)_h) \ar[rr]_{f_{big, h}} \ar[d]_{a_X} & &
\Sh((\Sch/Y)_h) \ar[d]^{a_Y} \\
\Sh(X_\etale) \ar[rr]^{f_{small}} & &
\Sh(Y_\etale)
}
$$가 있고 $a_X = \pi_X \circ \epsilon_X$와 $a_Y = \pi_X \circ \epsilon_X$가 있다. \end{lemma}

\begin{proof} 다이어그램의 교환성은 위상 절 \ref{topologies-section-change-topologies}에서 설명한 것과 유사한다. \end{proof}

\begin{lemma} \label{etale-cohomology-lemma-proper-push-pull-h-etale} 보조정리 \ref{etale-cohomology-lemma-push-pull-h-etale}에서 $f$가 적절하다면 $a_Y^{-1} \circ f_{small, *} = f_{big, h, *} \circ a_X^{-1}$가 있다. \end{lemma}

\begin{proof}
다음의 증명을 반복하여 이를 증명할 수 있다.
보조정리 \ref{etale-cohomology-lemma-compare-higher-direct-image-proper} 부분 (1);
대신 우리는 이것으로부터 결과를 추론할 것이다.
$\epsilon_{Y, *}$는 기본 전층의 ID 함자이므로,
동형사상을 반영한다. 설명
보조정리 \ref{etale-cohomology-lemma-describe-pullback-pi-h}
$\epsilon_{Y, *} \circ a_Y^{-1} = \pi_Y^{-1}$을 보여줍니다.
$X$의 경우에도 마찬가지이다. 표준지도를 보여주기 위해
$a_Y^{-1}f_{small, *}\mathcal{F} \to f_{big, h, *}a_X^{-1}\mathcal{F}$
동형사상이면 다음을 보여주는 것으로 충분한다.
\begin{align*}
\pi_Y^{-1}f_{small, *}\mathcal{F}
& =
\epsilon_{Y, *}a_Y^{-1}f_{small, *}\mathcal{F} \\
& \to 
\epsilon_{Y, *}f_{big, h, *}a_X^{-1}\mathcal{F} \\
& =
f_{big, \etale, *} \epsilon_{X, *}a_X^{-1}\mathcal{F} \\
& =
f_{big, \etale, *}\pi_X^{-1}\mathcal{F}
\end{align*}
동형사상이다. 이것은 일부입니다
(1) 중 보조정리 \ref{etale-cohomology-lemma-compare-higher-direct-image-proper}이다.
\end{proof}

\begin{lemma}
\label{etale-cohomology-lemma-compare-h-etale}
사상 $\epsilon : (\Sch/S)_h \to (\Sch/S)_\etale$ 비교를 고려해보세요.
$\mathcal{P}$는 고유 사상의 클래스를 나타냅니다.
을 위한$X$~에$(\Sch/S)_\etale$나타내다
$\mathcal{A}'_X \subset \textit{Ab}((\Sch/X)_\etale)$
형식의 층으로 구성된 전체 하위 카테고리
$\pi_X^{-1}\mathcal{F}$ 여기서 $\mathcal{F}$는
비틀림 아벨 층 $X_\etale$
그런 다음 사이트, 성질의 코호몰로지
(\ref{sites-cohomology-item-base-change-P}),
(\ref{sites-cohomology-item-restriction-A}),
(\ref{sites-cohomology-item-A-sheaf}),
(\ref{sites-cohomology-item-A-and-P}) 및
(\ref{sites-cohomology-item-refine-tau-by-P})
사이트의 코호몰로지, 상황
\ref{sites-cohomology-situation-compare} 보류하세요.
\end{lemma}

\begin{proof}
먼저 $\mathcal{A}'_X \subset \textit{Ab}((\Sch/X)_\etale)$를 보여줍니다.
(1), (2), (3) 및 (4) 조건을 확인하여 약한 Serre 하위 범주이다.
호몰로지, 보조정리 \ref{homology-lemma-characterize-weak-serre-subcategory}.
부분 (1), (2), (3)는 $\pi_X^{-1}$가 정확하고
보조정리 \ref{etale-cohomology-lemma-cohomological-descent-etale}에 의해 예에 대해 완전히 충실한다. 만약에
$0 \to \pi_X^{-1}\mathcal{F} \to \mathcal{G} \to \pi_X^{-1}\mathcal{F}' \to 0$
$\textit{Ab}((\Sch/X)_\etale)$의 짧은 완전열이다.
그런 다음 $0 \to \mathcal{F} \to \pi_{X, *}\mathcal{G} \to \mathcal{F}' \to 0$
보조정리 \ref{etale-cohomology-lemma-cohomological-descent-etale}로 정확한다.
특히 우리는 $\pi_{X, *}\mathcal{G}$가 아벨리안임을 알 수 있다.
꼬임층 - $X_\etale$.
따라서 $\mathcal{G} = \pi_X^{-1}\pi_{X, *}\mathcal{G}$은
$\mathcal{A}'_X$ 최종 조건을 확인한다.

\medskip\noindent 코호몰로지 위의 사이트, 성질 (\ref{sites-cohomology-item-base-change-P})는 스킴의 올곱이 존재함으로써 유지된다. 고유 사상 의 스킴은 고유 사상 의 스킴이다. 사상, 보조정리 \ref{morphisms-lemma-base-change-proper}를 참조하라. 그리고 유한 표현의 사상의 기본 변경은 다음과 같습니다. 유한 표현의 사상, 사상, 보조정리 \ref{morphisms-lemma-base-change-finite-presentation}를 참조하라.

\medskip\noindent 코호몰로지 위의 사이트, 성질 (\ref{sites-cohomology-item-restriction-A})는 위상의 가환 도식 (3)에서 따릅니다. 보조정리 \ref{topologies-lemma-morphism-big-small-etale}.

\medskip\noindent 코호몰로지는 사이트, 성질 (\ref{sites-cohomology-item-A-sheaf})는 보조정리 \ref{etale-cohomology-lemma-describe-pullback-pi-h}이다.

\medskip\noindent
코호몰로지 사이트, 성질 (\ref{sites-cohomology-item-A-and-P})에 의해 유지된다.
보조정리 \ref{etale-cohomology-lemma-compare-higher-direct-image-proper} 부분 (2)
그리고 $R^if_{small}\mathcal{F}$
$\mathcal{F}$가 아벨 비틀림인 경우 비틀림이다.
층 $X_\etale$에 대한 내용은 보조정리 \ref{etale-cohomology-lemma-torsion-direct-image}를 참조하라.

\medskip\noindent 코호몰로지 위의 사이트, 성질 (\ref{sites-cohomology-item-refine-tau-by-P})는 사상, 보조정리에 대한 추가 정보에 의해 암시된다. \ref{more-morphisms-lemma-dominate-fppf-etale-locally} 전사 유한 국부적 자유 사상은 전사, 고유 및 유한 표현이라는 사실과 결합되어 평탄성에 관한 추가 사항, 보조정리 \ref{flat-lemma-surjective-proper-finite-presentation-h}에 의해 h-피복을 정의한다. \end{proof}

\begin{lemma}
\label{etale-cohomology-lemma-V-C-all-n-etale-h}
위와 같이 표기한다.
\begin{enumerate}
\item $X \in \Ob((\Sch/S)_{h})$ 및 아벨 꼬임층 $\mathcal{F}$
$X_\etale$에 우리는
$\epsilon_{X, *}a_X^{-1}\mathcal{F} = \pi_X^{-1}\mathcal{F}$
그리고$R^i\epsilon_{X, *}(a_X^{-1}\mathcal{F}) = 0$~을 위한$i > 0$.
\item 고유 사상 $f : X \to Y$의 $(\Sch/S)_h$
그리고 abelian 꼬임층 $\mathcal{F}$ 위의 $X$ 우리는
$a_Y^{-1}(R^if_{small, *}\mathcal{F}) =
R^if_{big, h, *}(a_X^{-1}\mathcal{F})$
모든 $i$에 대해.
\item 스킴 $X$와 대상 $K$ 중 꼬임 코호몰로지층을 갖는 $D^+(X_\etale)$의 것에 대하여,
코호몰로지 층 지도
$\pi_X^{-1}K \to R\epsilon_{X, *}(a_X^{-1}K)$은 동형사상이다.
\item 고유 사상 $f : X \to Y$의 스킴
그리고 $K$ 에서 $D^+(X_\etale)$ 비틀림 코호몰로지 층
$a_Y^{-1}(Rf_{small, *}K) = Rf_{big, h, *}(a_X^{-1}K)$.
\end{enumerate}
\end{lemma}

\begin{proof}
보조정리 \ref{etale-cohomology-lemma-compare-h-etale}에 의해 기본정리
코호몰로지 사이트, 절 \ref{sites-cohomology-section-compare-general}
모두 현재 설정에 적용된다. 결과를 번역하려면
범주 $\mathcal{A}_X$
코호몰로지 사이트, 보조정리 \ref{sites-cohomology-lemma-A}
$\textit{Ab}((\Sch/X)_h)$의 전체 하위 범주이다.
$a_X^{-1}\mathcal{F}$ 형식의 층으로 구성됨
여기서 $\mathcal{F}$는 $X_\etale$의 아벨 꼬임층이다.

\medskip\noindent
부분 (1)는 다음과 같습니다.$(V_n)$모두를 위해$n$이는 다음과 같이 유지된다.
코호몰로지 사이트, 보조정리 \ref{sites-cohomology-lemma-V-C-all-n-general}.

\medskip\noindent
부분(2)은 다음의 결론에 $\epsilon_Y^{-1}$을 적용하여 이어집니다.
코호몰로지 사이트, 보조정리 \ref{sites-cohomology-lemma-V-implies-C-general}.

\medskip\noindent 부분(3)은 사이트, 보조정리 \ref{sites-cohomology-lemma-V-C-all-n-general} 부분(1)의 코호몰로지에서 이어집니다. $\pi_X^{-1}K$는 $D^+_{\mathcal{A}'_X}((\Sch/X)_\etale)$ 및 $a_X^{-1} = \epsilon_X^{-1} \circ a_X^{-1}$에 있다.

\medskip\noindent 부분(4)은 같은 이유로 사이트, 보조정리 \ref{sites-cohomology-lemma-V-C-all-n-general} 부분(2)의 코호몰로지에서 이어집니다. \end{proof}

\begin{lemma} \label{etale-cohomology-lemma-cohomological-descent-etale-h} $X$를 스킴으로 둡니다. 비틀림 코호몰로지 층이 있는 $K \in D^+(X_\etale)$의 경우 맵 $$
K \longrightarrow Ra_{X, *}a_X^{-1}K
$$는 위와 같이 $a_X : \Sh((\Sch/X)_h) \to \Sh(X_\etale)$가 있는 동형사상이다. \end{lemma}

\begin{proof}
먼저 진술을 다음과 같은 경우로 축소한다.
$K$는 단일 아벨 층으로 제공된다. 즉, $K$를 나타낸다.
비틀림의 복합체 $\mathcal{F}^\bullet$ 아래 경계에 의해
아벨 층. 이는 사이트, 보조정리에서 코호몰로지로 가능한다.
\ref{sites-cohomology-lemma-torsion}. 의 경우
층 알겠습니다
$\mathcal{F}^n = a_{X, *} a_X^{-1} \mathcal{F}^n$
그리고 층 $R^qa_{X, *}a_X^{-1}\mathcal{F}^n$
$q > 0$에 대해서는 0이다. Leray의 비주기성 보조정리
(유도 범주, 보조정리 \ref{derived-lemma-leray-acyclicity})
$a_X^{-1}\mathcal{F}^\bullet$에 적용됨
그리고 함자 $a_{X, *}$ 결론을 내립니다. 이제부터 $K = \mathcal{F}$로 가정한다.
여기서 $\mathcal{F}$는 비틀림 아벨 층이다.

\medskip\noindent
보조정리 \ref{etale-cohomology-lemma-describe-pullback-pi-h}에 의해 우리는
$a_{X, *}a_X^{-1}\mathcal{F} = \mathcal{F}$. 따라서 다음을 보여주는 것으로 충분한다.
$R^qa_{X, *}a_X^{-1}\mathcal{F} = 0$~을 위한$q > 0$.
이를 위해 $a_X = \epsilon_X \circ \pi_X$를 사용할 수 있으며
르레이 스펙트럼 열
(코호몰로지 사이트, 보조정리 \ref{sites-cohomology-lemma-relative-Leray}).
보조정리 \ref{etale-cohomology-lemma-V-C-all-n-etale-h}에 의해
우리는$R^i\epsilon_{X, *}(a_X^{-1}\mathcal{F}) = 0$~을 위한$i > 0$
및 $\epsilon_{X, *}a_X^{-1}\mathcal{F} = \pi_X^{-1}\mathcal{F}$.
보조정리 \ref{etale-cohomology-lemma-cohomological-descent-etale}에 의해 우리는
$R^j\pi_{X, *}(\pi_X^{-1}\mathcal{F}) = 0$~을 위한$j > 0$.
이로써 증명이 종료된다.
\end{proof}

\begin{lemma} \label{etale-cohomology-lemma-compare-cohomology-etale-h} 위와 같이 스킴 $X$ 및 $a_X : \Sh((\Sch/X)_h) \to \Sh(X_\etale)$의 경우: \begin{enumerate} \item $H^q(X_\etale, \mathcal{F}) = H^q_h(X, a_X^{-1}\mathcal{F})$ 비틀림에 대한 아벨 층 $\mathcal{F}$ $X_\etale$, \item $H^q(X_\etale, K) = H^q_h(X, a_X^{-1}K)$에 대한 $K \in D^+(X_\etale)$ 비틀림 코호몰로지 층. \end{enumerate} 예: $A$가 비틀림 아벨 군이면 $H^q_\etale(X, \underline{A}) = H^q_h(X, \underline{A})$이다. \end{lemma}

\begin{proof}
이는 보조정리 \ref{etale-cohomology-lemma-cohomological-descent-etale-h}에서 따온 것이다.
사이트, 비고 \ref{sites-cohomology-remark-before-Leray}에서 코호몰로지에 의해.
\end{proof}











\section{내림차순 에탈 층} \label{etale-cohomology-section-glueing-etale}

\noindent 우리는 fppf와 h 위상에서 에탈 층 ``붙이기'' 및 관련 결과를 증명했다. 우리는 이미 다음과 같은 관련 결과를 보여주었습니다. \begin{enumerate} \item 보조정리 \ref{etale-cohomology-lemma-describe-pullback}는 스킴 $S$의 작은 에탈 사이트에 대한 층이 결정한다는 것을 알려줍니다. fpqc 피복에 대한 층 조건을 만족하는 $S$의 큰 에탈 사이트에 대한 층 (그리고 더욱이 Zariski, 에탈, 매끄러운, 신토믹 및 fppf 피복), \item 보조정리 \ref{etale-cohomology-lemma-describe-pullback-pi-fppf}는 fppf 위상, \item 보조정리 \ref{etale-cohomology-lemma-describe-pullback-pi-ph}에 대한 이전 요점을 다시 설명한 것이다. ph 위상, \item 보조정리 \ref{etale-cohomology-lemma-describe-pullback-pi-h}는 h 위상, \item 보조정리 \ref{etale-cohomology-lemma-descent-sheaf-fppf-etale}에 대해 동일한 것을 증명한다. 에탈 층, 그리고 \item 비고 \ref{etale-cohomology-remark-cohomological-descent-finite}는 유한 전사 사상에 대한 에탈 층의 하강을 가지고 있음을 알려줍니다(아래에서 이를 명확히 하고 강화하겠습니다). \end{enumerate} 단순 공간에 관한 장에서 이에 대한 몇 가지 추가 결과를 증명할 것이다. 예 단순 공간, 절 \ref{spaces-simplicial-section-fppf-hypercovering} 및 \ref{spaces-simplicial-section-proper-hypercovering-spaces}를 참조하라.

\medskip\noindent
결과를 편리하게 표현하려면 몇 가지 표기법이 필요한다.
$\mathcal{U} = \{f_i : X_i \to X\}$
대상이 고정된 스킴의 사상 계열이어야 한다.
에탈 층에 대한 {\it 하강 기준}
$\mathcal{U}$은(는) 가족입니다
$((\mathcal{F}_i)_{i \in I}, (\varphi_{ij})_{i, j \in I})$ 여기서
\begin{enumerate}
\item $\mathcal{F}_i$는 $\Sh(X_{i, \etale})$에 있고
\item $\varphi_{ij} :
\text{pr}_{0, small}^{-1} \mathcal{F}_i
\longrightarrow
\text{pr}_{1, small}^{-1} \mathcal{F}_j$
$\Sh((X_i \times_X X_j)_\etale)$의 동형사상이다.
\end{enumerate}
{\it 코사이클 조건}는 다음 다이어그램을 보유한다.
$$
\xymatrix{
\text{pr}_{0, small}^{-1}\mathcal{F}_i
\ar[dr]_{\text{pr}_{02, small}^{-1}\varphi_{ik}}
\ar[rr]^{\text{pr}_{01, small}^{-1}\varphi_{ij}} & &
\text{pr}_{1, small}^{-1}\mathcal{F}_j
\ar[dl]^{\text{pr}_{12, small}^{-1}\varphi_{jk}} \\
& \text{pr}_{2, small}^{-1}\mathcal{F}_k
}
$$
$\Sh((X_i \times_X X_j \times_X X_k)_\etale)$로 출퇴근하세요.
분명한 개념이 있다.{\it 사상하강 데이터}
그리고 우리는 하강 데이터의 범주를 얻습니다.
하강 기준점
$((\mathcal{F}_i)_{i \in I}, (\varphi_{ij})_{i, j \in I})$
존재하는 경우 {\it 유효}라고 한다.
$\mathcal{F}$~에$\Sh(X_\etale)$그리고동형사상
$\varphi_i : f_{i, small}^{-1} \mathcal{F} \to \mathcal{F}_i$
$\Sh(X_{i, \etale})$는 $\varphi_{ij}$와 호환된다. 즉,
그렇게
$$
\varphi_{ij} =
\text{pr}_{1, small}^{-1} (\varphi_j) \circ
\text{pr}_{0, small}^{-1} (\varphi_i^{-1})
$$
이것을 말하는 또 다른 방법은 다음과 같습니다. 주어진 객체 $\mathcal{F}$
$\Sh(X_\etale)$에서 {\it 정식 하강 데이터}를 얻습니다.
$(f_{i, small}^{-1}\mathcal{F}_i, c_{ij})$ 여기서 $c_{ij}$
정식 동형사상
$$
c_{ij} : \text{pr}_{0, small}^{-1} f_{i, small}^{-1}\mathcal{F}
\longrightarrow
\text{pr}_{1, small}^{-1} f_{j, small}^{-1}\mathcal{F}
$$
하강 데이텀
$((\mathcal{F}_i)_{i \in I}, (\varphi_{ij})_{i, j \in I})$
효과적이다. 필요충분조건은 이는 표준과 동형이다.
일부와 연관된 하강 데이텀$\mathcal{F}$~에$\Sh(X_\etale)$.

\medskip\noindent
패밀리가 단일 사상 $\{X \to Y\}$로 구성된 경우,
그런 다음 하강 기준점을 $(\mathcal{F}, \varphi)$ 쌍으로 생각한다.
여기서 $\mathcal{F}$는 $\Sh(X_\etale)$의 객체이고
$\varphi$는 동형사상이다.
$$
\text{pr}_{0, small}^{-1} \mathcal{F}
\longrightarrow
\text{pr}_{1, small}^{-1} \mathcal{F}
$$
$\Sh((X \times_Y X)_\etale)$에서 보조주기 조건은 다음을 유지한다.
$$
\xymatrix{
\text{pr}_{0, small}^{-1}\mathcal{F}
\ar[dr]_{\text{pr}_{02, small}^{-1}\varphi}
\ar[rr]^{\text{pr}_{01, small}^{-1}\varphi} & &
\text{pr}_{1, small}^{-1}\mathcal{F}
\ar[dl]^{\text{pr}_{12, small}^{-1}\varphi} \\
& \text{pr}_{2, small}^{-1}\mathcal{F}
}
$$
$\Sh((X \times_Y X \times_Y X)_\etale)$에서 통근한다.
하강 데이터와 유효성에 사상라는 개념이 있다.
이전과 똑같습니다.

\medskip\noindent 먼저 전사 정역 사상에 대한 유효 하강을 증명한다.

\begin{lemma}
\label{etale-cohomology-lemma-glue-etale-sheaf-section}
$f : X \to Y$는 절을 갖는 스킴의 사상라고 가정한다. 그런 다음
함자
$$
\Sh(Y_\etale)
\longrightarrow
\text{에탈 층의 강하 데이터(대상족: }\{X \to Y\}\text{)}
$$
배상$\mathcal{G}$~에$\Sh(Y_\etale)$정식 하강 데이텀으로
범주의 동치이다.
\end{lemma}

\begin{proof} 이것은 형식적이며 당김 함자의 기능에만 의존한다. 자세한 내용은 생략한다. 힌트: $s : Y \to X$가 절인 경우 준역전은 함자가 $(\mathcal{F}, \varphi)$를 $s_{small}^{-1}\mathcal{F}$로 보내는 것이다. \end{proof}

\begin{lemma} \label{etale-cohomology-lemma-glue-etale-sheaf-integral-surjective} $f : X \to Y$를 스킴의 전사 정역 사상으로 둡니다. 함자 $$
\Sh(Y_\etale)
\longrightarrow
\text{에탈 층의 강하 데이터(대상족: }\{X \to Y\}\text{)}
$$는 범주의 동치이다. \end{lemma}

\begin{proof}
이 증명에서는 당김에서 아래 첨자 ${}_{small}$를 삭제한다.
함자를 앞으로 밀어냅니다.
$X_1 = X \times_Y X$을 표시하고 $f_1 : X_1 \to Y$를 표시한다.
사상 $f \circ \text{pr}_0 = f \circ \text{pr}_1$.
$(\mathcal{F}, \varphi)$를 $\{X \to Y\}$의 하강 기준점으로 설정한다.
집합 $\mathcal{F}_1 = \text{pr}_0^{-1}\mathcal{F}$로 합시다.
$\varphi$을 동형사상을 정의하는 것으로 생각할 수 있다.
$\mathcal{F}_1 \to \text{pr}_1^{-1}\mathcal{F}$.
우리는 하강 데이텀을 보내는 규칙을 주장
$(\mathcal{F}, \varphi)$
층으로
$$
\mathcal{G} =
\text{이퀄라이저}\left(
\xymatrix{
f_*\mathcal{F}
\ar@<1ex>[r] \ar@<-1ex>[r] &
f_{1, *}\mathcal{F}_1
}
\right)
$$
은 보조정리의 서술문에 있는 함자에 준반대이다.
두 개의 화살표 중 첫 번째 화살표는 지도에서 나옵니다.
$$
f_*\mathcal{F} \to
f_*\text{pr}_{0, *}\text{pr}_0^{-1}\mathcal{F} =
f_{1, *}\mathcal{F}_1
$$
두 번째 화살표는 지도에서 나옵니다.
$$
f_*\mathcal{F} \to
f_* \text{pr}_{1, *}\text{pr}_1^{-1}\mathcal{F}
\xleftarrow{\varphi}
f_* \text{pr}_{0, *} \text{pr}_0^{-1}\mathcal{F} =
f_{1, *}\mathcal{F}_1
$$
여기서 왼쪽을 가리키는 화살표는 반전이 가능한다.
이 작업을 증명하려면 보여줘야 한다.
표준 맵 $f^{-1}\mathcal{G} \to \mathcal{F}$
동형사상이다. 자세한 내용은 생략했다. 이를 증명하기 위해서는
사상 컬렉션을 다시 가져온 후 확인하면 충분한다.
$\Spec(k) \to Y$ 여기서 $k$는 대수적으로 닫힌 체이다.
즉, 해당 기본 변경 사항 $X_k \to X$이 공동으로 적용된다.
전사적으로 $X_\etale$에 층의 지도가 있는지 확인할 수 있다.
기하점에서 줄기를 보면 동형사상이다.
정리 \ref{etale-cohomology-theorem-exactness-stalks}.
보조정리 \ref{etale-cohomology-lemma-integral-pushforward-commutes-with-base-change}에 의해
하강 기준점에서 $\mathcal{G}$ 구성
$(\mathcal{F}, \varphi)$ 베이스 변경 시 통근한다.
따라서 우리는 $Y$가 대수적으로
닫힌 체 (기본 변경은 성질을 유지한다.
사상 $f$의 내용을 참조하라.
사상, 보조정리 \ref{morphisms-lemma-base-change-surjective}
및 \ref{morphisms-lemma-base-change-finite}).
이 경우 사상 $X \to Y$에는 절이 있으므로
함자는 다음과 동등하다는 것을 알고 있다.
보조정리 \ref{etale-cohomology-lemma-glue-etale-sheaf-section}.
그러나 독자는 함자가 동등하다는 것을 보여줄 수 있다.
필요충분조건은 위의 구성은 준역(quasi-inverse)이다.
자세한 내용은 생략했다. 이로써 증명이 완료된다.
\end{proof}

\begin{lemma} \label{etale-cohomology-lemma-glue-etale-sheaf-proper-surjective} $f : X \to Y$를 스킴의 전사 고유 사상으로 둡니다. 함자 $$
\Sh(Y_\etale)
\longrightarrow
\text{에탈 층의 강하 데이터(대상족: }\{X \to Y\}\text{)}
$$는 범주의 동치이다. \end{lemma}

\begin{proof} 보조정리 \ref{etale-cohomology-lemma-glue-etale-sheaf-integral-surjective}에 주어진 것과 똑같은 증명이 작동하지만, 보조정리 \ref{etale-cohomology-lemma-integral-pushforward-commutes-with-base-change}에 대한 호소는 보조정리 \ref{etale-cohomology-lemma-proper-base-change-f-star}에 대한 호소로 대체되어야 한다. \end{proof}

\begin{lemma} \label{etale-cohomology-lemma-glue-etale-sheaf-check-after-base-change} $f : X \to Y$를 스킴의 사상으로 둡니다. $Z \to Y$를 스킴의 전사 정역 사상 또는 스킴의 전사 고유 사상으로 둡니다. 함자 $$
\Sh(Z_\etale)
\longrightarrow
\text{에탈 층의 강하 데이터(대상족: }\{X \times_Y Z \to Z\}\text{)}
$$ 및 $$
\Sh((Z \times_Y Z)_\etale)
\longrightarrow
\text{에탈 층의 강하 데이터(대상족: }
\{X \times_Y (Z \times_Y Z) \to Z \times_Y Z\}\text{)}
$$가 범주의 등가인 경우 $$
\Sh(Y_\etale)
\longrightarrow
\text{에탈 층의 강하 데이터(대상족: }\{X \to Y\}\text{)}
$$는 등가이다. \end{lemma}

\begin{proof} 정의와 기본정리 \ref{etale-cohomology-lemma-glue-etale-sheaf-integral-surjective} 및 \ref{etale-cohomology-lemma-glue-etale-sheaf-proper-surjective}의 형식적 결과이다. 자세한 내용은 생략했다. \end{proof}

\begin{lemma} \label{etale-cohomology-lemma-glue-etale-sheaf-fppf-cover} $f : X \to Y$를 스킴의 사상으로 가정한다. 이는 전사적이고 평면적이며 국소적으로 유한한 표현이다. 함자 $$
\Sh(Y_\etale)
\longrightarrow
\text{에탈 층의 강하 데이터(대상족: }\{X \to Y\}\text{)}
$$는 범주의 동치이다. \end{lemma}

\begin{proof}
증명 와 정확히 같습니다.
보조정리 \ref{etale-cohomology-lemma-glue-etale-sheaf-integral-surjective}
우리는 주장 준역전이 함자 전송에 의해 제공된다.
$(\mathcal{F}, \varphi)$ 에
$$
\mathcal{G} =
\text{이퀄라이저}\left(
\xymatrix{
f_*\mathcal{F}
\ar@<1ex>[r] \ar@<-1ex>[r] &
f_{1, *}\mathcal{F}_1
}
\right)
$$
그리고 이것을 증명하기 위해서는 다음을 보여주는 것으로 충분한다.
$f^{-1}\mathcal{G} \to \mathcal{F}$은 동형사상이다.
이를 로컬에서 확인할 수 있으므로 $Y$라고 가정할 수 있다.
친하다. 그러면 유한 전사 사상을 찾을 수 있다.
$Z \to Y$ 열려 있는 피복
$Z = \bigcup W_i$ $W_i \to Y$가 $X$를 통해 인수가 된다.
보다
사상, 보조정리 \ref{more-morphisms-lemma-dominate-fppf-finite}에 대해 자세히 알아보세요.
보조정리 적용 중
\ref{etale-cohomology-lemma-glue-etale-sheaf-check-after-base-change}
우리는 증명하는 것으로 충분하다고 봅니다
$Y$을 $Z$ 및 $Z \times_Y Z$ 및 $f$로 바꾼 후 보조정리
기본 변경으로. 따라서 우리는 $f$가 절 Zariski를 로컬로 가지고 있다고 가정할 수 있다.
물론, 문제가 $Y$에 국한되어 있다는 사실을 이용하여 우리는
절이 있는 경우
보조정리 \ref{etale-cohomology-lemma-glue-etale-sheaf-section}.
\end{proof}

\begin{lemma} \label{etale-cohomology-lemma-glue-etale-sheaf-fppf} $\{f_i : X_i \to X\}$를 스킴의 fppf 피복으로 설정한다. 함자 $$
\Sh(X_\etale)
\longrightarrow
\text{에탈 층의 강하 데이터(대상족: }\{f_i : X_i \to X\}\text{)}
$$는 범주의 동치이다. \end{lemma}

\begin{proof}
우리는 보조정리 \ref{etale-cohomology-lemma-glue-etale-sheaf-fppf-cover}
사상 $f : \coprod X_i \to X$에 대해.
그런 다음 공식적인 인수는 $f$에 대한 하강 데이터를 보여줍니다.
피복에 대한 하강 데이터와 동일한다. 비교
하강의 경우 보조정리 \ref{descent-lemma-family-is-one}.
자세한 내용은 생략했다.
\end{proof}

\begin{lemma} \label{etale-cohomology-lemma-glue-etale-sheaf-modification} $f : X' \to X$를 스킴의 고유 사상으로 둡니다. $i : Z \to X$를 닫힌 몰입으로 설정한다. 집합 $E = Z \times_X X'$. 그림 $$
\xymatrix{
E \ar[d]_g \ar[r]_j & X' \ar[d]^f \\
Z \ar[r]^i & X
}
$$ $f$가 $X \setminus Z$에 대한 동형사상인 경우 함자 $$
\Sh(X_\etale)
\longrightarrow
\Sh(X'_\etale) \times_{\Sh(E_\etale)} \Sh(Z_\etale)
$$는 범주의 동치이다. \end{lemma}

\begin{proof}
우리는 $2$-올곱 범주를 다음과 같이 작업할 것이다.
범주론, 예 \ref{categories-example-2-fibre-product-categories}.
함자는 $\mathcal{F}$을 트리플로 보냅니다.
$(f^{-1}\mathcal{F}, i^{-1}\mathcal{F}, c)$ 여기서
$c : j^{-1}f^{-1}\mathcal{F} \to g^{-1}i^{-1}\mathcal{F}$
표준 동형사상이다. 준역전 함자를 구성하겠습니다. 허락하다
$(\mathcal{F}', \mathcal{G}, \alpha)$ 객체가 됨
화살표 오른쪽 부분.
우리는 동형사상을 얻습니다.
$$
i^{-1}f_*\mathcal{F}' = g_*j^{-1}\mathcal{F}'
\xrightarrow{g_*\alpha}
g_*g^{-1}\mathcal{G}
$$
첫 번째 동등성은 보조정리 \ref{etale-cohomology-lemma-proper-base-change-f-star}이다.
이를 사용하여 지도를 얻습니다.
$i_*\mathcal{G} \to i_*g_*g^{-1}\mathcal{G}$
및 $f'_*\mathcal{F}' \to i_*g_*g^{-1}\mathcal{G}$. 우리는 집합
$$
\mathcal{F} = f_*\mathcal{F}' \times_{i_*g_*g^{-1}\mathcal{G}} i_*\mathcal{G}
$$
그리고 우리는 주장 $\mathcal{F}$가 왼쪽의 객체라는 것을
오른쪽의 이미지가 동형인 화살표의
우리가 시작한 트리플. $\mathcal{F}$의 줄기를 계산해 보자.
에기하점 $\overline{x}$~의$X$.

\medskip\noindent
$\overline{x}$이 아닌 경우
$Z$에서 한편으로는 $\overline{x}$는 고유한
기하점 $\overline{x}'$ $X'$ 및
$\mathcal{F}'_{\overline{x}'} = (f_*\mathcal{F}')_{\overline{x}}$
반면에 우리는 $(i_*\mathcal{G})_{\overline{x}}$
및 $(i_*g_*g^{-1}\mathcal{G})_{\overline{x}}$은 싱글톤이다.
따라서 $\mathcal{F}_{\overline{x}}$는 다음과 같습니다.
$\mathcal{F}'_{\overline{x}'}$.

\medskip\noindent $\overline{x}$가 $Z$에 있는 경우, 즉 $\overline{x}$가 기하점 $\overline{z}$의 $Z$ 이미지인 경우 다음을 얻습니다. $(i_*\mathcal{G})_{\overline{x}} = \mathcal{G}_{\overline{z}}$ 및 $$
(i_*g_*g^{-1}\mathcal{G})_{\overline{x}} =
(g_*g^{-1}\mathcal{G})_{\overline{z}} =
\Gamma(E_{\overline{z}}, g^{-1}\mathcal{G}|_{E_{\overline{z}}})
$$ (위에서 사용된 pushforward에 대한 고유 기저변환에 의해) 및 유사하게 $$
(f_*\mathcal{F}')_{\overline{x}} =
\Gamma(X'_{\overline{x}}, \mathcal{F}'|_{X'_{\overline{x}}})
$$ ID $E_{\overline{z}} = X'_{\overline{x}}$가 있고 $\alpha$가 동형사상을 정의하므로 층 $\mathcal{F}'|_{X'_{\overline{x}}}$ 및 $g^{-1}\mathcal{G}|_{E_{\overline{z}}}$ 이 경우에는 $$
\mathcal{F}_{\overline{x}} = \mathcal{G}_{\overline{z}}
$$를 얻는다고 결론을 내립니다.

\medskip\noindent
증명을 끝내기 위해 표준 맵이 있음을 관찰한다.
$i^{-1}\mathcal{F} \to \mathcal{G}$ 및 $f^{-1}\mathcal{F} \to \mathcal{F}'$
줄기에서 동형사상을 생성하는 $\alpha$과 호환된다.
우리는 위에서 보았습니다. 우리는 이 지도의 신중한 구성을 생략했다.
\end{proof}

\begin{lemma}
\label{etale-cohomology-lemma-h-descent-etale-sheaves}
$S$를 스킴으로 설정한다. 그런 다음 그룹형으로 섬유화된 범주
$$
p : \mathcal{S} \longrightarrow (\Sch/S)_h
$$
올 범주 위의 $U$는 범주 $\Sh(U_\etale)$이다.
의 층 위의 작은 에탈 사이트 의 $U$ 은 그룹형의 스택이다.
\end{lemma}

\begin{proof}
보조정리를 증명하기 위해 조건 (1), (2) 및 (3)를 확인한다.
평탄도에 대한 자세한 내용은 보조정리 \ref{flat-lemma-refine-check-h-stack}이다.

\medskip\noindent 조건 (1)는 층에 대한 접착이 있기 때문에 유지됩니다(그리고 Zariski 피복은 에탈 피복). 사이트, 보조정리 \ref{sites-lemma-glue-sheaves}를 참조하라.

\medskip\noindent
조건 (2)를 보려면 $f : X \to Y$가 전사라고 가정하고,
플랫, 고유 사상 $S$에 대한 유한 표현의 $Y$ 아핀.
그런 다음 다음 중 하나에 의해 $\{X \to Y\}$에 대한 하강이 있다.
보조정리 \ref{etale-cohomology-lemma-glue-etale-sheaf-fppf-cover} 또는
보조정리 \ref{etale-cohomology-lemma-glue-etale-sheaf-proper-surjective}.

\medskip\noindent 조건 (3)는 보다 일반적인 보조정리 \ref{etale-cohomology-lemma-glue-etale-sheaf-modification}에서 바로 이어집니다. \end{proof}



















\section{사각형을 날려버리고 에탈 코호몰로지} \label{etale-cohomology-section-blow-up-square}

\noindent 팽창 사각형은 평탄도에 대한 추가 정보, 절 \ref{flat-section-blow-up-ph}에 소개되어 있다. 고유 기저변환 정리를 사용하면 사각형 폭발에 대한 Mayer-Vietoris 유형의 결과를 볼 수 있다.

\begin{lemma}
\label{etale-cohomology-lemma-blow-up-square-cohomological-descent}
$X$를 스킴으로 두고 $Z \subset X$를 폐쇄형 하위 체계로 둡니다.
유한 유형의 준일관성 아이디얼로 잘라냅니다. 다음을 고려하라.
해당 폭발 사각형
$$
\xymatrix{
E \ar[d]_\pi \ar[r]_j & X' \ar[d]^b \\
Z \ar[r]^i & X
}
$$
비틀림이 있는 $K \in D^+(X_\etale)$의 경우 코호몰로지 층
우리에겐 뚜렷한 삼각형이 있다
$$
K \to Ri_*(K|_Z) \oplus Rb_*(K|_{X'}) \to Rc_*(K|_E) \to K[1]
$$
$D(X_\etale)$ 여기서 $c = i \circ \pi = b \circ j$이다.
\end{lemma}

\begin{proof}
$K|_{X'}$ 표기는 $b_{small}^{-1}K$을 의미한다.
아래 경계를 선택하세요. 복합체 $\mathcal{F}^\bullet$
$K$을 나타내는 아벨 층이다. 그것을 관찰하십시오
$i_*(\mathcal{F}^\bullet|_Z)$는 $Ri_*(K|_Z)$을 나타냅니다.
$i_*$이 정확하기 때문이다.
(명제 \ref{etale-cohomology-proposition-finite-higher-direct-image-zero}).
준동형사상을 선택하세요.
$b_{small}^{-1}\mathcal{F}^\bullet \to \mathcal{I}^\bullet$
여기서 $\mathcal{I}^\bullet$는 단사의 복합체 아래 경계이다.
아벨 층 - $X'_\etale$. 이 지도는 지도와 인접해 있다.
$\mathcal{F}^\bullet \to b_*(\mathcal{I}^\bullet)$ 그리고
$b_*(\mathcal{I}^\bullet)$는 $Rb_*(K|_{X'})$을 나타냅니다.
우리는 $\pi_*(\mathcal{I}^\bullet|_E) = (b_*\mathcal{I}^\bullet)|_Z$
보조정리 \ref{etale-cohomology-lemma-proper-base-change-f-star} 및
보조정리 \ref{etale-cohomology-lemma-proper-base-change} 이 복합체는 다음을 나타냅니다.
$R\pi_*(K|_E)$. 그러므로 지도
$$
Ri_*(K|_Z) \oplus Rb_*(K|_{X'}) \to Rc_*(K|_E)
$$
아래 경계 복합체의 전사 맵으로 표시된다.
$$
i_*(\mathcal{F}^\bullet|_Z) \oplus
b_*(\mathcal{I}^\bullet)
\to
i_*\left(b_*(\mathcal{I}^\bullet)|_Z\right)
$$
우리의 구별되는 삼각형을 얻으려면 다음을 보여주면 충분한다.
정식 지도
$\mathcal{F}^\bullet \to i_*(\mathcal{F}^\bullet|_Z) \oplus
b_*(\mathcal{I}^\bullet)$
지도의 핵에 준동형적으로 매핑된다.
위에 표시된 복합체(즉, 짧은 완전열
복합체의 파생된 삼각형을 결정합니다
범주, 참조
유도 범주, 절 \ref{derived-section-canonical-delta-functor}).
우리는 이것을 확인할 수 있습니다줄기에기하점 $\overline{x}$~의$X$.
$\overline{x}$가 $Z$에 없으면 $X' \to X$는 동형사상이다.
$\overline{x}$의 열린 동네 위에 있다. 따라서 $\overline{x}'$
이 경우에는 $X'$의 해당 기하점을 나타냅니다.
우리는 그것을 보여줘야 해
$$
\mathcal{F}^\bullet_{\overline{x}} \to \mathcal{I}^\bullet_{\overline{x}'}
$$
준동형사상이다. $\mathcal{I}^\bullet$을 선택하면 이는 사실이다.
$\overline{x}$가 $Z$에 있는 경우
$b_(\mathcal{I}^\bullet)_{\overline{x}} \to 
i_*\left(b_*(\mathcal{I}^\bullet)|_Z\right)_{\overline{x}}$
은 아벨 군의 복합체의 동형사상이다. 따라서
핵은 다음과 같습니다.
$i_*(\mathcal{F}^\bullet|_Z)_{\overline{x}} =
\mathcal{F}^\bullet_{\overline{x}}$ 원하는 대로.
\end{proof}

\begin{lemma} \label{etale-cohomology-lemma-blow-up-square-etale-cohomology} $X$는 스킴이고 $K \in D^+(X_\etale)$는 비틀림 코호몰로지 층을 갖습니다. $Z \subset X$를 유한 유형의 준일관성 아이디얼로 잘라낸 폐쇄형 하위 구성이라고 가정한다. 해당 폭발 정사각형 $$
\xymatrix{
E \ar[d] \ar[r] & X' \ar[d]^b \\
Z \ar[r] & X
}
$$을 고려하라. 그런 다음 정식 긴 완전열 $$
H^p_\etale(X, K) \to
H^p_\etale(X', K|_{X'}) \oplus
H^p_\etale(Z, K|_Z) \to
H^p_\etale(E, K|_E) \to
H^{p + 1}_\etale(X, K)
$$ \end{lemma}가 있다.

\begin{proof}[첫 번째 증명] 이는 보조정리 \ref{etale-cohomology-lemma-blow-up-square-cohomological-descent}와 $$
R\Gamma(X, Rb_*(K|_{X'})) = R\Gamma(X', K|_{X'})
$$에서 바로 이어집니다. (코호몰로지 위의 사이트, 절 참조) \ref{sites-cohomology-section-leray}) 다른 것들도 마찬가지이다. \end{proof}

\begin{proof}[두 번째 증명] 보조정리 \ref{etale-cohomology-lemma-compare-cohomology-etale-ph}에 의해 이 코호몰로지 군은 $X, X', E, Z$의 코호몰로지이며 일부 복합체의 값을 갖습니다. 아벨 층은 사이트 $(\Sch/X)_{ph}$에 있다. (즉, 유도 범주의 객체 $a_X^{-1}K$는 위의 보조정리 \ref{etale-cohomology-lemma-describe-pullback-pi-ph}를 참조하고 $K|_{X'} = b_{small}^{-1}K$를 상기하라.) 평탄성에 대한 자세한 내용은 보조정리 \ref{flat-lemma-blow-up-square-ph} ph 표현 가능한 전층의 다이어그램 중 층화는 코카르테시안이다. 따라서 보조정리는 사이트, 보조정리 \ref{sites-cohomology-lemma-square-triangle}의 사이트 $(\Sch/X)_{ph}$ 및 가환 도식에 적용되는 매우 일반적인 코호몰로지에서 따릅니다. 보조정리. \end{proof}

\begin{lemma}
\label{etale-cohomology-lemma-blow-up-square-equivalence}
$X$를 스킴으로 두고 $Z \subset X$를 폐쇄형 하위 체계로 둡니다.
유한 유형의 준일관성 아이디얼로 잘라냅니다. 다음을 고려하라.
해당 폭발 사각형
$$
\xymatrix{
E \ar[d]_\pi \ar[r]_j & X' \ar[d]^b \\
Z \ar[r]^i & X
}
$$
주어진다고 가정하자
\begin{enumerate}
\item 대상 $K'$ 중 꼬임 코호몰로지층을 갖는 $D^+(X'_\etale)$의 것,
\item 대상 $L$ 중 꼬임 코호몰로지층을 갖는 $D^+(Z_\etale)$의 것, 그리고
\item 동형사상 $\gamma : K'|_E \to L|_E$.
\end{enumerate}
그러면 객체가 존재합니다$K$~의$D^+(X_\etale)$
및 동형사상 $f : K|_{X'} \to K'$, $g : K|_Z \to L$ 등
그 $\gamma = g|_E \circ f^{-1}|_E$.
게다가 주어진
\begin{enumerate}
\item 대상 $M$ 중 꼬임 코호몰로지층을 갖는 $D^+(X_\etale)$의 것,
\item 사상 $\alpha : K' \to M|_{X'}$ $D(X'_\etale)$,
\item 사상 $\beta : L \to M|_Z$ $D(Z_\etale)$,
\end{enumerate}
그렇게
$$
\alpha|_E  = \beta|_E \circ \gamma.
$$
그런 다음 존재한다.사상 $M \to K$~에$D(X_\etale)$
$X'$에 대한 제한은 $a \circ f$이다.
$Z$에 대한 제한은 $b \circ g$이다.
\end{lemma}

\begin{proof}
$K$이 존재하는 경우 보조정리 \ref{etale-cohomology-lemma-blow-up-square-cohomological-descent}
그것이 들어맞는 뚜렷한 삼각형을 알려줍니다. 따라서 우리는 간단히 선택한다.
구별되는 삼각형
$$
K \to Ri_*(L) \oplus Rb_*(K') \to Rc_*(L|_E) \to K[1]
$$
여기서 $c = i \circ \pi = b \circ j$. 여기 지도가 있다. $Ri_*(L) \to Rc_*(L|_E)$
~이다$Ri_*$부속물 매핑에 적용됨$E \to R\pi_*(L|_E)$.
맵 $Rb_*(K') \to Rc_*(L|_E)$은 표준 맵의 구성이다.
$Rb_*(K') \to Rc_*(K'|_E)) = R$ 및 $Rc_*(\gamma)$.
보조정리 명령문의 $g$ 및 $f$ 맵은 인접 항목이다.
이 지도들 중. 이 구별된 삼각형을 $Z$로 제한하면
그러면 지도 $Rb_*(K) \to Rc_*(L|_E)$는 동형사상이 된다.
기본 변경 정리 (보조정리 \ref{etale-cohomology-lemma-proper-base-change})에 의해
지도 $g : K|_Z \to L$는 동형사상이다.
구별된 삼각형을 보면 $f : K|_{X'} \to K'$
$X' \setminus E = X \setminus Z$에 대한 동형사상이다.
게다가, 우리는 $\gamma \circ f|_E = g|_E$을 구축하고 있다.
그런 다음 $\gamma$ 및 $g$가 동형사상이므로 결론을 내립니다.
$f$는 $E$의 기하점에서 줄기에 동형사상을 유도한다.
또한. 따라서 $f$은 동형사상이다.

\medskip\noindent 마지막 문장의 경우 $K'$를 $K|_{X'}$로, $L$를 $K|_Z$로, $\gamma$을 정식 식별로 바꿀 수 있다. $\alpha$ 및 $\beta$는 교환 사각형 $$
\xymatrix{
K \ar[r] \ar@{..>}[d] &
Ri_*(K|_Z) \oplus Rb_*(K|_{X'}) \ar[r] \ar[d]_{\beta \oplus \alpha} &
Rc_*(K|_E) \ar[r] \ar[d]_{\alpha|_E} &
K[1] \ar@{..>}[d] \\
M \ar[r] &
Ri_*(M|_Z) \oplus Rb_*(M|_{X'}) \ar[r] &
Rc_*(M|_E) \ar[r] &
M[1]
}
$$을 유도한다. 따라서 유도 범주의 공리에 의해 구별되는 삼각형의 사상을 생성하는 점선 화살표를 얻습니다. \end{proof}














\section{거의 폭발할 뻔한 사각형과 h 위상} \label{etale-cohomology-section-blow-up-h}

\noindent
이 절에서 우리는 다음과 같은 논의를 계속한다.
평탄도에 대한 자세한 내용은 절 \ref{flat-section-blow-up-h}이다.
독자의 편의를 위해 우리는
거의 폭발할 뻔한 사각형은 가환 도식
\begin{equation}
\label{etale-cohomology-equation-almost-blow-up-square}
\vcenter{
\xymatrix{
E \ar[d] \ar[r] & X' \ar[d]^b \\
Z \ar[r] & X
}
}
\end{equation}
다음 조건을 만족하는 스킴의 경우:
\begin{enumerate}
\item $Z \to X$는 유한 표현의 닫힌 몰입이다.
\item $E = b^{-1}(Z)$는 $X'$의 로컬 주요 폐쇄 하위 체계이다.
\item $b$는 적절하고 유한한 표현이다.
\item 준일관된 방식으로 잘라낸 폐쇄형 하위 체계 $X'' \subset X'$
아이디얼의 절의 $\mathcal{O}_{X'}$는 $E$에서 지원된다.
(성질, 보조정리 \ref{properties-lemma-sections-supported-on-closed-subset})
폭발이다$X$~에$Z$.
\end{enumerate}
사상 $b$는 동형사상을 유도한다.
$X' \setminus E \to X \setminus Z$.

\medskip\noindent
우리는 "h-층 조건"에 대한 기준을 제시할 것이다.
큰 fppf 사이트의 유도 범주에 있는 개체
$(\Sch/S)_{fppf}$의 스킴 $S$. 동일한 기본 범주
두 번째 위상, 즉 h 위상이 있다.
(평탄도에 대한 자세한 내용, 절 \ref{flat-section-h}).
fppf 피복은 h 피복라는 점을 기억하세요.
(평탄도에 대한 자세한 내용, 보조정리 \ref{flat-lemma-zariski-h}). 그러므로 우리는
사상을 고려해보세요
$$
\epsilon : (\Sch/S)_h \longrightarrow (\Sch/S)_{fppf}
$$
사이트, 절 \ref{sites-cohomology-section-compare}의 코호몰로지를 참조하라.
특히, 우리는 완전히 충실한 함자
$$
R\epsilon_* : D((\Sch/S)_h) \longrightarrow D((\Sch/S)_{fppf})
$$
그리고 우리는 질문할 수 있습니다: 이 함자의 본질적인 이미지는 무엇입니까?

\begin{lemma} \label{etale-cohomology-lemma-blow-up-square-h} 위와 같이 표기하면 $K$가 $R\epsilon_*$의 필수 이미지에 있는 경우 코호몰로지의 $c^K_{X, Z, X', E}$ 맵은 사이트, 보조정리에 있다. \ref{sites-cohomology-lemma-c-square}는 준동형사상이다. \end{lemma}

\begin{proof} h 위상에 ${}^\#$ 층화를 나타냅니다. 우리는 평탄도에 대한 추가 정보, 보조정리 \ref{flat-lemma-blow-up-square-h}, 즉 $h_X^\# = h_Z^\# \amalg_{h_E^\#} h_{X'}^\#$를 보았습니다. 반면, 지도 $h_E^\# \to h_{X'}^\#$는 $E \to X'$가 단사 사상이므로 단사적이다. 따라서 이 보조정리는 사이트, 보조정리 \ref{sites-cohomology-lemma-descent-squares-helper}에 대한 코호몰로지의 특별한 경우입니다(이 자체는 사이트, 보조정리에 대한 코호몰로지의 공식적인 결과이다. \ref{sites-cohomology-lemma-square-triangle}). \end{proof}

\begin{proposition}
\label{etale-cohomology-proposition-check-h}
$K$를 $D^+((\Sch/S)_{fppf})$의 객체로 둡니다.
그런 다음 $K$는 다음의 필수 이미지에 있다.
$R\epsilon_* : D((\Sch/S)_h) \to D((\Sch/S)_{fppf})$
필요충분조건은 $c^K_{X, X', Z, E}$는 준동형사상이다.
거의 폭발할 뻔한 사각형마다 (\ref{etale-cohomology-equation-almost-blow-up-square})
$(\Sch/S)_h$에서 $X$ 아핀을 사용한다.
\end{proposition}

\begin{proof} 코호몰로지를 사이트, 보조정리 \ref{sites-cohomology-lemma-descent-squares}에 적용하여 이를 증명한다. 이 가설은 보조정리 \ref{etale-cohomology-lemma-blow-up-square-h} 및 평탄성에 대한 추가 정보 명제에 의해 유지된다. \ref{flat-proposition-check-h}. \end{proof}

\begin{lemma} \label{etale-cohomology-lemma-refine-check-h} $K$를 $D^+((\Sch/S)_{fppf})$의 객체로 둡니다. 그런 다음 $K$는 $R\epsilon_* : D((\Sch/S)_h) \to D((\Sch/S)_{fppf})$의 필수 이미지에 있다. 필요충분조건은 $c^K_{X, X', Z, E}$는 평탄도에 대한 추가 정보, 예 \ref{flat-example-one-generator} 및 \ref{flat-example-two-generators}에서와 같이 거의 폭발하는 모든 정사각형에 대해 준동형사상이다. \end{lemma}

\begin{proof} 코호몰로지를 사이트, 보조정리 \ref{sites-cohomology-lemma-descent-squares}에 적용하여 이를 증명한다. 이 가설은 보조정리 \ref{etale-cohomology-lemma-blow-up-square-h} 및 평탄성에 대한 추가 정보 보조정리에 의해 유지된다. \ref{flat-lemma-refine-check-h} \end{proof}






\section{코호몰로지 의 구조층 h 위상} \label{etale-cohomology-section-cohomology-O-h}

\noindent
$p$를 소수로 둡니다. $(\mathcal{C}, \mathcal{O})$을 사이트로 설정한다.
$p\mathcal{O} = 0$로. 그런 다음 집합 $\colim_F \mathcal{O}$
시스템의 환의 층의 범주에 있는 여극한과 같습니다.
$$
\mathcal{O} \xrightarrow{F} \mathcal{O} \xrightarrow{F}
\mathcal{O} \xrightarrow{F} \ldots
$$
여기서 $F : \mathcal{O} \to \mathcal{O}$, $f \mapsto f^p$
프로베니우스 내형성(자기준동형)이다.

\begin{lemma}
\label{etale-cohomology-lemma-h-sheaf-colim-F}
$p$를 소수로 둡니다. $S$를 $\mathbf{F}_p$에 대한 스킴으로 설정한다.
층 $\mathcal{O}^{perf} = \colim_F \mathcal{O}$를 고려하라.
$(\Sch/S)_{fppf}$에 있다. 그러면 $\mathcal{O}^{perf}$가 필수 항목이다.
$R\epsilon_* : D((\Sch/S)_h) \to D((\Sch/S)_{fppf})$의 이미지.
\end{lemma}

\begin{proof}
이를 보조정리 \ref{etale-cohomology-lemma-refine-check-h} 기준을 사용하여 증명한다.
조건을 확인하기 전에
준컴팩트하고 준분리된 물체 $X$
$(\Sch/S)_{fppf}$ 우리는
$$
H^i_{fppf}(X, \mathcal{O}^{perf}) = \colim_F H^i_{fppf}(X, \mathcal{O})
$$
사이트의 코호몰로지를 참조하라.
보조정리 \ref{sites-cohomology-lemma-colim-works-over-collection}.
$H^i_{fppf}(X, \mathcal{O}) = H^i(X, \mathcal{O})$도 사용할 것이다.
하강, 명제 \ref{descent-proposition-same-cohomology-quasi-coherent}.

\medskip\noindent $A, f, J$를 평탄도에 대한 추가 정보, 예 \ref{flat-example-one-generator}로 두고 관련된 거의 폭발하는 사각형을 고려한다. $X$, $X'$, $Z$, $E$는 유사하므로 $\mathcal{O}$ 중 더 높은 코호몰로지는 없다. 따라서 $$
0 \to
\mathcal{O}^{perf}(X) \to
\mathcal{O}^{perf}(X') \oplus \mathcal{O}^{perf}(Z) \to
\mathcal{O}^{perf}(E) \to 0
$$가 짧은 완전열인지 확인하면 된다. 이는 평탄도에 대한 추가 정보 보조정리 \ref{flat-lemma-h-sheaf-colim-F}의 (증명에 나와 있다.

\medskip\noindent
$X, X', Z, E$는 다음과 같습니다.
평탄도에 대한 자세한 내용은 예 \ref{flat-example-two-generators}이다.
$X$와 $Z$는 유사하므로 우리는
$H^p(X, \mathcal{O}_X) = H^p(Z, \mathcal{O}_X) = 0$~을 위한$p > 0$.
평탄도에 대한 추가 정보, 보조정리 \ref{flat-lemma-funny-blow-up}
우리는$H^p(X', \mathcal{O}_{X'}) = 0$~을 위한$p > 0$.
$E = \mathbf{P}^1_Z$와 $Z$는 유사하므로 우리는 또한
$H^p(E, \mathcal{O}_E) = 0$~을 위한$p > 0$. 이전과 마찬가지로
단락을 확인하는 것으로 줄이다.
$$
0 \to
\mathcal{O}^{perf}(X) \to
\mathcal{O}^{perf}(X') \oplus \mathcal{O}^{perf}(Z) \to
\mathcal{O}^{perf}(E) \to 0
$$
짧은 완전열이다. 이는 (증명 의)에 표시되어 있다.
평탄도에 대한 자세한 내용은 보조정리 \ref{flat-lemma-h-sheaf-colim-F}이다.
\end{proof}

\begin{proposition} \label{etale-cohomology-proposition-h-cohomology-structure-sheaf} $p$를 소수로 둡니다. $S$를 $\mathbf{F}_p$에 걸쳐 준간소하고 준분리된 스킴으로 설정한다. 그런 다음 $$
H^i((\Sch/S)_h, \mathcal{O}^h) =
\colim_F H^i(S, \mathcal{O})
$$ 여기서 $\mathcal{O}^h$의 왼쪽은 구조층의 h 층화를 의미한다. \end{proposition}

\begin{proof}
이것은 단지 보조정리 \ref{etale-cohomology-lemma-h-sheaf-colim-F}를 재구성한 것이다.
그것을 기억해
$\mathcal{O}^h = \mathcal{O}^{perf} = \colim_F \mathcal{O}$, 참조
평탄도에 대한 자세한 내용은 보조정리 \ref{flat-lemma-char-p}이다.
보조정리 \ref{etale-cohomology-lemma-h-sheaf-colim-F}에 의해 우리는
$\mathcal{O}^{perf}$는 $D((\Sch/S)_{fppf})$의 객체로 간주된다.
형태이다$R\epsilon_*K$어떤 사람들에게는$K \in D((\Sch/S)_h)$.
그러면 실제로는 $K = \epsilon^{-1}\mathcal{O}^{perf}$
같음$\mathcal{O}^{perf}$왜냐하면$\mathcal{O}^{perf}$h는 h야층. 보다
코호몰로지 사이트, 절 \ref{sites-cohomology-section-compare}.
그러므로 $R\epsilon_*\mathcal{O}^{perf} = \mathcal{O}^{perf}$
(혼란스러운 표기법에 대해 사과드립니다).
따라서 보조정리는 이제 Leray를 따릅니다.
$$
R\Gamma_h(S, \mathcal{O}^{perf}) =
R\Gamma_{fppf}(S, R\epsilon_*\mathcal{O}^{perf}) =
R\Gamma_{fppf}(S, \mathcal{O}^{perf})
$$
그리고 그 사실은
$$
H^i_{fppf}(S, \mathcal{O}^{perf}) =
H^i_{fppf}(S, \colim_F \mathcal{O}) =
\colim_F H^i_{fppf}(S, \mathcal{O})
$$
$S$은 준컴팩트하고 준분리되어 있기 때문이다.
(보조정리 \ref{etale-cohomology-lemma-h-sheaf-colim-F}의 증명 참조).
\end{proof}













% Bibliography is emitted once by the cumulative reader.
